\section{{\v C}ech computation of higher direct images $R^i f_*$ (unconditional)}
\label{mk-chap:Cohomology_CechHigherDirectImage}

\section{The {\v C}ech nerve and {\v C}ech complex of an affine cover}

\begin{definition}[{\v C}ech nerve of an affine open cover]\label{mk-def:cech_nerve}
  \lean{AlgebraicGeometry.CechNerve}
  \leanok
  \dcref{custom}

  \uses{mk-def:cover_cech_nerve, mk-def:push_pull_functor, mk-def:cech_nerve_cosimplicial, mk-lem:push_pull_id, mk-lem:push_pull_comp}
  \textit{Source: Stacks Project, Cohomology of Schemes, \S{\v C}ech cohomology of
  quasi-coherent sheaves.}
  Let \(X\) be a scheme and let
  \(\mathcal{U} : X = \bigcup_{i \in I} U_i\) be a finite open cover of \(X\) by
  affine opens, with index set \(I\). For a finite nonempty tuple
  \(i_0, \ldots, i_p \in I\) write
  \(U_{i_0 \ldots i_p} = U_{i_0} \cap \cdots \cap U_{i_p}\) for the corresponding
  intersection. The \emph{{\v C}ech nerve} of \(\mathcal{U}\) is the augmented
  \emph{cosimplicial} object in \(\operatorname{QCoh}(X)\) whose object in
  cosimplicial degree \(p\) is the direct image, along the open immersions
  \(j_{i_0 \ldots i_p} : U_{i_0 \ldots i_p} \hookrightarrow X\), of the restriction
  of \(\mathcal{O}_X\) (or of any chosen quasi-coherent sheaf) to the
  \((p+1)\)-fold intersections,
  \[
    [p] \;\longmapsto\;
      \prod_{(i_0, \ldots, i_p) \in I^{p+1}}
        (j_{i_0 \ldots i_p})_* \,\bigl(\,\cdot\,|_{U_{i_0 \ldots i_p}}\bigr),
  \]
  with cofaces given by the restriction maps that omit one index and
  codegeneracies by the maps that repeat one index; the augmentation in degree
  \(-1\) is the object on all of \(X\). The variance is essential: the geometric
  nerve of the cover (the iterated fibre powers of \(\coprod_i U_i\) over \(X\),
  introduced in \S\ref{mk-sec:cech_three_part} below) is a \emph{simplicial} object
  in \emph{schemes}, but applying the covariant direct-image functor turns it into
  a \emph{cosimplicial} object in \(\operatorname{QCoh}(X)\); consequently the
  associated alternating-coface complex is a \emph{cochain} complex (differentials
  raising degree), not a chain complex. When \(X\) is separated each intersection
  \(U_{i_0 \ldots i_p}\) is itself affine. The local affine model is the
  \emph{standard open covering} \(U = \bigcup_{i=1}^n D(f_i)\) of an affine open
  \(U = \operatorname{Spec}(A)\) by basic opens, where \(f_1, \ldots, f_n \in A\)
  generate the unit ideal; on such a cover the intersections are again basic opens
  \(D(f_{i_0} \cdots f_{i_p})\).
\end{definition}

\begin{definition}[{\v C}ech complex of a quasi-coherent sheaf]\label{mk-def:cech_complex}
  \lean{AlgebraicGeometry.CechComplex}
  \leanok
  \dcref{custom}

  \uses{mk-def:cech_nerve, mk-def:relative_cech_complex_of_nerve}
  \textit{Source: Stacks Project, Cohomology of Schemes,
  \texttt{lemma-cech-cohomology-quasi-coherent-trivial} (standard-cover {\v C}ech
  vanishing).}
  Let \(\mathcal{F}\) be a quasi-coherent \(\mathcal{O}_X\)-module and let
  \(\mathcal{U} : X = \bigcup_{i \in I} U_i\) be a finite affine open cover with
  all intersections \(U_{i_0 \ldots i_p}\) affine (e.g.\ \(X\) separated). The
  \emph{{\v C}ech complex} \(\check{\mathcal{C}}^\bullet(\mathcal{U}, \mathcal{F})\)
  is the cochain complex obtained by applying the global-sections functor to the
  {\v C}ech nerve of Definition~\ref{mk-def:cech_nerve} evaluated at \(\mathcal{F}\);
  in degree \(p\) it is
  \[
    \check{\mathcal{C}}^p(\mathcal{U}, \mathcal{F})
      = \prod_{(i_0, \ldots, i_p) \in I^{p+1}}
        \mathcal{F}\bigl(U_{i_0 \ldots i_p}\bigr),
  \]
  with differential \(d : \check{\mathcal{C}}^p \to \check{\mathcal{C}}^{p+1}\) the
  alternating sum \((d s)_{i_0 \ldots i_{p+1}} = \sum_{j=0}^{p+1} (-1)^j\,
  s_{i_0 \ldots \widehat{i_j} \ldots i_{p+1}}|_{U_{i_0 \ldots i_{p+1}}}\) of the
  restriction maps. Over an affine \(U = \operatorname{Spec}(A)\) with
  \(\mathcal{F}|_U = \widetilde{M}\) and a standard cover by the
  \(D(f_i)\), this is the complex
  \(\prod_{i_0} M_{f_{i_0}} \to \prod_{i_0 i_1} M_{f_{i_0} f_{i_1}} \to \cdots\) of
  localisations. The \emph{relative {\v C}ech complex} \(\check{\mathcal{C}}^\bullet
  (\mathfrak{U}, \mathcal{F})\) for \(f : X \to S\) is the analogous complex in
  \(\operatorname{QCoh}(S)\) obtained by applying \(f_*\) over each finite
  intersection instead of global sections; its terms are
  \(\prod_{i_0 \ldots i_p} (f|_{U_{i_0 \ldots i_p}})_*
  \bigl(\mathcal{F}|_{U_{i_0 \ldots i_p}}\bigr)\), and it is a complex in
  \(\operatorname{QCoh}(S)\) because each \(U_{i_0 \ldots i_p}\) is affine and the
  pushforward of a quasi-coherent sheaf along a quasi-compact quasi-separated
  morphism is quasi-coherent.
\end{definition}

\subsection{The three-part construction of the relative {\v C}ech complex}
\label{mk-sec:cech_three_part}

It is convenient to factor the construction of
Definitions~\ref{mk-def:cech_nerve}--\ref{mk-def:cech_complex} into three independent
pieces. Two of them are purely formal and unconditional; the sole non-trivial
mathematical content is isolated in the middle piece.

\paragraph{(1) Geometric backbone — the augmented {\v C}ech nerve of an arrow.}
Package the affine cover as a single morphism
\(\coprod_{i \in I} U_i \longrightarrow X\) (the canonical map out of the
coproduct induced by the inclusions \(U_i \hookrightarrow X\)). The augmented
{\v C}ech nerve of this arrow is an augmented \emph{simplicial scheme} over
\(X\): its object in simplicial degree \(p\) is the \((p+1)\)-fold fibre power
\[
  \Bigl(\coprod_{i} U_i\Bigr)^{\times_X (p+1)}
    \;=\;
    \coprod_{(i_0, \ldots, i_p) \in I^{p+1}}
      U_{i_0} \times_X \cdots \times_X U_{i_p},
\]
with the augmentation the cover map to \(X\). Because the \(U_i\) are open
subschemes of \(X\), each fibre product \(U_{i_0} \times_X \cdots \times_X U_{i_p}\)
is canonically the intersection \(U_{i_0 \ldots i_p}\). This backbone exists
\emph{unconditionally}: it uses only that the category of schemes has coproducts
and finite limits (hence the wide pullbacks underlying the {\v C}ech nerve of an
arrow). It carries no sheaf-theoretic data yet — it is the simplicial diagram of
spaces over which the construction is indexed.

\paragraph{(2) Push--pull functor — lifting the backbone to modules.}
The geometric backbone is turned into a cosimplicial \(\mathcal{O}_X\)-module by
the \emph{relative-direct-image functor on the over-category}
\[
  G : (X/\mathbf{Sch})^{\mathrm{op}} \longrightarrow \operatorname{QCoh}(X),
  \qquad
  (Y \xrightarrow{\,p\,} X) \;\longmapsto\; p_*\, p^*\mathcal{F},
\]
sending an \(X\)-scheme \(p : Y \to X\) to the pushforward along \(p\) of the
pullback \(p^*\mathcal{F}\). Composing \(G\) with the (opposite of the) simplicial
backbone of (1) produces exactly the augmented cosimplicial object of
Definition~\ref{mk-def:cech_nerve}: on the degree-\(p\) piece
\(p : U_{i_0 \ldots i_p} \hookrightarrow X\) one has
\(p_* p^*\mathcal{F} = (j_{i_0 \ldots i_p})_*(\mathcal{F}|_{U_{i_0 \ldots i_p}})\),
and a morphism of \(X\)-schemes \(h : (Y_1, p_1) \to (Y_2, p_2)\) (so
\(h \circ p_2 = p_1\)\,, reading \(p\) as the structure map) induces the
restriction comparison \(G(p_2) \to G(p_1)\) via the unit/counit of the
pull-push adjunction together with the over-triangle identifying the two
structure maps. This functor \(G\) is the \emph{one non-formal step} in the whole
construction. Its functoriality --- that \(G\) respects identities and composition
of morphisms over \(X\) --- is precisely a statement about the coherence of the
comparison isomorphisms \((p \circ q)_* \cong p_* q_*\) and
\((p \circ q)^* \cong q^* p^*\) for pushforward and pullback of sheaves of
modules, together with their compatibility with the adjunction units. Concretely,
pullback and pushforward of sheaves of modules are organised as a
pseudofunctor \(\mathrm{LocallyDiscrete}(\mathbf{Sch}^{\mathrm{op}}) \to
\mathbf{Adj}\,\mathbf{Cat}\) (sending a morphism to the adjoint pair
\(p^* \dashv p_*\)), and the comparison isomorphisms
\((p \circ q)^* \cong q^* p^*\), \((p \circ q)_* \cong p_* q_*\) are its structure
\(2\)-cells. The functoriality of \(G\) needs exactly three facts: the mate identity
expressing the conjugate of the inverse pullback comparison as the pushforward
comparison, its normalisation at the identity morphism, and the pseudofunctor
unitor and pentagon (associativity) coherences. The laws are nonetheless not formal
trivialities: each is a multi-step mate calculation that transports the comparison
isomorphisms along the over-triangle \(g \circ p_2 = p_1\) of structure maps and
past the adjunction unit, transporting along the over-triangle equality; the
obstacle is the bookkeeping of those transports, not a missing geometric ingredient.

\paragraph{(3) Coherence-free plumbing — from the nerve to the cochain complex.}
The passage from the augmented cosimplicial \(\mathcal{O}_X\)-module of (1)--(2)
to the relative {\v C}ech complex in \(\operatorname{QCoh}(S)\) is again purely
formal and uses \emph{no} pushforward/pullback coherence. One forgets the
augmentation, pushes the cosimplicial object forward along \(f : X \to S\) (a
covariant operation applied degreewise), and takes the alternating-coface-map
cochain complex of the resulting cosimplicial object of \(\mathcal{O}_S\)-modules.
This step relies only on the preadditivity of \(\operatorname{QCoh}(S)\): the
alternating sum \(\sum_{j} (-1)^j d^j\) of the cofaces is defined in any
preadditive category and squares to zero by the cosimplicial identities. Thus the
relative {\v C}ech complex \(\check{\mathcal{C}}^\bullet(\mathfrak{U}, \mathcal{F})\)
of Definition~\ref{mk-def:cech_complex} is genuinely \emph{defined} in terms of the
{\v C}ech nerve of Definition~\ref{mk-def:cech_nerve}: once the nerve is constructed,
the complex follows with no further input. In summary, of the three pieces only
the push--pull functor (2) requires content beyond formal category theory, and
that content is the pseudofunctor coherence of pushforward and pullback: the
comparison \(2\)-cells of the pseudofunctor
\(\mathrm{LocallyDiscrete}(\mathbf{Sch}^{\mathrm{op}}) \to \mathbf{Adj}\,\mathbf{Cat}\)
and their unitor/pentagon identities.

\paragraph{Formal bricks of the three-part construction.}
The geometric backbone (1), the push--pull functor (2), and the plumbing (3) are
recorded as the following declarations. The backbone and plumbing are
unconditional; the push--pull object and morphism maps are likewise unconditional
bricks, and the functor laws assembling them are the lone coherence statement.

\begin{definition}[Arrow of an affine cover]\label{mk-def:cover_arrow}
  \lean{AlgebraicGeometry.coverArrow}
  \leanok
  \dcref{custom}

  Let \(\mathcal{U} : X = \bigcup_{i \in I} U_i\) be an affine open cover. Its
  \emph{cover arrow} is the single morphism
  \(\coprod_{i \in I} U_i \longrightarrow X\) obtained as the universal map out of
  the coproduct induced by the open immersions \(U_i \hookrightarrow X\). Packaging
  the cover as one arrow lets the generic {\v C}ech-nerve-of-an-arrow machinery
  apply.
\end{definition}

\begin{definition}[Augmented {\v C}ech nerve of a cover]
  \label{mk-def:cover_cech_nerve}
  \lean{AlgebraicGeometry.coverCechNerve, AlgebraicGeometry.coverCechNerveOver, AlgebraicGeometry.coverCechNerveOverAug}
  \dcref{custom}

  \uses{mk-def:cover_arrow}
  The augmented {\v C}ech nerve of the cover arrow of
  Definition~\ref{mk-def:cover_arrow}: the augmented \emph{simplicial scheme} over
  \(X\) whose degree-\(p\) object is the \((p+1)\)-fold fibre power
  \[
    \Bigl(\coprod_{i} U_i\Bigr)^{\times_X (p+1)}
      = \coprod_{(i_0, \ldots, i_p) \in I^{p+1}}
        U_{i_0} \times_X \cdots \times_X U_{i_p},
  \]
  with augmentation the cover map to \(X\). It exists unconditionally because the
  category of schemes has all finite limits (hence the wide pullbacks underlying the
  {\v C}ech nerve of an arrow). This is the geometric backbone (1).
\end{definition}

\begin{definition}[Push--pull object map]\label{mk-def:push_pull_obj}
  \lean{AlgebraicGeometry.pushPullObj}
  \leanok
  \dcref{custom}

  The object map of the push--pull functor \(G\): to an \(X\)-scheme
  \(p : Y \to X\) it assigns
  \[
    G(Y, p) \;=\; p_*\, p^*\mathcal{F},
  \]
  the pushforward along \(p\) of the pullback \(p^*\mathcal{F}\).
\end{definition}

\begin{definition}[Push--pull comparison map]
  \label{mk-def:push_pull_map}
  \lean{AlgebraicGeometry.pushPullMap, AlgebraicGeometry.rawPushPullMap, AlgebraicGeometry.pushPullMap_eq_raw, AlgebraicGeometry.rawPushPullMap_self, AlgebraicGeometry.rawPushPullMap_self_gen, AlgebraicGeometry.pushforwardComp_hom_app_id, AlgebraicGeometry.pushPull_unit_comp}
  \dcref{custom}

  \uses{mk-def:push_pull_obj}
  The morphism map of \(G\). A morphism \(g : (Y_2, p_2) \to (Y_1, p_1)\) of
  \(X\)-schemes has underlying map \(\bar g : Y_2 \to Y_1\) satisfying the
  over-triangle \(p_1 \circ \bar g = p_2\); the contravariant functor produces the
  restriction comparison \(G(Y_1, p_1) \to G(Y_2, p_2)\) as the five-step composite
  \[
    p_{1*}\,p_1^*\mathcal{F}
      \xrightarrow{\,p_{1*}(\eta)\,}
    p_{1*}\,\bar g_*\,\bar g^*\,p_1^*\mathcal{F}
      \xrightarrow{\,\sim\,}
    (p_1 \circ \bar g)_*\,\bar g^*\,p_1^*\mathcal{F}
      \;=\;
    p_{2*}\,\bar g^*\,p_1^*\mathcal{F}
      \xrightarrow{\,p_{2*}(\sim)\,}
    p_{2*}\,p_2^*\mathcal{F},
  \]
  where \(\eta\) is the unit of the adjunction \(\bar g^* \dashv \bar g_*\), the
  second arrow is the pushforward comparison \(p_{1*}\bar g_* \cong (p_1 \bar g)_*\),
  the equality is the substitution \(p_2 = p_1 \circ \bar g\) of the over-triangle,
  and the last arrow applies \(p_{2*}\) to the pullback comparison
  \(\bar g^* p_1^* \cong (p_1 \bar g)^* = p_2^*\). The two over-triangle
  substitutions are realised as coercions along \(p_1 \circ \bar g = p_2\).
\end{definition}

\begin{lemma}[Push--pull functor — identity law]\label{mk-lem:push_pull_id}
  \lean{AlgebraicGeometry.pushPullMap_id}
  \leanok
  \dcref{custom}

  \uses{mk-def:push_pull_obj, mk-def:push_pull_map}
  The morphism map of Definition~\ref{mk-def:push_pull_map} respects identities:
  \(G(\mathrm{id}_Y) = \mathrm{id}_{G(Y)}\) for every \(X\)-scheme \(Y\).
\end{lemma}

\begin{proof}

  A direct calculation with comparison isomorphisms, with no sectionwise content.
  Rewrite the head \(p_*(\eta)\) followed by the pushforward comparison as the
  conjugate-mate of the inverse pullback comparison; the identity leg then collapses
  by the identity normalisation of the pseudofunctor together with right unitality,
  and the two over-triangle transport coercions cancel against the unitality coercion.
\end{proof}

\begin{lemma}
[Push--pull functor — composition law]
  \label{mk-lem:push_pull_comp}
  \lean{AlgebraicGeometry.pushPullMap_comp, AlgebraicGeometry.rawPushPullMap_comp, AlgebraicGeometry.pushPull_pentagon}
  \dcref{custom}

  \uses{mk-def:push_pull_obj, mk-def:push_pull_map, mk-lem:push_pull_unit_mate, mk-lem:push_pull_transport_cancel}
  The morphism map of Definition~\ref{mk-def:push_pull_map} is contravariantly functorial:
  for composable morphisms \(h : Y_3 \to Y_2\), \(g : Y_2 \to Y_1\) of \(X\)-schemes,
  \[
    G(g \circ h) = G(h) \circ G(g).
  \]
  Together with Lemma~\ref{mk-lem:push_pull_id} this assembles
  Definitions~\ref{mk-def:push_pull_obj}--\ref{mk-def:push_pull_map} into a functor
  \(G : (X/\mathbf{Sch})^{\mathrm{op}} \to \operatorname{QCoh}(X)\); composing \(G\) with
  the backbone of Definition~\ref{mk-def:cover_cech_nerve} yields the {\v C}ech nerve of
  Definition~\ref{mk-def:cech_nerve}.
\end{lemma}

\begin{proof}
  \uses{mk-lem:push_pull_unit_mate, mk-lem:push_pull_transport_cancel}
  The composition identity \(G(g \circ h) = G(h) \circ G(g)\) follows from three
  coherence facts. First, the composite adjunction unit \(\eta^{g \circ h}\) decomposes into
  the iterated units \(\eta^{g}\), \(\eta^{h}\) via the mate identity of
  Lemma~\ref{mk-lem:push_pull_unit_mate} and the naturality of the adjunction unit. Second,
  the resulting pullback-side identity is the \emph{pentagon coherence} of the pullback
  pseudofunctor — the associativity \(2\)-cell relating the two bracketings of the comparison
  isomorphism \((p \circ q \circ r)^* \cong r^* q^* p^*\). Third, the over-triangle transport
  coherences are absorbed by Lemma~\ref{mk-lem:push_pull_transport_cancel}.
\end{proof}

\begin{lemma}[Base-change unit (mate) identity for the push--pull head]\label{mk-lem:push_pull_unit_mate}
  \lean{AlgebraicGeometry.pushPull_unit_mate}
  \leanok
  \dcref{custom}

  \uses{mk-def:push_pull_map}
  For composable scheme morphisms \(f : A \to B\), \(p : B \to Z\) and a
  \(\mathcal{O}_Z\)-module \(N\), the adjunction unit \(\eta^p_N\) of
  \(p^* \dashv p_*\) at \(N\), followed by \(p_*\) of the unit \(\eta^f\) for
  \(f^* \dashv f_*\) at \(p^*N\), followed by the pushforward comparison
  \(p_* f_* \cong (f \circ p)_*\), equals the unit \(\eta^{f \circ p}_N\) for the
  composite followed by \((f \circ p)_*\) applied to the inverse pullback
  comparison \((f \circ p)^* \cong f^* p^*\):
  \[
    \eta^p_N \;\circ\; p_*(\eta^f_{p^*N}) \;\circ\;
      \bigl(p_* f_* \cong (f \circ p)_*\bigr)
    \;=\;
    \eta^{f \circ p}_N \;\circ\; (f \circ p)_*\bigl((f\circ p)^* \cong f^* p^*\bigr)^{-1}.
  \]
  This is the mate-calculus core that converts the single-morphism unit
  \(\eta^{f \circ p}\) into the iterated units \(\eta^p, \eta^f\); it is the
  reusable ingredient consumed by the functoriality (pentagon) law of the
  push--pull comparison map of Definition~\ref{mk-def:push_pull_map}.
\end{lemma}

\begin{proof}

  A direct mate-calculus computation: unfold both sides using the definition of
  the adjunction unit for the composite and the pseudofunctor structure maps, then
  verify the equality by the mate identity for the pushforward/pullback comparison.
\end{proof}

\begin{lemma}[Over-triangle transport cancellation for the push--pull tail]\label{mk-lem:push_pull_transport_cancel}
  \lean{AlgebraicGeometry.pushPull_transport_cancel}
  \leanok
  \dcref{custom}

  \uses{mk-def:push_pull_map}
  Let \(\bar g : Y_2 \to Y_1\), \(p_1 : Y_1 \to X\), \(p_2 : Y_2 \to X\) satisfy the
  over-triangle \(\bar g \circ p_1 = p_2\), and let \(\mathcal{F}\) be a
  \(\mathcal{O}_X\)-module. Then the pullback-comparison leg of the push--pull map,
  conjugated by the two \(\mathrm{eqToHom}\) coercions along the over-triangle,
  collapses to the transport-light form
  \[
    e_1 \;\circ\; p_{2*}\bigl((\bar g^* p_1^* \cong (\bar g p_1)^*)_{\mathcal{F}}\bigr)
      \;\circ\; e_2
    \;=\;
    (\bar g p_1)_*\bigl((\bar g^* p_1^* \cong (\bar g p_1)^*)_{\mathcal{F}}\bigr)
      \;\circ\; e_3,
  \]
  where \(e_1, e_2, e_3\) are the \(\mathrm{eqToHom}\) coercions induced by
  \(\bar g \circ p_1 = p_2\). It states the over-triangle cancellation in the
  push--pull comparison map of Definition~\ref{mk-def:push_pull_map} with the
  over-triangle equality as a free hypothesis, so a single substitution turns the
  transports into identities; it is the reusable pre-coherence brick for the
  composition (pentagon) law.
\end{lemma}

\begin{proof}

  Substitute the over-triangle equality \(\bar g \circ p_1 = p_2\) to replace
  transport coercions by identities, then verify the cancellation by direct
  calculation with the pushforward comparison.
\end{proof}

\begin{definition}[Push--pull functor]\label{mk-def:push_pull_functor}
  \lean{AlgebraicGeometry.pushPullFunctor}
  \leanok
  \dcref{custom}

  \uses{mk-def:push_pull_obj, mk-def:push_pull_map, mk-lem:push_pull_id, mk-lem:push_pull_comp}
  The \emph{push--pull functor}
  \[
    G : (X/\mathbf{Sch})^{\mathrm{op}} \longrightarrow \operatorname{QCoh}(X),
  \]
  assembling the object map \(p \mapsto p_*\, p^*\mathcal{F}\) of
  Definition~\ref{mk-def:push_pull_obj} and the comparison morphism map of
  Definition~\ref{mk-def:push_pull_map} into a genuine functor. Functoriality is exactly the two
  laws established above: the identity law
  Lemma~\ref{mk-lem:push_pull_id} and the composition law Lemma~\ref{mk-lem:push_pull_comp}. The
  functor is contravariant in the over-category \((X/\mathbf{Sch})\); it is the lift of the
  geometric backbone (1) to \(\mathcal{O}_X\)-modules described as step (2) above.
\end{definition}

\begin{definition}[{\v C}ech nerve cosimplicial object]\label{mk-def:cech_nerve_cosimplicial}
  \lean{AlgebraicGeometry.cechNerveCosimplicial}
  \leanok
  \dcref{custom}

  \uses{mk-def:cover_cech_nerve, mk-def:push_pull_functor}
  The (non-augmented) cosimplicial \(\mathcal{O}_X\)-module obtained by transporting the
  geometric {\v C}ech backbone in \((X/\mathbf{Sch})\) of
  Definition~\ref{mk-def:cover_cech_nerve} through the push--pull functor
  \(G\) of Definition~\ref{mk-def:push_pull_functor}: degreewise, the \(p\)-th term is the
  product over \((i_0, \ldots, i_p) \in I^{p+1}\) of the \((p+1)\)-fold-intersection
  pushforwards \((j_{i_0 \ldots i_p})_*(\mathcal{F}|_{U_{i_0 \ldots i_p}})\). This is the
  underlying cosimplicial object that the {\v C}ech nerve of
  Definition~\ref{mk-def:cech_nerve} augments by the structure map \(\mathcal{F} \to
  \mathcal{C}^0\).
\end{definition}

\begin{definition}[Relative {\v C}ech complex from a nerve]\label{mk-def:relative_cech_complex_of_nerve}
  \lean{AlgebraicGeometry.relativeCechComplexOfNerve}
  \leanok
  \dcref{custom}

  \uses{mk-def:cech_nerve}
  Given \(f : X \to S\) and an augmented cosimplicial \(\mathcal{O}_X\)-module
  \(N\), the \emph{relative {\v C}ech cochain complex} in \(\operatorname{QCoh}(S)\)
  obtained by forgetting the augmentation of \(N\), pushing the resulting
  cosimplicial object forward degreewise along \(f\), and taking the
  alternating-coface-map cochain complex. This passage uses only the preadditivity
  of \(\operatorname{QCoh}(S)\) and \emph{no} pushforward/pullback coherence; it is
  the coherence-free plumbing (3) that turns the nerve of
  Definition~\ref{mk-def:cech_nerve} into the {\v C}ech complex.
\end{definition}

\section{Affine acyclicity of the {\v C}ech complex}

\begin{definition}
[Standard affine cover from a spanning family]
  \label{mk-def:standard_affine_cover}
  \lean{AlgebraicGeometry.Scheme.affineOpenCoverOfSpanRangeEqTop}
  \mathlibok
  \dcref{custom}

  Let \(R\) be a commutative ring (so \(R = \Gamma(\operatorname{Spec} R,
  \mathcal{O})\)) and let \(s : \iota \to R\) be a family of elements that generate
  the unit ideal, \(\operatorname{Ideal.span}(\operatorname{range} s) = \top\). The
  \emph{standard affine open cover} of \(\operatorname{Spec} R\) associated with
  \(s\) is the affine open cover whose index set is \(\iota\), whose \(i\)-th member
  is the basic open \(D(s_i)\), and whose \(i\)-th structure morphism is the map of
  affine schemes induced by the localisation homomorphism \(R \to R_{s_i}\); that
  is, in Mathlib notation the defining equation
  \(\bigl(\text{affineOpenCoverOfSpanRangeEqTop}\bigr)_f\, i =
  \operatorname{Spec.map}\bigl(\text{algebraMap}\; R\; (\operatorname{Localization.Away}(s_i))\bigr)\),
  so that the sections over the \(i\)-th piece are the \emph{away} localisation
  \(R_{s_i} = \operatorname{Localization.Away}(s_i)\) and the geometric piece is
  \(D(s_i)\). This is exactly the standard open covering
  \(\operatorname{Spec} R = \bigcup_i D(s_i)\) by basic opens, packaged as a single
  cover object. It is the spanning-family bundle \((s, hs)\) that the affine
  {\v C}ech vanishing of Lemma~\ref{mk-lem:cech_acyclic_affine} is stated over.
\end{definition}

\subsection{The categorical--module bridge (L1)}
\label{mk-sec:cech_l1_bridge}

The affine vanishing of Lemma~\ref{mk-lem:cech_acyclic_affine} concerns the
\emph{section} {\v C}ech complex \(\check{\mathcal{C}}^\bullet(\mathcal{U},
\mathcal{F})\), a cochain complex of abelian groups, while its contracting homotopy is
most naturally described on a concrete complex of localised \(R\)-modules. The
passage between them has two parts: identify sections and restriction maps with away
localisations and their canonical maps, then translate vanishing of homology into
exactness of the resulting module complex. Definition~\ref{mk-def:qcoh_sections_localized}
and Lemma~\ref{mk-lem:section_cech_homology_exact} give these two parts.

\begin{definition}

[Quasi-coherent sections over a basic open are an away localisation]
  \label{mk-def:qcoh_sections_localized}
  \lean{AlgebraicGeometry.qcohSectionsAwayLocalized,
    AlgebraicGeometry.basicOpen_sprod,
    AlgebraicGeometry.iInf_fin_succ,
    AlgebraicGeometry.qcohRestriction_eq_comparison,
    AlgebraicGeometry.tP,
    AlgebraicGeometry.tU,
    AlgebraicGeometry.phiL,
    AlgebraicGeometry.phi,
    AlgebraicGeometry.phi_eq_phiL,
    AlgebraicGeometry.restr_bridge}
  \dcref{01HV,custom}

  \uses{mk-def:section_cech_complex, mk-def:standard_affine_cover}
  \textit{Source: Stacks Project, Schemes, \texttt{lemma-spec-sheaves} (Tag 01HV),
  items (4)--(5).}
  Let \(R\) be a commutative ring, \(\mathcal{F}\) a quasi-coherent
  \(\mathcal{O}_{\operatorname{Spec} R}\)-module, and write
  \(M = \Gamma(\operatorname{Spec} R, \mathcal{F})\) for its module of global
  sections. For every \(g \in R\), the restriction-of-sections map
  \(M = \mathcal{F}(\operatorname{Spec} R) \to \mathcal{F}(D(g))\) exhibits the
  section group \(\mathcal{F}(D(g))\) as the away localisation \(M_g\) of \(M\); that
  is, this map is an \(\operatorname{IsLocalizedModule}\) for the multiplicative
  submonoid \(\operatorname{Submonoid.powers} g\) of powers of \(g\). Moreover, for
  \(D(g_2) \subseteq D(g_1)\) the restriction map
  \(\mathcal{F}(D(g_1)) \to \mathcal{F}(D(g_2))\) is, under these identifications, the
  canonical localisation map \(M_{g_1} \to M_{g_2}\) (the unique \(R\)-linear map
  compatible with the localisation units). Item~(4) of Tag~01HV gives the first
  assertion for \(\mathcal{F} = \widetilde{M}\) (\(\Gamma(D(g), \widetilde{M}) = M_g\)
  as an \(R_g\)-module) and item~(5) gives the second (the restriction along
  \(D(g) \subset D(f)\) is the localisation map \(M_f \to M_g\)).

  For a standard sheaf \(\widetilde M\), the distinguished-open comparison identifies
  \(\Gamma(D(g),\widetilde M)\) with \(M_g\), and its naturality identifies restriction
  maps with localization maps. For arbitrary quasi-coherent \(\mathcal F\), the affine
  equivalence
  \(\mathcal{F} \cong \widetilde{\Gamma(\operatorname{Spec} R,\mathcal{F})}\)
  reduces the assertion to this standard-sheaf case.

  Two combinatorial--geometric ingredients package the identification into the form
  consumed by Lemma~\ref{mk-lem:section_cech_homology_exact}. First, the multi-index
  intersection identity
  \(\bigl(\bigcap_k D(s_{\sigma(k)})\bigr) = D\bigl(\prod_k s_{\sigma(k)}\bigr)\) in
  \((\operatorname{Spec} R).\mathrm{Opens}\) (\(\operatorname{basicOpen\_sprod}\), proved
  by induction on the index length using the infimum-splitting helper
  \(\operatorname{iInf\_fin\_succ}\) and \(\operatorname{PrimeSpectrum.basicOpen\_mul}\)):
  it rewrites the \((p+1)\)-fold intersection over a multi-index \(\sigma\) as the single
  basic open \(D(s_\sigma)\) of the product \(s_\sigma = \prod_k s_{\sigma(k)}\), so the
  section group there is the away localisation \(M_{s_\sigma}\). Second, the restriction
  compatibility of item~(5) is made explicit: for a basic-open inclusion
  \(D(b) \subseteq D(a)\) the presheaf restriction map of \(\widetilde{M}\) between the
  two section groups \emph{is} the away-localisation comparison
  \(\operatorname{AwayComparison.comparison}\) of the two localisation units
  (\(\operatorname{qcohRestriction\_eq\_comparison}\), by the uniqueness
  \(\operatorname{AwayComparison.comparison\_unique}\) of a map of localised modules);
  this is the single differential brick that the {\v C}ech coface match of
  Lemma~\ref{mk-lem:section_cech_homology_exact} feeds on.

  Consequently the canonical localization isomorphisms commute with every face
  restriction. They identify the section {\v C}ech differential with the alternating
  sum of the corresponding localization maps.
\end{definition}

\begin{lemma}
[Section {\v C}ech homology vanishes iff the localised complex is exact]
  \label{mk-lem:section_cech_homology_exact}
  \lean{AlgebraicGeometry.sectionCech_homology_exact,
    AlgebraicGeometry.AwayComparison.Inverts,
    AlgebraicGeometry.AwayComparison.Inverts.isUnit_powers,
    AlgebraicGeometry.AwayComparison.Inverts.mul,
    AlgebraicGeometry.AwayComparison.Inverts.of_dvd,
    AlgebraicGeometry.AwayComparison.comparison,
    AlgebraicGeometry.AwayComparison.comparison_apply,
    AlgebraicGeometry.AwayComparison.comparison_comp,
    AlgebraicGeometry.AwayComparison.comparison_comp_apply,
    AlgebraicGeometry.AwayComparison.comparison_comp_structure,
    AlgebraicGeometry.AwayComparison.comparison_self,
    AlgebraicGeometry.AwayComparison.comparison_unique,
    AlgebraicGeometry.CechLocalized.cechCoeff,
    AlgebraicGeometry.CechLocalized.cechCoeff_transport_eq_comparison,
    AlgebraicGeometry.CechLocalized.cechCoface,
    AlgebraicGeometry.CechLocalized.cechLocalized_exact,
    AlgebraicGeometry.CechLocalized.cechPrepend,
    AlgebraicGeometry.CechLocalized.cech_hcomm,
    AlgebraicGeometry.CechLocalized.cech_hsh,
    AlgebraicGeometry.CechLocalized.cech_hu,
    AlgebraicGeometry.CechLocalized.sprod,
    AlgebraicGeometry.CechLocalized.sprod_cons,
    AlgebraicGeometry.CechLocalized.sprod_succAbove_dvd}
  \dcref{custom}

  \uses{mk-def:section_cech_complex, mk-def:qcoh_sections_localized, mk-lem:section_cech_module_exact,
    mk-lem:section_cech_product_equiv, mk-lem:section_cech_objd_apply,
    mk-lem:section_cech_coface_match, mk-lem:section_cech_ab_exact}
  \textit{Source: Stacks Project, Cohomology of Schemes,
  \texttt{lemma-cech-cohomology-quasi-coherent-trivial} (identification with the
  complex of localisations).}
  Let \(\mathcal{F}\) be a quasi-coherent \(\mathcal{O}_{\operatorname{Spec}
  R}\)-module and \(\mathcal{U}\) the standard cover associated with a spanning family
  \(s : \iota \to R\). Under the section identification of
  Definition~\ref{mk-def:qcoh_sections_localized}, the section {\v C}ech complex
  \(\check{\mathcal{C}}^\bullet(\mathcal{U}, \mathcal{F})\) of
  Definition~\ref{mk-def:section_cech_complex} is isomorphic, as a cochain complex, to
  the concrete complex of localised \(R\)-modules
  \[
    D^\bullet : \quad
    \prod_{\sigma : \{0\} \to \iota} M_{s_\sigma} \;\longrightarrow\;
    \prod_{\sigma : \{0,1\} \to \iota} M_{s_\sigma} \;\longrightarrow\;
    \cdots,
    \qquad
    s_\sigma = \prod_{k} s_{\sigma(k)},
  \]
  whose differential is the alternating sum of the canonical localisation maps. Its
  positive-degree homology vanishes:
  \[
    1 \leq p
    \;\Longrightarrow\;
    \operatorname{IsZero}\bigl(\check{\mathcal{C}}^\bullet(\mathcal{U},
      \mathcal{F}).\mathrm{homology}\, p\bigr),
  \]
  This follows from positive-degree exactness of the underlying group
  homomorphisms,
  \(\operatorname{Function.Exact}\bigl(D^{p-1} \to D^p \to D^{p+1}\bigr)\), which is
  Lemma~\ref{mk-lem:section_cech_module_exact}.

  The degree-\(p\) object of
  \(\check{\mathcal{C}}^\bullet(\mathcal{U}, \mathcal{F})\) is the \emph{categorical
  product}
  \(\prod^{c}_{\sigma : \mathrm{Fin}(p+1) \to \iota}
  \mathcal{F}(D(s_\sigma))\) in the category \(\mathrm{Ab}\) of abelian groups (the
  alternating-coface-map complex of the \(\mathrm{Ab}\)-valued cosimplicial object
  \(\operatorname{sectionCechCosimplicial}\)), not an \(R\)-module object, so the
  appropriate exactness criterion is therefore the one in \(\mathrm{Ab}\). Vanishing
  of homology is equivalent to exactness at degree \(p\), which in \(\mathrm{Ab}\) is
  the underlying-group condition
  \(\forall x_2,\; g(x_2)=0\Rightarrow\exists x_1,\;f(x_1)=x_2\)
  (Lemma~\ref{mk-lem:section_cech_ab_exact}). The categorical product is identified with
  per-\(\sigma\) localised modules by the element-level product equivalence
  \(\operatorname{ToType}(\prod^{c} F_\sigma) \simeq \prod_\sigma
  \operatorname{ToType}(F_\sigma)\)
  (Lemma~\ref{mk-lem:section_cech_product_equiv}), under which the {\v C}ech coface
  differential matches the localisation differential \(\operatorname{dDiff}\) of
  \(D^\bullet\) (the abstract cosimplicial unfold
  Lemma~\ref{mk-lem:section_cech_objd_apply} followed by the tilde-bridge identification
  Lemma~\ref{mk-lem:section_cech_coface_match}). Transporting the exactness of
  Lemma~\ref{mk-lem:section_cech_module_exact} across the product equivalence then
  proves the group-level exactness.
\end{lemma}

\begin{proof}
  \uses{mk-def:qcoh_sections_localized, mk-def:section_cech_complex}
  The term identification is the degreewise statement of
  Definition~\ref{mk-def:qcoh_sections_localized}: the degree-\(p\) section group
  \(\prod_\sigma \mathcal{F}(D(s_\sigma))\) is \(\prod_\sigma M_{s_\sigma}\), since
  each \((p+1)\)-fold intersection \(D(s_{\sigma(0)}) \cap \cdots \cap
  D(s_{\sigma(p)})\) is the basic open \(D(s_\sigma)\) of the product
  \(s_\sigma = \prod_k s_{\sigma(k)}\) and the sections there are the away
  localisation \(M_{s_\sigma}\). The differential identification is item~(5) of
  Definition~\ref{mk-def:qcoh_sections_localized}: each face restriction
  \(\mathcal{F}(D(s_{\sigma \circ d_j})) \to \mathcal{F}(D(s_\sigma))\) is the
  canonical localisation map \(M_{s_{\sigma \circ d_j}} \to M_{s_\sigma}\), so the
  alternating sum of restrictions is the alternating sum of localisation maps, the
  differential of \(D^\bullet\) (the coface match of
  Lemma~\ref{mk-lem:section_cech_coface_match}). The homology vanishing then follows in
  three steps. First, \(\operatorname{exactAt\_iff\_isZero\_homology}\) turns
  \(\operatorname{IsZero}\) of the homology at \(p\) into \(\operatorname{ExactAt}\) of
  the short complex at \(p\). Second, because the degree-\(p\) objects are
  \emph{categorical products in \(\mathrm{Ab}\)} rather than \(R\)-modules, the abelian
  analogue \(\operatorname{ShortComplex.ab\_exact\_iff}\)
  (Lemma~\ref{mk-lem:section_cech_ab_exact}) --- not the module-category criterion ---
  reduces \(\operatorname{ExactAt}\) to the underlying-group exactness
  \(\forall x_2,\; g(x_2) = 0 \Rightarrow \exists x_1,\; f(x_1) = x_2\). Third, the
  element-level product equivalence
  \(\operatorname{ToType}(\prod^{c} F_\sigma) \simeq \prod_\sigma
  \operatorname{ToType}(F_\sigma)\) of Lemma~\ref{mk-lem:section_cech_product_equiv}
  identifies each product group with the per-\(\sigma\) localised module \(D^p\), so the
  group-level exactness is exactly \(\operatorname{Function.Exact}(D^{p-1} \to D^p \to
  D^{p+1})\), supplied by \(\operatorname{dDiff\_exact}\)
  (Lemma~\ref{mk-lem:section_cech_module_exact}).

  This argument uses the equivalence between localized modules and its essential
  image: global sections is an inverse there and hence reflects zero homology. It does
  not require an a priori assertion that global sections preserves homology.
\end{proof}

\begin{lemma}

[Element-level product equivalence for the section {\v C}ech terms]
  \label{mk-lem:section_cech_product_equiv}
  \lean{AlgebraicGeometry.sectionCechProductEquiv,
    AlgebraicGeometry.sectionCechProductEquiv_apply,
    AlgebraicGeometry.sectionProdAddEquiv}
  \dcref{custom}

  \uses{mk-def:section_cech_complex, mk-def:qcoh_sections_localized}
  Let \(\{F_\sigma\}_{\sigma : \mathrm{Fin}(p+1) \to \iota}\) be the finite family of
  abelian groups \(F_\sigma = \mathcal{F}(D(s_\sigma))\) appearing in degree \(p\) of
  the section {\v C}ech complex of Definition~\ref{mk-def:section_cech_complex}. There is
  an equivalence of underlying types
  \[
    \operatorname{ToType}\Bigl(\prod^{c}_{\sigma} F_\sigma\Bigr)
    \;\simeq\;
    \prod_{\sigma} \operatorname{ToType}(F_\sigma),
  \]
  between the elements of the categorical product \(\prod^{c}_{\sigma} F_\sigma\) in
  \(\mathrm{Ab}\) and the dependent functions \(\sigma \mapsto\) (element of
  \(F_\sigma\)). It is supplied by the Mathlib concrete-limit comparison
  \(\operatorname{CategoryTheory.Limits.Concrete.productEquiv}\), which holds because the
  forgetful functor \(\operatorname{forget}\,\mathrm{Ab}\) preserves the discrete
  product (\(\operatorname{PreservesLimit}(\operatorname{Discrete.functor} F)\)). Under
  this equivalence each coordinate
  \(\operatorname{ToType}(\mathcal{F}(D(s_\sigma)))\) is identified with the away
  localisation \(M_{s_\sigma} = \operatorname{dCoeff}\) of
  Definition~\ref{mk-def:section_cech_module_complex} by the
  \(\operatorname{IsLocalizedModule}\) uniqueness of
  Definition~\ref{mk-def:qcoh_sections_localized} (both
  \(\mathcal{F}(D(s_\sigma))\) and \(M_{s_\sigma}\) localise \(M\) at \(s_\sigma\), so
  the comparison is the canonical isomorphism). This is the move between the
  categorical product of \(\mathrm{Ab}\) and the per-\(\sigma\) localised modules.
  Its packaging as an additive isomorphism --- the degreewise vertical map of the
  ladder transport in Lemma~\ref{mk-lem:section_cech_ab_exact} --- is recorded as
  \(\operatorname{sectionProdAddEquiv}\), the \(\operatorname{AddEquiv}\) underlying
  the coordinatewise comparison.
  The coordinate projection of the equivalence is recorded separately as the
  definitional restatement
  \(\operatorname{sectionCechProductEquiv\_apply}\): the \(\sigma\)-component of the
  image of \(y\) under the equivalence is the underlying group map of the categorical
  projection \(\pi_\sigma\) applied to \(y\).
\end{lemma}

\begin{proof}
  \uses{mk-def:section_cech_complex, mk-def:qcoh_sections_localized}
  The type-level equivalence is \(\operatorname{Concrete.productEquiv}\) applied to the
  family \(\sigma \mapsto F_\sigma\); its hypothesis is that
  \(\operatorname{forget}\,\mathrm{Ab}\) preserves the discrete-diagram limit, which is
  the standard instance for \(\mathrm{Ab}\). The per-coordinate identification
  \(\operatorname{ToType}(\mathcal{F}(D(s_\sigma))) \cong M_{s_\sigma}\) is the
  uniqueness clause of \(\operatorname{IsLocalizedModule}\): by
  Definition~\ref{mk-def:qcoh_sections_localized} the restriction map
  \(M \to \mathcal{F}(D(s_\sigma))\) and the localisation unit
  \(M \to M_{s_\sigma}\) are both localisations of \(M\) at
  \(\operatorname{Submonoid.powers}(s_\sigma)\)
  (\(\operatorname{qcohSectionsAwayLocalized}\)), hence canonically isomorphic over
  \(M\). Assembling the coordinatewise isomorphisms after the product equivalence gives
  the stated identification.
\end{proof}

\begin{lemma}
[Abstract cosimplicial unfold of the section {\v C}ech differential]
  \label{mk-lem:section_cech_objd_apply}
  \lean{AlgebraicGeometry.sectionCech_objD_apply,
    AlgebraicGeometry.sectionCechFaceRestr}
  \dcref{custom}

  \uses{mk-def:section_cech_complex, mk-lem:section_cech_product_equiv}
  Let \(\mathcal{F}\) be a presheaf of \(\mathcal{O}_X\)-modules and
  \(\mathcal{U} : X = \bigcup_{i \in I} U_i\) an open cover. For a multi-index
  \(\sigma : \mathrm{Fin}(q+2) \to I\) and a face index
  \(i \in \mathrm{Fin}(q+2)\), write
  \[
    \operatorname{sectionCechFaceRestr}(\sigma, i) :
      \mathcal{F}\Bigl(\,\bigcap_{l} U_{(\sigma \circ d_i)(l)}\Bigr)
      \;\longrightarrow\;
      \mathcal{F}\Bigl(\,\bigcap_{k} U_{\sigma(k)}\Bigr)
  \]
  for the \(i\)-th presheaf face restriction --- the value of \(\mathcal{F}\) on the
  inclusion of opens
  \(\bigcap_{k} U_{\sigma(k)} \subseteq \bigcap_{l} U_{(\sigma \circ d_i)(l)}\)
  obtained from the order-preserving coface \(d_i\) of the simplex category. Then,
  read through the product equivalence
  \(\operatorname{sectionCechProductEquiv}\) of
  Lemma~\ref{mk-lem:section_cech_product_equiv}, the underlying group action of the
  degree-\(q\) differential
  \(\operatorname{objD} q\) of the alternating-coface-map complex
  \(\operatorname{AlgebraicTopology.alternatingCofaceMapComplex}
  (\operatorname{sectionCechCosimplicial}\,\mathcal{U}\,\mathcal{F})\) is the alternating
  sum of the face restrictions applied to the deleted-face coordinates:
  \[
    \bigl(\operatorname{sectionCechProductEquiv}_{q+1}\,(\operatorname{objD} q\, t)\bigr)_\sigma
      \;=\;
      \sum_{i = 0}^{q+1} (-1)^i \,
        \operatorname{sectionCechFaceRestr}(\sigma, i)
          \Bigl(\bigl(\operatorname{sectionCechProductEquiv}_{q}\, t\bigr)_{\sigma \circ d_i}\Bigr).
  \]
  This is \emph{purely cosimplicial}: it records the shape of the section {\v C}ech
  differential as an alternating sum of presheaf restrictions, with \emph{no}
  sheaf-theoretic or localisation identification of those restrictions. It is the
  abstract shape that the localisation differential \(\operatorname{dDiff}\) of
  Definition~\ref{mk-def:section_cech_module_complex} matches once the tilde-bridge
  identification of Lemma~\ref{mk-lem:section_cech_coface_match} is supplied.
\end{lemma}

\begin{proof}
  \uses{mk-lem:section_cech_product_equiv}
  Unfold the degree-\(q\) map of \(\operatorname{alternatingCofaceMapComplex}\) on
  \(\operatorname{sectionCechCosimplicial}\): by construction \(\operatorname{objD} q\)
  is the alternating sum \(\sum_{i} (-1)^i\, d^i\) of the cosimplicial cofaces, each
  coface \(d^i\) being the \(\operatorname{Pi.lift}\) whose \(\sigma\)-component
  projects to the coordinate \(\sigma \circ d_i\) and post-composes with the presheaf
  restriction \(\operatorname{sectionCechFaceRestr}(\sigma, i)\). Reading the
  \(\sigma\)-coordinate through the product equivalence (the coordinate-projection
  restatement \(\operatorname{sectionCechProductEquiv\_apply}\)) commutes the projection
  past the categorical sum: applying the sum of \(\mathrm{Ab}\)-morphisms to an element
  distributes over the sum (the finite-sum-application identity
  \(\operatorname{ab\_hom\_finsetSum\_apply}\)), and each summand becomes
  \((-1)^i\) times the face restriction of the \((\sigma \circ d_i)\)-coordinate. The
  per-coordinate projection of the cosimplicial coface is the \(\operatorname{Pi.lift}\)
  computation \(\operatorname{Pi.lift\_π}\), which identifies the projected coface with
  \(\operatorname{sectionCechFaceRestr}(\sigma, i)\) precomposed with the projection to
  \(\sigma \circ d_i\). Assembling the signs yields the stated alternating sum.
\end{proof}

\begin{lemma}

[Coface match: the section {\v C}ech differential is the localisation differential]
  \label{mk-lem:section_cech_coface_match}
  \lean{AlgebraicGeometry.sectionCechCofaceMatch,
    AlgebraicGeometry.phiL_naturality,
    AlgebraicGeometry.phi_naturality}
  \dcref{custom}

  \uses{mk-lem:section_cech_objd_apply, mk-def:qcoh_sections_localized, mk-def:section_cech_module_complex}
  Take \(X = \operatorname{Spec} R\), the standard affine cover \(\mathcal{U}\)
  associated with a spanning family \(s : \iota \to R\), and \(\mathcal{F} =
  \widetilde{M}\) the quasi-coherent sheaf associated with an \(R\)-module \(M\)
  (for a general quasi-coherent \(\mathcal{F}\), see the note following the statement). In
  the situation of Lemma~\ref{mk-lem:section_cech_product_equiv}, the degree-\(p\)
  differential of the section {\v C}ech complex
  \(\check{\mathcal{C}}^\bullet(\mathcal{U}, \mathcal{F})\), read through the product
  equivalence, is exactly the localisation differential \(\operatorname{dDiff}\) of the
  module complex \(D^\bullet\) of Definition~\ref{mk-def:section_cech_module_complex}.
  Concretely, under the per-coordinate identifications each section group
  \(\mathcal{F}(D(s_\sigma))\) is the away localisation \(M_{s_\sigma}\) and each face
  restriction \(\operatorname{sectionCechFaceRestr}(\sigma, i)\) is the away-localisation
  coface \(\operatorname{dCoface}\), so the alternating sum of face restrictions of
  Lemma~\ref{mk-lem:section_cech_objd_apply} becomes the alternating sum of localisation
  cofaces, i.e.\ \(\operatorname{dDiff}\).
\end{lemma}

\begin{proof}
  \uses{mk-lem:section_cech_objd_apply, mk-def:section_cech_module_complex, mk-def:qcoh_sections_localized}
  The proof layers in two steps.

  \emph{(i) Abstract unfold.} By Lemma~\ref{mk-lem:section_cech_objd_apply}, the
  \(\sigma\)-coordinate of the degree-\(p\) differential, read through the product
  equivalence, is the alternating sum
  \(\sum_{i} (-1)^i\, \operatorname{sectionCechFaceRestr}(\sigma, i)\) of presheaf face
  restrictions applied to the deleted-face coordinates. This fixes the shape of the
  differential without any sheaf identification of the restrictions.

  \emph{(ii) Tilde \(\mathcal{F}\)-bridge.} For \(\mathcal{F} = \widetilde{M}\) each
  face restriction is identified with the localisation coface
  \(\operatorname{SectionCechModule.dCoface}\) by transporting along a per-coordinate
  isomorphism. The section group
  \(\mathcal{F}(D(s_\sigma))\) and the localised module
  \(M_{s_\sigma} = \operatorname{LocalizedModule}(\operatorname{Submonoid.powers}
  s_\sigma)\, M\) both localise \(M\) at the submonoid of powers of
  \(s_\sigma = \prod_k s_{\sigma(k)}\): the former by
  \(\operatorname{qcohSectionsAwayLocalized}\) together with the multi-index
  intersection identity \(\operatorname{basicOpen\_sprod}\) (which rewrites
  \(\bigcap_k D(s_{\sigma(k)})\) as the single basic open \(D(s_\sigma)\)), the latter
  by construction. Hence there is a canonical \(R\)-linear isomorphism
  \[
    \varphi_\sigma :
      \operatorname{ToType}\bigl(\mathcal{F}(D(s_\sigma))\bigr)
      \;\xrightarrow{\;\sim\;}\;
      M_{s_\sigma},
  \]
  the comparison of the two \(\operatorname{IsLocalizedModule}\) structures, unique by
  \(\operatorname{AwayComparison.comparison\_unique}\) and supplied by
  \(\operatorname{IsLocalizedModule.iso}\). The naturality square
  \[
    \varphi_\sigma \;\circ\; \operatorname{sectionCechFaceRestr}(\sigma, i)
      \;=\;
      \operatorname{dCoface} \;\circ\; \varphi_{\sigma \circ d_i}
  \]
  is exactly \(\operatorname{qcohRestriction\_eq\_comparison}\) of
  Definition~\ref{mk-def:qcoh_sections_localized}: the presheaf restriction of
  \(\widetilde{M}\) between the two section groups is the away-localisation comparison
  of the two localisation units. This face-naturality is recorded in two forms: at the
  bare localisation level as \(\operatorname{phiL\_naturality}\) and at the sheaf level
  (with \(\varphi_\sigma = \operatorname{phi}\)) as \(\operatorname{phi\_naturality}\).
  Substituting (ii) into (i) and summing with signs
  identifies the section {\v C}ech differential with \(\operatorname{dDiff}\).

  The isomorphism \(\varphi_\sigma\) is obtained by matching the two descriptions of
  the sections of \(\widetilde{M}\) over \(D(s_\sigma)\): the cosimplicial functor
  provides one description and the localisation property of
  Definition~\ref{mk-def:qcoh_sections_localized} provides the other. Their agreement,
  together with the naturality square, is the content of this bridge.
\end{proof}

\begin{lemma}
[Abelian-group exactness criterion for the section {\v C}ech complex]
  \label{mk-lem:section_cech_ab_exact}
  \lean{AlgebraicGeometry.sectionCechAbExact,
    AlgebraicGeometry.sectionCech_isZero_homology_of_objD_exact,
    AlgebraicGeometry.ab_hom_finsetSum_apply,
    AlgebraicGeometry.sectionToModuleAddEquiv,
    AlgebraicGeometry.sectionToModuleAddEquiv_apply,
    AlgebraicGeometry.sectionProdEquiv_symm_apply}
  \dcref{custom}

  \uses{mk-def:section_cech_complex, mk-lem:section_cech_objd_apply,
    mk-lem:section_cech_product_equiv, mk-lem:section_cech_coface_match,
    mk-lem:section_cech_module_exact}
  For \(1 \leq p\) the homology of the section {\v C}ech complex
  \(\check{\mathcal{C}}^\bullet(\mathcal{U}, \mathcal{F})\) of
  Definition~\ref{mk-def:section_cech_complex} vanishes,
  \(\operatorname{IsZero}(\check{\mathcal{C}}^\bullet(\mathcal{U},
  \mathcal{F}).\mathrm{homology}\, p)\). The argument splits into a homological
  reduction and a ladder transport:
  \begin{itemize}
    \item \emph{Reduction.} By standard homological algebra, vanishing of the
      homology at degree \(q+1\) is equivalent to exactness of the short complex
      around that node, and hence to function exactness of the two consecutive
      coface differentials \(d^q\), \(d^{q+1}\) as maps of abelian groups.
    \item \emph{Ladder transport.} The degreewise additive isomorphism
      (Lemma~\ref{mk-lem:section_cech_product_equiv}) between the section complex and
      the localised-module complex of Lemma~\ref{mk-lem:section_cech_module_exact}
      transports the \(R\)-module exactness across an additive ladder, yielding
      function exactness for the group-level coface differentials.
  \end{itemize}
\end{lemma}

\begin{proof}
  \uses{mk-lem:section_cech_objd_apply, mk-lem:section_cech_product_equiv,
    mk-lem:section_cech_coface_match, mk-lem:section_cech_module_exact}
  Fix \(p = q + 1 \geq 1\). By standard homological algebra it suffices to show that
  the two consecutive coface differentials \(d^q\), \(d^{q+1}\) of the section
  complex are exact as maps of abelian groups (exactness at a node iff vanishing
  homology, via the short complex around it).

  The function exactness is established by an additive ladder transport. The
  vertical maps of the ladder are the degreewise additive isomorphisms
  \(e_p\): the product equivalence of Lemma~\ref{mk-lem:section_cech_product_equiv}
  composed with the coordinatewise comparison \(\varphi_\sigma\) that identifies
  each section group with the localised module \(M_{s_\sigma}\). The two
  commuting squares of the ladder follow from
  Lemma~\ref{mk-lem:section_cech_objd_apply} (the differential unfolds as an
  alternating sum of face restrictions) and Lemma~\ref{mk-lem:section_cech_coface_match}
  (each face restriction equals the localisation coface), giving
  \(e_{p+1} \circ d^p = d_M^{\,p} \circ e_p\). The horizontal exactness is
  the \(R\)-module exactness of Lemma~\ref{mk-lem:section_cech_module_exact};
  transporting it across the additive ladder yields function exactness for the
  coface differentials of the section complex.
\end{proof}

\begin{definition}
[The un-localised section {\v C}ech module complex \(D^\bullet\)]
  \label{mk-def:section_cech_module_complex}
  \lean{AlgebraicGeometry.SectionCechModule.dCoeff,
    AlgebraicGeometry.SectionCechModule.dCoface,
    AlgebraicGeometry.SectionCechModule.dDiff,
    AlgebraicGeometry.SectionCechModule.dDiff_apply,
    AlgebraicGeometry.SectionCechModule.dToCech,
    AlgebraicGeometry.SectionCechModule.dToCech_isLocalizedModule,
    AlgebraicGeometry.SectionCechModule.cechCoface_dToCech,
    AlgebraicGeometry.SectionCechModule.dToCech_comm,
    AlgebraicGeometry.SectionCechModule.cechCofaceLin,
    AlgebraicGeometry.SectionCechModule.cechCoface_apply,
    AlgebraicGeometry.SectionCechModule.locDiff,
    AlgebraicGeometry.SectionCechModule.locDiff_apply,
    AlgebraicGeometry.SectionCechModule.locDiff_eq_depDiff,
    AlgebraicGeometry.SectionCechModule.locDiff_exact,
    AlgebraicGeometry.SectionCechModule.fLoc,
    AlgebraicGeometry.SectionCechModule.fLoc_apply,
    AlgebraicGeometry.SectionCechModule.fLoc_isLocalizedModule,
    AlgebraicGeometry.SectionCechModule.locDiff_fLoc,
    AlgebraicGeometry.SectionCechModule.map_dDiff_eq_locDiff,
    AlgebraicGeometry.SectionCechModule.spanIdx,
    AlgebraicGeometry.SectionCechModule.spanIdx_spec,
    AlgebraicGeometry.AwayComparison.comparison_isLocalizedModule,
    AlgebraicGeometry.AwayComparison.Inverts.smul_pow_cancel}
  \dcref{custom}

  \uses{mk-def:standard_affine_cover, mk-def:qcoh_sections_localized}
  Let \(R\) be a commutative ring, \(M = \Gamma(\operatorname{Spec} R, \mathcal{F})\)
  the global sections of a quasi-coherent \(\mathcal{O}_{\operatorname{Spec}
  R}\)-module \(\mathcal{F}\), and \(s : \iota \to R\) (\(\iota\) finite) a spanning
  family of the unit ideal, \(\operatorname{Ideal.span}(\operatorname{range} s) =
  \top\) (Definition~\ref{mk-def:standard_affine_cover}). The \emph{un-localised section
  {\v C}ech module complex} \(D^\bullet\) is the cochain complex of \(R\)-modules
  whose degree-\(p\) term is the product over multi-indices
  \(\sigma : \{0, \ldots, p\} \to \iota\) of the away localisations
  \[
    D^p \;=\; \prod_{\sigma} M_{s_\sigma},
    \qquad s_\sigma = \prod_{k} s_{\sigma(k)},
    \qquad M_{s_\sigma} = \operatorname{LocalizedModule}
      \bigl(\operatorname{Submonoid.powers}(s_\sigma)\bigr)\, M,
  \]
  with differential \(d^p : D^p \to D^{p+1}\) the alternating sum
  \(\sum_{j} (-1)^j\, \delta^j\) of the coface comparison maps \(\delta^j\), each of
  which is the canonical away-localisation comparison
  \(M_{s_{\sigma \circ d_j}} \to M_{s_\sigma}\) sending a fraction to its image under
  the further localisation at the omitted factor. The single-coefficient comparison
  map and its localisation property are recorded by the \(\mathtt{dCoface}\) and
  \(\mathtt{dCoeff}\) data and the differential by \(\mathtt{dDiff}\) (with action
  \(\mathtt{dDiff\_apply}\)); the auxiliary \(\mathtt{spanIdx}\)/\(\mathtt{spanIdx\_spec}\)
  select a witnessing index of the spanning relation.

  The complex \(D^\bullet\) is tied to the section {\v C}ech complex of
  Definition~\ref{mk-def:section_cech_complex} and to its localisations by three
  comparison families. First, the maps \(\mathtt{dToCech}\) identify each term
  \(D^p\) with the degree-\(p\) section group
  \(\prod_\sigma \mathcal{F}(D(s_\sigma))\) of
  Definition~\ref{mk-def:qcoh_sections_localized} as a localised module
  (\(\mathtt{dToCech\_isLocalizedModule}\)), intertwining the differentials
  (\(\mathtt{cechCoface\_dToCech}\), \(\mathtt{dToCech\_comm}\)); the section-side
  cofaces are \(\mathtt{cechCofaceLin}\) with action \(\mathtt{cechCoface\_apply}\).
  Second, localising \(D^\bullet\) at a spanning element \(s_r\) reproduces the
  \emph{localised} {\v C}ech complex: the localised differential \(\mathtt{locDiff}\)
  (action \(\mathtt{locDiff\_apply}\)) agrees, via
  \(\mathtt{locDiff\_eq\_depDiff}\), with the dependent-coefficient combinatorial
  differential, and its exactness is recorded as \(\mathtt{locDiff\_exact}\). The
  localisation comparison maps \(\mathtt{fLoc}\) (action \(\mathtt{fLoc\_apply}\))
  exhibit each localised term as a localised module
  (\(\mathtt{fLoc\_isLocalizedModule}\)) intertwining \(\mathtt{dDiff}\) with
  \(\mathtt{locDiff}\) (\(\mathtt{locDiff\_fLoc}\), \(\mathtt{map\_dDiff\_eq\_locDiff}\)).
  The key algebraic fact underlying all three families is the localisation
  transitivity \(M_a[1/b] = M_{ab}\): localising the away module \(M_{s_\sigma}\)
  further at \(s_r\) gives \(M_{s_r s_\sigma}\). This is the keystone
  \(\mathtt{AwayComparison.comparison\_isLocalizedModule}\) (the composite of two away
  localisations is again the away localisation at the product), with the cancellation
  bookkeeping \(\mathtt{AwayComparison.Inverts.smul\_pow\_cancel}\). Thus localising
  \(D^\bullet\) at \(s_r\) (an \(\operatorname{IsLocalizedModule.Away}\)) reproduces
  the localised Čech complex of Lemma~\ref{mk-lem:section_cech_homology_exact}, and
  exactness of \(D^\bullet\) can be checked one localisation at a time.
\end{definition}

\begin{lemma}[Positive-degree exactness of the section {\v C}ech module complex]\label{mk-lem:section_cech_module_exact}
  \lean{AlgebraicGeometry.SectionCechModule.dDiff_exact}
  \leanok
  \dcref{custom}

  \uses{mk-def:section_cech_module_complex}
  With the notation of Definition~\ref{mk-def:section_cech_module_complex}, the
  un-localised section {\v C}ech module complex \(D^\bullet\) is exact in positive
  degrees: for every \(p\),
  \[
    \operatorname{Function.Exact}\bigl(d^{p+1} : D^{p+1} \to D^{p+2},\;
      d^{p+2} : D^{p+2} \to D^{p+3}\bigr).
  \]
  The homology vanishing of the section {\v C}ech complex follows from this exactness
  by Lemma~\ref{mk-lem:section_cech_homology_exact}.
\end{lemma}

\begin{proof}
  \leanok
  \uses{mk-def:section_cech_module_complex}
  By the local-to-global criterion for a spanning family, it suffices to prove
  exactness after localizing at each \(s_r\). Localization transitivity identifies the
  localized differential with the {\v C}ech differential whose coefficients are all
  localized at \(s_r\). Since \(s_r\) is then invertible, prepending the fixed index
  \(r\) defines a contracting homotopy and gives \(dh+hd=\operatorname{id}\). Thus the
  complex is exact after every localization in the spanning family, and hence exact
  before localization.
\end{proof}

\begin{lemma}

[Standard-cover {\v C}ech-complex vanishing on affines]
  \label{mk-lem:cech_acyclic_affine}
  \lean{AlgebraicGeometry.sectionCech_affine_vanishing,
    AlgebraicGeometry.CombinatorialCech.combDifferential,
    AlgebraicGeometry.CombinatorialCech.combHomotopy,
    AlgebraicGeometry.CombinatorialCech.combHomotopy_zero,
    AlgebraicGeometry.CombinatorialCech.cons_comp_succAbove_succ,
    AlgebraicGeometry.CombinatorialCech.combHomotopy_spec,
    AlgebraicGeometry.CombinatorialCech.combDifferential_eq_of_cocycle,
    AlgebraicGeometry.CombinatorialCech.combSign_flip,
    AlgebraicGeometry.CombinatorialCech.combDifferential_comp,
    AlgebraicGeometry.CombinatorialCech.combDifferential_exact,
    AlgebraicGeometry.CombinatorialCech.depDiff,
    AlgebraicGeometry.CombinatorialCech.depHomotopy,
    AlgebraicGeometry.CombinatorialCech.depHomotopy_spec,
    AlgebraicGeometry.CombinatorialCech.depTransport,
    AlgebraicGeometry.CombinatorialCech.cons_comp_zero_succAbove,
    AlgebraicGeometry.CombinatorialCech.comp_succAbove_swap,
    AlgebraicGeometry.CombinatorialCech.depDiff_eq_of_cocycle,
    AlgebraicGeometry.CombinatorialCech.depDiff_comp,
    AlgebraicGeometry.CombinatorialCech.depDiff_exact}
  \dcref{custom}

  \uses{mk-def:section_cech_complex, mk-def:standard_affine_cover}
  \textit{Source: Stacks Project, Cohomology of Schemes,
  \texttt{lemma-cech-cohomology-quasi-coherent-trivial} (standard-cover {\v C}ech
  vanishing).}
  Let \(R\) be a commutative ring, \(U = \operatorname{Spec}(R)\) the associated
  affine scheme, and \(\mathcal{F}\) a quasi-coherent \(\mathcal{O}_U\)-module
  (\(\mathcal{F} : (\operatorname{Spec} R).\mathrm{Modules}\)). Let
  \(s : \iota \to R\) (\(\iota\) finite) be a spanning family of the unit ideal,
  \(\operatorname{Ideal.span}(\operatorname{range} s) = \top\), and let
  \(\mathcal{U} : U = \bigcup_{i \in \iota} D(s_i)\) be the associated standard
  affine cover of Definition~\ref{mk-def:standard_affine_cover}; the index set and the
  affine pieces \(D(s_i)\) with sections \(R_{s_i}\) are read off directly from the
  spanning family. The statement is \emph{intrinsic} to \(\operatorname{Spec} R\):
  there is no relative morphism \(f\) and no pushforward in the conclusion. Then the
  \emph{section {\v C}ech complex} \(\check{\mathcal{C}}^\bullet(\mathcal{U},
  \mathcal{F})\) of Definition~\ref{mk-def:section_cech_complex} --- the cochain
  complex of abelian groups whose degree-\(p\) term is
  \(\prod_{\sigma}\mathcal{F}(D(s_\sigma))\) --- has vanishing cohomology in all
  positive degrees:
  \[
    \check{H}^p(\mathcal{U}, \mathcal{F}) = 0 \qquad \text{for all } p > 0,
  \]
  where \(\check{H}^p\) denotes the cohomology of the section complex
  \(\check{\mathcal{C}}^\bullet(\mathcal{U}, \mathcal{F})\). Equivalently, the
  extended complex
  \(0 \to M \to \prod_{i_0} M_{s_{i_0}} \to \prod_{i_0 i_1} M_{s_{i_0} s_{i_1}}
  \to \cdots\) of localisations of \(M = \Gamma(\operatorname{Spec} R, \mathcal{F})\)
  is exact in positive degrees. This is the standard-cover {\v C}ech vanishing
  only; the sheaf-cohomology Serre vanishing \(H^p(U, \mathcal{F}) = 0\) is the
  separate Lemma~\ref{mk-lem:affine_serre_vanishing} below, obtained from this complex
  vanishing by the basis-comparison Lemma~\ref{mk-lem:cech_to_cohomology_on_basis}.
\end{lemma}
\begin{proof}
  \uses{mk-def:section_cech_complex, mk-def:qcoh_sections_localized, mk-lem:section_cech_homology_exact, mk-lem:section_cech_module_exact, mk-lem:section_cech_objd_apply, mk-lem:section_cech_coface_match}
  Write \(R\) for the ring, \(M = \Gamma(\operatorname{Spec} R, \mathcal{F})\) for the
  global sections, and \((s : \iota \to R,\; hs : \operatorname{Ideal.span}
  (\operatorname{range} s) = \top)\) for the spanning family of
  Definition~\ref{mk-def:standard_affine_cover}.

  \paragraph{Reduction to the localised complex.}
  By the section identification of Definition~\ref{mk-def:qcoh_sections_localized}, each
  \((p+1)\)-fold intersection \(D(s_{\sigma(0)}) \cap \cdots \cap D(s_{\sigma(p)})\)
  is the basic open \(D(s_\sigma)\) of the product
  \(s_\sigma = \prod_{k} s_{\sigma(k)}\), and its sections are the away localisation
  \(\mathcal{F}(D(s_\sigma)) = M_{s_\sigma}\); under these identifications each face
  restriction is the canonical localisation map. Hence
  Lemma~\ref{mk-lem:section_cech_homology_exact} identifies the section {\v C}ech complex
  \(\check{\mathcal{C}}^\bullet(\mathcal{U}, \mathcal{F})\) of
  Definition~\ref{mk-def:section_cech_complex} with the concrete complex of localised
  \(R\)-modules
  \[
    \prod_{i_0} M_{s_{i_0}} \;\to\;
    \prod_{i_0 i_1} M_{s_{i_0} s_{i_1}} \;\to\;
    \prod_{i_0 i_1 i_2} M_{s_{i_0} s_{i_1} s_{i_2}} \;\to\; \cdots,
  \]
  and reduces the vanishing
  \(\operatorname{IsZero}(\check{\mathcal{C}}^\bullet(\mathcal{U},
  \mathcal{F}).\mathrm{homology}\, p)\) for \(p > 0\) to positive-degree
  \(\operatorname{Function.Exact}\) of the underlying localisation maps. Equivalently,
  the extended complex \(0 \to M \to \prod_{i_0} M_{s_{i_0}} \to \cdots\) (whose
  truncation is studied in Algebra, Lemma \texttt{algebra-lemma-cover-module}) is
  exact in positive degrees. For \(\mathcal{F}=\widetilde M\), the coface
  identification is Lemma~\ref{mk-lem:section_cech_coface_match}. The general case follows
  from the affine equivalence
  \(\mathcal{F}\cong\widetilde{\Gamma\mathcal{F}}\).

  \paragraph{Local exactness.}
  By the local-to-global exactness criterion for the spanning family \((s,hs)\), it
  remains to prove exactness after localization at each \(s_r\). In that localization,
  \(s_r\) is invertible. Prepending the fixed index \(r\) therefore defines a homotopy
  \(h\) satisfying \(dh+hd=\operatorname{id}\): the term involving the first coface is
  the identity, while the remaining terms cancel in pairs after shifting the omitted
  index. Hence every localized complex is contractible in positive degrees. The
  local-to-global criterion proves exactness of the original complex, and
  Lemma~\ref{mk-lem:section_cech_homology_exact} gives the asserted vanishing.
\end{proof}

\subsection{Presheaf-level {\v C}ech machinery}
\label{mk-sec:presheaf_cech_machinery}

The {\v C}ech-acyclicity of injective \(\mathcal{O}_X\)-modules
(Lemma~\ref{mk-lem:injective_cech_acyclic} below) is not a statement about a single
affine cover: it is proved on the abelian category \(\mathrm{PMod}(\mathcal{O}_X)\)
of presheaves of \(\mathcal{O}_X\)-modules, where the {\v C}ech complex of sections
is the \(\operatorname{Hom}\)-complex of an explicit free resolution. The argument
rests on two presheaf-level facts, recorded here, and one transfer of injectivity.
First, a free complex \(K(\mathcal{U})_\bullet\) of presheaves of modules whose
\(\operatorname{Hom}\)-complex into \(\mathcal{F}\) \emph{is} the {\v C}ech complex
of sections (Definition~\ref{mk-def:cech_free_presheaf_complex},
Definition~\ref{mk-def:section_cech_complex} and
Lemma~\ref{mk-lem:cech_complex_hom_identification}). Second, the resolution property of
that complex: \(K(\mathcal{U})_\bullet\) resolves \(\mathcal{O}_{\mathcal{U}}\)
concentrated in degree \(0\) (Lemma~\ref{mk-lem:cech_free_complex_quasi_iso}). Together
these give: an injective \(\mathcal{O}_X\)-module is injective as a presheaf
(Lemma~\ref{mk-lem:injective_cech_acyclic} via
Lemma~\ref{mk-lem:injective_of_adjoint} applied to the sheafification adjunction
Lemma~\ref{mk-lem:mod_pmod_adjunction}), so \(\operatorname{Hom}(-, \mathcal{I})\) is
exact and carries the augmented resolution to an exact section complex, i.e.\
\(\check{H}^p(\mathcal{U}, \mathcal{I}) = 0\) for \(p > 0\). This direct route uses
neither enough-injectives in \(\mathrm{PMod}(\mathcal{O}_X)\) nor a right-derived
\(\delta\)-functor identification.

\begin{definition}

[Free presheaf complex of a cover]
  \label{mk-def:cech_free_presheaf_complex}
  \lean{AlgebraicGeometry.cechFreePresheafComplex,
    AlgebraicGeometry.freeYoneda,
    AlgebraicGeometry.coverOpen,
    AlgebraicGeometry.coverInterOpen,
    AlgebraicGeometry.coverInterOpen_comp_le,
    AlgebraicGeometry.cechFreeSimplicial,
    AlgebraicGeometry.cechFreePresheafComplex_X,
    AlgebraicGeometry.sigma_ι_eqToHom_transport,
    AlgebraicGeometry.cechFreePresheafComplexFam,
    AlgebraicGeometry.coverInterOpenFam,
    AlgebraicGeometry.coverInterOpen_comp_leFam,
    AlgebraicGeometry.cechFreeSimplicialFam,
    AlgebraicGeometry.cechFreePresheafComplex_XFam}
  \dcref{custom}

  \textit{Source: Stacks Project, Cohomology,
  \texttt{lemma-cech-map-into} (the free complex \(K(\mathcal{U})_\bullet\)).}
  Let \(X\) be a ringed space and \(\mathcal{U} : U = \bigcup_{i \in I} U_i\) an
  open covering, with \(j_{i_0 \ldots i_p} : U_{i_0 \ldots i_p} \hookrightarrow X\)
  the open immersion of each finite intersection. The Stacks formulation uses, for
  each intersection \(U_{i_0 \ldots i_p}\), the \emph{extension-by-zero} presheaf
  \((j_{i_0 \ldots i_p})_{!}\,\mathcal{O}_X|_{U_{i_0 \ldots i_p}}\), the presheaf of
  modules sending an open \(W\) to \(\mathcal{O}_X(W)\) when
  \(W \subset U_{i_0 \ldots i_p}\) and to \(0\) otherwise. This extension-by-zero
  presheaf is, canonically, the \emph{free presheaf of modules on the representable}
  \(\operatorname{free}(\mathbf{y}\,U_{i_0 \ldots i_p})\): writing
  \(\mathbf{y}\,V = \operatorname{Hom}(-, V)\) for the Yoneda presheaf of the open
  \(V\) on the site of opens of \(X\), the free presheaf of modules
  \(\operatorname{free}(\mathbf{y}\,V)\) sends \(W\) to the free
  \(\mathcal{O}_X(W)\)-module on the set \(\operatorname{Hom}(W, V)\), which is
  \(\mathcal{O}_X(W)\) if \(W \subset V\) and \(0\) otherwise — exactly
  extension-by-zero of \(\mathcal{O}_X|_V\). Its defining universal property
  \[
    \operatorname{Hom}_{\mathrm{PMod}(\mathcal{O}_X)}\bigl(
      \operatorname{free}(\mathbf{y}\,V),\, \mathcal{F}\bigr)
    \;\cong\; \mathcal{F}(V)
  \]
  is the free–forgetful adjunction for presheaves of modules composed with the
  Yoneda isomorphism \(\operatorname{Hom}(\mathbf{y}\,V, G) \cong G(V)\). The
  \emph{free presheaf complex} of \(\mathcal{U}\) is the chain complex
  \(K(\mathcal{U})_\bullet\) in the abelian category
  \(\mathrm{PMod}(\mathcal{O}_X)\) of presheaves of \(\mathcal{O}_X\)-modules whose
  term in degree \(p\) is the direct sum
  \[
    K(\mathcal{U})_p
      = \bigoplus_{(i_0, \ldots, i_p) \in I^{p+1}}
        \operatorname{free}\bigl(\mathbf{y}\,U_{i_0 \ldots i_p}\bigr),
  \]
  with \(K(\mathcal{U})_p\) placed in degree \(p \geq 0\) (and \(0\) in negative
  degrees), and whose differential \(K(\mathcal{U})_{p+1} \to K(\mathcal{U})_p\) is
  the alternating sum \(\sum_{j=0}^{p+1} (-1)^j\) of the maps
  \(\operatorname{free}(\mathbf{y}\,U_{i_0 \ldots i_{p+1}}) \to
  \operatorname{free}(\mathbf{y}\,U_{i_0 \ldots \widehat{i_j} \ldots i_{p+1}})\)
  obtained by applying \(\operatorname{free}\) to the representable index-dropping
  maps \(\mathbf{y}\,U_{i_0 \ldots i_{p+1}} \to \mathbf{y}\,U_{i_0 \ldots
  \widehat{i_j} \ldots i_{p+1}}\) induced by the inclusions of opens
  \(U_{i_0 \ldots i_{p+1}} \subset U_{i_0 \ldots \widehat{i_j} \ldots i_{p+1}}\)
  that drop the index \(i_j\). Each summand is a \emph{projective} object of
  \(\mathrm{PMod}(\mathcal{O}_X)\), because the universal property above identifies
  \(\operatorname{Hom}_{\mathcal{O}_X}(\operatorname{free}(\mathbf{y}\,U_{i_0 \ldots
  i_p}), -)\) with the exact ``sections over \(U_{i_0 \ldots i_p}\)'' functor
  \((-)(U_{i_0 \ldots i_p})\); this is the structural feature exploited below.

  We assume the index set \(I\) is \emph{finite} (so the required coproducts exist in
  \(\mathrm{PMod}(\mathcal{O}_X)\); this matches the finiteness hypothesis of the
  comparison target). Here the direct sum \(\bigoplus\) is the categorical coproduct
  \(\coprod\) (the indexed \(\operatorname{Sigma}\) colimit), which for a finite index
  set coincides with the finite biproduct.

  The parallel \texttt{...Fam} declarations pinned above are the cover-agnostic
  raw-finite-family mirror of this construction — indexed by an arbitrary finite
  family \((U_i)_{i \in I}\) of opens carrying \emph{no} covering hypothesis — and
  feed the family form of the {\v C}ech acyclicity lemma
  (\(\operatorname{injective\_cech\_acyclicFam}\)).
\end{definition}

\begin{definition}

[Section {\v C}ech complex of a presheaf of modules]
  \label{mk-def:section_cech_complex}
  \lean{AlgebraicGeometry.sectionCechComplex,
    AlgebraicGeometry.sectionCechCosimplicial}
  \dcref{custom}

  \uses{mk-def:cech_complex}
  \textit{Source: Stacks Project, Cohomology, \S{\v C}ech cohomology of presheaves.}
  Let \(X\) be a ringed space and \(\mathcal{U} : U = \bigcup_{i \in I} U_i\) an open
  covering. For a presheaf of \(\mathcal{O}_X\)-modules \(\mathcal{F}\), the
  \emph{section {\v C}ech complex} \(\check{\mathcal{C}}^\bullet(\mathcal{U},
  \mathcal{F})\) is the cochain complex of \emph{abelian groups} whose term in
  degree \(p\) is the product of sections over the \((p+1)\)-fold intersections,
  \[
    \check{\mathcal{C}}^p(\mathcal{U}, \mathcal{F})
      = \prod_{(i_0, \ldots, i_p) \in I^{p+1}}
        \mathcal{F}\bigl(U_{i_0 \ldots i_p}\bigr),
  \]
  each factor being the value of the ``sections over \(U_{i_0 \ldots i_p}\)'' functor
  (the evaluation of \(\mathcal{F}\) at the open \(U_{i_0 \ldots i_p}\)), with
  differential the alternating sum of restriction maps
  \[
    d(s)_{i_0 \ldots i_{p+1}}
      = \sum_{j=0}^{p+1} (-1)^j\,
        s_{i_0 \ldots \widehat{i_j} \ldots i_{p+1}}
        \big|_{U_{i_0 \ldots i_{p+1}}}.
  \]
  This object is \emph{distinct} from the relative {\v C}ech complex
  \(\check{\mathcal{C}}^\bullet(\mathfrak{U}, \mathcal{F})\) of
  Definition~\ref{mk-def:cech_complex}: the latter is the pushforward complex with terms
  \(\prod (f|_{U_{i_0 \ldots i_p}})_*(\mathcal{F}|_{U_{i_0 \ldots i_p}})\) in
  \(\operatorname{QCoh}(S)\), an object of relative sheaves over the base, whereas the
  section {\v C}ech complex here is the simpler complex of \emph{global sections} of
  \(\mathcal{F}\) over the intersections, an object of abelian groups (the underlying
  additive structure of the sections). The two must not be conflated: the
  presheaf-level acyclicity argument below runs entirely on this section complex.
\end{definition}

\begin{lemma}
[The free complex represents the {\v C}ech complex]
  \label{mk-lem:cech_complex_hom_identification}
  \lean{AlgebraicGeometry.cechComplex_hom_identification,
    AlgebraicGeometry.freeYonedaHomEquiv,
    AlgebraicGeometry.freeYonedaHomAddEquiv,
    AlgebraicGeometry.freeYonedaHomEquiv_apply,
    AlgebraicGeometry.homCechComplex,
    AlgebraicGeometry.homCechCosimplicial,
    AlgebraicGeometry.homCechSectionIsoApp,
    AlgebraicGeometry.homCechSectionIsoApp_hom_π,
    AlgebraicGeometry.pi_mapIso_hom_eq,
    AlgebraicGeometry.freeYonedaHomAddEquiv_naturality,
    AlgebraicGeometry.homCechSectionCosimplicialIso,
    AlgebraicGeometry.cechComplex_hom_identificationFam,
    AlgebraicGeometry.homCechCosimplicialFam,
    AlgebraicGeometry.homCechComplexFam,
    AlgebraicGeometry.homCechSectionIsoAppFam,
    AlgebraicGeometry.homCechSectionIsoApp_hom_πFam,
    AlgebraicGeometry.homCechCosimplicial_δFam,
    AlgebraicGeometry.homCechSectionCosimplicialIsoFam}
  \dcref{custom}

  \uses{mk-def:cech_free_presheaf_complex, mk-def:section_cech_complex}
  \textit{Source: Stacks Project, Cohomology,
  \texttt{lemma-cech-map-into}.}
  Let \(X\) be a ringed space and \(\mathcal{U} : U = \bigcup_{i \in I} U_i\) an
  open covering. For every presheaf of \(\mathcal{O}_X\)-modules \(\mathcal{F}\)
  there is an isomorphism of cochain complexes of \emph{abelian groups}
  \[
    \operatorname{Hom}_{\mathcal{O}_X}\bigl(K(\mathcal{U})_\bullet, \mathcal{F}\bigr)
      \;=\;
    \check{\mathcal{C}}^\bullet(\mathcal{U}, \mathcal{F}),
  \]
  natural in \(\mathcal{F} \in \mathrm{PMod}(\mathcal{O}_X)\). Here the left side is
  the cochain complex obtained by applying the contravariant
  \(\operatorname{Hom}_{\mathcal{O}_X}(-, \mathcal{F})\) to the chain complex
  \(K(\mathcal{U})_\bullet\) of Definition~\ref{mk-def:cech_free_presheaf_complex},
  and the right side is the section {\v C}ech complex of
  Definition~\ref{mk-def:section_cech_complex}.
\end{lemma}
\begin{proof}
  \uses{mk-def:cech_free_presheaf_complex, mk-def:section_cech_complex}
  The free–forgetful adjunction for presheaves of modules gives, composed with the
  Yoneda isomorphism \(\operatorname{Hom}(\mathbf{y}\,V, \mathcal{F}) \cong
  \mathcal{F}(V)\), the natural identification, for each multi-index,
  \[
    \operatorname{Hom}_{\mathcal{O}_X}\bigl(
      \operatorname{free}(\mathbf{y}\,U_{i_0 \ldots i_p}),\, \mathcal{F}
    \bigr)
    \;\cong\;
    \operatorname{Hom}\bigl(\mathbf{y}\,U_{i_0 \ldots i_p},\,
      \mathcal{F}\bigr)
    \;\cong\;
    \mathcal{F}(U_{i_0 \ldots i_p}),
  \]
  where the last group is the evaluation of \(\mathcal{F}\) at the open
  \(U_{i_0 \ldots i_p}\). Taking the product over multi-indices,
  \(\operatorname{Hom}_{\mathcal{O}_X}(K(\mathcal{U})_p, \mathcal{F}) =
  \prod_{i_0 \ldots i_p} \mathcal{F}(U_{i_0 \ldots i_p}) =
  \check{\mathcal{C}}^p(\mathcal{U}, \mathcal{F})\); the direct sum
  \(K(\mathcal{U})_p\) therefore represents the functor \(\mathcal{F} \mapsto
  \prod_{i_0 \ldots i_p} \mathcal{F}(U_{i_0 \ldots i_p})\). Under this
  identification the \((-1)^j\) maps \(\operatorname{free}(\mathbf{y}\,U_{i_0
  \ldots i_{p+1}}) \to \operatorname{free}(\mathbf{y}\,U_{i_0 \ldots \widehat{i_j}
  \ldots i_{p+1}})\) defining the differential of \(K(\mathcal{U})_\bullet\) become
  exactly the alternating-sum restriction differential of
  \(\check{\mathcal{C}}^\bullet(\mathcal{U}, \mathcal{F})\), so the two cochain
  complexes coincide; naturality in \(\mathcal{F}\) is immediate from the
  naturality of the free–forgetful adjunction and of Yoneda. The per-multi-index
  bijection \(\operatorname{Hom}_{\mathcal{O}_X}(\operatorname{free}(\mathbf{y}\,
  U_{i_0 \ldots i_p}), \mathcal{F}) \cong \mathcal{F}(U_{i_0 \ldots i_p})\) is additive,
  so it is promoted to an isomorphism of abelian groups (an \emph{additive}
  equivalence); taking the product over multi-indices and assembling these degreewise
  additive isomorphisms — which commute with the alternating-sum differentials — yields
  the complex isomorphism by matching components degree by degree.
\end{proof}

\begin{definition}

[The cover structure presheaf \(\mathcal{O}_{\mathcal{U}}\)]
  \label{mk-def:cover_structure_presheaf}
  \lean{AlgebraicGeometry.coverStructurePresheaf,
    AlgebraicGeometry.cechFreeAug,
    AlgebraicGeometry.cechFreeComplexAug,
    AlgebraicGeometry.cechFreeComplexAug_f_zero,
    AlgebraicGeometry.cechFreeSimplicial_δ_comp_aug,
    AlgebraicGeometry.cechFree_d_comp_aug,
    AlgebraicGeometry.cechFree_d_comp_factorThruImage,
    AlgebraicGeometry.freeYonedaAug,
    AlgebraicGeometry.freeYonedaAug_app_freeMk,
    AlgebraicGeometry.freeYonedaHomEquiv_freeYonedaAug,
    AlgebraicGeometry.freeYoneda_map_comp_aug,
    AlgebraicGeometry.coverStructurePresheafFam,
    AlgebraicGeometry.cechFreeAugFam,
    AlgebraicGeometry.cechFreeComplexAugFam,
    AlgebraicGeometry.cechFreeComplexAug_f_zeroFam,
    AlgebraicGeometry.cechFreeSimplicial_δ_comp_augFam,
    AlgebraicGeometry.cechFree_d_comp_augFam,
    AlgebraicGeometry.cechFree_d_comp_factorThruImageFam}
  \dcref{custom}

  \uses{mk-def:cech_free_presheaf_complex}
  \textit{Source: Stacks Project, Cohomology, \texttt{lemma-homology-complex}.}
  Let \(X\) be a ringed space and \(\mathcal{U} : U = \bigcup_{i \in I} U_i\) an open
  covering. The \emph{cover structure presheaf} \(\mathcal{O}_{\mathcal{U}} :
  X.\mathrm{PresheafOfModules}\) is the image presheaf of the augmentation
  \[
    \coprod_i \operatorname{free}(\mathbf{y}\,U_i) \;\longrightarrow\; \mathbf{1},
  \]
  where the target \(\mathbf{1} = \operatorname{PresheafOfModules.unit}(
  X.\mathrm{ringCatSheaf}.\mathrm{obj})\) is the structure presheaf \(\mathcal{O}_X\)
  regarded as a presheaf of modules over itself. Concretely
  \(\mathcal{O}_{\mathcal{U}}(W) = \mathcal{O}_X(W)\) when \(W\) is contained in some
  \(U_i\), and \(0\) otherwise. The degree-\(0\) term
  \(K(\mathcal{U})_0 = \coprod_i \operatorname{free}(\mathbf{y}\,U_i)\) of the free
  presheaf complex of Definition~\ref{mk-def:cech_free_presheaf_complex} carries the
  augmentation map \(K(\mathcal{U})_0 \to \mathcal{O}_{\mathcal{U}}\) onto this image
  presheaf.

  The parallel \texttt{...Fam} declarations pinned above are the cover-agnostic
  raw-finite-family mirror of this construction (an arbitrary finite family of opens,
  no covering hypothesis), feeding the family form of the {\v C}ech acyclicity lemma.
\end{definition}

\begin{lemma}
[Objectwise detection of quasi-isomorphisms of presheaf-of-module complexes]
  \label{mk-lem:quasiIso_of_evaluation}
  \lean{AlgebraicGeometry.quasiIso_of_evaluation,
    AlgebraicGeometry.isIso_Fmap_homologyMap,
    AlgebraicGeometry.isIso_of_evaluation}
  \dcref{custom}

  \uses{mk-def:cech_free_presheaf_complex}
  \textit{Source: Stacks Project, Cohomology, \texttt{lemma-homology-complex}
  (sectionwise reduction).}
  Let \(R\) be a presheaf of rings on a category \(C\) and let
  \(\varphi : K \to L\) be a morphism of chain complexes in the abelian category
  \(\mathrm{PMod}(R)\) of presheaves of \(R\)-modules. Suppose that for \emph{every}
  object \(V\) the evaluated morphism — the image of \(\varphi\) under the functor
  \(\bigl(\operatorname{PresheafOfModules.evaluation} R\, V\bigr).
  \operatorname{mapHomologicalComplex}\), a morphism of complexes of
  \(R(V)\)-modules — is a quasi-isomorphism. Then \(\varphi\) is a
  quasi-isomorphism. This is the objectwise-reduction step that turns the
  resolution claim of Lemma~\ref{mk-lem:cech_free_complex_quasi_iso} into a family of
  sectionwise (combinatorial) statements: homology in \(\mathrm{PMod}(R)\) is
  computed open-by-open, because the evaluation functors
  \(\operatorname{PresheafOfModules.evaluation} R\, V\) are exact (hence preserve
  homology) and jointly conservative: a morphism in \(\mathrm{PMod}(R)\) is an isomorphism
  iff each of its evaluations is an isomorphism, and each evaluation functor preserves
  homology.
\end{lemma}

\begin{proof}
  \uses{mk-def:cech_free_presheaf_complex}
  A morphism of complexes is a quasi-isomorphism iff it induces an isomorphism on
  homology in every degree (\(\operatorname{quasiIso\_iff}\),
  \(\operatorname{quasiIsoAt\_iff\_isIso\_homologyMap}\)). Fix a degree \(i\). The
  homology map \(H_i(\varphi)\) is an isomorphism in \(\mathrm{PMod}(R)\) iff each of
  its evaluations is an isomorphism, since the underlying
  presheaf-of-abelian-groups functor reflects isomorphisms and a morphism of
  presheaves is an isomorphism iff it is so objectwise (joint conservativity). For
  each \(V\), evaluation preserves homology, so the evaluation of \(H_i(\varphi)\)
  is the homology map of the evaluated complex morphism; the latter is an
  isomorphism because the evaluated morphism is a quasi-isomorphism by hypothesis.
  Hence \(H_i(\varphi)\) is an isomorphism for all \(i\), i.e.\ \(\varphi\) is a
  quasi-isomorphism.
\end{proof}

\begin{lemma}

[Sectionwise description of the evaluated free complex]
  \label{mk-lem:cech_free_eval_sectionwise}
  \lean{AlgebraicGeometry.cechFreeEval_X,
    AlgebraicGeometry.freeYonedaEval_iso_of_le,
    AlgebraicGeometry.freeYonedaEval_isZero_of_not_le,
    AlgebraicGeometry.cechFreeEval_XFam}
  \dcref{custom}

  \uses{mk-def:cech_free_presheaf_complex, mk-def:cover_structure_presheaf}
  \textit{Source: Stacks Project, Cohomology, \texttt{lemma-homology-complex}
  (description of the evaluated complex).}
  Fix an open \(V \subset X\) and write \(I_1 = I_1(V) = \{\,i \in I : V \leq U_i\,\}\)
  for the set of cover members containing \(V\). Evaluating the free presheaf complex
  \(K(\mathcal{U})_\bullet\) of Definition~\ref{mk-def:cech_free_presheaf_complex} at
  \(V\) gives, degreewise,
  \[
    K(\mathcal{U})_p(V)
      = \bigoplus_{\sigma : \mathrm{Fin}(p+1) \to I}
        \operatorname{free}(\mathbf{y}\,U_\sigma)(V)
      = \bigoplus_{\substack{\sigma : \mathrm{Fin}(p+1) \to I \\ V \leq U_\sigma}}
        \mathcal{O}_X(V)
      = \bigoplus_{\sigma : \mathrm{Fin}(p+1) \to I_1} \mathcal{O}_X(V),
  \]
  because \(\operatorname{free}(\mathbf{y}\,W)(V) = \mathcal{O}_X(V)\) when
  \(V \leq W\) and \(0\) otherwise (poset hom-set is a subsingleton), and
  \(V \leq U_\sigma = U_{\sigma(0)} \cap \cdots \cap U_{\sigma(p)}\) holds iff
  \(\sigma\) takes all its values in \(I_1\). Under this identification the
  differential of \(K(\mathcal{U})_\bullet(V)\) is the alternating-sum combinatorial
  {\v C}ech differential of the full simplex on the index set \(I_1\) with
  coefficients in the single module \(\mathcal{O}_X(V)\). Moreover the evaluated
  augmentation target is \(\mathcal{O}_{\mathcal{U}}(V) = \mathcal{O}_X(V)\) when
  \(I_1 \neq \varnothing\) and \(\mathcal{O}_{\mathcal{U}}(V) = 0\) when
  \(I_1 = \varnothing\). Thus the evaluation at \(V\) of the augmented complex
  \(K(\mathcal{U})_\bullet \to \mathcal{O}_{\mathcal{U}}[0]\)
  is precisely the augmented combinatorial
  {\v C}ech complex of the simplex on \(I_1\) with coefficients \(\mathcal{O}_X(V)\).
\end{lemma}

\begin{proof}
  \uses{mk-def:cech_free_presheaf_complex, mk-def:cover_structure_presheaf}
  The degreewise identity is the value at \(V\) of the explicit description of the
  free presheaf \(\operatorname{free}(\mathbf{y}\,W)\) recorded in
  Definition~\ref{mk-def:cech_free_presheaf_complex}: it sends \(V\) to the free
  \(\mathcal{O}_X(V)\)-module on \(\operatorname{Hom}(V, W)\), which is
  \(\mathcal{O}_X(V)\) if \(V \leq W\) and \(0\) otherwise. Applying
  \(\operatorname{PresheafOfModules.evaluation}\;V\) to the coproduct decomposition
  \(\operatorname{cechFreePresheafComplex\_X}\) commutes with the direct sum and
  collapses each summand to \(\mathcal{O}_X(V)\) or \(0\) according to whether
  \(V \leq U_\sigma\), i.e.\ whether \(\sigma\) factors through \(I_1\); the surviving
  summands are exactly those indexed by \(\sigma : \mathrm{Fin}(p+1) \to I_1\). The
  evaluated faces are the index-dropping restriction maps, which on the surviving
  summands are identities of \(\mathcal{O}_X(V)\), so the alternating sum is the
  combinatorial {\v C}ech differential. The augmentation value follows from
  \(\mathcal{O}_{\mathcal{U}}(V) = \mathcal{O}_X(V)\) exactly when some \(U_i\)
  contains \(V\) (\(I_1 \neq \varnothing\)), recorded in
  Definition~\ref{mk-def:cover_structure_presheaf}.
\end{proof}

\begin{lemma}
[Constant-coefficient combinatorial {\v C}ech contracting-homotopy engine]
  \label{mk-lem:free_cech_engine}
  \lean{AlgebraicGeometry.FreeCechEngine.combDifferential,
    AlgebraicGeometry.FreeCechEngine.combDifferential_comp,
    AlgebraicGeometry.FreeCechEngine.combDifferential_eq_of_cocycle,
    AlgebraicGeometry.FreeCechEngine.combDifferential_exact,
    AlgebraicGeometry.FreeCechEngine.combHomotopy,
    AlgebraicGeometry.FreeCechEngine.combHomotopy_spec,
    AlgebraicGeometry.FreeCechEngine.combHomotopy_zero,
    AlgebraicGeometry.FreeCechEngine.combSign_flip,
    AlgebraicGeometry.FreeCechEngine.cons_comp_succAbove_succ,
    AlgebraicGeometry.isZero_sigma_of_forall_isZero}
  \dcref{custom}

  \textit{Source: Stacks Project, Cohomology, \texttt{lemma-homology-complex}
  (differential and contracting homotopy of the evaluated complex).}
  Fix a nonempty index set \(J\) (in the application \(J = I_1(V)\)), a chosen index
  \(i_{\mathrm{fix}} \in J\), and a single coefficient module \(M\). The
  \emph{constant-coefficient combinatorial {\v C}ech engine} consists of: the
  alternating-sum {\v C}ech differential \(\operatorname{combDifferential}\) on the
  full simplex over \(J\) with coefficients \(M\),
  \[
    (\operatorname{combDifferential} t)_{i_0 \ldots i_{p+1}}
      = \sum_{j=0}^{p+1} (-1)^j\, t_{i_0 \ldots \widehat{i_j} \ldots i_{p+1}},
  \]
  which squares to zero, \(\operatorname{combDifferential} \circ
  \operatorname{combDifferential} = 0\) (\(\operatorname{combDifferential\_comp}\)); the
  \emph{prepend-\(i_{\mathrm{fix}}\)} homotopy \(\operatorname{combHomotopy}\),
  \((\operatorname{combHomotopy}\, i_{\mathrm{fix}}\, u)_{i_0 \ldots i_{p+1}}
  = u_{i_1 \ldots i_{p+1}}\) when \(i_0 = i_{\mathrm{fix}}\) and \(0\) otherwise; and
  the contracting identity \(d \circ h + h \circ d = \mathrm{id}\)
  (\(\operatorname{combHomotopy\_spec}\)), which yields positive-degree exactness
  \(\operatorname{combDifferential\_exact}\) and the cocycle-splitting form
  \(\operatorname{combDifferential\_eq\_of\_cocycle}\). The remaining declarations are
  sign- and index-bookkeeping consumed by the contracting identity: the
  succ-shift reindexing \(\operatorname{cons\_comp\_succAbove\_succ}\) and the
  alternating-sign flip \(\operatorname{combSign\_flip}\), together with the
  degenerate-case simplification \(\operatorname{combHomotopy\_zero}\). The generic
  category-theory helper \(\operatorname{isZero\_sigma\_of\_forall\_isZero}\) (a
  coproduct of zero objects is a zero object) is recorded here as well, used to collapse
  the empty-index direct sums of Lemma~\ref{mk-lem:cech_free_eval_empty}. This engine is
  self-contained combinatorics on \(\operatorname{Fin.cons}\),
  \(\operatorname{Fin.succAbove}\), \(\operatorname{Fin.predAbove}\), and
  \(\operatorname{Finset.sum\_involution}\); it mirrors the contracting homotopy of the
  Stacks \texttt{lemma-homology-complex} proof.
\end{lemma}

\begin{lemma}
[The constant-coefficient {\v C}ech engine \emph{chain} complex \(C_\bullet\)]
  \label{mk-lem:cech_engine_complex}
  \lean{AlgebraicGeometry.le_coverInterOpen_iff,
    AlgebraicGeometry.survivingEquiv,
    AlgebraicGeometry.cechFreeEvalDropZeros,
    AlgebraicGeometry.cechFreeEvalEngine_X,
    AlgebraicGeometry.coverSectionModule,
    AlgebraicGeometry.cechEngineX,
    AlgebraicGeometry.cechEngineD,
    AlgebraicGeometry.cechEngineD_ι,
    AlgebraicGeometry.cechEngineD_comp,
    AlgebraicGeometry.cechEngineComplex,
    AlgebraicGeometry.cechEnginePrepend,
    AlgebraicGeometry.cechEnginePrepend_ι,
    AlgebraicGeometry.cechEnginePrepend_spec,
    AlgebraicGeometry.cechEngineD_exact,
    AlgebraicGeometry.cechEngineAug0,
    AlgebraicGeometry.cechEngineAug0_split,
    AlgebraicGeometry.cechEngineAug0_ι,
    AlgebraicGeometry.cechEngineComplexAug,
    AlgebraicGeometry.cechEngineComplexAug_f_zero,
    AlgebraicGeometry.cechEngineComplex_exactAt,
    AlgebraicGeometry.cechEngineD_comp_aug,
    AlgebraicGeometry.le_coverInterOpen_iffFam,
    AlgebraicGeometry.survivingEquivFam,
    AlgebraicGeometry.cechFreeEvalDropZerosFam,
    AlgebraicGeometry.cechFreeEvalEngine_XFam,
    AlgebraicGeometry.cechEngineXFam,
    AlgebraicGeometry.cechEngineDFam,
    AlgebraicGeometry.cechEngineD_ιFam,
    AlgebraicGeometry.cechEngineD_compFam,
    AlgebraicGeometry.cechEngineComplexFam,
    AlgebraicGeometry.cechEnginePrependFam,
    AlgebraicGeometry.cechEnginePrepend_ιFam,
    AlgebraicGeometry.cechEnginePrepend_specFam,
    AlgebraicGeometry.cechEngineD_exactFam,
    AlgebraicGeometry.cechEngineAug0Fam,
    AlgebraicGeometry.cechEngineAug0_splitFam,
    AlgebraicGeometry.cechEngineAug0_ιFam,
    AlgebraicGeometry.cechEngineComplexAugFam,
    AlgebraicGeometry.cechEngineComplexAug_f_zeroFam,
    AlgebraicGeometry.cechEngineComplex_exactAtFam,
    AlgebraicGeometry.cechEngineD_comp_augFam}
  \dcref{custom}

  \uses{mk-lem:free_cech_engine}
  \textit{Source: Stacks Project, Cohomology, \texttt{lemma-homology-complex}
  (the evaluated complex \(K^{ext}_\bullet(W)\) and its differential).}
  Fix an open \(V\) with \(I_1(V) = \{\, i : V \leq U_i \,\} \neq \varnothing\), and
  let \(\mathcal{O}_X(V)\) be the single coefficient module
  (\(\operatorname{coverSectionModule}\)). The \emph{engine chain complex}
  \(C_\bullet = \operatorname{cechEngineComplex}\, \mathcal{U}\, V\) is the chain
  complex of \(\mathcal{O}_X(V)\)-modules with
  \[
    C_p \;=\; \coprod_{\sigma : \mathrm{Fin}(p+1) \to I_1(V)} \mathcal{O}_X(V)
    \qquad (\operatorname{cechEngineX}),
  \]
  the coproduct taken in \(\operatorname{ModuleCat}(\mathcal{O}_X(V))\), and with
  differential \(d = \operatorname{cechEngineD}\) given on the injection
  \(\iota_\sigma\) of the index \(\sigma : \mathrm{Fin}(p+1) \to I_1(V)\) by the
  alternating sum of index-drop reindexings,
  \[
    d \circ \iota_\sigma
      \;=\; \sum_{i=0}^{p} (-1)^i\, \iota_{\sigma \circ \operatorname{Fin.succAbove} i}
    \qquad (\operatorname{cechEngineD\_\iota}),
  \]
  i.e.\ the chain transpose of the index-\emph{raise} differential
  \(\operatorname{FreeCechEngine.combDifferential}\) of
  Lemma~\ref{mk-lem:free_cech_engine}. Then:
  \begin{itemize}
    \item \(d \circ d = 0\) (\(\operatorname{cechEngineD\_comp}\)), so the family
      assembles into the chain complex \(C_\bullet = \operatorname{ChainComplex.of}\)
      (\(\operatorname{cechEngineComplex}\)).
    \item The \emph{prepend-\(i_{\mathrm{fix}}\)} map
      \(s \mapsto \iota_{\operatorname{Fin.cons}\, i_{\mathrm{fix}}\, \sigma}\)
      (\(\operatorname{cechEnginePrepend}\), with its injection computation
      \(\operatorname{cechEnginePrepend\_\iota}\)) is a contracting homotopy in
      positive degree: \(s \circ d + d \circ s = \mathrm{id}\)
      (\(\operatorname{cechEnginePrepend\_spec}\)).
    \item Hence \(C_\bullet\) is exact in positive degree,
      \(\operatorname{Function.Exact}(\operatorname{cechEngineD}(n+1),
      \operatorname{cechEngineD} n)\) (\(\operatorname{cechEngineD\_exact}\)).
  \end{itemize}
  Four object-half bricks realise the degreewise identification
  \((\operatorname{eval} V)(K(\mathcal{U})_p) \cong C_p\) consumed by
  Lemma~\ref{mk-lem:cech_free_eval_engine_iso}: evaluation commutes with the coproduct;
  the non-surviving summands (those \(\sigma\) with \(V \not\leq U_\sigma\)) drop to the
  zero object (\(\operatorname{cechFreeEvalDropZeros}\),
  \(\operatorname{cechFreeEvalEngine\_X}\)); and the surviving index type is reindexed
  by the equivalence
  \(\operatorname{survivingEquiv} : (\mathrm{Fin}(p+1) \to I_1(V)) \simeq
  \{\, \sigma \;//\; V \leq \operatorname{coverInterOpen}\, \mathcal{U}\, \sigma \,\}\),
  where the membership criterion
  \(V \leq \operatorname{coverInterOpen}\, \mathcal{U}\, \sigma
  \iff \forall k,\ V \leq U_{\sigma(k)}\) is
  \(\operatorname{le\_coverInterOpen\_iff}\).

  The parallel \texttt{...Fam} engine declarations pinned above are the cover-agnostic
  raw-finite-family mirror of this construction (an arbitrary finite family of opens,
  no covering hypothesis), feeding the family form of the {\v C}ech acyclicity lemma.
\end{lemma}

\begin{proof}
  \uses{mk-lem:free_cech_engine}
  The complex \(C_\bullet\) is the chain/coproduct dual of the constant-coefficient
  core of Lemma~\ref{mk-lem:free_cech_engine}. The differential
  \(\operatorname{cechEngineD}\) is defined on coproduct injections by the
  alternating-sign index-drop \(\sigma \mapsto \sigma \circ \operatorname{Fin.succAbove}
  i\); both \(d \circ d = 0\) and the contracting identity
  \(s \circ d + d \circ s = \mathrm{id}\) are proved by extending along the
  injections with \(\operatorname{Limits.Sigma.hom\_ext}\) and reducing each
  injection-component to the corresponding combinatorial cancellation of
  \(\operatorname{FreeCechEngine.combDifferential\_comp}\) respectively
  \(\operatorname{FreeCechEngine.combHomotopy\_spec}\) (the prepend
  \(\operatorname{cechEnginePrepend}\) is the coproduct dual of
  \(\operatorname{combHomotopy}\)). Positive-degree exactness
  \(\operatorname{cechEngineD\_exact}\) reads off the contracting identity. The four
  object-half bricks are immediate: evaluation preserves \(\coprod\) and the surviving
  summands are picked out by \(\operatorname{survivingEquiv}\) under the criterion
  \(\operatorname{le\_coverInterOpen\_iff}\), the rest being zero objects
  (\(\operatorname{cechFreeEvalDropZeros}\)).
\end{proof}

\begin{lemma}

[Degreewise identification of the evaluated free complex with the engine complex]
  \label{mk-lem:cech_free_eval_engine_iso}
  \lean{AlgebraicGeometry.cechFreeEvalEngineIso,
    AlgebraicGeometry.cechFreeAug_eval_eq,
    AlgebraicGeometry.cechFreeEvalEngineIso_hom_f,
    AlgebraicGeometry.cechFreeEvalEngine_X_inv_hom_ι,
    AlgebraicGeometry.cechFreeEvalEngine_comm,
    AlgebraicGeometry.cechFreeEvalEngine_map_ι,
    AlgebraicGeometry.cechFreeEval_X_ι_inv,
    AlgebraicGeometry.cechFree_d_ι,
    AlgebraicGeometry.freeYonedaAug_app_comp,
    AlgebraicGeometry.freeYonedaEval_iso_of_le_hom_eq_aug,
    AlgebraicGeometry.freeYonedaEval_iso_of_le_natural,
    AlgebraicGeometry.cechFreeEvalEngineIsoFam,
    AlgebraicGeometry.cechFreeAug_eval_eqFam,
    AlgebraicGeometry.cechFreeEvalEngineIso_hom_fFam,
    AlgebraicGeometry.cechFreeEvalEngine_X_inv_hom_ιFam,
    AlgebraicGeometry.cechFreeEvalEngine_commFam,
    AlgebraicGeometry.cechFreeEvalEngine_map_ιFam,
    AlgebraicGeometry.cechFreeEval_X_ι_invFam,
    AlgebraicGeometry.cechFree_d_ιFam}
  \dcref{custom}

  \uses{mk-lem:cech_free_eval_sectionwise, mk-lem:free_cech_engine, mk-lem:cech_engine_complex}
  \textit{Source: Stacks Project, Cohomology, \texttt{lemma-homology-complex}
  (the evaluated complex and its differential).}
  Fix an open \(V\) with \(I_1(V) = \{\, i : V \leq U_i \,\} \neq \varnothing\). There
  is an isomorphism of chain complexes of \(\mathcal{O}_X(V)\)-modules
  \[
    \bigl((\operatorname{eval} V).\operatorname{mapHomologicalComplex}\bigr)
      \bigl(\operatorname{cechFreePresheafComplex} \mathcal{U}\bigr)
    \;\cong\;
    C_\bullet,
    \qquad
    C_p = \coprod_{\sigma : \mathrm{Fin}(p+1) \to I_1(V)} \mathcal{O}_X(V),
  \]
  where the coproduct is taken in \(\operatorname{ModuleCat}(\mathcal{O}_X(V))\) and the
  differential of \(C_\bullet\) is the constant-coefficient {\v C}ech differential
  \(\operatorname{FreeCechEngine.combDifferential}\) on the index set \(J = I_1(V)\)
  (Lemma~\ref{mk-lem:free_cech_engine}). The isomorphism is \emph{degreewise} and
  \emph{intertwines the evaluated alternating-face differential with}
  \(\operatorname{combDifferential}\). It is the chain/coproduct dual of the section-side
  cochain identification of Lemma~\ref{mk-lem:section_cech_homology_exact}; isolating it
  here gives the empty and nonempty cases (Lemmas~\ref{mk-lem:cech_free_eval_empty}
  and~\ref{mk-lem:cech_free_eval_nonempty}) a common, ready degreewise model.

  The parallel \texttt{...Fam} declarations pinned above are the cover-agnostic
  raw-finite-family mirror of this degreewise identification (an arbitrary finite
  family of opens, no covering hypothesis), feeding the family form of the {\v C}ech
  acyclicity lemma.
\end{lemma}

\begin{proof}
  \uses{mk-lem:cech_free_eval_sectionwise, mk-lem:free_cech_engine}
  \emph{Degreewise.} By \(\operatorname{cechFreeEval\_X}\) of
  Lemma~\ref{mk-lem:cech_free_eval_sectionwise} evaluation at \(V\) commutes with the
  coproduct,
  \((\operatorname{eval} V)(K(\mathcal{U})_p) \cong \coprod_{\sigma :
  \mathrm{Fin}(p+1) \to \mathcal{U}.I_0} (\operatorname{eval} V)(\operatorname{free}
  (\mathbf{y}\,U_\sigma))\). Split the index set into the \(\sigma\) landing in
  \(I_1(V)\) and the rest. Each summand with \(\sigma\) valued in \(I_1(V)\) (so
  \(V \leq U_\sigma\)) is \(\cong \mathcal{O}_X(V)\) by
  \(\operatorname{freeYonedaEval\_iso\_of\_le}\); each remaining summand is the zero
  object by \(\operatorname{freeYonedaEval\_isZero\_of\_not\_le}\). Hence the coproduct
  collapses to \(\coprod_{\sigma : \mathrm{Fin}(p+1) \to I_1(V)} \mathcal{O}_X(V) =
  C_p\).

  \emph{Differential match.} The evaluated differential is \((\operatorname{eval} V)\)
  applied to the alternating face sum
  \(\operatorname{AlternatingFaceMapComplex.objD}(\operatorname{cechFreeSimplicial}
  \mathcal{U})\). The face \(\delta_i\) reindexes the \(\sigma\)-summand by
  \(\sigma \mapsto \sigma \circ \operatorname{Fin.succAbove} i\) (the face action of
  \(\operatorname{cechFreeSimplicial}\)); restricted to the
  \(I_1\)-summands this is exactly the index-drop that the engine encodes, and the sign
  \((-1)^i\) is supplied by \(\operatorname{objD}\). The match is verified on the
  coproduct injections by \(\operatorname{Limits.Sigma.hom\_ext}\), reusing the
  naturality of \(\operatorname{Limits.PreservesCoproduct.iso}\); summing with signs
  identifies the evaluated differential with
  \(\operatorname{FreeCechEngine.combDifferential}\). This is the hard content of the
  identification and the chain/coproduct dual of the cochain coface match in
  Lemma~\ref{mk-lem:section_cech_homology_exact}.

  \emph{The three layers commuted past.} The degreewise iso is the \(3\)-fold
  composite
  \[
    (\operatorname{eval} V)(K(\mathcal{U})_p)
      \xrightarrow{\;\operatorname{cechFreeEval\_X}\;}
      \coprod_{\sigma : \mathrm{Fin}(p+1) \to \mathcal{U}.I_0}
        (\operatorname{eval} V)(\operatorname{free}(\mathbf{y}\,U_\sigma))
      \xrightarrow{\;\operatorname{cechFreeEvalDropZeros}\;}
      \coprod_{V \leq U_\sigma} \mathcal{O}_X(V)
      \xrightarrow{\;(\operatorname{Limits.Sigma.whiskerEquiv}\ \ldots)^{-1}\;}
      C_p,
  \]
  whose layers are: (1) \(\operatorname{cechFreeEval\_X}\) = the comparison
  \(\operatorname{Limits.PreservesCoproduct.iso}\) expressing that evaluation at
  \(V\) commutes with the coproduct \(\coprod\); (2) \(\operatorname{cechFreeEvalDropZeros}\),
  which kills the summands with \(V \not\leq U_\sigma\) (these evaluate to zero objects
  by \(\operatorname{freeYonedaEval\_isZero\_of\_not\_le}\)); (3)
  \((\operatorname{Limits.Sigma.whiskerEquiv}\ \operatorname{survivingEquiv}\
  (\operatorname{freeYonedaEval\_iso\_of\_le}\, \ldots))^{-1}\), which reindexes the
  surviving summands by \(\operatorname{survivingEquiv}\) of
  Lemma~\ref{mk-lem:cech_engine_complex} and identifies each evaluated summand with the
  single coefficient \(\mathcal{O}_X(V)\) via \(\operatorname{freeYonedaEval\_iso\_of\_le}\).

  \emph{Naturality past each layer.} The reindex \(\operatorname{survivingEquiv}\) is
  natural in the face map \(\sigma \mapsto \sigma \circ \operatorname{Fin.succAbove} i\):
  a surviving index \(\sigma\) (i.e.\ \(V \leq U_\sigma\), equivalently
  \(V \leq \operatorname{coverInterOpen}\, \mathcal{U}\, \sigma\) by
  \(\operatorname{le\_coverInterOpen\_iff}\)) is carried to a surviving
  \(\sigma \circ \operatorname{Fin.succAbove} i\) (still \(V \leq U_{\sigma \circ
  \operatorname{succAbove}\, i}\), since dropping a coordinate only weakens the meet),
  so the drop-zeros layer (2) and the whisker layer (3) commute with every face. On a
  surviving summand the evaluated transition \(\operatorname{freeYoneda}.\mathrm{map}
  (\operatorname{homOfLE}\, \ldots)\) collapses to the identity under
  \(\operatorname{freeYonedaEval\_iso\_of\_le}\); hence each face \(\delta_i\) matches the
  engine reindexing \(\sigma \mapsto \sigma \circ \operatorname{Fin.succAbove} i\)
  summand-for-summand, and summing with the \(\operatorname{objD}\) signs \((-1)^i\)
  identifies the evaluated alternating-face differential with
  \(\operatorname{cechEngineD}\) (\(= \) the chain transpose of
  \(\operatorname{combDifferential}\)).

  \emph{The genuine difficulty: variance.} The free \emph{chain} differential
  \emph{lowers} \(\mathrm{Fin}\)-arity (\(\sigma \mapsto \sigma \circ
  \operatorname{Fin.succAbove} i\), faces), whereas
  \(\operatorname{FreeCechEngine.combDifferential}\) \emph{raises} it (it inserts an
  index, summing over \(i \in I_1\)). The two are reconciled because on the finite
  biproduct \(\coprod = \prod\) the insertion/deletion adjunction makes the index-drop
  chain differential the transpose of the index-raise cochain differential; this
  transpose is precisely what \(\operatorname{cechEngineD}\) of
  Lemma~\ref{mk-lem:cech_engine_complex} packages, and verifying the commutation amounts to
  the injectionwise check \(\operatorname{Limits.Sigma.hom\_ext}\) outlined above.
\end{proof}

\begin{lemma}

[Evaluated free complex over an open meeting no cover member]
  \label{mk-lem:cech_free_eval_empty}
  \lean{AlgebraicGeometry.cechFreeEval_quasiIso_of_isEmpty,
    AlgebraicGeometry.cechFreeEval_isZero_of_isEmpty,
    AlgebraicGeometry.coverStructurePresheaf_eval_isZero_of_isEmpty,
    AlgebraicGeometry.isZero_homology_of_isZero_X,
    AlgebraicGeometry.cechFreeEval_quasiIso_of_isEmptyFam,
    AlgebraicGeometry.cechFreeEval_isZero_of_isEmptyFam,
    AlgebraicGeometry.coverStructurePresheaf_eval_isZero_of_isEmptyFam}
  \dcref{custom}

  \uses{mk-lem:cech_free_eval_sectionwise, mk-lem:free_cech_engine}
  \textit{Source: Stacks Project, Cohomology, \texttt{lemma-homology-complex}
  (empty case).}
  With the notation of Lemma~\ref{mk-lem:cech_free_eval_sectionwise}, suppose
  \(I_1(V) = \varnothing\), i.e.\ \(V\) is contained in no cover member \(U_i\). Then
  the evaluated augmented complex is identically zero: every degree-\(p\) term
  \(\bigoplus_{\sigma : \mathrm{Fin}(p+1) \to I_1} \mathcal{O}_X(V)\) is the empty
  direct sum \(0\) (there is no map \(\mathrm{Fin}(p+1) \to \varnothing\)) and the
  augmentation target is \(\mathcal{O}_{\mathcal{U}}(V) = 0\). Consequently the
  evaluation at \(V\) of \(\operatorname{cechFreeComplexAug}\) is a morphism of zero
  complexes, hence a quasi-isomorphism. The same argument applies to an arbitrary
  finite family of opens, without a covering hypothesis.
\end{lemma}

\begin{proof}
  \uses{mk-lem:cech_free_eval_sectionwise}
  By Lemma~\ref{mk-lem:cech_free_eval_sectionwise} each evaluated term is
  \(\bigoplus_{\sigma : \mathrm{Fin}(p+1) \to I_1} \mathcal{O}_X(V)\); when
  \(I_1 = \varnothing\) the index type \(\mathrm{Fin}(p+1) \to I_1\) is empty, so the
  term is the zero object, and likewise \(\mathcal{O}_{\mathcal{U}}(V) = 0\). A
  morphism whose source and target complexes are both zero has isomorphisms on all
  homology objects and is therefore a quasi-isomorphism.
\end{proof}

\begin{lemma}
[Prepend homotopy on the evaluated free complex, by transport]
  \label{mk-lem:cech_free_eval_prepend_homotopy}
  \dcref{custom}
  \uses{mk-lem:cech_free_eval_sectionwise, mk-lem:cech_free_eval_engine_iso, mk-lem:free_cech_engine, mk-lem:cech_engine_complex}
  \textit{Source: Stacks Project, Cohomology, \texttt{lemma-homology-complex}
  (the contracting map \(h\)).}
  With the notation of Lemma~\ref{mk-lem:cech_free_eval_sectionwise}, suppose
  \(I_1(V) \neq \varnothing\) and fix an index \(i_{\mathrm{fix}} \in I_1\). The
  contracting homotopy is built at the engine-complex level as
  \(\operatorname{cechEnginePrepend}\) of Lemma~\ref{mk-lem:cech_engine_complex}; its
  evaluated-complex avatar is obtained by transporting along the degreewise iso
  \(\operatorname{cechFreeEvalEngineIso}\) of
  Lemma~\ref{mk-lem:cech_free_eval_engine_iso}. Concretely the transported
  \emph{prepend homotopy} is the family of \(\mathcal{O}_X(V)\)-linear maps
  \[
    h_p : K(\mathcal{U})_p(V) \longrightarrow K(\mathcal{U})_{p+1}(V),
    \qquad
    h_p(s)_{i_0 \ldots i_{p+1}} =
      \begin{cases}
        s_{i_1 \ldots i_{p+1}} & i_0 = i_{\mathrm{fix}}, \\
        0 & i_0 \neq i_{\mathrm{fix}},
      \end{cases}
  \]
  defined on the \(I_1\)-indexed direct sum of
  Lemma~\ref{mk-lem:cech_free_eval_sectionwise} as the \(\operatorname{Sigma.desc}\) of
  the per-summand prepend-\(i_{\mathrm{fix}}\) maps. This is exactly the
  combinatorial content of \(\operatorname{FreeCechEngine.combHomotopy}\) of
  Lemma~\ref{mk-lem:free_cech_engine}, here with the single
  coefficient module \(\mathcal{O}_X(V)\), transported across the degreewise
  identification of Lemma~\ref{mk-lem:cech_free_eval_engine_iso}. The maps \(h_p\) (together with the
  degree \((-1)\!\to\!0\) map onto \(\mathcal{O}_{\mathcal{U}}(V) = \mathcal{O}_X(V)\))
  are the homotopy hom maps assembled into a
  \(\operatorname{HomologicalComplex.Homotopy}\) in
  Lemma~\ref{mk-lem:cech_free_eval_nonempty}.
\end{lemma}

\begin{proof}
  \uses{mk-lem:cech_free_eval_sectionwise}
  Under the identification of Lemma~\ref{mk-lem:cech_free_eval_sectionwise} each
  evaluated term is a finite direct sum of copies of \(\mathcal{O}_X(V)\) indexed by
  \(\sigma : \mathrm{Fin}(p+1) \to I_1\). The prepend-\(i_{\mathrm{fix}}\) prescription
  sends the basis component at \((i_1, \ldots, i_{p+1})\) to the component at
  \((i_{\mathrm{fix}}, i_1, \ldots, i_{p+1})\) and is \(\mathcal{O}_X(V)\)-linear; it is
  packaged as the \(\operatorname{Sigma.desc}\) of these per-summand maps, exactly as
  \(\operatorname{FreeCechEngine.combHomotopy}\) is defined. This is a definition,
  so there is nothing to prove beyond well-definedness, which is the linearity of
  \(\operatorname{Sigma.desc}\).
\end{proof}

\begin{lemma}
[The prepend homotopy contracts the augmented evaluated complex, by transport]
  \label{mk-lem:cech_free_eval_prepend_homotopy_spec}
  \dcref{custom}
  \uses{mk-lem:cech_free_eval_prepend_homotopy, mk-lem:free_cech_engine, mk-lem:cech_engine_complex, mk-lem:cech_free_eval_engine_iso}
  \textit{Source: Stacks Project, Cohomology, \texttt{lemma-homology-complex}
  (the contracting identity).}
  In the situation of Lemma~\ref{mk-lem:cech_free_eval_prepend_homotopy}
  (\(I_1(V) \neq \varnothing\), \(i_{\mathrm{fix}} \in I_1\) fixed), the transported
  prepend homotopy \(h\) satisfies the contracting identity
  \[
    d \circ h + h \circ d = \mathrm{id}
  \]
  on the augmented evaluated complex
  \(\cdots \to K(\mathcal{U})_1(V) \to K(\mathcal{U})_0(V) \to
  \mathcal{O}_{\mathcal{U}}(V) \to 0\) (including the degree-\(0\)/degree-\((-1)\)
  augmentation term). Here \(d\) is the alternating-sum {\v C}ech differential of
  Lemma~\ref{mk-lem:cech_free_eval_sectionwise}.
\end{lemma}

\begin{proof}
  \uses{mk-lem:cech_free_eval_prepend_homotopy, mk-lem:free_cech_engine}
  This is the same alternating-sum cancellation as
  \(\operatorname{FreeCechEngine.combHomotopy\_spec}\)
  (\(\operatorname{combDifferential}(\operatorname{combHomotopy} r\, t) +
  \operatorname{combHomotopy} r\, (\operatorname{combDifferential} t) = t\)),
  specialised to coefficients \(\mathcal{O}_X(V)\). Expanding
  \((dh + hd)(s)_{i_0 \ldots i_p}\): the \(j = 0\) term of \(dh\) prepends and then
  removes \(i_{\mathrm{fix}}\), returning \(s_{i_0 \ldots i_p}\); every remaining term
  of \(dh\) cancels in pairs against the corresponding term of \(hd\) (the
  prepend-\(i_{\mathrm{fix}}\) of the differential), because removing an index
  \(j \geq 1\) after prepending \(i_{\mathrm{fix}}\) equals prepending
  \(i_{\mathrm{fix}}\) after removing index \(j-1\), with opposite signs. The net
  result is \(s_{i_0 \ldots i_p}\), i.e.\ \(dh + hd = \mathrm{id}\); the index
  bookkeeping is the \(\operatorname{FreeCechEngine.cons\_comp\_succAbove\_succ}\) /
  \(\operatorname{FreeCechEngine.combSign\_flip}\) content of
  Lemma~\ref{mk-lem:free_cech_engine}.
\end{proof}

\begin{lemma}

[Evaluated free complex over an open meeting a cover member]
  \label{mk-lem:cech_free_eval_nonempty}
  \lean{AlgebraicGeometry.cechFreeEval_quasiIso_of_nonempty,
    AlgebraicGeometry.cechEngineComplexAug_quasiIso,
    AlgebraicGeometry.coverStructurePresheafEval_iso,
    AlgebraicGeometry.coverStructurePresheafEval_iso_hom,
    AlgebraicGeometry.epi_cechEngineAug0,
    AlgebraicGeometry.cechFreeEval_quasiIso_of_nonemptyFam,
    AlgebraicGeometry.cechEngineComplexAug_quasiIsoFam,
    AlgebraicGeometry.coverStructurePresheafEval_isoFam,
    AlgebraicGeometry.coverStructurePresheafEval_iso_homFam,
    AlgebraicGeometry.epi_cechEngineAug0Fam}
  \dcref{custom}

  \uses{mk-lem:cech_free_eval_sectionwise, mk-lem:cech_free_eval_engine_iso,
    mk-lem:cech_free_eval_prepend_homotopy, mk-lem:cech_free_eval_prepend_homotopy_spec,
    mk-lem:cech_engine_complex}
  \textit{Source: Stacks Project, Cohomology, \texttt{lemma-homology-complex}
  (homotopy equivalence to zero).}
  With the notation of Lemma~\ref{mk-lem:cech_free_eval_sectionwise}, suppose
  \(I_1(V) \neq \varnothing\). Then the evaluation at \(V\) of
  \(\operatorname{cechFreeComplexAug}\) is a quasi-isomorphism: the augmented
  evaluated complex is contractible, hence its homology is concentrated in degree
  \(0\) where it equals \(\mathcal{O}_{\mathcal{U}}(V) = \mathcal{O}_X(V)\).

  The parallel \texttt{...Fam} declarations pinned above are the cover-agnostic
  raw-finite-family mirror of this nonempty-index case (an arbitrary finite family of
  opens, no covering hypothesis), feeding the family form of the {\v C}ech
  acyclicity lemma.
\end{lemma}

\begin{proof}
  \uses{mk-lem:cech_free_eval_prepend_homotopy, mk-lem:cech_free_eval_prepend_homotopy_spec}
  By Lemmas~\ref{mk-lem:cech_free_eval_prepend_homotopy}
  and~\ref{mk-lem:cech_free_eval_prepend_homotopy_spec} the degreewise prepend maps
  \(h_p\) satisfy \(dh + hd = \mathrm{id}\) on the augmented evaluated complex, so they
  assemble into a \(\operatorname{HomologicalComplex.Homotopy}\) between the identity
  and the zero endomorphism of that complex. A complex whose identity is homotopic to
  zero is contractible; equivalently the augmented evaluated complex is a
  \(\operatorname{HomotopyEquiv}\) of \(K(\mathcal{U})_\bullet(V)\) with
  \(\mathcal{O}_{\mathcal{U}}(V)[0] = \mathcal{O}_X(V)[0]\). A homotopy equivalence is a
  quasi-isomorphism (\(\operatorname{Homotopy.toQuasiIso}\), or
  \(\operatorname{HomotopyEquiv.toQuasiIso}\) /
  \(\operatorname{QuasiIso.ofHomotopyEquiv}\)), so the evaluated augmentation map is a
  quasi-isomorphism.
\end{proof}

\begin{lemma}

[The free complex resolves \(\mathcal{O}_U\)]
  \label{mk-lem:cech_free_complex_quasi_iso}
  \lean{AlgebraicGeometry.cechFreeComplex_quasiIso,
    AlgebraicGeometry.cechFreeComplex_quasiIsoFam}
  \dcref{custom}

  \uses{mk-def:cech_free_presheaf_complex, mk-def:cover_structure_presheaf,
    mk-lem:quasiIso_of_evaluation, mk-lem:cech_free_eval_sectionwise,
    mk-lem:cech_free_eval_engine_iso, mk-lem:cech_free_eval_empty, mk-lem:cech_free_eval_nonempty}
  \textit{Source: Stacks Project, Cohomology,
  \texttt{lemma-homology-complex}.}
  Let \(X\) be a ringed space and \(\mathcal{U} : U = \bigcup_{i \in I} U_i\) an
  open covering. Let \(\mathcal{O}_{\mathcal{U}}\) be the cover structure presheaf of
  Definition~\ref{mk-def:cover_structure_presheaf}, the image presheaf of the
  augmentation \(\bigoplus_i (j_i)_{!}\,\mathcal{O}_X|_{U_i} \to \mathcal{O}_X\) (so
  \(\mathcal{O}_{\mathcal{U}}(W) = \mathcal{O}_X(W)\) when \(W\) is contained in some
  \(U_i\), and \(0\) otherwise). Then the homology
  presheaves of the free presheaf complex
  \(K(\mathcal{U})_\bullet\) of Definition~\ref{mk-def:cech_free_presheaf_complex}
  vanish away from degree \(0\) and equal \(\mathcal{O}_{\mathcal{U}}\) in degree
  \(0\):
  \[
    H_i\bigl(K(\mathcal{U})_\bullet\bigr) =
    \begin{cases}
      \mathcal{O}_{\mathcal{U}} & i = 0, \\
      0 & i \neq 0.
    \end{cases}
  \]
  Equivalently, the augmented complex \(\cdots \to K(\mathcal{U})_1 \to
  K(\mathcal{U})_0 \to \mathcal{O}_{\mathcal{U}} \to 0\) is exact as a complex of
  presheaves; that is, \(K(\mathcal{U})_\bullet\) is a (projective) resolution of
  \(\mathcal{O}_{\mathcal{U}}[0]\), so the augmentation
  \(K(\mathcal{U})_\bullet \to \mathcal{O}_{\mathcal{U}}[0]\) is a
  quasi-isomorphism — \(K(\mathcal{U})_\bullet\) is quasi-isomorphic to
  \(\mathcal{O}_{\mathcal{U}}\) concentrated in degree \(0\).

  The companion \(\operatorname{cechFreeComplex\_quasiIsoFam}\) is the cover-agnostic
  raw-finite-family mirror of this resolution statement — for an arbitrary finite
  family \((U_i)_{i \in I}\) of opens with \emph{no} covering hypothesis — and is the
  quasi-isomorphism consumed by the family form of the {\v C}ech acyclicity lemma.
\end{lemma}
\begin{proof}
  \uses{mk-lem:quasiIso_of_evaluation, mk-lem:cech_free_eval_sectionwise,
    mk-lem:cech_free_eval_engine_iso, mk-lem:cech_free_eval_empty, mk-lem:cech_free_eval_nonempty}
  The claim is that the augmentation chain map
  \(K(\mathcal{U})_\bullet \to \mathcal{O}_{\mathcal{U}}[0]\)
  (\(\operatorname{cechFreeComplexAug}\)) is a quasi-isomorphism. By the
  objectwise-reduction Lemma~\ref{mk-lem:quasiIso_of_evaluation}
  (\(\operatorname{quasiIso\_of\_evaluation}\)) it suffices to show that for every open
  \(V \subset X\) the evaluation
  \(\bigl((\operatorname{PresheafOfModules.evaluation}\, V).
  \operatorname{mapHomologicalComplex}\bigr)(\operatorname{cechFreeComplexAug})\) is a
  quasi-isomorphism of complexes of \(\mathcal{O}_X(V)\)-modules.

  Fix \(V\) and set \(I_1 = I_1(V) = \{\,i : V \leq U_i\,\}\). By
  Lemma~\ref{mk-lem:cech_free_eval_sectionwise} the evaluated augmented complex is the
  augmented combinatorial {\v C}ech complex of the simplex on \(I_1\) with
  coefficients \(\mathcal{O}_X(V)\), with augmentation target
  \(\mathcal{O}_{\mathcal{U}}(V)\). Split into two cases on \(I_1\). If
  \(I_1 = \varnothing\), Lemma~\ref{mk-lem:cech_free_eval_empty} shows both sides vanish
  and the evaluated augmentation is a quasi-isomorphism. If \(I_1 \neq \varnothing\),
  Lemma~\ref{mk-lem:cech_free_eval_nonempty} provides the prepend-\(i_{\mathrm{fix}}\)
  contracting homotopy, exhibiting the augmented evaluated complex as contractible and
  hence the evaluated augmentation as a quasi-isomorphism. In either case the
  hypothesis of Lemma~\ref{mk-lem:quasiIso_of_evaluation} holds for \(V\), so
  \(\operatorname{cechFreeComplexAug}\) is a quasi-isomorphism; equivalently the
  homology presheaves of \(K(\mathcal{U})_\bullet\) vanish away from degree \(0\) and
  equal \(\mathcal{O}_{\mathcal{U}}\) in degree \(0\).

  The contracting homotopy is built directly at the \(\mathcal{O}_X(V)\)-module
  (section) level as the prepend-\(i_{\mathrm{fix}}\) map — the same combinatorial
  content as \(\operatorname{FreeCechEngine.combHomotopy}\) /
  \(\operatorname{FreeCechEngine.combHomotopy\_spec}\) of
  Lemma~\ref{mk-lem:free_cech_engine}.
\end{proof}

\begin{lemma}[Right adjoint of a mono-preserving left adjoint preserves injectives]
  \label{mk-lem:injective_of_adjoint}
  \lean{CategoryTheory.Injective.injective_of_adjoint}
  \mathlibok
  \dcref{custom}

  Let \(L \dashv R\) be an adjunction between abelian (or, more generally, suitable)
  categories such that the left adjoint \(L\) preserves monomorphisms. Then for every
  injective object \(J\) of the target, \(R(J)\) is an injective object of the source.
\end{lemma}

\begin{lemma}[Sheafification is left adjoint to the inclusion of sheaves of modules]
  \label{mk-lem:mod_pmod_adjunction}
  \lean{PresheafOfModules.sheafificationAdjunction}
  \mathlibok
  \dcref{custom}

  For the identity ring map \(\alpha = \mathbf{1}\) over a scheme \(X\),
  sheafification is left adjoint to the inclusion
  \(\mathrm{Mod}(\mathcal{O}_X) \hookrightarrow \mathrm{PMod}(\mathcal{O}_X)\) of
  sheaves of \(\mathcal{O}_X\)-modules into presheaves of \(\mathcal{O}_X\)-modules:
  \[
    \operatorname{sheafification} \;\dashv\; \iota,
    \qquad
    \iota = \text{the inclusion } X.\mathrm{Modules} \hookrightarrow
    X.\mathrm{PresheafOfModules}.
  \]
  The left adjoint sheafification is exact, in particular it preserves
  monomorphisms.
\end{lemma}

\begin{lemma}
[Opposite-transport identification of the hom-{\v C}ech complex]
  \label{mk-lem:cech_complex_op_identification}
  \lean{AlgebraicGeometry.homCechComplexMapOpIso,
    AlgebraicGeometry.homCechComplex_d_eq,
    AlgebraicGeometry.homCechCosimplicial_δ,
    AlgebraicGeometry.homCechComplexMapOpIsoFam,
    AlgebraicGeometry.homCechComplex_d_eqFam}
  \dcref{custom}

  \uses{mk-lem:cech_complex_hom_identification, mk-def:cech_free_presheaf_complex}
  Let \(\mathcal{U} : U = \bigcup_{i \in I} U_i\) be a finite open covering and
  \(\mathcal{F}\) a presheaf of \(\mathcal{O}_X\)-modules. There is an isomorphism of
  cochain complexes of abelian groups
  \[
    \operatorname{homCechComplex} \mathcal{U}\, \mathcal{F}
    \;\cong\;
    \bigl((\operatorname{preadditiveYoneda}.\mathrm{obj}\,\mathcal{F}).
      \operatorname{mapHomologicalComplex}(\mathrm{up}\,\mathbb{N})\bigr)
      .\mathrm{obj}\bigl(\operatorname{HomologicalComplex.op}
        (\operatorname{cechFreePresheafComplex} \mathcal{U})\bigr),
  \]
  exhibiting the {\v C}ech hom-complex
  \(\operatorname{Hom}_{\mathcal{O}_X}(K(\mathcal{U})_\bullet, \mathcal{F})\) of
  Lemma~\ref{mk-lem:cech_complex_hom_identification} as the image of the
  \emph{opposite} of the free presheaf complex
  (Definition~\ref{mk-def:cech_free_presheaf_complex}) under the covariant
  \(\operatorname{Hom}(-, \mathcal{F})\) functor
  \(\operatorname{preadditiveYoneda}.\mathrm{obj}\,\mathcal{F}\).
\end{lemma}

\begin{lemma}

[Opposite-transport identification of the section {\v C}ech complex]
  \label{mk-lem:section_cech_complex_mapop_iso}
  \lean{AlgebraicGeometry.sectionCechComplexMapOpIso,
    AlgebraicGeometry.sectionCechComplexMapOpIsoFam}
  \dcref{custom}

  \uses{mk-lem:cech_complex_op_identification, mk-lem:cech_complex_hom_identification}
  With the notation above, composing the opposite-transport iso
  \(\operatorname{homCechComplexMapOpIso}\) of
  Lemma~\ref{mk-lem:cech_complex_op_identification} with the {\v C}ech hom-identification
  \(\operatorname{cechComplex\_hom\_identification}\) of
  Lemma~\ref{mk-lem:cech_complex_hom_identification} gives an isomorphism of cochain
  complexes
  \[
    \bigl((\operatorname{preadditiveYoneda}.\mathrm{obj}\,\mathcal{F}).
      \operatorname{mapHomologicalComplex}\bigr).\mathrm{obj}
      \bigl(\operatorname{HomologicalComplex.op}
        (\operatorname{cechFreePresheafComplex} \mathcal{U})\bigr)
    \;\cong\;
    \operatorname{sectionCechComplex}(\operatorname{coverOpen} \mathcal{U})\,
      \mathcal{F},
  \]
  i.e.\ \(\operatorname{sectionCechComplexMapOpIso}\) is
  \(\operatorname{homCechComplexMapOpIso}^{-1}\) followed by
  \(\operatorname{cechComplex\_hom\_identification}\). This is the bridge that lets the
  mapped-opposite of the free resolution be transported onto the section {\v C}ech
  complex of Definition~\ref{mk-def:section_cech_complex}.

  The companion \(\operatorname{sectionCechComplexMapOpIsoFam}\) is the cover-agnostic
  raw-finite-family mirror of this transport, stated over an arbitrary finite family
  \((U_i)_{i \in \iota}\) of opens with \emph{no} covering hypothesis.
\end{lemma}

\begin{lemma}
[Hom into an injective preserves exactness and quasi-isomorphisms]
  \label{mk-lem:hom_into_injective_exact}
  \lean{AlgebraicGeometry.preadditiveYoneda_obj_preservesFiniteColimits_of_injective,
    AlgebraicGeometry.quasiIso_map_preadditiveYoneda_of_injective}
  \dcref{custom}

  \uses{mk-lem:injective_of_adjoint}
  Let \(\mathcal{A}\) be an abelian category and \(I\) an injective object. Then the
  contravariant hom functor \(\operatorname{Hom}(-, I) : \mathcal{A}^{\mathrm{op}} \to
  \mathrm{Ab}\) is exact: it is automatically left exact (preserves finite limits), and
  injectivity of \(I\) makes it carry epimorphisms to epimorphisms, so it preserves
  finite colimits and hence preserves homology. Consequently \(\operatorname{Hom}(-, I)\)
  carries quasi-isomorphisms of chain complexes in \(\mathcal{A}^{\mathrm{op}}\) to
  quasi-isomorphisms of the \(\operatorname{Hom}\)-cochain complexes. Both statements
  are for a general abelian category, the right generality for reuse in
  Lemma~\ref{mk-lem:injective_cech_acyclic}. The mechanism --- a right adjoint / exact
  contravariant functor preserves the relevant exactness --- is the same homological
  principle behind Lemma~\ref{mk-lem:injective_of_adjoint}.
\end{lemma}

\begin{proof}
  \uses{mk-lem:injective_of_adjoint}
  For injective \(I\), injectivity gives that \(\operatorname{Hom}(-, I)\) carries
  epimorphisms to epimorphisms. Combined with the automatic preservation of finite limits
  (left exactness), the short-exact-sequence criterion shows it preserves finite colimits;
  hence, being also left exact, it preserves homology. A functor that preserves homology
  sends quasi-isomorphisms of chain complexes to quasi-isomorphisms degreewise, which is
  the second statement.
\end{proof}

\begin{lemma}
[Injective modules are {\v C}ech-acyclic]
  \label{mk-lem:injective_cech_acyclic}
  \lean{AlgebraicGeometry.injective_cech_acyclic,
    AlgebraicGeometry.injective_cech_acyclicFam,
    AlgebraicGeometry.injective_toPresheafOfModules}
  \dcref{01XB,custom}

  \uses{mk-def:cech_complex, mk-def:section_cech_complex, mk-lem:cech_complex_hom_identification, mk-lem:cech_free_complex_quasi_iso, mk-lem:injective_of_adjoint, mk-lem:mod_pmod_adjunction, mk-lem:cech_complex_op_identification, mk-lem:section_cech_complex_mapop_iso, mk-lem:hom_into_injective_exact}
  \textit{Source: Stacks Project, Cohomology,
  \texttt{lemma-injective-trivial-cech}.}
  Let \(X\) be a ringed space, let \(\mathcal{U} : U =
  \bigcup_{i \in I} U_i\) be any open covering, and let \(\mathcal{I}\) be an
  \emph{injective} \(\mathcal{O}_X\)-module. Then the {\v C}ech cohomology of
  \(\mathcal{I}\) is concentrated in degree zero:
  \[
    \check{H}^p(\mathcal{U}, \mathcal{I}) =
    \begin{cases}
      \mathcal{I}(U) & p = 0, \\
      0 & p > 0.
    \end{cases}
  \]
  The positive-degree vanishing holds over an \emph{arbitrary finite family of opens}
  \(\{U_i\}_{i \in \iota}\), with no covering hypothesis: the augmented free {\v C}ech
  resolution is exact \emph{stalkwise} — at a point \(x\) it is the augmented simplicial
  chain complex of the full simplex on \(\{i \mid x \in U_i\}\) with coefficients
  \(\mathcal{O}_{X,x}\), which is contractible — and a covering hypothesis only fixes the
  identity of the augmentation target \(\mathcal{O}_{\mathcal{U}}\), never the exactness.
  Consequently the lemma applies to any standard covering data of the affine cover system,
  with no restriction to coverings of the whole space.

  This cover-agnostic form is recorded as \(\operatorname{injective\_cech\_acyclicFam}\),
  parameterized by a raw finite family \(\{U_i\}_{i \in \iota}\) of opens
  (\(\iota\) finite) carrying \emph{no} covering hypothesis; it consumes the family-form
  resolution quasi-isomorphism \(\operatorname{cechFreeComplex\_quasiIsoFam}\) and is the
  form invoked over the standard covers of an arbitrary distinguished open \(D(f)\) by the
  cover-agnostic affine vanishing statement (Tag 01XB).
\end{lemma}
\begin{proof}
  \uses{mk-def:cech_complex, mk-def:section_cech_complex, mk-lem:cech_complex_hom_identification, mk-lem:cech_free_complex_quasi_iso, mk-lem:injective_of_adjoint, mk-lem:mod_pmod_adjunction, mk-lem:cech_complex_op_identification, mk-lem:section_cech_complex_mapop_iso, mk-lem:hom_into_injective_exact}
  The argument has two independent parts.

  \emph{Part 1: an injective sheaf is injective as a presheaf.} Sheafification is
  left adjoint to the inclusion
  \(\iota : \mathrm{Mod}(\mathcal{O}_X) \hookrightarrow \mathrm{PMod}(\mathcal{O}_X)\)
  of sheaves of \(\mathcal{O}_X\)-modules into presheaves of
  \(\mathcal{O}_X\)-modules (Lemma~\ref{mk-lem:mod_pmod_adjunction}), and its left
  adjoint sheafification is exact, hence preserves monomorphisms. By
  Lemma~\ref{mk-lem:injective_of_adjoint} the right adjoint \(\iota\) carries injective
  objects to injective objects; thus an injective \(\mathcal{O}_X\)-module
  \(\mathcal{I}\) is injective as an object of \(\mathrm{PMod}(\mathcal{O}_X)\). In
  particular the contravariant functor \(\operatorname{Hom}_{\mathcal{O}_X}(-,
  \mathcal{I})\) is exact on \(\mathrm{PMod}(\mathcal{O}_X)\).

  \emph{Part 2: vanishing.} By Lemma~\ref{mk-lem:cech_free_complex_quasi_iso} the free
  presheaf complex \(K(\mathcal{U})_\bullet\) of
  Definition~\ref{mk-def:cech_free_presheaf_complex} is a resolution of
  \(\mathcal{O}_{\mathcal{U}}\) concentrated in degree \(0\); equivalently the
  augmented complex \(\cdots \to K(\mathcal{U})_1 \to K(\mathcal{U})_0 \to
  \mathcal{O}_{\mathcal{U}} \to 0\) is exact in \(\mathrm{PMod}(\mathcal{O}_X)\).
  Applying the exact functor \(\operatorname{Hom}_{\mathcal{O}_X}(-, \mathcal{I})\)
  of Part 1 to this exact augmented complex yields an exact complex; by the
  identification \(\operatorname{Hom}_{\mathcal{O}_X}(K(\mathcal{U})_\bullet,
  \mathcal{I}) = \check{\mathcal{C}}^\bullet(\mathcal{U}, \mathcal{I})\) of
  Lemma~\ref{mk-lem:cech_complex_hom_identification} (with the section {\v C}ech complex
  of Definition~\ref{mk-def:section_cech_complex}), this is precisely the statement that
  the section {\v C}ech complex of \(\mathcal{I}\) is exact in positive degrees, with
  \(\operatorname{Hom}_{\mathcal{O}_X}(\mathcal{O}_{\mathcal{U}}, \mathcal{I}) =
  \mathcal{I}(U)\) supplying degree \(0\).

  \emph{Categorical assembly.} Concretely, the exactness is produced by the
  opposite-transport bridge. The free {\v C}ech augmentation is a quasi-isomorphism
  (Lemma~\ref{mk-lem:cech_free_complex_quasi_iso}); forming its opposite complex and
  applying \(\operatorname{Hom}(-, \mathcal{I})\) keeps it a quasi-isomorphism, since
  \(\mathcal{I}\) is injective as a presheaf of modules (Part 1) and
  \(\operatorname{Hom}(-, \mathcal{I})\) then preserves homology
  (Lemma~\ref{mk-lem:hom_into_injective_exact}). Transporting source and target of this
  opposite quasi-isomorphism across the opposite-transport isomorphisms
  (Lemma~\ref{mk-lem:cech_complex_op_identification} and
  Lemma~\ref{mk-lem:section_cech_complex_mapop_iso}) identifies it with the augmentation
  of the section {\v C}ech complex \(\check{\mathcal{C}}^\bullet(\mathcal{U},
  \mathcal{I})\) onto the degree-\(0\)-concentrated target \(\mathcal{O}_{\mathcal{U}}[0]\).
  The latter has vanishing homology in positive degrees, so the section {\v C}ech complex
  does too. Hence \(\check{H}^p(\mathcal{U}, \mathcal{I}) = 0\) for all \(p > 0\) and
  \(\check{H}^0(\mathcal{U}, \mathcal{I}) = \mathcal{I}(U)\). This direct route uses
  neither enough-injectives in \(\mathrm{PMod}(\mathcal{O}_X)\) nor a right-derived
  \(\delta\)-functor identification.
\end{proof}

\begin{lemma}

[Surjectivity on sections from cofinal {\v C}ech-$H^1$ vanishing]
  \label{mk-lem:ses_cech_h1}
  \lean{AlgebraicGeometry.ses_cech_h1,
    AlgebraicGeometry.sectionCech_objD_exact_of_isZero_homology,
    AlgebraicGeometry.sectionCech_one_coboundary_of_isZero_homology,
    AlgebraicGeometry.restr_trans,
    AlgebraicGeometry.restr_inj_of_eq,
    AlgebraicGeometry.restr_op_unique,
    AlgebraicGeometry.restr_g'_transport,
    AlgebraicGeometry.fι_sectionCechFaceRestr,
    AlgebraicGeometry.coverConst_iInf,
    AlgebraicGeometry.coverPair_iInf,
    AlgebraicGeometry.pair_comp_δ0,
    AlgebraicGeometry.pair_comp_δ1}
  \dcref{custom}

  \uses{mk-def:cech_complex}
  \textit{Source: Stacks Project, Cohomology,
  \texttt{lemma-ses-cech-h1}.}
  Let \(X\) be a ringed space and let
  \[
    0 \to \mathcal{F} \to \mathcal{G} \to \mathcal{H} \to 0
  \]
  be a short exact sequence of \(\mathcal{O}_X\)-modules. Let \(U \subset X\) be
  open. If there is a cofinal system of open coverings \(\mathcal{U}\) of \(U\) with
  \(\check{H}^1(\mathcal{U}, \mathcal{F}) = 0\), then the map on sections
  \(\mathcal{G}(U) \to \mathcal{H}(U)\) is surjective. Equivalently, the short exact
  sequence is exact on sections over such \(U\).
\end{lemma}
\begin{proof}
  \uses{mk-def:cech_complex}
  Let \(s \in \mathcal{H}(U)\). Choose an open covering \(\mathcal{U} : U =
  \bigcup_{i \in I} U_i\) fine enough that (a) \(\check{H}^1(\mathcal{U},
  \mathcal{F}) = 0\) and (b) each restriction \(s|_{U_i}\) lifts to a section
  \(s_i \in \mathcal{G}(U_i)\); both can be arranged simultaneously because the
  coverings with property (a) are cofinal and (b) holds on any sufficiently fine
  cover (surjectivity of \(\mathcal{G} \to \mathcal{H}\) is a local condition).
  On the double overlaps form
  \[
    s_{i_0 i_1} = s_{i_1}|_{U_{i_0 i_1}} - s_{i_0}|_{U_{i_0 i_1}}.
  \]
  Because the \(s_i\) all lift the single section \(s\), the difference
  \(s_{i_0 i_1}\) maps to \(0\) in \(\mathcal{H}\), hence lies in the subsheaf
  \(\mathcal{F}(U_{i_0 i_1})\); the family \((s_{i_0 i_1})\) is a {\v C}ech
  \(1\)-cocycle for \(\mathcal{F}\). By \(\check{H}^1(\mathcal{U}, \mathcal{F}) = 0\)
  it is a coboundary: there are \(t_i \in \mathcal{F}(U_i)\) with
  \(s_{i_0 i_1} = t_{i_1}|_{U_{i_0 i_1}} - t_{i_0}|_{U_{i_0 i_1}}\). Then the
  corrected sections \(s_i - t_i \in \mathcal{G}(U_i)\) agree on overlaps, so by the
  sheaf condition they glue to a global section of \(\mathcal{G}\) over \(U\) whose
  image in \(\mathcal{H}\) is \(s\). Hence \(\mathcal{G}(U) \to \mathcal{H}(U)\) is
  surjective.
\end{proof}

\section{Absolute sheaf cohomology as Ext of the corepresenting object}

The two vanishing lemmas below speak about the \emph{absolute} sheaf cohomology
\(H^p(U, \mathcal{F})\) of an \(\mathcal{O}_X\)-module over an open \(U\). Their
dimension-shift proof needs three structural facts about \(H^p(U, -)\): a degree-zero
identification with global sections, vanishing on injective modules, and a long exact
sequence attached to a short exact sequence of modules. We realize \(H^p(U, -)\) as the
\emph{Ext of a corepresenting object} in the single abelian category
\(X.\mathrm{Modules} = \mathrm{Mod}(\mathcal{O}_X)\): there is an
\(\mathcal{O}_X\)-module \(j_!\mathcal{O}_U\) corepresenting the
sections-over-\(U\) functor \(\mathcal{F} \mapsto \Gamma(U, \mathcal{F})\), and we set
\(H^p(U, \mathcal{F}) = \operatorname{Ext}^p_{\mathrm{Mod}(\mathcal{O}_X)}(j_!\mathcal{O}_U,
\mathcal{F})\). The corepresenting object is built \emph{without} the extension-by-zero
functor — it is the sheafification of the free module on the representable presheaf
\(\mathbf{y}\,U\) — and the realization keeps the injective module in the \emph{second}
Ext argument, where Mathlib supplies all three facts off the shelf with \emph{no} appeal
to restriction of module sheaves. The corepresenting object and its corepresentability
are fixed first, then the following definition fixes the realization, and the Mathlib
dependency anchors after it pin the exact results it rests on.

\begin{definition}[Extension-by-zero structure sheaf as a corepresenting object]\label{mk-def:jshriek_ou}
  \lean{AlgebraicGeometry.jShriekOU}
  \leanok
  \dcref{custom}

  \uses{mk-def:cech_free_presheaf_complex, mk-lem:mod_pmod_adjunction,
    mk-lem:cech_complex_hom_identification}
  \textit{Corepresenting object for sections over an open, built without the
  extension-by-zero functor.}
  Let \(X\) be a scheme and \(U \subseteq X\) an open subscheme, regarded through the
  representable presheaf \(\mathbf{y}\,U\) on the site of opens of \(X\). The
  \emph{extension-by-zero structure sheaf} \(j_!\mathcal{O}_U\) is the object of
  \(X.\mathrm{Modules} = \mathrm{Mod}(\mathcal{O}_X)\) obtained by sheafifying the free
  presheaf of modules on \(\mathbf{y}\,U\):
  \[
    j_!\mathcal{O}_U \;:=\; \operatorname{sheafification}\bigl(
      \operatorname{free}(\mathbf{y}\,U)\bigr)
    \;\in\; X.\mathrm{Modules},
  \]
  where \(\operatorname{free}\) is the free presheaf-of-modules functor on the
  representable (the degree-zero generator of the free presheaf complex of
  Definition~\ref{mk-def:cech_free_presheaf_complex}) and
  \(\operatorname{sheafification} \dashv \iota\) is the sheafification adjunction of
  Lemma~\ref{mk-lem:mod_pmod_adjunction}. Up to isomorphism this is the classical
  extension-by-zero \(j_!\mathcal{O}_U\) of the structure sheaf of \(U\) along the open
  immersion \(j : U \hookrightarrow X\); crucially, we build only this \emph{object} as a
  sheafified free module, never the extension-by-zero \emph{functor} \(j_!\). The object
  exists unconditionally: the free presheaf of modules on a representable and the
  sheafification of presheaves of \(\mathcal{O}_X\)-modules are both supplied by Mathlib.
\end{definition}

\begin{lemma}

[Corepresentability of global sections over \(U\)]
  \label{mk-lem:jshriek_corepr}
  \lean{AlgebraicGeometry.jShriekOU_homEquiv, AlgebraicGeometry.sheafificationHomAddEquiv}
  \dcref{custom}

  \uses{mk-def:jshriek_ou, mk-lem:mod_pmod_adjunction, mk-lem:cech_complex_hom_identification}
  \textit{The object \(j_!\mathcal{O}_U\) corepresents the sections-over-\(U\) functor.}
  For every \(\mathcal{O}_X\)-module \(\mathcal{F}\) there is a natural additive
  isomorphism
  \[
    \operatorname{Hom}_{X.\mathrm{Modules}}\bigl(j_!\mathcal{O}_U,\, \mathcal{F}\bigr)
    \;\cong\;
    \Gamma(U, \mathcal{F}) \;=\; \mathcal{F}(U),
  \]
  natural in \(\mathcal{F}\). Thus \(j_!\mathcal{O}_U\) corepresents the global-sections
  functor \(\mathcal{F} \mapsto \Gamma(U, \mathcal{F})\) on \(X.\mathrm{Modules}\).
\end{lemma}

\begin{proof}
  \uses{mk-def:jshriek_ou, mk-lem:mod_pmod_adjunction, mk-lem:cech_complex_hom_identification}
  Compose two natural additive isomorphisms. The sheafification adjunction
  \(\operatorname{sheafification} \dashv \iota\) of
  Lemma~\ref{mk-lem:mod_pmod_adjunction} transports corepresentability through
  sheafification:
  \[
    \operatorname{Hom}_{X.\mathrm{Modules}}\bigl(
      \operatorname{sheafification}(\operatorname{free}(\mathbf{y}\,U)),\, \mathcal{F}
    \bigr)
    \;\cong\;
    \operatorname{Hom}_{\mathrm{PMod}(\mathcal{O}_X)}\bigl(
      \operatorname{free}(\mathbf{y}\,U),\, \iota\mathcal{F}\bigr).
  \]
  The free–Yoneda additive equivalence \(\operatorname{freeYonedaHomAddEquiv}\) of
  Lemma~\ref{mk-lem:cech_complex_hom_identification} then identifies the right-hand
  presheaf Hom-group with the evaluation \((\iota\mathcal{F})(U) = \mathcal{F}(U) =
  \Gamma(U, \mathcal{F})\):
  \[
    \operatorname{Hom}_{\mathrm{PMod}(\mathcal{O}_X)}\bigl(
      \operatorname{free}(\mathbf{y}\,U),\, \iota\mathcal{F}\bigr)
    \;\cong\;
    \mathcal{F}(U).
  \]
  Both isomorphisms are additive and natural in \(\mathcal{F}\), so their composite is
  the asserted natural additive isomorphism.
\end{proof}

\begin{definition}[Absolute sheaf cohomology as Ext of the corepresenting object]\label{mk-def:absolute_cohomology}
  \lean{AlgebraicGeometry.absoluteCohomology}
  \leanok
  \dcref{custom}

  \uses{mk-def:jshriek_ou, mk-lem:jshriek_corepr, mk-lem:ext_bifunctor_mathlib,
    mk-lem:ext_covariant_les_mathlib, mk-lem:ext_eq_zero_of_injective_mathlib,
    mk-lem:ext_homequiv_zero_mathlib, mk-lem:hasext_standard_mathlib}
  \textit{Realization of absolute sheaf cohomology in the module-sheaf category.}
  Let \(X\) be a scheme, \(\mathcal{F}\) an \(\mathcal{O}_X\)-module, and \(U \subseteq
  X\) an open subscheme with open immersion \(j : U \hookrightarrow X\). Let
  \(j_!\mathcal{O}_U \in X.\mathrm{Modules}\) be the corepresenting object of
  Definition~\ref{mk-def:jshriek_ou}. The \emph{absolute sheaf cohomology of \(\mathcal{F}\)
  over \(U\)} is defined to be the Ext of this corepresenting object in the single
  abelian category \(X.\mathrm{Modules} = \mathrm{Mod}(\mathcal{O}_X)\) of sheaves of
  \(\mathcal{O}_X\)-modules:
  \[
    H^p(U, \mathcal{F}) \;:=\; \operatorname{Ext}^p_{\mathrm{Mod}(\mathcal{O}_X)}
      \bigl(j_!\mathcal{O}_U,\; \mathcal{F}\bigr),
    \qquad p \geq 0.
  \]
  Both Ext arguments live in \(X.\mathrm{Modules}\); no restriction of \(\mathcal{F}\)
  along \(j\) is taken. Here \(\operatorname{Ext}\) is Mathlib's Ext bifunctor
  (Lemma~\ref{mk-lem:ext_bifunctor_mathlib}) on \(X.\mathrm{Modules}\); the
  \(\mathrm{HasExt}\) structure it requires is available unconditionally by
  Lemma~\ref{mk-lem:hasext_standard_mathlib}. The three structural facts the downstream
  dimension-shift arguments consume are then exactly the corresponding Mathlib results:
  \begin{enumerate}
  \item \textbf{Degree zero is global sections.} In degree \(0\), Ext is the Hom-group,
    \(\operatorname{Ext}^0(j_!\mathcal{O}_U, \mathcal{F}) \cong
    \operatorname{Hom}_{X.\mathrm{Modules}}(j_!\mathcal{O}_U, \mathcal{F})\), by the
    degree-zero Ext comparison (Lemma~\ref{mk-lem:ext_homequiv_zero_mathlib}); composing
    with the corepresentability isomorphism
    \(\operatorname{Hom}_{X.\mathrm{Modules}}(j_!\mathcal{O}_U, \mathcal{F}) \cong
    \Gamma(U, \mathcal{F})\) of Lemma~\ref{mk-lem:jshriek_corepr} gives \(H^0(U,
    \mathcal{F}) \cong \Gamma(U, \mathcal{F})\).
  \item \textbf{Injective vanishing.} When the \emph{second} Ext argument is an
    injective \(\mathcal{O}_X\)-module \(\mathcal{I}\), one has
    \(\operatorname{Ext}^{n+1}(-, \mathcal{I}) = 0\) by
    Lemma~\ref{mk-lem:ext_eq_zero_of_injective_mathlib}. Since \(\mathcal{I}\) is exactly
    the second argument of \(\operatorname{Ext}^{n+1}(j_!\mathcal{O}_U, \mathcal{I})\),
    this yields \(H^{n+1}(U, \mathcal{I}) = 0\) directly: the short exact sequence stays
    in \(X.\mathrm{Modules}\), no restriction is taken, and
    restriction-preserves-injectives is not needed.
  \item \textbf{Long exact sequence.} For a short exact sequence \(0 \to \mathcal{F}
    \to \mathcal{I} \to \mathcal{Q} \to 0\) in \(X.\mathrm{Modules}\) and the
    \emph{fixed} first argument \(j_!\mathcal{O}_U\), the covariant Ext long exact
    sequence (Lemma~\ref{mk-lem:ext_covariant_les_mathlib}) is the long exact sequence of
    the groups \(H^p(U, -)\), with connecting maps given by composition with the Ext
    class of the extension.
  \end{enumerate}
\end{definition}

\begin{definition}
[Ext bifunctor in an abelian category]
  \label{mk-lem:ext_bifunctor_mathlib}
  \lean{CategoryTheory.Abelian.Ext}
  \mathlibok
  \dcref{custom}

  For an abelian category \(C\) with a \(\mathrm{HasExt}\) structure and objects \(X, Y
  \in C\) and \(n \in \mathbb{N}\), Mathlib provides the Ext group
  \(\operatorname{Ext}^n_C(X, Y)\), functorial (contravariantly in \(X\), covariantly in
  \(Y\)) and equipped with a composition pairing.
\end{definition}

\begin{lemma}
[Unconditional availability of \(\mathrm{HasExt}\)]
  \label{mk-lem:hasext_standard_mathlib}
  \lean{CategoryTheory.HasExt.standard}
  \mathlibok
  \dcref{custom}

  Every abelian category \(C\) carries a \(\mathrm{HasExt}\) structure; in particular the
  module-sheaf category \(U.\mathrm{Modules}\) admits the Ext bifunctor of
  Definition~\ref{mk-lem:ext_bifunctor_mathlib} unconditionally.
\end{lemma}

\begin{lemma}
[Degree-zero Ext is the Hom-group]
  \label{mk-lem:ext_homequiv_zero_mathlib}
  \lean{CategoryTheory.Abelian.Ext.homEquiv₀,
    CategoryTheory.Abelian.Ext.addEquiv₀}
  \mathlibok
  \dcref{custom}

  For objects \(A, Y\) of an abelian category with \(\mathrm{HasExt}\), there is a
  natural (additive) isomorphism \(\operatorname{Ext}^0(A, Y) \cong (A \to Y)\) between
  degree-zero Ext and the Hom-group. Applied to the first argument \(A =
  j_!\mathcal{O}_U\), this identifies \(\operatorname{Ext}^0(j_!\mathcal{O}_U,
  \mathcal{F})\) with
  \(\operatorname{Hom}_{X.\mathrm{Modules}}(j_!\mathcal{O}_U, \mathcal{F})\), which the
  corepresentability isomorphism of Lemma~\ref{mk-lem:jshriek_corepr} in turn identifies
  with \(\Gamma(U, \mathcal{F})\).
\end{lemma}

\begin{lemma}
[Ext vanishes on injectives in positive degree]
  \label{mk-lem:ext_eq_zero_of_injective_mathlib}
  \lean{CategoryTheory.Abelian.Ext.eq_zero_of_injective}
  \mathlibok
  \dcref{custom}

  Let \(C\) be an abelian category with \(\mathrm{HasExt}\) and \(I \in C\) an injective
  object. Then every element of \(\operatorname{Ext}^{n+1}(X, I)\) is zero, i.e.
  \(\operatorname{Ext}^{n+1}(X, I) = 0\) for all \(X\) and all \(n \geq 0\). The
  vanishing requires only the \emph{second} argument to be injective.
\end{lemma}

\begin{lemma}
[Covariant Ext long exact sequence of a short exact sequence]
  \label{mk-lem:ext_covariant_les_mathlib}
  \lean{CategoryTheory.Abelian.Ext.covariantSequence_exact,
    CategoryTheory.Abelian.Ext.covariant_sequence_exact₁,
    CategoryTheory.Abelian.Ext.covariant_sequence_exact₂,
    CategoryTheory.Abelian.Ext.covariant_sequence_exact₃}
  \mathlibok
  \dcref{custom}

  Let \(C\) be an abelian category with \(\mathrm{HasExt}\), let \(0 \to S_1 \to S_2 \to
  S_3 \to 0\) be a short exact sequence in \(C\), and fix a first argument \(X \in C\).
  Then the groups \(\operatorname{Ext}^n(X, S_i)\) assemble into a long exact sequence,
  with connecting homomorphism \(\operatorname{Ext}^{n_0}(X, S_3) \to
  \operatorname{Ext}^{n_1}(X, S_1)\) (for \(n_0 + 1 = n_1\)) given by composition with
  the Ext class of the extension. Mathlib packages the exactness both as the
  five-term \(\operatorname{ComposableArrows}\) exactness
  \(\operatorname{covariantSequence\_exact}\) and as the three step-lemmas
  three individual step-exactness lemmas.
\end{lemma}

\subsection{Project wrappers around the Ext realization}

The following are the thin project-side specializations of the Mathlib Ext anchors above to the
fixed first argument \(j_!\mathcal{O}_U\); they are the declarations the dimension-shift sub-lemmas
of Lemma~\ref{mk-lem:cech_to_cohomology_on_basis} actually invoke.

\begin{lemma}[Degree-zero absolute cohomology is global sections]\label{mk-lem:absolute_cohomology_zero}
  \lean{AlgebraicGeometry.absoluteCohomologyZeroAddEquiv}
  \leanok
  \dcref{custom}

  \uses{mk-def:absolute_cohomology, mk-lem:jshriek_corepr, mk-lem:ext_homequiv_zero_mathlib}
  \textit{Project specialization of Lemma~\ref{mk-lem:ext_homequiv_zero_mathlib} composed with
  corepresentability.}
  For every \(\mathcal{O}_X\)-module \(\mathcal{F}\) and open \(U \subseteq X\) there is an additive
  isomorphism
  \[
    H^0(U, \mathcal{F}) \;=\; \operatorname{Ext}^0(j_!\mathcal{O}_U, \mathcal{F})
      \;\cong\; \Gamma(U, \mathcal{F}).
  \]
\end{lemma}
\begin{proof}
  \leanok
  \uses{mk-lem:jshriek_corepr, mk-lem:ext_homequiv_zero_mathlib}
  Compose the degree-zero Ext comparison
  \(\operatorname{Ext}^0(j_!\mathcal{O}_U, \mathcal{F}) \cong
  \operatorname{Hom}_{X.\mathrm{Modules}}(j_!\mathcal{O}_U, \mathcal{F})\)
  (Lemma~\ref{mk-lem:ext_homequiv_zero_mathlib}) with the corepresentability isomorphism
  \(\operatorname{Hom}_{X.\mathrm{Modules}}(j_!\mathcal{O}_U, \mathcal{F}) \cong \Gamma(U,
  \mathcal{F})\) of Lemma~\ref{mk-lem:jshriek_corepr}. Both are additive, so is the composite.
\end{proof}

\begin{lemma}

[Naturality of \(H^0 \cong \Gamma\) in the coefficient sheaf]
  \label{mk-lem:absolute_cohomology_zero_natural}
  \lean{AlgebraicGeometry.absoluteCohomologyZeroAddEquiv_naturality,
    AlgebraicGeometry.homEquiv₀_comp_mk₀,
    AlgebraicGeometry.freeYonedaHomEquiv_naturality,
    AlgebraicGeometry.sheafificationHomAddEquiv_naturality,
    AlgebraicGeometry.jShriekOU_homEquiv_nat,
    AlgebraicGeometry.jShriekOU_homEquiv_naturality}
  \dcref{custom}

  \uses{mk-lem:absolute_cohomology_zero, mk-lem:jshriek_corepr}
  \textit{Project-bespoke obligation.}
  The isomorphism \(H^0(U, -) \cong \Gamma(U, -)\) of Lemma~\ref{mk-lem:absolute_cohomology_zero} is
  natural in the coefficient sheaf: for a morphism \(g : \mathcal{F}_1 \to \mathcal{F}_2\) in
  \(X.\mathrm{Modules}\) the square
  \[
    \begin{array}{ccc}
      H^0(U, \mathcal{F}_1) & \longrightarrow & H^0(U, \mathcal{F}_2)\\
      \cong\downarrow & & \downarrow\cong\\
      \Gamma(U, \mathcal{F}_1) & \xrightarrow{\;g_U\;} & \Gamma(U, \mathcal{F}_2)
    \end{array}
  \]
  commutes, where the top arrow is \(H^0(U, g)\), the functorial action of \(g\) on
  \(\operatorname{Ext}^0\) in the second variable, and the bottom arrow is the sections map
  \(g_U = g(U)\).
\end{lemma}
\begin{proof}
  \uses{mk-lem:absolute_cohomology_zero, mk-lem:jshriek_corepr}
  Under the degree-zero Ext comparison, post-composition by the degree-zero Ext class of \(g\)
  corresponds to post-composition by \(g\) in \(\operatorname{Hom}_{X.\mathrm{Modules}}(j_!
  \mathcal{O}_U, -)\); under the corepresentability isomorphism of Lemma~\ref{mk-lem:jshriek_corepr}
  (which is natural in the second variable, being a composite of the sheafification adjunction and
  the free–Yoneda evaluation, both natural), this in turn corresponds to applying \(g\) on sections
  over \(U\), i.e. \(g_U\). Hence the square commutes. In particular, when \(g_U\) is surjective so
  is \(H^0(U, g)\) — the transfer used by the surjectivity step of
  Lemma~\ref{mk-lem:absolute_cohomology_one_vanishing}.
\end{proof}

\begin{lemma}[Injective vanishing for absolute cohomology]\label{mk-lem:absolute_cohomology_injective_vanishing}
  \lean{AlgebraicGeometry.absoluteCohomology_eq_zero_of_injective}
  \leanok
  \dcref{custom}

  \uses{mk-def:absolute_cohomology, mk-lem:ext_eq_zero_of_injective_mathlib}
  \textit{Project specialization of Lemma~\ref{mk-lem:ext_eq_zero_of_injective_mathlib}.}
  For an injective \(\mathcal{O}_X\)-module \(\mathcal{I}\) and any \(n \geq 0\),
  \[
    H^{n+1}(U, \mathcal{I}) \;=\; \operatorname{Ext}^{n+1}(j_!\mathcal{O}_U, \mathcal{I}) \;=\; 0.
  \]
  The injective \(\mathcal{I}\) is the \emph{second} Ext argument, so no restriction is taken.
\end{lemma}
\begin{proof}
  \leanok
  \uses{mk-lem:ext_eq_zero_of_injective_mathlib}
  Immediate from Lemma~\ref{mk-lem:ext_eq_zero_of_injective_mathlib} with \(X = j_!\mathcal{O}_U\) and
  \(I = \mathcal{I}\).
\end{proof}

\begin{lemma}
[Covariant long exact sequence for absolute cohomology]
  \label{mk-lem:absolute_cohomology_covariant_les}
  \lean{AlgebraicGeometry.absoluteCohomology_covariant_exact₁,
    AlgebraicGeometry.absoluteCohomology_covariant_exact₂,
    AlgebraicGeometry.absoluteCohomology_covariant_exact₃}
  \dcref{custom}

  \uses{mk-def:absolute_cohomology, mk-lem:ext_covariant_les_mathlib}
  \textit{Project specialization of Lemma~\ref{mk-lem:ext_covariant_les_mathlib} at fixed first
  argument \(j_!\mathcal{O}_U\).}
  For a short exact sequence \(0 \to \mathcal{F}_1 \to \mathcal{F}_2 \to \mathcal{F}_3 \to 0\) in
  \(X.\mathrm{Modules}\), the groups \(H^n(U, \mathcal{F}_i) = \operatorname{Ext}^n(j_!
  \mathcal{O}_U, \mathcal{F}_i)\) assemble into the covariant long exact sequence, recorded in
  three step-exactness lemmas (for positions 1, 2, and 3 of the sequence).
\end{lemma}
\begin{proof}
  \uses{mk-lem:ext_covariant_les_mathlib}
  Each wrapper fixes the first Ext argument to \(j_!\mathcal{O}_U\) and threads the short exact
  sequence and degree hypotheses unchanged into the corresponding step-exactness lemma of
  Lemma~\ref{mk-lem:ext_covariant_les_mathlib}.
\end{proof}

\begin{lemma}
[Serre vanishing on affines]
  \label{mk-lem:affine_serre_vanishing}
  \lean{AlgebraicGeometry.affine_serre_vanishing,
    AlgebraicGeometry.affine_serre_vanishing_of_tildeVanishing}
  \dcref{01XB,custom}

  \uses{mk-def:affine_cover_system, mk-lem:affine_cech_vanishing_qcoh,
    mk-lem:affine_cech_vanishing_tilde_subcover,
    mk-lem:cech_to_cohomology_on_basis, mk-def:absolute_cohomology}
  \textit{Source: Stacks Project, Cohomology of Schemes, Tag 01XB
  (\texttt{lemma-quasi-coherent-affine-cohomology-zero}, Serre vanishing on
  affines).}
  Let \(U\) be an affine scheme and \(\mathcal{F}\) a quasi-coherent
  \(\mathcal{O}_U\)-module. Then the sheaf cohomology of \(\mathcal{F}\) vanishes in
  all positive degrees:
  \[
    H^p(U, \mathcal{F}) = 0 \qquad \text{for all } p > 0.
  \]
  The formal statement carries \([\,\mathrm{EnoughInjectives}\ X.\mathrm{Modules}\,]\) as an
  explicit instance hypothesis, inherited from Lemma~\ref{mk-lem:cech_to_cohomology_on_basis}.
\end{lemma}
\begin{proof}
  \uses{mk-def:affine_cover_system, mk-lem:affine_cech_vanishing_qcoh,
    mk-lem:cech_to_cohomology_on_basis, mk-def:absolute_cohomology}
  The proof instantiates the basis-comparison criterion
  (Lemma~\ref{mk-lem:cech_to_cohomology_on_basis}) at the affine cover system. Assemble the affine cover
  system \(s\) of Definition~\ref{mk-def:affine_cover_system}: its basis \(\mathcal{B}\) is the
  standard (distinguished) opens of the affine \(U\) and its admissible coverings \(\mathrm{Cov}\) are
  the standard open covers, the three nontrivial fields being discharged by
  Lemmas~\ref{mk-lem:affine_faces_mem}, \ref{mk-lem:affine_surj_of_vanishing}
  and~\ref{mk-lem:affine_injective_acyclic}.

  For the quasi-coherent coefficient \(\mathcal{F}\), Lemma~\ref{mk-lem:affine_cech_vanishing_qcoh}
  supplies the seed: \(\mathcal{F}\) has vanishing higher {\v C}ech cohomology for \(s\)
  (Definition~\ref{mk-def:has_vanishing_higher_cech}), i.e.\ condition (3) of
  Lemma~\ref{mk-lem:cech_to_cohomology_on_basis} holds for \(\mathcal{F}\). Applying the
  basis-comparison lemma to \(s\) and \(\mathcal{F}\) yields, for every standard open
  \(V \in \mathcal{B}\) and every \(p > 0\), the vanishing \(H^p(V, \mathcal{F}) = 0\). Taking
  \(V = U\) (the whole affine, itself a standard open) gives \(H^p(U, \mathcal{F}) = 0\) for all
  \(p > 0\). Throughout, the sheaf cohomology \(H^p(U, -) = \operatorname{Ext}^p(j_!\mathcal{O}_U, -)\)
  is the Ext realization of Definition~\ref{mk-def:absolute_cohomology}.
\end{proof}

\subsection{The affine instantiation (Tag 01XB)}

The Serre vanishing of Lemma~\ref{mk-lem:affine_serre_vanishing} is obtained by feeding a concrete
\emph{affine} cover system into the basis-comparison criterion. We construct that cover system
(Definition~\ref{mk-def:affine_cover_system}) and discharge its three nontrivial fields, then supply the
quasi-coherent seed that meets condition (3). The sub-lemmas below are the formalize-ready cut points
of this instantiation.

\begin{lemma}[Distinguished opens are closed under finite intersection]\label{mk-lem:affine_faces_mem}
  \lean{AlgebraicGeometry.affine_faces_mem}
  \leanok
  \dcref{custom}

  \uses{mk-def:standard_affine_cover}
  \textit{Source: Stacks Project, Schemes, \texttt{lemma-standard-open}
  (standard opens).}
  Let \(R\) be a commutative ring and \(\operatorname{Spec} R\) the associated affine scheme. The
  distinguished opens \(D(f)\), \(f \in R\), form a basis closed under intersection: \(D(f) \cap D(g)
  = D(fg)\). Consequently, for any standard open cover \(\mathcal{U} : V = \bigcup_{i \in \iota}
  D(g_i)\) of a distinguished open \(V\), every finite intersection \(D(g_{i_0}) \cap \cdots \cap
  D(g_{i_p}) = D(g_{i_0} \cdots g_{i_p})\) of its members is again a distinguished open. This is the
  \(\mathrm{faces\_mem}\) field (condition (1)) of the affine cover system.
\end{lemma}
\begin{proof}
  \leanok
  \uses{mk-def:standard_affine_cover}
  The identity \(D(f) \cap D(g) = D(fg)\) holds because a prime \(\mathfrak{p}\) avoids both \(f\) and
  \(g\) iff it avoids the product \(fg\). Iterating over the indices \(i_0, \ldots, i_p\) of a face
  gives \(\bigcap_k D(g_{i_k}) = D\bigl(\prod_k g_{i_k}\bigr)\), a distinguished open; the base set
  \(V = D(f)\) is distinguished by hypothesis. Hence every face of a standard cover lies in the basis.
\end{proof}

\begin{lemma}[Standard covers are cofinal among open covers of a distinguished open]\label{mk-lem:standard_cover_cofinal}
  \lean{AlgebraicGeometry.standard_cover_cofinal}
  \leanok
  \dcref{009L,custom}

  \uses{mk-def:standard_affine_cover, mk-lem:primeSpectrum_isBasis_basicOpen}
  \textit{Source: Stacks Project, Schemes, \texttt{lemma-standard-open}~(2), and
  Sheaves on Spaces, Tag 009L, \texttt{lemma-cofinal-systems-coverings-standard-case}.}
  Let \(V = D(f)\) be a distinguished open of \(\operatorname{Spec} R\), and let
  \(\{W_\alpha\}_{\alpha \in A}\) be an arbitrary open covering of \(V\) (indexed by a type \(A\),
  with each \(W_\alpha\) an open of \(\operatorname{Spec} R\) contained in \(V\)). Then there exist a
  finite index \(n \in \mathbb{N}\), a family \(g : \{1,\dots,n\} \to R\) of ring elements, and a
  reindexing \(\varphi : \{1,\dots,n\} \to A\) such that
  \[
    D(f) = \bigcup_{i=1}^n D(g_i)
    \qquad\text{and}\qquad
    D(g_i) \subseteq W_{\varphi(i)} \quad (1 \le i \le n).
  \]
  That is, every open covering of \(V\) is \emph{refined} by a finite covering by distinguished opens,
  each member of which sits inside a member of the original covering. This concrete
  indexed-cover/refinement statement is exactly the cofinality input the sheaf-on-a-basis principle
  (Tag 009L) requires; together with the closure of the basis under intersection
  (Lemma~\ref{mk-lem:affine_faces_mem}) it makes the standard covers a cofinal system among the open
  coverings of \(V\). It is the refinement step invoked in
  Lemma~\ref{mk-lem:affine_surj_of_vanishing}: a covering of \(V\) produced by local surjectivity is
  refined to a standard cover \(\mathcal{U} \in \mathrm{Cov}\) carrying the local lifts, with the
  reindexing \(\varphi\) transporting the local sections.
\end{lemma}
\begin{proof}
  \leanok
  \uses{mk-def:standard_affine_cover, mk-lem:affine_faces_mem, mk-lem:primeSpectrum_isBasis_basicOpen}
  The basic opens \(D(g)\), \(g \in R\), form a topological basis of \(\operatorname{Spec} R\)
  (Lemma~\ref{mk-lem:primeSpectrum_isBasis_basicOpen}), so each member \(W_\alpha\) of the given covering
  is a union of basic opens contained in it; choosing one basic open inside each \(W_\alpha\) through
  each point of \(V\) yields a basic-open covering of \(D(f)\) refining \(\{W_\alpha\}\). The
  distinguished open \(D(f)\) is quasi-compact, hence a finite subfamily
  \(D(g_1), \dots, D(g_n)\) already covers \(D(f)\); recording for each \(i\) the index
  \(\varphi(i) \in A\) of a member containing \(D(g_i)\) gives the asserted finite refinement
  \(D(f) = \bigcup_{i=1}^n D(g_i)\) with \(D(g_i) \subseteq W_{\varphi(i)}\). The basis is closed
  under the intersections required by Tag 009L by Lemma~\ref{mk-lem:affine_faces_mem}, so the standard
  covers form a cofinal system in the sense of the sheaf-on-a-basis cofinality principle.
\end{proof}

\begin{lemma}[Section surjectivity for the affine cover system]\label{mk-lem:affine_surj_of_vanishing}
  \lean{AlgebraicGeometry.affine_surj_of_vanishing}
  \leanok
  \dcref{custom}

  \uses{mk-lem:ses_cech_h1, mk-lem:standard_cover_cofinal, mk-lem:to_sheaf_preserves_epi,
    mk-lem:presheaf_is_locally_surjective, mk-lem:sheaf_locally_surjective_iff_epi,
    mk-lem:scheme_isBasis_affineOpens}
  \textit{Source: Stacks Project, Cohomology, \texttt{lemma-ses-cech-h1}.}
  Let \(V\) be a distinguished open of the affine and let \(S : 0 \to S_1 \to S_2 \to S_3 \to 0\) be a
  short exact sequence of \(\mathcal{O}_X\)-modules whose left term \(S_1\) has vanishing positive
  {\v C}ech cohomology over the standard covers of \(V\), i.e.\ \(\check{H}^p(\mathcal{U}, S_1) = 0\)
  for every standard cover \(\mathcal{U}\) of \(V\) and every \(p > 0\). Then the section map
  \(S_2(V) \to S_3(V)\) is surjective. This discharges the \(\mathrm{surj\_of\_vanishing}\) field of
  the affine cover system.
\end{lemma}
\begin{proof}
  \leanok
  \uses{mk-lem:ses_cech_h1, mk-lem:standard_cover_cofinal, mk-lem:to_sheaf_preserves_epi,
    mk-lem:presheaf_is_locally_surjective, mk-lem:sheaf_locally_surjective_iff_epi,
    mk-lem:scheme_isBasis_affineOpens}
    Fix a distinguished open \(V \in \mathcal{B}\) and a section \(t \in S_3(V)\); we produce a global
  lift of \(t\) over \(V\).

  \emph{Step 1 (local surjectivity of the section map).} The map \(g : S_2 \to S_3\) is an
  epimorphism in the category of \(\mathcal{O}_X\)-modules. By the gap-fill
  Lemma~\ref{mk-lem:to_sheaf_preserves_epi} the underlying-abelian-sheaf functor preserves this
  epimorphism, so the induced map of abelian sheaves is an epimorphism; by
  Lemma~\ref{mk-lem:sheaf_locally_surjective_iff_epi} an epimorphism of sheaves is locally surjective,
  hence \(g\) is a locally surjective morphism of presheaves
  (Lemma~\ref{mk-lem:presheaf_is_locally_surjective}).

  \emph{Step 2 (refine to a standard cover carrying local lifts).} Local surjectivity at \(t\)
  produces an open covering \(\{W_\alpha\}\) of \(V\) together with sections
  \(t_\alpha \in S_2(W_\alpha)\) mapping to \(t|_{W_\alpha}\). Since the affine opens form a basis
  (Lemma~\ref{mk-lem:scheme_isBasis_affineOpens}) and standard covers are cofinal among open coverings
  of \(V\) (Lemma~\ref{mk-lem:standard_cover_cofinal}), refine \(\{W_\alpha\}\) to a standard cover
  \(\mathcal{U} = \{D(g_i)\} \in \mathrm{Cov}\) of \(V\); restricting the \(t_\alpha\) along the
  refinement yields local lifts \(s^{\mathrm{loc}}_i \in S_2(D(g_i))\) with
  \(g(s^{\mathrm{loc}}_i) = t|_{D(g_i)}\).

  \emph{Step 3 (glue via \(\check{H}^1 = 0\)).} The left term \(S_1\) has vanishing positive
  {\v C}ech cohomology over every covering of \(V\) in \(\mathrm{Cov}\) (the field hypothesis); in
  particular \(\check{H}^1(\mathcal{U}, S_1) = 0\). Feeding the cover \(\mathcal{U}\), the local lifts
  \(s^{\mathrm{loc}}\), and this vanishing to the {\v C}ech \(\check{H}^1\)-surjectivity criterion
  (Lemma~\ref{mk-lem:ses_cech_h1}, which is stated over a raw finite family of opens) produces a global
  section \(s \in S_2(V)\) with \(g(s) = t\). Hence \(S_2(V) \to S_3(V)\) is surjective.
\end{proof}

\begin{lemma}[The sheafification--toSheaf square]
  \label{mk-lem:sheafificationCompToSheaf_mathlib}
  \lean{PresheafOfModules.sheafificationCompToSheaf}
  \mathlibok
  \dcref{custom}

  Let \(R\) be a sheaf of rings for the Grothendieck topology \(J\), and write
  \(L := \operatorname{sheafification}(\mathrm{id}_{R_0})
  : \mathrm{PreMod}(R_0) \to \mathrm{Mod}(R)\) for the sheafification functor from presheaves of
  \(R_0\)-modules to sheaves of \(R\)-modules (a left adjoint and a localization, whose counit is an
  isomorphism). Then there is a natural isomorphism of functors
  \[
    L \,\circ\, \mathrm{toSheaf}\,R \;\cong\;
    \mathrm{toPresheaf}\,R_0 \,\circ\, \operatorname{sheafify}_J^{\mathrm{Ab}},
  \]
  where \(\mathrm{toPresheaf}\,R_0\) forgets a presheaf of modules to its underlying presheaf of
  abelian groups and \(\operatorname{sheafify}_J^{\mathrm{Ab}} = \mathrm{presheafToSheaf}\,J\,
  \mathbf{Ab}\) is sheafification of presheaves of abelian groups. This is the commuting square
  relating module-sheafification to abelian-group-sheafification through the forgetful functors.
\end{lemma}

\begin{lemma}[Colimit preservation gives epimorphism preservation]
  \label{mk-lem:preservesEpi_of_preservesColimitsOfShape_mathlib}
  \lean{CategoryTheory.preservesEpimorphisms_of_preservesColimitsOfShape}
  \mathlibok
  \dcref{custom}

  Let \(F : \mathcal{C} \to \mathcal{D}\) be a functor that preserves colimits of shape
  \(\mathrm{WalkingSpan}\) (equivalently, preserves pushouts). Then \(F\) preserves epimorphisms: if
  \(f\) is an epimorphism in \(\mathcal{C}\) then \(F(f)\) is an epimorphism in \(\mathcal{D}\). In
  particular a functor preserving finite colimits preserves epimorphisms, since
  \(\mathrm{WalkingSpan}\) is a finite shape.
\end{lemma}

\begin{lemma}[The underlying-abelian-sheaf functor preserves finite colimits]
  \label{mk-lem:toSheaf_preservesFiniteColimits}
  \lean{AlgebraicGeometry.toSheaf_preservesFiniteColimits}
  \leanok
  \dcref{custom}
  \uses{mk-lem:sheafificationCompToSheaf_mathlib, mk-lem:mod_pmod_adjunction}
  \textit{Project-bespoke gap-fill (the missing right-exact dual of \(\mathrm{toSheaf}\)).}
  The forgetful functor \(\mathrm{toSheaf}\,R : \mathrm{Mod}(R) \to \mathrm{Sh}(J, \mathbf{Ab})\)
  from sheaves of \(\mathcal{O}_X\)-modules to the underlying sheaves of abelian groups preserves
  finite colimits. This is the right-exact counterpart of the limit-side fact that
  \(\mathrm{toSheaf}\,R\) preserves finite limits; Mathlib supplies the latter but not the former, so
  it is built here.
\end{lemma}
\begin{proof}
  \leanok
  \uses{mk-lem:sheafificationCompToSheaf_mathlib, mk-lem:mod_pmod_adjunction}
  Write \(L := \operatorname{sheafification}(\mathrm{id}_{R_0}) : \mathrm{PreMod}(R_0) \to
  \mathrm{Mod}(R)\) for the module-sheafification functor; \(L\) is a left adjoint whose counit is an
  isomorphism, so it is a localization that preserves all colimits. Throughout, \(\circ\) denotes
  diagrammatic composition: \(F \circ G\) means ``apply \(F\), then \(G\)''.

  \emph{Step 1 (the composite \(L \circ \mathrm{toSheaf}\) preserves finite colimits).} By the
  square isomorphism (Lemma~\ref{mk-lem:sheafificationCompToSheaf_mathlib}),
  \(L \circ \mathrm{toSheaf}\,R \cong \mathrm{toPresheaf}\,R_0 \circ
  \operatorname{sheafify}_J^{\mathrm{Ab}}\). The forgetful functor \(\mathrm{toPresheaf}\,R_0\) on
  presheaves of modules preserves finite colimits (colimits of presheaves of modules are computed
  objectwise on the underlying abelian groups), and abelian-group sheafification
  \(\operatorname{sheafify}_J^{\mathrm{Ab}}\) is a left adjoint, hence preserves all colimits.
  Therefore their composite, and so the isomorphic functor \(L \circ \mathrm{toSheaf}\,R\), preserves
  finite colimits.

  \emph{Step 2 (descend from the composite to \(\mathrm{toSheaf}\,R\) alone).} Let
  \(\iota : \mathrm{Mod}(R) \hookrightarrow \mathrm{PreMod}(R_0)\) be the inclusion of sheaves of
  modules into presheaves of modules, so that \(L \dashv \iota\) is the sheafification adjunction
  (Lemma~\ref{mk-lem:mod_pmod_adjunction}). The descent proceeds in three clauses.

  \emph{(i) The counit is an isomorphism.} Every sheaf of modules is already a sheaf, so its
  sheafification recovers it: the counit of \(L \dashv \iota\) is a natural isomorphism
  \[
    \varepsilon : \iota \circ L \;\xrightarrow{\;\cong\;}\; \mathrm{id}_{\mathrm{Mod}(R)} .
  \]

  \emph{(ii) \(\mathrm{toSheaf}\,R\) is a retract of \(L \circ \mathrm{toSheaf}\,R\).} Whiskering the
  counit isomorphism \(\varepsilon\) with \(\mathrm{toSheaf}\,R\) yields a natural isomorphism
  \[
    \mathrm{toSheaf}\,R \;\cong\; \iota \circ \bigl(L \circ \mathrm{toSheaf}\,R\bigr),
  \]
  which exhibits \(\mathrm{toSheaf}\,R\) as a retract of the composite \(L \circ \mathrm{toSheaf}\,R\):
  the section is precomposition with \(\iota\), and the retraction is \(\varepsilon\) whiskered with
  \(\mathrm{toSheaf}\,R\). By Step 1 the composite \(L \circ \mathrm{toSheaf}\,R\) preserves finite
  colimits.

  \emph{(iii) A retract of a finite-colimit-preserving functor preserves finite colimits.} If \(G\)
  sends a finite colimit cocone to a colimit cocone and \(F\) is a retract of \(G\), then the colimit
  comparison map for \(F\) is a retract of the (isomorphism) comparison map for \(G\), hence itself an
  isomorphism, so \(F\) preserves that colimit. Applying this with \(F = \mathrm{toSheaf}\,R\) and
  \(G = L \circ \mathrm{toSheaf}\,R\), we conclude that \(\mathrm{toSheaf}\,R\) preserves finite
  colimits.
\end{proof}

\begin{lemma}
[The underlying-abelian-sheaf functor preserves epimorphisms]
  \label{mk-lem:to_sheaf_preserves_epi}
  \lean{AlgebraicGeometry.toSheaf_preservesEpimorphisms}
  \leanok
  \dcref{custom}
  \uses{mk-lem:toSheaf_preservesFiniteColimits, mk-lem:preservesEpi_of_preservesColimitsOfShape_mathlib}
  \textit{Project-bespoke gap-fill (instance).}
  The forgetful functor from sheaves of \(\mathcal{O}_X\)-modules to the underlying sheaves of
  abelian groups preserves epimorphisms: if \(g : S_2 \to S_3\) is an epimorphism of
  \(\mathcal{O}_X\)-modules then the induced morphism of underlying abelian sheaves is an
  epimorphism. This is the one small instance unlocking the passage from a module epimorphism to
  local surjectivity in Lemma~\ref{mk-lem:affine_surj_of_vanishing}.
\end{lemma}
\begin{proof}
  \leanok
  \uses{mk-lem:toSheaf_preservesFiniteColimits, mk-lem:preservesEpi_of_preservesColimitsOfShape_mathlib}
  The functor \(\mathrm{toSheaf}\,R\) preserves finite colimits
  (Lemma~\ref{mk-lem:toSheaf_preservesFiniteColimits}); in particular it preserves colimits of shape
  \(\mathrm{WalkingSpan}\), which is a finite shape. A functor that preserves pushouts preserves
  epimorphisms (Lemma~\ref{mk-lem:preservesEpi_of_preservesColimitsOfShape_mathlib}), so
  \(\mathrm{toSheaf}\,R\) carries the epimorphism \(g\) to an epimorphism of abelian sheaves.
\end{proof}

\begin{lemma}[Locally surjective morphisms of presheaves]
  \label{mk-lem:presheaf_is_locally_surjective}
  \lean{CategoryTheory.Presheaf.IsLocallySurjective}
  \mathlibok
  \dcref{custom}

  For a Grothendieck topology \(J\), a morphism \(\alpha : \mathcal{G} \to \mathcal{H}\) of presheaves
  is \emph{locally surjective} when every section of \(\mathcal{H}\) over an object is, locally on a
  covering sieve, in the image of \(\alpha\). On a topological space this is the predicate
  \(\operatorname{TopCat.Presheaf.IsLocallySurjective}\): for each open \(V\) and each section
  \(t \in \mathcal{H}(V)\) there is an open covering \(\{W_\alpha\}\) of \(V\) with \(t|_{W_\alpha}\)
  in the image of \(\alpha_{W_\alpha}\).
\end{lemma}

\begin{lemma}[An epimorphism of sheaves is locally surjective]
  \label{mk-lem:sheaf_locally_surjective_iff_epi}
  \lean{CategoryTheory.Sheaf.isLocallySurjective_iff_epi'}
  \mathlibok
  \dcref{custom}

  \uses{mk-lem:presheaf_is_locally_surjective}
  A morphism of sheaves (of abelian groups, valued in a suitable category) is an epimorphism in the
  category of sheaves if and only if it is locally surjective. In particular an epimorphism of
  abelian sheaves on the space \(X\) is locally surjective, supplying the covering with local lifts
  used in Step 2 of Lemma~\ref{mk-lem:affine_surj_of_vanishing}.
\end{lemma}

\begin{lemma}[Affine opens form a basis]
  \label{mk-lem:scheme_isBasis_affineOpens}
  \lean{AlgebraicGeometry.Scheme.isBasis_affineOpens}
  \mathlibok
  \dcref{custom}

  The affine opens of a scheme \(X\) form a basis for its topology. On \(\operatorname{Spec} R\) this
  refines any open covering of a distinguished open to one by basic/affine opens, the input that
  turns a local-surjectivity covering into a standard cover lying in \(\mathrm{Cov}\)
  (Lemma~\ref{mk-lem:standard_cover_cofinal}).
\end{lemma}

\begin{lemma}[Basic opens form a topological basis of the prime spectrum]
  \label{mk-lem:primeSpectrum_isBasis_basicOpen}
  \lean{PrimeSpectrum.isTopologicalBasis_basic_opens}
  \mathlibok
  \dcref{custom}

  For a commutative (semi)ring \(R\), the distinguished/basic opens \(D(r)\), \(r \in R\), form a
  topological basis of \(\operatorname{Spec} R\): the sets \(\{D(r) : r \in R\}\) are open and every
  open of \(\operatorname{Spec} R\) is a union of them. This is the basis input used to refine an
  arbitrary open covering of a distinguished open to one by basic opens
  (Lemma~\ref{mk-lem:standard_cover_cofinal}).
\end{lemma}

\begin{lemma}[Standard-cover opens are the distinguished opens \(D(s_i)\)]\label{mk-lem:cover_datum_bridge}
  \lean{AlgebraicGeometry.coverOpen_affineOpenCoverOfSpan}
  \leanok
  \dcref{custom}

  \uses{mk-def:standard_affine_cover}
  \textit{Project-bespoke compatibility lemma.}
  Let \(\mathcal{U}\) be the standard affine cover associated to a spanning family
  \(s : \iota \to R\) with \(\operatorname{span}(\operatorname{range} s) = \top\)
  (Definition~\ref{mk-def:standard_affine_cover}). For each index \(i\) the \(i\)-th open of
  \(\mathcal{U}\) is exactly the distinguished open \(D(s_i)\):
  \[
    \operatorname{coverOpen} \mathcal{U}\, i \;=\; D(s_i)
    \quad\text{in } \mathrm{Opens}(\operatorname{Spec} R).
  \]
  This identifies the opens of a standard cover with distinguished opens, which is the form in
  which a standard cover is recognised inside the basis \(\mathcal{B}\) of the affine cover system.
\end{lemma}
\begin{proof}
  \leanok
  \uses{mk-def:standard_affine_cover}
  By construction the \(i\)-th member of the standard affine cover attached to \(s\) is the basic
  open \(D(s_i)\) of \(\operatorname{Spec} R\); the spanning hypothesis
  \(\operatorname{span}(\operatorname{range} s) = \top\) guarantees these distinguished opens cover
  \(\operatorname{Spec} R\), and the open carried at index \(i\) is \(D(s_i)\) on the nose.
\end{proof}

\begin{lemma}[Injective acyclicity for the affine cover system]\label{mk-lem:affine_injective_acyclic}
  \lean{AlgebraicGeometry.affine_injective_acyclic}
  \leanok
  \dcref{custom}

  \uses{mk-lem:injective_cech_acyclic, mk-lem:cover_datum_bridge}
  \textit{Source: Stacks Project, Cohomology, \texttt{lemma-injective-trivial-cech}.}
  Every injective \(\mathcal{O}_X\)-module \(\mathcal{I}\) has vanishing positive-degree {\v C}ech
  cohomology over a standard cover \emph{of the whole affine}: if
  \(s : \iota \to R\) spans, \(\operatorname{span}(\operatorname{range} s) = \top\), so that the
  distinguished opens \(D(s_i)\) cover \(\operatorname{Spec} R\), then
  \(\check{H}^p(\mathcal{U}, \mathcal{I}) = 0\) for every \(p > 0\). This is the special case of the
  family-form injective {\v C}ech-acyclicity (Lemma~\ref{mk-lem:injective_cech_acyclic}) at the
  spanning family \(s\); the general \(\mathrm{injective\_acyclic}\) field of the affine cover system
  is discharged by that family-form lemma directly.
\end{lemma}
\begin{proof}
  \leanok
  \uses{mk-lem:injective_cech_acyclic, mk-lem:cover_datum_bridge}
  The opens of the standard cover \(\mathcal{U}\) attached to the spanning family \(s\) are the
  distinguished opens \(D(s_i)\) (Lemma~\ref{mk-lem:cover_datum_bridge}). Applying the family-form
  injective {\v C}ech-acyclicity (Lemma~\ref{mk-lem:injective_cech_acyclic}) to the finite family
  \(i \mapsto D(s_i)\) gives \(\check{H}^p(\mathcal{U}, \mathcal{I}) = 0\) for all \(p > 0\).
\end{proof}

\begin{definition}[The affine cover system]\label{mk-def:affine_cover_system}
  \lean{AlgebraicGeometry.affineCoverSystem}
  \leanok
  \dcref{01XB,custom}

  \uses{mk-def:basis_cov_system, mk-lem:affine_faces_mem, mk-lem:affine_surj_of_vanishing,
    mk-lem:injective_cech_acyclic}
  \textit{Source: Stacks Project, Cohomology of Schemes, Tag 01XB
  (\texttt{lemma-quasi-coherent-affine-cohomology-zero}, proof).}
  The \emph{affine cover system} for an affine scheme \(X = \operatorname{Spec} R\) (or an affine open
  of a scheme) is the instance of Definition~\ref{mk-def:basis_cov_system} whose basis \(\mathcal{B}\) is
  the distinguished opens \(D(f)\), \(f \in R\), and whose admissible coverings \(\mathrm{Cov}\) are
  the standard open covers (finite coverings of a distinguished open by distinguished opens). Beyond
  the basis \(\mathcal{B}\) and the covering set \(\mathrm{Cov}\) just described, it carries three
  proof fields: the \(\mathrm{faces\_mem}\) field by Lemma~\ref{mk-lem:affine_faces_mem}; the
  \(\mathrm{surj\_of\_vanishing}\) field by Lemma~\ref{mk-lem:affine_surj_of_vanishing}; and the
  \(\mathrm{injective\_acyclic}\) field by the family-form injective {\v C}ech-acyclicity
  Lemma~\ref{mk-lem:injective_cech_acyclic} (applied to the distinguished opens of each standard cover).
\end{definition}

\begin{lemma}[Sections--tilde equivalence from an invertible counit]\label{mk-lem:qcoh_iso_tilde_sections}
  \lean{AlgebraicGeometry.qcoh_iso_tilde_sections,
    AlgebraicGeometry.qcoh_iso_tilde_sections_hom,
    AlgebraicGeometry.qcoh_iso_tilde_sections_inv}
  \leanok
  \dcref{01HV,custom}

  \textit{Source: Stacks Project, Schemes, Tag 01HV, \texttt{lemma-spec-sheaves}.}
  Let \(X = \operatorname{Spec} R\) and let \(\mathcal{F}\) be an
  \(\mathcal{O}_X\)-module for which the tilde--global-sections counit
  \[
    \operatorname{fromTilde\Gamma} :
      \widetilde{\Gamma(X,\mathcal F)} \longrightarrow \mathcal F
  \]
  is an isomorphism. Then its inverse gives a canonical isomorphism
  \(\mathcal F\cong\widetilde{\Gamma(X,\mathcal F)}\); the forward map is
  the inverse of the counit and the reverse map is the counit itself.
\end{lemma}
\begin{proof}
  \leanokTake the inverse of
  \(\operatorname{fromTilde\Gamma}\). The inverse laws give the two
  isomorphism identities, and identify its inverse component with the
  original counit.
\end{proof}

\begin{lemma}[Affine structure theorem from a global presentation]\label{mk-lem:qcoh_iso_tilde_sections_of_presentation}
  \lean{AlgebraicGeometry.qcoh_iso_tilde_sections_of_presentation}
  \leanok
  \dcref{01I8,custom}

  \uses{mk-lem:qcoh_iso_tilde_sections, mk-lem:isIso_fromTildeGamma_of_presentation}
  \textit{Source: Stacks Project, Schemes, Tag 01I8, \texttt{lemma-quasi-coherent-affine}.}
  Let \(X = \operatorname{Spec} R\) and let \(\mathcal{F}\) be an \(\mathcal{O}_X\)-module that admits
  a \emph{global} presentation (a global exact sequence
  \(\mathcal{O}_X^{(J)} \to \mathcal{O}_X^{(I)} \to \mathcal{F} \to 0\), the data
  \(\mathcal{F}.\mathrm{Presentation}\)). Then \(\mathcal{F}\) is isomorphic to the sheaf
  \(\widetilde{\Gamma(X, \mathcal{F})}\) associated to its module of global sections. This is the
  presentation-driven discharge of the conditional hypothesis
  \(\operatorname{IsIso}(\operatorname{fromTilde\Gamma})\) of Lemma~\ref{mk-lem:qcoh_iso_tilde_sections}.
\end{lemma}
\begin{proof}
  \leanok
  \uses{mk-lem:qcoh_iso_tilde_sections, mk-lem:isIso_fromTildeGamma_of_presentation}
  A global presentation of \(\mathcal{F}\) makes the tilde–\(\Gamma\) counit
  \(\operatorname{fromTilde\Gamma}\) an isomorphism
  (Lemma~\ref{mk-lem:isIso_fromTildeGamma_of_presentation}): \(\Gamma\) and \(\widetilde{(-)}\) are exact
  enough to carry the presentation of \(\mathcal{F}\) to a presentation of
  \(\widetilde{\Gamma(X, \mathcal{F})}\), and the counit is then an isomorphism by the five lemma on
  the two presentations. Apply
  \ref{mk-lem:qcoh_iso_tilde_sections} to this invertible counit.
\end{proof}

\begin{lemma}[Counit is iso from a presentation]
  \label{mk-lem:isIso_fromTildeGamma_of_presentation}
  \lean{AlgebraicGeometry.isIso_fromTildeΓ_of_presentation}
  \mathlibok
  \dcref{custom}

  For \(R\) a ring and \(\mathcal{F}\) an \(\mathcal{O}_{\operatorname{Spec} R}\)-module equipped with
  a global presentation \(\mathcal{F}.\mathrm{Presentation}\), the tilde–\(\Gamma\) counit
  \(\operatorname{fromTilde\Gamma} : \widetilde{\Gamma(\operatorname{Spec} R, \mathcal{F})} \to
  \mathcal{F}\) is an isomorphism.
\end{lemma}

\section{The 01I8 counit isomorphism via section localization}

Let \(X=\operatorname{Spec}R\) and let \(\mathcal F\) be quasi-coherent. The
distinguished-open component of the counit
\(\operatorname{fromTilde\Gamma}:\widetilde{\Gamma(X,\mathcal F)}\to\mathcal F\)
is the canonical map between two localizations of \(\Gamma(X,\mathcal F)\) at the
powers of \(f\). Thus it is an isomorphism once the restriction map
\(\Gamma(X,\mathcal F)\to\Gamma(D(f),\mathcal F)\) is known to have the localization
property. Lemma~\ref{mk-lem:qcoh_section_isLocalizedModule} establishes this property;
checking it on the basis of distinguished opens then proves the counit is an
isomorphism.

\begin{lemma}[The tilde--\(\Gamma\) counit]
  \label{mk-lem:fromTildeGamma_mathlib}
  \lean{AlgebraicGeometry.Scheme.Modules.fromTildeΓ}
  \mathlibok
  \dcref{custom}

  For a ring \(R\) and an \(\mathcal{O}_{\operatorname{Spec} R}\)-module \(\mathcal{F}\), the
  tilde--\(\Gamma\) counit
  \(\operatorname{fromTilde\Gamma} : \widetilde{\Gamma(\operatorname{Spec} R, \mathcal{F})} \to
  \mathcal{F}\) is the natural morphism of sheaves of modules whose component over a distinguished
  open \(D(f)\) is the localization lift \(\operatorname{IsLocalizedModule.lift}\) of the
  section-restriction map \(\Gamma(\operatorname{Spec} R, \mathcal{F}) \to \Gamma(D(f), \mathcal{F})\)
  along the structure-sheaf localization \(\widetilde{(-)}.\operatorname{toOpen}\) (the latter
  exhibiting \(\Gamma(D(f), \widetilde{M}) = M_f\) as a localization at the powers of \(f\)).
\end{lemma}

\begin{lemma}[Criterion for the counit to be an isomorphism]
  \label{mk-lem:isIso_fromTildeGamma_iff_mathlib}
  \lean{AlgebraicGeometry.isIso_fromTildeΓ_iff}
  \mathlibok
  \dcref{custom}

  For an \(\mathcal{O}_{\operatorname{Spec} R}\)-module \(\mathcal{F}\), the counit
  \(\operatorname{fromTilde\Gamma}\) is an isomorphism if and only if \(\mathcal{F}\) lies in the
  essential image of the tilde functor \(M \mapsto \widetilde{M}\); equivalently
  \(\mathcal{F} \cong \widetilde{M}\) for some \(R\)-module \(M\).
\end{lemma}

\begin{lemma}[The forgetful functor reflects isomorphisms]
  \label{mk-lem:forget_reflectsIso_mathlib}
  \lean{SheafOfModules.fullyFaithfulForget}
  \mathlibok
  \dcref{custom}

  The forgetful functor \(\operatorname{forget}\) from sheaves of \(\mathcal{O}_X\)-modules to the
  underlying sheaves of abelian groups is fully faithful, hence reflects isomorphisms; the same holds
  for \(\operatorname{toSheaf}\). Consequently a morphism of sheaves of modules is an isomorphism as
  soon as its underlying morphism of sheaves of abelian groups is, and — since a morphism of sheaves
  which is an isomorphism on every member of a basis of the topology is an isomorphism — it suffices
  to check \(\operatorname{fromTilde\Gamma}\) on the basis of distinguished opens \(\{D(f)\}\).
\end{lemma}

\begin{lemma}[Two localizations at the same set are canonically isomorphic]
  \label{mk-lem:isLocalizedModule_linearEquiv_mathlib}
  \lean{IsLocalizedModule.linearEquiv}
  \mathlibok
  \dcref{custom}

  Let \(S \subseteq R\) be a submonoid of a commutative ring and let \(\varphi : M \to M'\) and
  \(\psi : M \to M''\) be \(R\)-linear maps that each exhibit their target as the localization of
  \(M\) at \(S\) (i.e.\ \(\operatorname{IsLocalizedModule} S\, \varphi\) and
  \(\operatorname{IsLocalizedModule} S\, \psi\)). Then there is a canonical \(R\)-linear isomorphism
  \(M' \cong M''\) compatible with \(\varphi, \psi\). In particular the localization lift
  \(\operatorname{IsLocalizedModule.lift} S\, \varphi\, \psi\) is an isomorphism whenever both
  \(\varphi\) and \(\psi\) are localizations of \(M\) at \(S\).
\end{lemma}

\begin{lemma}[The structure-sheaf localization of \(\widetilde{M}\) over \(D(f)\)]
  \label{mk-lem:tilde_toOpen_isLocalizedModule_mathlib}
  \lean{AlgebraicGeometry.tilde.toOpen}
  \mathlibok
  \dcref{custom}

  For \(M\) an \(R\)-module and \(f \in R\), Mathlib's map
  \(\operatorname{toOpen} : M \to \Gamma(D(f), \widetilde{M})\) exhibits \(\Gamma(D(f), \widetilde{M}) =
  M_f\) as the localization of \(M\) at the powers of \(f\): it carries an
  \(\operatorname{IsLocalizedModule}(\operatorname{powers} f)\) instance. Moreover the global component
  \(\operatorname{toOpen}(M, \top)\) is an isomorphism (Mathlib's
  \(\operatorname{isIso\_toOpen\_top}\)), identifying \(\Gamma(\operatorname{Spec} R, \widetilde{M}) =
  M\).
\end{lemma}

\begin{lemma}[Transport of a presentation along a colimit-preserving functor]
  \label{mk-lem:presentation_map_mathlib}
  \lean{SheafOfModules.Presentation.map}
  \mathlibok
  \dcref{custom}

  Let \(F : \operatorname{SheafOfModules} R \to \operatorname{SheafOfModules} S\) be a functor between
  categories of sheaves of modules (possibly over different sites) that preserves colimits and carries
  the unit object to the unit object, \(F(\mathcal{O}_R) \cong \mathcal{O}_S\). Then any presentation
  of a sheaf of modules \(\mathcal{M}\) — a right-exact sequence
  \(\mathcal{O}_R^{(J)} \to \mathcal{O}_R^{(I)} \to \mathcal{M} \to 0\) of generators and relations —
  is carried by \(F\) to a presentation of \(F(\mathcal{M})\).
\end{lemma}

\begin{lemma}[Transport of a presentation across an isomorphism]
  \label{mk-lem:presentation_ofIsIso_mathlib}
  \lean{SheafOfModules.Presentation.ofIsIso}
  \mathlibok
  \dcref{custom}

  Let \(\phi : \mathcal{M} \to \mathcal{N}\) be an isomorphism of sheaves of modules. A presentation of
  \(\mathcal{M}\) induces a presentation of \(\mathcal{N}\), by composing its generators and relations
  with \(\phi\).
\end{lemma}

\begin{lemma}[Site-equivalence transport of sheaves of modules]
  \label{mk-lem:pushforwardPushforwardEquivalence_mathlib}
  \lean{SheafOfModules.pushforwardPushforwardEquivalence}
  \mathlibok
  \dcref{custom}

  Let \(e : \mathcal{C} \simeq \mathcal{D}\) be an equivalence of sites whose functor and inverse
  are continuous.  If \(\mathcal{O}_{\mathcal{C}}\) and \(\mathcal{O}_{\mathcal{D}}\) are sheaves of
  rings related by mutually inverse comparison maps along the two pushforwards, then \(e\) induces
  an equivalence of categories of sheaves of modules
  \(\operatorname{SheafOfModules} \mathcal{O}_{\mathcal{C}} \simeq
  \operatorname{SheafOfModules} \mathcal{O}_{\mathcal{D}}\), compatible with those pushforwards.
  This is the engine that bridges two pictures of ``restriction to an open''.
\end{lemma}

\begin{lemma}[The over-site of a topological space is equivalent to its subspace site]
  \label{mk-lem:overEquivalence_mathlib}
  \lean{TopologicalSpace.Opens.overEquivalence}
  \mathlibok
  \dcref{custom}

  For a topological space \(T\) and an open \(U \subseteq T\), the over-category
  \((\operatorname{Opens} T).\operatorname{over} U\) of opens contained in \(U\) is equivalent to the
  category \(\operatorname{Opens} U\) of opens of the subspace \(U\).
\end{lemma}

\begin{lemma}[Continuity of the over-site equivalence]
  \label{mk-lem:overEquivalence_isContinuous}
  \lean{AlgebraicGeometry.Opens.overEquivalence_functor_coverPreserving,
    AlgebraicGeometry.Opens.overEquivalence_inverse_coverPreserving,
    AlgebraicGeometry.Opens.overEquivalence_functor_isContinuous,
    AlgebraicGeometry.Opens.overEquivalence_inverse_isContinuous,
    AlgebraicGeometry.overEquivalence_functor_isContinuous_toScheme,
    AlgebraicGeometry.overEquivalence_inverse_isContinuous_toScheme}
  \dcref{custom}

  \uses{mk-lem:overEquivalence_mathlib}
  Both functors of the over-site equivalence
  \((\operatorname{Opens} T).\operatorname{over} U \simeq \operatorname{Opens} U\) of
  Lemma~\ref{mk-lem:overEquivalence_mathlib} are cover-preserving and continuous for the respective
  Grothendieck topologies of pointwise covers. The same holds after identifying the
  subspace \(D(g)\) with its associated open subscheme.
\end{lemma}
\begin{proof}
  \uses{mk-lem:overEquivalence_mathlib}
  The opens Grothendieck topology is the pointwise-cover predicate; the over-topology unfolds via
  the membership criterion for over-topologies, and the subspace-versus-image point correspondence is
  the obvious \(\langle x', hx'\rangle\) / underlying-value bookkeeping, giving cover-preservation of
  each functor. Continuity then follows from cover-preservation together with the cover-density,
  local-fullness and local-faithfulness that hold automatically for the fully-faithful, essentially
  surjective functors of an equivalence.
\end{proof}

\begin{lemma}[Quasi-coherence datum: a cover with a presentation on each member]
  \label{mk-lem:quasicoherentData_mathlib}
  \lean{SheafOfModules.QuasicoherentData}
  \mathlibok
  \dcref{custom}

  A \emph{quasi-coherence datum} for a sheaf of modules \(\mathcal{F}\) consists of a family of objects
  \(U_i\) of the site that covers the terminal object, together with, for each \(i\), a presentation of
  the over-picture restriction \(\mathcal{F}.\operatorname{over} U_i\). A sheaf of modules is
  quasi-coherent precisely when such a datum exists.
\end{lemma}

\begin{lemma}[Sections of a restricted module are sections over the image open]
  \label{mk-lem:restrict_obj_mathlib}
  \lean{AlgebraicGeometry.Scheme.Modules.restrict_obj}
  \mathlibok
  \dcref{custom}

  For an \(\mathcal{O}_Y\)-module \(\mathcal{M}\), an open immersion \(\iota : X \to Y\), and an open
  \(U \subseteq X\), the sections of the restricted module over \(U\) coincide definitionally with the
  sections of \(\mathcal{M}\) over the image open \(\iota(U)\):
  \(\Gamma(\mathcal{M}.\operatorname{restrict} \iota, U) = \Gamma(\mathcal{M}, \iota(U))\), and
  likewise for the restriction maps. The section comparison is therefore an equality, not merely an
  isomorphism.
\end{lemma}

\begin{lemma}[Section restriction of \(\widetilde{M}\) localizes]\label{mk-lem:tilde_section_isLocalizedModule}
  \lean{AlgebraicGeometry.tilde_section_isLocalizedModule}
  \leanok
  \dcref{custom}

  \uses{mk-lem:tilde_toOpen_isLocalizedModule_mathlib}
  Let \(M\) be an \(R\)-module and \(f \in R\). The section-restriction map
  \(\Gamma(\operatorname{Spec} R, \widetilde{M}) \to \Gamma(D(f), \widetilde{M})\) of the associated
  sheaf \(\widetilde{M}\) exhibits its target as the localization of its source at the powers of
  \(f\): one has \(\operatorname{IsLocalizedModule}(\operatorname{powers} f)\) for it.
\end{lemma}
\begin{proof}
  \leanok
  \uses{mk-lem:tilde_toOpen_isLocalizedModule_mathlib}
  By Lemma~\ref{mk-lem:tilde_toOpen_isLocalizedModule_mathlib} the distinguished-open component
  \(\operatorname{toOpen} : M \to \Gamma(D(f), \widetilde{M})\) is a localization at the powers of
  \(f\), and the global component \(M \to \Gamma(\operatorname{Spec} R, \widetilde{M})\) is an
  isomorphism. Up to this global isomorphism, the section-restriction map is exactly
  \(\operatorname{toOpen}\) over \(D(f)\); transporting the localization property along the
  global-sections isomorphism that identifies \(M\) with \(\Gamma(\operatorname{Spec} R,
  \widetilde{M})\) on the source shows that the section-restriction map is itself a localization at the
  powers of \(f\).
\end{proof}

\begin{lemma}[Section restriction localizes when the counit is an isomorphism]\label{mk-lem:section_isLocalizedModule_of_isIso_fromTildeGamma}
  \lean{AlgebraicGeometry.section_isLocalizedModule_of_isIso_fromTildeΓ}
  \leanok
  \dcref{custom}

  \uses{mk-lem:tilde_section_isLocalizedModule, mk-lem:qcoh_iso_tilde_sections}
  Let \(\mathcal{F}\) be an \(\mathcal{O}_{\operatorname{Spec} R}\)-module for which the
  tilde--\(\Gamma\) counit \(\operatorname{fromTilde\Gamma}\) is an isomorphism, and let \(f \in R\).
  Then the section-restriction map \(\Gamma(\operatorname{Spec} R, \mathcal{F}) \to \Gamma(D(f),
  \mathcal{F})\) is a localization at the powers of \(f\).
\end{lemma}
\begin{proof}
  \leanok
  \uses{mk-lem:tilde_section_isLocalizedModule, mk-lem:qcoh_iso_tilde_sections}
  When \(\operatorname{fromTilde\Gamma}\) is an isomorphism, \(\mathcal{F} \cong \widetilde{M}\) with
  \(M = \Gamma(\operatorname{Spec} R, \mathcal{F})\) (Lemma~\ref{mk-lem:qcoh_iso_tilde_sections}). The
  counit isomorphism, pushed through the forgetful functor, identifies the section-restriction map of
  \(\mathcal{F}\) with that of \(\widetilde{M}\) — compatibly on source and target by naturality. The
  latter is a localization by Lemma~\ref{mk-lem:tilde_section_isLocalizedModule}; conjugating by the
  isomorphism transports the localization property to the section-restriction map of \(\mathcal{F}\).
\end{proof}

\begin{lemma}[Section restriction localizes for a globally presented module]\label{mk-lem:section_isLocalizedModule_of_presentation}
  \lean{AlgebraicGeometry.section_isLocalizedModule_of_presentation}
  \leanok
  \dcref{custom}

  \uses{mk-lem:section_isLocalizedModule_of_isIso_fromTildeGamma,
    mk-lem:isIso_fromTildeGamma_of_presentation}
  Let \(\mathcal{F}\) be an \(\mathcal{O}_{\operatorname{Spec} R}\)-module that admits a global
  presentation, and let \(f \in R\). Then the section-restriction map
  \(\Gamma(\operatorname{Spec} R, \mathcal{F}) \to \Gamma(D(f), \mathcal{F})\) is a localization at the
  powers of \(f\).
\end{lemma}
\begin{proof}
  \leanok
  \uses{mk-lem:section_isLocalizedModule_of_isIso_fromTildeGamma,
    mk-lem:isIso_fromTildeGamma_of_presentation}
    A global presentation makes the counit \(\operatorname{fromTilde\Gamma}\) an isomorphism
  (Lemma~\ref{mk-lem:isIso_fromTildeGamma_of_presentation}); now apply
  Lemma~\ref{mk-lem:section_isLocalizedModule_of_isIso_fromTildeGamma}.
\end{proof}

\begin{lemma}[Finite presentation cover from quasi-coherence]
  \label{mk-lem:qcoh_finite_presentation_cover}
  \lean{AlgebraicGeometry.qcoh_finite_presentation_cover, AlgebraicGeometry.coversTop_iSup_eq_top}
  \dcref{custom}

  \uses{mk-lem:exists_finite_basicOpen_subcover, mk-lem:quasicoherentData_mathlib}
  Let \(X = \operatorname{Spec} R\) and let \(\mathcal{F}\) be a quasi-coherent
  \(\mathcal{O}_X\)-module. Then there exist a finite family \(g_1, \ldots, g_n \in R\) with
  \(\operatorname{span}\{g_1, \ldots, g_n\} = R\) and, for each \(j\), a member \(U_{\varphi(j)}\) of
  the quasi-coherence cover with \(D(g_j) \subseteq U_{\varphi(j)}\) carrying a presentation of the
  over-picture restriction \(\mathcal{F}.\operatorname{over} U_{\varphi(j)}\).
\end{lemma}
\begin{proof}
  \uses{mk-lem:exists_finite_basicOpen_subcover, mk-lem:quasicoherentData_mathlib}
  Quasi-coherence supplies a quasi-coherence datum (Lemma~\ref{mk-lem:quasicoherentData_mathlib}): a cover
  \((U_i)_i\) of \(X\) together with a presentation of \(\mathcal{F}.\operatorname{over} U_i\) for each
  \(i\). The cover-of-terminal-object condition translates to \(\bigcup_i U_i = \top\); feeding this to
  the finite basic-open refinement (Lemma~\ref{mk-lem:exists_finite_basicOpen_subcover}) yields finitely
  many \(g_j\) with \(D(g_j) \subseteq U_{\varphi(j)}\) and \(\operatorname{span}\{g_j\} = R\). The
  required presentation of \(\mathcal{F}.\operatorname{over} U_{\varphi(j)}\) is the one attached to
  \(U_{\varphi(j)}\) by the datum.
\end{proof}

\begin{lemma}[Restricting a presentation to a basic open]\label{mk-lem:presentation_over_basicOpen}
  \lean{AlgebraicGeometry.presentationOverBasicOpen}
  \leanok
  \dcref{custom}

  \uses{mk-lem:presentation_map_mathlib}
  Let \(\mathcal{M}\) be an \(\mathcal{O}_{\operatorname{Spec} R}\)-module, let \(U\) be an open
  carrying a presentation of the over-restriction \(\mathcal{M}.\operatorname{over} U\), and let
  \(g \in R\) with \(D(g) \subseteq U\). Then \(\mathcal{M}.\operatorname{over} D(g)\) admits a
  presentation.
\end{lemma}
\begin{proof}
  \leanok
  \uses{mk-lem:presentation_map_mathlib, mk-lem:pushforwardPushforwardEquivalence_mathlib,
    mk-lem:presentation_ofIsIso_mathlib}
  The further over-restriction from \(U\) to the smaller open \(D(g)\) is a left adjoint, hence
  preserves colimits, and carries the unit object to the unit object; the nested over-restriction
  instances make it a functor between the corresponding categories of sheaves of modules. Transport the
  given presentation of \(\mathcal{M}.\operatorname{over} U\) along this functor by
  Lemma~\ref{mk-lem:presentation_map_mathlib} to obtain a presentation of
  \(\mathcal{M}.\operatorname{over} D(g)\). Concretely the transport is realized by first building the
  relevant site equivalence via Lemma~\ref{mk-lem:pushforwardPushforwardEquivalence_mathlib}, transporting
  the presentation through its inverse with \(\operatorname{Presentation.map}\), and closing along the
  residual isomorphism with \(\operatorname{Presentation.ofIsIso}\)
  (Lemma~\ref{mk-lem:presentation_ofIsIso_mathlib}).
\end{proof}

\begin{definition}[Module equivalence over a basic open]
  \label{mk-def:modules_over_basicOpen_equivalence}
  \lean{AlgebraicGeometry.modulesOverBasicOpenEquivalence,
    AlgebraicGeometry.overForgetIso, AlgebraicGeometry.overForgetInvIso,
    AlgebraicGeometry.overBasicOpenRingHom, AlgebraicGeometry.overBasicOpenRingInvHom,
    AlgebraicGeometry.specBasicOpen_ι_image_overEquivalence_functor}
  \dcref{custom}

  \uses{mk-lem:overEquivalence_isContinuous, mk-lem:pushforwardPushforwardEquivalence_mathlib}
  For \(g \in R\), there is an equivalence of categories of sheaves of
  modules
  \[
    \widehat{D(g)}.\operatorname{toScheme}.\mathrm{Modules} \;\simeq\;
    \operatorname{SheafOfModules}\bigl((\operatorname{Spec} R).\operatorname{ringCatSheaf}.\operatorname{over} D(g)\bigr)
  \]
  obtained by feeding the structure-sheaf comparison data
  \(\varphi = \operatorname{overBasicOpenRingHom} g\),
  \(\psi = \operatorname{overBasicOpenRingInvHom} g\) and their coherence witnesses into
  Lemma~\ref{mk-lem:pushforwardPushforwardEquivalence_mathlib} along the now-continuous site equivalence
  of Lemma~\ref{mk-lem:overEquivalence_isContinuous}. The comparison morphisms \(\varphi, \psi\) are built
  by whiskering the structure sheaf along the functor isomorphism
  \(\operatorname{overForgetIso} g : \operatorname{Over.forget} D(g) \cong
  \operatorname{overEquivalence}.\operatorname{functor} \circ \iota.\operatorname{opensFunctor}\)
  (whose object component is the image--preimage identity
  \(\iota(\iota^{-1} V) = V\) for opens \(V \le D(g)\)) and its definitional inverse
  \(\operatorname{overForgetInvIso} g\). Lemma~\ref{mk-lem:restrict_over_compat} gives its
  object-level comparison with restriction to the open subscheme.
\end{definition}

\begin{lemma}[Compatibility between over-restriction and open-subscheme restriction]\label{mk-lem:restrict_over_compat}
  \lean{AlgebraicGeometry.overBasicOpenIsoRestrict}
  \leanok
  \dcref{custom}

  \uses{mk-def:modules_over_basicOpen_equivalence, mk-lem:modules_restrict_basicOpen,
    mk-lem:pushforwardPushforwardEquivalence_mathlib, mk-lem:restrict_obj_mathlib}
  Let \(\mathcal{F}\) be an \(\mathcal{O}_{\operatorname{Spec} R}\)-module and \(g \in R\); write
  \(\iota : \widehat{D(g)} \hookrightarrow \operatorname{Spec} R\) for the open immersion of the basic
  open and \(\operatorname{e} = \operatorname{modulesOverBasicOpenEquivalence} g\)
  (Definition~\ref{mk-def:modules_over_basicOpen_equivalence}) for this equivalence. Then its
  inverse carries the over-picture object \(\mathcal{F}.\operatorname{over} D(g)\) — a
  sheaf of modules over the localized site \((\operatorname{Opens} \operatorname{Spec}
  R).\operatorname{over} D(g)\), where the quasi-coherence presentation lives — to the honest subspace
  restriction \(\mathcal{F}.\operatorname{restrict} \iota\):
  \[
    \operatorname{e}.\operatorname{inverse}.\operatorname{obj}
      \bigl(\mathcal{F}.\operatorname{over} D(g)\bigr)
      \;\cong\; \mathcal{F}.\operatorname{restrict} \iota
    \qquad \text{in } \widehat{D(g)}.\mathrm{Modules}.
  \]
  Thus the over-picture restriction agrees with restriction to the corresponding open
  subscheme.
\end{lemma}
\begin{proof}
  \leanok
  \uses{mk-def:modules_over_basicOpen_equivalence, mk-lem:modules_restrict_basicOpen,
    mk-lem:pushforwardPushforwardEquivalence_mathlib,
    mk-lem:restrict_obj_mathlib, mk-lem:overEquivalence_isContinuous}
    Both objects describe \(\mathcal{F}\) restricted to \(D(g)\), in two equivalent
  sites. The equivalence
  \[
    (\operatorname{Opens}\operatorname{Spec}R).\operatorname{over}D(g)
      \simeq \operatorname{Opens}(\widehat{D(g)})
  \]
  is continuous by Lemma~\ref{mk-lem:overEquivalence_isContinuous}. The structure sheaf
  on the over-site and the structure sheaf of the open subscheme agree under this
  equivalence: both assign to an open \(W\subseteq\widehat{D(g)}\) the sections over
  its image \(\iota(W)\). Lemma~\ref{mk-lem:pushforwardPushforwardEquivalence_mathlib}
  therefore induces an equivalence of their module categories. Under it,
  Lemma~\ref{mk-lem:restrict_obj_mathlib} identifies the subscheme restriction
  \(\mathcal F|_{\widehat{D(g)}}\) with the over-picture object
  \(\mathcal F.\operatorname{over}D(g)\). Reading this identification in the inverse
  direction gives the asserted isomorphism.
\end{proof}

\begin{lemma}[Presentation of the affine restriction]
  \label{mk-lem:presentation_modulesRestrictBasicOpen}
  \lean{AlgebraicGeometry.presentationModulesRestrictBasicOpen,
    AlgebraicGeometry.restrictBasicOpenUnitIso,
    AlgebraicGeometry.pullbackObjUnitToUnit_isIso_basicOpen}
  \dcref{custom}

  \uses{mk-lem:presentation_over_basicOpen, mk-lem:restrict_over_compat,
    mk-def:modules_over_basicOpen_equivalence, mk-lem:presentation_ofIsIso_mathlib,
    mk-lem:modules_restrict_basicOpen}
  Let \(\mathcal{F}\) be an \(\mathcal{O}_{\operatorname{Spec} R}\)-module, let \(U\) be an open
  carrying a presentation of \(\mathcal{F}.\operatorname{over} U\), and let \(g \in R\) with
  \(D(g) \subseteq U\). Then the affine restriction \(\mathcal{F}_{(g)} =
  \operatorname{modulesRestrictBasicOpen} g\, \mathcal{F}\), as an
  \((\operatorname{Spec} R_g).\mathrm{Modules}\)-object, admits a global presentation.
\end{lemma}
\begin{proof}
  \uses{mk-lem:presentation_over_basicOpen, mk-lem:restrict_over_compat,
    mk-def:modules_over_basicOpen_equivalence, mk-lem:presentation_ofIsIso_mathlib,
    mk-lem:modules_restrict_basicOpen}
  By Lemma~\ref{mk-lem:presentation_over_basicOpen}, \(\mathcal{F}.\operatorname{over} D(g)\)
  admits a presentation. Carry this presentation through the inverse equivalence
  \(\operatorname{e}.\operatorname{inverse}\) (Definition~\ref{mk-def:modules_over_basicOpen_equivalence})
  and across the isomorphism \(\operatorname{e}.\operatorname{inverse}.\operatorname{obj}
  (\mathcal{F}.\operatorname{over} D(g)) \cong \mathcal{F}.\operatorname{restrict} \iota\)
  (Lemma~\ref{mk-lem:restrict_over_compat}) using Lemma~\ref{mk-lem:presentation_ofIsIso_mathlib},
  obtaining a presentation of the subspace restriction \(\mathcal{F}.\operatorname{restrict} \iota\) on
  \(\widehat{D(g)}\). Finally transport along the canonical affine identification \(D(g) \cong
  \operatorname{Spec} R_g\): restriction along this isomorphism preserves colimits and
  carries the structure-sheaf unit of \(\widehat{D(g)}\) to the unit of \(\operatorname{Spec} R_g\)
  — this unit compatibility holds because the identification is an isomorphism of schemes.
  The transported presentation, under
  the comparison of Lemma~\ref{mk-lem:modules_restrict_basicOpen}, is a presentation of
  \(\mathcal{F}_{(g)}\).
\end{proof}

\begin{lemma}[Degree-0/1 sheaf-axiom equalizer of a sheaf on a finite standard cover]
  \label{mk-lem:qcoh_section_equalizer}
  \lean{AlgebraicGeometry.qcoh_section_equalizer, AlgebraicGeometry.res_trans_apply}
  \dcref{custom}

  Let \(X = \Spec R\), let \(\mathcal{F}\) be an \(\mathcal{O}_X\)-module (a sheaf of modules on the
  spectral site), let \(W \subseteq X\) be an open, and let \(g_1, \ldots, g_n \in R\) be a finite family
  with \(W = \bigcup_{j=1}^{n} \bigl(W \cap D(g_j)\bigr)\). Then the augmented two-term {\v C}ech sequence
  \[
    0 \longrightarrow \Gamma(W, \mathcal{F})
      \xrightarrow{\ \rho\ } \prod_{j} \Gamma\bigl(W \cap D(g_j), \mathcal{F}\bigr)
      \xrightarrow{\ \delta\ } \prod_{j,k} \Gamma\bigl(W \cap D(g_j g_k), \mathcal{F}\bigr)
  \]
  is exact at both nonzero terms: \(\rho\) — the product of the restriction maps — is injective, and a
  family \((s_j)_j\) lies in the image of \(\rho\) iff \(\delta\bigl((s_j)_j\bigr) = 0\), where
  \(\delta\) is the difference of the two overlap-restriction maps and
  \(W \cap D(g_j g_k) = W \cap D(g_j) \cap D(g_k)\).
\end{lemma}
\begin{proof}

  This is nothing but the sheaf axiom for \(\mathcal{F}\), specialised to the finite open cover
  \(\{W \cap D(g_j)\}_{j}\) of \(W\) and read off in degrees \(0\) and \(1\). The sheaf of modules
  \(\mathcal{F}\) on \(\Spec R\) is a sheaf, so its underlying presheaf satisfies the gluing axiom for
  every open cover. For the cover \(\{W \cap D(g_j)\}_j\) of \(W\) the gluing axiom says exactly two
  things. First, a section of \(\mathcal{F}\) over \(W\) is determined by its restrictions to the
  \(W \cap D(g_j)\): this is the injectivity of \(\rho\). Second, a \emph{matching family} — a tuple
  \((s_j)_j\) whose members agree on the pairwise intersections
  \(W \cap D(g_j) \cap D(g_k)\), which for basic opens are again basic opens \(W \cap D(g_j g_k)\) —
  glues to a (necessarily unique) section over \(W\): this is the surjectivity of \(\rho\) onto
  \(\ker \delta\). Reading the two restriction-to-overlap maps as the two cofaces of the {\v C}ech nerve
  and \(\delta\) as their difference turns ``matching family'' into the equation
  \(\delta\bigl((s_j)_j\bigr) = 0\), giving the asserted exactness. Only the sheaf condition on the cover
  \(\{W \cap D(g_j)\}_j\) is used; no positive-degree {\v C}ech vanishing enters.
\end{proof}

\begin{lemma}[Localisation preserves exactness]
  \label{mk-lem:localized_module_map_exact_mathlib}
  \lean{IsLocalizedModule.map_exact}
  \mathlibok
  \dcref{custom}

  Let \(S \subseteq R\) be a submonoid of a commutative ring and let
  \(M_0 \xrightarrow{\,a\,} M_1 \xrightarrow{\,b\,} M_2\) be \(R\)-linear maps forming an exact
  sequence at \(M_1\) (the image of \(a\) equals the kernel of \(b\)). Fix, for each \(i\), a
  localisation map \(\ell_i : M_i \to M_i'\) at \(S\) (so each \(M_i'\) is the localised module
  \(S^{-1} M_i\)). Then the induced maps \(S^{-1} a : M_0' \to M_1'\) and
  \(S^{-1} b : M_1' \to M_2'\) again form an exact sequence at \(M_1'\). In other words,
  localisation at \(S\) is an exact functor; in particular it carries short exact sequences to short
  exact sequences and commutes
  with the formation of kernels.
\end{lemma}

\begin{lemma}[Base-ring descent of a localisation along an algebra map]\label{mk-lem:isLocalizedModule_powers_restrictScalars_of_algebraMap}
  \lean{AlgebraicGeometry.isLocalizedModule_powers_restrictScalars_of_algebraMap}
  \leanok
  \dcref{custom}

  Let \(A\) be an \(R\)-algebra, let \(M\) and \(N\) be \(A\)-modules that are also \(R\)-modules
  compatibly with the algebra structure (\(R \to A \to M\) and \(R \to A \to N\) are scalar towers),
  and let \(\varphi : M \to N\) be \(A\)-linear. Fix \(f \in R\) and write \(f_A\) for its image in
  \(A\) under the algebra structure map \(R \to A\). If \(\varphi\) exhibits \(N\) as the localisation of
  \(M\) at the powers of \(f_A\) (as \(A\)-modules), then
  \(\varphi\), viewed as an \(R\)-linear map via restriction of scalars along \(R \to A\), exhibits
  \(N\) as the localisation of \(M\) at the powers of \(f\) (as \(R\)-modules).
\end{lemma}
\begin{proof}
  \leanok

The three defining clauses of being a localised module transfer directly.
  Surjectivity of quotients
  and injectivity-up-to-torsion are statements about the underlying additive maps and are unchanged by
  restricting scalars. For the unit clause one must see that multiplication by \(f^k\) on \(N\) is
  bijective as an \(R\)-endomorphism: by the scalar tower \(f^k \cdot x = f_A^{\,k} \cdot x\) for every
  \(x\), the \(R\)-endomorphism \(\,x \mapsto f^k \cdot x\,\) coincides as a function with the
  \(A\)-endomorphism \(\,x \mapsto f_A^{\,k} \cdot x\,\), and bijectivity is a property of the
  underlying map independent of the choice of base ring. Hence all three clauses hold over \(R\).
\end{proof}

\begin{lemma}[Image opens of the affine tile identification]\label{mk-lem:tile_image_opens_identities}
  \lean{AlgebraicGeometry.tile_image_opens_identities}
  \leanok
  \dcref{custom}

  \uses{mk-lem:presentation_modulesRestrictBasicOpen}
  Let \(g, f \in R\), let \(R_g = R[g^{-1}]\) be the localisation of \(R\) away from \(g\), and
  let \(\iota : \Spec R_g \xrightarrow{\sim} D(g) \hookrightarrow X\) be the open immersion
  identifying \(\Spec R_g\) with the basic open \(D(g) \subseteq X = \Spec R\) (the composite of the
  canonical homeomorphism \(\Spec R_g \xrightarrow{\sim} D(g)\) with the open inclusion). Then the image
  opens of \(\iota\) of the two relevant opens of \(\Spec R_g\) are
  \[
    \iota\bigl(\Spec R_g\bigr) = D(g),
    \qquad
    \iota\bigl(D(\bar f)\bigr) = D(gf) = D(g) \cap D(f),
  \]
  where \(\bar f\) is the image of \(f\) under the localisation map \(R \to R_g\).
\end{lemma}
\begin{proof}
  \leanok
  \uses{mk-lem:presentation_modulesRestrictBasicOpen}
  The canonical homeomorphism \(\Spec R_g \xrightarrow{\sim} D(g)\) carries \(\Spec R_g\)
  isomorphically onto \(D(g)\), so the first identity is just the image of the whole space under that
  homeomorphism followed by the open inclusion. For the second, \(\bar f\) is the image of \(f\) under
  \(R \to R_g\); its basic open \(D(\bar f) \subseteq \Spec R_g\) maps under the homeomorphism to
  \(D(g) \cap D(f)\), since restricting to \(D(g)\) and then taking the locus where \(f\) is invertible
  is the locus where both \(g\) and \(f\) are invertible. Finally \(D(g) \cap D(f) = D(gf)\) since the
  basic open of a product is the intersection of the basic opens. These are equalities of opens, proved
  set-theoretically; they are not definitional.
\end{proof}

\begin{lemma}[Native action of the global-ring section functor is the structure-sheaf action]\label{mk-lem:modulesSpecToSheaf_smul_eq}
  \lean{AlgebraicGeometry.modulesSpecToSheaf_smul_eq}
  \leanok
  \dcref{custom}

  Let \(\mathcal{F}\) be an \(\mathcal{O}_{\Spec R}\)-module, let \(W \subseteq \Spec R\) be open, let
  \(r \in R\), and let \(x\) be a section over \(W\) of the global-ring section functor
  \(\Gamma_R(-,\mathcal{F})\). Then the native \(R\)-action on \(x\) coincides with the structure-sheaf
  action of the restricted global-sections image of \(r\):
  \[
    r \cdot x \;=\; c_R \cdot x_{\mathcal{F}},
    \qquad
    c_R \;=\; \rho^{W}\bigl(\theta_R(r)\bigr),
  \]
  where \(\theta_R : R \xrightarrow{\sim} \Gamma(\Spec R, \mathcal{O})\) is the global-sections
  identification, \(\rho^{W}\) is restriction of sections from \(\Spec R\) to \(W\), and \(x_{\mathcal{F}}\)
  denotes \(x\) regarded as a section of \(\mathcal{F}\) over \(W\) for its structure-sheaf module
  structure (the underlying carrier is unchanged).
\end{lemma}
\begin{proof}
  \leanok

The action of the global-ring section functor \(\Gamma_R(-,\mathcal{F})\) is, by definition, the
  restriction of scalars of \(\mathcal{F}\)'s structure-sheaf module action along the global-sections
  identification \(\theta_R\). Unfolding that definition, the \(R\)-scalar \(r\) acts as the structure-sheaf
  scalar \(\rho^{W}(\theta_R(r))\) on the unchanged underlying section; the equality therefore holds by
  definitional unfolding.
\end{proof}

\begin{lemma}[Tile action transports definitionally to the ambient structure-sheaf action]\label{mk-lem:modulesRestrictBasicOpen_smul_eq}
  \lean{AlgebraicGeometry.modulesRestrictBasicOpen_smul_eq}
  \leanok
  \dcref{custom}

  \uses{mk-lem:tile_image_opens_identities}
  Let \(\mathcal{F}\) be an \(\mathcal{O}_{\Spec R}\)-module and let \(g \in R\). Write
  \(\mathcal{F}_{(g)}\) for the restricted module (tile) over \(\Spec R_g\) obtained as the iterated
  restriction of \(\mathcal{F}\) along the affine identification \(\Spec R_g \cong D(g)\) followed by the
  open inclusion \(D(g) \hookrightarrow \Spec R\). Let \(c \in \Gamma(\top, \mathcal{O}_{\Spec R_g})\) and
  let \(m\) be a section of \(\mathcal{F}_{(g)}\) over \(\top\). Then the tile action of \(c\) on \(m\)
  transports to the ambient structure-sheaf action of \(\mathcal{F}\):
  \[
    c \cdot m \;=\; \bigl(\alpha_g^{-1}\bigl(\beta_g^{-1}(c)\bigr)\bigr) \cdot m_{\mathcal{F}},
  \]
  where \(\beta_g\) and \(\alpha_g\) are the ring isomorphisms on global sections induced respectively by
  the affine identification \(\Spec R_g \cong D(g)\) and by the open inclusion \(D(g) \hookrightarrow
  \Spec R\), and \(m_{\mathcal{F}}\) denotes \(m\) regarded as a section of \(\mathcal{F}\) over the
  iterated-image open \(W \subseteq \Spec R\) (which equals \(D(g)\) by
  Lemma~\ref{mk-lem:tile_image_opens_identities}).
\end{lemma}
\begin{proof}
  \leanok
  \uses{mk-lem:tile_image_opens_identities}
  Restriction of a module along an open immersion is the pushforward whose scalar action is reindexed by
  the inverse of the open immersion's global-sections ring map. The tile \(\mathcal{F}_{(g)}\) is the
  composite of two such restrictions --- along the affine identification \(\Spec R_g \cong D(g)\) and along
  the open inclusion \(D(g) \hookrightarrow \Spec R\) --- so its action of \(c\) is, by definition, the
  ambient \(\mathcal{F}\)-action of \(\alpha_g^{-1}(\beta_g^{-1}(c))\) on the unchanged underlying section.
  The equality holds by definitional unfolding; keeping the \(\mathcal{F}\)-side open in the
  iterated-image form \(W\) (rather than rewriting it to \(D(g)\)) is what makes the two carriers
  coincide on the nose.
\end{proof}

\begin{lemma}[Tile action transports to the ambient action over an arbitrary open]\label{mk-lem:modulesRestrictBasicOpen_smul_eq_genV}
  \lean{AlgebraicGeometry.modulesRestrictBasicOpen_smul_eq'}
  \leanok
  \dcref{custom}

  \uses{mk-lem:tile_image_opens_identities}
  The general-open companion of Lemma~\ref{mk-lem:modulesRestrictBasicOpen_smul_eq}. Let \(\mathcal{F}\)
  be an \(\mathcal{O}_{\Spec R}\)-module, let \(g \in R\), and let \(V \subseteq \Spec R_g\) be an
  \emph{arbitrary} open (the sibling treats only \(V = \top\)). For \(c \in \Gamma(V,
  \mathcal{O}_{\Spec R_g})\) and a section \(m\) of the tile \(\mathcal{F}_{(g)}\) over \(V\), the tile
  action of \(c\) on \(m\) transports to the ambient structure-sheaf action of \(\mathcal{F}\) over the
  iterated-image open \(\iota(V) \subseteq \Spec R\):
  \[
    c \cdot m \;=\; \bigl(\alpha_g^{-1}\bigl(\beta_g^{-1}(c)\bigr)\bigr) \cdot m_{\mathcal{F}},
  \]
  where \(\beta_g\) and \(\alpha_g\) are the section ring isomorphisms induced by the affine
  identification \(\Spec R_g \cong D(g)\) and by the open inclusion \(D(g) \hookrightarrow \Spec R\)
  (now read at the open \(V\) rather than at \(\top\)), and \(m_{\mathcal{F}}\) denotes \(m\) regarded
  as a section of \(\mathcal{F}\) over the iterated-image open \(\iota(V)\).
\end{lemma}
\begin{proof}
  \leanok
  \uses{mk-lem:tile_image_opens_identities}
  As in Lemma~\ref{mk-lem:modulesRestrictBasicOpen_smul_eq}, the tile \(\mathcal{F}_{(g)}\) is the
  composite of the two open-immersion restrictions, whose scalar action is by definition reindexed by
  the inverses of the two section ring maps. At an arbitrary open \(V\) this reindexing is identical, so
  the equality holds by definitional unfolding; keeping the \(\mathcal{F}\)-side open in the
  iterated-image form \(\iota(V)\) (rather than rewriting it to a basic open) is what makes the two
  carriers coincide on the nose.
\end{proof}

\begin{lemma}[Open-immersion section iso reads as a restriction]\label{mk-lem:appTop_appIso_inv_eq_res}
  \lean{AlgebraicGeometry.appTop_appIso_inv_eq_res}
  \leanok
  \dcref{custom}

  Let \(f : X \to Y\) be an open immersion of schemes. On global sections \(f\) induces the comparison
  ring map \(f^{\sharp}_{\top} : \Gamma(Y, \mathcal{O}_Y) \to \Gamma(X, \mathcal{O}_X)\), and the
  open-immersion section isomorphism
  \(\beta_f : \Gamma(X, \mathcal{O}_X) \xrightarrow{\sim} \Gamma\bigl(f(\top), \mathcal{O}_Y\bigr)\)
  identifies the global sections of \(X\) with the sections of \(Y\) over the image open
  \(f(\top) \subseteq Y\). Then the composite of \(f^{\sharp}_{\top}\) with \(\beta_f^{-1}\) is the
  structure-sheaf restriction of \(Y\) from \(\top\) down to the image open:
  \[
    \beta_f^{-1} \circ f^{\sharp}_{\top} \;=\; \rho^{f(\top)}_{Y}.
  \]
  Equivalently, under the section identification \(\beta_f\) the global-sections comparison map of \(f\)
  is nothing but restriction along \(f(\top) \hookrightarrow Y\).
\end{lemma}
\begin{proof}
  \leanok

The forward map of \(\beta_f\) is, by construction, the comparison map \(f^{\sharp}\) read on the image
  open. Rewriting the identity in the form \(f^{\sharp}_{\top} = \rho^{f(\top)}_Y \circ \beta_f\) and
  applying the naturality of the comparison maps \(f^{\sharp}\) with respect to the restriction maps of
  \(\mathcal{O}_X\) and \(\mathcal{O}_Y\), the right-hand side reduces to \(f^{\sharp}_{\top}\)
  precomposed with the identity restriction of \(\mathcal{O}_X\) on \(\top\); the source open \(\top\)
  and its image are matched by the thin-category uniqueness of open inclusions. This leaves exactly
  \(\rho^{f(\top)}_Y\).
\end{proof}

\begin{lemma}[$\Gamma$--$\Spec$ naturality of the localisation immersion, section form]\label{mk-lem:key_morph}
  \lean{AlgebraicGeometry.key_morph}
  \leanok
  \dcref{custom}

  \uses{mk-lem:appTop_appIso_inv_eq_res}
  Let \(g \in R\), write \(R_g = R[g^{-1}]\), let \(\lambda : R \to R_g\) be the localisation map, and let
  \(\iota = \Spec(\lambda) : \Spec R_g \to \Spec R\) be the induced affine morphism (it is the open
  immersion identifying \(\Spec R_g\) with \(D(g) = \iota(\top)\)). Then the restriction to the image
  open \(D(g)\) of the global-sections identification
  \(\theta_R : R \xrightarrow{\sim} \Gamma(\Spec R, \mathcal{O})\) factors as the localisation map
  followed by the global-sections identification of \(\Spec R_g\) and the (inverse) section isomorphism
  of \(\iota\):
  \[
    \rho^{D(g)} \circ \theta_R
    \;=\;
    \beta_\iota^{-1} \circ \theta_{R_g} \circ \lambda,
  \]
  where \(\theta_{R_g} : R_g \xrightarrow{\sim} \Gamma(\Spec R_g, \mathcal{O})\) is the global-sections
  identification of \(\Spec R_g\) and \(\beta_\iota\) is the open-immersion section isomorphism of
  \(\iota\) at \(\top\).
\end{lemma}
\begin{proof}
  \leanok
  \uses{mk-lem:appTop_appIso_inv_eq_res}
  This is the naturality of the \(\Spec\)-to-global-sections comparison applied to the ring map
  \(\lambda\). Concretely, the Mathlib naturality square for the global-sections identification (the
  \(\Gamma\)-\(\Spec\) adjunction unit, \texttt{ΓSpecIso} naturality) gives
  \(\lambda \) followed by \(\theta_{R_g}\) equals \(\theta_R\) followed by the comparison map
  \(\iota^{\sharp}_{\top}\) of \(\iota = \Spec(\lambda)\). Substituting this into the right-hand side and
  reading the composite of \(\iota^{\sharp}_{\top}\) with \(\beta_\iota^{-1}\) as the restriction to the
  image open \(D(g)\) by Lemma~\ref{mk-lem:appTop_appIso_inv_eq_res} yields \(\rho^{D(g)} \circ \theta_R\).
\end{proof}

\begin{lemma}[Composition of the two tile section isomorphisms]\label{mk-lem:tile_appIso_comp}
  \lean{AlgebraicGeometry.tile_appIso_comp}
  \leanok
  \dcref{custom}

  In the situation of Lemma~\ref{mk-lem:key_morph} the localisation immersion \(\iota\) factors as the
  composite of the affine identification \(\Spec R_g \xrightarrow{\sim} D(g)\) (of schemes) with the open
  inclusion \(D(g) \hookrightarrow \Spec R\). The (inverse) open-immersion section isomorphisms of these
  two factors compose into the single (inverse) section isomorphism \(\beta_\iota^{-1}\) of \(\iota\) at
  \(\top\), up to a structure-sheaf transport along the equality of image opens
  \(\iota(\top) = D(g)\):
  \[
    \beta_{\hookrightarrow}^{-1} \circ \beta_{\mathrm{aff}}^{-1}
    \;=\;
    \tau \circ \beta_\iota^{-1},
  \]
  where \(\beta_{\mathrm{aff}}\) and \(\beta_{\hookrightarrow}\) are the section isomorphisms of the
  affine identification and of the open inclusion respectively, and \(\tau\) is the structure-sheaf
  transport induced by the equality of the composed image open with \(D(g)\).
\end{lemma}
\begin{proof}
  \leanok

The functoriality of the open-immersion section isomorphism under composition of open immersions
  (\texttt{comp\_appIso}) expresses \(\beta_\iota\) as the composite of the two factor isomorphisms, with
  the section isomorphism of the open inclusion being the identity. Folding the resulting transport
  through the equality of image opens \(\iota(\top) = D(g)\) gives the displayed identity, with \(\tau\)
  the structure-sheaf map of that equality of opens.
\end{proof}

\begin{lemma}[Structure-sheaf ring identity of the tile comparison, morphism form]\label{mk-lem:tile_section_ring_identity}
  \lean{AlgebraicGeometry.tile_section_ring_identity}
  \leanok
  \dcref{custom}

  \uses{mk-lem:key_morph, mk-lem:tile_appIso_comp}
  In the situation of Lemma~\ref{mk-lem:key_morph}, the restriction to the tile image open \(D(g)\) of the
  global-sections identification \(\theta_R\) of \(\Spec R\) equals the localisation map followed by the
  global-sections identification of \(\Spec R_g\) and the two (inverse) open-immersion section
  isomorphisms of the tile:
  \[
    \rho^{D(g)} \circ \theta_R
    \;=\;
    \beta_{\hookrightarrow}^{-1} \circ \beta_{\mathrm{aff}}^{-1} \circ \theta_{R_g} \circ \lambda.
  \]
  This is the morphism-level (ring-map) form of the residual ring identity of
  Lemma~\ref{mk-lem:tile_section_comparison}.
\end{lemma}
\begin{proof}
  \leanok
  \uses{mk-lem:key_morph, mk-lem:tile_appIso_comp}
  Substitute the \(\Gamma\)-\(\Spec\) naturality identity of Lemma~\ref{mk-lem:key_morph}, which rewrites
  \(\rho^{D(g)} \circ \theta_R\) as \(\beta_\iota^{-1} \circ \theta_{R_g} \circ \lambda\), and then
  replace the single section isomorphism \(\beta_\iota^{-1}\) by the composite of the two tile factor
  isomorphisms using Lemma~\ref{mk-lem:tile_appIso_comp}. The structure-sheaf transport \(\tau\) of
  Lemma~\ref{mk-lem:tile_appIso_comp} merges with the restriction maps in the thin category of opens
  (functoriality of the structure presheaf), and the two displayed sides coincide.
\end{proof}

\begin{lemma}
[Structure-sheaf ring identity of the tile comparison at an arbitrary open]
  \label{mk-lem:tile_section_ring_identity_genV}
  \lean{AlgebraicGeometry.tile_section_ring_identity', AlgebraicGeometry.appIso_inv_res,
    AlgebraicGeometry.appIso_inv_res_assoc}
  \dcref{custom}

  \uses{mk-lem:tile_section_ring_identity}
  The general-open companion of Lemma~\ref{mk-lem:tile_section_ring_identity}. For an \emph{arbitrary}
  open \(V \subseteq \Spec R_g\), the restriction to the iterated-image open \(\iota(V)\) of the
  global-sections identification \(\theta_R\) of \(\Spec R\) equals the localisation map followed by the
  global-sections identification of \(\Spec R_g\), the restriction to \(V\), and the two (inverse)
  open-immersion section isomorphisms of the tile read at \(V\):
  \[
    \rho^{\iota(V)} \circ \theta_R
    \;=\;
    \beta_{\hookrightarrow,V}^{-1} \circ \beta_{\mathrm{aff},V}^{-1} \circ \rho^{V}_{R_g} \circ \theta_{R_g} \circ \lambda.
  \]
  This supplies the ring identity at the target open \(V = D(\bar f)\) consumed by the general-open
  scalar reconciliation Lemma~\ref{mk-lem:tile_scalar_compat_genV}.
\end{lemma}
\begin{proof}
  \uses{mk-lem:tile_section_ring_identity}
  Derive from the \(V = \top\) case (Lemma~\ref{mk-lem:tile_section_ring_identity}) by post-composing both
  sides with the structure-sheaf restriction \(\iota(V) \le \iota(\top)\) and pushing that restriction
  through the two open-immersion section isomorphisms. The push-through is the section-restriction form
  of the open-immersion section-iso naturality: for an open immersion \(f\) and opens \(U' \le U\), the
  inverse section iso at \(U\) followed by the \(Y\)-restriction \(f(U') \le f(U)\) equals the
  \(X\)-restriction \(U' \le U\) followed by the inverse section iso at \(U'\). The residual restriction
  \(\rho^{\top}_{R_g}\) on \(\Spec R_g\) is the identity by
  uniqueness of morphisms in the thin category of opens.
\end{proof}

\begin{lemma}[Scalar compatibility of the tile comparison at the source open]\label{mk-lem:tile_scalar_compat}
  \lean{AlgebraicGeometry.tile_scalar_compat}
  \leanok
  \dcref{custom}

  \uses{mk-lem:modulesSpecToSheaf_smul_eq, mk-lem:modulesRestrictBasicOpen_smul_eq,
    mk-lem:tile_section_ring_identity}
  Let \(\mathcal{F}\) be a quasi-coherent \(\mathcal{O}_{\Spec R}\)-module, let \(g, r \in R\), and let
  \(x\) be a section of the affine tile \(\mathcal{F}_{(g)}\) over \(\top\) --- equivalently, regarding
  the underlying carrier, a section of \(\mathcal{F}\) over the tile image open \(D(g)\). Then the native
  \(R\)-action of \(r\) on \(x\) (through \(\Gamma_R(-,\mathcal{F})\)) coincides with the \(R_g\)-action
  of the localised scalar \(\lambda(r) \in R_g\) on \(x\) (through the tile
  \(\Gamma_{R_g}(-, \mathcal{F}_{(g)})\)):
  \[
    r \cdot x \;=\; \lambda(r) \cdot x,
    \qquad \lambda : R \to R_g \text{ the localisation map}.
  \]
  This is the scalar core, at the source open \(V = \top\), of the section comparison of
  Lemma~\ref{mk-lem:tile_section_comparison}; it is the scalar-tower compatibility \(R \to R_g\) on the
  common underlying section carrier. It is \emph{not} the full natural isomorphism, only the equality of
  the two scalar actions on the single open \(V = \top\).
\end{lemma}
\begin{proof}
  \leanok
  \uses{mk-lem:modulesSpecToSheaf_smul_eq, mk-lem:modulesRestrictBasicOpen_smul_eq,
    mk-lem:tile_section_ring_identity}
    By Lemma~\ref{mk-lem:modulesSpecToSheaf_smul_eq} the native \(R\)-action of \(r\) on \(x\) is the
  structure-sheaf action of \(\rho^{D(g)}(\theta_R(r))\), and by
  Lemma~\ref{mk-lem:modulesRestrictBasicOpen_smul_eq} the \(R_g\)-action of \(\lambda(r)\) on \(x\) (as a
  tile section) transports to the structure-sheaf action of
  \(\beta_{\hookrightarrow}^{-1}(\beta_{\mathrm{aff}}^{-1}(\theta_{R_g}(\lambda(r))))\) on the same
  underlying carrier. Both actions therefore reduce to a structure-sheaf scalar action, and the two
  scalars agree by the morphism-level ring identity of Lemma~\ref{mk-lem:tile_section_ring_identity}
  evaluated at \(r\).
\end{proof}

\begin{lemma}[Scalar compatibility of the tile comparison at an arbitrary open]\label{mk-lem:tile_scalar_compat_genV}
  \lean{AlgebraicGeometry.tile_scalar_compat'}
  \leanok
  \dcref{custom}

  \uses{mk-lem:modulesSpecToSheaf_smul_eq, mk-lem:modulesRestrictBasicOpen_smul_eq_genV,
    mk-lem:tile_section_ring_identity_genV}
  The general-open companion of Lemma~\ref{mk-lem:tile_scalar_compat}. Let \(\mathcal{F}\) be a
  quasi-coherent \(\mathcal{O}_{\Spec R}\)-module, let \(g, r \in R\), let \(V \subseteq \Spec R_g\) be
  an \emph{arbitrary} open, and let \(x\) be a section of the affine tile \(\mathcal{F}_{(g)}\) over
  \(V\) --- equivalently, on the underlying carrier, a section of \(\mathcal{F}\) over the
  iterated-image open \(\iota(V)\). Then the native \(R\)-action of \(r\) on \(x\) coincides with the
  \(R_g\)-action of the localised scalar \(\lambda(r) \in R_g\) on \(x\):
  \[
    r \cdot x \;=\; \lambda(r) \cdot x,
    \qquad \lambda : R \to R_g \text{ the localisation map}.
  \]
  The instance at \(V = D(\bar f)\) is the scalar-tower compatibility \(R \to R_g\) at the \emph{target}
  open of the per-tile section localization (Lemma~\ref{mk-lem:tile_section_localization}); the
  \(V = \top\) case is Lemma~\ref{mk-lem:tile_scalar_compat}.
\end{lemma}
\begin{proof}
  \leanok
  \uses{mk-lem:modulesSpecToSheaf_smul_eq, mk-lem:modulesRestrictBasicOpen_smul_eq_genV,
    mk-lem:tile_section_ring_identity_genV}
    The same argument as Lemma~\ref{mk-lem:tile_scalar_compat}, now at an arbitrary open \(V\). By
  Lemma~\ref{mk-lem:modulesSpecToSheaf_smul_eq} the native \(R\)-action of \(r\) on \(x\) reduces to a
  structure-sheaf scalar action, and by Lemma~\ref{mk-lem:modulesRestrictBasicOpen_smul_eq_genV} the
  \(R_g\)-action of \(\lambda(r)\) on \(x\) (as a tile section over \(V\)) reduces to a structure-sheaf
  scalar action on the same underlying carrier. The two scalars then agree by the morphism-level ring
  identity of Lemma~\ref{mk-lem:tile_section_ring_identity_genV} evaluated at \(r\).
\end{proof}

\begin{lemma}
[Natural section comparison between the tile and the ambient sheaf]
  \label{mk-lem:tile_section_comparison}
  \dcref{custom}
  \uses{mk-lem:presentation_modulesRestrictBasicOpen, mk-lem:restrict_obj_mathlib,
    mk-lem:tile_image_opens_identities, mk-lem:modulesSpecToSheaf_smul_eq,
    mk-lem:modulesRestrictBasicOpen_smul_eq, mk-lem:tile_section_ring_identity,
    mk-lem:tile_scalar_compat}
  In the situation of Lemma~\ref{mk-lem:tile_image_opens_identities}, with
  \(\mathcal{F}_{(g)}\) the restriction of the quasi-coherent module
  \(\mathcal{F}\) along \(\iota : \Spec R_g \cong D(g) \hookrightarrow X\), there is an \(R_g\)-linear
  isomorphism, natural in the open \(V \subseteq \Spec R_g\) and intertwining the restriction maps of the
  two sheaves,
  \[
    \Gamma_{R_g}\bigl(V, \mathcal{F}_{(g)}\bigr) \;\cong\; \Gamma_R\bigl(\iota(V), \mathcal{F}\bigr).
  \]
  Here \(\Gamma_{R_g}(-, \mathcal{F}_{(g)})\) is the global-ring section functor attached to the
  presented tile over \(R_g\) (valued in \(R_g\)-modules), and \(\Gamma_R(-, \mathcal{F})\) the
  corresponding functor for \(\mathcal{F}\) over \(R\).
\end{lemma}
\begin{proof}
  \uses{mk-lem:presentation_modulesRestrictBasicOpen, mk-lem:restrict_obj_mathlib,
    mk-lem:tile_image_opens_identities, mk-lem:modulesSpecToSheaf_smul_eq,
    mk-lem:modulesRestrictBasicOpen_smul_eq, mk-lem:tile_section_ring_identity,
    mk-lem:tile_scalar_compat}
  By Lemma~\ref{mk-lem:restrict_obj_mathlib}, both sides have the same underlying additive
  group, and these identifications commute with restriction maps. It remains to compare
  scalar multiplication. For \(r\in R\), naturality of global sections for the
  localization morphism gives
  \[
    \rho^{D(g)}\bigl(\theta_R(r)\bigr)
      \;=\;
    \beta_g^{-1}\bigl(\theta_{R_g}(\bar r)\bigr),
    \qquad \bar r \;=\; \text{image of } r \text{ in } R_g,
  \]
  where \(\theta_R\) and \(\theta_{R_g}\) are the affine global-section
  identifications and \(\beta_g\) is induced by
  \(\operatorname{Spec}R_g\cong D(g)\). This is precisely the compatibility of
  \(\Gamma(D(g),\mathcal O)\cong R_g\) with the localization map \(R\to R_g\), as
  recorded in Lemma~\ref{mk-lem:tile_section_ring_identity}. Together with
  Lemmas~\ref{mk-lem:modulesSpecToSheaf_smul_eq} and
  \ref{mk-lem:modulesRestrictBasicOpen_smul_eq}, it shows that the additive comparison
  respects the \(R_g\)-action. Hence it is an \(R_g\)-linear isomorphism, and its
  compatibility with restrictions makes it natural in \(V\).
\end{proof}

\begin{lemma}[Scalar tower on a bundled restriction-of-scalars module]\label{mk-lem:isScalarTower_restrictScalars_obj}
  \lean{AlgebraicGeometry.isScalarTower_restrictScalars_obj}
  \leanok
  \dcref{custom}

  Let \(R \to S\) be a ring map and let \(M\) be an \(S\)-module. Then the \(R\)-module obtained from
  \(M\) by restriction of scalars along \(R \to S\), written
  \((\mathrm{restrictScalars}_{R \to S})\,M\), carries a scalar-tower structure: the restricted
  \(R\)-action factors through the \(S\)-action via the algebra map \(R \to S\).
\end{lemma}
\begin{proof}
  \leanok

By definition of restriction of scalars, the \(R\)-action on \(M\) is given by
  \(r \cdot m = \phi(r) \cdot m\) where \(\phi : R \to S\) is the algebra map. This is precisely the
  scalar-tower identity: the \(R\)-action factors through the \(S\)-action via \(\phi\). Hence \(R\),
  \(S\), and \(\phi\) form a scalar tower on \((\mathrm{restrictScalars}_{R \to S})\,M\).
\end{proof}

\begin{lemma}[Reconciliation equivalence for the tile section comparison]
  \label{mk-lem:tileReconcileEquiv}
  \lean{AlgebraicGeometry.tileReconcileEquiv, AlgebraicGeometry.tileReconcileEquiv_apply,
    AlgebraicGeometry.tileReconcileEquiv_symm_apply}
  \dcref{custom}

  \uses{mk-lem:tile_scalar_compat_genV, mk-lem:modulesSpecToSheaf_smul_eq}
  Let \(\mathcal{F}\) be a quasi-coherent \(\mathcal{O}_{\Spec R}\)-module, let \(g \in R\), write \(R_g\)
  for the localisation of \(R\) away from \(g\), and let \(V \subseteq \Spec R_g\) be an open. Let
  \(\iota : \Spec R_g \xrightarrow{\sim} D(g) \hookrightarrow \Spec R\) be the localisation open
  immersion, so that \(\iota(V)\) is the iterated image of \(V\) under the affine identification
  \(\Spec R_g \cong D(g)\) followed by the open inclusion \(D(g) \hookrightarrow \Spec R\). Then there is
  an \(R\)-linear isomorphism
  \[
    (\mathrm{restrictScalars}_{R \to R_g})\, \Gamma_{R_g}\bigl(V, \mathcal{F}_{(g)}\bigr)
      \;\xrightarrow{\ \sim\ }\;
      \Gamma_R\bigl(\iota(V),\ \mathcal{F}\bigr),
  \]
  which is the identity on underlying elements. The forward and inverse maps are each the identity on
  elements, so the only content is the reconciliation of the two module structures carried by the
  common section type.
\end{lemma}
\begin{proof}
  \uses{mk-lem:tile_scalar_compat_genV, mk-lem:modulesSpecToSheaf_smul_eq}
  The two section modules share one underlying carrier: the tile sections over \(V\) coincide as
  sets with the sections of \(\mathcal{F}\) over the image open \(\iota(V)\), by construction of the
  restriction functor. The underlying map is therefore the identity, and it is additive. Its \(R\)-linearity is
  exactly the general-open scalar-tower compatibility of Lemma~\ref{mk-lem:tile_scalar_compat_genV}: the
  native \(R\)-action on the \(\mathcal{F}\)-side (described by
  Lemma~\ref{mk-lem:modulesSpecToSheaf_smul_eq}) coincides with the \(R_g\)-action read through the
  localisation map \(R \to R_g\), i.e.\ the restriction-of-scalars action carried by the tile side. The
  identity map is thus \(R\)-linear with \(R\)-linear inverse, hence an \(R\)-linear isomorphism.
\end{proof}

\begin{lemma}

[Per-tile section localisation at \(f\)]
  \label{mk-lem:tile_section_localization}
  \lean{AlgebraicGeometry.tile_section_localization, AlgebraicGeometry.tile_restrict_map_apply}
  \dcref{custom}

  \uses{mk-lem:presentation_modulesRestrictBasicOpen,
    mk-lem:section_isLocalizedModule_of_presentation, mk-lem:restrict_obj_mathlib,
    mk-lem:isLocalizedModule_powers_restrictScalars_of_algebraMap,
    mk-lem:tile_image_opens_identities, mk-lem:tile_scalar_compat, mk-lem:tile_scalar_compat_genV,
    mk-lem:tileReconcileEquiv, mk-lem:isScalarTower_restrictScalars_obj,
    mk-lem:modulesSpecToSheaf_smul_eq, mk-lem:modulesRestrictBasicOpen_smul_eq}
  Let \(X = \Spec R\), let \(\mathcal{F}\) be a quasi-coherent \(\mathcal{O}_X\)-module, let \(f \in R\),
  and let \(g \in R\) be either a cover element \(g = g_j\) or an overlap product \(g = g_j g_k\) of the
  finite presentation cover of Lemma~\ref{mk-lem:qcoh_finite_presentation_cover}, so that the restriction \(\mathcal{F}_{(g)}\) of
  \(\mathcal{F}\) to \(\Spec R_g\) admits a global presentation. Then the restriction-of-sections map
  \(\Gamma(D(g), \mathcal{F}) \to \Gamma(D(gf), \mathcal{F})\) exhibits its target as the localisation of
  its source at the powers of \(f\); equivalently, there is a canonical \(R\)-linear isomorphism
  \(\Gamma(D(g), \mathcal{F})_f \cong \Gamma(D(gf), \mathcal{F})\).
\end{lemma}
\begin{proof}
  \uses{mk-lem:presentation_modulesRestrictBasicOpen, mk-lem:section_isLocalizedModule_of_presentation,
    mk-lem:restrict_obj_mathlib, mk-lem:isLocalizedModule_powers_restrictScalars_of_algebraMap,
    mk-lem:tile_image_opens_identities, mk-lem:tile_scalar_compat_genV, mk-lem:tileReconcileEquiv,
    mk-lem:isScalarTower_restrictScalars_obj,
    mk-lem:modulesSpecToSheaf_smul_eq, mk-lem:modulesRestrictBasicOpen_smul_eq}
  Let \(\iota : \Spec R_g \xrightarrow{\sim} D(g) \hookrightarrow X\) be the open immersion of
  Lemma~\ref{mk-lem:tile_image_opens_identities}, with \(\bar f\) the image of \(f\) under the
  localisation map \(R \to R_g\). We compare the section maps after restriction of
  scalars from \(R_g\) to \(R\).

  \emph{Step 1 (global presentation of the tile).} Since \(D(g)\) lies in a presentation member of the
  cover, Lemma~\ref{mk-lem:presentation_modulesRestrictBasicOpen} presents the tile
  \(\mathcal{F}_{(g)}\) globally over \(\Spec R_g\).

  \emph{Step 2 (localisation over \(R_g\)).} Apply
  Lemma~\ref{mk-lem:section_isLocalizedModule_of_presentation} to the globally presented
  \(\mathcal{F}_{(g)}\) and \(\bar f\): the section-restriction map
  \(\Gamma_{R_g}(\Spec R_g, \mathcal{F}_{(g)}) \to \Gamma_{R_g}(D(\bar f), \mathcal{F}_{(g)})\)
  exhibits its target as a localisation at the powers of \(\bar f\), \emph{as \(R_g\)-modules}. Note this
  is over \(R_g\), not yet over \(R\).

  \emph{Step 3 (opens identities).} By Lemma~\ref{mk-lem:tile_image_opens_identities} the two image opens of
  \(\iota\) are \(\iota(\Spec R_g) = D(g)\) and \(\iota(D(\bar f)) = D(gf) = D(g) \cap D(f)\).

  \emph{Step 4 (the section comparison, as a restriction of scalars).} The localisation of Step 2 lives
  over \(R_g\); it must be related to the section-restriction map of \(\mathcal{F}\) over \(R\). The two
  maps share the same underlying map of section sets: this is the carrier identification of
  Lemma~\ref{mk-lem:restrict_obj_mathlib}, now read as an equality of \(R\)-linear maps (with the
  \(\mathcal{F}\)-side open kept in its iterated-image form \(W = \iota(V)\)). The one genuine difference
  is the base ring: the tile sections carry an \(R_g\)-module structure, while the
  \(\mathcal{F}\)-sections carry an \(R\)-module structure. The descent that reconciles them is precisely
  the \emph{restriction of scalars} along the localisation map \(R \to R_g\): each tile section module,
  viewed through restriction of scalars, becomes an \(R\)-module, and the Step-2 map so viewed is the
  section-restriction map \(\Gamma_R(D(g), \mathcal{F}) \to \Gamma_R(D(gf), \mathcal{F})\).

  For this restriction of scalars to be meaningful, the native \(R\)-action on each section must agree
  with its \(R_g\)-action read through \(R \to R_g\); this is the scalar-tower compatibility. At the
  source open \(V = \top\) (image \(D(g)\)) it is exactly Lemma~\ref{mk-lem:tile_scalar_compat}, whose two
  scalar prescriptions agree by the scalar-reduction identities
  (Lemmas~\ref{mk-lem:modulesSpecToSheaf_smul_eq} and \ref{mk-lem:modulesRestrictBasicOpen_smul_eq}). At the
  target open \(V = D(\bar f)\) (image \(D(gf)\)) it is the general-open companion
  Lemma~\ref{mk-lem:tile_scalar_compat_genV}, the \(V = D(\bar f)\) instance of the same ring-naturality
  content. Both opens of the Step-2 map are thus covered. In the descent these two scalar
  reconciliations are packaged uniformly as the \(R\)-linear reconciliation isomorphism of
  Lemma~\ref{mk-lem:tileReconcileEquiv} (instantiated at the source and target opens), whose
  \(R\)-linearity is supplied by the general-open compatibility Lemma~\ref{mk-lem:tile_scalar_compat_genV};
  the descent transports the Step-2 localisation through these isomorphisms.

  Reading Step 2's localisation-at-\(\bar f\) property through restriction of scalars along
  \(R \to R_g\), and matching source and target opens by Step 3, yields: the map
  \(\Gamma_R(D(g), \mathcal{F}) \to \Gamma_R(D(gf), \mathcal{F})\) exhibits its target as a localisation
  at the powers of \(\bar f\), \emph{still as \(R_g\)-modules}.

  \emph{Step 5 (base-ring descent to \(R\)).} Finally descend the base ring from \(R_g\) to \(R\) by
  Lemma~\ref{mk-lem:isLocalizedModule_powers_restrictScalars_of_algebraMap}, applied with \(A = R_g\) and the
  scalar tower \(R \to R_g\) on the two section modules (supplied structurally on the bundled
  restriction-of-scalars carriers by Lemma~\ref{mk-lem:isScalarTower_restrictScalars_obj}): since \(\bar f\) is the image of \(f\) under
  \(R \to R_g\), the powers of \(\bar f\) descend to the powers of \(f\), and the restriction of scalars
  (from \(R_g\) to \(R\)) of the Step-4 map exhibits \(\Gamma_R(D(gf), \mathcal{F})\) as the localisation
  of \(\Gamma_R(D(g), \mathcal{F})\) at the powers of \(f\) \emph{over \(R\)}. Equivalently there is a
  canonical \(R\)-linear isomorphism \(\Gamma(D(g), \mathcal{F})_f \cong \Gamma(D(gf), \mathcal{F})\).

  Crucially this localisation is carried out entirely on the tile \(\mathcal{F}_{(g)}\), where
  \(\mathcal{F}\) is associated to a module (the tile is globally presented, hence of the form
  \(\widetilde{M_g}\)), and never on
  the global object \(\mathcal{F}\). This is exactly the structural feature that keeps the keystone
  Lemma~\ref{mk-lem:qcoh_section_isLocalizedModule} non-circular: the keystone-at-\(g\) statement is
  \emph{not} invoked; only the presented-tile lemma is.
\end{proof}

\begin{lemma}[Per-overlap section localisation at \(f\)]
  \label{mk-lem:overlap_section_localization}
  \lean{AlgebraicGeometry.overlap_section_localization, AlgebraicGeometry.overlap_target_eq,
    AlgebraicGeometry.presheaf_map_comp₂_apply}
  \dcref{custom}

  \uses{mk-lem:tile_section_localization}
  Let \(X = \Spec R\), let \(\mathcal{F}\) be a quasi-coherent \(\mathcal{O}_X\)-module, let \(f \in R\),
  and let \(a, b \in R\) be such that \(D(a)\) lies in a cover member of
  Lemma~\ref{mk-lem:qcoh_finite_presentation_cover} carrying a global presentation (so the product \(a b\) is
  an overlap of that cover). Then the section-restriction map
  \(\Gamma(D(a) \cap D(b), \mathcal{F}) \to \Gamma(D(af) \cap D(bf), \mathcal{F})\) exhibits its target as
  the localisation of its source at the powers of \(f\); equivalently, there is a canonical \(R\)-linear
  isomorphism \(\Gamma(D(a) \cap D(b), \mathcal{F})_f \cong \Gamma(D(af) \cap D(bf), \mathcal{F})\).
\end{lemma}
\begin{proof}
  \uses{mk-lem:tile_section_localization}
  This is the per-tile localisation Lemma~\ref{mk-lem:tile_section_localization} taken at the single element
  \(g = a b\), transported along the two basic-open lattice identities \(D(ab) = D(a) \cap D(b)\) and
  \(D((ab)f) = D(af) \cap D(bf)\). Applying the per-tile lemma to \(g = a b\) gives that the restriction
  \(\Gamma(D(ab), \mathcal{F}) \to \Gamma(D(abf), \mathcal{F})\) localises at the powers of \(f\). The
  source identity rewrites \(D(ab)\) as \(D(a) \cap D(b)\); the target identity rewrites \(D(abf) =
  D((ab)f)\) as \(D(af) \cap D(bf)\) (both are the closed-under-products meet of the corresponding basic
  opens). Transporting the localisation map along the induced isomorphisms of section modules — which
  fold the three intervening restriction maps into the single restriction along
  \(D(af) \cap D(bf) \subseteq D(a) \cap D(b)\) — preserves the
  \(\operatorname{IsLocalizedModule}\) property, yielding the claim.
\end{proof}

\begin{lemma}[Kernel comparison: a left-exact ladder localises on the left]\label{mk-lem:isLocalizedModule_of_exact}
  \lean{AlgebraicGeometry.isLocalizedModule_of_exact}
  \leanok
  \dcref{custom}

  \uses{mk-lem:localized_module_map_exact_mathlib}
  Let \(S \subseteq R\) be a submonoid of a commutative ring, and consider a commutative ladder of
  \(R\)-modules
  \[
    \begin{array}{ccccc}
      A & \xrightarrow{\ i\ } & B & \xrightarrow{\ p\ } & C \\
      \big\downarrow{\,a} & & \big\downarrow{\,b} & & \big\downarrow{\,c} \\
      A' & \xrightarrow{\ i'\ } & B' & \xrightarrow{\ p'\ } & C'
    \end{array}
  \]
  in which both rows are left-exact (\(i\) and \(i'\) injective, with \(\operatorname{im} i = \ker p\) and
  \(\operatorname{im} i' = \ker p'\)). Suppose the two right-hand verticals \(b : B \to B'\) and
  \(c : C \to C'\) are localisation maps at \(S\) (\(\operatorname{IsLocalizedModule} S\, b\) and
  \(\operatorname{IsLocalizedModule} S\, c\)). Then the left vertical \(a : A \to A'\) is a localisation map
  at \(S\) (\(\operatorname{IsLocalizedModule} S\, a\)). This is the converse direction of the Mathlib fact
  that localisation preserves exactness (Lemma~\ref{mk-lem:localized_module_map_exact_mathlib}): there, the
  verticals being localisations forces the rows to stay exact; here, the rows being exact lets a
  localisation on the two right columns be transported to the left column.
\end{lemma}
\begin{proof}
  \leanok

We verify the three defining clauses of \(\operatorname{IsLocalizedModule} S\, a\) by a diagram chase,
  using only the left-exactness of the two rows and the localisation hypotheses on \(b\) and \(c\).

  \emph{(i) Each \(s \in S\) acts invertibly on \(A'\).} Injectivity of \(s \cdot -\) on \(A'\): if
  \(s \cdot x = s \cdot y\) then \(s \cdot i'(x) = s \cdot i'(y)\), so \(i'(x) = i'(y)\) because
  \(s \cdot -\) is bijective on \(B'\); injectivity of
  \(i'\) gives \(x = y\). Surjectivity of \(s \cdot -\) onto \(A'\): given \(x \in A'\), pick
  \(w \in B'\) with \(s \cdot w = i'(x)\) (bijectivity of \(s \cdot -\) on \(B'\)); then
  \(s \cdot p'(w) = p'(i'(x)) = 0\), and since \(s \cdot -\) is injective on \(C'\) we get \(p'(w) = 0\),
  so left-exactness of the bottom row yields \(x' \in A'\) with \(i'(x') = w\), whence
  \(i'(s \cdot x') = s \cdot w = i'(x)\) and \(s \cdot x' = x\). Thus \(s \cdot -\) is bijective on \(A'\).

  \emph{(ii) Every \(y \in A'\) is \(s^{-1} \cdot a(x)\) for some \(x \in A\), \(s \in S\).} Map \(y\) into
  \(B'\) via \(i'\); the surjectivity clause of \(b\) gives \(w \in B\) and \(s \in S\) with
  \(b(w) = s \cdot i'(y)\). Then \(c(p(w)) = p'(b(w)) = s \cdot p'(i'(y)) = 0 = c(0)\), so since
  \(c\) is a localisation map there exists \(t \in S\) with \(t \cdot p(w) = 0\), i.e.\
  \(p(t \cdot w) = 0\). Left-exactness of the top row yields \(x \in A\) with \(i(x) = t \cdot w\). Then
  \(i'(a(x)) = b(i(x)) = t \cdot b(w) = (t s) \cdot i'(y) = i'((t s) \cdot y)\), and injectivity of \(i'\)
  gives \(a(x) = (t s) \cdot y\), i.e.\ \(y = (ts)^{-1} \cdot a(x)\).

  \emph{(iii) If \(a(x) = a(x')\) then \(s \cdot x = s \cdot x'\) for some \(s \in S\).} From
  \(a(x) = a(x')\) we get \(b(i(x)) = i'(a(x)) = i'(a(x')) = b(i(x'))\); since \(b\) is a
  localisation map there exists \(s \in S\) with \(s \cdot i(x) = s \cdot i(x')\),
  i.e.\ \(i(s \cdot x) = i(s \cdot x')\), and injectivity of \(i\) gives \(s \cdot x = s \cdot x'\).

  The three clauses are exactly \(\operatorname{IsLocalizedModule} S\, a\).
\end{proof}

\begin{lemma}[Kernel comparison: the localised global equalizer matches the \(D(f)\)-cover equalizer]\label{mk-lem:qcoh_section_kernel_comparison}
  \lean{AlgebraicGeometry.qcoh_section_kernel_comparison}
  \leanok
  \dcref{custom}

  \uses{mk-lem:qcoh_section_isLocalizedModule, mk-lem:isLocalizedModule_linearEquiv_mathlib}
  Let \(X = \Spec R\), let \(\mathcal{F}\) be a quasi-coherent \(\mathcal{O}_X\)-module, let \(f \in R\),
  and let \(g_1, \ldots, g_n\) be the finite presentation cover of
  Lemma~\ref{mk-lem:qcoh_finite_presentation_cover}, with \(\operatorname{span}\{g_j\} = R\). Then the
  canonical \(R\)-linear map \(\Gamma(X, \mathcal{F})_f \to \Gamma(D(f), \mathcal{F})\) induced by
  \(\rho_f : \Gamma(X, \mathcal{F}) \to \Gamma(D(f), \mathcal{F})\) (the localisation lift of \(\rho_f\)
  at the powers of \(f\)) is an isomorphism.
\end{lemma}
\begin{proof}
  \leanok
  \uses{mk-lem:qcoh_section_isLocalizedModule, mk-lem:isLocalizedModule_linearEquiv_mathlib}
    By Lemma~\ref{mk-lem:qcoh_section_isLocalizedModule}, the
  section-restriction map \(\rho_f : \Gamma(X, \mathcal{F}) \to \Gamma(D(f), \mathcal{F})\) is shown to
  exhibit \(\Gamma(D(f), \mathcal{F})\) as the localisation of \(\Gamma(X, \mathcal{F})\) at the powers of
  \(f\). Packaging this localisation map as a linear equivalence
  (Lemma~\ref{mk-lem:isLocalizedModule_linearEquiv_mathlib})
  yields the canonical \(R\)-linear isomorphism \(\Gamma(X, \mathcal{F})_f \cong \Gamma(D(f),
  \mathcal{F})\); its underlying map is exactly the localisation lift of \(\rho_f\), so that lift is an
  isomorphism.
\end{proof}

\begin{lemma}[Sections over a basic open localize the global sections]\label{mk-lem:qcoh_section_isLocalizedModule}
  \lean{AlgebraicGeometry.qcoh_section_isLocalizedModule}
  \leanok
  \dcref{01HV,custom}

  \uses{mk-lem:qcoh_finite_presentation_cover, mk-lem:qcoh_section_equalizer,
    mk-lem:localized_module_map_exact_mathlib, mk-lem:tile_section_localization,
    mk-lem:isLocalizedModule_of_exact, mk-lem:section_isLocalizedModule_of_presentation}
  \textit{Source: Stacks Project, Schemes, Tag 01HV, \texttt{lemma-spec-sheaves}~(4) and
  \texttt{lemma-widetilde-pullback}.}
  Let \(X = \operatorname{Spec} R\), let \(\mathcal{F}\) be a quasi-coherent
  \(\mathcal{O}_X\)-module, and let \(f \in R\). Then the restriction-of-sections map
  \(\rho_f : \Gamma(X, \mathcal{F}) \to \Gamma(D(f), \mathcal{F})\) exhibits
  \(\Gamma(D(f), \mathcal{F})\) as the localization of \(\Gamma(X, \mathcal{F})\) at the powers of
  \(f\). This is the
  generalization of Tag 01HV~(4) (\(\Gamma(D(f), \widetilde{M}) = M_f\)) from \(\widetilde{M}\) to an
  arbitrary quasi-coherent \(\mathcal{F}\).
\end{lemma}
\begin{proof}
  \leanok
  \uses{mk-lem:qcoh_finite_presentation_cover, mk-lem:qcoh_section_equalizer,
    mk-lem:localized_module_map_exact_mathlib, mk-lem:tile_section_localization,
    mk-lem:overlap_section_localization, mk-lem:isLocalizedModule_of_exact,
    mk-lem:section_isLocalizedModule_of_presentation}
    By Lemma~\ref{mk-lem:qcoh_finite_presentation_cover}, refine the quasi-coherence cover of
  \(\mathcal{F}\) to a finite standard cover \(\operatorname{Spec} R = \bigcup_{j=1}^{n} D(g_j)\) with
  \(\operatorname{span}\{g_1, \ldots, g_n\} = R\), where each \(D(g_j)\) — and hence each overlap
  \(D(g_j g_k)\) — lies in a cover member carrying a presentation of the over-picture restriction, so
  that the tiles \(\mathcal{F}_{(g_j)}\) and \(\mathcal{F}_{(g_j g_k)}\) are globally presented.

  Write the degree-\(0/1\) sheaf-axiom equalizer (Lemma~\ref{mk-lem:qcoh_section_equalizer}) twice: once for
  \(\mathcal{F}\) on the standard cover \(X = \bigcup_j D(g_j)\) (take \(W = X\)), and once for the induced
  cover \(D(f) = \bigcup_j \bigl(D(f) \cap D(g_j)\bigr)\) of \(D(f)\) (take \(W = D(f)\), using
  \(D(f) \cap D(g_j) = D(g_j f)\)):
  \[
    0 \to \Gamma(X, \mathcal{F}) \xrightarrow{\rho} \prod_j \Gamma(D(g_j), \mathcal{F})
      \xrightarrow{\delta} \prod_{j,k} \Gamma(D(g_j g_k), \mathcal{F}),
  \]
  \[
    0 \to \Gamma(D(f), \mathcal{F}) \xrightarrow{\rho'} \prod_j \Gamma(D(g_j f), \mathcal{F})
      \xrightarrow{\delta'} \prod_{j,k} \Gamma(D(g_j g_k f), \mathcal{F}).
  \]
  Both rows are left-exact (the equalizer presents the first map as the kernel of the second). They
  assemble into a commutative ladder
  \[
    \begin{array}{ccccc}
      \Gamma(X, \mathcal{F}) & \xrightarrow{\rho} & \prod_j \Gamma(D(g_j), \mathcal{F})
        & \xrightarrow{\delta} & \prod_{j,k} \Gamma(D(g_j g_k), \mathcal{F}) \\
      \big\downarrow{\,\rho_f} & & \big\downarrow{\,b} & & \big\downarrow{\,c} \\
      \Gamma(D(f), \mathcal{F}) & \xrightarrow{\rho'} & \prod_j \Gamma(D(g_j f), \mathcal{F})
        & \xrightarrow{\delta'} & \prod_{j,k} \Gamma(D(g_j g_k f), \mathcal{F})
    \end{array}
  \]
  whose three verticals are the section-restriction maps along the inclusions \(D(f) \subseteq X\),
  \(D(g_j f) \subseteq D(g_j)\), and \(D(g_j g_k f) \subseteq D(g_j g_k)\). Both squares commute by
  functoriality of restriction: restricting a global section first to \(D(g_j)\) and then to
  \(D(g_j f)\) equals restricting first to \(D(f)\) and then to \(D(g_j f)\), and likewise on the
  overlaps.

  The two right-hand verticals are localisations at the powers of \(f\). Componentwise, the per-tile
  localisation (Lemma~\ref{mk-lem:tile_section_localization}) exhibits \(\Gamma(D(g_j f), \mathcal{F})\) as
  the localisation of \(\Gamma(D(g_j), \mathcal{F})\) at the powers of \(f\), and the per-overlap
  localisation (Lemma~\ref{mk-lem:overlap_section_localization}) — the transport of
  Lemma~\ref{mk-lem:tile_section_localization} along \(D(g_j g_k) = D(g_j) \cap D(g_k)\) and
  \(D(g_j g_k f) = D(g_j f) \cap D(g_k f)\) — exhibits \(\Gamma(D(g_j g_k f), \mathcal{F})\) as the
  localisation of \(\Gamma(D(g_j g_k), \mathcal{F})\). Since a finite product of localisation maps at the
  same set is again a localisation, the product maps \(b\) and \(c\) are
  \(\operatorname{IsLocalizedModule}\) at the powers of \(f\).

  Apply the abstract left-exact-ladder kernel comparison
  Lemma~\ref{mk-lem:isLocalizedModule_of_exact} to this ladder, with \(S = \operatorname{powers} f\), the
  two left-exact rows, and the two localised right-hand verticals \(b, c\). It concludes that the left
  vertical \(\rho_f : \Gamma(X, \mathcal{F}) \to \Gamma(D(f), \mathcal{F})\) is itself a localisation map
  at the powers of \(f\); equivalently, its canonical lift \(\Gamma(X, \mathcal{F})_f \cong
  \Gamma(D(f), \mathcal{F})\) is an isomorphism. (The exactness of localisation,
  Lemma~\ref{mk-lem:localized_module_map_exact_mathlib}, is the converse statement: it certifies that the
  localised \(X\)-row is again left-exact, the consistency under which the kernel comparison operates.)

  \emph{Crucial non-circularity point.} The only ``sections-localise'' inputs are the per-tile
  localisations (Lemma~\ref{mk-lem:tile_section_localization}), which live on the tiles
  \(\mathcal{F}_{(g_j)}\) and \(\mathcal{F}_{(g_j g_k)}\). There \(\mathcal{F}\) is associated to a
  module — the tile is globally presented, hence \(\mathcal{F}_{(g)} \cong \widetilde{M_g}\) (since
  \(\widetilde{(-)}\) preserves colimits), \textbf{not} by computing \(\Gamma(D(g), -)\) of an abstract
  cokernel — and the localisation is given by
  Lemma~\ref{mk-lem:section_isLocalizedModule_of_presentation} via the canonical section
  equality \(\Gamma(\mathcal{F}_{(g)}, V) = \Gamma(\mathcal{F}, \iota(V))\)
  (Lemma~\ref{mk-lem:restrict_obj_mathlib}). The global statement for \(\mathcal{F}\) is then recovered
  from the tiles by the sheaf-axiom equalizer and the exactness of localisation, \textbf{not} by
  identifying \(\Gamma(X, \mathcal{F})\) with an abstract localised module of its restrictions. In
  particular no instance of the keystone at the \(g_j\) is invoked, so the regress that an algebraic
  span-cover descent on the global sections would create (reducing the keystone-at-\(f\) to the
  keystone-at-\(g_j\), the same statement) is avoided.
\end{proof}

\begin{lemma}[Quasi-coherence implies the counit is an isomorphism, via section-localization (01I8)]
  \label{mk-lem:qcoh_isIso_fromTildeGamma}
  \lean{AlgebraicGeometry.isIso_fromTildeΓ_of_quasicoherent,
    AlgebraicGeometry.isIso_fromTildeΓ_app_basicOpen}
  \dcref{01I8,custom}

  \uses{mk-lem:qcoh_section_isLocalizedModule, mk-lem:fromTildeGamma_mathlib,
    mk-lem:isIso_fromTildeGamma_iff_mathlib, mk-lem:forget_reflectsIso_mathlib,
    mk-lem:isIso_fromTildeGamma_iff_isLocalizing_mathlib,
    mk-lem:isLocalizedModule_linearEquiv_mathlib}
  \textit{Source: Stacks Project, Schemes, Tag 01I8, \texttt{lemma-quasi-coherent-affine}.}
  Let \(X = \operatorname{Spec} R\) and let \(\mathcal{F}\) be a quasi-coherent
  \(\mathcal{O}_X\)-module. Then the tilde--\(\Gamma\) counit
  \(\operatorname{fromTilde\Gamma} : \widetilde{\Gamma(X, \mathcal{F})} \to \mathcal{F}\) is an
  isomorphism. This is the affine structure theorem 01I8, proved here by checking the counit on the
  basis of distinguished opens. It is the unconditional instance that discharges the
  \(\operatorname{IsIso}(\operatorname{fromTilde\Gamma})\) hypothesis of
  Lemma~\ref{mk-lem:qcoh_iso_tilde_sections}.
\end{lemma}
\begin{proof}
  \uses{mk-lem:qcoh_section_isLocalizedModule, mk-lem:fromTildeGamma_mathlib,
    mk-lem:isIso_fromTildeGamma_iff_mathlib, mk-lem:forget_reflectsIso_mathlib,
    mk-lem:isLocalizedModule_linearEquiv_mathlib}
  The forgetful functor reflects isomorphisms and the distinguished opens \(\{D(f)\}\) form a basis
  of \(\operatorname{Spec} R\) (Lemma~\ref{mk-lem:forget_reflectsIso_mathlib}), so it suffices to show
  that the component of \(\operatorname{fromTilde\Gamma}\) over each \(D(f)\) is an isomorphism. By
  Lemma~\ref{mk-lem:fromTildeGamma_mathlib} that component is the localization lift
  \(\operatorname{IsLocalizedModule.lift}(\operatorname{powers} f)\) of the section-restriction map
  \(\rho_f : \Gamma(X, \mathcal{F}) \to \Gamma(D(f), \mathcal{F})\) along the structure-sheaf
  localization \(\widetilde{(-)}.\operatorname{toOpen}\), the latter being an
  \(\operatorname{IsLocalizedModule}\) at the powers of \(f\). By the keystone
  Lemma~\ref{mk-lem:qcoh_section_isLocalizedModule}, \(\rho_f\) is itself an
  \(\operatorname{IsLocalizedModule}\) at the powers of \(f\); hence both the source and the target of
  the component are localizations of \(\Gamma(X, \mathcal{F})\) at the powers of \(f\), and the lift
  between two such localizations is an isomorphism
  (Lemma~\ref{mk-lem:isLocalizedModule_linearEquiv_mathlib}). Thus
  \(\operatorname{fromTilde\Gamma}\) is an isomorphism on every \(D(f)\), and therefore an
  isomorphism. (Equivalently, this exhibits \(\mathcal{F}\) in the essential image of
  \(\widetilde{(-)}\), the criterion of Lemma~\ref{mk-lem:isIso_fromTildeGamma_iff_mathlib}.)
\end{proof}

\begin{remark}[The 01I8 decomposition: quasi-coherence \(\Rightarrow\) counit is iso]
  \label{mk-rem:o1i8_decomposition}
  \dcref{01I8,custom}
  \uses{mk-lem:qcoh_iso_tilde_sections, mk-lem:qcoh_section_isLocalizedModule,
    mk-lem:qcoh_isIso_fromTildeGamma, mk-lem:qcoh_finite_presentation_cover, mk-lem:qcoh_section_equalizer,
    mk-lem:localized_module_map_exact_mathlib, mk-lem:tile_section_localization,
    mk-lem:qcoh_section_kernel_comparison,
    mk-lem:fromTildeGamma_mathlib, mk-lem:isIso_fromTildeGamma_iff_mathlib,
    mk-lem:forget_reflectsIso_mathlib, mk-lem:isLocalizedModule_linearEquiv_mathlib,
    mk-lem:exists_finite_basicOpen_subcover, mk-lem:qcoh_iso_tilde_sections_of_presentation,
    mk-lem:isIso_fromTildeGamma_of_presentation}
  \textit{Source: Stacks Project, Schemes, Tag 01I8, \texttt{lemma-quasi-coherent-affine}
  (proof).}
  Let \(X=\operatorname{Spec}R\) and let \(\mathcal F\) be quasi-coherent. Choose a
  finite basic-open cover \(\{D(g_j)\}\) on which \(\mathcal F\) admits presentations.
  The sheaf axiom presents both \(\Gamma(X,\mathcal F)\) and
  \(\Gamma(D(f),\mathcal F)\) as equalizers for the covers \(\{D(g_j)\}\) and
  \(\{D(fg_j)\}\). Exactness of localization, together with the section-localization
  isomorphisms on each tile and overlap, identifies these equalizers after localizing
  at \(f\). Hence
  \[
    \Gamma(D(f),\mathcal F)\cong\Gamma(X,\mathcal F)_f.
  \]
  The component of \(\operatorname{fromTilde\Gamma}\) on \(D(f)\) is the canonical map
  between these two localizations, so it is an isomorphism. Since distinguished opens
  form a basis, the counit is an isomorphism and
  \(\mathcal F\cong\widetilde{\Gamma(X,\mathcal F)}\).

  Alternatively, one may choose global generators for \(\mathcal F\), prove that the
  kernel of the resulting epimorphism from a free module is quasi-coherent, and compare
  the two resulting presentations after applying \(\widetilde{(-)}\). The lemmas below
  record this global-generation argument.
\end{remark}

We now give the global-generation argument. It uses quasi-coherence of free modules
and finite basic-open refinements of covers, as in the Stacks Project section
``Quasi-coherent sheaves on affines''.

\begin{lemma}[Free \(\mathcal{O}\)-modules are quasi-coherent]\label{mk-lem:free_isQuasicoherent}
  \lean{AlgebraicGeometry.free_isQuasicoherent}
  \leanok
  \dcref{custom}

  \textit{Source: Stacks Project, Schemes, \texttt{lemma-colimit-quasi-coherent}.}
  Let \(X = \operatorname{Spec} R\). For any index type \(\iota\) the free
  \(\mathcal{O}_X\)-module \(\mathcal{O}_X^{(\iota)}\) on \(\iota\) generators is
  quasi-coherent. Indeed it is the sheaf \(\widetilde{R^{(\iota)}}\) associated to the
  free \(R\)-module \(R^{(\iota)} = \bigoplus_{\iota} R\) (the finitely-supported
  functions \(\iota \to R\)): the tilde functor carries \(R^{(\iota)}\) to a direct sum
  of copies of the structure sheaf \(\mathcal{O}_X = \widetilde{R}\), and every sheaf of
  the form \(\widetilde{M}\) is quasi-coherent. Quasi-coherence is preserved under the
  isomorphism \(\mathcal{O}_X^{(\iota)} \cong \widetilde{R^{(\iota)}}\); concretely it is
  obtained from the identification of the free \(\mathcal{O}_X\)-module with the
  tilde of the free \(R\)-module.
\end{lemma}

\begin{lemma}[Finite basic-open refinement of a cover of \(\operatorname{Spec} R\)]\label{mk-lem:exists_finite_basicOpen_subcover}
  \lean{AlgebraicGeometry.exists_finite_basicOpen_subcover}
  \leanok
  \dcref{custom}

  \textit{Source: Stacks Project, Schemes, \texttt{lemma-standard-open} and
  \texttt{lemma-quasi-coherent-affine} (proof).}
  Let \(R\) be a ring and \((U_i)_{i \in \iota}\) a family of opens of
  \(\operatorname{Spec} R\) with \(\bigsqcup_i U_i = \operatorname{Spec} R\) (i.e.\
  \(\bigcup_i U_i = \top\)). Then there exist finitely many elements
  \(f_1, \ldots, f_n \in R\) and indices \(\varphi : \{1, \ldots, n\} \to \iota\) such
  that each basic open satisfies \(D(f_j) \subseteq U_{\varphi(j)}\) and the \(f_j\)
  generate the unit ideal, \(\operatorname{span}\{f_1, \ldots, f_n\} = R\) (equivalently
  \(\bigcup_j D(f_j) = \operatorname{Spec} R\)).
\end{lemma}
\begin{proof}
  \leanok

The basic opens \(D(f)\), \(f \in R\), form a basis of the Zariski topology on
  \(\operatorname{Spec} R\); hence every \(U_i\) is a union of basic opens, and refining
  the cover \((U_i)\) yields a cover of \(\operatorname{Spec} R\) by basic opens
  \(D(f)\), each contained in some \(U_{\varphi}\). The space \(\operatorname{Spec} R\)
  is quasi-compact, so finitely many of these basic opens \(D(f_1), \ldots, D(f_n)\)
  already cover \(\operatorname{Spec} R\), with \(D(f_j) \subseteq U_{\varphi(j)}\) by
  construction. Finally a finite family of basic opens covers all of
  \(\operatorname{Spec} R\) if and only if the \(f_j\) generate the unit ideal:
  \(\bigcup_j D(f_j) = \operatorname{Spec} R \iff \operatorname{span}\{f_j\} = R\), since
  the closed complement \(V(\operatorname{span}\{f_j\})\) is empty exactly when the ideal
  is the whole ring.
\end{proof}

\begin{lemma}[Restriction of an \(\mathcal{O}_{\operatorname{Spec} R}\)-module to \(D(f)\)
  as an \(\mathcal{O}_{\operatorname{Spec} R_f}\)-module]
  \label{mk-lem:modules_restrict_basicOpen}
  \lean{AlgebraicGeometry.modulesRestrictBasicOpen, AlgebraicGeometry.modulesRestrictBasicOpenIso}
  \dcref{01I8,custom}

  \textit{Source: Stacks Project, Schemes, \texttt{lemma-widetilde-pullback}, Tag 01I8.}
  Let \(X = \operatorname{Spec} R\), let \(\mathcal{F}\) be an \(\mathcal{O}_X\)-module
  (an object of \((\operatorname{Spec} R).\mathrm{Modules}\)), and let \(f \in R\). Then,
  transporting along the canonical isomorphism of affine schemes
  \(D(f) \cong \operatorname{Spec} R_f\), the restriction \(\mathcal{F}|_{D(f)}\) is naturally
  an \(\mathcal{O}_{\operatorname{Spec} R_f}\)-module: there is an object
  \(\mathcal{F}_{(f)} \in (\operatorname{Spec} R_f).\mathrm{Modules}\) together with an
  isomorphism \(\mathcal{F}|_{D(f)} \cong \mathcal{F}_{(f)}\) of sheaves of modules over the
  identification \(D(f) \cong \operatorname{Spec} R_f\). This assignment is functorial in
  \(\mathcal{F}\).
\end{lemma}
\begin{proof}
  The basic open \(D(f) \subseteq \operatorname{Spec} R\) is an affine scheme, and the
  canonical comparison \(D(f) \cong \operatorname{Spec} R_f\) is an isomorphism of affine
  schemes (it corresponds to the localisation ring map \(R \to R_f\)). Pulling the restricted
  module \(\mathcal{F}|_{D(f)}\) back along the inverse of this isomorphism transports it to a
  sheaf of \(\mathcal{O}_{\operatorname{Spec} R_f}\)-modules \(\mathcal{F}_{(f)}\); the unit of
  the transport supplies the comparison isomorphism. Restriction of sheaves of modules and
  pullback along an isomorphism are both functorial, so the assignment
  \(\mathcal{F} \mapsto \mathcal{F}_{(f)}\) is functorial in \(\mathcal{F}\).
\end{proof}

\begin{definition}[The distinguished open \(D(f)\) as an open of \(\operatorname{Spec} R\)]\label{mk-def:spec_basic_open}
  \lean{AlgebraicGeometry.specBasicOpen}
  \leanok
  \dcref{custom}

  Let \(R\) be a commutative ring and \(f \in R\). The \emph{distinguished open}
  \(D(f)\) is the open subset \(\operatorname{specBasicOpen} f = \{\mathfrak{p} : f \notin
  \mathfrak{p}\}\) of \(\operatorname{Spec} R\), packaged as an object of the frame of opens
  \((\operatorname{Spec} R).\mathrm{Opens}\); it is the basic open \(\operatorname{PrimeSpectrum.basicOpen} f\).
\end{definition}

\begin{lemma}[The localisation morphism \(\operatorname{Spec} R_f \to \operatorname{Spec} R\)]\label{mk-lem:spec_away_to_spec}
  \lean{AlgebraicGeometry.specAwayToSpec}
  \leanok
  \dcref{custom}

  \uses{mk-def:spec_basic_open, mk-lem:modules_restrict_basicOpen}
  Let \(R\) be a commutative ring and \(f \in R\). The morphism of schemes
  \(\operatorname{specAwayToSpec} f : \operatorname{Spec} R_f \to \operatorname{Spec} R\) is the
  composite of the inverse of the comparison isomorphism \(D(f) \cong \operatorname{Spec} R_f\)
  (\(\operatorname{basicOpenIsoSpecAway} f\)) with the open immersion
  \(D(f) \hookrightarrow \operatorname{Spec} R\) of Definition~\ref{mk-def:spec_basic_open}; it is the
  geometric morphism underlying the restriction transport of
  Lemma~\ref{mk-lem:modules_restrict_basicOpen}.
\end{lemma}

\begin{lemma}[The localisation morphism is the spectrum of \(R \to R_f\)]\label{mk-lem:spec_away_to_spec_eq}
  \lean{AlgebraicGeometry.specAwayToSpec_eq}
  \leanok
  \dcref{custom}

  \uses{mk-lem:spec_away_to_spec}
  With notation as in Lemma~\ref{mk-lem:spec_away_to_spec}, the morphism
  \(\operatorname{specAwayToSpec} f\) equals the spectrum of the localisation ring map, i.e.\
  \(\operatorname{specAwayToSpec} f = \operatorname{Spec.map}(\operatorname{algebraMap} R\,(R_f))\)
  where \(R_f = \operatorname{Localization.Away} f\). Hence the restriction-to-\(D(f)\) transport is
  pullback along \(\operatorname{Spec}\) of the localisation \(R \to R_f\).
\end{lemma}
\begin{proof}
  \leanok
  \uses{mk-lem:spec_away_to_spec}
  By definition \(\operatorname{specAwayToSpec} f\) is
  \((\operatorname{basicOpenIsoSpecAway} f)^{-1} \circ \iota_{D(f)}\); rewriting with
  \(\operatorname{Iso.inv\_comp\_eq}\) reduces the claim to the defining property of the comparison
  isomorphism \(\operatorname{basicOpenIsoSpecAway} f\), which identifies \(D(f)\) with
  \(\operatorname{Spec} R_f\) via the localisation map \(R \to R_f\).
\end{proof}

\begin{lemma}[Restricting \(\widetilde{M}\) to \(D(f)\) gives \(\widetilde{M_f}\)]
  \label{mk-lem:tilde_restrict_basicOpen}
  \lean{AlgebraicGeometry.tilde_restrict_basicOpen}
  \dcref{01I8,custom}

  \uses{mk-lem:modules_restrict_basicOpen}
  \textit{Source: Stacks Project, Schemes, \texttt{lemma-widetilde-pullback}, Tag 01I8.}
  Let \(R\) be a ring, \(M\) an \(R\)-module, and \(f \in R\). Under the identification
  \(D(f) \cong \operatorname{Spec} R_f\) of Lemma~\ref{mk-lem:modules_restrict_basicOpen}, the
  restriction of the associated sheaf \(\widetilde{M}\) to the standard open \(D(f)\) is
  isomorphic to the associated sheaf of the away-localised module:
  \(\widetilde{M}\,|_{D(f)} \cong \widetilde{M_f}\) as \(\mathcal{O}_{\operatorname{Spec}
  R_f}\)-modules, where \(M_f = (\operatorname{powers} f)^{-1} M\).
\end{lemma}
\begin{proof}
  \uses{mk-lem:modules_restrict_basicOpen}
  Transport \(\widetilde{M}\,|_{D(f)}\) to an \(\mathcal{O}_{\operatorname{Spec}
  R_f}\)-module by Lemma~\ref{mk-lem:modules_restrict_basicOpen}. On a standard open
  \(D(g) \subseteq D(f)\) the sections of \(\widetilde{M}\,|_{D(f)}\) are the localisation
  \(M_{gf}\), which is exactly the localisation of \(M_f\) at (the image of) \(g\); since
  \(\psi^{-1}(D(f)) = D(\psi^\sharp(f))\) for the localisation map \(\psi^\sharp : R \to R_f\),
  the value on \(D(f)\) itself is \(M_f\). These identifications are compatible with
  restriction, hence patch to an isomorphism of sheaves
  \(\widetilde{M}\,|_{D(f)} \cong \widetilde{M_f}\) over \(\operatorname{Spec} R_f\).
\end{proof}

\begin{lemma}[Presentations restrict along an open immersion of basic opens]
  \label{mk-lem:presentation_restrict_basicOpen}
  \lean{AlgebraicGeometry.presentation_restrict_basicOpen}
  \dcref{01I8,custom}

  \uses{mk-lem:modules_restrict_basicOpen, mk-lem:tilde_restrict_basicOpen}
  \textit{Source: Stacks Project, Schemes, \texttt{lemma-quasi-coherent-affine} (proof),
  Tag 01I8.}
  Let \(X = \operatorname{Spec} R\), let \(\mathcal{F}\) be an \(\mathcal{O}_X\)-module, let
  \(U \subseteq X\) be open, and suppose \(\mathcal{F}|_U\) admits a presentation
  \(\mathcal{O}_U^{(J)} \to \mathcal{O}_U^{(I)} \to \mathcal{F}|_U \to 0\). Then for every
  basic open \(D(g) \subseteq U\), the restricted module
  \(\mathcal{F}|_{D(g)}\) — viewed as an \(\mathcal{O}_{\operatorname{Spec} R_g}\)-module via
  Lemma~\ref{mk-lem:modules_restrict_basicOpen} — admits a presentation
  \(\mathcal{O}^{(J)} \to \mathcal{O}^{(I)} \to \mathcal{F}|_{D(g)} \to 0\) over
  \(\operatorname{Spec} R_g\); in particular \(\mathcal{F}|_{D(g)}\) is quasi-coherent over
  \(R_g\).
\end{lemma}
\begin{proof}
  \uses{mk-lem:modules_restrict_basicOpen, mk-lem:tilde_restrict_basicOpen}
  Restriction of sheaves of modules along the open immersion
  \(D(g) \hookrightarrow U\) is an additive functor that preserves the structure sheaf, hence
  preserves cokernels and arbitrary direct sums of \(\mathcal{O}\); applying it to the given
  right-exact sequence yields a right-exact sequence
  \(\mathcal{O}_{D(g)}^{(J)} \to \mathcal{O}_{D(g)}^{(I)} \to \mathcal{F}|_{D(g)} \to 0\).
  Transporting along \(D(g) \cong \operatorname{Spec} R_g\)
  (Lemma~\ref{mk-lem:modules_restrict_basicOpen}) turns the free modules
  \(\mathcal{O}_{D(g)}^{(I)}\), \(\mathcal{O}_{D(g)}^{(J)}\) into the free
  \(\mathcal{O}_{\operatorname{Spec} R_g}\)-modules \(\widetilde{R_g}^{(I)}\),
  \(\widetilde{R_g}^{(J)}\) (the structure sheaf restricts to the structure sheaf, and free
  modules restrict to free modules by Lemma~\ref{mk-lem:tilde_restrict_basicOpen}). This is a
  presentation of \(\mathcal{F}|_{D(g)}\) over \(\operatorname{Spec} R_g\), exhibiting it as
  quasi-coherent.
\end{proof}

\begin{lemma}[Restriction of a quasi-coherent sheaf to a basic open is quasi-coherent]
  \label{mk-lem:isQuasicoherent_restrict_basicOpen}
  \lean{AlgebraicGeometry.isQuasicoherent_restrict_basicOpen}
  \dcref{01I8,custom}

  \uses{mk-lem:exists_finite_basicOpen_subcover, mk-lem:qcoh_iso_tilde_sections_of_presentation,
    mk-lem:modules_restrict_basicOpen, mk-lem:tilde_restrict_basicOpen,
    mk-lem:presentation_restrict_basicOpen}
  \textit{Source: Stacks Project, Schemes, \texttt{lemma-quasi-coherent-affine} (proof),
  Tag 01I8.}
  Let \(X = \operatorname{Spec} R\), let \(\mathcal{F}\) be a quasi-coherent
  \(\mathcal{O}_X\)-module, and let \(f \in R\). Then, transporting along the canonical
  isomorphism of affine schemes \(D(f) \cong \operatorname{Spec} R_f\)
  (Lemma~\ref{mk-lem:modules_restrict_basicOpen}), the restriction \(\mathcal{F}|_{D(f)}\) is a
  quasi-coherent \(\mathcal{O}_{\operatorname{Spec} R_f}\)-module, and in particular it admits a
  presentation \(\mathcal{O}^{(J)} \to \mathcal{O}^{(I)} \to \mathcal{F}|_{D(f)} \to 0\) over
  \(\operatorname{Spec} R_f\). Consequently \(\mathcal{F}|_{D(f)} \cong \widetilde{M}\) for the
  \(R_f\)-module \(M = \Gamma(D(f), \mathcal{F})\).
\end{lemma}
\begin{proof}
  \uses{mk-lem:exists_finite_basicOpen_subcover, mk-lem:qcoh_iso_tilde_sections_of_presentation,
    mk-lem:modules_restrict_basicOpen, mk-lem:tilde_restrict_basicOpen,
    mk-lem:presentation_restrict_basicOpen}
  Refine the quasi-coherence cover of \(\mathcal{F}\) to a finite standard-open cover
  \(\operatorname{Spec} R = \bigcup_j D(g_j)\), \(D(g_j) \subseteq U_{\varphi(j)}\), by
  Lemma~\ref{mk-lem:exists_finite_basicOpen_subcover}. Transport \(\mathcal{F}|_{D(f)}\) to an
  \(\mathcal{O}_{\operatorname{Spec} R_f}\)-module via
  Lemma~\ref{mk-lem:modules_restrict_basicOpen}. The basic opens
  \(D(g_j f) = D(g_j) \cap D(f)\) cover \(D(f) = \operatorname{Spec} R_f\); on each,
  Lemma~\ref{mk-lem:presentation_restrict_basicOpen} restricts the presentation over
  \(U_{\varphi(j)}\) to one of \(\mathcal{F}|_{D(g_j f)}\), and
  Lemma~\ref{mk-lem:tilde_restrict_basicOpen} identifies the restriction with the associated sheaf
  of an away-localised module. Hence \(\mathcal{F}|_{D(f)}\) is locally of the form
  \(\widetilde{(-)}\), and the finitely many local presentations assemble into a presentation of
  \(\mathcal{F}|_{D(f)}\) over \(\operatorname{Spec} R_f\). Feeding it to
  Lemma~\ref{mk-lem:qcoh_iso_tilde_sections_of_presentation} yields
  \(\mathcal{F}|_{D(f)} \cong \widetilde{\Gamma(D(f), \mathcal{F})}\) as an \(R_f\)-module.
\end{proof}

\begin{lemma}[\(\operatorname{IsLocalizedModule}\) is local on a finite spanning cover]
  \label{mk-lem:isLocalizedModule_of_span_cover}
  \lean{AlgebraicGeometry.isLocalizedModule_of_span_cover, AlgebraicGeometry.exists_sum_pow_eq_one,
    AlgebraicGeometry.mem_range_of_span_pow, AlgebraicGeometry.eq_zero_of_span_pow,
    AlgebraicGeometry.map_smul_endFun, AlgebraicGeometry.bump_eq, AlgebraicGeometry.per_j_surj,
    AlgebraicGeometry.per_j_eq}
  \dcref{custom}

  Let \(R\) be a commutative ring, let \(M\) and \(N\) be \(R\)-modules, let
  \(g : M \to N\) be an \(R\)-linear map, let \(f \in R\), and let
  \(s : \{1, \ldots, n\} \to R\) be a finite family with
  \(\operatorname{span}\{s_1, \ldots, s_n\} = R\) (the unit ideal). For each \(j\) write
  \(M_{s_j} = (\operatorname{powers} s_j)^{-1} M\) and \(N_{s_j} = (\operatorname{powers}
  s_j)^{-1} N\) for the localisations at the powers of \(s_j\), and let
  \(g_{s_j} : M_{s_j} \to N_{s_j}\) be the induced \(R\)-linear (equivalently
  \(R_{s_j}\)-linear) map. Suppose that for every \(j\) the localised map \(g_{s_j}\)
  exhibits \(N_{s_j}\) as the localisation of \(M_{s_j}\) at the powers of (the image of)
  \(f\), i.e.\ \(\operatorname{IsLocalizedModule}(\operatorname{powers} f)\, g_{s_j}\). Then
  \(g\) itself exhibits \(N\) as the localisation of \(M\) at the powers of \(f\):
  \(\operatorname{IsLocalizedModule}(\operatorname{powers} f)\, g\).
\end{lemma}
\begin{proof}

  This is a purely commutative-algebra statement: the predicate
  \(\operatorname{IsLocalizedModule}(\operatorname{powers} f)\, g\) is the conjunction of its
  three defining conditions, and each is verified by descent along the spanning cover. Throughout we use the standard partition
  of unity: since \(\operatorname{span}\{s_j\} = R\), for any chosen exponents \(N_j\) the
  powers \(s_j^{N_j}\) again span the unit ideal, so there are \(r_j \in R\) with
  \(\sum_j r_j s_j^{N_j} = 1\); and an element of an \(R\)-module is zero iff its image in
  \(M_{s_j}\) vanishes for every \(j\) (locality of being zero on a spanning cover, the
  \(s_j\)-torsion form of the same partition of unity).

  \emph{\(f\) acts invertibly on \(N\).} For each \(j\), the hypothesis
  \(\operatorname{IsLocalizedModule}(\operatorname{powers} f)\, g_{s_j}\) makes multiplication
  by \(f\) on \(N_{s_j}\) bijective. Multiplication by \(f\) on \(N\) is an isomorphism iff it
  is so after localising at each \(s_j\) (an \(R\)-linear endomorphism is bijective iff it is
  bijective on the cover \(\{D(s_j)\}\), since injectivity and surjectivity are both local on
  a spanning cover); hence \(f \cdot (-) : N \to N\) is bijective, and likewise every power
  \(f^k\).

  \emph{Every \(y \in N\) satisfies \(f^k y = g(x)\) for some \(x \in M\), \(k\).}
  Fix \(y \in N\). For each \(j\), surjectivity of \(g_{s_j}\) as a localisation at the powers
  of \(f\) yields \(m_j \in M\), exponents \(a_j\) (clearing the \(s_j\)-denominator) and
  \(k_j\) with \(s_j^{a_j} f^{k_j} y = g(m_j)\) holding after localising at \(s_j\); clearing
  the remaining \(s_j\)-torsion (enlarging \(a_j\)) we may take this as an honest equality
  \(s_j^{a_j} f^{k_j} y = g(m_j)\) in \(N\). Choose a common \(k = \max_j k_j\) (replacing
  \(m_j\) by \(f^{k - k_j} m_j\)) so that \(s_j^{a_j} f^{k} y = g(m_j)\) for all \(j\). Pick,
  by the partition of unity, \(r_j\) with \(\sum_j r_j s_j^{a_j} = 1\) (possible because the
  \(s_j^{a_j}\) span the unit ideal). Then
  \(f^{k} y = \sum_j r_j s_j^{a_j} f^{k} y = \sum_j r_j\, g(m_j) = g\bigl(\sum_j r_j m_j\bigr)\),
  so \(x = \sum_j r_j m_j\) and the exponent \(k\) work.

  \emph{If \(g(x) = g(x')\) then \(f^k(x - x') = 0\) for some \(k\).} Put
  \(z = x - x'\), so \(g(z) = 0\). For each \(j\), the equaliser clause of
  \(\operatorname{IsLocalizedModule}(\operatorname{powers} f)\, g_{s_j}\) gives, from the
  vanishing of \(g_{s_j}(z) = 0\) in \(N_{s_j}\), an exponent \(k_j\) with
  \(f^{k_j} z = 0\) after localising at \(s_j\); clearing the \(s_j\)-denominator there is
  \(b_j\) with \(s_j^{b_j} f^{k_j} z = 0\) in \(M\). Take \(k = \max_j k_j\), so
  \(s_j^{b_j} f^{k} z = 0\) for all \(j\). With \(r_j\) chosen so that
  \(\sum_j r_j s_j^{b_j} = 1\),
  \(f^{k} z = \sum_j r_j s_j^{b_j} f^{k} z = 0\). Hence \(f^{k}(x - x') = 0\).

  The three clauses are exactly the data of
  \(\operatorname{IsLocalizedModule}(\operatorname{powers} f)\, g\).
\end{proof}

\begin{lemma}[Localised sections of a quasi-coherent sheaf]
  \label{mk-lem:qcoh_localized_sections}
  \lean{AlgebraicGeometry.qcoh_localized_sections}
  \dcref{custom}

  \uses{mk-lem:exists_finite_basicOpen_subcover, mk-lem:isQuasicoherent_restrict_basicOpen, mk-lem:isLocalizedModule_of_span_cover}
  \textit{Source: Stacks Project, Schemes, \texttt{lemma-widetilde-pullback}.}
  Let \(X = \operatorname{Spec} R\), let \(\mathcal{F}\) be a quasi-coherent
  \(\mathcal{O}_X\)-module, and let \(f \in R\). Then the restriction of sections
  \(\Gamma(X, \mathcal{F}) \to \Gamma(D(f), \mathcal{F})\) exhibits
  \(\Gamma(D(f), \mathcal{F})\) as the localisation of \(\Gamma(X, \mathcal{F})\) at the
  powers of \(f\): it is an \(\operatorname{IsLocalizedModule}\) for the submonoid
  \(\{1, f, f^2, \ldots\}\). In particular every section over \(D(f)\) is, after
  multiplication by a suitable power \(f^N\), the restriction of a global section, and a
  global section restricting to \(0\) on \(D(f)\) is annihilated by a power of \(f\).
\end{lemma}
\begin{proof}
  \uses{mk-lem:exists_finite_basicOpen_subcover, mk-lem:isQuasicoherent_restrict_basicOpen, mk-lem:isLocalizedModule_of_span_cover}
  Choose, by Lemma~\ref{mk-lem:exists_finite_basicOpen_subcover}, a finite basic-open cover
  \(\operatorname{Spec} R = \bigcup_{j=1}^{n} D(s_j)\) with
  \(\operatorname{span}\{s_1, \ldots, s_n\} = R\). On each member
  \(D(s_j) = \operatorname{Spec} R_{s_j}\) the restriction \(\mathcal{F}|_{D(s_j)}\) is, by
  Lemma~\ref{mk-lem:isQuasicoherent_restrict_basicOpen}, a quasi-coherent
  \(\mathcal{O}_{\operatorname{Spec} R_{s_j}}\)-module admitting a presentation, hence (by the
  affine structure theorem) \(\mathcal{F}|_{D(s_j)} \cong \widetilde{M_j}\) for an
  \(R_{s_j}\)-module \(M_j = \Gamma(D(s_j), \mathcal{F})\). For the standard sheaf
  \(\widetilde{M_j}\) the sections over a distinguished open \(D(g) \subseteq D(s_j)\) are
  the away localisation \((M_j)_g\) (\texttt{lemma-widetilde-pullback}:
  \(\Gamma(D(g), \widetilde{M}) = M_g\)). Applying this with \(g\) the image of \(f\) shows
  that, after localising at \(s_j\), the restriction map
  \(\Gamma(X, \mathcal{F}) \to \Gamma(D(f), \mathcal{F})\) becomes the localisation map
  \((M_j) \to (M_j)_f\), i.e.\ the localised map
  \(g_{s_j} : \Gamma(X, \mathcal{F})_{s_j} \to \Gamma(D(f), \mathcal{F})_{s_j}\) is an
  \(\operatorname{IsLocalizedModule}\) for the powers of \(f\).

  Now apply Lemma~\ref{mk-lem:isLocalizedModule_of_span_cover} with \(M = \Gamma(X, \mathcal{F})\),
  \(N = \Gamma(D(f), \mathcal{F})\), the restriction map \(g\), the element \(f\), and the
  spanning family \(s_1, \ldots, s_n\): since \(g_{s_j}\) is an
  \(\operatorname{IsLocalizedModule}\) for the powers of \(f\) for every \(j\) and the
  \(s_j\) generate the unit ideal, the patching primitive concludes that \(g\) itself is an
  \(\operatorname{IsLocalizedModule}\) for \(\{1, f, f^2, \ldots\}\). Concretely this packages
  the two clauses: every section over \(D(f)\) is, after multiplication by some \(f^N\), the
  restriction of a global section, and a global section vanishing on \(D(f)\) is annihilated
  by a power of \(f\).
\end{proof}

\begin{lemma}[Global generation of a quasi-coherent sheaf on an affine]
  \label{mk-lem:qcoh_global_generation}
  \lean{AlgebraicGeometry.qcoh_global_generation}
  \dcref{01I8,custom}

  \uses{mk-lem:qcoh_localized_sections, mk-lem:exists_finite_basicOpen_subcover}
  \textit{Source: Stacks Project, Schemes, Tag 01I8, \texttt{lemma-quasi-coherent-affine}
  (proof); cf.\ Hartshorne II.5.16--17.}
  Let \(X = \operatorname{Spec} R\) and \(\mathcal{F}\) a quasi-coherent
  \(\mathcal{O}_X\)-module. Then \(\mathcal{F}\) is generated by its global sections:
  there is an index set \(I\) of global sections and an epimorphism of
  \(\mathcal{O}_X\)-modules \(\mathcal{O}_X^{(I)} \twoheadrightarrow \mathcal{F}\).
  Equivalently, \(\mathcal{F}\) carries a global generating family (the data
  \(\mathcal{F}.\mathrm{GeneratingSections}\)).
\end{lemma}
\begin{proof}
  \uses{mk-lem:qcoh_localized_sections, mk-lem:exists_finite_basicOpen_subcover}
  By Lemma~\ref{mk-lem:exists_finite_basicOpen_subcover} pick a finite basic-open cover
  \(\operatorname{Spec} R = \bigcup_{j=1}^{n} D(f_j)\) on each of which
  \(\mathcal{F}|_{D(f_j)} \cong \widetilde{M_j}\). Each \(M_j\) is generated, as an
  \(R_{f_j}\)-module, by local sections \(t \in \Gamma(D(f_j), \mathcal{F})\). By the
  surjectivity half of Lemma~\ref{mk-lem:qcoh_localized_sections}, after multiplying by a
  suitable power \(f_j^{N}\) each such local section is the restriction of a global
  section \(s \in \Gamma(X, \mathcal{F})\); since \(f_j\) is a unit on \(D(f_j)\), these
  global sections still generate \(\mathcal{F}\) over \(D(f_j)\). Collecting the finitely
  many resulting global sections over all \(j\) gives a family that generates
  \(\mathcal{F}\) on every member of the cover, hence everywhere: the induced map
  \(\mathcal{O}_X^{(I)} \to \mathcal{F}\) is surjective on each \(D(f_j)\) and so is an
  epimorphism of sheaves.
\end{proof}

\begin{lemma}[The tilde functor preserves kernels]\label{mk-lem:tilde_preserves_kernels}
  \lean{AlgebraicGeometry.tildePreservesFiniteLimits, AlgebraicGeometry.tilde_preservesMonomorphisms, AlgebraicGeometry.tilde_preservesHomology, AlgebraicGeometry.sectMapₗ, AlgebraicGeometry.sectMapₗ_eq, AlgebraicGeometry.tilde_injective_app_basicOpen, AlgebraicGeometry.tilde_preservesFiniteColimits, AlgebraicGeometry.tilde_toStalk_map_injective, AlgebraicGeometry.stalkMapₗ, AlgebraicGeometry.stalkMapₗ_eq, AlgebraicGeometry.stalkMapₗ_injective}
  \leanok
  \dcref{01HV,custom}

  \uses{mk-lem:right_exact_mono_left_exact}
  \textit{Source: Stacks Project, Schemes,
  \texttt{lemma-kernel-cokernel-quasi-coherent} (and Tag 01HV, exactness of
  \(\widetilde{(-)}\)).}
  Let \(X = \operatorname{Spec} R\). The functor \(M \mapsto \widetilde{M}\) from
  \(R\)-modules to \(\mathcal{O}_X\)-modules is exact; in particular it preserves
  kernels: for an \(R\)-module map \(\varphi : M \to N\) the induced
  \(\widetilde{\varphi} : \widetilde{M} \to \widetilde{N}\) has
  \(\operatorname{Ker}(\widetilde{\varphi}) = \widetilde{\operatorname{Ker}(\varphi)}\),
  and more generally \(\widetilde{(-)}\) carries finite limits of \(R\)-modules to the
  corresponding limits of associated sheaves.
\end{lemma}

\begin{proof}
  \leanok\uses{mk-lem:right_exact_mono_left_exact}
  Right exactness is free: \(\widetilde{(-)}\) is a left adjoint. By
  Lemma~\ref{mk-lem:right_exact_mono_left_exact} it therefore suffices to show that
  \(\widetilde{(-)}\) preserves monomorphisms, and monomorphy may be checked on the sections
  over a basis of \(\operatorname{Spec} R\) --- basis injectivity gives injectivity on stalks,
  and a morphism of sheaves with injective stalk maps is a monomorphism.

  Take the basis of basic opens. Over \(D(r)\) the sections of \(\widetilde{M}\) are the
  localisation \(M_r\): the canonical map \(M \to \Gamma(D(r), \widetilde{M})\) exhibits the
  target as a localisation of \(M\) at the powers of \(r\). The section map of
  \(\widetilde{\varphi}\) over \(D(r)\) agrees with \(\varphi\) after composing with that
  canonical map --- which is the naturality of \(M \mapsto \Gamma(D(r), \widetilde{M})\) in
  \(M\) --- so, both being \(R\)-linear maps out of a localisation, the section map *is* the
  localisation \(\varphi_r\) of \(\varphi\). Localisation is flat, so \(\varphi_r\) is injective
  whenever \(\varphi\) is, which is the required basis injectivity.

  Note that no stalk is formed in the argument: the flatness input is applied over the basic
  opens, and the passage to stalks is only the standard basis-to-stalk implication used to
  recognise a monomorphism.
\end{proof}
\begin{proof}
  \textit{Source: Stacks Project, Schemes, \texttt{lemma-spec-sheaves} (item~(7) and
  proof), Tag 01HV.}
  Kernels in the category of sheaves of \(\mathcal{O}_X\)-modules are computed
  objectwise and sheafified, and isomorphisms of sheaves may be checked on stalks: a
  morphism of \(\mathcal{O}_X\)-modules is an isomorphism if and only if it induces an
  isomorphism on every stalk. So it suffices to compare stalks. By the stalk computation
  for the tilde functor (\texttt{lemma-spec-sheaves} item~(7)) the stalk of
  \(\widetilde{M}\) at a point \(x \in \operatorname{Spec} R\) corresponding to the prime
  \(\mathfrak{p}\) is the localisation
  \(\widetilde{M}_x = M_{\mathfrak{p}}\), naturally in \(M\): the stalk functor at \(x\)
  applied to \(\widetilde{(-)}\) is the localisation functor \(M \mapsto M_{\mathfrak{p}}\).

  Now let \(\varphi : M \to N\) be an \(R\)-module map with kernel
  \(K = \operatorname{Ker}(\varphi)\), giving the exact sequence
  \(0 \to K \to M \xrightarrow{\varphi} N\). Localisation \(R \to R_{\mathfrak{p}}\) is
  flat, hence the functor \(M \mapsto M_{\mathfrak{p}}\) is exact and in particular
  preserves kernels: \(K_{\mathfrak{p}} = \operatorname{Ker}(\varphi_{\mathfrak{p}} :
  M_{\mathfrak{p}} \to N_{\mathfrak{p}})\). Translating through the stalk identification,
  the natural comparison map
  \(\widetilde{K} \to \operatorname{Ker}(\widetilde{\varphi})\) induces, at every point
  \(x \leftrightarrow \mathfrak{p}\), the isomorphism \(K_{\mathfrak{p}} \xrightarrow{\sim}
  \operatorname{Ker}(\varphi_{\mathfrak{p}})\) (using that the stalk functor, being a
  filtered colimit, is itself exact and commutes with forming the kernel). Being a
  stalkwise isomorphism, the comparison map is an isomorphism of sheaves, so
  \(\operatorname{Ker}(\widetilde{\varphi}) = \widetilde{\operatorname{Ker}(\varphi)}\).
  The same stalkwise-flatness argument applied to a finite diagram of \(R\)-modules shows
  that \(\widetilde{(-)}\) carries finite limits to finite limits; equivalently
  \(\widetilde{(-)}\) is left exact, i.e.\ preserves finite limits (kernels).

\end{proof}

\begin{lemma}[The kernel of a generating epimorphism is quasi-coherent]
  \label{mk-lem:qcoh_kernel_qcoh}
  \lean{AlgebraicGeometry.qcoh_kernel_qcoh}
  \dcref{custom}

  \uses{mk-lem:qcoh_global_generation, mk-lem:tilde_preserves_kernels, mk-lem:free_isQuasicoherent}
  \textit{Source: Stacks Project, Schemes,
  \texttt{lemma-kernel-cokernel-quasi-coherent}.}
  Let \(X = \operatorname{Spec} R\), let \(\mathcal{F}\) be a quasi-coherent
  \(\mathcal{O}_X\)-module, and let \(\sigma : \mathcal{F}.\mathrm{GeneratingSections}\)
  be the generating family of Lemma~\ref{mk-lem:qcoh_global_generation}, with epimorphism
  \(\sigma.\pi : \mathcal{O}_X^{(I)} \twoheadrightarrow \mathcal{F}\). Then
  \(\operatorname{Ker}(\sigma.\pi)\) is again quasi-coherent on \(\operatorname{Spec} R\),
  and therefore — by Lemma~\ref{mk-lem:qcoh_global_generation} applied to it — globally
  generated, yielding a second generating family
  \(\tau : (\operatorname{Ker} \sigma.\pi).\mathrm{GeneratingSections}\).
\end{lemma}
\begin{proof}
  \uses{mk-lem:qcoh_global_generation, mk-lem:tilde_preserves_kernels, mk-lem:free_isQuasicoherent}
  The free sheaf \(\mathcal{O}_X^{(I)}\) is quasi-coherent
  (Lemma~\ref{mk-lem:free_isQuasicoherent}) and \(\mathcal{F}\) is quasi-coherent by
  hypothesis. On each member \(D(f_j)\) of a finite trivialising cover both source and
  target are of the form \(\widetilde{(-)}\), and the map \(\sigma.\pi\) corresponds to an
  \(R_{f_j}\)-module map; since \(\widetilde{(-)}\) preserves kernels
  (Lemma~\ref{mk-lem:tilde_preserves_kernels}),
  \(\operatorname{Ker}(\sigma.\pi)|_{D(f_j)}\) is the associated sheaf of the kernel of
  that module map, hence quasi-coherent locally on the cover. Quasi-coherence is local,
  so \(\operatorname{Ker}(\sigma.\pi)\) is quasi-coherent on all of
  \(\operatorname{Spec} R\). Applying Lemma~\ref{mk-lem:qcoh_global_generation} to the
  quasi-coherent \(\operatorname{Ker}(\sigma.\pi)\) produces the generating family
  \(\tau\).
\end{proof}

\begin{lemma}[Counit is iso from two global generating families]\label{mk-lem:isIso_fromTildeGamma_of_genSections}
  \lean{AlgebraicGeometry.isIso_fromTildeΓ_of_genSections}
  \leanok
  \dcref{custom}

  \uses{mk-lem:isIso_fromTildeGamma_of_presentation}
  Let \(X = \operatorname{Spec} R\) and \(\mathcal{F}\) an \(\mathcal{O}_X\)-module.
  Suppose given a global generating family
  \(\sigma : \mathcal{F}.\mathrm{GeneratingSections}\) (an epimorphism
  \(\sigma.\pi : \mathcal{O}_X^{(I)} \twoheadrightarrow \mathcal{F}\)) and a global
  generating family \(\tau : (\operatorname{Ker} \sigma.\pi).\mathrm{GeneratingSections}\)
  of its kernel. Then the tilde--\(\Gamma\) counit
  \(\operatorname{fromTilde\Gamma} : \widetilde{\Gamma(X, \mathcal{F})} \to \mathcal{F}\)
  is an isomorphism.
\end{lemma}
\begin{proof}
  \leanok
  \uses{mk-lem:isIso_fromTildeGamma_of_presentation}
  The two generating families assemble into a global presentation of \(\mathcal{F}\): the
  composite \(\mathcal{O}_X^{(J)} \xrightarrow{\tau.\pi} \operatorname{Ker} \sigma.\pi
  \hookrightarrow \mathcal{O}_X^{(I)} \xrightarrow{\sigma.\pi} \mathcal{F}\) is a global
  exact sequence \(\mathcal{O}_X^{(J)} \to \mathcal{O}_X^{(I)} \to \mathcal{F} \to 0\),
  i.e.\ the data \(\mathcal{F}.\mathrm{Presentation}\). Feeding this presentation to
  Lemma~\ref{mk-lem:isIso_fromTildeGamma_of_presentation} makes
  \(\operatorname{fromTilde\Gamma}\) an isomorphism.
\end{proof}

\begin{lemma}[Affine structure theorem from two generating families]\label{mk-lem:qcoh_iso_tilde_sections_of_genSections}
  \lean{AlgebraicGeometry.qcoh_iso_tilde_sections_of_genSections}
  \leanok
  \dcref{custom}

  \uses{mk-lem:isIso_fromTildeGamma_of_genSections, mk-lem:qcoh_iso_tilde_sections}
  Let \(X = \operatorname{Spec} R\) and \(\mathcal{F}\) an \(\mathcal{O}_X\)-module
  equipped with a global generating family \(\sigma\) and a global generating family
  \(\tau\) of \(\operatorname{Ker} \sigma.\pi\). Then \(\mathcal{F} \cong
  \widetilde{\Gamma(X, \mathcal{F})}\), naturally identifying
  \(\Gamma(D(f), \mathcal{F})\) with \(\Gamma(X, \mathcal{F})_f\) for every \(f \in R\).
\end{lemma}
\begin{proof}
  \leanok
  \uses{mk-lem:isIso_fromTildeGamma_of_genSections, mk-lem:qcoh_iso_tilde_sections}
  By Lemma~\ref{mk-lem:isIso_fromTildeGamma_of_genSections} the counit
  \(\operatorname{fromTilde\Gamma}\) is an isomorphism, which is exactly the conditional
  hypothesis of Lemma~\ref{mk-lem:qcoh_iso_tilde_sections}; its conclusion gives the
  isomorphism \(\mathcal{F} \cong \widetilde{\Gamma(X, \mathcal{F})}\) and the
  identification of distinguished-open sections with away localisations.
\end{proof}

\begin{lemma}[Quasi-coherence implies the counit is an isomorphism by global generation]
  \label{mk-lem:isIso_fromTildeGamma_of_quasicoherent}
  \dcref{01I8,custom}
  \uses{mk-lem:qcoh_global_generation, mk-lem:qcoh_kernel_qcoh,
    mk-lem:isIso_fromTildeGamma_of_genSections}
  \textit{Source: Stacks Project, Schemes, Tag 01I8,
  \texttt{lemma-quasi-coherent-affine}.}
  Let \(X = \operatorname{Spec} R\) and \(\mathcal{F}\) a quasi-coherent
  \(\mathcal{O}_X\)-module. Then the tilde--\(\Gamma\) counit
  \(\operatorname{fromTilde\Gamma} : \widetilde{\Gamma(X, \mathcal{F})} \to \mathcal{F}\)
  is an isomorphism. This is the affine global-generation instance
  (\([{\,}\mathrm{IsQuasicoherent}\,\mathcal{F}{\,}] \Rightarrow
  \operatorname{IsIso}(\operatorname{fromTilde\Gamma})\), Stacks Tag 01I8) that
  discharges the conditional hypothesis of Lemma~\ref{mk-lem:qcoh_iso_tilde_sections}.
\end{lemma}
\begin{proof}
  \uses{mk-lem:qcoh_global_generation, mk-lem:qcoh_kernel_qcoh,
    mk-lem:isIso_fromTildeGamma_of_genSections}
  Quasi-coherence of \(\mathcal{F}\) gives, by Lemma~\ref{mk-lem:qcoh_global_generation}, a
  global generating family \(\sigma\). Its kernel \(\operatorname{Ker} \sigma.\pi\) is
  again quasi-coherent and hence globally generated
  (Lemma~\ref{mk-lem:qcoh_kernel_qcoh}), yielding a generating family \(\tau\). Feeding the
  pair \((\sigma, \tau)\) to Lemma~\ref{mk-lem:isIso_fromTildeGamma_of_genSections} makes
  \(\operatorname{fromTilde\Gamma}\) an isomorphism.
\end{proof}

\begin{lemma}
[Standard-cover {\v C}ech vanishing for quasi-coherent coefficients]
  \label{mk-lem:affine_cech_vanishing_qcoh}
  \lean{AlgebraicGeometry.affine_cech_vanishing_qcoh,
    AlgebraicGeometry.affine_cech_vanishing_qcoh_of_tildeVanishing,
    AlgebraicGeometry.affine_tildeVanishing}
  \dcref{01XB,custom}

  \uses{mk-lem:cech_acyclic_affine, mk-lem:qcoh_iso_tilde_sections, mk-lem:affine_cech_vanishing_tilde_subcover,
    mk-lem:cechCohomology_isZero_of_iso, mk-def:affine_cover_system, mk-def:has_vanishing_higher_cech}
  \textit{Source: Stacks Project, Cohomology of Schemes, Tag 01XB
  (\texttt{lemma-quasi-coherent-affine-cohomology-zero}, proof).}
  Let \(X = \operatorname{Spec} R\) be an affine scheme and \(\mathcal{F}\) a quasi-coherent
  \(\mathcal{O}_X\)-module. Then \(\mathcal{F}\) has vanishing higher {\v C}ech cohomology for the
  affine cover system (Definition~\ref{mk-def:affine_cover_system},
  Definition~\ref{mk-def:has_vanishing_higher_cech}): for every standard cover \(\mathcal{U}\) and every
  \(p > 0\) one has \(\check{H}^p(\mathcal{U}, \mathcal{F}) = 0\). This supplies condition (3) of the
  basis-comparison Lemma~\ref{mk-lem:cech_to_cohomology_on_basis}, via the tilde-case residual
  Lemma~\ref{mk-lem:affine_cech_vanishing_tilde_subcover}.
\end{lemma}
\begin{proof}
  \uses{mk-lem:cech_acyclic_affine, mk-lem:qcoh_iso_tilde_sections,
    mk-lem:affine_cech_vanishing_tilde_subcover, mk-lem:cechCohomology_isZero_of_iso}
  \emph{Step 1 (reduction to the tilde case).} By Lemma~\ref{mk-lem:qcoh_iso_tilde_sections} the
  quasi-coherent \(\mathcal{F}\) is isomorphic to \(\widetilde{M}\) for
  \(M = \Gamma(X, \mathcal{F})\), with sections over each face \(D(g_\sigma)\) the away localisation
  \(M_{g_\sigma}\). Since {\v C}ech cohomology transports along an isomorphism of coefficient
  presheaves (Lemma~\ref{mk-lem:cechCohomology_isZero_of_iso}), it suffices to prove the vanishing for
  the tilde sheaf \(\widetilde{M}\): for every standard cover \(\mathcal{U}\) and every \(p > 0\),
  \(\check{H}^p(\mathcal{U}, \widetilde{M}) = 0\).

  \emph{Step 2 (the tilde residual).} A standard cover of the affine \(X = \operatorname{Spec} R\) in
  the affine cover system is a cover of a distinguished open \(D(f) = \bigsqcup_i D(g_i)\) by smaller
  distinguished opens; for a proper \(D(f)\) the family \(\{g_i\}\) does \emph{not} span the unit
  ideal of \(R\), so the whole-space standard-cover vanishing of
  Lemma~\ref{mk-lem:cech_acyclic_affine} (stated under \(\operatorname{span}\{g_i\} = R\)) does not
  apply directly. The required positive-degree vanishing \(\check{H}^p(\{D(g_i)\}, \widetilde{M}) = 0\)
  over such a proper cover is exactly Lemma~\ref{mk-lem:affine_cech_vanishing_tilde_subcover}, proved by
  change of base to \(R_f = \operatorname{Localization.Away} f\) (where the cover \emph{does} span the
  unit ideal). Combining Steps 1 and 2 gives \(\check{H}^p(\mathcal{U}, \mathcal{F}) = 0\) for all
  \(p > 0\) and every standard cover \(\mathcal{U}\); hence \(\mathcal{F}\) has vanishing higher
  {\v C}ech cohomology for the affine cover system.
\end{proof}

\begin{lemma}[{\v C}ech cohomology transports along an isomorphism of coefficients]\label{mk-lem:cechCohomology_isZero_of_iso}
  \lean{AlgebraicGeometry.cechCohomology_isZero_of_iso}
  \leanok
  \dcref{custom}

  \uses{mk-def:cech_cohomology_accessor, mk-def:section_cech_functoriality}
  \textit{Project-local transport lemma.}
  Let \(U : \iota \to \mathrm{Opens}(X)\) be an index family and let \(e : \mathcal{F} \cong
  \mathcal{G}\) be an isomorphism of presheaves of \(\mathcal{O}_X\)-modules. If
  \(\check{H}^p(\mathcal{U}, \mathcal{F}) = 0\) for some \(p\), then
  \(\check{H}^p(\mathcal{U}, \mathcal{G}) = 0\).
\end{lemma}
\begin{proof}
  \leanok
  \uses{mk-def:section_cech_functoriality}
  The section {\v C}ech complex is functorial in the coefficient presheaf
  (Definition~\ref{mk-def:section_cech_functoriality}), so \(e\) induces an isomorphism of section
  {\v C}ech complexes \(\check{\mathcal{C}}^\bullet(\mathcal{U}, \mathcal{F}) \cong
  \check{\mathcal{C}}^\bullet(\mathcal{U}, \mathcal{G})\), and applying the degree-\(p\) homology
  functor yields an isomorphism \(\check{H}^p(\mathcal{U}, \mathcal{F}) \cong
  \check{H}^p(\mathcal{U}, \mathcal{G})\). A group isomorphic to a zero object is itself zero.
\end{proof}

\begin{lemma}[A finite family of distinguished opens covers \(\operatorname{Spec} R\)
  iff its elements span the unit ideal]
  \label{mk-lem:iSup_basicOpen_eq_top_iff_span}
  \lean{PrimeSpectrum.iSup_basicOpen_eq_top_iff}
  \mathlibok
  \dcref{custom}

  For a family \(s : \iota \to R\), one has \(\bigsqcup_i D(s_i) = \operatorname{Spec} R\) if and only
  if \(\operatorname{span}(\{s_i : i \in \iota\}) = R\).
\end{lemma}

\begin{lemma}[A standard cover of \(D(f)\) spans the unit ideal of \(R_f\)]\label{mk-lem:affine_cover_span_localizationAway}
  \lean{AlgebraicGeometry.affine_cover_span_localizationAway}
  \leanok
  \dcref{custom}

  \uses{mk-lem:iSup_basicOpen_eq_top_iff_span}
  \textit{Project-local infrastructure.}
  Let \(R\) be a ring, \(f \in R\), and \(g : \iota \to R\) a finite family with \(D(f) =
  \bigsqcup_i D(g_i)\) in \(\operatorname{Spec} R\). Then the images
  \(\operatorname{algebraMap} R\, R_f\, (g_i)\) span the unit ideal of \(R_f =
  \operatorname{Localization.Away} f\):
  \[
    \operatorname{span}\bigl(\{\, g_i/1 : i \in \iota \,\}\bigr) = R_f.
  \]
\end{lemma}
\begin{proof}
  \leanok


Pull the equality \(D(f) = \bigsqcup_i D(g_i)\) back along the canonical open immersion
  \(\operatorname{Spec} R_f \cong D(f) \hookrightarrow \operatorname{Spec} R\), whose underlying
  continuous map is the comap of \(\operatorname{algebraMap} R\, R_f\). Comap commutes with
  distinguished opens and with suprema, so the pullback of \(\bigsqcup_i D(g_i)\) is
  \(\bigsqcup_i D(g_i/1)\), while the pullback of \(D(f)\) is all of \(\operatorname{Spec} R_f\)
  because \(f\) is a unit in \(R_f\). Thus \(\bigsqcup_i D(g_i/1) = \operatorname{Spec} R_f\), which
  by the span criterion (a family of distinguished opens covers the whole spectrum iff the elements
  generate the unit ideal) is equivalent to \(\operatorname{span}\{g_i/1\} = R_f\).
\end{proof}

\begin{lemma}[Localisation transitivity for away modules]\label{mk-lem:away_comparison_isLocalizedModule}
  \lean{AlgebraicGeometry.AwayComparison.comparison_isLocalizedModule}
  \leanok
  \dcref{custom}

  \textit{Project-local infrastructure.}
  Let \(R\) be a commutative ring, \(M\) an \(R\)-module, and \(a, b \in R\). Write
  \(M_a = \operatorname{LocalizedModule}(\operatorname{Submonoid.powers} a)\, M\) for the away
  localisation of \(M\) at \(a\). Then the canonical transitivity comparison \(M_a \to M_{ab}\)
  exhibits \(M_{ab}\) as the localisation of \(M_a\) at (the image of) \(b\):
  \[
    M_a\bigl[1/b\bigr] \;=\; M_{ab},
  \]
  i.e.\ the composite of the two away-localisation units is again the away localisation at the product
  \(ab\). In particular, it yields the degreewise identification
  \(M_{g_\sigma} \cong (M_f)_{g_\sigma}\) used in
  Lemma~\ref{mk-lem:affine_cech_vanishing_tilde_subcover}: when \(f\) is invertible in \(M_{g_\sigma}\)
  (because \(g_\sigma \in \sqrt{(f)}\)), localising \(M\) away from \(g_\sigma\) and localising
  \(M_f\) away from \(g_\sigma\) present the same module.
\end{lemma}

\begin{lemma}[Two-step away localisation localises at the product factor]\label{mk-lem:isLocalizedModule_comp_away}
  \lean{AlgebraicGeometry.AwayComparison.isLocalizedModule_comp_away}
  \leanok
  \dcref{custom}

  \textit{Project-local infrastructure.}
  Let \(R \to R_f\) be the away localisation of \(R\) at \(f\), let \(M\) be an \(R\)-module, and let
  \(\mathrm{mk}_f : M \to M_f\) exhibit \(M_f\) as the localisation of \(M\) at the powers of \(f\).
  Suppose \(a \in R\) lies in the radical of \((f)\), say \(a^j = f \cdot h\) for some \(h \in R\) and
  \(j \geq 1\) (equivalently \(f \mid a^j\)), and let \(g : M_f \to N\) be an \(R_f\)-linear map
  exhibiting \(N\) as the localisation of \(M_f\) at the powers of the image
  \(\operatorname{algebraMap} R\, R_f\, (a)\). Then the \(R\)-linear composite
  \[
    g \circ \mathrm{mk}_f \;:\; M \longrightarrow N
  \]
  exhibits \(N\) as the localisation of \(M\) at the powers of \(a\); that is, \(N = M\bigl[1/a\bigr]\).
\end{lemma}
\begin{proof}
  \leanok

Verify the three defining clauses of a localisation module directly. \emph{Units act invertibly:} a
  power \(a^n\) acts on \(N\) through its image \(\operatorname{algebraMap} R\, R_f\, (a)^n\), which is
  invertible because it is the localisation base for \(g\). \emph{Surjectivity:} any \(y \in N\) is
  \(g(m_0)\) divided by a power of \(a\) for some \(m_0 \in M_f\), and \(m_0\) is \(\mathrm{mk}_f(m)\)
  divided by a power \(f^l\) of \(f\); writing \(a^{jl} = (f h)^l = f^l h^l\) converts the \(f^l\) in the
  denominator into a power of \(a\), so \(y\) is the image of \(h^l \cdot m\) up to a power of \(a\).
  \emph{Cancellation:} if \(g(\mathrm{mk}_f(m)) = g(\mathrm{mk}_f(m'))\), chaining the cancellation
  property of \(g\) and then of \(\mathrm{mk}_f\) produces a single power of \(a\) annihilating
  \(m - m'\). This route avoids the converse of ``a localisation survives restriction
  of scalars''.
\end{proof}

\begin{lemma}[Section {\v C}ech module exactness from a localisation-away spanning datum]\label{mk-lem:section_cech_module_exact_of_localizationAway}
  \lean{AlgebraicGeometry.SectionCechModule.dDiff_exact_of_localizationAway}
  \leanok
  \dcref{custom}

  \uses{mk-lem:section_cech_module_exact, mk-lem:isLocalizedModule_comp_away}
  \textit{Project-local infrastructure.}
  Let \(M\) be an \(R\)-module, \(s : \iota \to R\) a finite family, and \(f \in R\). Assume
  \begin{enumerate}
  \item each \(s_i\) lies in the radical of \((f)\) (so \(f \mid s_i^{k_i}\) for some \(k_i\)), and
  \item the images \(\{\operatorname{algebraMap} R\, R_f\, (s_i)\}\) span the unit ideal of
    \(R_f = \operatorname{Localization.Away} f\).
  \end{enumerate}
  Then the un-localised section {\v C}ech module complex \(D^\bullet(s, M)\) of
  Definition~\ref{mk-def:section_cech_module_complex} is exact in every positive degree --- \emph{without}
  assuming that \(s\) spans the unit ideal of \(R\) itself.
\end{lemma}
\begin{proof}
  \leanok
  \uses{mk-lem:section_cech_module_exact, mk-lem:isLocalizedModule_comp_away}
  Run the computation over \(R_f\). Instantiate the polymorphic positive-degree exactness of
  Lemma~\ref{mk-lem:section_cech_module_exact} over the base ring \(R_f\), with module
  \(M_f = \operatorname{LocalizedModule}(\operatorname{Submonoid.powers} f)\, M\) and the family
  \(s/1 : \iota \to R_f\); hypothesis (2) is exactly the spanning datum that lemma requires, so the
  \(R_f\)-side complex \(D^\bullet(s/1, M_f)\) is exact in positive degrees. Transport this exactness
  back to the \(R\)-side along the degreewise \(R\)-linear isomorphisms
  \(M_{s_\sigma} \cong (M_f)_{s_\sigma}\) of Lemma~\ref{mk-lem:isLocalizedModule_comp_away} --- valid since
  each \(s_\sigma = \prod_k s_{\sigma k}\) lies in the radical of \((f)\) by (1). These isomorphisms are
  the vertical maps of an isomorphism of cochain complexes, commuting with the alternating-sum cofaces
  by naturality of the localisation comparison, and exactness transfers across an isomorphism of
  complexes.
\end{proof}

\begin{lemma}[Standard subcover {\v C}ech vanishing for the tilde sheaf (the residual)]
  \label{mk-lem:affine_cech_vanishing_tilde_subcover}
  \lean{AlgebraicGeometry.sectionCech_homology_exact_of_localizationAway,
    AlgebraicGeometry.sectionCechAbExact_loc}
  \dcref{custom}

  \uses{mk-lem:section_cech_module_exact, mk-lem:section_cech_homology_exact,
    mk-lem:affine_cover_span_localizationAway, mk-lem:away_comparison_isLocalizedModule,
    mk-def:cech_cohomology_accessor}
  \textit{Source: Stacks Project, Cohomology of Schemes,
  \texttt{lemma-cech-cohomology-quasi-coherent-trivial}.}
  Let \(R\) be a ring, \(M\) an \(R\)-module, \(f \in R\), and \(g : \{1, \ldots, n\} \to R\) a finite
  family with \(D(f) = \bigsqcup_i D(g_i)\) in \(\operatorname{Spec} R\). Then for every \(p > 0\),
  \[
    \check{H}^p\bigl(\{D(g_i)\}, \widetilde{M}\bigr) = 0.
  \]
\end{lemma}
\begin{proof}
  \uses{mk-lem:section_cech_module_exact, mk-lem:section_cech_homology_exact,
    mk-lem:affine_cover_span_localizationAway, mk-lem:away_comparison_isLocalizedModule,
    mk-lem:section_cech_module_exact_of_localizationAway, mk-lem:isLocalizedModule_comp_away}
  Set \(R_f = \operatorname{Localization.Away} f\). Following the Stacks proof's ``Write
  \(U = \operatorname{Spec}(A)\)'' step, we run the whole computation over the ring \(R_f\), where the
  cover \emph{does} span the unit ideal. The two algebraic inputs this needs are read off from the
  geometry of the cover.

  \emph{(a) Each \(g_i\) lies in the radical of \((f)\).} From \(D(f) = \bigsqcup_i D(g_i)\) we have in
  particular \(D(g_i) \subseteq D(f)\) for every \(i\), and the inclusion of basic opens
  \(D(g_i) \subseteq D(f)\) is equivalent to \(g_i \in \sqrt{(f)}\) (a power of \(g_i\) is a multiple of
  \(f\)). \emph{(b) The cover spans over \(R_f\).} Because \(f\) is a unit in \(R_f\), the equality
  \(D(f) = \bigsqcup_i D(g_i)\) pulls back to \(\operatorname{Spec} R_f = \bigsqcup_i D(g_i/1)\), so the
  images \(g_i/1\) span the unit ideal of \(R_f\) (Lemma~\ref{mk-lem:affine_cover_span_localizationAway}).

  With (a) and (b) in hand, Lemma~\ref{mk-lem:section_cech_module_exact_of_localizationAway} supplies the
  positive-degree exactness of the un-localised section {\v C}ech module complex \(D^\bullet(g, M)\):
  internally it instantiates the polymorphic exactness of
  Lemma~\ref{mk-lem:section_cech_module_exact} over \(R_f\) with the module
  \(M_f = \operatorname{LocalizedModule}(\operatorname{Submonoid.powers} f)\, M\), and transports the
  result back along the degreewise isomorphisms \(M_{g_\sigma} \cong (M_f)_{g_\sigma}\) of
  Lemma~\ref{mk-lem:isLocalizedModule_comp_away} (for a tuple \(\sigma\), \(g_\sigma = \prod_k g_{\sigma k}
  \in \sqrt{(f)}\), so both sides present \(M\) localised away from \(g_\sigma f\)).

  Finally, the tilde-bridge ladder identifies the section {\v C}ech complex of \(\widetilde{M}\) over
  \(\{D(g_i)\}\) with this localised-module complex --- via the same section-to-module identification
  and coface matching used in Lemma~\ref{mk-lem:section_cech_homology_exact} --- and wraps positive-degree
  exactness as the vanishing of the degree-\(p\) homology, giving
  \(\check{H}^p(\{D(g_i)\}, \widetilde{M}) = 0\) for all \(p > 0\).
\end{proof}

The {\v C}ech-to-cohomology comparison on a basis
(Lemma~\ref{mk-lem:cech_to_cohomology_on_basis}) is proved by a dimension shift on the
coefficient module, decomposed below into four sub-lemmas: the short exact sequence of
{\v C}ech complexes coming from a basis (Lemma~\ref{mk-lem:cech_ses_of_basis}), the
resulting stability of vanishing higher {\v C}ech cohomology under the injective
quotient (Lemma~\ref{mk-lem:quotient_vanishing_cech}), the sheaf-cohomology base case
\(H^1(U, \mathcal{F}) = 0\) (Lemma~\ref{mk-lem:absolute_cohomology_one_vanishing}), and the
  dimension-shift induction \(H^p(U, \mathcal{F}) = 0\) for \(p > 0\)
  (Lemma~\ref{mk-lem:absolute_cohomology_pos_vanishing}).

\subsection{Section {\v C}ech infrastructure and the cover-system encoding}

The sub-lemmas below have a \emph{cover-local}, presheaf-level form: each lemma fixes
a single index family \(U : \iota \to \mathrm{Opens}(X)\) and works with the
section {\v C}ech complexes of that one covering, taking the geometric inputs (section-level
exactness on faces, positive-degree {\v C}ech vanishing) as hypotheses. The cover-global
\((\mathcal{B}, \mathrm{Cov})\) data of Lemma~\ref{mk-lem:cech_to_cohomology_on_basis} re-enters only
in the inductive step (Lemma~\ref{mk-lem:absolute_cohomology_pos_vanishing}), where it is packaged by
the cover-system encoding (Definition~\ref{mk-def:basis_cov_system},
Definition~\ref{mk-def:has_vanishing_higher_cech}) and discharges those hypotheses via the per-face
short-exact-sequence derivation (Lemma~\ref{mk-lem:face_ses_of_sheaf_ses}). The supporting
definitions and the degreewise product engine are collected first.

\begin{definition}[{\v C}ech cohomology accessor]\label{mk-def:cech_cohomology_accessor}
  \lean{AlgebraicGeometry.cechCohomology}
  \leanok
  \dcref{custom}

  \uses{mk-def:section_cech_complex}
  \textit{Project-local accessor.}
  For an index family \(U : \iota \to \mathrm{Opens}(X)\) (giving the covering \(\mathcal{U}\)), a
  presheaf of \(\mathcal{O}_X\)-modules \(\mathcal{F}\), and \(p \geq 0\), the \emph{{\v C}ech
  cohomology} \(\check{H}^p(\mathcal{U}, \mathcal{F})\) is the degree-\(p\) homology of the section
  {\v C}ech complex of Definition~\ref{mk-def:section_cech_complex},
  \[
    \check{H}^p(\mathcal{U}, \mathcal{F}) \;:=\;
      H^p\bigl(\check{\mathcal{C}}^\bullet(\mathcal{U}, \mathcal{F})\bigr),
  \]
  taken as an object of \(\mathrm{Ab}\). This is the named wrapper that the 01EO chain refers to
  uniformly.
\end{definition}

\begin{definition}
[Functoriality of the section {\v C}ech complex]
  \label{mk-def:section_cech_functoriality}
  \lean{AlgebraicGeometry.sectionCechCosimplicialMap,
    AlgebraicGeometry.sectionCechCosimplicialFunctor,
    AlgebraicGeometry.sectionCechComplexFunctor,
    AlgebraicGeometry.sectionCechComplexMap}
  \dcref{custom}

  \uses{mk-def:section_cech_complex}
  \textit{Project-local.}
  A morphism \(\varphi : \mathcal{F} \to \mathcal{G}\) of presheaves of \(\mathcal{O}_X\)-modules
  acts on the section {\v C}ech data of Definition~\ref{mk-def:section_cech_complex} coordinatewise:
  on the basic-open section group of each intersection it is the map
  \(\varphi(U_{i_0 \ldots i_p}) : \mathcal{F}(U_{i_0 \ldots i_p}) \to \mathcal{G}(U_{i_0 \ldots
  i_p})\). These assemble into a cosimplicial morphism (the \emph{section {\v C}ech cosimplicial
  map}), functorial in \(\varphi\) (the \emph{section {\v C}ech cosimplicial functor}). Composing
  with the alternating-coface-map complex functor packages the section {\v C}ech complex as a
  functor in the coefficient presheaf (the \emph{section {\v C}ech complex functor}) and yields,
  for each \(\varphi\), the induced chain map of section {\v C}ech complexes
  \(\check{\mathcal{C}}^\bullet(\mathcal{U}, \varphi)\). These are the maps used to build the short
  complex of Definition~\ref{mk-def:section_cech_short_complex}.
\end{definition}

\begin{definition}[Per-face section short complex]\label{mk-def:face_short_complex}
  \lean{AlgebraicGeometry.faceShortComplex}
  \leanok
  \dcref{custom}

  \uses{mk-def:section_cech_complex}
  Let \(U : \iota \to \mathrm{Opens}(X)\) be an index family and let
  \(P : \mathcal{P}_1 \to \mathcal{P}_2 \to \mathcal{P}_3\) be a short complex of presheaves of
  \(\mathcal{O}_X\)-modules. For a degree \(p\) and a tuple \(\sigma : \{0, \ldots, p\} \to \iota\)
  write \(V_\sigma = \bigcap_k U(\sigma k)\) for the corresponding basic open. The \emph{per-face
  short complex} \(\mathrm{face}(U, P, \sigma)\) is the short complex of abelian groups obtained by
  evaluating \(P\) on sections over \(V_\sigma\),
  \[
    \mathcal{P}_1(V_\sigma) \longrightarrow \mathcal{P}_2(V_\sigma) \longrightarrow
      \mathcal{P}_3(V_\sigma),
  \]
  the maps being the section-level actions over \(V_\sigma\) of the two maps of \(P\). It is the
  \((p, \sigma)\)-indexed factor whose product over all tuples \(\sigma\) is the degree-\(p\) term
  of the section {\v C}ech complex.
\end{definition}

\begin{definition}[Short complex of section {\v C}ech complexes]\label{mk-def:section_cech_short_complex}
  \lean{AlgebraicGeometry.sectionCechComplexShortComplex}
  \leanok
  \dcref{custom}

  \uses{mk-def:section_cech_complex, mk-def:section_cech_functoriality}
  \textit{Project-local.}
  For an index family \(U\) and a short complex
  \(P : \mathcal{P}_1 \to \mathcal{P}_2 \to \mathcal{P}_3\) of presheaves of
  \(\mathcal{O}_X\)-modules, the \emph{short complex of section {\v C}ech complexes} is
  \[
    0 \to \check{\mathcal{C}}^\bullet(\mathcal{U}, \mathcal{P}_1) \to
      \check{\mathcal{C}}^\bullet(\mathcal{U}, \mathcal{P}_2) \to
      \check{\mathcal{C}}^\bullet(\mathcal{U}, \mathcal{P}_3) \to 0,
  \]
  a short complex in the category of cochain complexes of abelian groups, with the two maps the
  chain maps induced by the maps of \(P\) through the functoriality of
  Definition~\ref{mk-def:section_cech_functoriality}. Its short-exactness is the conclusion of
  Lemma~\ref{mk-lem:cech_ses_of_basis} and the input to Lemma~\ref{mk-lem:quotient_vanishing_cech}.
\end{definition}

\begin{definition}

[Cover system on a basis]
  \label{mk-def:basis_cov_system}
  \lean{AlgebraicGeometry.BasisCovSystem, AlgebraicGeometry.CovDatum}
  \dcref{custom}

  \uses{mk-def:cech_complex, mk-lem:ses_cech_h1, mk-lem:injective_cech_acyclic}
  \textit{Project-bespoke encoding of the hypotheses of Stacks 01EO.}
  A \emph{cover system on a basis} for a scheme \(X\) is a record with five fields. A
  \emph{covering datum} (the helper \(\mathrm{CovDatum}\)) is a pair \((\iota, U)\) of an index type
  \(\iota\) together with an index family \(U : \iota \to \mathrm{Opens}(X)\); it is the raw shape
  \(\Sigma\,\iota,\ \iota \to \mathrm{Opens}(X)\) over which the section {\v C}ech complex is built.
  The cover system bundles:
  \begin{enumerate}
  \item a basis \(\mathcal{B} \subseteq \mathrm{Opens}(X)\) for the topology;
  \item a set \(\mathrm{Cov}\) of admissible coverings, each a \(\mathrm{CovDatum}\);
  \item \emph{faces-in-basis} (\(\mathrm{faces\_mem}\)): for every covering in \(\mathrm{Cov}\) every
    finite intersection \(U_{i_0 \ldots i_p}\) of its opens lies in \(\mathcal{B}\) (condition (1) of
    Lemma~\ref{mk-lem:cech_to_cohomology_on_basis});
  \item \emph{section surjectivity} (\(\mathrm{surj\_of\_vanishing}\)): for every \(V \in \mathcal{B}\)
    and every short exact sequence \(S : 0 \to S_1 \to S_2 \to S_3 \to 0\) of \(\mathcal{O}_X\)-modules
    whose left term \(S_1\) has vanishing positive {\v C}ech cohomology over the coverings of \(V\) in
    \(\mathrm{Cov}\), the section map \(S_2(V) \to S_3(V)\) is surjective;
  \item \emph{injective acyclicity} (\(\mathrm{injective\_acyclic}\)): for every injective
    \(\mathcal{O}_X\)-module \(\mathcal{I}\) and every covering \(\mathcal{U} \in \mathrm{Cov}\) the
    positive-degree {\v C}ech cohomology vanishes, \(\check{H}^p(\mathcal{U}, \mathcal{I}) = 0\) for
    all \(p > 0\).
  \end{enumerate}
  Unlike a bare combinatorial record, the structure carries two sheaf-theoretic ({\v C}ech-cohomology)
  fields. The field \(\mathrm{surj\_of\_vanishing}\) is \emph{not} the raw cofinality datum: it is the
  section-surjectivity \emph{output} that cofinality together with Lemma~\ref{mk-lem:ses_cech_h1}
  produces, packaged directly as the consequence the 01EO chain consumes. The field
  \(\mathrm{injective\_acyclic}\) is the injective {\v C}ech-acyclicity of
  Lemma~\ref{mk-lem:injective_cech_acyclic} for the system's own coverings. These are exactly the data
  that the dimension-shift induction Lemma~\ref{mk-lem:absolute_cohomology_pos_vanishing} consumes.
\end{definition}

\begin{definition}[Vanishing higher {\v C}ech cohomology for a cover system]\label{mk-def:has_vanishing_higher_cech}
  \lean{AlgebraicGeometry.HasVanishingHigherCech}
  \leanok
  \dcref{01XB,custom}

  \uses{mk-def:basis_cov_system, mk-def:cech_cohomology_accessor}
  \textit{Project-bespoke inductive predicate (condition (3) of Stacks 01EO, per module).}
  Given a cover system \(s\) on a basis (Definition~\ref{mk-def:basis_cov_system}), an
  \(\mathcal{O}_X\)-module \(\mathcal{F}\) \emph{has vanishing higher {\v C}ech cohomology for
  \(s\)} when
  \[
    \check{H}^p(\mathcal{U}, \mathcal{F}) = 0 \qquad
      \text{for every } \mathcal{U} \in s.\mathrm{Cov} \text{ and every } p > 0,
  \]
  with \(\check{H}^p\) the accessor of Definition~\ref{mk-def:cech_cohomology_accessor}. This is the
  per-module form of condition (3). It is deliberately stated for an \emph{arbitrary}
  \(\mathcal{O}_X\)-module: the dimension-shift induction of
  Lemma~\ref{mk-lem:absolute_cohomology_pos_vanishing} feeds the quotient \(\mathcal{Q} =
  \mathcal{I}/\mathcal{F}\) back as the coefficient, and \(\mathcal{Q}\) need not be quasi-coherent
  even when \(\mathcal{F}\) is. The inductive class is therefore ``\(\mathcal{O}_X\)-modules with
  vanishing higher {\v C}ech cohomology for \(s\)'', a quasi-coherent \(\mathcal{F}\) entering only
  as the seed at the affine (Tag 01XB) instantiation; the predicate must \emph{not} be specialized
  to quasi-coherent modules.
\end{definition}

\begin{lemma}
[Product of short exact sequences in \(\mathrm{Ab}\) (AB4*)]
  \label{mk-lem:short_exact_pi_map}
  \lean{AlgebraicGeometry.shortExact_piMap, AlgebraicGeometry.pi_π_map_apply}
  \dcref{custom}

  \uses{mk-def:cech_complex}
  Let \((S_j)_{j \in J}\) be a family, indexed by an arbitrary type \(J\), of short exact sequences
  of abelian groups, \(S_j : 0 \to A_j \xrightarrow{f_j} B_j \xrightarrow{g_j} C_j \to 0\). Then
  the short complex assembled from the product maps,
  \[
    0 \to \prod_{j \in J} A_j \xrightarrow{\;\prod_j f_j\;} \prod_{j \in J} B_j
      \xrightarrow{\;\prod_j g_j\;} \prod_{j \in J} C_j \to 0,
  \]
  is again short exact. This is the AB4* exactness-of-products property of \(\mathrm{Ab}\); it is
  the degreewise engine of Lemma~\ref{mk-lem:cech_ses_of_basis}, since each {\v C}ech term is a product
  of basic-open section groups (Definition~\ref{mk-def:cech_complex}).
\end{lemma}
\begin{proof}
  Monomorphy of the product map \(\prod_j f_j\) is automatic, as products preserve monomorphisms.
  Exactness at the middle is checked componentwise: an element of \(\prod_j B_j\) annihilated by
  \(\prod_j g_j\) is annihilated in each coordinate \(j\), hence lies in the image of \(f_j\) by
  exactness of \(S_j\), and the resulting coordinatewise preimages assemble to a preimage in
  \(\prod_j A_j\). The one nontrivial point is epimorphy of \(\prod_j g_j\), which is \emph{not} an
  automatic instance in \(\mathrm{Ab}\): given a target tuple in \(\prod_j C_j\), surjectivity of
  each \(g_j\) supplies a coordinatewise preimage, and these assemble — through the canonical
  identification of an \(\mathrm{Ab}\)-valued product with the set-theoretic product of its
  components — to a preimage of the whole tuple. Hence \(\prod_j g_j\) is an epimorphism and the
  displayed sequence is short exact.
\end{proof}

\begin{lemma}[Short exact sequence of {\v C}ech complexes from a basis]\label{mk-lem:cech_ses_of_basis}
  \lean{AlgebraicGeometry.cechComplex_shortExact_of_basis}
  \leanok
  \dcref{01EO,custom}

  \uses{mk-def:cech_complex, mk-def:face_short_complex, mk-def:section_cech_short_complex,
    mk-lem:ses_cech_h1, mk-lem:short_exact_pi_map}
  \textit{Source: Stacks Project, Cohomology, Tag 01EO,
  \texttt{lemma-cech-vanish-basis} (proof).}
  Fix an index family \(U : \iota \to \mathrm{Opens}(X)\) (giving the covering \(\mathcal{U}\)) and a
  short complex
  \[
    P : \mathcal{P}_1 \to \mathcal{P}_2 \to \mathcal{P}_3
  \]
  of presheaves of \(\mathcal{O}_X\)-modules. Assume that for every degree \(p\) and every tuple
  \(\sigma : \{0, \ldots, p\} \to \iota\) the per-face short complex of
  Definition~\ref{mk-def:face_short_complex} — the section sequence
  \[
    \mathcal{P}_1(V_\sigma) \longrightarrow \mathcal{P}_2(V_\sigma) \longrightarrow
      \mathcal{P}_3(V_\sigma), \qquad V_\sigma = \bigcap_k U(\sigma k)
  \]
  — is short exact. Then the induced short complex of section {\v C}ech complexes of
  Definition~\ref{mk-def:section_cech_short_complex},
  \[
    0 \to \check{\mathcal{C}}^\bullet(\mathcal{U}, \mathcal{P}_1) \to
      \check{\mathcal{C}}^\bullet(\mathcal{U}, \mathcal{P}_2) \to
      \check{\mathcal{C}}^\bullet(\mathcal{U}, \mathcal{P}_3) \to 0,
  \]
  is short exact (as a short complex of cochain complexes of abelian groups).

  \emph{Remark.} In the cover-global \((\mathcal{B}, \mathrm{Cov})\) setting of
  Lemma~\ref{mk-lem:cech_to_cohomology_on_basis}, where \(P\) is the section short complex of a short
  exact sequence \(0 \to \mathcal{F} \to \mathcal{I} \to \mathcal{Q} \to 0\) of sheaves, the per-face
  hypothesis is discharged by the per-face short-exact-sequence derivation of
  Lemma~\ref{mk-lem:face_ses_of_sheaf_ses}: left-exactness of the section sequence over \(V_\sigma\) is
  automatic, and surjectivity \(\mathcal{I}(V_\sigma) \to \mathcal{Q}(V_\sigma)\) is the output of
  Lemma~\ref{mk-lem:ses_cech_h1} at the basic open \(V_\sigma\) (condition (1) placing \(V_\sigma \in
  \mathcal{B}\)).
\end{lemma}

\begin{proof}
  \leanok
  \uses{mk-def:section_cech_short_complex, mk-def:face_short_complex, mk-lem:short_exact_pi_map}
  Short-exactness of a short complex of cochain complexes is checked degreewise: it suffices to
  verify short-exactness in each cohomological degree \(p\). In degree \(p\) the term of the
  section {\v C}ech complex is the product \(\prod_{\sigma}(-)(V_\sigma)\) over all tuples
  \(\sigma : \{0, \ldots, p\} \to \iota\) of section groups over the basic opens \(V_\sigma\), and
  the two maps are the products of the corresponding per-face maps; thus the degree-\(p\) short
  complex is the product over \(\sigma\) of the per-face short complexes
  (Definition~\ref{mk-def:face_short_complex}). Each of these is short exact by hypothesis, so their
  product is short exact by the AB4* product-of-short-exact-sequences lemma
  (Lemma~\ref{mk-lem:short_exact_pi_map}). As this holds in every degree, the short complex of section
  {\v C}ech complexes (Definition~\ref{mk-def:section_cech_short_complex}) is short exact.
\end{proof}

\begin{lemma}

[Per-face short exact sequence from a sheaf short exact sequence]
  \label{mk-lem:face_ses_of_sheaf_ses}
  \lean{AlgebraicGeometry.faceShortComplex_shortExact_of_sheaf_ses,
    AlgebraicGeometry.sectionsFunctor}
  \dcref{01EO,custom}

  \uses{mk-lem:ses_cech_h1, mk-def:face_short_complex}
  \textit{Source: Stacks Project, Cohomology, Tag 01EO,
  \texttt{lemma-cech-vanish-basis} (proof).}
  Let \(S : 0 \to \mathcal{F} \to \mathcal{I} \to \mathcal{Q} \to 0\) be a short exact sequence of
  sheaves of \(\mathcal{O}_X\)-modules, let \(U : \iota \to \mathrm{Opens}(X)\) be an index family,
  and suppose that on every face \(V_\sigma = \bigcap_k U(\sigma k)\) the section map
  \(\mathcal{I}(V_\sigma) \to \mathcal{Q}(V_\sigma)\) is surjective — the basis-surjectivity datum
  supplied by Lemma~\ref{mk-lem:ses_cech_h1}. Writing \(P\) for the section short complex of \(S\)
  regarded as a short complex of presheaves of modules, the per-face short complex of
  Definition~\ref{mk-def:face_short_complex} is then short exact for every degree \(p\) and tuple
  \(\sigma\). This lemma produces the per-face hypothesis consumed by
  Lemma~\ref{mk-lem:cech_ses_of_basis}, and is the bridge by which
  Lemma~\ref{mk-lem:absolute_cohomology_pos_vanishing} instantiates the cover-local lemmas from a
  genuine sheaf short exact sequence.
\end{lemma}
\begin{proof}
  \uses{mk-lem:ses_cech_h1, mk-def:face_short_complex}
  Fix \(p\) and \(\sigma\), and put \(V = V_\sigma = \bigcap_k U(\sigma k)\). Monomorphy of
  \(\mathcal{F}(V) \to \mathcal{I}(V)\) and exactness at the middle of \(\mathcal{F}(V) \to
  \mathcal{I}(V) \to \mathcal{Q}(V)\) are the left-exactness of the sections functor applied to the
  short exact sequence of sheaves \(S\): the sections functor \(\mathcal{G} \mapsto \mathcal{G}(V)\)
  is the composite of the inclusion of sheaves of \(\mathcal{O}_X\)-modules into presheaves of
  modules (a right adjoint, hence preserving finite limits) with evaluation at \(V\) (which
  preserves limits), so it is left exact and carries the exactness of \(S\) to exactness of the
  section sequence at \(\mathcal{I}(V)\). The remaining map \(\mathcal{I}(V) \to \mathcal{Q}(V)\) is
  surjective by hypothesis, i.e. by the output of Lemma~\ref{mk-lem:ses_cech_h1} at the basic open
  \(V\). Hence the per-face short complex is short exact.
\end{proof}

\begin{lemma}[Quotient preserves vanishing positive homology (homological core)]\label{mk-lem:cech_homology_quotient_vanishing}
  \lean{AlgebraicGeometry.cechHomology_quotient_vanishing}
  \leanok
  \dcref{custom}

  \uses{mk-lem:homology_long_exact_sequence}
  \textit{Abstract homological core of Lemma~\ref{mk-lem:quotient_vanishing_cech}.}
  Let \(T : 0 \to T_1 \to T_2 \to T_3 \to 0\) be a short exact sequence of cochain complexes
  (cohomological \(\mathbb{N}\)-grading) of abelian groups, and suppose the middle term \(T_2\) and
  the left term \(T_1\) both have vanishing positive homology: \(H^p(T_2) = 0\) and \(H^p(T_1) = 0\)
  for all \(p > 0\). Then the right term \(T_3\) also has vanishing positive homology, \(H^p(T_3) =
  0\) for all \(p > 0\). This is the section-Čech-free homological core consumed by
  Lemma~\ref{mk-lem:quotient_vanishing_cech}.
\end{lemma}
\begin{proof}
  \leanok
  \uses{mk-lem:homology_long_exact_sequence}
  Fix \(p > 0\). In the homology long exact sequence of \(T\)
  (Lemma~\ref{mk-lem:homology_long_exact_sequence}) the two homology groups of the middle term
  adjacent to the connecting map at \(T_3\) in degree \(p\), namely \(H^p(T_2)\) and \(H^{p+1}(T_2)\),
  both vanish (positive degrees, by the hypothesis on \(T_2\)); hence the connecting homomorphism is
  an isomorphism \(H^p(T_3) \cong H^{p+1}(T_1)\). The right-hand side vanishes since \(p + 1 > 0\)
  and \(T_1\) has vanishing positive homology. Therefore \(H^p(T_3) = 0\).
\end{proof}

\begin{lemma}[Quotient preserves vanishing higher {\v C}ech cohomology]\label{mk-lem:quotient_vanishing_cech}
  \lean{AlgebraicGeometry.quotient_cech_vanishing_of_basis}
  \leanok
  \dcref{01EO,custom}

  \uses{mk-def:cech_complex, mk-def:cech_cohomology_accessor, mk-lem:cech_ses_of_basis,
    mk-lem:injective_cech_acyclic, mk-lem:cech_homology_quotient_vanishing}
  \textit{Source: Stacks Project, Cohomology, Tag 01EO,
  \texttt{lemma-cech-vanish-basis} (proof).}
  Fix an index family \(U : \iota \to \mathrm{Opens}(X)\) (giving the covering \(\mathcal{U}\)) and a
  short complex \(P : \mathcal{P}_1 \to \mathcal{P}_2 \to \mathcal{P}_3\) of presheaves of
  \(\mathcal{O}_X\)-modules, as in Lemma~\ref{mk-lem:cech_ses_of_basis}. Suppose:
  \begin{enumerate}
  \item the short complex of section {\v C}ech complexes of \(P\) is short exact (the conclusion of
    Lemma~\ref{mk-lem:cech_ses_of_basis});
  \item the middle term \(\mathcal{P}_2\) has vanishing positive {\v C}ech cohomology,
    \(\check{H}^p(\mathcal{U}, \mathcal{P}_2) = 0\) for all \(p > 0\); and
  \item the left term \(\mathcal{P}_1\) has vanishing positive {\v C}ech cohomology,
    \(\check{H}^p(\mathcal{U}, \mathcal{P}_1) = 0\) for all \(p > 0\).
  \end{enumerate}
  Then the right term \(\mathcal{P}_3\) (the quotient) also has vanishing positive {\v C}ech
  cohomology: \(\check{H}^p(\mathcal{U}, \mathcal{P}_3) = 0\) for all \(p > 0\). Here
  \(\check{H}^p\) is the accessor of Definition~\ref{mk-def:cech_cohomology_accessor}. At instantiation
  hypothesis (2) is supplied for \(\mathcal{P}_2 = \mathcal{I}\) injective by
  Lemma~\ref{mk-lem:injective_cech_acyclic}, and hypothesis (3) for \(\mathcal{P}_1 = \mathcal{F}\) by
  condition (3) of Lemma~\ref{mk-lem:cech_to_cohomology_on_basis}.
\end{lemma}

\begin{proof}
  \leanok
  \uses{mk-lem:cech_homology_quotient_vanishing, mk-lem:injective_cech_acyclic}
  Apply the abstract homological core (Lemma~\ref{mk-lem:cech_homology_quotient_vanishing}) to the
  short exact sequence \(T\) of cochain complexes furnished by hypothesis (1) — the short complex of
  section {\v C}ech complexes of \(P\). Its homology in degree \(p\) is by definition
  \(\check{H}^p(\mathcal{U}, -)\) (Definition~\ref{mk-def:cech_cohomology_accessor}), so hypotheses (2)
  and (3) are exactly the vanishing of the positive homology of the middle and left terms of \(T\).
  The core lemma then yields the vanishing of the positive homology of the right term, that is
  \(\check{H}^p(\mathcal{U}, \mathcal{P}_3) = 0\) for all \(p > 0\). Concretely, the connecting
  isomorphism of the homology long exact sequence realizes the dimension shift
  \(\check{H}^p(\mathcal{U}, \mathcal{P}_3) \cong \check{H}^{p+1}(\mathcal{U}, \mathcal{P}_1)\), the
  step used in the 01EO proof; at instantiation \(\check{H}^{>0}(\mathcal{U}, \mathcal{P}_2) = 0\) is
  the injective acyclicity of Lemma~\ref{mk-lem:injective_cech_acyclic}.
\end{proof}

\begin{lemma}[Sheaf-cohomology base case \(H^1(U, \mathcal{F}) = 0\)]\label{mk-lem:absolute_cohomology_one_vanishing}
  \lean{AlgebraicGeometry.absoluteCohomology_one_eq_zero_of_basis}
  \leanok
  \dcref{01EO,custom}

  \uses{mk-def:absolute_cohomology, mk-lem:cech_ses_of_basis,
    mk-lem:absolute_cohomology_covariant_les, mk-lem:absolute_cohomology_injective_vanishing,
    mk-lem:absolute_cohomology_zero, mk-lem:absolute_cohomology_zero_natural}
  \textit{Source: Stacks Project, Cohomology, Tag 01EO,
  \texttt{lemma-cech-vanish-basis} (proof).}
  Let \(U \subseteq X\) be an open, and let
  \(S : 0 \to \mathcal{F} \to \mathcal{I} \to \mathcal{Q} \to 0\) be a short exact sequence in
  \(X.\mathrm{Modules}\) with \(\mathcal{I}\) injective. Assume that the section map
  \(\mathcal{I}(U) \to \mathcal{Q}(U)\) — the action on sections over \(U\) of the surjection of
  \(S\) — is surjective. Then the first absolute cohomology vanishes,
  \[
    H^1(U, \mathcal{F}) = 0,
  \]
  where \(H^p(U, -)\) is the Ext realization of Definition~\ref{mk-def:absolute_cohomology}. The only
  geometric input is the section-surjectivity hypothesis; at instantiation it is supplied, at a basic
  open \(U \in \mathcal{B}\), by the per-face surjectivity of Lemma~\ref{mk-lem:ses_cech_h1} (through
  Lemma~\ref{mk-lem:cech_ses_of_basis}).
\end{lemma}

\begin{proof}
  \leanok
  \uses{mk-def:absolute_cohomology, mk-lem:cech_ses_of_basis,
    mk-lem:absolute_cohomology_covariant_les, mk-lem:absolute_cohomology_injective_vanishing,
    mk-lem:absolute_cohomology_zero, mk-lem:absolute_cohomology_zero_natural}
    Run the covariant Ext long exact sequence
  (Lemma~\ref{mk-lem:ext_covariant_les_mathlib}) of
  \(0 \to \mathcal{F} \to \mathcal{I} \to \mathcal{Q} \to 0\) at the fixed first
  argument \(j_!\mathcal{O}_U\), giving the exact strand
  \[
    H^0(U, \mathcal{I}) \xrightarrow{\ g\ } H^0(U, \mathcal{Q})
      \xrightarrow{\ \delta\ } H^1(U, \mathcal{F}) \to H^1(U, \mathcal{I}).
  \]
  By Lemma~\ref{mk-lem:ext_eq_zero_of_injective_mathlib} (the injective \(\mathcal{I}\) is
  the \emph{second} Ext argument, Lemma~\ref{mk-lem:absolute_cohomology_injective_vanishing})
  the term \(H^1(U, \mathcal{I}) =
  \operatorname{Ext}^1(j_!\mathcal{O}_U, \mathcal{I}) = 0\). In degree \(0\) the
  identification \(H^0(U, -) = \Gamma(U, -)\) of Lemma~\ref{mk-lem:absolute_cohomology_zero}
  is natural in the coefficient sheaf
  (Lemma~\ref{mk-lem:absolute_cohomology_zero_natural}), so it carries \(g = H^0(U,
  \mathcal{I} \to \mathcal{Q})\) to the section map
  \(\mathcal{I}(U) \to \mathcal{Q}(U)\); the latter is surjective by hypothesis, hence so is
  \(g\). Surjectivity of \(g\) makes the connecting map
  \(\delta = 0\); since also \(H^1(U, \mathcal{I}) = 0\), the strand reads
  \(0 \to H^1(U, \mathcal{F}) \to 0\), whence \(H^1(U, \mathcal{F}) = 0\).
\end{proof}

\begin{lemma}

[Dimension-shift induction \(H^p(U, \mathcal{F}) = 0\) for \(p > 0\)]
  \label{mk-lem:absolute_cohomology_pos_vanishing}
  \lean{AlgebraicGeometry.absoluteCohomology_eq_zero_of_basis,
    AlgebraicGeometry.injSES, AlgebraicGeometry.injSES_shortExact}
  \dcref{01EO,custom}

  \uses{mk-def:absolute_cohomology, mk-def:basis_cov_system, mk-def:has_vanishing_higher_cech,
    mk-lem:quotient_vanishing_cech, mk-lem:absolute_cohomology_one_vanishing,
    mk-lem:face_ses_of_sheaf_ses, mk-lem:ext_covariant_les_mathlib,
    mk-lem:ext_eq_zero_of_injective_mathlib}
  \textit{Source: Stacks Project, Cohomology, Tag 01EO,
  \texttt{lemma-cech-vanish-basis} (proof).}
  Let \(s\) be a cover system on a basis for \(X\) (Definition~\ref{mk-def:basis_cov_system}). For every
  \(\mathcal{O}_X\)-module \(\mathcal{F}\) that has vanishing higher {\v C}ech cohomology for \(s\)
  (Definition~\ref{mk-def:has_vanishing_higher_cech}), every basic open \(U \in s.\mathcal{B}\), and
  every \(p > 0\), the absolute cohomology vanishes,
  \[
    H^p(U, \mathcal{F}) = 0,
  \]
  where \(H^p(U, -)\) is the Ext realization of Definition~\ref{mk-def:absolute_cohomology}. The
  formal statement carries \([\,\mathrm{EnoughInjectives}\ X.\mathrm{Modules}\,]\) as an explicit
  instance hypothesis: enough injectives in the category of sheaves of
  \(\mathcal{O}_X\)-modules is needed to embed \(\mathcal{F}\) into an injective \(\mathcal{I}\), and
  that instance is not currently available in the ambient library for sheaves of modules (it would
  follow from \(\mathrm{IsGrothendieckAbelian}(\mathrm{SheafOfModules}\ R)\)).
\end{lemma}

\begin{proof}
  \uses{mk-def:basis_cov_system, mk-def:has_vanishing_higher_cech, mk-lem:quotient_vanishing_cech,
    mk-lem:absolute_cohomology_one_vanishing, mk-lem:face_ses_of_sheaf_ses, mk-lem:cech_ses_of_basis,
    mk-lem:ses_cech_h1, mk-lem:ext_covariant_les_mathlib, mk-lem:ext_eq_zero_of_injective_mathlib}
  Induct on \(p \ge 1\) with the statement quantified over \emph{all} \(\mathcal{O}_X\)-modules
  \(\mathcal{F}\) having vanishing higher {\v C}ech cohomology for \(s\)
  (Definition~\ref{mk-def:has_vanishing_higher_cech}) simultaneously, so that the quotient may take the
  role of \(\mathcal{F}\) at the next stage. Fix such an \(\mathcal{F}\), embed
  \(\mathcal{F} \hookrightarrow \mathcal{I}\) into an injective \(\mathcal{O}_X\)-module, and put
  \(\mathcal{Q} = \mathcal{I}/\mathcal{F}\), giving a short exact sequence \(S : 0 \to \mathcal{F}
  \to \mathcal{I} \to \mathcal{Q} \to 0\).

  We first record the geometric input that both the quotient-closure and the base case need. For
  every \(\mathcal{U} \in s.\mathrm{Cov}\), the faces-in-basis field \(\mathrm{faces\_mem}\) (condition
  (1)) places every face \(V_\sigma\) in \(s.\mathcal{B}\); the section-surjectivity field
  \(\mathrm{surj\_of\_vanishing}\) of the cover system, applied to the short exact sequence \(S\) at
  \(V_\sigma\), then gives section surjectivity \(\mathcal{I}(V_\sigma) \to \mathcal{Q}(V_\sigma)\).
  Its left-term {\v C}ech-vanishing hypothesis \(\check{H}^{>0}(-,\mathcal{F}) = 0\) over the coverings
  of \(V_\sigma\) is supplied by \(\mathcal{F}\)'s vanishing higher {\v C}ech cohomology
  (Definition~\ref{mk-def:has_vanishing_higher_cech}); mathematically this field packages
  Lemma~\ref{mk-lem:ses_cech_h1} (cofinality plus the {\v C}ech \(\check{H}^1\)-surjectivity criterion).
  By the per-face derivation
  (Lemma~\ref{mk-lem:face_ses_of_sheaf_ses}) the per-face short complexes of \(S\) are short exact, so
  Lemma~\ref{mk-lem:cech_ses_of_basis} gives the short exact sequence of section {\v C}ech complexes for
  \(\mathcal{U}\); in particular the section surjectivity \(\mathcal{I}(U) \to \mathcal{Q}(U)\) holds
  at every basic open \(U \in s.\mathcal{B}\) (take \(\mathcal{U}\) a covering of \(U\) in
  \(s.\mathrm{Cov}\) and the length-one face). Feeding the {\v C}ech short exact sequence into
  Lemma~\ref{mk-lem:quotient_vanishing_cech} — with \(\check{H}^{>0}(\mathcal{U}, \mathcal{I}) = 0\)
  from injective acyclicity and \(\check{H}^{>0}(\mathcal{U}, \mathcal{F}) = 0\) from
  \(\mathcal{F}\)'s hypothesis — shows \(\check{H}^{>0}(\mathcal{U}, \mathcal{Q}) = 0\) for every
  \(\mathcal{U} \in s.\mathrm{Cov}\); that is, \(\mathcal{Q}\) again has vanishing higher {\v C}ech
  cohomology for \(s\), so it is itself a valid instance of the inductive statement.

  \emph{Base case \(p = 1\):} apply Lemma~\ref{mk-lem:absolute_cohomology_one_vanishing} to \(S\) at
  the basic open \(U \in s.\mathcal{B}\), its section-surjectivity hypothesis discharged by the
  paragraph above.

  \emph{Step \(p \Rightarrow p+1\):} the covariant Ext long exact sequence
  (Lemma~\ref{mk-lem:ext_covariant_les_mathlib}) of
  \(0 \to \mathcal{F} \to \mathcal{I} \to \mathcal{Q} \to 0\) at \(j_!\mathcal{O}_U\)
  gives the exact strand
  \[
    H^p(U, \mathcal{Q}) \xrightarrow{\ \delta\ } H^{p+1}(U, \mathcal{F}) \to
      H^{p+1}(U, \mathcal{I}).
  \]
  The right term vanishes since \(\mathcal{I}\) is injective
  (Lemma~\ref{mk-lem:ext_eq_zero_of_injective_mathlib}), and the left term
  \(H^p(U, \mathcal{Q}) = 0\) by the inductive hypothesis applied to \(\mathcal{Q}\)
  (admissible by the previous paragraph). Exactness forces
  \(H^{p+1}(U, \mathcal{F}) = 0\). This completes the induction, giving
  \(H^p(U, \mathcal{F}) = 0\) for all \(p > 0\) and all \(U \in s.\mathcal{B}\).
\end{proof}

\begin{lemma}[{\v C}ech-to-cohomology comparison on a basis]\label{mk-lem:cech_to_cohomology_on_basis}
  \lean{AlgebraicGeometry.cech_eq_cohomology_of_basis}
  \leanok
  \dcref{01EO,custom}

  \uses{mk-def:cech_complex, mk-def:absolute_cohomology, mk-lem:cech_ses_of_basis,
    mk-lem:quotient_vanishing_cech, mk-lem:absolute_cohomology_one_vanishing,
    mk-lem:absolute_cohomology_pos_vanishing}
  \textit{Source: Stacks Project, Cohomology, Tag 01EO, Lemma
  \texttt{lemma-cech-vanish-basis}.}
  Let \(X\) be a ringed space, let \(\mathcal{B}\) be a basis for its topology, let
  \(\mathcal{F}\) be an \(\mathcal{O}_X\)-module, and suppose there is a set
  \(\mathrm{Cov}\) of open coverings such that: (1) every \(\mathcal{U} :
  U = \bigcup_{i \in I} U_i\) in \(\mathrm{Cov}\) has \(U, U_i \in \mathcal{B}\) and
  all intersections \(U_{i_0 \ldots i_p} \in \mathcal{B}\); (2) for each
  \(U \in \mathcal{B}\) the coverings of \(U\) occurring in \(\mathrm{Cov}\) are
  cofinal among all open coverings of \(U\); and (3) for every \(\mathcal{U} \in
  \mathrm{Cov}\) the {\v C}ech cohomology vanishes,
  \(\check{H}^p(\mathcal{U}, \mathcal{F}) = 0\) for all \(p > 0\). Then
  \[
    H^p(U, \mathcal{F}) = 0 \qquad \text{for all } p > 0 \text{ and any }
      U \in \mathcal{B}.
  \]
  This is the basis-comparison vehicle that upgrades the standard-cover {\v C}ech
  vanishing of Lemma~\ref{mk-lem:cech_acyclic_affine} to the Serre vanishing of
  Lemma~\ref{mk-lem:affine_serre_vanishing}: instantiating \(\mathcal{B}\) as the affine
  opens and \(\mathrm{Cov}\) as the standard affine covers, condition (3) is exactly
  Lemma~\ref{mk-lem:cech_acyclic_affine}. The formal target of this lemma is the
  \emph{vanishing} conclusion \(H^p(U, \mathcal{F}) = 0\) for \(p > 0\) under
  hypotheses (1)--(3), not an explicit isomorphism \(\check{H}^p(\mathcal{U},
  \mathcal{F}) \cong H^p(U, \mathcal{F})\) of {\v C}ech and sheaf cohomology; only
  the positive-degree vanishing is needed downstream, and it is all the statement
  establishes. The formal statement carries
  \([\,\mathrm{EnoughInjectives}\ X.\mathrm{Modules}\,]\) as an explicit instance hypothesis,
  inherited from the dimension-shift induction (Lemma~\ref{mk-lem:absolute_cohomology_pos_vanishing}):
  it is needed to embed the coefficient module into an injective, and that instance is not currently
  available in the ambient library for sheaves of \(\mathcal{O}_X\)-modules (it would follow from
  \(\mathrm{IsGrothendieckAbelian}(\mathrm{SheafOfModules}\ R)\)).
\end{lemma}
\begin{proof}
  \leanok
  \uses{mk-def:basis_cov_system, mk-def:has_vanishing_higher_cech,
    mk-lem:absolute_cohomology_pos_vanishing, mk-def:absolute_cohomology}
    Here \(H^p(U, -)\) denotes the Ext realization of absolute sheaf cohomology fixed in
  Definition~\ref{mk-def:absolute_cohomology}. The proof is the thin assembly of the sub-lemmas through
  the cover-system encoding. Conditions (1)--(2) of the hypotheses are exactly the faces-in-basis and
  cofinality data of a cover system on a basis, so \(\mathcal{B}\) and \(\mathrm{Cov}\) assemble into
  an instance \(s\) of Definition~\ref{mk-def:basis_cov_system} with \(s.\mathcal{B} = \mathcal{B}\) and
  \(s.\mathrm{Cov} = \mathrm{Cov}\); condition (3) is exactly the assertion that \(\mathcal{F}\) has
  vanishing higher {\v C}ech cohomology for \(s\) (Definition~\ref{mk-def:has_vanishing_higher_cech}).

  The conclusion is then the dimension-shift induction
  Lemma~\ref{mk-lem:absolute_cohomology_pos_vanishing} applied to this instance: for every basic open
  \(U \in s.\mathcal{B} = \mathcal{B}\) and every \(p > 0\) one has \(H^p(U, \mathcal{F}) = 0\). That
  induction internally embeds \(\mathcal{F}\) into an injective \(\mathcal{I}\), forms the quotient
  \(\mathcal{Q} = \mathcal{I}/\mathcal{F}\), keeps \(\mathcal{Q}\) inside the inductive class via the
  {\v C}ech short exact sequence (Lemma~\ref{mk-lem:cech_ses_of_basis}) and its quotient-stability
  (Lemma~\ref{mk-lem:quotient_vanishing_cech}), and bottoms out at the base case
  (Lemma~\ref{mk-lem:absolute_cohomology_one_vanishing}).

  The argument never uses affine sheaf-cohomology vanishing
  (Lemma~\ref{mk-lem:affine_serre_vanishing}); its only {\v C}ech-vanishing input is condition (3),
  which on the affine instantiation is supplied by the standard-cover {\v C}ech vanishing of
  Lemma~\ref{mk-lem:cech_acyclic_affine}. Thus the dependency
  \(\text{affine Serre vanishing} \to \text{this lemma}\) is one-directional and the graph is
  acyclic.
\end{proof}

\section{The {\v C}ech complex computes $R^i f_*$}

The comparison is proved by the \emph{acyclic-resolution} (Cartan--Leray) route. The
augmented {\v C}ech complex is a resolution of \(\mathcal{F}\) whose terms are
acyclic for \(f_*\); the abstract homological-algebra
Lemma~\ref{mk-lem:acyclic_resolution_computes_derived} then identifies the derived
functor with the cohomology of the complex obtained by applying \(f_*\). The two
geometric inputs are isolated in the following sub-lemmas.

\subsection{Object layer for the augmented {\v C}ech complex}

Before stating the resolution lemma we name the object it is about: the augmented {\v C}ech complex
\(0 \to \mathcal{F} \to \mathcal{C}^0 \to \mathcal{C}^1 \to \cdots\) as a cochain complex in the
category \(X.\mathrm{Modules}\) of sheaves of \(\mathcal{O}_X\)-modules, together with its augmentation
map and the identity \(\varepsilon \cdot d^0 = 0\) that makes the augmentation legitimate. This object
layer is purely formal --- it uses only the {\v C}ech nerve of Definition~\ref{mk-def:cech_nerve} and the
unitors of the push--pull adjunction --- and is independent of the exactness (acyclicity) statement.

\begin{definition}[Un-augmented {\v C}ech cochain complex on \(X\)]\label{mk-def:cech_complex_on_X}
  \lean{AlgebraicGeometry.cechComplexOnX}
  \leanok
  \dcref{custom}

  \uses{mk-def:cech_nerve, mk-def:relative_cech_complex_of_nerve}
  \textit{Project-local.}
  Given an open cover \(\mathcal{U}\) of \(X\) and a sheaf of \(\mathcal{O}_X\)-modules \(\mathcal{F}\),
  the \emph{un-augmented {\v C}ech cochain complex}
  \[
    \mathcal{C}^0 \to \mathcal{C}^1 \to \mathcal{C}^2 \to \cdots
  \]
  on \(X\) is the alternating-coface-map cochain complex of the cosimplicial \(\mathcal{O}_X\)-module
  obtained by forgetting the augmentation of the {\v C}ech nerve of Definition~\ref{mk-def:cech_nerve}
  (its \(\mathrm{Augmented.drop}\)). Its degree-\(p\) term is \(\mathcal{C}^p =
  \prod_{(i_0, \ldots, i_p)} (j_{i_0 \ldots i_p})_*\bigl(\mathcal{F}|_{U_{i_0 \ldots i_p}}\bigr)\),
  exactly as in the relative {\v C}ech complex of
  Definition~\ref{mk-def:relative_cech_complex_of_nerve} specialised to \(f = \mathrm{id}_X\).
\end{definition}

\begin{definition}[Augmentation point isomorphism]\label{mk-def:cech_nerve_point_iso}
  \lean{AlgebraicGeometry.cechNervePointIso}
  \leanok
  \dcref{custom}

  \uses{mk-def:cech_nerve}
  \textit{Project-local.}
  The augmentation point of the {\v C}ech nerve of Definition~\ref{mk-def:cech_nerve} is the value of the
  nerve at the initial augmentation object, namely \((\mathrm{id}_X)_* (\mathrm{id}_X)^* \mathcal{F}\).
  The \emph{augmentation point isomorphism} is the canonical isomorphism
  \[
    (\mathrm{CechNerve}\ \mathcal{U}\ \mathcal{F}).\mathrm{left} \;\cong\; \mathcal{F}
  \]
  obtained by composing the pullback unitor \((\mathrm{id}_X)^* \mathcal{F} \cong \mathcal{F}\) with the
  pushforward unitor \((\mathrm{id}_X)_* \mathcal{G} \cong \mathcal{G}\) of the push--pull adjunction.
\end{definition}

\begin{definition}[{\v C}ech augmentation map]\label{mk-def:cech_augmentation}
  \lean{AlgebraicGeometry.cechAugmentation}
  \leanok
  \dcref{custom}

  \uses{mk-def:cech_nerve_point_iso, mk-def:cech_complex_on_X}
  \textit{Project-local.}
  The \emph{{\v C}ech augmentation} is the morphism \(\varepsilon : \mathcal{F} \to \mathcal{C}^0\) of
  sheaves of \(\mathcal{O}_X\)-modules obtained by composing the inverse of the augmentation point
  isomorphism of Definition~\ref{mk-def:cech_nerve_point_iso} with the degree-\(0\) component of the
  nerve's augmentation natural transformation (which lands in the degree-\(0\) term \(\mathcal{C}^0\) of
  the complex of Definition~\ref{mk-def:cech_complex_on_X}).
\end{definition}

\begin{lemma}
[The augmentation kills the first differential]
  \label{mk-lem:cech_augmentation_comp_d}
  \lean{AlgebraicGeometry.cechAugmentation_comp_d,
    AlgebraicGeometry.augmentation_comp_alternatingCofaceMap_objD_zero}
  \dcref{custom}

  \uses{mk-def:cech_augmentation}
  \textit{Project-local.}
  The composite of the augmentation \(\varepsilon : \mathcal{F} \to \mathcal{C}^0\) of
  Definition~\ref{mk-def:cech_augmentation} with the first {\v C}ech differential vanishes,
  \[
    \varepsilon \cdot d^0 \;=\; 0.
  \]
\end{lemma}
\begin{proof}
  \uses{mk-def:cech_augmentation}
  This is the standard augmented-cosimplicial identity. The differential \(d^0 = \delta^0 - \delta^1\)
  is the alternating sum of the two cofaces, and naturality of the augmentation natural transformation
  of the cosimplicial object intertwines \(\varepsilon\) with \(\delta^0\) and with \(\delta^1\) in the
  same way; the two contributions therefore cancel.
\end{proof}

\begin{definition}[Augmented {\v C}ech complex]\label{mk-def:cech_augmented_complex}
  \lean{AlgebraicGeometry.cechAugmentedComplex}
  \leanok
  \dcref{custom}

  \uses{mk-def:cech_complex_on_X, mk-def:cech_augmentation, mk-lem:cech_augmentation_comp_d}
  \textit{Project-local.}
  The \emph{augmented {\v C}ech complex} is the cochain complex
  \[
    0 \to \mathcal{F} \to \mathcal{C}^0 \to \mathcal{C}^1 \to \cdots
  \]
  in \(X.\mathrm{Modules}\) obtained by prepending the coefficient sheaf \(\mathcal{F}\) in degree \(0\)
  to the complex of Definition~\ref{mk-def:cech_complex_on_X} along the augmentation
  \(\varepsilon\) of Definition~\ref{mk-def:cech_augmentation}, using the augmentation identity
  \(\varepsilon \cdot d^0 = 0\) of Lemma~\ref{mk-lem:cech_augmentation_comp_d}. Concretely its degree-\(0\)
  term is \(\mathcal{F}\) and its degree-\((p+1)\) term is \(\mathcal{C}^p\). This is the object whose
  exactness is the content of Lemma~\ref{mk-lem:cech_augmented_resolution}.
\end{definition}

\begin{lemma}
[Sheafification kills a presheaf that is locally bijective to a locally-zero one]
  \label{mk-lem:sheafify_kills_locally_zero}
  \lean{CategoryTheory.GrothendieckTopology.isZero_presheafToSheaf_obj_of_W,
    CategoryTheory.GrothendieckTopology.isZero_presheafToSheaf_obj_of_W_isZero,
    CategoryTheory.GrothendieckTopology.isZero_presheafToSheaf_obj_of_isLocallyBijective}
  \dcref{custom}

  \textit{Project-bespoke site-theory lemma (pure Mathlib site theory).}
  Let \(J\) be a Grothendieck topology on a site \(\mathcal{C}\) and let \(J.W\) denote the
  class of \emph{locally bijective} maps of presheaves --- the local-isomorphism class for
  \(J\), i.e.\ the maps inverted by sheafification. Then:
  \begin{enumerate}
    \item Sheafification \(\operatorname{presheafToSheaf}\) carries every map in \(J.W\) to an
      isomorphism; consequently it sends a presheaf that is locally bijective to a presheaf
      with zero sheafification to a presheaf with zero sheafification. Precisely, if
      \(\varphi : P \to Q\) lies in \(J.W\) and the sheafification of \(Q\) is the zero
      sheaf, then the sheafification of \(P\) is the zero sheaf.
    \item In particular, if \(P\) admits a locally bijective map to the zero presheaf --- for
      instance if \(P\) is locally zero, i.e.\ vanishes on a basis of the topology --- then
      the sheafification of \(P\) is the zero sheaf.
    \item For an abelian-valued presheaf the locally-bijective hypothesis is the conjunction
      of local injectivity and local surjectivity of the comparison map; a map that is both
      locally injective and locally surjective lies in \(J.W\).
  \end{enumerate}
\end{lemma}
\begin{proof}
  Sheafification \(\operatorname{presheafToSheaf} : \widehat{\mathcal{C}} \to
  \operatorname{Sh}(J)\) inverts exactly the maps of the class \(J.W\): by definition,
  \(\varphi \in J.W\) if and only if \(\operatorname{presheafToSheaf}(\varphi)\) is an
  isomorphism. Hence for \(\varphi : P \to Q\) in \(J.W\) the induced map
  \(\operatorname{sheafify}(P) \to \operatorname{sheafify}(Q)\) is an isomorphism, and
  being a zero object transports along an isomorphism: if \(\operatorname{sheafify}(Q)\)
  is zero then so is \(\operatorname{sheafify}(P)\). This proves (1); part (2) is the
  special case \(Q = 0\) (the zero presheaf, whose sheafification is the zero sheaf),
  and a presheaf that vanishes on a basis admits such a locally bijective map to \(0\).
  For (3), sheafification is additive, and a map of abelian presheaves that is
  simultaneously locally injective and locally surjective is a local isomorphism,
  i.e.\ lies in \(J.W\); applying (1) then collapses its sheafification.
\end{proof}

\begin{lemma}

[A faithful zero-preserving functor reflects zero objects]
  \label{mk-lem:isZero_of_faithful_preservesZeroMorphisms}
  \lean{AlgebraicGeometry.isZero_of_faithful_preservesZeroMorphisms,
    CategoryTheory.Functor.isZero_of_faithful_preservesZeroMorphisms}
  \dcref{custom}

  Let \(F : \mathcal{C} \to \mathcal{D}\) be a functor between categories with zero
  morphisms that is \emph{faithful} and \emph{preserves zero morphisms}. If the image
  \(F(Y)\) of an object \(Y\) is a zero object of \(\mathcal{D}\), then \(Y\) is a zero
  object of \(\mathcal{C}\).
\end{lemma}
\begin{proof}

  An object \(Y\) is a zero object if and only if its identity morphism equals its zero
  endomorphism, \(\mathrm{id}_Y = 0\). Applying this characterization to \(F(Y)\), the
  hypothesis gives \(\mathrm{id}_{F(Y)} = 0\). Since \(F\) preserves identities and zero
  morphisms, \(F(\mathrm{id}_Y) = \mathrm{id}_{F(Y)} = 0 = F(0)\); faithfulness of \(F\)
  then makes the map on hom-sets injective, so \(\mathrm{id}_Y = 0\), whence \(Y\) is a
  zero object.
\end{proof}

\begin{lemma}[Sheafification of a presheaf that is locally a zero object vanishes]\label{mk-lem:isZero_presheafToSheaf_of_locally_isZero}
  \lean{AlgebraicGeometry.isZero_presheafToSheaf_of_locally_isZero}
  \leanok
  \dcref{custom}

  \uses{mk-lem:sheafify_kills_locally_zero}
  Let \(J\) be a Grothendieck topology on a site \(\mathcal{C}\) and let \(Q\) be a
  presheaf of abelian groups on \(\mathcal{C}\). Suppose that for every object \(U\) there
  is a covering sieve \(S \in J(U)\) all of whose members \(g : V \to U\) land on an object
  \(V\) where \(Q(V)\) is a zero object. Then the sheafification of \(Q\) is the zero
  sheaf.
\end{lemma}
\begin{proof}
  \leanok
  \uses{mk-lem:sheafify_kills_locally_zero}
  This is the ``locally zero'' packaging of
  Lemma~\ref{mk-lem:sheafify_kills_locally_zero}. Let \(Z\) be the constant presheaf with
  value the zero group; its sheafification is the zero sheaf. The zero map \(0 : Q \to Z\)
  is locally injective --- on each covering sieve furnished by the hypothesis the source
  sections \(Q(V)\) are subsingletons, so the comparison is injective locally --- and it
  is locally surjective because the target is objectwise a zero object. Hence
  \(0 : Q \to Z\) is locally bijective, i.e.\ lies in the local-isomorphism class
  \(J.W\), and part (3) followed by part (1) of
  Lemma~\ref{mk-lem:sheafify_kills_locally_zero} collapses the sheafification of \(Q\) to the
  zero sheaf.
\end{proof}

\begin{lemma}[The augmented {\v C}ech complex is a resolution]\label{mk-lem:cech_augmented_resolution}
  \lean{AlgebraicGeometry.cechAugmented_exact}
  \leanok
  \dcref{custom}

  \uses{mk-def:cech_nerve, mk-def:cech_augmented_complex, mk-lem:cech_free_eval_prepend_homotopy,
    mk-lem:cech_free_eval_prepend_homotopy_spec, mk-lem:cohomology_sheaf_is_sheafify_homology,
    mk-lem:sheafify_kills_locally_zero, mk-lem:isZero_of_faithful_preservesZeroMorphisms,
    mk-lem:isZero_presheafToSheaf_of_locally_isZero, mk-lem:cechSection_isZero_homology}
  \textit{Source: Stacks Project, Cohomology of Schemes,
  \texttt{lemma-cech-cohomology-quasi-coherent-trivial} (proof).}
  Let \(\mathcal{F}\) be a quasi-coherent \(\mathcal{O}_X\)-module and let
  \(\mathfrak{U} : X = \bigcup_{i \in I} U_i\) be a finite affine open cover with all
  intersections affine. Form the augmented {\v C}ech complex on \(X\),
  \[
    0 \;\to\; \mathcal{F} \;\to\; \mathcal{C}^0 \;\to\; \mathcal{C}^1 \;\to\;
    \mathcal{C}^2 \;\to\; \cdots,
    \qquad
    \mathcal{C}^p \;=\; \prod_{(i_0, \ldots, i_p) \in I^{p+1}}
      (j_{i_0 \ldots i_p})_*\bigl(\mathcal{F}|_{U_{i_0 \ldots i_p}}\bigr),
  \]
  the degree-\(p\) term being the degree-\(p\) object of the {\v C}ech nerve of
  Definition~\ref{mk-def:cech_nerve}, with augmentation the canonical map
  \(\mathcal{F} \to \mathcal{C}^0\). This augmented complex is exact; that is, the
  {\v C}ech nerve is a resolution of \(\mathcal{F}\) in \(\operatorname{QCoh}(X)\).
\end{lemma}
\begin{proof}
  \leanok
  \uses{mk-def:cech_augmented_complex, mk-lem:cech_free_eval_prepend_homotopy,
    mk-lem:cech_free_eval_prepend_homotopy_spec, mk-lem:cech_free_eval_sectionwise,
    mk-lem:cech_free_eval_engine_iso, mk-lem:cech_free_eval_nonempty,
    mk-lem:cohomology_sheaf_is_sheafify_homology, mk-lem:sheafify_kills_locally_zero,
    mk-lem:isZero_of_faithful_preservesZeroMorphisms,
    mk-lem:isZero_presheafToSheaf_of_locally_isZero, mk-lem:cechSection_isZero_homology}
    \textit{Source: Stacks Project, Cohomology of Schemes,
  \texttt{lemma-cech-cohomology-quasi-coherent-trivial} (proof).}
  We prove exactness of the augmented {\v C}ech complex of
  Definition~\ref{mk-def:cech_augmented_complex} in \(X.\mathrm{Modules}\) by showing that each homology
  object \(\mathcal{H}^p\) (\(p \geq 1\), together with the augmentation node) is the zero object. The
  argument works entirely with sections and sheafification, never with stalks.

  \emph{Step 1: reflect through the forgetful functor to abelian sheaves.} The reflection is carried
  out through the forgetful functor \(\operatorname{SheafOfModules.toSheaf} : X.\mathrm{Modules} \to
  \operatorname{Sh}(X, \mathbf{Ab})\) from sheaves of \(\mathcal{O}_X\)-modules to sheaves of abelian
  groups. This functor is faithful and additive, hence preserves zero morphisms; a faithful additive
  functor reflects the property of an object being zero. It therefore suffices to prove that the image
  \(\operatorname{SheafOfModules.toSheaf}(\mathcal{H}^p)\) of each homology object is zero in the
  category of abelian sheaves.

  \emph{Step 2: the homology sheaf is the sheafification of the presheaf homology.} By the
  sheafification-commutes-with-homology engine of
  Lemma~\ref{mk-lem:cohomology_sheaf_is_sheafify_homology}, the degree-\(p\) homology of the complex of
  abelian sheaves is the sheafification of the degree-\(p\) homology of the underlying complex of
  \emph{presheaves}. The presheaf homology evaluated at an open \(V\) is the homology of the section
  complex \(\check{\mathcal{C}}^\bullet(\mathcal{U}, \mathcal{F})(V)\) --- the {\v C}ech cochain complex
  of \(\mathcal{F}\) over the cover restricted to \(V\) --- so the presheaf in question is
  \(V \mapsto \check{H}^p(V, \mathcal{F})\).

  \emph{Step 3: the presheaf homology vanishes locally, via the prepend-\(i_{\mathrm{fix}}\) homotopy.}
  This presheaf is not zero as a presheaf, but it vanishes on a basis. The relevant basis is the family
  of opens \(V\) contained in some single cover element \(U_i\); since \(\{U_i\}\) covers \(X\), every
  open is covered by such \(V\) (intersect with the \(U_i\)), so this family is cofinal among all opens.
  Fix such a \(V\) with \(V \subseteq U_{i_{\mathrm{fix}}}\). The section complex
  \(\check{\mathcal{C}}^\bullet(\mathcal{U}, \mathcal{F})(V)\) is the {\v C}ech cochain complex of
  \(\mathcal{F}\) over the \emph{restricted} cover \(\{U_s \cap V\}_s\) of \(V\); this restricted cover
  has the member \(U_{i_{\mathrm{fix}}} \cap V = V\) equal to the whole of \(V\). A cover possessing a
  member equal to the total space carries a \emph{contracting homotopy}: the prepend-\(i_{\mathrm{fix}}\)
  map \((h\,s)_{i_0 \ldots i_p} = s_{i_{\mathrm{fix}}\, i_0 \ldots i_p}\) (with the evident
  restriction along \(U_{i_0 \ldots i_p} \cap V \subseteq U_{i_{\mathrm{fix}}\, i_0 \ldots i_p} \cap V\))
  satisfies \(d \circ h + h \circ d = \mathrm{id}\). This is the section-level objectwise analogue of the
  combinatorial homotopy \(\operatorname{combHomotopy}\) /
  \(\operatorname{combHomotopy\_spec}\) of Lemma~\ref{mk-lem:cech_free_eval_prepend_homotopy}; it is
  \emph{cover-agnostic and coefficient-agnostic} --- it uses neither quasi-coherence of \(\mathcal{F}\)
  nor affineness of the cover, only that \(V\) sits inside a cover member. Hence the section complex over
  every such \(V\) is contractible, so its homology vanishes in every degree. Consequently the canonical
  map \(0 \to (\text{presheaf homology})\) is a local isomorphism on this basis --- it is locally
  bijective. (This is a different mechanism from the affine {\v C}ech vanishing of
  Lemma~\ref{mk-lem:affine_cech_vanishing_tilde_subcover}, which concerns the \emph{standard} cover
  \(\{D(f_i)\}\) of an affine and feeds the 01XB Serre vanishing, not the present local check.)

  To convert this into a vanishing statement about the abelian-sheaf homology of Step 1, the
  module-level sheafification of the homology (the comparison isomorphism
  \(\operatorname{homologyIsoSheafify}\) underlying
  Lemma~\ref{mk-lem:cohomology_sheaf_is_sheafify_homology}) is transported to the abelian-sheaf side
  along the \emph{sheafification square}
  \[
    \operatorname{SheafOfModules.toSheaf} \circ \operatorname{sheafify}
      \;\cong\;
    \operatorname{presheafToSheaf} \circ \operatorname{forget},
  \]
  the natural isomorphism expressing that sheafifying a presheaf of modules and then forgetting to
  abelian sheaves agrees with forgetting to abelian presheaves and then sheafifying. Under this
  square the image \(\operatorname{SheafOfModules.toSheaf}(\mathcal{H}^p)\) becomes the abelian
  sheafification \(\operatorname{presheafToSheaf}\) of the abelian presheaf homology
  \(V \mapsto \check{H}^p(V, \mathcal{F})\). Lemma~\ref{mk-lem:sheafify_kills_locally_zero} now applies:
  the locally bijective comparison \(0 \to (\text{presheaf homology})\) is inverted by
  \(\operatorname{presheafToSheaf}\) (in the abelian setting it is simultaneously locally injective
  and locally surjective, hence in the local-isomorphism class), so its abelian sheafification is the
  zero sheaf. Combined with Steps 1--2 this gives \(\mathcal{H}^p = 0\) for all \(p \geq 1\).

  \emph{Step 3, proof mechanism (from the section homotopy to the vanishing of homology).} The prose
  above describes the contracting homotopy mathematically; we now record the precise mechanism by which it
  produces the required vanishing of the section-complex homology, since this is the step a formal proof
  must instantiate.
  \begin{enumerate}
    \item[(a)] After Steps 1--2 and the sheafification/site reduction, the residual obligation is exactly
      that, for every \(V \le U_{i_{\mathrm{fix}}}\) and every degree \(p\), the homology of the
      \emph{section cochain complex} \(\Gamma\bigl(V,\, \mathcal{C}^\bullet_{\mathrm{aug}}(\mathfrak{U},
      \mathcal{F})\bigr)\) --- viewed as a complex of abelian groups --- is a zero object. In the categorical
      packaging this is the homology of the complex obtained by applying the evaluation-at-\(V\) functor
      (through \(\operatorname{toPresheaf}\)) to the augmented {\v C}ech complex.
    \item[(b)] This abstract evaluated complex is identified with the \emph{concrete} {\v C}ech section
      cochain complex \(D^\bullet\) whose degree-\(p\) term is the product
      \(\prod_{\sigma} \Gamma(U_\sigma \cap V,\, \mathcal{F})\) with the alternating coface differential. The
      identification is the naturality of the map-of-homological-complexes functor together with a
      per-degree isomorphism of the product-of-sections term; it is the same
      categorical-to-combinatorial section-complex identification used for the free evaluated complex
      (Lemmas~\ref{mk-lem:cech_free_eval_sectionwise} and~\ref{mk-lem:cech_free_eval_engine_iso}). This per-degree
      identification is carried out in
      Lemma~\ref{mk-lem:cechSection_complex_iso}: its degreewise object isomorphism
      \(\Gamma(V, \operatorname{pushPullObj}\mathcal{F}\,Y_p) \cong \prod_\sigma \Gamma(U_\sigma \cap V,
      \mathcal{F})\) is Lemma~\ref{mk-lem:pushPull_eval_prod_iso} (which rests on the coproduct-to-product
      decomposition of the backbone push--pull object, Lemma~\ref{mk-lem:pushPull_sigma_iso}, the geometric
      backbone identification Lemma~\ref{mk-lem:cech_backbone_left_sigma}, and the per-leg section calculation
      Lemma~\ref{mk-lem:pushPull_leg_sections}), and its differential match reuses
      Lemma~\ref{mk-lem:section_cech_objd_apply}.
    \item[(c)] Because \(V \le U_{i_{\mathrm{fix}}}\), prepending the index \(i_{\mathrm{fix}}\) is the
      \emph{identity} on each section term --- \(U_{i_{\mathrm{fix}}\,\sigma} \cap V = U_\sigma \cap V\)
      since \(V \subseteq U_{i_{\mathrm{fix}}}\) --- so the prepend map gives a contracting homotopy from the
      identity of \(D^\bullet\) to the zero map, i.e.\ a homotopy \(\mathrm{id}_{D^\bullet} \simeq 0\), whose
      component maps are the combinatorial prepend maps \(\operatorname{combHomotopy} i_{\mathrm{fix}}\) and
      whose homotopy relation \(d \circ h + h \circ d = \mathrm{id}\) is exactly
      \(\operatorname{combHomotopy\_spec}\) (Lemma~\ref{mk-lem:cech_free_eval_prepend_homotopy} and its
      specification Lemma~\ref{mk-lem:cech_free_eval_prepend_homotopy_spec}). On the concrete section complex
      \(D'^\bullet\) this contracting homotopy is supplied by Lemma~\ref{mk-lem:cechSection_contractible}:
      since \(V \le U_{i_{\mathrm{fix}}}\) makes \(U_{i_{\mathrm{fix}}} \cap U_\sigma \cap V = U_\sigma \cap
      V\), prepending \(i_{\mathrm{fix}}\) is the identity on each coefficient and the dependent
      combinatorial engine yields \(\mathrm{id}_{D'^\bullet} \simeq 0\).
    \item[(d)] A homotopy \(\mathrm{id}_{D^\bullet} \simeq 0\) forces the homology to vanish in every degree:
      by homotopy-invariance the map induced on degree-\(p\) homology by the identity equals the map induced
      by the zero map; the former is the identity of \(D^\bullet\)'s degree-\(p\) homology and the latter is
      zero, so that identity equals zero, which is the criterion for the homology to be a zero object. (This
      composes the homotopy-invariance of the induced homology map with the identity/zero computations and
      the ``identity equals zero'' characterization of a zero object --- the same characterization used in
      Lemma~\ref{mk-lem:isZero_of_faithful_preservesZeroMorphisms}, packaged once and for all as
      Lemma~\ref{mk-lem:isZero_homology_of_homotopy_id_zero}.) Items (a)--(d) together are exactly the content
      of Lemma~\ref{mk-lem:cechSection_isZero_homology}, which is the single residual obligation discharging
      this step.
    \item[(e)] One route is recorded as a dead end: the homotopy must be built directly as a homotopy of
      homological complexes, following the construction of
      Lemma~\ref{mk-lem:cech_free_eval_nonempty}. It cannot be obtained from an extra-degeneracy package,
      because the only such package available is the \emph{simplicial} one producing a homotopy equivalence
      of \emph{chain} complexes (face maps), whereas the augmented {\v C}ech complex here is a \emph{cochain}
      complex (coface maps); there is no cosimplicial/cochain dual to invoke.
  \end{enumerate}

  \emph{Step 4: the augmentation node.} The augmentation \(\varepsilon : \mathcal{F} \to
  \mathcal{C}^0\) is handled by the \emph{same} prepend-\(i_{\mathrm{fix}}\) homotopy of Step 3. Over an
  open \(V \subseteq U_{i_{\mathrm{fix}}}\), the contracting homotopy extends to the augmentation term:
  the degree-\((-1)\)/degree-\(0\) component of \(h\) sends a section \(s \in \prod_s \mathcal{F}(U_s
  \cap V)\) to its \(i_{\mathrm{fix}}\)-component \(s_{i_{\mathrm{fix}}} \in \mathcal{F}(U_{i_{\mathrm{fix}}}
  \cap V) = \mathcal{F}(V)\), so the identity \(d \circ h + h \circ d = \mathrm{id}\) holds at the
  augmentation node as well. In the proof mechanism above this is the degree-\(0\) instance: the
  same per-degree identification (b) and the same homotopy (c) cover the augmentation term, the
  \(i_{\mathrm{fix}}\)-component map \(s \mapsto s_{i_{\mathrm{fix}}}\) landing in
  \(\mathcal{F}(U_{i_{\mathrm{fix}}} \cap V) = \mathcal{F}(V)\). Thus the augmented section complex over each
  such \(V\) is exact at degree
  \(0\) (and the augmentation is the kernel inclusion there). Sheafifying this local exactness as in
  Step 3 makes the augmentation-node homology zero. Together with Steps 1--3, every homology object of
  the augmented complex is zero, so the complex is exact and the {\v C}ech nerve of
  Definition~\ref{mk-def:cech_nerve} is a resolution of \(\mathcal{F}\) in \(X.\mathrm{Modules}\). In
  particular the statement holds for an arbitrary \(\mathcal{O}_X\)-module \(\mathcal{F}\) and an
  arbitrary open cover; quasi-coherence and affineness are needed only by the downstream consumers
  (termwise \(f_*\)-acyclicity), not by this exactness.

  \emph{Remark.} An alternative proof uses the stalk-at-prime criterion directly: the complex is
  exact if and only if it is exact after localizing at each prime \(\mathfrak{p}\), and for each
  \(\mathfrak{p}\) one fixes an index \(i_{\text{fix}}\) with \(f_{i_{\text{fix}}} \notin \mathfrak{p}\)
  and defines a contracting homotopy of the localized complex by
  \(h(s)_{i_0 \ldots i_p} = s_{i_{\text{fix}} i_0 \ldots i_p}\), as in the Stacks source proof above.
  This approach is equivalent to the sections-and-sheafification argument of Steps 1--4.
\end{proof}

\begin{lemma}[A contracting homotopy of the identity makes homology vanish]\label{mk-lem:isZero_homology_of_homotopy_id_zero}
  \lean{AlgebraicGeometry.isZero_homology_of_homotopy_id_zero}
  \leanok
  \dcref{custom}

  Let \(\mathcal{A}\) be a preadditive category and \(D^\bullet\) a cochain complex in
  \(\mathcal{A}\). If there is a homotopy \(\mathrm{id}_{D^\bullet} \simeq 0\) between the
  identity endomorphism of \(D^\bullet\) and the zero endomorphism, then for every degree
  \(p\) the homology object \(H^p(D^\bullet)\) is a zero object.
\end{lemma}
\begin{proof}
  \leanok

The map induced on homology is a homotopy invariant: homotopic chain maps induce equal
  maps on each homology object. Applying this to the homotopy \(\mathrm{id}_{D^\bullet}
  \simeq 0\), the map induced by \(\mathrm{id}_{D^\bullet}\) on \(H^p(D^\bullet)\) equals the
  map induced by \(0\). The former is \(\mathrm{id}_{H^p(D^\bullet)}\) (homology is a
  functor) and the latter is the zero endomorphism of \(H^p(D^\bullet)\); hence
  \(\mathrm{id}_{H^p(D^\bullet)} = 0\). An object whose identity morphism equals its zero
  endomorphism is a zero object, so \(H^p(D^\bullet)\) is zero.
\end{proof}

\begin{lemma}[Homology vanishes after transporting a null-homotopic complex across an isomorphism]\label{mk-lem:isZero_homology_of_iso_homotopy_id_zero}
  \lean{AlgebraicGeometry.isZero_homology_of_iso_homotopy_id_zero}
  \leanok
  \dcref{custom}

  \uses{mk-lem:isZero_homology_of_homotopy_id_zero}
  Let \(\mathcal{A}\) be a preadditive category with homology and let \(D^\bullet, D'^\bullet\) be
  cochain complexes in \(\mathcal{A}\). If \(D^\bullet \cong D'^\bullet\) and \(D'^\bullet\) admits
  a contracting homotopy \(\mathrm{id}_{D'^\bullet} \simeq 0\), then for every degree \(p\) the
  homology \(H^p(D^\bullet)\) is a zero object. This is the consumer glue that turns the
  contractibility of the model complex \(D'^\bullet\) into the vanishing of the homology of the
  isomorphic complex \(D^\bullet\) actually occurring in Step~3 of
  Lemma~\ref{mk-lem:cech_augmented_resolution}.
\end{lemma}
\begin{proof}
  \leanok
  \uses{mk-lem:isZero_homology_of_homotopy_id_zero}
  By Lemma~\ref{mk-lem:isZero_homology_of_homotopy_id_zero} the null-homotopy \(\mathrm{id}_{D'^\bullet}
  \simeq 0\) forces \(H^p(D'^\bullet)\) to be a zero object in every degree. Homology is a functor,
  so the isomorphism \(D^\bullet \cong D'^\bullet\) induces an isomorphism \(H^p(D^\bullet) \cong
  H^p(D'^\bullet)\); an object isomorphic to a zero object is itself a zero object, whence
  \(H^p(D^\bullet)\) is zero.
\end{proof}

\begin{lemma}[The degree-\(p\) {\v C}ech-nerve object is a wide pullback]\label{mk-lem:cechBackbone_obj_widePullback}
  \lean{AlgebraicGeometry.cechBackbone_obj_widePullback}
  \leanok
  \dcref{custom}

  \uses{mk-def:cover_cech_nerve}
  Let \(\mathcal{U} : X = \bigcup_{i \in I} U_i\) be an open cover with cover arrow
  \(q = \operatorname{Sigma.desc}\mathcal{U}.f : \coprod_{i \in I} U_i \to X\). For every
  \(p\), the underlying scheme of the degree-\(p\) object
  \((\operatorname{coverCechNerveOver}\mathcal{U}).\mathrm{obj}\,[p]\) of the {\v C}ech nerve of
  Definition~\ref{mk-def:cover_cech_nerve} is the wide pullback
  \[
    \operatorname{widePullback} X \;\bigl(\,\lambda\,k : \mathrm{Fin}(p+1).\ \textstyle\coprod_{i} U_i\,\bigr)
      \;\bigl(\,\lambda\,k.\ q\,\bigr),
  \]
  the \((p+1)\)-fold fibre power of \(q\) over \(X\), with structure map to \(X\) the canonical
  \(\operatorname{WidePullback.base}\). Equivalently, in \(\operatorname{Over} X\) the object is
  \(\operatorname{Over.mk}(\operatorname{WidePullback.base})\).
\end{lemma}
\begin{proof}
  \uses{mk-def:cover_cech_nerve}
  Pure unfolding of the definitions. By Definition~\ref{mk-def:cover_cech_nerve},
  \(\operatorname{coverCechNerveOver}\mathcal{U}\) is the lift into \(\operatorname{Over} X\) of
  \(\operatorname{coverCechNerve}\mathcal{U} = (\operatorname{coverArrow}\mathcal{U}).
  \operatorname{augmentedCechNerve}\) along its augmentation. The augmented {\v C}ech nerve of an
  arrow \(f\) has, in degree \(n\), the object \(\operatorname{Arrow.cechNerve} f\,.\mathrm{obj}\,
  n = \operatorname{widePullback} f.\mathrm{right}\,(\lambda\,k : \mathrm{Fin}(n+1).\ f.\mathrm{left})
  \,(\lambda\,k.\ f.\mathrm{hom})\). Here \(\operatorname{coverArrow}\mathcal{U} =
  \operatorname{Arrow.mk}(q)\) with \(\mathrm{right} = X\), \(\mathrm{left} = \coprod_i U_i\) and
  \(\mathrm{hom} = q\), so the degree-\(p\) object is exactly the displayed wide pullback, and the
  \(\operatorname{Over.lift}\) records its structure map to \(X\) as
  \(\operatorname{WidePullback.base}\).
\end{proof}

\begin{lemma}

[Wide fibre power over a base is the iterated product in the slice]
  \label{mk-lem:widePullback_overX_eq_prod}
  \lean{CategoryTheory.widePullback_overX_eq_prod,
    CategoryTheory.widePullback_overX_isLimit}
  \dcref{custom}

  \uses{mk-lem:widePullbackCone_isLimitOfFan_mathlib, mk-lem:over_mkIdTerminal_mathlib}
  Let \(\mathcal{C}\) be a category with pullbacks, \(S \in \mathcal{C}\), and
  \((Z_k \to S)_{k \in \mathrm{Fin}(n)}\) a finite family over \(S\). Then the wide pullback over
  \(S\),
  \[
    \operatorname{widePullback} S\;(\lambda\,k.\ Z_k)\;(\lambda\,k.\ Z_k \to S),
  \]
  is the \(n\)-fold product \(\prod_{k} Z_k\) of the same family taken in the slice category
  \(\operatorname{Over} S\) --- i.e.\ the \(n\)-fold fibre power over \(S\) is the \(n\)-fold product
  in \(\operatorname{Over} S\). In particular the fibre power has the clean recursion
  \(\prod_{k : \mathrm{Fin}(p+2)} Z_k \cong Z_0 \times_S \prod_{k : \mathrm{Fin}(p+1)} Z_{k+1}\)
  along \(\mathrm{Fin}(p+2) \cong \mathrm{Fin}(1) \oplus \mathrm{Fin}(p+1)\).
\end{lemma}
\begin{proof}
  \uses{mk-lem:widePullbackCone_isLimitOfFan_mathlib, mk-lem:over_mkIdTerminal_mathlib}
  In the slice \(\operatorname{Over} S\) the object \(\operatorname{Over.mk}(\mathrm{id}_S)\) is
  terminal by Lemma~\ref{mk-lem:over_mkIdTerminal_mathlib}, and the legs of the wide-pullback diagram
  all land in this terminal base. By Lemma~\ref{mk-lem:widePullbackCone_isLimitOfFan_mathlib} a
  limiting fan of the legs is then a limit of the wide-pullback diagram, so the wide pullback over
  \(S\) coincides with the product of the legs computed in \(\operatorname{Over} S\). The recursion
  is the standard splitting of a finite product along \(\mathrm{Fin}(p+2) \cong \mathrm{Fin}(1)
  \oplus \mathrm{Fin}(p+1)\) (a binary product with the head factor).
\end{proof}

\begin{lemma}[Base case of the wide-fibre-power decomposition]\label{mk-lem:coproduct_distrib_fibrePower_zero}
  \lean{CategoryTheory.FinitaryPreExtensive.widePullback_coproduct_iso_zero}
  \leanok
  \dcref{custom}

  \uses{mk-lem:widePullback_overX_eq_prod}
  Let \(\mathcal{C}\) be a category with pullbacks, \(\iota\) a finite index type, \(S \in
  \mathcal{C}\), and \((X_i \to S)_{i \in \iota}\) a finite family over \(S\). Then the \(1\)-fold
  wide fibre power of \(\coprod_i X_i \to S\) over \(S\) decomposes as
  \[
    \bigl(\textstyle\coprod_{i} X_i\bigr)^{\times_S 1}
      \;\cong\;
    \coprod_{\sigma : \mathrm{Fin}(1) \to \iota} X_{\sigma(0)},
  \]
  the \(\sigma\)-component being the \(1\)-fold fibre power \(X_{\sigma(0)}\).
\end{lemma}
\begin{proof}
  \uses{mk-lem:widePullback_overX_eq_prod}
  By Lemma~\ref{mk-lem:widePullback_overX_eq_prod} the \(1\)-fold fibre power is the \(1\)-fold product
  in \(\operatorname{Over} S\), which is the object itself: \((\coprod_i X_i)^{\times_S 1} \cong
  \coprod_i X_i\). On the right, the evaluation-at-\(0\) bijection \(\mathrm{Fin}(1) \to \iota\)
  \(\;\cong\;\) \(\iota\) identifies \(\coprod_{\sigma : \mathrm{Fin}(1)\to\iota} X_{\sigma(0)}\)
  with \(\coprod_i X_i\) (reindexing a coproduct along a bijection of index types). Composing the
  two gives the claim --- the trivial reindexing.
\end{proof}

\begin{lemma}

[One-sided binary distributivity in the slice]
  \label{mk-lem:prod_coproduct_distrib}
  \lean{CategoryTheory.FinitaryPreExtensive.prod_coproduct_distrib,
    CategoryTheory.FinitaryPreExtensive.coprodFirst_distrib}
  \dcref{custom}

  \uses{mk-lem:isIso_sigmaDesc_map_mathlib, mk-lem:isIso_sigmaDesc_fst_mathlib}
  Let \(\mathcal{C}\) be a finitary pre-extensive category with pullbacks, \(S \in \mathcal{C}\), and
  fix objects over \(S\). For a finite family \((Y_\tau \to S)_\tau\) and a single \(A \to S\) the
  fibre product distributes over the coproduct on each side:
  \[
    A \times_S \bigl(\textstyle\coprod_{\tau} Y_\tau\bigr)
      \;\cong\;
    \coprod_{\tau} \bigl(A \times_S Y_\tau\bigr),
    \qquad
    \bigl(\textstyle\coprod_{i} X_i\bigr) \times_S B
      \;\cong\;
    \coprod_{i} \bigl(X_i \times_S B\bigr),
  \]
  with each isomorphism the \(\operatorname{Sigma.desc}\) of the pullback maps, hence compatible with
  the projections to \(S\).
\end{lemma}
\begin{proof}
  \uses{mk-lem:isIso_sigmaDesc_map_mathlib, mk-lem:isIso_sigmaDesc_fst_mathlib}
  This is the binary distributivity instance Lemma~\ref{mk-lem:isIso_sigmaDesc_map_mathlib}
  (\(\coprod_{(i,j)}(X_i \times_S Y_j) \cong (\coprod_i X_i)\times_S(\coprod_j Y_j)\)) specialised by
  taking one of the two families to be a singleton: with the first family the singleton \(\{A\}\)
  one gets \(A \times_S \coprod_\tau Y_\tau \cong \coprod_\tau (A \times_S Y_\tau)\), and with the
  second family the singleton \(\{B\}\) one gets \((\coprod_i X_i)\times_S B \cong \coprod_i (X_i
  \times_S B)\).
\end{proof}

\begin{lemma}[One-sided distributivity in \(\operatorname{Over} S\) binary-product form]
  \label{mk-lem:overProd_coproduct_distrib}
  \lean{CategoryTheory.FinitaryPreExtensive.overProd_coproduct_distrib,
    CategoryTheory.FinitaryPreExtensive.pcd_hom_fst,
    CategoryTheory.FinitaryPreExtensive.pcd_hom_snd,
    CategoryTheory.FinitaryPreExtensive.cf_hom_fst,
    CategoryTheory.FinitaryPreExtensive.overSigma_hom_eq}
  \dcref{custom}

  \uses{mk-lem:prod_coproduct_distrib, mk-lem:overProdLeftIsoPullback_mathlib}
  Let \(\mathcal{C}\) be a finitary pre-extensive category with pullbacks and \(S \in \mathcal{C}\).
  For a finite index type \(\iota : \operatorname{Type} 0\), a family \((A_i)_{i \in \iota}\) of
  objects of \(\operatorname{Over} S\), and a single \(B \in \operatorname{Over} S\), the
  \emph{binary product in the slice} distributes over the coproduct:
  \[
    \bigl(\textstyle\coprod_{i} A_i\bigr) \times B
      \;\cong\;
    \coprod_{i} \bigl(A_i \times B\bigr)
      \qquad\text{in } \operatorname{Over} S,
  \]
  where \(\times\) is the categorical binary product of \(\operatorname{Over} S\). This is the form in
  which the inductive step of Lemma~\ref{mk-lem:coproduct_distrib_fibrePower} consumes distributivity:
  the recursion lives entirely in \(\operatorname{Over} S\), so it needs the slice binary product, not
  the \(\mathcal{C}\)-level fibre product of Lemma~\ref{mk-lem:prod_coproduct_distrib}.
\end{lemma}
\begin{proof}
  \uses{mk-lem:prod_coproduct_distrib, mk-lem:overProdLeftIsoPullback_mathlib}
  The forgetful functor \(\operatorname{Over} S \to \mathcal{C}\) reflects isomorphisms: a morphism of
  the slice is an isomorphism iff its underlying \(\mathcal{C}\)-morphism (its \(.\mathrm{left}\)) is.
  It therefore suffices to identify the underlying objects. The coproduct in \(\operatorname{Over} S\)
  is computed on underlying objects, \((\coprod_i A_i).\mathrm{left} = \coprod_i A_i.\mathrm{left}\),
  and the underlying object of a binary product in \(\operatorname{Over} S\) is the fibre product of
  the two structure maps, \((Y \times Z).\mathrm{left} \cong \operatorname{pullback} Y.\mathrm{hom}\,
  Z.\mathrm{hom}\), via \(\operatorname{Over.prodLeftIsoPullback}\)
  (Lemma~\ref{mk-lem:overProdLeftIsoPullback_mathlib}). Under these identifications both sides become
  the \(\mathcal{C}\)-level objects of the one-sided distributivity
  Lemma~\ref{mk-lem:prod_coproduct_distrib}: \((\coprod_i A_i.\mathrm{left}) \times_S B.\mathrm{left}\)
  on the left and \(\coprod_i (A_i.\mathrm{left} \times_S B.\mathrm{left})\) on the right. Matching
  \(.\mathrm{left}\) with that \(\mathcal{C}\)-level isomorphism, and noting all comparison maps commute
  with the projection to \(S\), gives the slice isomorphism.
\end{proof}

\begin{lemma}[One-sided distributivity in \(\operatorname{Over} S\), product-on-the-right form]\label{mk-lem:overProd_coproduct_distrib_right}
  \lean{CategoryTheory.FinitaryPreExtensive.overProd_coproduct_distrib_right}
  \leanok
  \dcref{custom}

  \uses{mk-lem:overProd_coproduct_distrib}
  Let \(\mathcal{C}\) be a finitary pre-extensive category with pullbacks and \(S \in \mathcal{C}\).
  For a finite index type \(\iota : \operatorname{Type} 0\), a single \(A \in \operatorname{Over} S\)
  and a family \((Y_i)_{i \in \iota}\) of objects of \(\operatorname{Over} S\), the slice binary
  product distributes over the coproduct in the second argument:
  \[
    A \times \bigl(\textstyle\coprod_{i} Y_i\bigr)
      \;\cong\;
    \coprod_{i} \bigl(A \times Y_i\bigr)
      \qquad\text{in } \operatorname{Over} S.
  \]
  This is the braiding-twin of Lemma~\ref{mk-lem:overProd_coproduct_distrib}; the inductive step of
  Lemma~\ref{mk-lem:coproduct_distrib_fibrePower} uses it to distribute the head factor over the
  coproduct of multi-indexed fibre powers.
\end{lemma}
\begin{proof}
  \uses{mk-lem:overProd_coproduct_distrib}
  Braid the slice binary product, \(A \times (\coprod_i Y_i) \cong (\coprod_i Y_i) \times A\), apply
  the left-handed distributivity (Lemma~\ref{mk-lem:overProd_coproduct_distrib}) to get
  \(\coprod_i (Y_i \times A)\), and braid each summand back to \(A \times Y_i\). All three steps are
  isomorphisms of \(\operatorname{Over} S\), so the composite is the asserted distributivity.
\end{proof}

\begin{lemma}[The underlying object of a slice binary product is the pullback]
  \label{mk-lem:overProdLeftIsoPullback_mathlib}
  \lean{CategoryTheory.Over.prodLeftIsoPullback}
  \mathlibok
  \dcref{custom}

  For \(S \in \mathcal{C}\) and two objects \(Y, Z \in \operatorname{Over} S\) (with \(\mathcal{C}\)
  having the relevant pullback), the underlying object of their binary product in the slice is the
  fibre product of their structure maps: \((Y \times Z).\mathrm{left} \cong \operatorname{pullback}
  Y.\mathrm{hom}\, Z.\mathrm{hom}\), compatibly with the projections.
\end{lemma}

\begin{lemma}[Nested coproduct flattens and reindexes to the next fibre power]\label{mk-lem:coproduct_fibrePower_reindex}
  \lean{CategoryTheory.FinitaryPreExtensive.coproduct_fibrePower_reindex}
  \leanok
  \dcref{custom}

  \uses{mk-lem:sigmaSigmaIso_mathlib}
  Let \(\iota\) be a finite index type, \(S \in \mathcal{C}\), \((X_i \to S)_i\) a family over
  \(S\), and for \(\tau : \mathrm{Fin}(p+1)\to\iota\) write \(Y_\tau = X_{\tau(0)} \times_S \cdots
  \times_S X_{\tau(p)}\). Then the nested coproduct flattens and reindexes to a coproduct over
  \((p+2)\)-fold multi-indices:
  \[
    \coprod_{i}\coprod_{\tau} \bigl(X_i \times_S Y_\tau\bigr)
      \;\cong\;
    \coprod_{(i,\tau)} \bigl(X_i \times_S Y_\tau\bigr)
      \;\cong\;
    \coprod_{\sigma : \mathrm{Fin}(p+2) \to \iota}
      \bigl(X_{\sigma(0)} \times_S \cdots \times_S X_{\sigma(p+1)}\bigr),
  \]
  where the second isomorphism is the reindexing \((i, \tau) \mapsto \sigma =
  \operatorname{Fin.cons} i\,\tau\), a bijection \(\iota \times (\mathrm{Fin}(p+1)\to\iota) \cong
  (\mathrm{Fin}(p+2)\to\iota)\), under which \(X_i \times_S Y_\tau\) is exactly the \((p+2)\)-fold
  fibre power \(X_{\sigma(0)} \times_S \cdots \times_S X_{\sigma(p+1)}\).
\end{lemma}
\begin{proof}
  \uses{mk-lem:sigmaSigmaIso_mathlib}
  The first isomorphism is \(\operatorname{sigmaSigmaIso}\) of
  Lemma~\ref{mk-lem:sigmaSigmaIso_mathlib}, collapsing the doubly-indexed coproduct
  \(\coprod_i \coprod_\tau\) into a coproduct over pairs \((i,\tau)\). The
  \(\operatorname{Fin.cons}\) bijection \(\iota \times (\mathrm{Fin}(p+1)\to\iota) \cong
  (\mathrm{Fin}(p+2)\to\iota)\), \((i,\tau)\mapsto \operatorname{Fin.cons} i\,\tau\), reindexes that
  coproduct over \(\sigma\); and since \(\sigma(0) = i\) and \(\sigma(k+1) = \tau(k)\), the component
  \(X_i \times_S Y_\tau = X_i \times_S (X_{\tau(0)} \times_S \cdots)\) is term-for-term the
  \((p+2)\)-fold fibre power \(X_{\sigma(0)} \times_S \cdots \times_S X_{\sigma(p+1)}\).
\end{proof}

\begin{lemma}
[Coproduct distributes over the wide fibre power]
  \label{mk-lem:coproduct_distrib_fibrePower}
  \lean{CategoryTheory.FinitaryPreExtensive.widePullback_coproduct_iso,
    CategoryTheory.FinitaryPreExtensive.prodFinSuccIso}
  \dcref{custom}

  \uses{mk-lem:widePullback_overX_eq_prod, mk-lem:coproduct_distrib_fibrePower_zero,
    mk-lem:prod_coproduct_distrib, mk-lem:coproduct_fibrePower_reindex,
    mk-lem:overProd_coproduct_distrib, mk-lem:overProd_coproduct_distrib_right,
    mk-lem:finitaryExtensive_scheme_mathlib}
  Let \(\mathcal{C}\) be a category with pullbacks that is finitary pre-extensive
  (\(\operatorname{FinitaryPreExtensive}\mathcal{C}\)), let \(\iota\) be a finite index type,
  \(S \in \mathcal{C}\), and \((X_i \to S)_{i \in \iota}\) a finite family over \(S\). Then for every
  \(p\) the \((p+1)\)-fold wide fibre power of the coproduct map \(\coprod_i X_i \to S\) over \(S\)
  decomposes as a coproduct over multi-indices:
  \[
    \bigl(\textstyle\coprod_{i} X_i\bigr)^{\times_S (p+1)}
      \;\cong\;
    \coprod_{\sigma : \mathrm{Fin}(p+1) \to \iota}
      \bigl(X_{\sigma(0)} \times_S \cdots \times_S X_{\sigma(p)}\bigr),
  \]
  the \(\sigma\)-component being the \((p+1)\)-fold fibre power of the legs \(X_{\sigma(k)} \to S\)
  over \(S\), and the structure map of the coproduct being
  \(\operatorname{Sigma.desc}\) of the component structure maps.

  \emph{Slice-product normal form.} The proof carries out the recursion in
  \(\operatorname{Over} S\), so the \(\sigma\)-component is presented as the iterated \emph{slice
  product} \(\prod^{\mathrm{c}}_{k}\operatorname{Over.mk}\bigl(f(\sigma k)\bigr)\) over
  \(\operatorname{Over} S\) rather than as the wide fibre power directly; the two are identified by
  Lemma~\ref{mk-lem:widePullback_overX_eq_prod} (wide fibre power over \(S\) equals the iterated product in
  \(\operatorname{Over} S\)). The result is stated
  abstractly for \(\operatorname{FinitaryPreExtensive}\mathcal{C}\) (reusable)
  and then instantiated at \(\mathcal{C} = \operatorname{Scheme}\), which is finitary extensive by
  Lemma~\ref{mk-lem:finitaryExtensive_scheme_mathlib}.

  \emph{Universe constraint.} This statement requires the index type \(\iota : \operatorname{Type} 0\):
  the binary-distributivity primitive it rests on (Lemma~\ref{mk-lem:isIso_sigmaDesc_fst_mathlib}) does not
  hold for index types in higher universes. When the cover index is finite but lives in a higher
  universe, one first reindexes along an equivalence to \(\mathrm{Fin}(n)\) before applying this
  result; see Lemma~\ref{mk-lem:cech_backbone_left_sigma}.
\end{lemma}
\begin{proof}
  \uses{mk-lem:widePullback_overX_eq_prod, mk-lem:coproduct_distrib_fibrePower_zero,
    mk-lem:prod_coproduct_distrib, mk-lem:coproduct_fibrePower_reindex, mk-lem:overProd_coproduct_distrib,
    mk-lem:overProd_coproduct_distrib_right, mk-lem:finitaryExtensive_scheme_mathlib}
  Induction on \(p\). Throughout, Lemma~\ref{mk-lem:widePullback_overX_eq_prod} lets us treat the
  \((p+1)\)-fold fibre power over \(S\) as the \((p+1)\)-fold product in \(\operatorname{Over} S\),
  so the induction has the clean recursion \(P^{\times_S (p+2)} \cong P \times_S P^{\times_S (p+1)}\)
  with \(P = \coprod_i X_i\).

  \emph{Base step} \(p = 0\) is Lemma~\ref{mk-lem:coproduct_distrib_fibrePower_zero}: the \(1\)-fold
  fibre power is \(\coprod_i X_i\) itself, indexed by \(\sigma : \mathrm{Fin}(1) \to \iota\).

  \emph{Inductive step}: assume \(\bigl(\coprod_i X_i\bigr)^{\times_S (p+1)} \cong \coprod_{\tau :
  \mathrm{Fin}(p+1) \to \iota} Y_\tau\) with \(Y_\tau = X_{\tau(0)} \times_S \cdots \times_S
  X_{\tau(p)}\). By the recursion of Lemma~\ref{mk-lem:widePullback_overX_eq_prod} and the inductive
  hypothesis,
  \[
    \bigl(\textstyle\coprod_i X_i\bigr)^{\times_S (p+2)}
      \;\cong\;
    \bigl(\textstyle\coprod_i X_i\bigr) \times_S \bigl(\textstyle\coprod_\tau Y_\tau\bigr).
  \]
  Applying the one-sided distributivity of
  Lemma~\ref{mk-lem:prod_coproduct_distrib} on each side turns the right-hand side into
  \(\coprod_{i}\coprod_{\tau} \bigl(X_i \times_S Y_\tau\bigr)\). Since the recursion lives in
  \(\operatorname{Over} S\), this distributivity is consumed in its \(\operatorname{Over} S\)
  binary-product form (Lemma~\ref{mk-lem:overProd_coproduct_distrib}), which expresses
  \((\coprod_\tau Y_\tau) \times B \cong \coprod_\tau (Y_\tau \times B)\) for the binary product in the
  slice rather than the \(\mathcal{C}\)-level pullback. Finally
  Lemma~\ref{mk-lem:coproduct_fibrePower_reindex} flattens this nested coproduct
  (via \(\operatorname{sigmaSigmaIso}\)) and reindexes the pairs \((i,\tau)\) as
  \(\sigma = \operatorname{Fin.cons} i\,\tau : \mathrm{Fin}(p+2) \to \iota\), identifying \(X_i
  \times_S Y_\tau\) with \(X_{\sigma(0)} \times_S \cdots \times_S X_{\sigma(p+1)}\). The structure
  maps to \(S\) agree because each constituent isomorphism is built from \(\operatorname{Sigma.desc}\)
  / product comparisons that commute with the projection to \(S\) by construction. Instantiating at
  \(\mathcal{C} = \operatorname{Scheme}\) (Lemma~\ref{mk-lem:finitaryExtensive_scheme_mathlib}) gives the
  scheme-level statement.
\end{proof}

\begin{lemma}

[The wide fibre power of open immersions is the intersection open]
  \label{mk-lem:widePullback_openImm_inter}
  \lean{AlgebraicGeometry.widePullback_openImm_inter,
    AlgebraicGeometry.mem_iInf_opens_of_finite}
  \dcref{custom}

  \uses{mk-def:cech_free_presheaf_complex, mk-lem:isPullback_opens_inf_mathlib}
  Let \(\mathcal{U} : X = \bigcup_{i \in I} U_i\) be an open cover and
  \(\sigma : \mathrm{Fin}(p+1) \to I\) a multi-index. The \((p+1)\)-fold wide fibre power over
  \(X\) of the open immersions \((U_{\sigma(k)} \hookrightarrow X)_{k}\) is the intersection open:
  \[
    \operatorname{widePullback} X \;(\lambda\,k.\ U_{\sigma(k)})\;(\lambda\,k.\ (U_{\sigma(k)}).\iota)
      \;\cong\;
    (\operatorname{coverInterOpen}\mathcal{U}\,\sigma).\mathrm{toScheme}
      \;=\; \bigl(\textstyle\bigcap_k U_{\sigma(k)}\bigr).\mathrm{toScheme},
  \]
  over \(X\), and the structure map is the open immersion
  \(j_\sigma : \operatorname{coverInterOpen}\mathcal{U}\,\sigma \hookrightarrow X\) of
  Definition~\ref{mk-def:cech_free_presheaf_complex}.
\end{lemma}
\begin{proof}
  \uses{mk-lem:isPullback_opens_inf_mathlib, mk-def:cech_free_presheaf_complex}
  Induction on the family, iterating the binary case. The pullback over \(X\) of two open
  immersions \(U.\iota, V.\iota\) is the intersection open \((U \sqcap V).\mathrm{toScheme}\) with
  the two \(\operatorname{homOfLE}\) inclusions --- this is
  \(\operatorname{isPullback\_opens\_inf}\) of Lemma~\ref{mk-lem:isPullback_opens_inf_mathlib}. The
  wide pullback of the family \((U_{\sigma(k)})_{k}\) is the iterated binary pullback over \(X\)
  along \(\mathrm{Fin}(p+1)\), so iterating the binary intersection gives
  \(\bigcap_k U_{\sigma(k)}\), which is the indexed infimum
  \(\operatorname{coverInterOpen}\mathcal{U}\,\sigma = \bigcap_k \operatorname{coverOpen}
  \mathcal{U}\,(\sigma(k))\) of Definition~\ref{mk-def:cech_free_presheaf_complex} (recall
  \(\operatorname{coverOpen}\mathcal{U}\,i = (\mathcal{U}.f\,i).\operatorname{opensRange}\) is the
  image open \(U_i\)). The structure map to \(X\) is the inclusion of the intersection, namely
  \(j_\sigma\).
\end{proof}

\begin{lemma}

[The degree-\(p\) {\v C}ech backbone is the coproduct of intersection opens]
  \label{mk-lem:cech_backbone_left_sigma}
  \lean{AlgebraicGeometry.cechBackbone_left_sigma,
    AlgebraicGeometry.coverInterProdIso,
    AlgebraicGeometry.widePullbackBaseCongr}
  \dcref{custom}

  \uses{mk-lem:cechBackbone_obj_widePullback, mk-lem:coproduct_distrib_fibrePower,
    mk-lem:widePullback_overX_eq_prod, mk-lem:widePullback_openImm_inter, mk-def:cover_cech_nerve,
    mk-def:cech_free_presheaf_complex}
  Let \(\mathcal{U} : X = \bigcup_{i \in I} U_i\) be an open cover with cover arrow
  \(q : \coprod_{i \in I} U_i \to X\), and let \(Y_p\) denote the degree-\(p\) object of the
  {\v C}ech nerve \(\operatorname{coverCechNerveOver}\mathcal{U}\) of
  Definition~\ref{mk-def:cover_cech_nerve} --- the \((p+1)\)-fold fibre power of \(q\) over
  \(X\). Then \(Y_p\) is canonically the coproduct
  \[
    Y_p \;\cong\; \coprod_{\sigma : \mathrm{Fin}(p+1) \to I} U_\sigma,
    \qquad
    U_\sigma \;=\; \operatorname{coverInterOpen}\mathcal{U}\,\sigma \;=\;
      \bigcap_{k} U_{\sigma(k)},
  \]
  indexed by the multi-indices \(\sigma\), and the structure map \(q_p : Y_p \to X\) is
  \(\operatorname{Sigma.desc}(\sigma \mapsto j_\sigma)\), where
  \(j_\sigma : U_\sigma \hookrightarrow X\) is the open immersion of the intersection open
  \(U_\sigma\) of Definition~\ref{mk-def:cech_free_presheaf_complex}. In words: the fibre power
  of the open immersions \(U_i \hookrightarrow X\) is the intersection open, and the disjoint
  union over \(\sigma\) assembles the {\v C}ech backbone.
\end{lemma}
\begin{proof}
  \uses{mk-lem:cechBackbone_obj_widePullback, mk-lem:coproduct_distrib_fibrePower,
    mk-lem:widePullback_overX_eq_prod, mk-lem:widePullback_openImm_inter}
  Assemble the chain in \(\operatorname{Over} X\). By
  Lemma~\ref{mk-lem:cechBackbone_obj_widePullback} the degree-\(p\) object \(Y_p\) is the wide
  pullback \(\bigl(\coprod_i U_i\bigr)^{\times_X (p+1)}\), the \((p+1)\)-fold fibre power of
  \(q\) over \(X\). By Lemma~\ref{mk-lem:coproduct_distrib_fibrePower} (instantiated at
  \(\mathcal{C} = \operatorname{Scheme}\), with \(S = X\) and the family
  \(U_i \hookrightarrow X\)) this fibre power is the coproduct
  \(\coprod_{\sigma : \mathrm{Fin}(p+1) \to I} \bigl(U_{\sigma(0)} \times_X \cdots \times_X
  U_{\sigma(p)}\bigr)\) over the multi-indices. This decomposition delivers each \(\sigma\)-component in
  its slice-product normal form, namely the iterated product
  \(\prod^{\mathrm{c}}_{k}\operatorname{Over.mk}(U_{\sigma(k)} \hookrightarrow X)\) in
  \(\operatorname{Over} X\); before it can be identified as an intersection open it must be converted
  back to the wide-pullback form, which is exactly Lemma~\ref{mk-lem:widePullback_overX_eq_prod} (the
  iterated product over \(X\) in \(\operatorname{Over} X\) equals the wide fibre power over \(X\)). With
  the component thus in wide-pullback form, Lemma~\ref{mk-lem:widePullback_openImm_inter} applies: each
  \(\sigma\)-component, being the wide fibre power over \(X\) of the open immersions
  \(U_{\sigma(k)} \hookrightarrow X\), is the intersection open
  \(U_\sigma = \operatorname{coverInterOpen}\mathcal{U}\,\sigma = \bigcap_k U_{\sigma(k)}\), with
  structure map the open immersion \(j_\sigma\). The per-\(\sigma\) component identification
  packaging this slice-product\(\to\)intersection-open isomorphism is
  \(\operatorname{coverInterProdIso}\), and the transport of a wide fibre power along an
  equality of its base object (used to align the two wide-pullback forms over \(X\)) is
  \(\operatorname{widePullbackBaseCongr}\). Composing these
  isomorphisms gives \(Y_p \cong \coprod_\sigma U_\sigma\); the structure-map agreement is
  automatic, since each step's isomorphism commutes with the projection to \(X\), so the
  universal map out of the coproduct is \(\operatorname{Sigma.desc}(\sigma \mapsto j_\sigma)\).

  \emph{Universe reduction.} The distributivity Lemma~\ref{mk-lem:coproduct_distrib_fibrePower} and the
  binary-distributivity leaf it rests on (Lemma~\ref{mk-lem:isIso_sigmaDesc_fst_mathlib}) require the
  index type to lie in \(\operatorname{Type} 0\); the cover index \(I\) of \(\mathcal{U}\) is a
  finite type that need not live in \(\operatorname{Type} 0\). The reduction is therefore performed
  here: from finiteness of \(I\) choose an equivalence \(e : I \cong \mathrm{Fin}(n)\) (Mathlib's
  \(\operatorname{Finite.equivFin}\), or \(\operatorname{Fintype.equivFin}\) once a \(\operatorname{Fintype}\)
  instance is fixed), and reindex the cover family \((U_i \hookrightarrow X)\) and every coproduct- and
  wide-pullback-index along it to land at \(\mathrm{Fin}(n) : \operatorname{Type} 0\). Two index
  transports are involved. On the left, the source coproduct is reindexed along \(e\),
  \[
    \coprod_{i \in I} \mathcal{U}.X
      \;\cong\;
    \coprod_{j \in \mathrm{Fin}(n)} \bigl(\mathcal{U}.X \circ e^{-1}\bigr),
  \]
  the coproduct comparison induced by the index equivalence (\(\operatorname{Sigma.whiskerEquiv}\)).
  On the right, the multi-index type is reindexed on its codomain,
  \[
    \bigl(\mathrm{Fin}(p+1) \to \mathrm{Fin}(n)\bigr)
      \;\cong\;
    \bigl(\mathrm{Fin}(p+1) \to I\bigr),
  \]
  the arrow-congruence \(\operatorname{Equiv.arrowCongr}\) applied to \(e^{-1}\) in the codomain, under
  which the coproduct over \(\sigma\)-components of the \(\mathrm{Fin}(n)\)-decomposition matches the
  coproduct over the original multi-indices. Apply the \(\operatorname{Type} 0\) decomposition of
  Lemma~\ref{mk-lem:coproduct_distrib_fibrePower} at \(\mathrm{Fin}(n)\), then transport the resulting
  isomorphism back along these two equivalences. Reindexing a coproduct (resp.\ a wide fibre power)
  along a bijection of index types is an isomorphism compatible with the structure maps, so the
  transported isomorphism is again \(\operatorname{Sigma.desc}(\sigma \mapsto j_\sigma)\) over \(X\).
\end{proof}

\begin{lemma}

[The cover arrow is the coproduct of its member arrows in the slice]
  \label{mk-lem:coverArrow_over_sigma}
  \lean{AlgebraicGeometry.coverArrowOverSigmaIso,
    AlgebraicGeometry.coverArrowOverCofan,
    AlgebraicGeometry.coverArrowOverIsColimit,
    CategoryTheory.overSigmaDescCofan,
    CategoryTheory.overSigmaDescIsColimit,
    CategoryTheory.overSigmaDescIso}
  \dcref{custom}

  \uses{mk-def:cover_cech_nerve}
  Let \(\mathcal{U} : X = \bigcup_{i \in I} U_i\) be an open cover with cover arrow
  \(q = \operatorname{Sigma.desc}\mathcal{U}.f : \coprod_{i \in I} U_i \to X\). Then in the slice
  category \(\operatorname{Over} X\) the object \(\operatorname{Over.mk}(q)\) carrying the cover arrow is
  the coproduct of the member arrows \(\operatorname{Over.mk}(\mathcal{U}.f\,i)\):
  \[
    \operatorname{Over.mk}\bigl(\operatorname{Sigma.desc}\mathcal{U}.f\bigr)
      \;\cong\;
    \coprod_{i \in I} \operatorname{Over.mk}(\mathcal{U}.f\,i)
      \qquad\text{in } \operatorname{Over} X.
  \]
  This is the degree-\(0\) instance of the {\v C}ech backbone identification
  (Lemma~\ref{mk-lem:cech_backbone_left_sigma}): the source of the cover arrow is the disjoint union of the
  cover members, structurally over \(X\).
\end{lemma}
\begin{proof}
  \uses{mk-def:cover_cech_nerve}
  The object \(\operatorname{Over.mk}(q)\) is the coproduct of the family
  \((\operatorname{Over.mk}(\mathcal{U}.f\,i))_i\) in \(\operatorname{Over} X\). The \(i\)-th
  coproduct inclusion is the slice morphism lifting the inclusion
  \(\iota_i : U_i \hookrightarrow \coprod_{j \in I} U_j\); it is well-defined as a morphism over
  \(X\) because \(q \circ \iota_i = \mathcal{U}.f\,i\). The universal property holds: any family
  of slice morphisms \(\operatorname{Over.mk}(\mathcal{U}.f\,i) \to Z\) into a common slice object
  \(Z\) factors uniquely through \(\operatorname{Over.mk}(q)\) via the universal map of the
  underlying coproduct \(\coprod_{j \in I} U_j\), lifted to the slice since the forgetful functor
  \(\operatorname{Over} X \to \operatorname{Scheme}\) reflects the required equalities.
  Comparing with the standard coproduct in \(\operatorname{Over} X\) yields the asserted isomorphism.
\end{proof}

\begin{lemma}[Reflect an isomorphism of module sheaves to the underlying presheaf of abelian groups]\label{mk-lem:isIso_modules_of_toPresheaf}
  \lean{AlgebraicGeometry.isIso_modules_of_toPresheaf}
  \leanok
  \dcref{custom}

  \uses{mk-lem:forget_reflectsIso_mathlib}
  Let \(\varphi : \mathcal{M} \to \mathcal{N}\) be a morphism in \(X.\mathrm{Modules}\). If the
  image of \(\varphi\) under the forgetful functor
  \(\operatorname{Scheme.Modules.toPresheaf} X : X.\mathrm{Modules} \to \mathrm{Presheaf}(\mathrm{Ab})\)
  to the underlying presheaf of abelian groups is an isomorphism, then \(\varphi\) is an
  isomorphism in \(X.\mathrm{Modules}\).
\end{lemma}
\begin{proof}
  \uses{mk-lem:forget_reflectsIso_mathlib}
  The forgetful functor from sheaves of \(\mathcal{O}_X\)-modules to the underlying (pre)sheaf of
  abelian groups is fully faithful, hence reflects isomorphisms
  (Lemma~\ref{mk-lem:forget_reflectsIso_mathlib}); the inclusion of sheaves of abelian groups into
  presheaves of abelian groups is likewise fully faithful, so the composite to presheaves of abelian
  groups reflects isomorphisms as well. A functor reflecting isomorphisms turns the hypothesis that
  the image of \(\varphi\) is an isomorphism into the conclusion that \(\varphi\) itself is one.
\end{proof}

\begin{lemma}
[Restrict--stalk comparison along an open immersion]
  \label{mk-lem:restrictStalkNatIso_mathlib}
  \lean{AlgebraicGeometry.Scheme.Modules.restrictStalkNatIso}
  \mathlibok
  \dcref{custom}

  Let \(f : U \hookrightarrow X\) be an open immersion of schemes and \(y \in U\) a point with
  image \(x = f(y)\). For a module sheaf \(\mathcal{M} \in X.\mathrm{Modules}\), restricting
  \(\mathcal{M}\) along \(f\) and then taking the stalk of the underlying presheaf of abelian
  groups at \(y\) is naturally isomorphic to the stalk of \(\mathcal{M}\) (as a presheaf of
  abelian groups) at \(x\); i.e.\ there is a natural isomorphism between the functors
  \(\mathcal{M} \mapsto (\operatorname{restrict} f\,\mathcal{M})_y\) and \(\mathcal{M} \mapsto
  \mathcal{M}_x\).
\end{lemma}

\begin{lemma}
[Stalkwise conservativity for sheaves of abelian groups]
  \label{mk-lem:isIso_iff_stalkFunctor_map_iso_mathlib}
  \lean{TopCat.Presheaf.isIso_iff_stalkFunctor_map_iso}
  \mathlibok
  \dcref{custom}

  A morphism \(\varphi\) of sheaves of abelian groups on a topological space \(X\) is an
  isomorphism if and only if for every point \(x \in X\) the induced map on stalks
  \(\varphi_x\) is an isomorphism. (Stalk functors are jointly conservative on
  \(\mathrm{Sheaf}(\mathrm{Ab}, X)\).)
\end{lemma}

\begin{lemma}
[Underlying presheaf of a module sheaf is a sheaf]
  \label{mk-lem:modules_isSheaf_mathlib}
  \lean{AlgebraicGeometry.Scheme.Modules.isSheaf}
  \mathlibok
  \dcref{custom}

  For \(\mathcal{M} \in X.\mathrm{Modules}\), the underlying presheaf of abelian groups
  \(\mathcal{M}.\mathrm{presheaf}\) satisfies the sheaf condition, so it lifts to an object of
  \(\mathrm{Sheaf}(\mathrm{Ab}, X)\).
\end{lemma}

\begin{lemma}[A module-sheaf morphism is an isomorphism iff it is one on every member of an open cover]\label{mk-lem:isIso_iff_isIso_restrict}
  \lean{AlgebraicGeometry.Scheme.Modules.Hom.isIso_iff_isIso_restrict}
  \leanok
  \dcref{custom}

  \uses{mk-lem:isIso_modules_of_toPresheaf, mk-lem:restrictStalkNatIso_mathlib,
    mk-lem:isIso_iff_stalkFunctor_map_iso_mathlib, mk-lem:modules_isSheaf_mathlib}
  Let \(X\) be a scheme, \(\varphi : \mathcal{M} \to \mathcal{N}\) a morphism in
  \(X.\mathrm{Modules}\), and \(\mathfrak{U}\) an open cover of \(X\). Then \(\varphi\) is an
  isomorphism if and only if, for every index \(j\), its restriction along the open immersion
  \(\mathfrak{U}.\mathrm{map}\,j\) to the cover member \(\mathfrak{U}.\mathrm{obj}\,j\) is
  an isomorphism.
\end{lemma}
\begin{proof}
  \uses{mk-lem:isIso_modules_of_toPresheaf, mk-lem:restrictStalkNatIso_mathlib,
    mk-lem:isIso_iff_stalkFunctor_map_iso_mathlib, mk-lem:modules_isSheaf_mathlib}
  \((\Rightarrow)\) Functors preserve isomorphisms: if \(\varphi\) is an isomorphism then each
  restriction functor \(\operatorname{restrictFunctor}(\mathfrak{U}.\mathrm{map}\,j)\) sends it to
  an isomorphism.

  \((\Leftarrow)\) Assume every restriction is an isomorphism. It suffices to show that the image
  of \(\varphi\) under the forgetful functor \(\operatorname{Scheme.Modules.toPresheaf} X\) to the
  underlying presheaf of abelian groups is an isomorphism, because that functor reflects
  isomorphisms (\ref{mk-lem:isIso_modules_of_toPresheaf}). The underlying presheaves are
  sheaves of abelian groups by \ref{mk-lem:modules_isSheaf_mathlib}; by stalk
  conservativity (\ref{mk-lem:isIso_iff_stalkFunctor_map_iso_mathlib}), it suffices to
  show that \(\varphi\) is an isomorphism on every stalk.

  Fix \(x \in X\) and choose an index \(j\) such that \(x\) lies in the image of the cover map
  \(\mathfrak{U}.\mathrm{map}\,j\) (the covering property); choose \(y\) in
  \(\mathfrak{U}.\mathrm{obj}\,j\) mapping to \(x\). The restrict-versus-stalk comparison
  \(\operatorname{restrictStalkNatIso}(\mathfrak{U}.\mathrm{map}\,j)\,y\)
  (\ref{mk-lem:restrictStalkNatIso_mathlib}) --- a natural isomorphism between ``restrict along the
  open immersion \(\mathfrak{U}.\mathrm{map}\,j\) then take the stalk at \(y\)'' and ``take the
  stalk at \(x = (\mathfrak{U}.\mathrm{map}\,j)\,y\)'' --- conjugates the stalk at \(y\) of the
  restricted morphism \((\operatorname{restrictFunctor}(\mathfrak{U}.\mathrm{map}\,j)).\mathrm{map}\,\varphi\)
  with the stalk of \(\varphi\) at \(x\). Since the restricted morphism is an
  isomorphism by hypothesis, its stalk at \(y\) is an isomorphism, hence so is the stalk of
  \(\varphi\) at \(x\). As \(x\) was arbitrary, \(\varphi\) is a stalkwise isomorphism, hence an
  isomorphism.
\end{proof}

\begin{lemma}
[Push--pull on a binary coproduct of two legs is the product of the two leg push--pulls]
  \label{mk-lem:pushPull_binary_coprod_prod}
  \lean{AlgebraicGeometry.pushPull_binary_coprod_prod,
    AlgebraicGeometry.isIso_prodLift_of_isLimit,
    AlgebraicGeometry.isIso_map_prodLift_of_isLimit,
    AlgebraicGeometry.coprodDecompMap,
    AlgebraicGeometry.isIso_coprodDecompMap}
  \dcref{custom}

  \uses{mk-def:push_pull_obj, mk-lem:pushPull_binary_leg_coherence, mk-lem:isIso_modules_of_toPresheaf,
    mk-lem:coprodPresheafObjIso_mathlib, mk-lem:isProductOfDisjoint_mathlib,
    mk-lem:evaluation_preserves_products_mathlib}
  Let \(Y_0, Y_1 : \operatorname{Over} X\) be two objects of the slice category over a scheme \(X\),
  and form the binary coproduct of their underlying schemes \(Y_0.\mathrm{left} \amalg Y_1.\mathrm{left}\)
  with structure map
  \(\operatorname{coprod.desc}\,Y_0.\mathrm{hom}\,Y_1.\mathrm{hom} : Y_0.\mathrm{left} \amalg
  Y_1.\mathrm{left} \to X\). Then the push--pull object on this binary coproduct is the binary product
  of the two per-leg push--pull objects:
  \[
    \operatorname{pushPullObj}\mathcal{F}\,\Bigl(\operatorname{Over.mk}\bigl(
      \operatorname{coprod.desc}\,Y_0.\mathrm{hom}\,Y_1.\mathrm{hom}\bigr)\Bigr)
      \;\cong\;
    \operatorname{pushPullObj}\mathcal{F}\,Y_0 \;\times\; \operatorname{pushPullObj}\mathcal{F}\,Y_1
    \qquad \text{in } X.\mathrm{Modules}.
  \]
\end{lemma}
\begin{proof}
  \uses{mk-lem:pushPull_binary_leg_coherence, mk-lem:isIso_modules_of_toPresheaf,
    mk-lem:coprodPresheafObjIso_mathlib, mk-lem:isProductOfDisjoint_mathlib,
    mk-lem:evaluation_preserves_products_mathlib}
  Write \(q = \operatorname{coprod.desc}\,Y_0.\mathrm{hom}\,Y_1.\mathrm{hom}\) and \(M =
  q^*\mathcal{F}\); let \(\operatorname{inl} : Y_0.\mathrm{left} \to Y_0.\mathrm{left} \amalg
  Y_1.\mathrm{left}\) and \(\operatorname{inr} : Y_1.\mathrm{left} \to Y_0.\mathrm{left} \amalg
  Y_1.\mathrm{left}\) be the coproduct inclusions, each an arrow over \(X\) (since \(q \circ
  \operatorname{inl} = Y_0.\mathrm{hom}\) and \(q \circ \operatorname{inr} = Y_1.\mathrm{hom}\)), and
  write \(\overline{\operatorname{inl}}, \overline{\operatorname{inr}}\) for these over-arrows.

  \emph{The canonical comparison (mandatory framing).} The asserted isomorphism is taken to be the
  \emph{canonical} map
  \[
    \operatorname{prod.lift}\bigl(\operatorname{pushPullMap}\mathcal{F}\,\overline{\operatorname{inl}},\,
      \operatorname{pushPullMap}\mathcal{F}\,\overline{\operatorname{inr}}\bigr) :
    \operatorname{pushPullObj}\mathcal{F}\,(\operatorname{Over.mk} q) \longrightarrow
    \operatorname{pushPullObj}\mathcal{F}\,Y_0 \times \operatorname{pushPullObj}\mathcal{F}\,Y_1,
  \]
  the product-lift of the two push--pull maps of the inclusions (Definition~\ref{mk-def:push_pull_map}).
  This framing is mandatory, not cosmetic: the downstream section-identification and {\v C}ech-%
  contractibility steps require the forward map of this isomorphism to be \emph{literally} this
  \(\operatorname{prod.lift}\) of push--pull maps; a non-canonical chain isomorphism with the same
  source and target would not serve them. So the task is precisely to prove this canonical map is an
  isomorphism.

  \emph{The reference chain.} Compare the canonical map against a manifestly invertible composite
  built from the binary disjoint-union decomposition. Write \(P = \operatorname{inl}_*
  (M|_{\operatorname{inl}})\) and \(Q = \operatorname{inr}_*(M|_{\operatorname{inr}})\) for the two
  restriction-to-component pushforwards, and let \(\operatorname{coprodDecompMap} M : M \to P \times
  Q\) be their product-lift of restriction-adjunction units. Then
  \[
    (\operatorname{pushforward} q).\mathrm{obj}\, M
    \xrightarrow{(\operatorname{pushforward} q)(\operatorname{coprodDecompMap} M)}
    (\operatorname{pushforward} q).\mathrm{obj}(P \times Q)
    \xrightarrow{\,\sim\,}
    (\operatorname{pushforward} q).\mathrm{obj}\, P \times (\operatorname{pushforward} q).\mathrm{obj}\, Q
    \xrightarrow{\,\sim\,}
    \operatorname{pushPullObj}\mathcal{F}\,Y_0 \times \operatorname{pushPullObj}\mathcal{F}\,Y_1
  \]
  is a chain of isomorphisms: the first arrow is \(q_*\) applied to the decomposition iso (an
  isomorphism by the leaf below); the second is that \(q_* = \operatorname{pushforward} q\) preserves
  the binary product \(P \times Q\); the third is \(\operatorname{prod.mapIso}\,\mathrm{idiso}_0\,
  \mathrm{idiso}_1\), the product of the two per-leg identifications \(\mathrm{idiso}_0 :
  (\operatorname{pushforward} q).\mathrm{obj}\, P \cong \operatorname{pushPullObj}\mathcal{F}\,Y_0\)
  and its \(\operatorname{inr}\)-twin. Each \(\mathrm{idiso}_0\) is assembled by pushing forward along
  \(q_*\) the inner isomorphism that chains the identification of restriction with open-immersion
  pullback (Lemma~\ref{mk-lem:restrictFunctorIsoPullback_mathlib}), the pullback-composition comparison
  for \(\operatorname{inl} \circ q\), and the over-triangle reindexing \(q \circ \operatorname{inl} =
  Y_0.\mathrm{hom}\), together with a transport of the codomain along this same triangle. That
  codomain transport is almost entirely definitional: pushforward-composition is identity-on-objects,
  so \(q_*(\operatorname{inl}_* N) = (\operatorname{inl} \circ q)_* N\) holds by definition
  and only the over-triangle relabelling remains.

  \emph{The leaf — binary disjoint-union decomposition.} The decomposition comparison
  \(\operatorname{coprodDecompMap} M\) is an isomorphism. By
  Lemma~\ref{mk-lem:isIso_modules_of_toPresheaf} it suffices to check this on sections over each open
  \(V \subseteq Y_0.\mathrm{left} \amalg Y_1.\mathrm{left}\). There the open splits as the disjoint
  union of its two component traces \(\operatorname{inl}^{-1}V\) and \(\operatorname{inr}^{-1}V\)
  (disjoint because the coproduct summands are), so the binary disjoint-union sheaf decomposition ---
  sections over a binary coproduct scheme are the product of the sections over the two preimages,
  Lemma~\ref{mk-lem:coprodPresheafObjIso_mathlib}, equivalently \(F(U \sqcup V') \cong F(U) \times
  F(V')\) for disjoint opens, Lemma~\ref{mk-lem:isProductOfDisjoint_mathlib} --- exhibits \(\Gamma(V, M)
  \cong \Gamma(\operatorname{inl}^{-1}V, M|_{\operatorname{inl}}) \times
  \Gamma(\operatorname{inr}^{-1}V, M|_{\operatorname{inr}})\), and this identification is exactly the
  pair of restriction-adjunction units making up \(\operatorname{coprodDecompMap} M\). (Evaluation at
  \(V\) preserves the binary product, Lemma~\ref{mk-lem:evaluation_preserves_products_mathlib}, so the
  right-hand sections are the corresponding product.) Hence \(\operatorname{coprodDecompMap} M\) is an
  isomorphism, and the reference chain above is an isomorphism.

  \emph{Matching the canonical map to the chain.} It remains to see the canonical
  \(\operatorname{prod.lift}\) of push--pull maps coincides with the (invertible) reference composite,
  whence it too is an isomorphism. Two maps into a binary product agree once they agree after
  composing with each projection, so it suffices to match the two legs. The
  \(\operatorname{inl}\)-leg of the reference composite is \((\operatorname{pushforward} q).\mathrm{map}
  \,u_0\) followed by \(\mathrm{idiso}_0\), where \(u_0\) is the \(\operatorname{inl}\)-component of
  \(\operatorname{coprodDecompMap} M\); the per-leg coherence identity
  (Lemma~\ref{mk-lem:pushPull_binary_leg_coherence}) says exactly that this equals
  \(\operatorname{pushPullMap}\mathcal{F}\,\overline{\operatorname{inl}}\), the \(\operatorname{inl}\)-leg
  of the canonical map, and symmetrically for \(\operatorname{inr}\). Therefore the canonical map
  equals the reference isomorphism and is itself an isomorphism in \(X.\mathrm{Modules}\).
\end{proof}

\begin{lemma}[Unit of a left-adjoint-uniqueness isomorphism]
  \label{mk-lem:unit_leftAdjointUniq_mathlib}
  \lean{CategoryTheory.Adjunction.unit_leftAdjointUniq_hom_app}
  \mathlibok
  \dcref{custom}

  If \(F, F'\) are both left adjoint to \(G\) (adjunctions \(\mathrm{adj}_1, \mathrm{adj}_2\)), the
  uniqueness isomorphism \(\operatorname{leftAdjointUniq}\mathrm{adj}_1\,\mathrm{adj}_2 : F \cong F'\)
  intertwines the units: composing the unit of \(\mathrm{adj}_1\) with \(G\) applied to the
  isomorphism recovers the unit of \(\mathrm{adj}_2\), i.e.\ at each object
  \(\mathrm{adj}_2.\mathrm{unit} = \mathrm{adj}_1.\mathrm{unit} \circ G\bigl(
  (\operatorname{leftAdjointUniq}\mathrm{adj}_1\,\mathrm{adj}_2).\mathrm{hom}\bigr)\).
\end{lemma}

\begin{lemma}[Per-leg coherence of the binary push--pull comparison]\label{mk-lem:pushPull_binary_leg_coherence}
  \lean{AlgebraicGeometry.pushPull_binary_leg_coherence}
  \leanok
  \dcref{custom}

  \uses{mk-def:push_pull_map, mk-lem:restrictFunctorIsoPullback_mathlib,
    mk-lem:unit_leftAdjointUniq_mathlib, mk-lem:pushPullCoprodLegIso}
  In the setting of Lemma~\ref{mk-lem:pushPull_binary_coprod_prod}, with \(q =
  \operatorname{coprod.desc}\,Y_0.\mathrm{hom}\,Y_1.\mathrm{hom}\) and \(M = q^*\mathcal{F}\), write
  \(u_0\) for the \(\operatorname{inl}\)-component of \(\operatorname{coprodDecompMap} M\) (the unit of
  the restriction adjunction along \(\operatorname{inl}\) at \(M\)) and \(\mathrm{idiso}_0 :
  (\operatorname{pushforward} q).\mathrm{obj}\bigl(\operatorname{inl}_*(M|_{\operatorname{inl}})\bigr)
  \cong \operatorname{pushPullObj}\mathcal{F}\,Y_0\) for the per-leg identification of the reference
  chain. Then the \(\operatorname{inl}\)-leg of the canonical comparison factors through these:
  \[
    \operatorname{pushPullMap}\mathcal{F}\,\overline{\operatorname{inl}}
      \;=\; (\operatorname{pushforward} q).\mathrm{map}\, u_0 \;\circ\; \mathrm{idiso}_0.\mathrm{hom},
  \]
  and symmetrically for the \(\operatorname{inr}\)-leg.
\end{lemma}
\begin{proof}
  \uses{mk-def:push_pull_map, mk-lem:restrictFunctorIsoPullback_mathlib,
    mk-lem:unit_leftAdjointUniq_mathlib, mk-lem:pushPullCoprodLegIso}
  The two legs are symmetric; treat \(\operatorname{inl}\). Unfold the push--pull map on the
  left-hand side to its raw five-step form (Definition~\ref{mk-def:push_pull_map}):
  \(\operatorname{pushPullMap}\mathcal{F}\,\overline{\operatorname{inl}}\) is \(q_*\) applied to the
  composite of the restriction-adjunction unit \(\eta^{\operatorname{inl}}\) at \(M\) with
  \(\operatorname{inl}_*\) of the pullback-composition comparison for \(\operatorname{inl} \circ q\)
  acting on \(\mathcal{F}\), followed by the over-triangle transport that \(\mathrm{idiso}_0\)
  encodes. Now the open immersion \(\operatorname{inl}\) presents its restriction functor as a left
  adjoint, and the comparison \(\operatorname{restrictFunctorIsoPullback}\)
  (Lemma~\ref{mk-lem:restrictFunctorIsoPullback_mathlib}) is precisely the uniqueness isomorphism
  between this restriction functor and the open-immersion pullback. The unit-intertwining identity
  (Lemma~\ref{mk-lem:unit_leftAdjointUniq_mathlib}) rewrites the adjunction unit as
  \(\eta^{\operatorname{inl}} = u_0 \circ \operatorname{inl}_*\bigl(
  (\operatorname{restrictFunctorIsoPullback}\operatorname{inl}).\mathrm{hom}\bigr)\). Substituting and
  cancelling the comparison isomorphism against the matching inverse factor inside \(\mathrm{idiso}_0\),
  the residual identity is an equality of two canonical equality-transport morphisms along the single
  open-set equality \(q \circ \operatorname{inl} = Y_0.\mathrm{hom}\); these agree by proof-%
  irrelevance of that equality. Hence \(\operatorname{pushPullMap}\mathcal{F}\,
  \overline{\operatorname{inl}} = (\operatorname{pushforward} q).\mathrm{map}\, u_0 \circ
  \mathrm{idiso}_0.\mathrm{hom}\), as claimed.
\end{proof}

\begin{lemma}[Canonical leg identification of the binary push--pull chain]\label{mk-lem:pushPullCoprodLegIso}
  \lean{AlgebraicGeometry.pushPullCoprodLegIso}
  \leanok
  \dcref{custom}

  \uses{mk-def:push_pull_obj, mk-lem:restrictFunctorIsoPullback_mathlib}
  In the setting of Lemma~\ref{mk-lem:pushPull_binary_coprod_prod}, let \(c\) be one of the coproduct
  inclusions \(\operatorname{inl}, \operatorname{inr} : C \to Y_0.\mathrm{left} \amalg Y_1.\mathrm{left}\),
  with structure map \(p_C : C \to X\) and over-triangle \(w_C : c \circ q = p_C\). Writing \(M =
  q^*\mathcal{F}\), there is a canonical isomorphism identifying the corresponding factor of \(q_*(P
  \times Q)\) with the per-leg push--pull object:
  \[
    (\operatorname{pushforward} q).\mathrm{obj}\Bigl((\operatorname{pushforward} c).\mathrm{obj}
      \bigl((q^*\mathcal{F}).\operatorname{restrict} c\bigr)\Bigr)
      \;\cong\;
    \operatorname{pushPullObj}\mathcal{F}\,(\operatorname{Over.mk} p_C).
  \]
\end{lemma}
\begin{proof}
  \uses{mk-lem:restrictFunctorIsoPullback_mathlib}
  Push forward along \(q_*\) the inner isomorphism chaining the identification of restriction with
  open-immersion pullback (Lemma~\ref{mk-lem:restrictFunctorIsoPullback_mathlib}), the
  pullback-composition comparison for \(c \circ q\), and the over-triangle reindexing \(c \circ q = p_C\);
  then transport the codomain along the same triangle by an equality-of-objects rewrite. This is the
  per-leg identification \(\mathrm{idiso}\) of the reference chain in
  Lemma~\ref{mk-lem:pushPull_binary_coprod_prod}.
\end{proof}

\begin{lemma}[Splitting an $\operatorname{Option}$-indexed coproduct]\label{mk-lem:sigmaOptionIso}
  \lean{CategoryTheory.sigmaOptionIso}
  \leanok
  \dcref{custom}

  Let \(\mathcal{C}\) be a category and \(Z : \operatorname{Option}\alpha \to \mathcal{C}\) a family
  for which the relevant coproducts exist (the total coproduct \(\coprod Z\), the coproduct
  \(\coprod_{a : \alpha} Z(\operatorname{some} a)\) over the proper part, and their binary coproduct).
  Then the total coproduct splits off its \(\operatorname{none}\) summand:
  \[
    \coprod_{o : \operatorname{Option}\alpha} Z(o) \;\cong\; Z(\operatorname{none}) \amalg
      \Bigl(\coprod_{a : \alpha} Z(\operatorname{some} a)\Bigr).
  \]
\end{lemma}
\begin{proof}

  The forward map is \(\operatorname{Sigma.desc}\) of the cocone sending the \(\operatorname{none}\)
  component to the left binary inclusion and each \(\operatorname{some} a\) component to the composite
  of the \(a\)-th inclusion into \(\coprod_{a} Z(\operatorname{some} a)\) with the right binary
  inclusion. The inverse is \(\operatorname{coprod.desc}\) of the \(\operatorname{none}\) inclusion and
  \(\operatorname{Sigma.desc}\) of the \(\operatorname{some}\) inclusions. The two round-trip identities
  follow from the universal properties of the coproducts, checked componentwise on each
  \(\operatorname{none}\) and \(\operatorname{some}\) term. (Mathlib provides the associativity
  isomorphism for nested coproducts but no \(\operatorname{Option}\)-split, so this is supplied
  directly.)
\end{proof}

\begin{lemma}[Push--pull objects transport along a slice isomorphism]\label{mk-lem:pushPullObjCongr}
  \lean{AlgebraicGeometry.pushPullObjCongr}
  \leanok
  \dcref{custom}

  \uses{mk-def:push_pull_obj, mk-def:push_pull_map}
  Let \(\mathcal{F} : X.\mathrm{Modules}\) and let \(e : Y \cong Y'\) be an isomorphism in
  \(\operatorname{Over} X\). Then the push--pull objects on \(Y\) and \(Y'\) are isomorphic,
  \[
    \operatorname{pushPullObj}\mathcal{F}\,Y \;\cong\; \operatorname{pushPullObj}\mathcal{F}\,Y'.
  \]
\end{lemma}
\begin{proof}
  \uses{mk-def:push_pull_map}
  The push--pull object is a contravariant functor of its slice argument through the push--pull map
  (Definition~\ref{mk-def:push_pull_map}). Take the forward map to be
  \(\operatorname{pushPullMap}\mathcal{F}\,e.\mathrm{inv}\) and the backward map
  \(\operatorname{pushPullMap}\mathcal{F}\,e.\mathrm{hom}\) (contravariance: the \(\mathrm{hom}\)
  direction uses \(e.\mathrm{inv}\)). The two composites are identities by the functoriality of the
  push--pull map under composition and identity of over-morphisms.
\end{proof}

\begin{lemma}[Slice lift of the $\operatorname{Option}$-coproduct splitting]\label{mk-lem:over_sigmaOptionIso}
  \lean{AlgebraicGeometry.overSigmaOptionIso}
  \leanok
  \dcref{custom}

  \uses{mk-lem:sigmaOptionIso}
  Let \(\mathrm{legs} : \operatorname{Option}\alpha \to \operatorname{Over} X\) be a family of objects
  of the slice over \(X\). The splitting of Lemma~\ref{mk-lem:sigmaOptionIso} lifts to an isomorphism in
  \(\operatorname{Over} X\) between the coproduct slice object and the binary-split slice object:
  \[
    \operatorname{Over.mk}\bigl(\operatorname{Sigma.desc}(o \mapsto (\mathrm{legs}\,o).\mathrm{hom})\bigr)
      \;\cong\;
    \operatorname{Over.mk}\bigl(\operatorname{coprod.desc}\,(\mathrm{legs}\,\operatorname{none}).\mathrm{hom}\,
      (\operatorname{Sigma.desc}(a \mapsto (\mathrm{legs}\,(\operatorname{some} a)).\mathrm{hom}))\bigr).
  \]
\end{lemma}
\begin{proof}
  \uses{mk-lem:sigmaOptionIso}
  Build the slice isomorphism as \(\operatorname{Over.isoMk}\) of the underlying-scheme isomorphism
  \(\operatorname{sigmaOptionIso}(o \mapsto (\mathrm{legs}\,o).\mathrm{left})\)
  (Lemma~\ref{mk-lem:sigmaOptionIso}), together with the verification that its forward map commutes with
  the two structure arrows. The latter is checked componentwise: on each \(\operatorname{none}\) and
  \(\operatorname{some} a\) component both descent maps reduce to \((\mathrm{legs}\,o).\mathrm{hom}\) by
  the universal properties of each coproduct.
\end{proof}

\begin{lemma}[Splitting an $\operatorname{Option}$-indexed product]\label{mk-lem:piOptionIso}
  \lean{AlgebraicGeometry.piOptionIso}
  \leanok
  \dcref{custom}

  Dual to Lemma~\ref{mk-lem:sigmaOptionIso}: let \(W : \operatorname{Option}\alpha \to \mathcal{C}\) with
  the relevant products existing. Then the total product splits off its \(\operatorname{none}\) factor:
  \[
    \prod_{o : \operatorname{Option}\alpha} W(o) \;\cong\; W(\operatorname{none}) \times
      \Bigl(\prod_{a : \alpha} W(\operatorname{some} a)\Bigr).
  \]
\end{lemma}
\begin{proof}

  The same construction as Lemma~\ref{mk-lem:sigmaOptionIso} with all arrows reversed: the forward map is
  \(\operatorname{prod.lift}\) of the projection onto \(W(\operatorname{none})\) and the
  \(\operatorname{Pi.lift}\) of the \(\operatorname{some}\) projections; the inverse is
  \(\operatorname{Pi.lift}\) of the matching projections out of the binary product. The round-trip
  identities follow from the universal properties of the products.
\end{proof}

\begin{lemma}[Push--pull on the empty coproduct is terminal]
  \label{mk-lem:pushPull_coprod_prod_empty}
  \lean{AlgebraicGeometry.pushPull_coprod_prod_empty,
    AlgebraicGeometry.isZero_modules_of_isEmpty}
  \dcref{custom}

  \uses{mk-def:push_pull_obj, mk-def:coprodToProdMap}
  For the empty index type \(\iota = \operatorname{PEmpty}\) and the empty family of legs, the coproduct
  \(\coprod_{i} \mathrm{legs}\,i\) is the initial scheme \(\bot\); the push--pull object on its unique
  structure map is terminal in \(X.\mathrm{Modules}\), hence isomorphic to the empty product
  \(\prod_{i \in \operatorname{PEmpty}} \operatorname{pushPullObj}\mathcal{F}\,(\mathrm{legs}\,i)\),
  which is likewise terminal:
  \[
    \operatorname{pushPullObj}\mathcal{F}\,\Bigl(\operatorname{Over.mk}\bigl(
      \operatorname{Sigma.desc}(i \mapsto (\mathrm{legs}\,i).\mathrm{hom})\bigr)\Bigr)
      \;\cong\;
    \prod_{i \in \operatorname{PEmpty}} \operatorname{pushPullObj}\mathcal{F}\,(\mathrm{legs}\,i).
  \]
\end{lemma}
\begin{proof}
  \uses{mk-def:push_pull_obj, mk-def:coprodToProdMap}
  The coproduct over the empty index type is the initial scheme \(\bot\): the structure map
  \(\operatorname{Sigma.desc}\) out of \(\coprod_{i \in \operatorname{PEmpty}} (\mathrm{legs}\,i).
  \mathrm{left}\) is the unique morphism \(q : \bot \to X\) out of the initial object. Every sheaf of
  modules over the empty scheme is the zero object: the only open of \(\bot\) is the empty open, over
  which the structure sheaf has subsingleton (hence zero) sections, so every
  \(\mathcal{O}_\bot\)-module has zero sections over every open and is therefore the zero module sheaf.
  In particular the pulled-back module \((\operatorname{pullback} q).\mathrm{obj}\,\mathcal{F}\) is the
  zero object of \(\bot.\mathrm{Modules}\).

  The push--pull object is \(\operatorname{pushPullObj}\mathcal{F}\,(\operatorname{Over.mk} q) =
  (\operatorname{pushforward} q).\mathrm{obj}\bigl((\operatorname{pullback} q).\mathrm{obj}\,
  \mathcal{F}\bigr)\). The module pushforward is additive, so it sends the zero object to a zero
  object, which is terminal. The empty product \(\prod_{i \in \operatorname{PEmpty}}
  \operatorname{pushPullObj}\mathcal{F}\,(\mathrm{legs}\,i)\) is likewise terminal. A morphism between
  terminal objects is an isomorphism, so the comparison map
  \(\operatorname{coprodToProdMap}\mathcal{F}\,\mathrm{legs}\) (Definition~\ref{mk-def:coprodToProdMap}) is
  an isomorphism.

  The single remaining formal obligation is that the pulled-back module is zero, namely
  \(\operatorname{IsZero}\bigl((\operatorname{pullback} q).\mathrm{obj}\,\mathcal{F}\bigr)\) for \(q\)
  the map out of the empty scheme. This is a presheaf-of-modules \(\operatorname{IsZero}\) statement
  established pointwise through the faithful sections-to-presheaf functor, using the subsingleton
  sections of the structure sheaf over the empty scheme.
\end{proof}

\begin{definition}[The $i$-th coproduct inclusion as an over-morphism]\label{mk-def:coprodOverIncl}
  \lean{AlgebraicGeometry.coprodOverIncl}
  \leanok
  \dcref{custom}

  Let \((\mathrm{legs} : \iota \to \operatorname{Over} X)\) be a family of objects of the slice over a
  scheme \(X\), and suppose the coproduct \(\coprod_i (\mathrm{legs}\,i).\mathrm{left}\) exists. For each
  index \(i\), the canonical coproduct inclusion
  \(\iota_i : (\mathrm{legs}\,i).\mathrm{left} \to \coprod_i (\mathrm{legs}\,i).
  \mathrm{left}\) is promoted to a morphism of the slice category \(\operatorname{Over} X\),
  \[
    \overline{\iota_i} : \mathrm{legs}\,i \longrightarrow \operatorname{Over.mk}\bigl(
      \operatorname{Sigma.desc}(i \mapsto (\mathrm{legs}\,i).\mathrm{hom})\bigr),
  \]
  whose underlying scheme map is \(\iota_i\); the over-triangle commutes because
  \(\iota_i \circ \operatorname{Sigma.desc}(i \mapsto (\mathrm{legs}\,i).
  \mathrm{hom}) = (\mathrm{legs}\,i).\mathrm{hom}\) by the defining property of
  \(\operatorname{Sigma.desc}\).
\end{definition}

\begin{definition}[Canonical coproduct-to-product comparison of push--pull objects]\label{mk-def:coprodToProdMap}
  \lean{AlgebraicGeometry.coprodToProdMap}
  \leanok
  \dcref{custom}

  \uses{mk-def:push_pull_obj, mk-def:push_pull_map, mk-def:coprodOverIncl}
  In the setting of Definition~\ref{mk-def:coprodOverIncl}, assume in addition that the product
  \(\prod_i \operatorname{pushPullObj}\mathcal{F}\,(\mathrm{legs}\,i)\) exists. The \emph{canonical
  comparison map}
  \[
    \operatorname{coprodToProdMap}\mathcal{F}\,\mathrm{legs} :
      \operatorname{pushPullObj}\mathcal{F}\,\bigl(\operatorname{Over.mk}(\operatorname{Sigma.desc}
        (i \mapsto (\mathrm{legs}\,i).\mathrm{hom}))\bigr) \longrightarrow
      \prod_i \operatorname{pushPullObj}\mathcal{F}\,(\mathrm{legs}\,i)
  \]
  is the \(\operatorname{Pi.lift}\) whose \(i\)-th component is the push--pull map
  \(\operatorname{pushPullMap}\mathcal{F}\,\overline{\iota_i}\) (Definition~\ref{mk-def:push_pull_map}) of
  the \(i\)-th coproduct inclusion (Definition~\ref{mk-def:coprodOverIncl}). Equivalently, it is the unique
  map whose composite with the \(i\)-th product projection equals
  \(\operatorname{pushPullMap}\mathcal{F}\,\overline{\iota_i}\) for every \(i\). This is the canonical
  framing kept invariant throughout the finite induction of
  Lemma~\ref{mk-lem:pushPull_coprod_prod}.
\end{definition}

\begin{lemma}[The comparison map is invariant under index reindexing]\label{mk-lem:coprodToProd_isIso_of_equiv}
  \lean{AlgebraicGeometry.coprodToProd_isIso_of_equiv}
  \leanok
  \dcref{custom}

  \uses{mk-def:push_pull_obj, mk-def:coprodOverIncl, mk-def:coprodToProdMap, mk-lem:pushPullObjCongr}
  Let \(e : \alpha \simeq \beta\) be an equivalence of finite index types and let
  \(\mathrm{legs} : \beta \to \operatorname{Over} X\) be a family of slice objects. If the canonical
  comparison map \(\operatorname{coprodToProdMap}\mathcal{F}\,(\mathrm{legs} \circ e)\) for the reindexed
  family is an isomorphism, then so is \(\operatorname{coprodToProdMap}\mathcal{F}\,\mathrm{legs}\). In
  other words, the predicate ``\(\operatorname{coprodToProdMap}\mathcal{F}\,(\,\cdot\,)\) is an
  isomorphism'' is stable under reindexing the leg family along an index equivalence.
\end{lemma}
\begin{proof}
  \uses{mk-def:coprodOverIncl, mk-def:coprodToProdMap, mk-lem:pushPullObjCongr}
  The induction hypothesis supplies that
  \(\operatorname{coprodToProdMap}\mathcal{F}\,(\mathrm{legs} \circ e)\) is an isomorphism; it suffices
  to exhibit \(\operatorname{coprodToProdMap}\mathcal{F}\,\mathrm{legs}\) as the conjugate of that map by
  two reindexing isomorphisms, since conjugating an isomorphism by isomorphisms again yields an
  isomorphism.

  \emph{Transporting the source.} Reindexing the coproduct along \(e\) (the whiskered reindexing
  \(\operatorname{Sigma.whiskerEquiv} e\,(i \mapsto \operatorname{Iso.refl})\)) together with the
  identification of descent objects gives an isomorphism in \(\operatorname{Over} X\)
  \[
    \operatorname{Over.mk}\bigl(\operatorname{Sigma.desc}(o \mapsto (\mathrm{legs}\,o).\mathrm{hom})\bigr)
      \;\cong\;
    \operatorname{Over.mk}\bigl(\operatorname{Sigma.desc}(a \mapsto ((\mathrm{legs} \circ e)\,a).
      \mathrm{hom})\bigr).
  \]
  Carrying the push--pull object across this slice isomorphism by
  Lemma~\ref{mk-lem:pushPullObjCongr} identifies the source of
  \(\operatorname{coprodToProdMap}\mathcal{F}\,\mathrm{legs}\) with the source of
  \(\operatorname{coprodToProdMap}\mathcal{F}\,(\mathrm{legs} \circ e)\).

  \emph{Transporting the target.} Reindexing the product along \(e\) (the whiskered reindexing
  \(\operatorname{Pi.whiskerEquiv} e\,(a \mapsto \operatorname{Iso.refl})\), equivalently the
  product reindexing isomorphism \(\operatorname{Pi.mapIso}\) along \(e\)) identifies
  \(\prod_{o : \beta} \operatorname{pushPullObj}\mathcal{F}\,(\mathrm{legs}\,o)\) with
  \(\prod_{a : \alpha} \operatorname{pushPullObj}\mathcal{F}\,((\mathrm{legs} \circ e)\,a)\).

  \emph{Matching the canonical form.} It remains to check that conjugating
  \(\operatorname{coprodToProdMap}\mathcal{F}\,(\mathrm{legs} \circ e)\) by these two reindexing
  isomorphisms reproduces \(\operatorname{coprodToProdMap}\mathcal{F}\,\mathrm{legs}\). Since both maps
  are determined by their composites with the product projections, the check is performed projection by
  projection: by the defining projection law of the comparison map
  (\(\operatorname{coprodToProdMap}\mathcal{F}\,\mathrm{legs} \circ \pi_i =
  \operatorname{pushPullMap}\mathcal{F}\,\overline{\iota_i}\), Definition~\ref{mk-def:coprodToProdMap}),
  each projection of either side is the push--pull map of a coproduct inclusion, and the reindexing
  isomorphisms carry the \(i\)-th inclusion \(\overline{\iota_i}\) of the \(\beta\)-coproduct to the
  \(e^{-1}(i)\)-th inclusion of the \(\alpha\)-coproduct. The naturality of
  \(\operatorname{pushPullMap}\mathcal{F}\) under this reindexing of the inclusions makes the two
  projections agree. Hence the conjugate equals
  \(\operatorname{coprodToProdMap}\mathcal{F}\,\mathrm{legs}\), and the induction hypothesis yields
  \(\operatorname{IsIso}\bigl(\operatorname{coprodToProdMap}\mathcal{F}\,\mathrm{legs}\bigr)\).
\end{proof}

\begin{lemma}[Adjoining one index to the comparison-map induction]
  \label{mk-lem:coprodToProd_isIso_option}
  \lean{AlgebraicGeometry.coprodToProd_isIso_option, AlgebraicGeometry.pushPullObjCongr_hom,
    AlgebraicGeometry.pushPull_binary_coprod_prod_hom, AlgebraicGeometry.piOptionIso_inv_π_none,
    AlgebraicGeometry.piOptionIso_inv_π_some}
  \dcref{custom}

  \uses{mk-def:coprodToProdMap, mk-lem:pushPull_binary_coprod_prod, mk-lem:over_sigmaOptionIso,
    mk-lem:piOptionIso, mk-lem:pushPullObjCongr}
  Suppose that for every family \(\mathrm{legs} : \alpha \to \operatorname{Over} X\) the canonical
  comparison map \(\operatorname{coprodToProdMap}\mathcal{F}\,\mathrm{legs}\) is an isomorphism. Then for
  every family \(\mathrm{legs} : \operatorname{Option}\alpha \to \operatorname{Over} X\) the comparison
  map \(\operatorname{coprodToProdMap}\mathcal{F}\,\mathrm{legs}\) is an isomorphism.
\end{lemma}
\begin{proof}
  \uses{mk-def:coprodToProdMap, mk-lem:pushPull_binary_coprod_prod, mk-lem:over_sigmaOptionIso,
    mk-lem:piOptionIso, mk-lem:pushPullObjCongr}
  Write \(\mathrm{ls} = \mathrm{legs} \circ \operatorname{some}\) for the restriction to the proper
  part. The \(\operatorname{Option}\alpha\)-coproduct splits as the binary coproduct of the
  \(\operatorname{none}\) leg and the \(\alpha\)-coproduct (Lemma~\ref{mk-lem:over_sigmaOptionIso});
  transporting the push--pull object across this slice isomorphism (Lemma~\ref{mk-lem:pushPullObjCongr})
  and applying the binary decomposition (Lemma~\ref{mk-lem:pushPull_binary_coprod_prod}) together with the
  induction hypothesis on \(\mathrm{ls}\) presents the source as a binary product, which the inverse of
  the \(\operatorname{Option}\)-product splitting (Lemma~\ref{mk-lem:piOptionIso}) reassembles into the
  total product \(\prod_{o : \operatorname{Option}\alpha} \operatorname{pushPullObj}\mathcal{F}\,
  (\mathrm{legs}\,o)\). A per-\(\operatorname{Option}\)-case projection chase then matches this composite
  reference isomorphism with the canonical comparison map
  \(\operatorname{coprodToProdMap}\mathcal{F}\,\mathrm{legs}\) (Definition~\ref{mk-def:coprodToProdMap}),
  whence the latter is an isomorphism.
\end{proof}

\begin{lemma}
[Push--pull on a finite coproduct of legs is the product of the leg push--pulls]
  \label{mk-lem:pushPull_coprod_prod}
  \lean{AlgebraicGeometry.pushPull_coprod_prod, AlgebraicGeometry.coprodToProdMap_comp_π,
    AlgebraicGeometry.isIso_coprodToProdMap}
  \dcref{custom}

  \uses{mk-def:push_pull_obj, mk-def:coprodOverIncl, mk-def:coprodToProdMap, mk-lem:isIso_modules_of_toPresheaf,
    mk-lem:pushPull_binary_coprod_prod, mk-lem:pushPullObjCongr, mk-lem:over_sigmaOptionIso, mk-lem:piOptionIso,
    mk-lem:pushPull_coprod_prod_empty, mk-lem:coprodToProd_isIso_of_equiv, mk-lem:coprodToProd_isIso_option}
  Let \((\mathrm{legs} : \iota \to \operatorname{Over} X)\) be a finite family of objects of
  the slice category over a scheme \(X\) (\(\iota\) finite), and form the coproduct
  \(\coprod_{i \in \iota} \mathrm{legs}\,i\) with structure map
  \(\operatorname{Sigma.desc}(i \mapsto (\mathrm{legs}\,i).\mathrm{hom}) : \coprod_i
  \mathrm{legs}\,i \to X\). Then the push--pull object on this coproduct is the product of
  the per-leg push--pull objects:
  \[
    \operatorname{pushPullObj}\mathcal{F}\,\Bigl(\operatorname{Over.mk}\bigl(
      \operatorname{Sigma.desc}(i \mapsto (\mathrm{legs}\,i).\mathrm{hom})\bigr)\Bigr)
      \;\cong\;
    \prod_{i \in \iota} \operatorname{pushPullObj}\mathcal{F}\,(\mathrm{legs}\,i)
    \qquad \text{in } X.\mathrm{Modules}.
  \]
  This is the general indexed-coproduct\(\to\)product disjoint-union decomposition of a
  module sheaf: the legs may have overlapping images in \(X\), but they are disjoint inside
  the coproduct space \(\coprod_i \mathrm{legs}\,i\), so push--pull splits componentwise.
\end{lemma}
\begin{proof}
  \uses{mk-def:push_pull_obj, mk-def:coprodOverIncl, mk-def:coprodToProdMap, mk-lem:isIso_modules_of_toPresheaf,
    mk-lem:pushPull_binary_coprod_prod, mk-lem:pushPullObjCongr, mk-lem:over_sigmaOptionIso, mk-lem:piOptionIso,
    mk-lem:pushPull_coprod_prod_empty, mk-lem:coprodToProd_isIso_of_equiv, mk-lem:coprodToProd_isIso_option}
  Argue by induction on the finite index type \(\iota\), using the finite-index recursion principle: a
  property of finite types that is stable under index equivalence, holds for the empty type, and is
  preserved by adjoining one element (passing from \(\alpha\) to \(\operatorname{Option}\alpha\)) holds
  for every finite type. Throughout, the isomorphism is kept in its \emph{canonical} form as
  the product of push--pull maps \(\prod_{i} \operatorname{pushPullMap}\mathcal{F}\,\overline{\iota_i}\),
  where \(\overline{\iota_i}\) denotes the \(i\)-th coproduct inclusion viewed as an over-morphism
  into \(\operatorname{Over.mk}(\operatorname{Sigma.desc}(i \mapsto (\mathrm{legs}\,i).\mathrm{hom}))\).

  \emph{Empty base.} For \(\iota = \operatorname{PEmpty}\) the claim is
  Lemma~\ref{mk-lem:pushPull_coprod_prod_empty}: both sides are terminal.

  \emph{Reindexing.} The property is stable under index equivalence by
  Lemma~\ref{mk-lem:coprodToProd_isIso_of_equiv}.

  \emph{Adjoining one index.} Suppose the claim holds for \(\alpha\); the passage to
  \(\operatorname{Option}\alpha\) is Lemma~\ref{mk-lem:coprodToProd_isIso_option}: the
  \(\operatorname{Option}\)-coproduct splits as the binary coproduct of the \(\operatorname{none}\) leg
  and the \(\alpha\)-coproduct (Lemma~\ref{mk-lem:over_sigmaOptionIso}), the binary decomposition
  (Lemma~\ref{mk-lem:pushPull_binary_coprod_prod}) and the induction hypothesis present the source as a
  binary product, the inverse product splitting (Lemma~\ref{mk-lem:piOptionIso}) reassembles it into the
  total product, and a per-\(\operatorname{Option}\)-case projection chase matches the result with the
  canonical comparison map. (All isomorphisms produced live in \(X.\mathrm{Modules}\), verified through
  the underlying presheaf by Lemma~\ref{mk-lem:isIso_modules_of_toPresheaf} where needed.)
\end{proof}

\begin{lemma}

[Push--pull over the {\v C}ech backbone is the product over intersection opens]
  \label{mk-lem:pushPull_sigma_iso}
  \lean{AlgebraicGeometry.pushPull_sigma_iso, AlgebraicGeometry.pushPull_sigma_iso_π}
  \dcref{custom}

  \uses{mk-def:push_pull_obj, mk-lem:cech_backbone_left_sigma, mk-lem:pushPull_coprod_prod}
  With the notation of Lemma~\ref{mk-lem:cech_backbone_left_sigma}, the push--pull object on
  the degree-\(p\) backbone decomposes as a product over the multi-indices:
  \[
    \operatorname{pushPullObj}\mathcal{F}\,Y_p
      \;\cong\;
    \prod_{\sigma : \mathrm{Fin}(p+1) \to I}
      \operatorname{pushPullObj}\mathcal{F}\,(\operatorname{Over.mk} j_\sigma)
    \qquad \text{in } X.\mathrm{Modules},
  \]
  equivalently \((q_p)_*\,(q_p)^*\mathcal{F} \cong \prod_\sigma (j_\sigma)_*\,
  (j_\sigma)^*\mathcal{F}\).
\end{lemma}
\begin{proof}
  \uses{mk-lem:cech_backbone_left_sigma, mk-lem:pushPull_coprod_prod}
  By Lemma~\ref{mk-lem:cech_backbone_left_sigma} the backbone \(Y_p\) is the coproduct
  \(\coprod_\sigma U_\sigma\) with structure map \(q_p = \operatorname{Sigma.desc}(\sigma
  \mapsto j_\sigma)\). Apply Lemma~\ref{mk-lem:pushPull_coprod_prod} to this coproduct family,
  with the finite multi-index type \(\sigma : \mathrm{Fin}(p+1) \to I\) as the index and the
  intersection-open legs \(\operatorname{Over.mk} j_\sigma\), to obtain the asserted product
  decomposition.
\end{proof}

\begin{lemma}[Sections of one push--pull leg over an open]\label{mk-lem:pushPull_leg_sections}
  \lean{AlgebraicGeometry.pushPull_leg_sections}
  \leanok
  \dcref{custom}

  \uses{mk-def:push_pull_obj, mk-def:cech_free_presheaf_complex, mk-lem:pushforward_obj_obj_mathlib,
    mk-lem:restrictFunctorIsoPullback_mathlib, mk-lem:restrict_obj_mathlib}
  Let \(j_\sigma : U_\sigma \hookrightarrow X\) be the open immersion of the intersection
  open \(U_\sigma = \operatorname{coverInterOpen}\mathcal{U}\,\sigma\), and let
  \(V \subseteq X\) be an open. Then there is a natural isomorphism of section groups
  \[
    \Gamma\bigl(V,\, \operatorname{pushPullObj}\mathcal{F}\,(\operatorname{Over.mk}
      j_\sigma)\bigr)
      \;\cong\;
    \Gamma\bigl(U_\sigma \cap V,\, \mathcal{F}\bigr).
  \]
\end{lemma}
\begin{proof}
  \uses{mk-lem:pushforward_obj_obj_mathlib, mk-lem:restrictFunctorIsoPullback_mathlib,
    mk-lem:restrict_obj_mathlib}
  This follows from three standard identifications. First, the leg is
  \((j_\sigma)_*\,(j_\sigma)^*\mathcal{F}\); pushforward sections are computed by preimage,
  \(\Gamma\bigl(V, (j_\sigma)_* N\bigr) = \Gamma\bigl(j_\sigma^{-1}(V), N\bigr)\)
  (Lemma~\ref{mk-lem:pushforward_obj_obj_mathlib}, an equality), applied with \(N =
  (j_\sigma)^*\mathcal{F}\). Second, pullback along the open immersion \(j_\sigma\) is the
  restriction functor: \((j_\sigma)^*\mathcal{F} \cong \mathcal{F}.\operatorname{restrict}
  j_\sigma\) by Lemma~\ref{mk-lem:restrictFunctorIsoPullback_mathlib}. Third, the sections of a
  restricted module over \(j_\sigma^{-1}(V)\) are the sections of \(\mathcal{F}\) over the
  image open (Lemma~\ref{mk-lem:restrict_obj_mathlib}, an equality): \(\Gamma\bigl(j_\sigma^{-1}
  (V), \mathcal{F}.\operatorname{restrict} j_\sigma\bigr) = \Gamma\bigl(j_\sigma\bigl(
  j_\sigma^{-1}(V)\bigr), \mathcal{F}\bigr) = \Gamma\bigl(U_\sigma \cap V, \mathcal{F}\bigr)\),
  the last equality because the image of the preimage of \(V\) under the open immersion
  \(j_\sigma : U_\sigma \hookrightarrow X\) is exactly \(U_\sigma \cap V\). Composing the
  three gives the stated isomorphism.
\end{proof}

\begin{lemma}[Degreewise section identification of the {\v C}ech backbone]\label{mk-lem:pushPull_eval_prod_iso}
  \lean{AlgebraicGeometry.pushPull_eval_prod_iso}
  \leanok
  \dcref{custom}

  \uses{mk-lem:pushPull_sigma_iso, mk-lem:pushPull_leg_sections,
    mk-lem:evaluation_preserves_products_mathlib}
  For an open \(V \subseteq X\) and every degree \(p\) there is an isomorphism
  \[
    \Gamma\bigl(V,\, \operatorname{pushPullObj}\mathcal{F}\,Y_p\bigr)
      \;\cong\;
    \prod_{\sigma : \mathrm{Fin}(p+1) \to I}
      \Gamma\bigl(\operatorname{coverInterOpen}\mathcal{U}\,\sigma \cap V,\, \mathcal{F}\bigr).
  \]
\end{lemma}
\begin{proof}
  \uses{mk-lem:pushPull_sigma_iso, mk-lem:pushPull_leg_sections,
    mk-lem:evaluation_preserves_products_mathlib}
  Combine the product decomposition of the backbone push--pull object
  (Lemma~\ref{mk-lem:pushPull_sigma_iso}), \(\operatorname{pushPullObj}\mathcal{F}\,Y_p \cong
  \prod_\sigma \operatorname{pushPullObj}\mathcal{F}\,(\operatorname{Over.mk} j_\sigma)\),
  with the fact that the section (evaluation-at-\(V\)) functor on \(X.\mathrm{Modules}\)
  preserves products (Lemma~\ref{mk-lem:evaluation_preserves_products_mathlib}): \(\Gamma(V,
  \prod_\sigma N_\sigma) \cong \prod_\sigma \Gamma(V, N_\sigma)\). This reduces the
  left-hand side to \(\prod_\sigma \Gamma\bigl(V, \operatorname{pushPullObj}\mathcal{F}\,
  (\operatorname{Over.mk} j_\sigma)\bigr)\), and the per-leg section identification
  (Lemma~\ref{mk-lem:pushPull_leg_sections}) rewrites each factor as \(\Gamma(U_\sigma \cap V,
  \mathcal{F})\).
\end{proof}

\begin{lemma}[An additive functor commutes with augmentation, degreewise the identity]\label{mk-lem:mapHC_augment_iso}
  \lean{AlgebraicGeometry.mapHC_augment_iso}
  \leanok
  \dcref{custom}

  \uses{mk-lem:map_augment_cond}
  Let \(\Phi : \mathcal{A} \to \mathcal{B}\) be an additive functor between preadditive
  categories, \(C^\bullet\) a cochain complex in \(\mathcal{A}\), and \(f : Y \to C^0\) an
  augmentation map satisfying \(f \cdot d^0 = 0\). Then applying \(\Phi\) degreewise commutes
  with the augmentation construction: there is an isomorphism of cochain complexes
  \[
    \Phi\bigl(C^\bullet.\operatorname{augment} f\bigr)
      \;\cong\;
    \bigl(\Phi C^\bullet\bigr).\operatorname{augment}\bigl(\Phi f\bigr),
  \]
  every component of which is the identity. (Here \(\Phi f\) satisfies its own augmentation
  identity \(\Phi(f) \cdot d^0 = 0\) by Lemma~\ref{mk-lem:map_augment_cond}.)
\end{lemma}
\begin{proof}
  \uses{mk-lem:map_augment_cond}
  The two augmented complexes have, by construction, the same underlying graded object:
  in degree \(0\) both read \(\Phi Y\), and in degree \(n+1\) both read \(\Phi C^n\). Define
  the isomorphism by taking every degreewise component to be the identity \(\mathrm{id}\).
  The differential-compatibility square in degree \(0\) and in each
  degree \(n+1\) then closes by unfolding the definition of \(\operatorname{augment}\): the
  augmented differentials of \(\Phi(C^\bullet.\operatorname{augment} f)\) and of
  \((\Phi C^\bullet).\operatorname{augment}(\Phi f)\) are literally the same map
  (\(\Phi f\) in the bottom slot, \(\Phi(d^{n})\) above), so each square commutes on the nose.
\end{proof}

\begin{lemma}[Additive functors preserve the augmentation condition]\label{mk-lem:map_augment_cond}
  \lean{AlgebraicGeometry.map_augment_cond}
  \leanok
  \dcref{custom}

  Let \(\Phi : \mathcal{A} \to \mathcal{B}\) be an additive functor between preadditive
  categories, \(C^\bullet\) a cochain complex in \(\mathcal{A}\), and \(f : Y \to C^0\) a map
  with \(f \cdot d^0 = 0\). Then the image map \(\Phi f : \Phi Y \to \Phi C^0\) again satisfies
  the augmentation condition \(\Phi(f) \cdot d^0 = 0\), where \(d^0\) is the lowest differential
  of the degreewise image complex \(\Phi C^\bullet\).
\end{lemma}
\begin{proof}

  The lowest differential of \(\Phi C^\bullet\) is, by definitional reshaping, \(\Phi(d^0)\)
  applied to the original lowest differential of \(C^\bullet\). Hence
  \(\Phi(f) \cdot \Phi(d^0) = \Phi(f \cdot d^0) = \Phi(0) = 0\) by functoriality of \(\Phi\)
  on composites and its preservation of zero.
\end{proof}

\begin{lemma}[Gluing an iso of augmented complexes from base, node, and compatibility data]\label{mk-lem:augmentCochainIso}
  \lean{AlgebraicGeometry.augmentCochainIso}
  \leanok
  \dcref{custom}

  Let \(C^\bullet, C'^\bullet\) be cochain complexes in a preadditive category, with augmentation
  maps \(f : Y \to C^0\) (\(f \cdot d^0 = 0\)) and \(f' : Y' \to C'^0\) (\(f' \cdot d'^0 = 0\)).
  Given an isomorphism \(\varphi : C^\bullet \cong C'^\bullet\) of the base complexes, an
  isomorphism \(e_Y : Y \cong Y'\) of the augmentation objects, and a compatibility square
  \(f \cdot \varphi_0 = e_Y \cdot f'\) intertwining the augmentation maps through \(e_Y\) and the
  degree-\(0\) component \(\varphi_0\) of \(\varphi\), there is an isomorphism of augmented
  cochain complexes
  \[
    C^\bullet.\operatorname{augment} f \;\cong\; C'^\bullet.\operatorname{augment} f'.
  \]
\end{lemma}
\begin{proof}

  Define the isomorphism degreewise: take the degree-\(0\) component to be \(e_Y\) and the
  degree-\(n+1\) component to be \(\varphi_n\), the \(n\)-th component of \(\varphi\). The
  degree-\(0\) differential square is exactly the supplied compatibility square
  \(f \cdot \varphi_0 = e_Y \cdot f'\). For each degree \(n+1\), the augmented differential is the
  original differential of the base complex (the augmentation only modifies the bottom slot), so
  the square reduces to the commutativity condition of \(\varphi\) at degree \(n\).
  Hence all squares commute and the components assemble to an isomorphism.
\end{proof}

\begin{lemma}[Open-meet identity: intersection with \(V\) distributes over the cover meet]\label{mk-lem:coverInterOpen_inf_distrib}
  \lean{AlgebraicGeometry.coverInterOpen_inf_eq_iInf_inf}
  \leanok
  \dcref{custom}

  \uses{mk-def:cech_free_presheaf_complex}
  Let \(\mathcal{U}\) be an open cover of \(X\), let \(V \subseteq X\) be an open, and let
  \(\sigma : \mathrm{Fin}(p+1) \to I\) be a multi-index. Then, as opens of \(X\),
  \[
    \operatorname{coverInterOpen}\mathcal{U}\,\sigma \cap V
      \;=\;
    \bigcap_{k}\bigl(\operatorname{coverOpen}\mathcal{U}\,(\sigma k) \cap V\bigr).
  \]
\end{lemma}
\begin{proof}

  By definition \(\operatorname{coverInterOpen}\mathcal{U}\,\sigma = \bigcap_k
  \operatorname{coverOpen}\mathcal{U}\,(\sigma k)\), the \((p+1)\)-fold meet of the
  cover-member opens. The index type \(\mathrm{Fin}(p+1)\) is nonempty, and in the frame of
  opens a binary meet distributes over a \emph{nonempty} indexed meet:
  \(\bigl(\bigcap_k A_k\bigr) \cap V = \bigcap_k (A_k \cap V)\). (Nonemptiness is essential:
  over the empty index the left side is \(\top \cap V = V\) while the right side is \(\top\).)
  Taking \(A_k = \operatorname{coverOpen}\mathcal{U}\,(\sigma k)\) gives the claim.
\end{proof}

\begin{lemma}[Degreewise object iso of the non-augmented core comparison]\label{mk-lem:coreIso_obj_iso}
  \lean{AlgebraicGeometry.coreIso_objIso}
  \leanok
  \dcref{custom}

  \uses{mk-lem:pushPull_eval_prod_iso, mk-lem:coverInterOpen_inf_distrib}
  Fix an open \(V \subseteq X\) and a degree \(p\). Write \(\Psi = \operatorname{forget}\cdot
  \operatorname{restrictScalars}(\mathrm{id})\) for the push--pull adapter and \(G_V =
  \operatorname{toPresheaf}\cdot\operatorname{ev}_V\) for the section-at-\(V\) functor. Then
  the degree-\(p\) terms of the evaluated non-augmented backbone complex and of the concrete
  restricted section {\v C}ech complex are isomorphic as abelian groups:
  \[
    \bigl((G_V\circ\Psi)\,\check{\mathcal{C}}^\bullet(\mathcal{U},\mathcal{F})\bigr).X\,p
      \;\cong\;
    \bigl(\check{\mathcal{C}}^\bullet(\mathcal{U}',\mathcal{F})\bigr).X\,p,
  \]
  i.e.\ \(\Gamma\bigl(V, \operatorname{pushPullObj}\mathcal{F}\,Y_p\bigr) \cong
  \prod_{\sigma : \mathrm{Fin}(p+1)\to I}
  \Gamma\bigl(\bigcap_k(\operatorname{coverOpen}\mathcal{U}\,(\sigma k)\cap V),\,
  \mathcal{F}\bigr)\), where \(\mathcal{U}'\) is the restricted family \(U'_i =
  \operatorname{coverOpen}\mathcal{U}\,i \cap V\).
\end{lemma}
\begin{proof}
  \uses{mk-lem:pushPull_eval_prod_iso, mk-lem:coverInterOpen_inf_distrib}
  Lemma~\ref{mk-lem:pushPull_eval_prod_iso} supplies the evaluated push--pull product iso
  \(\Gamma(V, \operatorname{pushPullObj}\mathcal{F}\,Y_p) \cong \prod_\sigma
  \Gamma(\operatorname{coverInterOpen}\mathcal{U}\,\sigma \cap V, \mathcal{F})\). Post-compose
  with the product reindexing \(\prod_\sigma \Gamma(\cdot, \mathcal{F})\) along the open-meet
  identity \(\operatorname{coverInterOpen}\mathcal{U}\,\sigma \cap V = \bigcap_k
  (\operatorname{coverOpen}\mathcal{U}\,(\sigma k)\cap V)\) of
  Lemma~\ref{mk-lem:coverInterOpen_inf_distrib} — factorwise the \(\operatorname{eqToIso}\) on the
  section group \(\Gamma(\cdot, \mathcal{F})\) obtained by applying \(\Gamma(\cdot, \mathcal{F})\)
  to that equality of opens. The composite lands on the degree-\(p\) term \(\prod_\sigma
  \Gamma(\bigcap_k(\operatorname{coverOpen}\mathcal{U}\,(\sigma k)\cap V), \mathcal{F})\) of
  \(\check{\mathcal{C}}^\bullet(\mathcal{U}', \mathcal{F})\), since for the restricted family
  \(U'_{\sigma k} = \operatorname{coverOpen}\mathcal{U}\,(\sigma k)\cap V\) the section
  {\v C}ech term at \(\sigma\) is exactly \(\Gamma(\bigcap_k U'_{\sigma k}, \mathcal{F})\).
\end{proof}

\begin{lemma}
[Per-leg coherence of the push--pull map of an open-immersion over-morphism]
  \label{mk-lem:pushPull_leg_coherence}
  \lean{AlgebraicGeometry.pushPull_leg_coherence, AlgebraicGeometry.pushPullLegIso}
  \dcref{custom}

  \uses{mk-def:push_pull_map, mk-lem:restrictFunctorIsoPullback_mathlib,
    mk-lem:pullbackPushforwardAdjunction_mathlib}
  Let \(q : A \to X\) be a morphism of schemes, let \(c : C' \hookrightarrow A\) be an open
  immersion, and let \(pC : C' \to X\) satisfy \(c \cdot q = pC\). For a sheaf of
  \(\mathcal{O}_X\)-modules \(\mathcal{F}\), the push--pull map of the over-\(X\) morphism
  \(\operatorname{Over.mk} c : \operatorname{Over.mk} pC \to \operatorname{Over.mk} q\) factors,
  through a canonical leg isomorphism, as the pushforward of the restriction-adjunction unit:
  \[
    \operatorname{pushPullMap}\mathcal{F}\,(\operatorname{Over.mk} c)
      \;=\;
    q_*\bigl(\eta^c_{\,q^*\mathcal{F}}\bigr)\cdot
      \bigl(\operatorname{pushPullLegIso} q\,c\,pC\,\mathcal{F}\bigr).\mathrm{hom},
  \]
  where \(\eta^c\) is the unit of the restriction adjunction along the open immersion \(c\),
  evaluated at the pullback \(q^*\mathcal{F}\), and \(\operatorname{pushPullLegIso}\) is the
  canonical leg isomorphism assembled from the restriction-versus-pullback comparison
  isomorphism \(\operatorname{restrictFunctorIsoPullback}\)
  (Lemma~\ref{mk-lem:restrictFunctorIsoPullback_mathlib}), the pullback-composition comparison
  \(\operatorname{pullbackComp}\) for the composite \(c \cdot q\), and the reindexing
  \(\operatorname{pullbackCongr}\) along the defining equality \(c \cdot q = pC\).
\end{lemma}
\begin{proof}
  \uses{mk-lem:restrictFunctorIsoPullback_mathlib, mk-lem:pullbackPushforwardAdjunction_mathlib}
  Unfold the raw description of the push--pull map of an open-immersion over-morphism. The key
  mate-calculus identity is that the restriction-adjunction unit \(\eta^c\), post-composed with the
  pushforward of the comparison isomorphism \(\operatorname{restrictFunctorIsoPullback}\), coincides
  with the unit of the pullback--pushforward adjunction along \(c\)
  (Lemma~\ref{mk-lem:pullbackPushforwardAdjunction_mathlib}); this is the general uniqueness of a left
  adjoint applied to the two presentations of pullback as restriction and as the genuine inverse
  image. Substituting the defining equality \(c \cdot q = pC\) collapses the reindexing comparison
  \(\operatorname{pullbackCongr}\) to the identity, and the surviving comparison isomorphisms compose
  to exactly the leg isomorphism \(\operatorname{pushPullLegIso}\). Hence the push--pull map equals
  the pushforward of the restriction unit followed by that leg isomorphism.
\end{proof}

\begin{lemma}
[Per-leg restriction naturality of the push--pull face]
  \label{mk-lem:pushPull_interLegHom_sections}
  \lean{AlgebraicGeometry.pushPull_interLegHom_sections,
    AlgebraicGeometry.pushPull_toRestrict_comm, AlgebraicGeometry.pullbackComp_rFIP_compat,
    AlgebraicGeometry.inner_beta_chain, AlgebraicGeometry.restrict_unit_comp,
    AlgebraicGeometry.unit_pushforward_rFIP_inv, AlgebraicGeometry.thin_resid5,
    AlgebraicGeometry.pls_eq}
  \dcref{custom}

  \uses{mk-lem:pushPull_leg_sections, mk-lem:pushPull_leg_coherence}
  Fix the finite cover \(\mathcal{U}\), a sheaf of \(\mathcal{O}_X\)-modules \(\mathcal{F}\), an
  open \(V \subseteq X\), a degree \(p\), a multi-index \(\sigma' : \mathrm{Fin}(p+2) \to I\), and a
  coface index \(k\). Write \(U_{\sigma'} = \operatorname{coverInterOpen}\mathcal{U}\,\sigma'\) and
  \(j_{\sigma'} : U_{\sigma'} \hookrightarrow X\) for the intersection-open immersion. Let
  \(\operatorname{interLegHom}\mathcal{U}\,\sigma'\,k : \operatorname{Over.mk} j_{\sigma'} \to
  \operatorname{Over.mk} j_{\sigma' \circ d_k}\) be the over-\(X\) morphism whose underlying map is
  the open immersion \(c : U_{\sigma'} \hookrightarrow U_{\sigma' \circ d_k}\) of intersection opens
  (deleting the \(k\)-th index can only shrink the meet). Through the per-leg section
  identifications \(\operatorname{pls} = \operatorname{pushPull\_leg\_sections}\)
  (Lemma~\ref{mk-lem:pushPull_leg_sections}) at source and target, the section over \(V\) of the
  push--pull map of \(\operatorname{interLegHom}\) equals the plain \(\mathcal{F}\)-restriction along
  \(U_{\sigma'} \cap V \subseteq U_{\sigma' \circ d_k} \cap V\):
  \[
    G_V\bigl(\Psi(\operatorname{interLegHom}\mathcal{U}\,\sigma'\,k)\bigr)\cdot
      \operatorname{pls}(\sigma').\mathrm{hom}
      \;=\;
    \operatorname{pls}(\sigma' \circ d_k).\mathrm{hom}\cdot
      \mathcal{F}\bigl(\operatorname{homOfLE}\,(U_{\sigma'} \cap V \subseteq
        U_{\sigma' \circ d_k} \cap V)\bigr).
  \]
\end{lemma}
\begin{proof}
  \uses{mk-lem:pushPull_leg_sections, mk-lem:pushPull_leg_coherence}
  The argument has four steps.

  \emph{(a) Per-leg coherence.} By Lemma~\ref{mk-lem:pushPull_leg_coherence}, the push--pull map of an
  over-morphism \(\operatorname{Over.mk} c\) with \(c\) an open immersion is, through the canonical
  leg isomorphism, the pushforward of the restriction-adjunction unit \(\eta^c\). Thus the only
  nontrivial content of \(\operatorname{pushPullMap}\mathcal{F}\,(\operatorname{interLegHom})\) is the
  restriction unit along the geometric face \(c\); the remaining factors are canonical comparison
  isomorphisms.

  \emph{(b) Unit versus pullback comparison.} The restriction-adjunction unit, composed with the
  pushforward of the comparison isomorphism \(\operatorname{restrictFunctorIsoPullback}\) (the
  uniqueness isomorphism between the restriction adjunction and the genuine pullback--pushforward
  adjunction), is itself the unit of the pullback--pushforward adjunction. This replaces the
  restriction unit by the inverse-image unit without changing the map.

  \emph{(c) Composition law along \(c\) and \(q\).} The unit of the composite restriction/pullback
  factors through the units of the two pieces via the comparison isomorphisms
  \(\operatorname{pullbackComp}\) and the analogous restriction-composition comparison: their
  naturality squares chain the inner comparison along the defining equality \(c \cdot q = pC\). This
  identifies \(\operatorname{pushPullMap}\mathcal{F}\,(\operatorname{interLegHom})\) post-composed
  with the inverse of the target leg isomorphism with the pushforward of the restriction unit,
  transported by the reindexing isomorphisms at \(c \cdot q = pC\).

  \emph{(d) On sections over \(V\).} At the level of sections the restriction-adjunction unit acts
  as the plain presheaf restriction \(\mathcal{F}(\operatorname{homOfLE}\,\cdots)\) — the value of
  \(\mathcal{F}\) on the inclusion of opens induced by \(c\). The leg isomorphism
  \(\operatorname{pls}.\mathrm{hom}\) reads, on sections, as the section-over-\(V\) of the
  pushforward of the comparison isomorphism followed by a canonical equality isomorphism. Composing the
  contributions of (a)--(c), every comparison and transport factor is a morphism between section
  groups indexed by opens; since the lattice of opens is a thin category, all such reindexing
  morphisms are uniquely determined and compose to a single restriction. The whole composite
  therefore collapses to the single \(\mathcal{F}\)-restriction along
  \(U_{\sigma'} \cap V \subseteq U_{\sigma' \circ d_k} \cap V\), which is exactly the claimed
  \(\operatorname{homOfLE}\) map.
\end{proof}

\begin{lemma}
[Backbone inclusion projects to the canonical component map]
  \label{mk-lem:backboneIncl_proj}
  \lean{AlgebraicGeometry.backboneIncl_proj, AlgebraicGeometry.backboneIncl,
    AlgebraicGeometry.backboneProj, AlgebraicGeometry.interProj,
    AlgebraicGeometry.coverInterToMember, AlgebraicGeometry.coverInterToMember_fac,
    AlgebraicGeometry.over_hom_ext_mono, AlgebraicGeometry.overSigmaHomEq,
    AlgebraicGeometry.ι_whiskerEquiv_inv, AlgebraicGeometry.cf_hom_snd,
    AlgebraicGeometry.ι_cf_hom, AlgebraicGeometry.ι_overProd_distrib_inv,
    AlgebraicGeometry.ι_overProd_distrib_right_inv, AlgebraicGeometry.overWPproj,
    AlgebraicGeometry.overDescIncl, AlgebraicGeometry.ι_overSigmaDescIso_hom,
    AlgebraicGeometry.wpEqProd_inv_overWPproj, AlgebraicGeometry.wpEqProd_hom_π,
    AlgebraicGeometry.prodFinSuccIso_hom_fst, AlgebraicGeometry.prodFinSuccIso_hom_snd_π,
    AlgebraicGeometry.prodFinSuccIso_inv_π_succ, AlgebraicGeometry.prodFinSuccIso_inv_π_zero,
    AlgebraicGeometry.coverArrowIncl, AlgebraicGeometry.coverReindexHom,
    AlgebraicGeometry.entry_chain, AlgebraicGeometry.glue_chain,
    AlgebraicGeometry.ι_fibrePower_reindex_inv, AlgebraicGeometry.ι_reindex_wpci_inv_overWPproj,
    AlgebraicGeometry.ι_sigmaMapIso_inv, AlgebraicGeometry.ι_wpci_inv_overWPproj}
  \dcref{custom}

  \uses{mk-lem:cech_backbone_left_sigma, mk-lem:pushPull_sigma_iso, mk-def:cover_cech_nerve}
  Fix a degree \(p\) and a multi-index \(\tau : \mathrm{Fin}(p+1) \to I\). The backbone inclusion
  \(\operatorname{backboneIncl}\mathcal{U}\,p\,\tau : \operatorname{Over.mk} j_\tau \to Y_p\) of the
  \(\tau\)-summand, followed by the \(l\)-th wide-pullback projection
  \(\operatorname{backboneProj}\mathcal{U}\,p\,l\), equals the canonical component map
  \(\operatorname{interProj}\mathcal{U}\,\tau\,l\) — the over-\(X\) inclusion of the intersection
  open \(U_\tau\) into the \((\tau l)\)-th cover member:
  \[
    \operatorname{backboneIncl}\mathcal{U}\,p\,\tau\cdot
      \operatorname{backboneProj}\mathcal{U}\,p\,l
      \;=\;
    \operatorname{interProj}\mathcal{U}\,\tau\,l.
  \]
\end{lemma}
\begin{proof}
  \uses{mk-lem:cech_backbone_left_sigma}
  The backbone inclusion is assembled, through the coproduct decomposition of the backbone
  (Lemma~\ref{mk-lem:cech_backbone_left_sigma}), from several canonical comparison isomorphisms: the
  distributivity isomorphisms reassociating the iterated fibre product, a canonical summand
  reindexing, and the projection of the wide-pullback comparison. Composing
  with the \(l\)-th projection and reassociating these layers reduces the left-hand side to a single
  lift into the \((\tau l)\)-th cover member \(\mathcal{U}.f(\tau l)\). Because that member arrow is
  an open immersion, hence a monomorphism, the lift is rigid: any two over-\(X\) maps agreeing after
  post-composition with it are equal. The canonical component map \(\operatorname{interProj}\) is,
  by its defining factorization (\(\operatorname{coverInterToMember\_fac}\)), exactly such a lift, so
  the two composites coincide.
\end{proof}

\begin{lemma}
[Backbone inclusion intertwines the nerve coface with the geometric face]
  \label{mk-lem:backboneIncl_nerveδ}
  \lean{AlgebraicGeometry.backboneIncl_nerveδ, AlgebraicGeometry.backbone_hom_ext,
    AlgebraicGeometry.nerveδ_backboneProj, AlgebraicGeometry.cechNerve_drop_δ,
    AlgebraicGeometry.interLegHom, AlgebraicGeometry.interLegHom_interProj}
  \dcref{custom}

  \uses{mk-lem:backboneIncl_proj, mk-def:cover_cech_nerve}
  For a multi-index \(\sigma' : \mathrm{Fin}(p+2) \to I\) and a coface index \(k\), the backbone
  inclusion intertwines the \(k\)-th {\v C}ech-nerve coface \(\delta^{\mathrm{nerve}}_k\) with the
  geometric face \(\operatorname{interLegHom}\):
  \[
    \operatorname{backboneIncl}\mathcal{U}\,(p+1)\,\sigma'\cdot \delta^{\mathrm{nerve}}_k
      \;=\;
    \operatorname{interLegHom}\mathcal{U}\,\sigma'\,k\cdot
      \operatorname{backboneIncl}\mathcal{U}\,p\,(\sigma' \circ d_k).
  \]
\end{lemma}
\begin{proof}
  \uses{mk-lem:backboneIncl_proj}
  Both sides are over-\(X\) maps into the degree-\(p\) backbone term \(Y_p\). By the wide-pullback
  extensionality principle \(\operatorname{backbone\_hom\_ext}\), it suffices to check equality after
  post-composition with each component projection \(\operatorname{backboneProj}\mathcal{U}\,p\,l\).
  On the left, reassociating and applying \(\operatorname{nerveδ\_backboneProj}\) (the nerve coface
  followed by a projection is again the projection at the dropped index) and then
  Lemma~\ref{mk-lem:backboneIncl_proj} yields the canonical component map
  \(\operatorname{interProj}\mathcal{U}\,\sigma'\,(d_k\,l)\). On the right, applying
  Lemma~\ref{mk-lem:backboneIncl_proj} and then \(\operatorname{interLegHom\_interProj}\) (the geometric
  face intertwines the canonical component maps) yields the same map. Hence the two composites agree
  after every projection, so they are equal.
\end{proof}

\begin{lemma}
[Coordinate formula for the degreewise object iso]
  \label{mk-lem:coreIso_objIso_π}
  \lean{AlgebraicGeometry.coreIso_objIso_π, AlgebraicGeometry.GVΨ_map_eq,
    AlgebraicGeometry.sectionFunctorV, AlgebraicGeometry.piMapIso_hom_π,
    AlgebraicGeometry.pushPull_sigma_iso_π_incl}
  \dcref{custom}

  \uses{mk-lem:coreIso_obj_iso, mk-lem:pushPull_eval_prod_iso, mk-lem:coverInterOpen_inf_distrib,
    mk-lem:pushPull_leg_sections, mk-lem:pushPull_sigma_iso}
  With the notation of Lemma~\ref{mk-lem:coreIso_obj_iso}, fix a degree \(p\), an open \(V \subseteq X\),
  and a multi-index \(\tau : \mathrm{Fin}(p+1) \to I\). The \(\tau\)-projection of the degreewise
  object isomorphism \((\mathrm{objIso}\,p).\mathrm{hom}\) is the evaluated push--pull map of the
  backbone inclusion \(\operatorname{backboneIncl}\mathcal{U}\,p\,\tau\), followed by the per-leg
  section identification \(\operatorname{pls}(\tau)\) (Lemma~\ref{mk-lem:pushPull_leg_sections}) and the
  open-meet reindex of Lemma~\ref{mk-lem:coverInterOpen_inf_distrib}:
  \[
    (\mathrm{objIso}\,p).\mathrm{hom}\cdot \pi_\tau
      \;=\;
    G_V\bigl(\operatorname{pushPullMap}\mathcal{F}\,
        (\operatorname{backboneIncl}\mathcal{U}\,p\,\tau)\bigr)\cdot
      \operatorname{pls}(\tau).\mathrm{hom}\cdot
      \bigl(\text{open-meet reindex on } \Gamma(\cdot,\mathcal{F})\bigr).
  \]
\end{lemma}
\begin{proof}
  \uses{mk-lem:coreIso_obj_iso, mk-lem:pushPull_eval_prod_iso, mk-lem:coverInterOpen_inf_distrib,
    mk-lem:pushPull_sigma_iso}
  Unfold \((\mathrm{objIso}\,p).\mathrm{hom}\) as in Lemma~\ref{mk-lem:coreIso_obj_iso}: it is the
  evaluated push--pull product isomorphism of Lemma~\ref{mk-lem:pushPull_eval_prod_iso}, reindexed
  factorwise by the open-meet identity of Lemma~\ref{mk-lem:coverInterOpen_inf_distrib}. Projecting
  onto the \(\tau\)-factor, the product comparison of the section functor composed with the
  \(\tau\)-projection of the product isomorphism is, by the leg of the coproduct-to-product
  decomposition (Lemma~\ref{mk-lem:pushPull_sigma_iso}), the push--pull map of the backbone inclusion
  of the \(\tau\)-summand. Reading off the three surviving layers term by term: the
  product-decomposition leg supplies \(G_V(\operatorname{pushPullMap}\mathcal{F}\,
  (\operatorname{backboneIncl}\mathcal{U}\,p\,\tau))\), the per-leg section identification supplies
  \(\operatorname{pls}(\tau).\mathrm{hom}\), and the factorwise object-equality isomorphism supplies
  the open-meet reindex. Composing these gives the asserted coordinate formula.
\end{proof}

\begin{lemma}
[Per-leg naturality of the core comparison coface]
  \label{mk-lem:coreIso_comm_leg}
  \lean{AlgebraicGeometry.coreIso_comm_leg, AlgebraicGeometry.map_op_eqToHom_swap}
  \dcref{custom}

  \uses{mk-lem:coreIso_obj_iso, mk-lem:pushPull_eval_prod_iso, mk-lem:pushPull_sigma_iso,
    mk-lem:pushPull_leg_sections, mk-lem:section_cech_product_equiv,
    mk-lem:coverInterOpen_inf_distrib}
  With the notation of Lemma~\ref{mk-lem:coreIso_obj_iso}, fix a degree \(p\), a coface index
  \(k \le p+1\), and a multi-index \(\sigma' : \mathrm{Fin}(p+2) \to I\). Read the degree-\(p\)
  and degree-\((p+1)\) terms of both complexes through the product equivalence
  \(\operatorname{sectionCechProductEquiv}\) (Lemma~\ref{mk-lem:section_cech_product_equiv}). Then
  the \(\sigma'\)-coordinate of the evaluated push--pull coface followed by the object iso,
  \(G_V\bigl(\Psi(\delta^{\mathrm{nerve}}_k)\bigr) \cdot (\mathrm{objIso}\,(p+1)).\mathrm{hom}\),
  equals the presheaf face restriction \(\operatorname{sectionCechFaceRestr}(\sigma', k)\) — the
  value of \(\mathcal{F}\) on the inclusion of intersection opens
  \[
    \bigcap_{j}\bigl(\operatorname{coverOpen}\mathcal{U}\,(\sigma' j) \cap V\bigr)
      \;\subseteq\;
    \bigcap_{l}\bigl(\operatorname{coverOpen}\mathcal{U}\,((\sigma' \circ d_k) l) \cap V\bigr)
  \]
  — applied to the \((\sigma' \circ d_k)\)-coordinate of the image of the source under
  \((\mathrm{objIso}\,p).\mathrm{hom}\). In words: the projection of the evaluated push--pull
  coface onto the \(\sigma'\)-component is exactly the \(\mathcal{F}\)-restriction along the
  \(k\)-th face inclusion.
\end{lemma}
\begin{proof}
  \uses{mk-lem:pushPull_eval_prod_iso, mk-lem:pushPull_sigma_iso, mk-lem:pushPull_leg_sections,
    mk-lem:coverInterOpen_inf_distrib, mk-lem:pushPull_interLegHom_sections, mk-lem:backboneIncl_proj,
    mk-lem:backboneIncl_nerveδ, mk-lem:coreIso_objIso_π}
  Unwind \((\mathrm{objIso}\,(p+1)).\mathrm{hom}\) through its construction in
  Lemma~\ref{mk-lem:coreIso_obj_iso}: it is the evaluated push--pull product iso
  \(\operatorname{pushPull\_eval\_prod\_iso}\) (Lemma~\ref{mk-lem:pushPull_eval_prod_iso}),
  reindexed factorwise by the open-meet identity (Lemma~\ref{mk-lem:coverInterOpen_inf_distrib}).
  The product iso factors through the coproduct-to-product decomposition of the backbone
  push--pull object (Lemma~\ref{mk-lem:pushPull_sigma_iso}), the preservation of products by the
  section functor, and the per-leg section identification
  \(\Gamma(V, \operatorname{pushPullObj}\mathcal{F}\,(\operatorname{Over.mk} j_{\sigma'})) \cong
  \Gamma(U_{\sigma'} \cap V, \mathcal{F})\) (Lemma~\ref{mk-lem:pushPull_leg_sections}). Projecting
  onto the \(\sigma'\)-leg, the evaluated {\v C}ech-nerve coface \(\delta^{\mathrm{nerve}}_k\)
  acts on sections as the restriction of \(\mathcal{F}\) along the inclusion of the
  \(\sigma'\)-intersection open into the \((\sigma' \circ d_k)\)-intersection open — precisely the
  coface induced by deleting the \(k\)-th index. The open-meet identity
  (Lemma~\ref{mk-lem:coverInterOpen_inf_distrib}) rewrites both intersection opens as the meets
  \(\bigcap_j(\operatorname{coverOpen}\mathcal{U}\,(\sigma' j) \cap V)\) and
  \(\bigcap_l(\operatorname{coverOpen}\mathcal{U}\,((\sigma' \circ d_k) l) \cap V)\), identifying
  this restriction with \(\operatorname{sectionCechFaceRestr}(\sigma', k)\) on the nose.
\end{proof}

\begin{lemma}[Per-coface square of the core comparison]\label{mk-lem:coreIso_comm_coface}
  \lean{AlgebraicGeometry.coreIso_comm_coface}
  \leanok
  \dcref{custom}

  \uses{mk-lem:coreIso_comm_leg, mk-lem:section_cech_product_equiv, mk-lem:section_cech_objd_apply}
  With the notation of Lemma~\ref{mk-lem:coreIso_obj_iso}, for each degree \(p\) and each coface
  index \(k \le p+1\) the individual cofaces are intertwined by the object isomorphisms:
  \[
    (\mathrm{objIso}\,p).\mathrm{hom} \cdot \delta^{\mathrm{sec}}_k
      \;=\;
    G_V\bigl(\Psi(\delta^{\mathrm{nerve}}_k)\bigr) \cdot (\mathrm{objIso}\,(p+1)).\mathrm{hom},
  \]
  where \(\delta^{\mathrm{sec}}_k\) is the \(k\)-th section-{\v C}ech coface and
  \(\delta^{\mathrm{nerve}}_k\) the \(k\)-th {\v C}ech-nerve coface of the backbone, both maps
  into the degree-\((p+1)\) term.
\end{lemma}
\begin{proof}
  \leanok
  \uses{mk-lem:coreIso_comm_leg, mk-lem:section_cech_product_equiv, mk-lem:section_cech_objd_apply}
  Both sides are maps into the degree-\((p+1)\) term, which is the product over multi-indices
  \(\sigma' : \mathrm{Fin}(p+2) \to I\); by the product equivalence
  \(\operatorname{sectionCechProductEquiv}\) (Lemma~\ref{mk-lem:section_cech_product_equiv}) it
  suffices to check equality coordinate by coordinate. Fix \(\sigma'\). The \(\sigma'\)-coordinate
  of the left-hand side is, by the per-coface coordinate computation underlying
  Lemma~\ref{mk-lem:section_cech_objd_apply}, the presheaf face restriction
  \(\operatorname{sectionCechFaceRestr}(\sigma', k)\) applied to the \((\sigma' \circ d_k)\)-
  coordinate of \((\mathrm{objIso}\,p).\mathrm{hom}\). The \(\sigma'\)-coordinate of the
  right-hand side is the same expression by Lemma~\ref{mk-lem:coreIso_comm_leg}. Hence all
  coordinates agree and the two cofaces are equal.
\end{proof}

\begin{lemma}
[Alternating-sum assembly of the core comparison square]
  \label{mk-lem:coreIso_comm_sum}
  \lean{AlgebraicGeometry.coreIso_comm_sum, AlgebraicGeometry.abHom_finsetSum_apply}
  \dcref{custom}

  \uses{mk-lem:coreIso_comm_coface, mk-lem:section_cech_objd_apply, mk-lem:section_cech_product_equiv}
  With the notation of Lemma~\ref{mk-lem:coreIso_obj_iso}, for each degree \(p\) the full
  alternating-coface differentials are intertwined by the object isomorphisms:
  \[
    (\mathrm{objIso}\,p).\mathrm{hom} \cdot d^p_{\check{\mathcal{C}}(\mathcal{U}')}
      \;=\;
    d^p_{(G_V\circ\Psi)\check{\mathcal{C}}(\mathcal{U})} \cdot (\mathrm{objIso}\,(p+1)).\mathrm{hom}.
  \]
\end{lemma}
\begin{proof}
  \uses{mk-lem:coreIso_comm_coface, mk-lem:section_cech_objd_apply, mk-lem:section_cech_product_equiv}
  Both differentials are the alternating coface sums \(d^p = \sum_{k=0}^{p+1} (-1)^k
  \delta_k\) (the \(\operatorname{objD}\) of the alternating-coface-map complex). Work
  elementwise: by Lemma~\ref{mk-lem:section_cech_objd_apply}, read through the product equivalence
  (Lemma~\ref{mk-lem:section_cech_product_equiv}), each side — evaluated at an element of the source
  and at a coordinate \(\sigma'\) — is the pointwise finite sum \(\sum_{k} (-1)^k\,(\cdots)_k\) of
  the same shape. The two sums agree summand by summand: the \(k\)-th summands carry the same sign
  \((-1)^k\) and, by the per-coface square of Lemma~\ref{mk-lem:coreIso_comm_coface}, the same coface
  contribution. A termwise comparison of the two alternating sums therefore yields equality of
  the differentials in every degree.
\end{proof}

\begin{lemma}[The core comparison intertwines the {\v C}ech differentials]\label{mk-lem:coreIso_comm}
  \lean{AlgebraicGeometry.coreIso_comm}
  \leanok
  \dcref{custom}

  \uses{mk-lem:coreIso_comm_sum, mk-lem:coreIso_obj_iso}
  With the notation of Lemma~\ref{mk-lem:coreIso_obj_iso}, for each \(p\) the degreewise object
  isomorphisms \(\mathrm{objIso}\,p\) intertwine the two differentials:
  \[
    (\mathrm{objIso}\,p).\mathrm{hom} \cdot
      d^{p}_{\check{\mathcal{C}}(\mathcal{U}')}
      \;=\;
    d^{p}_{(G_V\circ\Psi)\check{\mathcal{C}}(\mathcal{U})} \cdot
      (\mathrm{objIso}\,(p+1)).\mathrm{hom}.
  \]
  Consequently \(\mathrm{coreIso} := \operatorname{isoOfComponents}(\mathrm{objIso},\,\cdot)\)
  is an isomorphism of cochain complexes \((G_V\circ\Psi)\check{\mathcal{C}}^\bullet(\mathcal{U},
  \mathcal{F}) \cong \check{\mathcal{C}}^\bullet(\mathcal{U}', \mathcal{F})\).
\end{lemma}
\begin{proof}
  \leanok
  \uses{mk-lem:coreIso_comm_sum, mk-lem:coreIso_obj_iso}
  The intertwining square for each \(p\) is exactly the alternating-sum assembly of
  Lemma~\ref{mk-lem:coreIso_comm_sum}: both {\v C}ech differentials are the alternating coface sums
  \(\sum_k (-1)^k \delta_k\), and that lemma matches them by summing the per-coface squares of
  Lemma~\ref{mk-lem:coreIso_comm_coface} termwise. Holding for every \(p\), the degreewise object
  isomorphisms \(\mathrm{objIso}\,p\) together with these commuting squares assemble
  \(\mathrm{coreIso} := \operatorname{isoOfComponents}(\mathrm{objIso},\,\cdot)\) into an
  isomorphism of cochain complexes.

\end{proof}

\begin{definition}[Augmentation of the section {\v C}ech complex over \(V\)]\label{mk-def:sectionCechAugV}
  \lean{AlgebraicGeometry.sectionCechAugV}
  \leanok
  \dcref{custom}

  \uses{mk-def:section_cech_complex}
  Fix an open \(V \subseteq X\). The augmentation of the concrete section {\v C}ech complex for
  the restricted family \(U'_i=U_i\cap V\) is the product of the restriction maps
  \[
    \varepsilon_V : \Gamma(V, \mathcal{F}) \longrightarrow
      \prod_{\sigma:\operatorname{Fin}1\to I}
        \Gamma\!\left(\bigcap_k(U_{\sigma(k)}\cap V),\mathcal{F}\right),
    \qquad
    t\longmapsto
      \left(t\big|_{\cap_k(U_{\sigma(k)}\cap V)}\right)_{\sigma}.
  \]
\end{definition}

\begin{lemma}
[Coordinatewise value of the section augmentation]
  \label{mk-lem:sectionCechAugV_π}
  \lean{AlgebraicGeometry.sectionCechAugV_π}
  \leanok
  \dcref{custom}

  \uses{mk-def:sectionCechAugV}
  Fix an open \(V \subseteq X\) and a one-element multi-index \(\sigma : \mathrm{Fin}\,1 \to I\).
  The \(\sigma\)-coordinate of the augmentation \(\varepsilon_V\) of
  Definition~\ref{mk-def:sectionCechAugV}, namely its composite with the projection onto the
  \(\sigma\)-factor of
  \(\prod_\tau \Gamma\bigl(\bigcap_l(U_{\tau l}\cap V),\mathcal{F}\bigr)\), is the plain presheaf
  restriction of \(\mathcal{F}\) along the inclusion of intersection opens
  \(\bigcap_l(U_{\sigma l}\cap V) \subseteq V\):
  \[
    \operatorname{pr}_{\sigma}\circ\varepsilon_V
      \;=\;
    \mathcal{F}\bigl(\,\textstyle\bigcap_l(U_{\sigma l}\cap V)\subseteq V\,\bigr)
      \;:\; \Gamma(V,\mathcal{F}) \longrightarrow
      \Gamma\bigl(\textstyle\bigcap_l(U_{\sigma l}\cap V),\mathcal{F}\bigr).
  \]
\end{lemma}
\begin{proof}
  \leanok
  \uses{mk-def:sectionCechAugV}
  This is the defining property of a map into a product: composing the product of the
  restrictions with its \(\sigma\)-projection gives the restriction indexed by \(\sigma\).
\end{proof}

\begin{lemma}[The section augmentation kills the first differential]\label{mk-lem:sectionCechAugV_comp_d}
  \lean{AlgebraicGeometry.sectionCechAugV_comp_d}
  \leanok
  \dcref{custom}

  \uses{mk-def:sectionCechAugV, mk-lem:sectionCechAugV_π, mk-lem:section_cech_objd_apply}
  For the augmentation \(\varepsilon_V\) of Definition~\ref{mk-def:sectionCechAugV},
  \[
    \varepsilon_V \cdot d^0 \;=\; 0.
  \]
\end{lemma}
\begin{proof}
  \leanok
  \uses{mk-lem:sectionCechAugV_π, mk-lem:section_cech_objd_apply}
  It suffices to check every coordinate indexed by \(\sigma:\operatorname{Fin}2\to I\).
  The alternating differential has two terms. By
  Lemma~\ref{mk-lem:sectionCechAugV_π}, each is the composite restriction from \(V\), first to one
  of \(U_{\sigma(0)}\cap V\) and \(U_{\sigma(1)}\cap V\), and then to their intersection.
  Transitivity of restriction identifies both composites with the restriction from \(V\) to
  \((U_{\sigma(0)}\cap V)\cap(U_{\sigma(1)}\cap V)\). They occur with opposite signs, so their
  sum is zero.
\end{proof}

\begin{lemma}[Comparison of the evaluated and concrete section augmentations]\label{mk-lem:mappedSectionCechAugV_eq}
  \lean{AlgebraicGeometry.mappedSectionCechAugV_eq}
  \leanok
  \dcref{custom}

  \uses{mk-def:cech_augmentation, mk-def:sectionCechAugV, mk-lem:sectionCechAugV_π,
    mk-lem:coreIso_obj_iso, mk-lem:coreIso_objIso_π, mk-lem:pushPull_sigma_iso,
    mk-lem:pushPull_leg_sections, mk-lem:pushPull_eval_prod_iso}
  Let \(\Psi=\operatorname{forget}\cdot\operatorname{restrictScalars}(\mathrm{id})\), let
  \(G_V=\operatorname{toPresheaf}\cdot\operatorname{ev}_V\), and let
  \(\mathrm{objIso}\,0\) be the degree-zero isomorphism of
  Lemma~\ref{mk-lem:coreIso_obj_iso}. Then
  \[
    G_V\bigl(\Psi(\operatorname{cechAugmentation})\bigr)
      \cdot (\mathrm{objIso}\,0).\mathrm{hom}=\varepsilon_V.
  \]
\end{lemma}
\begin{proof}
  \leanok
  \uses{mk-lem:sectionCechAugV_π, mk-lem:coreIso_objIso_π, mk-lem:pushPull_sigma_iso,
    mk-lem:pushPull_leg_sections, mk-lem:pushPull_eval_prod_iso}
  Maps into the product are equal if all their coordinates are equal. Fix
  \(\sigma:\operatorname{Fin}1\to I\). The \(\sigma\)-projection of
  \(\mathrm{objIso}\,0\) identifies the corresponding push--pull term with sections over
  \(U_{\sigma(0)}\cap V\). The augmentation followed by this leg is the unit for the open
  immersion \(U_{\sigma(0)}\hookrightarrow X\); under the section identification, that unit is
  restriction from \(V\) to \(U_{\sigma(0)}\cap V\). This is the \(\sigma\)-coordinate of
  \(\varepsilon_V\) by Lemma~\ref{mk-lem:sectionCechAugV_π}.
\end{proof}

\begin{lemma}

[The evaluated augmented {\v C}ech complex is the concrete section {\v C}ech complex]
  \label{mk-lem:cechSection_complex_iso}
  \lean{AlgebraicGeometry.cechSection_complex_iso,
    AlgebraicGeometry.sectionCechComplexV}
  \dcref{custom}

  \uses{mk-lem:pushPull_eval_prod_iso, mk-lem:section_cech_objd_apply,
    mk-lem:section_cech_product_equiv, mk-def:cech_augmented_complex,
    mk-def:cech_augmentation, mk-lem:cech_augmentation_comp_d,
    mk-def:sectionCechAugV, mk-lem:sectionCechAugV_comp_d,
    mk-lem:mappedSectionCechAugV_eq, mk-lem:coreIso_obj_iso, mk-lem:coreIso_comm,
    mk-lem:coverInterOpen_inf_distrib,
    mk-lem:mapHC_augment_iso, mk-lem:map_augment_cond, mk-lem:augmentCochainIso}
  Fix an open \(V \subseteq X\). The augmented {\v C}ech complex of
  Definition~\ref{mk-def:cech_augmented_complex} evaluated at \(V\) (through
  \(\operatorname{toPresheaf}\)) --- call it \(D^\bullet\), the cochain complex of abelian
  groups whose augmentation node is \(\Gamma(V, \mathcal{F})\) and whose {\v C}ech degree-\(p\)
  term is \(\Gamma(V, \operatorname{pushPullObj}\mathcal{F}\,Y_p)\) --- is isomorphic, as a
  cochain complex, to the \emph{augmented} concrete section {\v C}ech complex
  \[
    D'^\bullet_{\mathrm{aug}} \;:=\;
    \bigl(\check{\mathcal{C}}^\bullet(\mathcal{U}', \mathcal{F})\bigr).\operatorname{augment}
      \varepsilon\, h_\varepsilon
  \]
  built from the concrete section {\v C}ech complex over the restricted family
  \(U'_\sigma = \operatorname{coverInterOpen}\mathcal{U}\,\sigma \cap V\) (degree-\(p\) term
  \(\prod_\sigma \Gamma(U_\sigma \cap V, \mathcal{F})\), alternating coface restriction
  differential) by prepending the augmentation node \(\Gamma(V, \mathcal{F})\) along
  \begin{itemize}
    \item the restriction product map \(\varepsilon : \Gamma(V, \mathcal{F}) \to
      \prod_i \Gamma(U_i \cap V, \mathcal{F})\), \(t \mapsto \bigl(t|_{U_i \cap V}\bigr)_i\)
      of Definition~\ref{mk-def:sectionCechAugV}; by
      Lemma~\ref{mk-lem:mappedSectionCechAugV_eq}, this is the evaluation at \(V\) of the {\v C}ech
      augmentation \(\mathcal{F} \to \mathcal{C}^0\) of Definition~\ref{mk-def:cech_augmentation}; and
    \item the augmentation identity \(h_\varepsilon : \varepsilon \cdot d^0 = 0\), the
      content of Lemma~\ref{mk-lem:sectionCechAugV_comp_d}.
  \end{itemize}
  Concretely: at the augmentation node both complexes carry \(\Gamma(V, \mathcal{F})\),
  identified by the identity map; at {\v C}ech degree \(p\) the object isomorphism
  \(\Gamma(V, \operatorname{pushPullObj}\mathcal{F}\,Y_p) \cong \prod_\sigma \Gamma(U_\sigma
  \cap V, \mathcal{F})\) is Lemma~\ref{mk-lem:pushPull_eval_prod_iso}; and these degreewise maps
  intertwine the two differentials, including the augmentation differential \(\varepsilon\).
\end{lemma}
\begin{proof}
  \uses{mk-lem:pushPull_eval_prod_iso, mk-lem:section_cech_objd_apply,
    mk-lem:section_cech_product_equiv, mk-def:cech_augmentation, mk-lem:cech_augmentation_comp_d,
    mk-lem:mapHC_augment_iso, mk-lem:map_augment_cond, mk-lem:augmentCochainIso,
    mk-lem:mappedSectionCechAugV_eq, mk-lem:sectionCechAugV_comp_d,
    mk-lem:coreIso_obj_iso, mk-lem:coreIso_comm, mk-lem:coverInterOpen_inf_distrib}
  The augmentation node of both complexes is \(\Gamma(V, \mathcal{F})\), matched by the
  identity. At {\v C}ech degree \(p\) the object isomorphism between
  \(\Gamma(V, \operatorname{pushPullObj}\mathcal{F}\,Y_p)\) and \(\prod_\sigma \Gamma(U_\sigma
  \cap V, \mathcal{F})\) is Lemma~\ref{mk-lem:pushPull_eval_prod_iso}. It remains to check the
  degreewise isomorphisms commute with the differentials.

  \emph{Augmentation bookkeeping (peeling the node).} Before matching differentials we strip the
  augmentation node off \(D^\bullet\). Writing \(\Psi = \operatorname{forget} \cdot
  \operatorname{restrictScalars}(\mathrm{id})\) for the push--pull adapter and \(G_V =
  \operatorname{toPresheaf} \cdot \operatorname{ev}_{V}\) for the section functor, the evaluated
  complex is \(D^\bullet = (G_V \circ \Psi)\bigl((\check{\mathcal{C}}^\bullet)
  .\operatorname{augment}\,(\operatorname{cechAugmentation})\bigr)\). Applying
  Lemma~\ref{mk-lem:mapHC_augment_iso} twice --- once for \(\Psi\) and once for \(G_V\), each an
  additive functor, with the augmentation condition transported by
  Lemma~\ref{mk-lem:map_augment_cond} --- commutes both functors past
  \(\operatorname{augment}\), so \(D^\bullet \cong \bigl((G_V \circ \Psi)
  \check{\mathcal{C}}^\bullet\bigr).\operatorname{augment}\,\bigl((G_V \circ \Psi)
  (\operatorname{cechAugmentation})\bigr)\) with identity components. Lemma~\ref{mk-lem:augmentCochainIso}
  then glues an iso of augmented complexes from three pieces: the \emph{non-augmented} core iso
  \[
    \mathrm{coreIso} : (G_V \circ \Psi)\bigl(\check{\mathcal{C}}^\bullet(\mathcal{U},
      \mathcal{F})\bigr) \;\cong\; \check{\mathcal{C}}^\bullet(\mathcal{U}', \mathcal{F}),
  \]
  the augmentation-object iso \(\Gamma(V, \mathcal{F}) \cong \Gamma(V, \mathcal{F})\) (the
  identity \(e_Y\)), and the degree-\(0\) compatibility square \(h_{\mathrm{compat}}\). This
  reduces the goal to the non-augmented \(\mathrm{coreIso}\), and the differential-match steps
  below apply to that reduced (non-augmented) goal --- not to the full augmented \(D^\bullet\),
  whose augmentation bookkeeping the three helper lemmas already discharge.

  \emph{The non-augmented core iso \(\mathrm{coreIso}\).} It is assembled componentwise from
  its degreewise object isos and a differential-match square. The degree-\(p\) object iso
  \(\mathrm{objIso}\,p : \Gamma(V,
  \operatorname{pushPullObj}\mathcal{F}\,Y_p) \cong \prod_\sigma \Gamma(U_\sigma \cap V,
  \mathcal{F})\) is Lemma~\ref{mk-lem:coreIso_obj_iso} (the evaluated push--pull product iso
  Lemma~\ref{mk-lem:pushPull_eval_prod_iso}, reindexed by the open-meet identity
  Lemma~\ref{mk-lem:coverInterOpen_inf_distrib}), and the differential-match square is
  Lemma~\ref{mk-lem:coreIso_comm}.

  \emph{{\v C}ech differentials (degrees \(p \geq 0\)).} This is the content of
  Lemma~\ref{mk-lem:coreIso_comm}: the differential of the section complex is, read through the
  product equivalence \(\operatorname{sectionCechProductEquiv}\)
  (Lemma~\ref{mk-lem:section_cech_product_equiv}), the alternating sum of presheaf face
  restrictions \(\sum_i (-1)^i\, \operatorname{sectionCechFaceRestr}(\sigma, i)\) computed in
  Lemma~\ref{mk-lem:section_cech_objd_apply}. The differential of \(D^\bullet\) is the
  evaluation at \(V\) of the {\v C}ech-nerve coface maps, assembled from the same push--pull
  restriction comparisons; under the degreewise identification each coface of \(D^\bullet\)
  becomes the corresponding face restriction, so the alternating sums agree. Reusing the
  alternating-sum bookkeeping of Lemma~\ref{mk-lem:section_cech_objd_apply} --- rather than
  rebuilding it --- matches these differentials.

  \emph{Augmentation differential.} The lowest differential of \(D^\bullet\) is the evaluation
  at \(V\) of the {\v C}ech augmentation \(\varepsilon : \mathcal{F} \to \mathcal{C}^0\)
  (Definition~\ref{mk-def:cech_augmentation}); read through the degree-\(0\) instance of
  Lemma~\ref{mk-lem:pushPull_eval_prod_iso} it is the restriction product map
  \(t \mapsto (t|_{U_i \cap V})_i\) by
  Lemma~\ref{mk-lem:mappedSectionCechAugV_eq}. This supplies the compatibility square
  \(h_{\mathrm{compat}}\) for Lemma~\ref{mk-lem:augmentCochainIso}. The source augmentation is a
  complex differential by Lemma~\ref{mk-lem:cech_augmentation_comp_d}, while the concrete
  restriction product is one by Lemma~\ref{mk-lem:sectionCechAugV_comp_d}. Hence the degreewise maps
  assemble to an
  isomorphism of cochain complexes \(D^\bullet \cong D'^\bullet_{\mathrm{aug}}\).
\end{proof}

\begin{lemma}

[The section {\v C}ech complex inside a cover member is contractible]
  \label{mk-lem:cechSection_contractible}
  \lean{AlgebraicGeometry.cechSection_contractible,
    AlgebraicGeometry.cechSectionAugComplex, AlgebraicGeometry.cechSectionCoeff,
    AlgebraicGeometry.cechSectionCoeff_transport, AlgebraicGeometry.cechSectionCoface,
    AlgebraicGeometry.cechSectionD_coord, AlgebraicGeometry.cechSectionDepDiff_zero,
    AlgebraicGeometry.cechSectionHomotopyComp, AlgebraicGeometry.cechSectionHomotopyComp_coord,
    AlgebraicGeometry.cechSectionHomotopyZero, AlgebraicGeometry.cechSectionPrepend,
    AlgebraicGeometry.cechSection_comm_one, AlgebraicGeometry.cechSection_comm_succ,
    AlgebraicGeometry.cechSection_comm_zero, AlgebraicGeometry.cechSection_hsh,
    AlgebraicGeometry.cechSection_hu, AlgebraicGeometry.cechSection_succ_step}
  \dcref{custom}

  \uses{mk-lem:cechSection_complex_iso, mk-lem:cech_acyclic_affine, mk-lem:cech_engine_complex,
    mk-def:cech_augmentation, mk-lem:cech_augmentation_comp_d}
  \textit{Source: Stacks Project, Cohomology of Schemes,
  \texttt{lemma-cech-cohomology-quasi-coherent-trivial} (proof, the projection homotopy).}
  Suppose \(V \le \operatorname{coverOpen}\mathcal{U}\,i_{\mathrm{fix}}\) for some fixed index
  \(i_{\mathrm{fix}}\), i.e.\ \(V\) lies inside one cover member. Then the \emph{augmented}
  concrete section {\v C}ech complex \(D'^\bullet_{\mathrm{aug}}\) of
  Lemma~\ref{mk-lem:cechSection_complex_iso} --- the section {\v C}ech complex over the family
  \(U'_j = \operatorname{coverInterOpen}\mathcal{U}\,j \cap V\), augmented at
  \(\Gamma(V, \mathcal{F})\) along the restriction product \(\varepsilon\) --- admits a
  contracting homotopy \(\mathrm{id}_{D'^\bullet_{\mathrm{aug}}} \simeq 0\).
\end{lemma}
\begin{proof}
  \uses{mk-lem:cech_acyclic_affine, mk-lem:cech_engine_complex, mk-def:cech_augmentation,
    mk-lem:cech_augmentation_comp_d}
  Because \(V \le U_{i_{\mathrm{fix}}}\), the restricted family \(U'_j = U_j \cap V\) has a
  \emph{maximum}: \(U'_{i_{\mathrm{fix}}} = U_{i_{\mathrm{fix}}} \cap V = V\), and every
  \(U'_j \le V\). Consequently, for each multi-index \(\sigma\),
  \(U'_{i_{\mathrm{fix}}\,\sigma} = U_{i_{\mathrm{fix}}} \cap U_\sigma \cap V = U_\sigma \cap
  V = U'_\sigma\): prepending the index \(i_{\mathrm{fix}}\) does not change the open, so the
  prepend map is the \emph{identity} on each coefficient \(\Gamma(U'_\sigma, \mathcal{F})\)
  (the membership criterion \(\operatorname{le\_coverInterOpen\_iff}\) of
  Lemma~\ref{mk-lem:cech_engine_complex}). The contracting homotopy is the prepend-\(i_{\mathrm{fix}}\)
  map, \(h(s)_\sigma = s_{i_{\mathrm{fix}}\,\sigma}\), exactly the projection homotopy of the
  Stacks proof; it splits into two parts.

  \emph{Positive {\v C}ech degrees \((p \ge 1\)).} On the section {\v C}ech part the
  dependent-coefficient combinatorial {\v C}ech engine --- \(\operatorname{depDiff}\),
  \(\operatorname{depHomotopy}\), \(\operatorname{depHomotopy\_spec}\),
  \(\operatorname{depDiff\_exact}\) of Lemma~\ref{mk-lem:cech_acyclic_affine} --- supplies the
  prepend-\(i_{\mathrm{fix}}\) homotopy \(h\) with
  \(d \circ h + h \circ d = \mathrm{id}\), since its unit, shift, and coface-commutation
  compatibilities all hold with the prepend map acting as the identity on each coefficient.

  \emph{Augmentation node (the piece the engine does not supply).} The engine governs only the
  {\v C}ech degrees; the contracting-homotopy identity at the augmentation node must be checked
  directly. The relevant homotopy component runs from the lowest {\v C}ech term
  \(\prod_i \Gamma(U_i \cap V, \mathcal{F})\) down to the augmentation term
  \(\Gamma(V, \mathcal{F})\), and is the \(i_{\mathrm{fix}}\)-th coordinate projection
  \(\pi_{i_{\mathrm{fix}}} : (s_j)_j \mapsto s_{i_{\mathrm{fix}}}\) --- well defined because
  \(U'_{i_{\mathrm{fix}}} = V\), so \(s_{i_{\mathrm{fix}}} \in \Gamma(U_{i_{\mathrm{fix}}} \cap V,
  \mathcal{F}) = \Gamma(V, \mathcal{F})\). We verify
  \(\varepsilon \circ \pi_{i_{\mathrm{fix}}} + (\text{degree-}0\ \text{engine term}) = \mathrm{id}\)
  on \(s \in \prod_i \Gamma(U_i \cap V, \mathcal{F})\). The augmentation differential
  \(\varepsilon\) is the restriction product, so
  \(\bigl(\varepsilon(\pi_{i_{\mathrm{fix}}} s)\bigr)_j = s_{i_{\mathrm{fix}}}|_{U_j \cap V}\); the
  degree-\(0\) contribution of the {\v C}ech homotopy supplies the complementary term
  \(s_j - s_{i_{\mathrm{fix}}}|_{U_j \cap V}\) coming from the prepend-\(i_{\mathrm{fix}}\) shift on
  the {\v C}ech part. Their sum is \(s_j\), i.e.\ the identity. (This is the bottom node of the
  Stacks projection homotopy: \(h\) on \(\prod_{i_0} M_{f_{i_0}} \to M\) is projection onto the
  factor \(M_{f_{i_{\mathrm{fix}}}} = M\).) Together the two parts give a homotopy
  \(\mathrm{id}_{D'^\bullet_{\mathrm{aug}}} \simeq 0\) of cochain complexes.
\end{proof}

\begin{lemma}[Local vanishing of the evaluated augmented {\v C}ech homology]\label{mk-lem:cechSection_isZero_homology}
  \lean{AlgebraicGeometry.cechSection_isZero_homology}
  \leanok
  \dcref{custom}

  \uses{mk-lem:cechSection_complex_iso, mk-lem:cechSection_contractible,
    mk-lem:isZero_homology_of_homotopy_id_zero}
  Let \(V \le \operatorname{coverOpen}\mathcal{U}\,i\) for some index \(i\). Then for every
  degree \(p\) the homology \(H^p(D^\bullet)\) of the augmented {\v C}ech complex evaluated at
  \(V\) (the complex \(D^\bullet\) of Lemma~\ref{mk-lem:cechSection_complex_iso}) is a zero
  object; equivalently \(D^\bullet\) carries a contracting homotopy \(\mathrm{id}_{D^\bullet}
  \simeq 0\). This is the residual obligation of Step~3 of
  Lemma~\ref{mk-lem:cech_augmented_resolution}.
\end{lemma}
\begin{proof}
  \leanok
  \uses{mk-lem:cechSection_complex_iso, mk-lem:cechSection_contractible,
    mk-lem:isZero_homology_of_homotopy_id_zero}
    By Lemma~\ref{mk-lem:cechSection_complex_iso} the evaluated complex \(D^\bullet\) is
  isomorphic to the augmented concrete section {\v C}ech complex \(D'^\bullet_{\mathrm{aug}}\).
  The hypothesis \(V \le \operatorname{coverOpen}\mathcal{U}\,i\) places \(V\) inside one cover
  member, so Lemma~\ref{mk-lem:cechSection_contractible} gives a contracting homotopy
  \(\mathrm{id}_{D'^\bullet_{\mathrm{aug}}} \simeq 0\); transporting it across the complex
  isomorphism yields \(\mathrm{id}_{D^\bullet} \simeq 0\).
  Lemma~\ref{mk-lem:isZero_homology_of_homotopy_id_zero} then makes
  \(H^p(D^\bullet)\) a zero object in every degree.
\end{proof}

\begin{lemma}[Pushforward sections are sections over the preimage]
  \label{mk-lem:pushforward_obj_obj_mathlib}
  \lean{AlgebraicGeometry.Scheme.Modules.pushforward_obj_obj}
  \mathlibok
  \dcref{custom}

  For a morphism \(f : X \to Y\), an \(\mathcal{O}_X\)-module \(\mathcal{M}\), and an open
  \(U \subseteq Y\), the sections of the pushforward \(f_*\mathcal{M}\) over \(U\) coincide
  definitionally with the sections of \(\mathcal{M}\) over the preimage:
  \(\Gamma(U, f_*\mathcal{M}) = \Gamma(f^{-1}(U), \mathcal{M})\).
\end{lemma}

\begin{lemma}[Pullback along an open immersion is the restriction functor]
  \label{mk-lem:restrictFunctorIsoPullback_mathlib}
  \lean{AlgebraicGeometry.Scheme.Modules.restrictFunctorIsoPullback}
  \mathlibok
  \dcref{custom}

  For an open immersion \(f : X \to Y\), the restriction functor
  \(\operatorname{restrictFunctor} f : Y.\mathrm{Modules} \to X.\mathrm{Modules}\) is
  naturally isomorphic to the module pullback functor \(f^* = \operatorname{pullback} f\). In
  particular \(f^*\mathcal{M} \cong \mathcal{M}.\operatorname{restrict} f\) for every
  \(\mathcal{O}_Y\)-module \(\mathcal{M}\).
\end{lemma}

\begin{lemma}[Evaluation of a sheaf of modules preserves products]
  \label{mk-lem:evaluation_preserves_products_mathlib}
  \lean{SheafOfModules.evaluationPreservesLimitsOfShape}
  \mathlibok
  \dcref{custom}

  For a sheaf of rings on a site and any object \(U\), the evaluation (sections-at-\(U\))
  functor \(\operatorname{evaluation}\,U : \operatorname{SheafOfModules} \to \mathrm{Ab}\)
  preserves limits of every shape; in particular it preserves products, so \(\Gamma(U,
  \prod_\sigma N_\sigma) \cong \prod_\sigma \Gamma(U, N_\sigma)\).
\end{lemma}

\begin{lemma}[Sections over a binary coproduct scheme are a product]
  \label{mk-lem:coprodPresheafObjIso_mathlib}
  \lean{AlgebraicGeometry.Scheme.coprodPresheafObjIso}
  \mathlibok
  \dcref{custom}

  For schemes \(X, Y\) and an open \(U\) of the coproduct \(X \amalg Y\), the structure-sheaf
  sections over \(U\) decompose as the product of the sections over the two preimages:
  \((X \amalg Y).\mathcal{O}(U) \cong X.\mathcal{O}(\operatorname{inl}^{-1} U) \times
  Y.\mathcal{O}(\operatorname{inr}^{-1} U)\). This is the binary disjoint-union sheaf
  decomposition that Lemma~\ref{mk-lem:pushPull_sigma_iso} ports to the indexed,
  \(\operatorname{SheafOfModules}\) setting.
\end{lemma}

\begin{lemma}[Sections of a sheaf over a disjoint union are a product]
  \label{mk-lem:isProductOfDisjoint_mathlib}
  \lean{TopCat.Sheaf.isProductOfDisjoint}
  \mathlibok
  \dcref{custom}

  For a sheaf \(F\) (valued in a category with products) on a topological space and two opens
  \(U, V\) with \(U \sqcap V = \bot\) (disjoint), the sections over the union are the product
  of the sections: \(F(U \sqcup V) \cong F(U) \times F(V)\). Iterated over a finite family,
  this is the disjoint-union decomposition underlying Lemma~\ref{mk-lem:pushPull_sigma_iso}.
\end{lemma}

\begin{lemma}[The category of schemes is finitary extensive]
  \label{mk-lem:finitaryExtensive_scheme_mathlib}
  \lean{AlgebraicGeometry.instFinitaryExtensiveScheme}
  \mathlibok
  \dcref{custom}

  The category \(\operatorname{Scheme}\) is finitary extensive: finite coproducts are disjoint
  and universal (stable under pullback). In particular it is finitary pre-extensive and has
  finite coproducts and pullbacks, so the abstract distributivity of
  Lemma~\ref{mk-lem:isIso_sigmaDesc_map_mathlib} applies at \(\mathcal{C} = \operatorname{Scheme}\).
\end{lemma}

\begin{lemma}[Coproducts distribute over the binary fibre product]
  \label{mk-lem:isIso_sigmaDesc_map_mathlib}
  \lean{CategoryTheory.FinitaryPreExtensive.isIso_sigmaDesc_map}
  \mathlibok
  \dcref{custom}

  In a finitary pre-extensive category with pullbacks, for finite index types \(\iota, \iota'\),
  a base \(S\), and families \((X_i \to S)_i\), \((Y_j \to S)_j\), the canonical map
  \(\coprod_{(i,j)} (X_i \times_S Y_j) \to (\coprod_i X_i) \times_S (\coprod_j Y_j)\),
  \(\operatorname{Sigma.desc}\) of the pullback maps, is an isomorphism. Taking one family to be a
  singleton gives the one-sided distributivity
  \((\coprod_i X_i) \times_S Y \cong \coprod_i (X_i \times_S Y)\) used in the induction of
  Lemma~\ref{mk-lem:coproduct_distrib_fibrePower}.
\end{lemma}

\begin{lemma}[The coproduct-of-pullback-fst comparison is an isomorphism]
  \label{mk-lem:isIso_sigmaDesc_fst_mathlib}
  \lean{CategoryTheory.FinitaryPreExtensive.isIso_sigmaDesc_fst}
  \mathlibok
  \dcref{custom}

  In a finitary pre-extensive category, for a finite index type \(\alpha\) in \(\operatorname{Type}
  0\), a base \(S\), a family \((X_a \to S)_{a \in \alpha}\) and a single \(Y \to S\), the comparison
  map \(\operatorname{Sigma.desc}\) of the pullback first-projections
  \(\coprod_a (X_a \times_S Y) \to (\coprod_a X_a) \times_S Y\) is an isomorphism. This is the
  primitive on which the binary distributivity Lemma~\ref{mk-lem:isIso_sigmaDesc_map_mathlib} and the
  one-sided form Lemma~\ref{mk-lem:prod_coproduct_distrib} rest. \emph{It is stated in Mathlib only for}
  \(\alpha : \operatorname{Type} 0\), which is why the whole distributivity chain
  (Lemma~\ref{mk-lem:coproduct_distrib_fibrePower}) is confined to \(\operatorname{Type} 0\) and the
  cover index is reduced to \(\mathrm{Fin}(n)\) at Lemma~\ref{mk-lem:cech_backbone_left_sigma}.
\end{lemma}

\begin{lemma}[Pullback of two open immersions is the intersection open]
  \label{mk-lem:isPullback_opens_inf_mathlib}
  \lean{AlgebraicGeometry.isPullback_opens_inf}
  \mathlibok
  \dcref{custom}

  For two opens \(U, V\) of a scheme \(X\), the square with vertex \((U \sqcap V).\mathrm{toScheme}\)
  and the two \(\operatorname{homOfLE}\) inclusions into \(U\) and \(V\) is a pullback of the open
  immersions \(U.\iota, V.\iota\). Equivalently \(\operatorname{pullback} U.\iota\,V.\iota \cong
  (U \sqcap V).\mathrm{toScheme}\): the fibre product of open immersions over \(X\) is the
  intersection of their images.
\end{lemma}

\begin{lemma}[A wide pullback over a terminal base is the product of its legs]
  \label{mk-lem:widePullbackCone_isLimitOfFan_mathlib}
  \lean{CategoryTheory.Limits.WidePullbackCone.isLimitOfFan}
  \mathlibok
  \dcref{custom}

  When the base of a wide pullback is a terminal object, a limiting fan of the legs is a limit of
  the wide-pullback diagram. Hence in \(\operatorname{Over} S\) --- where \(\operatorname{Over.mk}
  (\mathrm{id}_S)\) is terminal --- the wide pullback over \(S\) of a family is the product of the
  legs in the slice, which is what lets the fibre power over \(S\) be treated as an iterated
  product in Lemma~\ref{mk-lem:coproduct_distrib_fibrePower}.
\end{lemma}

\begin{lemma}[The identity arrow is terminal in the slice category]
  \label{mk-lem:over_mkIdTerminal_mathlib}
  \lean{CategoryTheory.Over.mkIdTerminal}
  \mathlibok
  \dcref{custom}

  For an object \(S\) of a category, \(\operatorname{Over.mk}(\mathrm{id}_S)\) is a terminal object
  of the slice category \(\operatorname{Over} S\). This is what makes the base of a wide pullback
  over \(S\) terminal in the slice, so that
  Lemma~\ref{mk-lem:widePullbackCone_isLimitOfFan_mathlib} exhibits the fibre power over \(S\) as the
  product in \(\operatorname{Over} S\) in Lemma~\ref{mk-lem:widePullback_overX_eq_prod}.
\end{lemma}

\begin{lemma}[Nested coproducts flatten to a coproduct over pairs]
  \label{mk-lem:sigmaSigmaIso_mathlib}
  \lean{CategoryTheory.Limits.sigmaSigmaIso}
  \mathlibok
  \dcref{custom}

  For a doubly-indexed family \((Z_{i,j})\), the nested coproduct is the coproduct over pairs:
  \(\coprod_i \coprod_j Z_{i,j} \cong \coprod_{(i,j)} Z_{i,j}\). This collapses the nested
  coproduct produced by one inductive step of the binary distributivity into a single coproduct
  over the multi-indices \(\sigma\) in Lemma~\ref{mk-lem:coproduct_distrib_fibrePower}.
\end{lemma}

\begin{lemma}
[Cohomology sheaf is the sheafification of the objectwise homology]
  \label{mk-lem:cohomology_sheaf_is_sheafify_homology}
  \lean{PresheafOfModules.homologyIsoSheafify,
    PresheafOfModules.counitComplexIso,
    PresheafOfModules.sheafificationAdditive,
    CategoryTheory.Functor.mapHomologyIso'}
  \dcref{01XJ,custom}

  Let \(K\) be a cochain complex of presheaves of modules over a site admitting
  sheafification. For every index \(i\) there is a natural
  isomorphism between the \(i\)-th homology of the sheafified complex and the
  sheafification of the \(i\)-th objectwise (presheaf-level) homology of \(K\),
  \[
    H^i\bigl(\operatorname{sheafify}(K)\bigr)
      \;\cong\;
    \operatorname{sheafify}\bigl(H^i(K)\bigr).
  \]
  Equivalently, sheafification commutes with the formation of homology. This is the
  reusable engine underlying the presheaf description of the higher direct images
  (Stacks Tag~01XJ): it is what turns the resolution-internal cohomology sheaf into a
  sheafification of an explicit cohomology presheaf.
\end{lemma}

\begin{proof}
  Sheafification is exact: as the left adjoint of the inclusion of sheaves into
  presheaves it preserves all colimits, and on a site admitting sheafification it
  additionally preserves finite limits, so it preserves both the kernels and the
  cokernels through which homology is computed. An exact additive functor commutes with
  the homology functor of a cochain complex; chasing this through the counit complex
  isomorphism \(\operatorname{sheafify}(\iota K) \cong K\) (the inclusion–sheafification
  adjunction counit, an isomorphism on sheaves) gives
  \(H^i(K) \cong H^i(\operatorname{sheafify} K) \cong
  \operatorname{sheafify}(H^i(K))\). The supporting bricks are the
  additivity of sheafification, the counit complex isomorphism, and the
  complex-level transport of the short-complex map–homology isomorphism along an
  additive homology-preserving functor.
\end{proof}

\begin{lemma}
[Presheaf description of the higher direct images]
  \label{mk-lem:higher_direct_image_presheaf}
  \lean{AlgebraicGeometry.higherDirectImage_iso_sheafify_presheafHomology,
    AlgebraicGeometry.pushforwardResolutionPresheafComplex}
  \dcref{01XJ,custom}

  \uses{mk-def:higher_direct_image, mk-lem:cohomology_sheaf_is_sheafify_homology}
  \textit{Source: Stacks Project, Cohomology, Tag 01XJ
  (\texttt{lemma-describe-higher-direct-images}).}
  Let \(f : X \to S\) be a morphism of schemes and \(\mathcal{G}\) an
  \(\mathcal{O}_X\)-module, and assume the category of \(\mathcal{O}_X\)-modules has
  enough injective objects, so that the derived pushforward \(R^k f_*\) is defined.
  Choose an injective resolution \(\mathcal{G} \to \mathcal{I}^\bullet\) in
  \(\mathcal{O}_X\)-modules and push it forward degreewise to obtain the cochain complex
  of presheaves of modules \(f_*\mathcal{I}^\bullet\) on \(S\). For every \(k \geq 0\)
  the higher direct image \(R^k f_*\mathcal{G}\) is isomorphic to the
  \emph{sheafification of the presheaf of objectwise homology} of this pushed complex,
  \[
    R^k f_*\mathcal{G}
      \;\cong\;
    \operatorname{sheafify}\Bigl(\,V \longmapsto
      H^k\bigl((f_*\mathcal{I}^\bullet)(V)\bigr)\,\Bigr),
    \qquad V \subseteq S \text{ open},
  \]
  the sheafification of the \(k\)-th presheaf-level homology of
  \(f_*\mathcal{I}^\bullet\). As a closing remark, this gives the affine-local
  vanishing criterion consumed downstream
  (Lemmas~\ref{mk-lem:open_immersion_pushforward_comp} and
  \ref{mk-lem:cech_term_pushforward_acyclic}): since the sheafification of a presheaf
  vanishes exactly when that presheaf vanishes on a basis of opens,
  \(R^k f_*\mathcal{G} = 0\) iff \(H^k\bigl((f_*\mathcal{I}^\bullet)(V)\bigr) = 0\)
  for every \(V\) in a basis of the topology of \(S\); the affine opens form such a
  basis, so it suffices to check the vanishing on affine \(V\).

  The identification of the displayed presheaf homology with the absolute cohomology
  of the preimage,
  \[
    H^k\bigl((f_*\mathcal{I}^\bullet)(V)\bigr)
      \;=\;
    H^k\bigl(f^{-1}(V), \mathcal{G}|_{f^{-1}(V)}\bigr),
  \]
  which uses that \(\mathcal{I}^\bullet|_{f^{-1}(V)}\) is again injective over
  \(f^{-1}(V)\), is supplied at point of use; this lemma records the
  resolution-internal sheafify-of-objectwise-homology form.
\end{lemma}

\begin{proof}
  \uses{mk-def:higher_direct_image, mk-lem:cohomology_sheaf_is_sheafify_homology}
  Choose an injective resolution \(\mathcal{G} \to \mathcal{I}^\bullet\) in
  \(\mathcal{O}_X\)-modules. By definition \(R^k f_*\mathcal{G}\) is the \(k\)-th
  cohomology sheaf of the complex \(f_*\mathcal{I}^\bullet\) of presheaves of modules
  over \(S\). The comparison engine of
  Lemma~\ref{mk-lem:cohomology_sheaf_is_sheafify_homology} now applies: taking homology in
  \(\mathcal{O}_S\)-modules is the sheafification of the objectwise (presheaf-level)
  homology, because the sheafification/inclusion adjunction is exact (it preserves the
  finite limits and colimits computing kernels and cokernels). Hence
  \[
    R^k f_*\mathcal{G}
      \;\cong\;
    \operatorname{sheafify}\Bigl(\,V \longmapsto
      H^k\bigl((f_*\mathcal{I}^\bullet)(V)\bigr)\,\Bigr),
  \]
  which is the stated isomorphism. Since the sheafification of a presheaf vanishes
  exactly when the underlying presheaf vanishes on a basis of opens, this yields the
  affine-local vanishing criterion.
\end{proof}

\begin{lemma}
[Ext as the cohomology of a Hom complex into an injective resolution]
  \label{mk-lem:ext_homcomplex_mathlib}
  \lean{CategoryTheory.InjectiveResolution.extAddEquivCohomologyClass,
    CochainComplex.HomComplex.homologyAddEquiv,
    CategoryTheory.InjectiveResolution.cochainComplexXIso}
  \mathlibok
  \dcref{custom}

  \uses{mk-lem:right_derived_injective_resolution}
  Let \(\mathcal{I}^\bullet\) be an injective resolution of \(Y\) in an abelian category with
  enough injectives. Then there is a natural identification, for every \(n\),
  \[
    \operatorname{Ext}^n(X, Y)
      \;\cong\;
    H^n\bigl(\operatorname{Hom}(X[0], \mathcal{I}^\bullet)\bigr),
  \]
  the composite of the Ext-to-cohomology-class equivalence
  \(\operatorname{extAddEquivCohomologyClass}\) with the cohomology-class-to-homology equivalence
  \(\operatorname{homologyAddEquiv}\) of the Hom complex; the degreewise comparison
  \(\operatorname{cochainComplexXIso}\) matches the \(\mathbb{N}\)- and \(\mathbb{Z}\)-indexed forms of
  the resolution. The companion right-derived-object comparison
  \((F.\mathrm{rightDerived}\,n)(X) \cong H^n(F \cdot \mathcal{I}^\bullet)\) for an additive functor
  \(F\) is recorded in Lemma~\ref{mk-lem:right_derived_injective_resolution}.
\end{lemma}

\begin{lemma}
[Ext-vanishing forces the coyoneda right-derived object to vanish]
  \label{mk-lem:isZero_coyoneda_rightDerived_of_forall_ext_eq_zero}
  \lean{AlgebraicGeometry.isZero_coyoneda_rightDerived_of_forall_ext_eq_zero,
    AlgebraicGeometry.preadditiveCoyoneda_mapHomologicalComplex_d_apply}
  \dcref{custom}

  \uses{mk-lem:ext_homcomplex_mathlib, mk-lem:right_derived_injective_resolution}
  Let \(\mathcal{C}\) be an abelian category with enough injectives, let \(P, H \in \mathcal{C}\), and
  let \(q \geq 1\). If every class \(e \in \operatorname{Ext}^q(P, H)\) is zero, then the right-derived
  object of the covariant hom-functor vanishes at \(H\):
  \[
    \bigl((\operatorname{Hom}(P, -)).\mathrm{rightDerived}\,q\bigr)(H) \;\text{ is a zero object.}
  \]
\end{lemma}
\begin{proof}
  \uses{mk-lem:ext_homcomplex_mathlib, mk-lem:right_derived_injective_resolution}
  Choose an injective resolution \(H \to \mathcal{I}^\bullet\). By the right-derived-object comparison
  (Lemma~\ref{mk-lem:right_derived_injective_resolution}) the value
  \(((\operatorname{Hom}(P,-)).\mathrm{rightDerived}\,q)(H)\) is identified with the degree-\(q\)
  homology of the \(\operatorname{Hom}\)-cochain complex \(n \mapsto (P \to \mathcal{I}^n)\), whose
  differential is post-composition with the differential of \(\mathcal{I}^\bullet\). A degree-\(q\)
  cocycle is a map \(f : P \to \mathcal{I}^q\) with \(f \circ d = 0\), and by the
  Ext--Hom-complex identification (Lemma~\ref{mk-lem:ext_homcomplex_mathlib}) such an \(f\) is a
  coboundary exactly when its class \(\operatorname{extMk} f \in \operatorname{Ext}^q(P, H)\) vanishes.
  The hypothesis makes every such class zero, so every \(q\)-cocycle is a coboundary; the homology at
  \(q\) is therefore zero and the right-derived object is a zero object.
\end{proof}

\begin{lemma}[An isomorphism in the first Ext-argument transfers subsingletons]\label{mk-lem:subsingleton_ext_of_iso_fst}
  \lean{AlgebraicGeometry.subsingleton_ext_of_iso_fst}
  \leanok
  \dcref{custom}

  \uses{mk-lem:ext_homcomplex_mathlib}
  Let \(\mathcal{C}\) be an abelian category with an Ext-theory, \(\varphi : A \cong B\) an
  isomorphism, \(Y \in \mathcal{C}\) an object and \(q\) a degree. If \(\operatorname{Ext}^q(B, Y)\)
  is a subsingleton (has at most one element), then so is \(\operatorname{Ext}^q(A, Y)\).
\end{lemma}
\begin{proof}
  \leanok
  \uses{mk-lem:ext_homcomplex_mathlib}
  Ext is contravariant in its first argument. For any \(z \in \operatorname{Ext}^q(A, Y)\),
  associativity of Yoneda composition with the degree-\(0\) classes of \(\varphi\) gives
  \[
    z = \operatorname{mk}_0(\varphi.\mathrm{hom}) \circ \bigl(\operatorname{mk}_0(\varphi.\mathrm{inv})
      \circ z\bigr),
  \]
  using \(\varphi.\mathrm{inv} \circ \varphi.\mathrm{hom} = \mathrm{id}\) and the identity law
  \(\operatorname{mk}_0(\mathrm{id}) \circ z = z\). The inner factor
  \(\operatorname{mk}_0(\varphi.\mathrm{inv}) \circ z\) lives in the subsingleton
  \(\operatorname{Ext}^q(B, Y)\), hence equals \(0\); composing with \(0\) gives \(z = 0\). Since every
  element is \(0\), \(\operatorname{Ext}^q(A, Y)\) is a subsingleton.
\end{proof}

\begin{lemma}[Injective resolutions provide enough injectives]\label{mk-lem:enoughInjectives_of_hasInjectiveResolutions}
  \lean{AlgebraicGeometry.enoughInjectives_of_hasInjectiveResolutions}
  \leanok
  \dcref{custom}

  \uses{mk-lem:right_derived_injective_resolution}
  An abelian category equipped with chosen injective resolutions has enough injectives: every object
  admits a monomorphism into an injective object.
\end{lemma}
\begin{proof}
  \leanok
  \uses{mk-lem:right_derived_injective_resolution}
  Given an object \(X\), take its chosen injective resolution \(X \to \mathcal{I}^\bullet\). The
  degree-\(0\) term \(\mathcal{I}^0\) is injective and the resolution's unit \(X \to \mathcal{I}^0\) is
  a monomorphism, so it is an injective presentation of \(X\). As every object has one, the category
  has enough injectives.
\end{proof}

\begin{lemma}
[Affine open immersions have vanishing higher direct images]
  \label{mk-lem:open_immersion_pushforward_acyclic}
  \lean{AlgebraicGeometry.higherDirectImage_openImmersion_acyclic,
    AlgebraicGeometry.isAffineHom_of_affine_separated}
  \dcref{custom}

  \uses{mk-lem:affine_serre_vanishing, mk-lem:higher_direct_image_presheaf,
    mk-lem:sheafify_kills_locally_zero,
    mk-lem:isZero_presheafToSheaf_of_locally_isZero, mk-lem:ext_homcomplex_mathlib,
    mk-lem:isZero_presheafToSheaf_of_sections_locally_zero, mk-lem:pushforward_sections_functor,
    mk-lem:rightDerivedNatIso, mk-lem:sectionsFunctorCorepIso,
    mk-lem:affine_serre_vanishing_general_open, mk-lem:modules_isoSpec_ext_transport,
    mk-lem:isoSpec_scheme_mathlib,
    mk-lem:isZero_coyoneda_rightDerived_of_forall_ext_eq_zero,
    mk-lem:ext_jShriekOU_eq_zero_of_specIso}
  \textit{Source: Stacks Project, Cohomology of Schemes,
  \texttt{lemma-relative-affine-vanishing}.}
  Let \(X\) be separated and \(j : U \hookrightarrow X\) the open immersion of an affine
  open \(U\). Then \(j\) is an affine morphism (an open immersion of an affine open into a
  separated scheme), and its higher direct images vanish:
  \[
    R^q j_*\mathcal{H} = 0 \qquad \text{for all } q \geq 1
  \]
  and every quasi-coherent \(\mathcal{O}_U\)-module \(\mathcal{H}\).
\end{lemma}
\begin{proof}
  \uses{mk-lem:affine_serre_vanishing, mk-lem:higher_direct_image_presheaf, mk-lem:sheafify_kills_locally_zero, mk-lem:isZero_presheafToSheaf_of_locally_isZero, mk-lem:ext_homcomplex_mathlib, mk-lem:isZero_presheafToSheaf_of_sections_locally_zero, mk-lem:pushforward_sections_functor, mk-lem:rightDerivedNatIso, mk-lem:sectionsFunctorCorepIso, mk-lem:affine_serre_vanishing_general_open, mk-lem:modules_isoSpec_ext_transport, mk-lem:isoSpec_scheme_mathlib, mk-lem:isZero_coyoneda_rightDerived_of_forall_ext_eq_zero, mk-lem:ext_jShriekOU_eq_zero_of_specIso}
  An open immersion of an affine open \(U\) into a separated scheme \(X\) is an
  affine morphism: separatedness makes the immersion \(j\) quasi-separated and the
  preimage of any affine open of \(X\) meets \(U\) in an affine, so \(j\) is affine.
  The relative affine vanishing for affine morphisms then gives \(R^q j_*\mathcal{H} =
  0\) for \(q \geq 1\); concretely, by the presheaf description
  (Lemma~\ref{mk-lem:higher_direct_image_presheaf}) \(R^q j_*\mathcal{H}\) is the
  sheafification of \(V \mapsto H^q(j^{-1}(V), \mathcal{H})\), and for affine \(V\) the
  preimage \(j^{-1}(V)\) is affine, on which the Serre vanishing of
  Lemma~\ref{mk-lem:affine_serre_vanishing} kills \(H^q\) for \(q \geq 1\). Since the affine
  opens form a basis of \(X\), the presheaf \(V \mapsto H^q(j^{-1}(V), \mathcal{H})\) is
  locally zero, so its sheafification \(R^q j_*\mathcal{H}\) vanishes for \(q \geq 1\).

  \emph{Proof details.} The argument rests on three identifications.

  \emph{(1) The cohomology-presheaf identification.} The objectwise homology in degree
  \(n\) of the pushed-forward resolution complex \(j_*\mathcal{I}^\bullet\) (for a chosen
  injective resolution \(\mathcal{H} \to \mathcal{I}^\bullet\) of the coefficient), evaluated at an open
  \(V\), is identified with the absolute cohomology presheaf
  \(V \mapsto H^n(j^{-1}(V), \mathcal{H}) = \operatorname{Ext}^n\bigl(j_!\mathcal{O}_{j^{-1}(V)},
  \mathcal{H}\bigr)\). Concretely: \(\operatorname{Ext}^n(X, Y)\) over an injective resolution of
  the second argument \(Y\) is identified with the \(n\)-th cohomology of the \(\operatorname{Hom}\)
  cochain complex \(H^n\bigl(\operatorname{Hom}(X[0], \mathcal{I}^\bullet)\bigr)\), via the composite
  of the Ext-to-cohomology-class equivalence and the cohomology-class-to-homology equivalence of the
  Hom complex (the equivalences of Lemma~\ref{mk-lem:ext_homcomplex_mathlib}). Taking \(X = j_!\mathcal{O}_{j^{-1}(V)}\) and \(Y =
  \mathcal{H}\), the degree-\(k\) Hom group \(\operatorname{Hom}(j_!\mathcal{O}_{j^{-1}(V)},
  \mathcal{I}^k)\) is identified with \(\Gamma(j^{-1}(V), \mathcal{I}^k) = (j_*\mathcal{I}^k)(V)\) by
  the corepresentability \(j_!\mathcal{O}_{j^{-1}(V)} \dashv \Gamma(j^{-1}(V), -)\) of the absolute-cohomology setup,
  degreewise and naturally; the \(\mathbb{N}\)/\(\mathbb{Z}\) indexing of the two cochain complexes is
  matched by the resolution's degreewise comparison isomorphism. The right-derived side is settled by
  the already-used identification of \(R^n j_*\) with the homology of \(j_*\mathcal{I}^\bullet\)
  (Lemma~\ref{mk-lem:higher_direct_image_presheaf}); no separate ``right-derived equals Ext'' comparison
  lemma is needed, since both compute the same homology over the same injective resolution and meet at
  corepresentability.

  \emph{(2) Serre vanishing for a general affine open.} The Serre vanishing of
  Lemma~\ref{mk-lem:affine_serre_vanishing} is pinned to the \emph{top} open \(\top\) of a spectrum
  \(\operatorname{Spec} R\): it kills \(\operatorname{Ext}^q(j_!\mathcal{O}_\top, -)\) for \(q \geq 1\).
  The cohomology that actually occurs here is over a \emph{general} affine open --- the preimage
  \(j^{-1}(V)\) of an affine open \(V\) of \(X\) that may straddle
  the boundary of \(j(U)\), so \(j^{-1}(V)\) is a general, possibly non-distinguished, affine open of
  \(U\). The residual is therefore \(\operatorname{Ext}^q(j_!\mathcal{O}_{j^{-1}(V)}, \mathcal{H}) = 0\)
  over the abstract affine scheme \(U\), for \(q \geq 1\). We discharge it by splitting it into two
  independent transports.

  \emph{(2a) The residual lives over an affine scheme.} The immersion \(j\) is an affine morphism:
  for \(U\) affine and \(X\) separated,
  \(j : U \hookrightarrow X\) is affine, so the preimage of every affine open of \(X\) is an affine
  open of \(U\). In particular \(j^{-1}(V) = U \cap V\) is an
  affine open of the affine scheme \(U\). After the reduction to the right-derived sections functor
  --- by the presheaf description (Lemma~\ref{mk-lem:higher_direct_image_presheaf}) and the
  corepresentability isomorphism of Lemma~\ref{mk-lem:sectionsFunctorCorepIso}, upgraded to right derived
  functors by Lemma~\ref{mk-lem:rightDerivedNatIso} --- the residual reads
  \(\operatorname{Ext}^q_{U}(j_!\mathcal{O}_{j^{-1}(V)}, \mathcal{H}) = 0\), an \(\operatorname{Ext}\)
  vanishing in the module category of the abstract affine scheme \(U\), for \(q \geq 1\).

  \emph{(2b) Transport along the whole-scheme spectrum isomorphism (Need~\#1).} The abstract affine
  scheme \(U\) carries the canonical \emph{whole-scheme} isomorphism
  \(U \cong \operatorname{Spec}\Gamma(U, \mathcal{O}_U)\) (\(\operatorname{Scheme.isoSpec}\),
  Lemma~\ref{mk-lem:isoSpec_scheme_mathlib}), a genuine isomorphism of schemes because \(U\) is affine.
  By Lemma~\ref{mk-lem:modules_isoSpec_ext_transport} it induces an equivalence of module categories
  \(U.\mathrm{Modules} \simeq (\operatorname{Spec}\Gamma U).\mathrm{Modules}\) and transports
  \(\operatorname{Ext}\) along it (via \(\operatorname{Ext.mapExactFunctor}\)), carrying the residual to
  \(\operatorname{Ext}^q_{\operatorname{Spec}\Gamma U}(j_!\mathcal{O}_{V'}, \mathcal{H}') = 0\), where
  \(V'\) is the image affine open of \(j^{-1}(V)\) and \(\mathcal{H}'\) the transported (still
  quasi-coherent) coefficient. Under the equivalence \(V'\) is in general a non-distinguished affine
  open of \(\operatorname{Spec}\Gamma U\).

  \emph{(2c) General-affine-open Serre vanishing (Need~\#2).} Over \(\operatorname{Spec}\Gamma U\) the
  open \(V'\) is a general affine open, and Lemma~\ref{mk-lem:affine_serre_vanishing_general_open} gives
  \(\operatorname{Ext}^q(j_!\mathcal{O}_{V'}, \mathcal{H}') = 0\) for \(q \geq 1\) and every affine
  open \(V'\) and quasi-coherent \(\mathcal{H}'\). This is read off from the basis-comparison criterion
  (Lemma~\ref{mk-lem:cech_to_cohomology_on_basis}) with the cover system's basis \(\mathcal{B}\) enlarged
  from the distinguished opens to \emph{all} affine opens; it requires no fresh cover system on the
  open \(V'\) and stays entirely inside \((\operatorname{Spec}\Gamma U).\mathrm{Modules}\). Combining
  (2b) and (2c) kills the residual.

  \textit{Remark (a blocked approach).} One might try to make \(j^{-1}(V)\) the \emph{whole} space by
  transporting along the spectrum isomorphism of the \emph{open subscheme}
  \(j^{-1}(V) \cong \operatorname{Spec}\Gamma(j^{-1}(V), \mathcal{O})\). This is a dead end: the
  inclusion \(j^{-1}(V) \hookrightarrow U\) is an open immersion, not an isomorphism, so there is no
  equivalence \(U.\mathrm{Modules} \simeq (j^{-1}(V)).\mathrm{Modules}\) to transport
  \(\operatorname{Ext}\) along; using the open-subscheme isomorphism forces a restriction functor
  \(\mathcal{H} \mapsto \mathcal{H}|_{j^{-1}(V)}\), whose derived comparison
  \(\operatorname{Ext}^q_{U}(j_!\mathcal{O}_{j^{-1}(V)}, \mathcal{H}) \cong H^q(j^{-1}(V),
  \mathcal{H}|_{j^{-1}(V)})\) is exactly ``restriction preserves injectives'' --- infrastructure not
  available in the present formalism. The sound route therefore keeps the spectrum isomorphism at the
  level of the \emph{whole} affine \(U\) as in~(2b) and handles the general affine open ambiently
  via~(2c).

  \emph{(3) Locally zero implies zero sheafification.} The
  passage from ``\(H^q(\cdots) = 0\) on a basis of affine opens'' to ``\(R^q j_* = 0\)'' is the
  module-sheafification analogue of Lemma~\ref{mk-lem:sheafify_kills_locally_zero} (and of its packaged
  form Lemma~\ref{mk-lem:isZero_presheafToSheaf_of_locally_isZero}). It is obtained by transporting the
  abelian-sheaf statement across the natural isomorphism
  \(\operatorname{toSheaf} \circ \operatorname{sheafify} \cong \operatorname{presheafToSheaf} \circ
  \operatorname{forget}\) between module-sheafification followed by the forgetful functor and abelian
  sheafification; the ``locally zero'' hypothesis it consumes is exactly the basis-of-affines vanishing
  supplied by identifications (1) and~(2) above.
\end{proof}

\begin{lemma}[Pushing forward an affine open immersion degenerates the composite direct image]\label{mk-lem:open_immersion_pushforward_comp}
  \lean{AlgebraicGeometry.higherDirectImage_openImmersion_comp}
  \leanok
  \dcref{custom}

  \uses{mk-lem:open_immersion_pushforward_acyclic, mk-lem:injective_of_adjoint,
    mk-def:right_acyclic, mk-lem:right_derived_vanishes_injective,
    mk-lem:acyclic_resolution_computes_derived, mk-lem:right_derived_injective_resolution,
    mk-lem:right_derived_zero_iso_self, mk-lem:scheme_pullbackPushforwardAdjunction_mathlib,
    mk-lem:restrictFunctorIsoPullback_mathlib, mk-lem:scheme_pushforwardComp_mathlib}
  Let \(X\) be separated, \(j : U \hookrightarrow X\) the open immersion of an affine
  open \(U\), and \(f : X \to S\) any morphism, with composite \(g := f \circ j : U \to S\).
  Then for every quasi-coherent \(\mathcal{O}_U\)-module \(\mathcal{H}\) and every \(k \geq 0\),
  \[
    R^k f_*\bigl(j_*\mathcal{H}\bigr) \;\cong\; R^k g_*\mathcal{H}.
  \]
\end{lemma}
\begin{proof}
  \leanok
  \uses{mk-lem:open_immersion_pushforward_acyclic, mk-lem:injective_of_adjoint,
    mk-def:right_acyclic, mk-lem:right_derived_vanishes_injective,
    mk-lem:acyclic_resolution_computes_derived, mk-lem:right_derived_injective_resolution,
    mk-lem:right_derived_zero_iso_self, mk-lem:scheme_pullbackPushforwardAdjunction_mathlib,
    mk-lem:restrictFunctorIsoPullback_mathlib, mk-lem:scheme_pushforwardComp_mathlib}
    Choose an injective resolution \(\mathcal{H} \to \mathcal{I}^\bullet\) in
  \(\mathcal{O}_U\)-modules and apply the pushforward \(j_*\) degreewise, forming the complex
  \(K := j_*\mathcal{I}^\bullet\) of \(\mathcal{O}_X\)-modules. We verify the hypotheses of the
  acyclic-resolution comparison (Lemma~\ref{mk-lem:acyclic_resolution_computes_derived}) for the
  functor \(G = f_*\) and the candidate \(f_*\)-acyclic resolution \(K\) of \(j_*\mathcal{H}\), then
  transport along the pushforward-composition isomorphism. The argument is the
  categorical route; it establishes the \(f_*\)-acyclicity of each
  \(j_*\mathcal{I}^n\) by an adjoint-preserves-injectives argument and is an equivalent, cleaner
  alternative to the presheaf-description argument of the Stacks Project (which would compute
  cohomology over the opens \(U \cap f^{-1}(V)\); these are general opens of the affine \(U\) and
  need not be affine, since \(f^{-1}(V)\) is open but not affine in \(X\), so that route re-enters
  the restriction-of-injectives obstruction we deliberately avoid).

  \emph{(b) Each \(j_*\mathcal{I}^n\) is \(f_*\)-acyclic.} The pushforward
  \(j_* = \operatorname{pushforward} j\) preserves injective objects. Indeed, it is the right
  adjoint in the adjunction \(\operatorname{pullback} j \dashv \operatorname{pushforward} j\)
  (Lemma~\ref{mk-lem:scheme_pullbackPushforwardAdjunction_mathlib}), and its left adjoint
  \(\operatorname{pullback} j\) preserves monomorphisms: for an open immersion \(j\) the pullback
  functor is isomorphic to the restriction functor,
  \(\operatorname{restrictFunctor} j \cong \operatorname{pullback} j\)
  (Lemma~\ref{mk-lem:restrictFunctorIsoPullback_mathlib}), and restriction of a sheaf of modules to an
  open subscheme is exact, hence preserves monos. By the criterion that the right adjoint of a
  mono-preserving left adjoint preserves injectives (Lemma~\ref{mk-lem:injective_of_adjoint}), each
  \(j_*\mathcal{I}^n\) is injective in \(X.\mathrm{Modules}\). An injective object is
  right-\(G\)-acyclic for \emph{any} additive functor \(G\)
  (Lemma~\ref{mk-lem:right_derived_vanishes_injective}, in the sense of
  Definition~\ref{mk-def:right_acyclic}), so \(j_*\mathcal{I}^n\) is \(f_*\)-acyclic:
  \(R^k f_*(j_*\mathcal{I}^n) = 0\) for all \(k \geq 1\). No Serre vanishing on
  \(U \cap f^{-1}(V)\) is needed, and the restriction-of-injectives wall is avoided.

  \emph{(a) \(K\) is exact in positive degrees.} The objectwise homology of the pushed-forward
  resolution \(K = j_*\mathcal{I}^\bullet\) in degree \(k\) computes the right-derived functor of
  \(j_*\) along the injective resolution \(\mathcal{I}^\bullet\): by
  Lemma~\ref{mk-lem:right_derived_injective_resolution},
  \[
    H^k\bigl(j_*\mathcal{I}^\bullet\bigr) \;\cong\; \bigl(R^k j_*\bigr)(\mathcal{H})
      \;=\; R^k j_*\mathcal{H}.
  \]
  By Lemma~\ref{mk-lem:open_immersion_pushforward_acyclic} (\(j\) affine, \(\mathcal{H}\) quasi-coherent
  on the affine \(U\)) this vanishes for \(k \geq 1\); hence \(K\) is exact at every positive degree
  \(n+1\).

  \emph{Augmentation: \(j_*\mathcal{H} \cong K.\mathrm{cycles}\,0\).} The pushforward \(j_*\) is left
  exact, being a right adjoint (Lemma~\ref{mk-lem:scheme_pullbackPushforwardAdjunction_mathlib}) and so
  finite-limit-preserving. For a left-exact functor the zeroth right-derived functor is the functor
  itself, \(R^0 j_* \cong j_*\) (Lemma~\ref{mk-lem:right_derived_zero_iso_self}), and the degree-zero
  cycles of the applied resolution compute \(R^0 j_*\); thus \(j_*\) carries the resolution
  augmentation \(\mathcal{H} \cong (\mathcal{I}^\bullet).\mathrm{cycles}\,0\) to
  \[
    j_*\mathcal{H} \;\cong\; H^0\bigl(j_*\mathcal{I}^\bullet\bigr)
      \;=\; K.\mathrm{cycles}\,0,
  \]
  the augmentation that makes \(K\) a resolution of \(j_*\mathcal{H}\).

  By (a), (b) and the augmentation, \(K = j_*\mathcal{I}^\bullet\) is an \(f_*\)-acyclic resolution
  of \(j_*\mathcal{H}\). The acyclic-resolution comparison
  (Lemma~\ref{mk-lem:acyclic_resolution_computes_derived}) applied with \(G = f_*\) gives, for every
  \(k\), an isomorphism
  \[
    R^k f_*\bigl(j_*\mathcal{H}\bigr) \;\cong\; H^k\bigl(f_*(K)\bigr)
      \;=\; H^k\bigl(f_*(j_*\mathcal{I}^\bullet)\bigr).
  \]

  \emph{Transport: \(H^k(f_*(K)) \cong R^k g_*\mathcal{H}\).} The pushforward-composition isomorphism
  \(\operatorname{pushforwardComp} j\,f : \operatorname{pushforward} j \circ \operatorname{pushforward}
  f \cong \operatorname{pushforward}(j \circ f)\) (Lemma~\ref{mk-lem:scheme_pushforwardComp_mathlib})
  identifies \(f_*(j_*\mathcal{I}^\bullet)\) with \((f \circ j)_*\mathcal{I}^\bullet = g_*\mathcal{I}^
  \bullet\); applying it degreewise and taking homology, while recognising \(H^k(g_*\mathcal{I}^
  \bullet) \cong R^k g_*\mathcal{H}\) via the same injective-resolution identification
  (Lemma~\ref{mk-lem:right_derived_injective_resolution}), yields
  \[
    H^k\bigl(f_*(j_*\mathcal{I}^\bullet)\bigr) \;\cong\; R^k g_*\mathcal{H}.
  \]
  Composing the acyclic-resolution comparison with this transport gives the stated isomorphism
  \(R^k f_*(j_*\mathcal{H}) \cong R^k g_*\mathcal{H}\).
\end{proof}

\begin{lemma}
[The pushforward-sections functor and its additivity]
  \label{mk-lem:pushforward_sections_functor}
  \lean{AlgebraicGeometry.pushforwardSectionsFunctor,
    AlgebraicGeometry.pushforwardSectionsFunctor_additive}
  \dcref{custom}

  \uses{mk-lem:higher_direct_image_presheaf}
  Let \(j : U \to X\) be a morphism and \(W \subseteq X\) an open. The
  \emph{pushforward-sections functor} \(\operatorname{pushforwardSectionsFunctor} j\,W :
  U.\mathrm{Modules} \to \mathrm{Ab}\) sends an \(\mathcal{O}_U\)-module \(M\) to its sections
  \(\Gamma(W, j_* M) = \Gamma(j^{-1}(W), M)\) --- the pushforward along \(j\), forgotten to
  abelian groups, viewed as a presheaf and evaluated at \(W\). It is an additive functor, and
  its \(q\)-th right derived functor applied to \(H\) computes the absolute cohomology
  \(H^q(j^{-1}(W), H)\).
\end{lemma}
\begin{proof}
  \uses{mk-lem:higher_direct_image_presheaf}
  The functor is the composite of pushforward \(\operatorname{pushforward} j\) (additive), the
  forgetful functor to presheaves of abelian groups (additive), and evaluation at \(W\)
  (additive). A composite of additive functors is additive; the additivity follows from the additivity of
  each factor. The identification of its right derived functor with
  absolute cohomology is the presheaf description of
  Lemma~\ref{mk-lem:higher_direct_image_presheaf}.
\end{proof}

\begin{lemma}[The presheaf-of-modules functor is additive]\label{mk-lem:toPresheafOfModules_additive}
  \lean{AlgebraicGeometry.toPresheafOfModules_additive}
  \leanok
  \dcref{custom}

  The functor \(\operatorname{toPresheafOfModules} X : X.\mathrm{Modules} \to
  \mathrm{PresheafOfModules}(\mathcal{O}_X)\), sending a sheaf of \(\mathcal{O}_X\)-modules to its
  underlying presheaf of modules, is additive: it preserves finite biproducts and the additive
  structure on hom-groups.
\end{lemma}
\begin{proof}
  \leanok

The functor is the composite of the forgetful functor
  \(\mathrm{SheafOfModules}(\mathcal{O}_X) \to \mathrm{PresheafOfModules}(\mathcal{O}_X)\) and the
  restriction of scalars along the identity \(\operatorname{restrictScalars}(\mathrm{id}_{\mathcal{O}_X})\).
  Both factors are additive, and a composite of additive functors is additive.
\end{proof}

\begin{lemma}[The sections functor is additive]\label{mk-lem:sectionsFunctor_additive}
  \lean{AlgebraicGeometry.sectionsFunctor_additive}
  \leanok
  \dcref{custom}

  \uses{mk-lem:toPresheafOfModules_additive}
  For an open \(V \subseteq X\) the global-sections functor
  \(\operatorname{sectionsFunctor} V = \Gamma(V, -) : X.\mathrm{Modules} \to \mathrm{Ab}\) is
  additive.
\end{lemma}
\begin{proof}
  \leanok
  \uses{mk-lem:toPresheafOfModules_additive}
  The functor factors as the additive functor to presheaves of modules
  (Lemma~\ref{mk-lem:toPresheafOfModules_additive}), followed by the underlying-presheaf functor and
  evaluation at \(V\), each of which is additive; the composite is therefore additive.
\end{proof}

\begin{lemma}[The sections functor is corepresented by \(j_!\mathcal{O}_V\)]\label{mk-lem:sectionsFunctorCorepIso}
  \lean{AlgebraicGeometry.sectionsFunctorCorepIso}
  \leanok
  \dcref{custom}

  \uses{mk-lem:jshriek_corepr, mk-lem:sectionsFunctor_additive}
  For an open \(V \subseteq X\) there is a natural isomorphism of additive functors
  \[
    \operatorname{sectionsFunctor} V
      \;\cong\;
    \operatorname{Hom}_{X.\mathrm{Modules}}\bigl(j_!\mathcal{O}_V,\, -\bigr)
  \]
  between \(\Gamma(V, -)\) and the covariant additive hom-functor corepresented by
  \(j_!\mathcal{O}_V\). It is the functorial upgrade of the corepresentability bijection of
  Lemma~\ref{mk-lem:jshriek_corepr}.
\end{lemma}
\begin{proof}
  \leanok
  \uses{mk-lem:jshriek_corepr, mk-lem:sectionsFunctor_additive}
  Assemble the natural transformation from the additive corepresentability isomorphism
  \(\operatorname{Hom}(j_!\mathcal{O}_V, \mathcal{F}) \cong \Gamma(V, \mathcal{F})\) of
  Lemma~\ref{mk-lem:jshriek_corepr}, taken componentwise; its naturality square in \(\mathcal{F}\) is
  exactly the naturality of that bijection in the coefficient sheaf. The components are
  isomorphisms, so the assembled transformation is a natural isomorphism. Both functors are
  additive (Lemma~\ref{mk-lem:sectionsFunctor_additive} and the additivity of a hom-functor).
\end{proof}

\begin{lemma}[Right derived functors transport along a natural isomorphism]\label{mk-lem:rightDerivedNatIso}
  \lean{AlgebraicGeometry.rightDerivedNatIso}
  \leanok
  \dcref{custom}

  Let \(\mathcal{A}\) be an abelian category with enough injectives, \(\mathcal{B}\) abelian, and
  \(F, G : \mathcal{A} \to \mathcal{B}\) additive functors. A natural isomorphism \(F \cong G\)
  induces, for every \(n\), a natural isomorphism of right derived functors
  \[
    F.\mathrm{rightDerived}\,n \;\cong\; G.\mathrm{rightDerived}\,n.
  \]
\end{lemma}
\begin{proof}
  \leanok

Right derivation is functorial in the functor argument: a natural transformation \(F \to G\)
  induces a natural transformation \(F.\mathrm{rightDerived}\,n \to G.\mathrm{rightDerived}\,n\),
  compatibly with composition and identities. Applying this to the hom and inverse of the given
  isomorphism \(F \cong G\) yields two derived natural transformations, and the composition and
  identity compatibilities show they are mutually inverse; hence
  \(F.\mathrm{rightDerived}\,n \cong G.\mathrm{rightDerived}\,n\).
\end{proof}

\begin{lemma}[The canonical isomorphism of an affine scheme with its spectrum]
  \label{mk-lem:isoSpec_scheme_mathlib}
  \lean{AlgebraicGeometry.Scheme.isoSpec}
  \mathlibok
  \dcref{custom}

  Let \(U\) be a scheme with \([\,\mathrm{IsAffine}\ U\,]\). Then there is a canonical
  \emph{whole-scheme} isomorphism of schemes
  \[
    U \;\cong\; \operatorname{Spec}\Gamma(U, \mathcal{O}_U),
  \]
  identifying \(U\) with the spectrum of its ring of global sections. This is a genuine isomorphism of
  schemes (not merely an open immersion), so it induces an equivalence of module categories.
\end{lemma}

\begin{lemma}[Ext transports along an exact functor]
  \label{mk-lem:ext_mapExactFunctor_mathlib}
  \lean{CategoryTheory.Abelian.Ext.mapExactFunctor}
  \mathlibok
  \dcref{custom}

  Let \(F : \mathcal{A} \to \mathcal{B}\) be an exact functor between abelian categories. Then for
  objects \(X, Y\) of \(\mathcal{A}\) and every \(n\) there is an induced additive map
  \[
    \operatorname{Ext}^n_{\mathcal{A}}(X, Y) \;\longrightarrow\;
      \operatorname{Ext}^n_{\mathcal{B}}(F X, F Y),
  \]
  compatible with the additive structure, with degree-zero \(\operatorname{Hom}\)-action, and with
  Yoneda composition. When \(F\) is one half of an equivalence the maps for \(F\) and its inverse are
  mutually inverse, so \(\operatorname{Ext}\) is carried isomorphically across the equivalence.
\end{lemma}

\begin{lemma}[Pushforward of sheaves of modules and its coherences]
  \label{mk-lem:modules_pushforward_mathlib}
  \lean{AlgebraicGeometry.Scheme.Modules.pushforward}
  \mathlibok
  \dcref{custom}

  For a morphism of schemes \(f : X \to Y\) there is an additive pushforward functor
  \(\operatorname{pushforward} f : X.\mathrm{Modules} \to Y.\mathrm{Modules}\), together with the
  coherence isomorphisms \(\operatorname{pushforward}(\mathrm{id}_X) \cong \operatorname{id}\),
  \(\operatorname{pushforward} f \circ \operatorname{pushforward} g \cong
  \operatorname{pushforward}(f \circ g)\), and \(\operatorname{pushforward} f \cong
  \operatorname{pushforward} g\) for \(f = g\). When \(\varphi : X \cong Y\) is an isomorphism, these
  assemble \(\operatorname{pushforward}\varphi.\mathrm{hom}\) and
  \(\operatorname{pushforward}\varphi.\mathrm{inv}\) into an equivalence
  \(X.\mathrm{Modules} \simeq Y.\mathrm{Modules}\).
\end{lemma}

\subsection{The enlarged affine instantiation (general affine opens, Need~\#2)}

The distinguished-open affine cover system (Definition~\ref{mk-def:affine_cover_system}) computes Serre
vanishing only on distinguished opens \(D(f)\). To reach \emph{every} affine open \(V\) we enlarge the
basis \(\mathcal{B}\) to all affine opens, keeping the admissible coverings \(\mathrm{Cov}\) the standard
covers \(i \mapsto D(g_i)\). The three cover-system fields and the quasi-coherent {\v C}ech-vanishing seed
are re-discharged below in their general-affine-open form; the only genuinely new content is the seed
(Lemma~\ref{mk-lem:affine_cech_vanishing_general_seed}), whose change-of-ring discharge replaces the single
localisation \(R_f\) of the distinguished case by the coordinate ring \(\Gamma(V)\).

\begin{lemma}[Distinguished opens of \(\operatorname{Spec} R\) are affine]\label{mk-lem:isAffineOpen_specBasicOpen}
  \lean{AlgebraicGeometry.isAffineOpen_specBasicOpen}
  \leanok
  \dcref{custom}

  \textit{Project-local infrastructure.}
  For a ring \(R\) and \(r \in R\), the distinguished open \(D(r) \subseteq \operatorname{Spec} R\) is an
  affine open: it is isomorphic, as a scheme, to \(\operatorname{Spec} R[1/r]\). Consequently every
  distinguished open lies in the enlarged basis of all affine opens, which discharges the
  \(\mathrm{faces\_mem}\) field of the enlarged affine cover system
  (Definition~\ref{mk-def:affine_cover_system_general}).
\end{lemma}
\begin{proof}
  \leanok


The canonical open immersion \(D(r) \cong \operatorname{Spec} R[1/r] \hookrightarrow \operatorname{Spec}
  R\) (the basic-open / away-localisation isomorphism \(\operatorname{basicOpenIsoSpecAway}\)) is an
  isomorphism onto \(D(r)\); since its source is affine and affineness transports across a scheme
  isomorphism, \(D(r)\) is affine.
\end{proof}

\begin{lemma}[Standard covers are cofinal among open covers of a general affine open]\label{mk-lem:standard_cover_cofinal_affine}
  \lean{AlgebraicGeometry.standard_cover_cofinal_affine}
  \leanok
  \dcref{009L,custom}

  \uses{mk-def:standard_affine_cover, mk-lem:primeSpectrum_isBasis_basicOpen, mk-lem:affine_faces_mem}
  \textit{Source: Stacks Project, Sheaves on Spaces, Tag 009L,
  \texttt{lemma-cofinal-systems-coverings-standard-case}.}
  Let \(V\) be \emph{any} affine open of \(\operatorname{Spec} R\), and let \(\{W_\alpha\}_{\alpha \in A}\)
  be an arbitrary open covering of \(V\) (each \(W_\alpha\) an open of \(\operatorname{Spec} R\) contained
  in \(V\)). Then there exist a finite \(n\), a family \(g : \{1, \dots, n\} \to R\), and a reindexing
  \(\varphi : \{1, \dots, n\} \to A\) such that
  \[
    V = \bigcup_{i=1}^n D(g_i)
    \qquad\text{and}\qquad
    D(g_i) \subseteq W_{\varphi(i)} \quad (1 \le i \le n).
  \]
  That is, every open covering of an arbitrary affine open \(V\) is refined by a finite standard cover by
  distinguished opens, each sitting inside a member of the original covering. This is the
  general-affine-open companion of Lemma~\ref{mk-lem:standard_cover_cofinal}, the only change being the source
  of quasi-compactness: the compactness of the affine open \(V\) (every affine scheme is quasi-compact) in
  place of the compactness of a distinguished \(D(f)\).
\end{lemma}
\begin{proof}
  \leanok
  \uses{mk-def:standard_affine_cover, mk-lem:primeSpectrum_isBasis_basicOpen, mk-lem:affine_faces_mem}
  The basic opens \(D(g)\) form a topological basis of \(\operatorname{Spec} R\)
  (Lemma~\ref{mk-lem:primeSpectrum_isBasis_basicOpen}), so each \(W_\alpha\) is a union of basic opens it
  contains; choosing one basic open inside a member through each point of \(V\) yields a basic-open covering
  of \(V\) refining \(\{W_\alpha\}\). An affine open \(V\) is quasi-compact, so a finite subfamily
  \(D(g_1), \dots, D(g_n)\) already covers \(V\); recording for each \(i\) an index \(\varphi(i)\) with
  \(D(g_i) \subseteq W_{\varphi(i)}\) gives the asserted finite refinement \(V = \bigcup_i D(g_i)\). The
  basis is closed under the intersections required by Tag 009L (Lemma~\ref{mk-lem:affine_faces_mem}), so the
  standard covers form a cofinal system among the open coverings of \(V\).
\end{proof}

\begin{lemma}[Section surjectivity for the enlarged affine cover system]\label{mk-lem:affine_surj_of_vanishing_affine}
  \lean{AlgebraicGeometry.affine_surj_of_vanishing_affine}
  \leanok
  \dcref{custom}

  \uses{mk-lem:ses_cech_h1, mk-lem:standard_cover_cofinal_affine, mk-lem:to_sheaf_preserves_epi,
    mk-lem:presheaf_is_locally_surjective, mk-lem:sheaf_locally_surjective_iff_epi}
  \textit{Project-local infrastructure.}
  Let \(V\) be \emph{any} affine open of \(\operatorname{Spec} R\) and let
  \(S : 0 \to S_1 \to S_2 \to S_3 \to 0\) be a short exact sequence of \(\mathcal{O}_X\)-modules whose
  left term \(S_1\) has vanishing positive {\v C}ech cohomology over every standard cover
  \(i \mapsto D(g_i)\) \emph{whose union is affine}. Then the section map \(S_2(V) \to S_3(V)\) is
  surjective. This is the \(\mathrm{surj\_of\_vanishing}\) field of the enlarged affine cover system, the
  general-affine-open companion of Lemma~\ref{mk-lem:affine_surj_of_vanishing}.
\end{lemma}
\begin{proof}
  \leanok
  \uses{mk-lem:ses_cech_h1, mk-lem:standard_cover_cofinal_affine, mk-lem:to_sheaf_preserves_epi,
    mk-lem:presheaf_is_locally_surjective, mk-lem:sheaf_locally_surjective_iff_epi}
    The proof of Lemma~\ref{mk-lem:affine_surj_of_vanishing} carries over verbatim. A module epimorphism
  \(S_2 \to S_3\) is locally surjective on sections (Lemmas~\ref{mk-lem:to_sheaf_preserves_epi},
  \ref{mk-lem:sheaf_locally_surjective_iff_epi}, \ref{mk-lem:presheaf_is_locally_surjective}); the resulting
  open covering of \(V\) carrying local lifts of a section \(t \in S_3(V)\) is refined to a finite standard
  cover \(i \mapsto D(g_i)\) of \(V\) by Lemma~\ref{mk-lem:standard_cover_cofinal_affine} (this is where the
  quasi-compactness of the affine open \(V\) replaces that of \(D(f)\)), whose union equals \(V\) and is
  therefore affine. The {\v C}ech \(H^1\)-gluing sequence (Lemma~\ref{mk-lem:ses_cech_h1}), fed the vanishing
  \(\check{H}^1(\{D(g_i)\}, S_1) = 0\) for that affine-union standard cover, glues the local lifts to a
  global lift of \(t\) over \(V\).
\end{proof}

\begin{definition}[The enlarged affine cover system]\label{mk-def:affine_cover_system_general}
  \lean{AlgebraicGeometry.affineCoverSystemGeneral}
  \leanok
  \dcref{custom}

  \uses{mk-def:basis_cov_system, mk-lem:affine_faces_mem, mk-lem:isAffineOpen_specBasicOpen,
    mk-lem:affine_surj_of_vanishing_affine, mk-lem:injective_cech_acyclic}
  \textit{Project-local: enlargement of Definition~\ref{mk-def:affine_cover_system}.}
  The \emph{enlarged affine cover system} for \(\operatorname{Spec} R\) is the instance of
  Definition~\ref{mk-def:basis_cov_system} whose basis \(\mathcal{B}\) is \emph{all} affine opens of
  \(\operatorname{Spec} R\) and whose admissible coverings \(\mathrm{Cov}\) are the standard finite covers
  \(i \mapsto D(g_i)\) whose union \(\bigcup_i D(g_i)\) is affine. It generalises
  Definition~\ref{mk-def:affine_cover_system}, which uses only the distinguished opens \(D(f)\) as its basis.
  Its three proof fields are discharged by: \(\mathrm{faces\_mem}\) --- each face
  \(\bigcap_k D(g_{\sigma k}) = D(\prod_k g_{\sigma k})\) is distinguished, hence affine
  (Lemma~\ref{mk-lem:isAffineOpen_specBasicOpen}), so lies in the enlarged basis
  (Lemma~\ref{mk-lem:affine_faces_mem}); \(\mathrm{surj\_of\_vanishing}\) by
  Lemma~\ref{mk-lem:affine_surj_of_vanishing_affine}; and \(\mathrm{injective\_acyclic}\) by the cover-agnostic
  family-form injective {\v C}ech-acyclicity (Lemma~\ref{mk-lem:injective_cech_acyclic}).
\end{definition}

\begin{lemma}[Localised base change is a localisation away]\label{mk-lem:isLocalizedModule_baseChange_away}
  \lean{AlgebraicGeometry.isLocalizedModule_baseChange_away}
  \leanok
  \dcref{custom}

  \uses{mk-lem:isLocalizedModule_comp_away}
  \textit{Project-local infrastructure, assembled from Mathlib base-change primitives.}
  Let \(\varphi : R \to S\) be a ring map, \(M\) an \(R\)-module, and \(g \in R\) with image
  \(\bar g = \varphi(g) \in S\). The base change \(M \otimes_R S\), further localised at the powers of
  \(\bar g\), is a localisation of \(M\): the composite
  \[
    M \;\longrightarrow\; M \otimes_R S \;\longrightarrow\; (M \otimes_R S)_{\bar g}
  \]
  exhibits \((M \otimes_R S)_{\bar g}\) as the base change of \(M\) along the composite
  \(R \to S \to S_{\bar g}\). In particular, whenever the localised ring \(S_{\bar g}\) is identified
  with the localisation \(R_g\) --- as happens when \(D_S(\bar g)\) and \(D_R(g)\) name the same basic
  open --- the composite presents \((M \otimes_R S)_{\bar g}\) as the away localisation \(M_g\).
\end{lemma}
\begin{proof}
  \leanok
  \uses{mk-lem:isLocalizedModule_comp_away}
  Tensoring with \(S\) exhibits \(M \otimes_R S\) as the base change of \(M\) along \(R \to S\)
  (\(\operatorname{TensorProduct.isBaseChange}\)); localising at the powers of \(\bar g\) exhibits
  \((M \otimes_R S)_{\bar g}\) as the base change of \(M \otimes_R S\) along \(S \to S_{\bar g}\)
  (\(\operatorname{isLocalizedModule\_iff\_isBaseChange}\)). Composing the two base changes
  (\(\operatorname{IsBaseChange.comp}\)) exhibits the composite \(M \to (M \otimes_R S)_{\bar g}\) as the
  base change of \(M\) along \(R \to S_{\bar g}\). Under the identification \(S_{\bar g} = R_g\) this is
  exactly the away localisation of \(M\) at \(g\); the two-step composite specialises to the localised
  module recipe already recorded in Lemma~\ref{mk-lem:isLocalizedModule_comp_away}.
\end{proof}

\begin{lemma}[Section {\v C}ech module exactness from an affine-cover base-change datum]
  \label{mk-lem:section_cech_module_exact_of_affineCover}
  \lean{AlgebraicGeometry.SectionCechModule.dDiff_exact_of_affineCover,
    AlgebraicGeometry.sectionCechAbExact_affine}
  \dcref{custom}

  \uses{mk-lem:section_cech_module_exact, mk-lem:isLocalizedModule_baseChange_away}
  \textit{Project-local infrastructure; mirror of
  Lemma~\ref{mk-lem:section_cech_module_exact_of_localizationAway}.}
  Let \(M\) be an \(R\)-module and \(g : \iota \to R\) a finite family whose union
  \(V = \bigcup_i D(g_i)\) is an affine open of \(\operatorname{Spec} R\), with coordinate ring
  \(S = \Gamma(V, \mathcal{O}_V)\) and restriction map \(\varphi : R \to S\), \(\bar g_i = \varphi(g_i)\).
  Since the \(D_S(\bar g_i)\) cover \(\operatorname{Spec} S = V\), the images \(\bar g_i\) span the unit
  ideal of \(S\). Then the un-localised section {\v C}ech module complex \(D^\bullet(g, M)\) of
  Definition~\ref{mk-def:section_cech_module_complex} is exact in every positive degree --- \emph{without}
  assuming that \(g\) spans the unit ideal of \(R\) itself.
\end{lemma}
\begin{proof}
  \uses{mk-lem:section_cech_module_exact, mk-lem:isLocalizedModule_baseChange_away}
  Run the computation over \(S\). Instantiate the polymorphic positive-degree exactness of
  Lemma~\ref{mk-lem:section_cech_module_exact} over the base ring \(S\), with module \(M \otimes_R S\) and
  the spanning family \(\bar g : \iota \to S\); spanning is exactly the affineness of the cover's union,
  so the \(S\)-side complex \(D^\bullet(\bar g, M \otimes_R S)\) is exact in positive degrees. Transport
  this exactness back to the \(R\)-side along the degreewise \(R\)-linear isomorphisms
  \(M_{g_\sigma} \cong (M \otimes_R S)_{\bar g_\sigma}\) of
  Lemma~\ref{mk-lem:isLocalizedModule_baseChange_away} (for a tuple \(\sigma\),
  \(g_\sigma = \prod_k g_{\sigma k}\) and \(\bar g_\sigma = \varphi(g_\sigma)\), with
  \(\Gamma(D(g_\sigma))\) the common localisation). Realised through \(\operatorname{IsLocalizedModule.iso}\),
  these comparisons are the vertical maps of a ladder, and exactness transfers across it by
  \(\operatorname{Function.Exact.of\_ladder\_addEquiv\_of\_exact}\), the cofaces commuting by naturality
  of the localisation comparison.
\end{proof}

\begin{lemma}[Tilde {\v C}ech vanishing over an affine-union standard cover]\label{mk-lem:sectionCech_homology_exact_of_affineCover}
  \lean{AlgebraicGeometry.sectionCech_homology_exact_of_affineCover}
  \leanok
  \dcref{custom}

  \uses{mk-lem:section_cech_module_exact_of_affineCover, mk-lem:section_cech_ab_exact,
    mk-lem:section_cech_homology_exact}
  \textit{Project-local infrastructure; mirror of
  Lemma~\ref{mk-lem:affine_cech_vanishing_tilde_subcover}.}
  Let \(R\) be a ring, \(M\) an \(R\)-module, and \(g : \iota \to R\) a finite family whose union
  \(V = \bigcup_i D(g_i)\) is an affine open of \(\operatorname{Spec} R\). Then for every \(p > 0\) the
  homology of the section {\v C}ech complex of \(\widetilde{M}\) over \(\{D(g_i)\}\) vanishes,
  \[
    \operatorname{IsZero}\bigl(\check{\mathcal{C}}^\bullet(\{D(g_i)\}, \widetilde{M}).\mathrm{homology}\, p\bigr).
  \]
\end{lemma}
\begin{proof}
  \leanok
  \uses{mk-lem:section_cech_module_exact_of_affineCover, mk-lem:section_cech_ab_exact,
    mk-lem:section_cech_homology_exact}
  Lemma~\ref{mk-lem:section_cech_module_exact_of_affineCover} supplies positive-degree exactness of the
  un-localised section {\v C}ech module complex \(D^\bullet(g, M)\) from the affineness of the cover's
  union. The abelian-group exactness criterion (the private \(\operatorname{sectionCechAbExact\_affine}\)
  packaging of Lemma~\ref{mk-lem:section_cech_ab_exact}) transports this module exactness across the
  degreewise additive ladder identifying the section complex with the localised-module complex
  (Lemma~\ref{mk-lem:section_cech_homology_exact}), and reads the resulting function exactness off as the
  vanishing of the degree-\(p\) homology via
  \(\operatorname{sectionCech\_isZero\_homology\_of\_objD\_exact}\).
\end{proof}

\begin{lemma}

[The general-affine-open change-of-ring {\v C}ech-vanishing seed]
  \label{mk-lem:affine_cech_vanishing_general_seed}
  \lean{AlgebraicGeometry.sectionCech_homology_exact_of_affineOpen,
    AlgebraicGeometry.basicOpen_algMap_section,
    AlgebraicGeometry.affine_tildeVanishing_general}
  \dcref{01HV,custom}

  \uses{mk-lem:isoSpec_scheme_mathlib, mk-lem:cech_acyclic_affine, mk-lem:affine_cech_vanishing_qcoh,
    mk-lem:isAffineOpen_specBasicOpen, mk-lem:sectionCech_homology_exact_of_affineCover,
    mk-lem:isLocalizedModule_baseChange_away}
  \textit{Source: Stacks Project, Schemes, Tag 01HV, \texttt{lemma-spec-sheaves}~(4), and
  Cohomology of Schemes, \texttt{lemma-cech-cohomology-quasi-coherent-trivial}.}
  Let \(R\) be a ring, \(M\) an \(R\)-module, and \(g : \{1, \dots, n\} \to R\) a finite family such that
  \(V := \bigcup_i D(g_i)\) is an affine open of \(\operatorname{Spec} R\) (a general affine open, \emph{not}
  assumed to be a single \(D(f)\)). Then for every \(p > 0\) the section {\v C}ech complex of
  \(\widetilde{M}\) over the cover \(\{D(g_i)\}\) has vanishing positive homology:
  \[
    \check{H}^p\bigl(\{D(g_i)\}, \widetilde{M}\bigr) = 0.
  \]
\end{lemma}
\begin{proof}
  \uses{mk-lem:isoSpec_scheme_mathlib, mk-lem:cech_acyclic_affine, mk-lem:affine_cech_vanishing_qcoh,
    mk-lem:isAffineOpen_specBasicOpen, mk-lem:sectionCech_homology_exact_of_affineCover,
    mk-lem:isLocalizedModule_baseChange_away}
  Because \(V\) is a proper affine open, the family \(\{g_i\}\) need not span the unit ideal of \(R\), so
  the standard-cover vanishing of Lemma~\ref{mk-lem:cech_acyclic_affine} does not apply over \(R\). As in the
  distinguished case (Lemma~\ref{mk-lem:affine_cech_vanishing_tilde_subcover}) we run the computation over a
  ring where the cover \emph{does} span the unit ideal --- here the coordinate ring of \(V\) itself, rather
  than a single localisation \(R_f\).

  \emph{Change of ring to \(S = \Gamma(V, \mathcal{O}_V)\).} Set \(S := \Gamma(V, \mathcal{O}_V)\), let
  \(\varphi : R \to S\) be the restriction map and \(\bar g_i := \varphi(g_i)\). Since \(V\) is affine,
  \(\operatorname{Scheme.isoSpec}\) (Lemma~\ref{mk-lem:isoSpec_scheme_mathlib}) gives a scheme isomorphism
  \(V \cong \operatorname{Spec} S\) under which each \(D(g_i) \subseteq V\) (an affine open, by
  Lemma~\ref{mk-lem:isAffineOpen_specBasicOpen}) corresponds to the distinguished open \(D_S(\bar g_i)\) of
  \(\operatorname{Spec} S\). The \(D_S(\bar g_i)\) cover \(\operatorname{Spec} S = V\), so the \(\bar g_i\)
  \emph{span the unit ideal of \(S\)}.

  \emph{Identification of the two section complexes.} The section {\v C}ech complex of \(\widetilde{M}\)
  over \(\{D(g_i)\}\) has degree-\(p\) term \(\prod_\sigma \Gamma(D(g_\sigma), \widetilde{M}) = \prod_\sigma
  M_{g_\sigma}\) (Tag 01HV item (4), \(g_\sigma = \prod_k g_{\sigma k}\)). We identify it, degreewise, with
  the standard-cover (full-span) section {\v C}ech complex over \(\operatorname{Spec} S\) of
  \(\widetilde{M \otimes_R S}\), whose degree-\(p\) term is \(\prod_\sigma (M \otimes_R S)_{\bar g_\sigma}\).
  The per-\(\sigma\) comparison \(M_{g_\sigma} \cong (M \otimes_R S)_{\bar g_\sigma}\) is a localised
  base-change isomorphism: \(\Gamma(D(g_\sigma)) = R_{g_\sigma}\) is simultaneously the localisation of \(R\)
  away from \(g_\sigma\) and the localisation of \(S\) away from \(\bar g_\sigma\) (since
  \((\operatorname{Spec} R).\operatorname{basicOpen}\bar g_\sigma = D(g_\sigma)\)), and \((M \otimes_R
  S)_{\bar g_\sigma}\) is the base change of \(M\) along the composite \(R \to S \to S_{\bar g_\sigma}\;\)
  (transitivity of base change for localised modules). These comparisons are the vertical maps of an
  isomorphism of cochain complexes, commuting with the alternating-sum cofaces by naturality of the
  localisation comparison.

  \emph{Conclusion.} Over \(S\) the family \(\bar g\) spans the unit ideal, so the full-span standard-cover
  vanishing --- the polymorphic core of Lemma~\ref{mk-lem:cech_acyclic_affine}, re-instantiated over \(S\) with
  the module \(M \otimes_R S\), exactly as it supplies the quasi-coherent seed of
  Lemma~\ref{mk-lem:affine_cech_vanishing_qcoh} on a whole spectrum --- gives positive-degree exactness of the
  \(\operatorname{Spec} S\) complex. This change-of-base computation and its transport back along the
  per-\(\sigma\) base-change isomorphisms \(M_{g_\sigma} \cong (M \otimes_R S)_{\bar g_\sigma}\)
  (Lemma~\ref{mk-lem:isLocalizedModule_baseChange_away}) is exactly
  Lemma~\ref{mk-lem:sectionCech_homology_exact_of_affineCover}, whose conclusion is the asserted vanishing
  \(\check{H}^p(\{D(g_i)\}, \widetilde{M}) = 0\) for all \(p > 0\).
\end{proof}

\begin{lemma}[General-affine-open quasi-coherent {\v C}ech vanishing from the tilde seed]\label{mk-lem:affine_cech_vanishing_general_of_tildeVanishing}
  \lean{AlgebraicGeometry.affine_cech_vanishing_qcoh_general_of_tildeVanishing}
  \leanok
  \dcref{custom}

  \uses{mk-lem:qcoh_iso_tilde_sections, mk-lem:affine_cech_vanishing_general_seed}
  \textit{Project-local: enlarged-system analogue of the reduction inside
  Lemma~\ref{mk-lem:affine_cech_vanishing_qcoh}.}
  Let \(\mathcal{F}\) be a quasi-coherent \(\mathcal{O}_{\operatorname{Spec} R}\)-module with
  \(M = \Gamma(\operatorname{Spec} R, \mathcal{F})\). Then \(\mathcal{F}\) has vanishing higher {\v C}ech
  cohomology for the enlarged affine cover system (Definition~\ref{mk-def:affine_cover_system_general}): for
  every standard cover \(i \mapsto D(g_i)\) whose union is affine and every \(p > 0\),
  \(\check{H}^p(\{D(g_i)\}, \mathcal{F}) = 0\). The reduction is transport along the affine structure
  isomorphism \(\mathcal{F} \cong \widetilde{M}\) (Lemma~\ref{mk-lem:qcoh_iso_tilde_sections}) to the tilde
  seed of Lemma~\ref{mk-lem:affine_cech_vanishing_general_seed}.
\end{lemma}
\begin{proof}
  \leanok
  \uses{mk-lem:qcoh_iso_tilde_sections, mk-lem:affine_cech_vanishing_general_seed}
  By Lemma~\ref{mk-lem:qcoh_iso_tilde_sections} the quasi-coherent \(\mathcal{F}\) is isomorphic to
  \(\widetilde{M}\). {\v C}ech cohomology transports along an isomorphism of coefficients
  (Lemma~\ref{mk-lem:cechCohomology_isZero_of_iso}), so it suffices to vanish
  \(\check{H}^p(\{D(g_i)\}, \widetilde{M})\) for every standard cover whose union is affine and every
  \(p > 0\) --- which is exactly Lemma~\ref{mk-lem:affine_cech_vanishing_general_seed}.
\end{proof}

\begin{lemma}[General-affine Serre vanishing from the enlarged {\v C}ech seed]\label{mk-lem:affine_serre_vanishing_general_of_seed}
  \lean{AlgebraicGeometry.affine_serre_vanishing_general_of_seed}
  \leanok
  \dcref{custom}

  \uses{mk-lem:cech_to_cohomology_on_basis, mk-def:affine_cover_system_general}
  \textit{Project-local: enlarged-system instantiation of
  Lemma~\ref{mk-lem:cech_to_cohomology_on_basis}.}
  Let \(\mathcal{F}\) be a quasi-coherent \(\mathcal{O}_{\operatorname{Spec} R}\)-module which has vanishing
  higher {\v C}ech cohomology for the enlarged affine cover system
  (Definition~\ref{mk-def:affine_cover_system_general}). Then for every affine open \(V\) of
  \(\operatorname{Spec} R\) and every \(p > 0\) the absolute cohomology
  \(H^p(V, \mathcal{F}) = \operatorname{Ext}^p(j_!\mathcal{O}_V, \mathcal{F})\) vanishes. The statement
  carries \([\,\mathrm{EnoughInjectives}\ (\operatorname{Spec} R).\mathrm{Modules}\,]\), inherited from
  Lemma~\ref{mk-lem:cech_to_cohomology_on_basis}.
\end{lemma}
\begin{proof}
  \leanok
  \uses{mk-lem:cech_to_cohomology_on_basis, mk-def:affine_cover_system_general}
  Apply the basis-comparison criterion (Lemma~\ref{mk-lem:cech_to_cohomology_on_basis}) to the enlarged
  affine cover system (Definition~\ref{mk-def:affine_cover_system_general}), whose basis is all affine opens.
  Conditions (1) and (2) are the system's \(\mathrm{faces\_mem}\) and cofinality data, and condition (3) is
  the supplied seed. Reading off the conclusion at the member \(V \in \mathcal{B}\) gives
  \(H^p(V, \mathcal{F}) = 0\) for \(p > 0\).
\end{proof}

\begin{lemma}[General-affine Serre vanishing from the tilde leaf]\label{mk-lem:affine_serre_vanishing_general_of_tildeVanishing}
  \lean{AlgebraicGeometry.affine_serre_vanishing_general_of_tildeVanishing}
  \leanok
  \dcref{custom}

  \uses{mk-lem:affine_cech_vanishing_general_of_tildeVanishing,
    mk-lem:affine_serre_vanishing_general_of_seed}
  \textit{Project-local: the fully-reduced working form behind
  Lemma~\ref{mk-lem:affine_serre_vanishing_general_open}.}
  Let \(\mathcal{F}\) be a quasi-coherent \(\mathcal{O}_{\operatorname{Spec} R}\)-module and \(V\) any affine
  open of \(\operatorname{Spec} R\). Then \(\operatorname{Ext}^p(j_!\mathcal{O}_V, \mathcal{F}) = 0\) for all
  \(p > 0\). This composes the seed reduction
  (Lemma~\ref{mk-lem:affine_cech_vanishing_general_of_tildeVanishing}) with the basis-comparison
  instantiation (Lemma~\ref{mk-lem:affine_serre_vanishing_general_of_seed}).
\end{lemma}
\begin{proof}
  \leanok
  \uses{mk-lem:affine_cech_vanishing_general_of_tildeVanishing,
    mk-lem:affine_serre_vanishing_general_of_seed}
    Lemma~\ref{mk-lem:affine_cech_vanishing_general_of_tildeVanishing} supplies the enlarged-system seed
  \(\operatorname{HasVanishingHigherCech}(\operatorname{affineCoverSystemGeneral} R)\, \mathcal{F}\) (its
  tilde-leaf hypothesis being discharged by Lemma~\ref{mk-lem:affine_cech_vanishing_general_seed}); feeding
  this seed to Lemma~\ref{mk-lem:affine_serre_vanishing_general_of_seed} yields
  \(\operatorname{Ext}^p(j_!\mathcal{O}_V, \mathcal{F}) = 0\) for \(p > 0\).
\end{proof}

\begin{lemma}[Serre vanishing for a general affine open (Need~\#2)]\label{mk-lem:affine_serre_vanishing_general_open}
  \lean{AlgebraicGeometry.affine_serre_vanishing_general_open}
  \leanok
  \dcref{custom}

  \uses{mk-lem:cech_to_cohomology_on_basis, mk-lem:affine_surj_of_vanishing_affine,
    mk-lem:standard_cover_cofinal_affine, mk-def:affine_cover_system_general, mk-lem:affine_faces_mem,
    mk-lem:injective_cech_acyclic, mk-lem:affine_cech_vanishing_general_of_tildeVanishing,
    mk-lem:affine_cech_vanishing_general_seed, mk-def:absolute_cohomology}
  Let \(R\) be a ring, \(V\) \emph{any} affine open of \(\operatorname{Spec} R\), and
  \(\mathcal{H}\) a quasi-coherent \(\mathcal{O}_{\operatorname{Spec} R}\)-module. Then for every
  \(q \geq 1\) the absolute cohomology vanishes:
  \[
    H^q(V, \mathcal{H}) \;=\; \operatorname{Ext}^q(j_!\mathcal{O}_V, \mathcal{H}) \;=\; 0.
  \]
  This is the general-affine-open companion of the distinguished-open Serre vanishing
  (Lemma~\ref{mk-lem:affine_serre_vanishing}): the basis \(\mathcal{B}\) of the affine cover system is
  enlarged from the distinguished opens \(D(f)\) to all affine opens, while the admissible coverings
  \(\mathrm{Cov}\) remain the standard covers \(i \mapsto D(g_i)\).
\end{lemma}
\begin{proof}
  \leanok
  \uses{mk-lem:cech_to_cohomology_on_basis, mk-lem:affine_surj_of_vanishing_affine,
    mk-lem:standard_cover_cofinal_affine, mk-def:affine_cover_system_general, mk-lem:affine_faces_mem,
    mk-lem:injective_cech_acyclic, mk-lem:affine_cech_vanishing_general_of_tildeVanishing,
    mk-lem:affine_cech_vanishing_general_seed, mk-def:absolute_cohomology}
    Instantiate the basis-comparison criterion (Lemma~\ref{mk-lem:cech_to_cohomology_on_basis}) at the
  enlarged affine cover system (Definition~\ref{mk-def:affine_cover_system_general}), whose basis
  \(\mathcal{B}\) is \emph{all} affine opens of \(\operatorname{Spec} R\) and whose admissible coverings
  \(\mathrm{Cov}\) are the standard covers of an affine open \(V = \bigcup_i D(g_i)\). Its three
  cover-system fields are:
  \begin{itemize}
    \item \(\mathrm{faces\_mem}\): the faces of a standard cover are the distinguished opens
      \(\bigcap_k D(g_{\sigma k}) = D\bigl(\prod_k g_{\sigma k}\bigr)\), still distinguished and a
      fortiori affine, so they lie in the enlarged \(\mathcal{B}\) (Lemmas~\ref{mk-lem:affine_faces_mem},
      \ref{mk-lem:isAffineOpen_specBasicOpen}).
    \item \(\mathrm{injective\_acyclic}\): the family-form injective {\v C}ech-acyclicity
      (Lemma~\ref{mk-lem:injective_cech_acyclic}) is cover-agnostic and applies to the standard covers of
      any affine open verbatim.
    \item \(\mathrm{surj\_of\_vanishing}\): the general-affine-open section surjectivity
      (Lemma~\ref{mk-lem:affine_surj_of_vanishing_affine}), whose refinement step
      (Lemma~\ref{mk-lem:standard_cover_cofinal_affine}) uses the quasi-compactness of the affine open
      \(V\) in place of that of a distinguished \(D(f)\).
  \end{itemize}
  With \(\mathcal{F} = \mathcal{H}\) quasi-coherent, condition (3) of
  Lemma~\ref{mk-lem:cech_to_cohomology_on_basis} (standard-cover {\v C}ech vanishing) is \emph{not} the
  distinguished-open seed of Lemma~\ref{mk-lem:affine_serre_vanishing}: a standard cover here covers a
  general affine open \(V = \bigcup_i D(g_i)\), whose coordinate ring \(\Gamma(V)\) is not a single
  localisation \(R_f\), so the \(D(f)\) seed does not apply. Condition (3) is instead supplied by the
  enlarged-system seed \(\operatorname{HasVanishingHigherCech}(\operatorname{affineCoverSystemGeneral}
  R)\, \mathcal{H}\) of Lemma~\ref{mk-lem:affine_cech_vanishing_general_of_tildeVanishing}, which transports
  along \(\mathcal{H} \cong \widetilde{M}\) to the change-of-ring tilde seed
  (Lemma~\ref{mk-lem:affine_cech_vanishing_general_seed}). Reading off the conclusion at the member
  \(V \in \mathcal{B}\) gives \(H^q(V, \mathcal{H}) = 0\) for \(q \geq 1\); the identification
  \(H^q(V, -) = \operatorname{Ext}^q(j_!\mathcal{O}_V, -)\) is the Ext realization of
  Definition~\ref{mk-def:absolute_cohomology}.
\end{proof}

\begin{lemma}[A scheme isomorphism carries the corepresenting object along]
  \label{mk-lem:jshriek_transport_along_iso}
  \lean{AlgebraicGeometry.jShriekOU_transport_along_iso,
    AlgebraicGeometry.sectionsCorep,
    AlgebraicGeometry.sectionsCorepPushforward}
  \dcref{custom}

  \uses{mk-def:jshriek_ou, mk-lem:sectionsFunctorCorepIso}
  Let \(\varphi : X \cong Y\) be an isomorphism of schemes, \(V \subseteq X\) an open, and
  \(\Phi = \operatorname{pushforwardEquivOfIso}\varphi : X.\mathrm{Modules} \simeq Y.\mathrm{Modules}\)
  the induced equivalence. Then the pushforward carries the corepresenting object of
  Definition~\ref{mk-def:jshriek_ou} along the isomorphism:
  \[
    \Phi.\mathrm{functor}.\mathrm{obj}\bigl(j_!\mathcal{O}_V\bigr)
      \;\cong\;
    j_!\mathcal{O}_{\varphi.\mathrm{inv}^{-1}\,V},
  \]
  where \(\varphi.\mathrm{inv}^{-1}\,V\) is the preimage of \(V\) under \(\varphi.\mathrm{inv}\), i.e.\
  the image open of \(V\) in \(Y\).
\end{lemma}
\begin{proof}
  \uses{mk-def:jshriek_ou, mk-lem:sectionsFunctorCorepIso}
  Both sides corepresent the functor \(\operatorname{sectionsFunctor}(\varphi.\mathrm{inv}^{-1}V)
  \,\circ\, \mathrm{forget}\): the object \(\Phi.\mathrm{functor}(j_!\mathcal{O}_V)\) via the
  adjunction underlying the equivalence \(\Phi\), and the object
  \(j_!\mathcal{O}_{\varphi.\mathrm{inv}^{-1}V}\) via the defining corepresentability of
  \(\Gamma(\varphi.\mathrm{inv}^{-1}V,{-})\) by \(j_!\mathcal{O}\)
  (Definition~\ref{mk-def:jshriek_ou}). Both corepresent the same functor, so uniqueness of
  corepresenting objects identifies them.

  In detail, the argument is a corepresentability transport; \(j_!\mathcal{O}_V\) is never unfolded.
  Write \(A = j_!\mathcal{O}_V\). We exhibit two corepresentations of one and the same
  \(\mathrm{Type}\)-valued functor and invoke the uniqueness-up-to-isomorphism of a corepresenting
  object to obtain the stated isomorphism, with no adjunction-mate transposition or naturality square
  computed by hand. The common functor is built by chaining three identifications.
  \begin{enumerate}
    \item The adjunction underlying the equivalence \(\Phi\) presents the functor corepresented by
      \(\Phi.\mathrm{functor}(A)\) as the composite of \(\Phi.\mathrm{inverse}\) with the functor
      corepresented by \(A\).
    \item The corepresentability isomorphism of Lemma~\ref{mk-lem:sectionsFunctorCorepIso}, whiskered by the
      forgetful functor to types, identifies the functor corepresented by \(A\) with the sections
      functor: it is \(\operatorname{sectionsFunctor} V\) as a \(\mathrm{Type}\)-valued functor. This
      rests on the corepresentability of \(\Gamma(V,-)\) by \(j_!\mathcal{O}_V\) of
      Definition~\ref{mk-def:jshriek_ou}.
    \item Sections of a pushforward are, by definition, sections over the preimage open:
      \(\Phi.\mathrm{inverse} \circ \operatorname{sectionsFunctor} V =
      \operatorname{sectionsFunctor}(\varphi.\mathrm{inv}^{-1}V)\), an equality of functors. Reading
      Lemma~\ref{mk-lem:sectionsFunctorCorepIso} at the open \(\varphi.\mathrm{inv}^{-1}V\) then exhibits
      this same functor as the one corepresented by \(j_!\mathcal{O}_{\varphi.\mathrm{inv}^{-1}V}\).
  \end{enumerate}
  Thus \(\Phi.\mathrm{functor}(A)\) and \(j_!\mathcal{O}_{\varphi.\mathrm{inv}^{-1}V}\) corepresent the
  same functor; the uniqueness of corepresenting objects up to isomorphism reflects this to the stated
  isomorphism of objects. The only geometric inputs are the two definitional equalities
  ``corepresented functor \(=\) sections functor'' and ``sections of a pushforward \(=\) sections over
  the preimage''.
\end{proof}

\begin{lemma}[Quasi-coherence from a cover with quasi-coherent restrictions]
  \label{mk-lem:isQuasicoherent_of_coversTop_mathlib}
  \lean{SheafOfModules.IsQuasicoherent.of_coversTop}
  \mathlibok
  \dcref{custom}

  Let \(\mathcal{F}\) be a sheaf of modules and \((U_i)_{i \in I}\) a family of objects of the site
  covering the terminal object. If each over-picture restriction \(\mathcal{F}.\operatorname{over} U_i\)
  is quasi-coherent, then \(\mathcal{F}\) is quasi-coherent.
\end{lemma}

\begin{lemma}[Extraction of the local quasi-coherence datum]
  \label{mk-lem:nonempty_quasicoherentData_mathlib}
  \lean{SheafOfModules.IsQuasicoherent.nonempty_quasicoherentData}
  \mathlibok
  \dcref{custom}

  A quasi-coherent sheaf of modules \(\mathcal{F}\) admits a quasi-coherence datum: there is a family
  \((U_i)_{i \in I}\) of objects of the site covering the terminal object together with, for each \(i\),
  a presentation of \(\mathcal{F}.\operatorname{over} U_i\).
\end{lemma}

\begin{lemma}[Pushforward along an iso commutes with restriction to an open]
  \label{mk-lem:pushforward_commutes_restriction}
  \dcref{custom}
  \uses{mk-lem:pushforwardPushforwardEquivalence_mathlib, mk-lem:restrict_obj_mathlib}
  Let \(\varphi : X \cong Y\) be an isomorphism of schemes,
  \(\Phi = \operatorname{pushforwardEquivOfIso}\varphi\) the induced equivalence of module categories,
  and \(V \subseteq X\) an open with image \(V' = \varphi.\mathrm{inv}^{-1}\,V \subseteq Y\). The
  homeomorphism underlying \(\varphi\) restricts to a homeomorphism \(V \cong V'\), hence to an
  equivalence of the opens-sites \(\operatorname{Opens} V \simeq \operatorname{Opens} V'\); let
  \[
    e : (X.\mathrm{Modules} \!\restriction\! V) \;\simeq\; (Y.\mathrm{Modules} \!\restriction\! V')
  \]
  be the induced equivalence of restricted-module categories obtained from the site-equivalence
  transport of Lemma~\ref{mk-lem:pushforwardPushforwardEquivalence_mathlib}. Then for every
  \(\mathcal{O}_X\)-module \(H\) there is a comparison isomorphism
  \[
    e.\mathrm{functor}.\mathrm{obj}\bigl(H.\operatorname{over} V\bigr)
      \;\cong\;
    \bigl(\Phi.\mathrm{functor}.\mathrm{obj}\, H\bigr).\operatorname{over} V',
  \]
  identifying ``restrict to \(V\) then push forward over \(V'\)'' with ``push forward then restrict to
  \(V'\)''.
\end{lemma}
\begin{proof}
  \uses{mk-lem:pushforwardPushforwardEquivalence_mathlib, mk-lem:restrict_obj_mathlib}
  Both sides are computed open-by-open on the subspace site of \(V'\). For an open
  \(W \subseteq V'\) with preimage \(W_0 = \varphi.\mathrm{hom}^{-1}\,W \subseteq V\), the sections of the
  right-hand restriction are the sections of \(\Phi.\mathrm{functor}.\mathrm{obj}\, H\) over the image
  open \(W\), which — pushforward acting on opens by the inverse-image homeomorphism — are the sections
  of \(H\) over \(W_0\) (the restriction-of-objects identity of Lemma~\ref{mk-lem:restrict_obj_mathlib},
  applied to the open immersions \(V \hookrightarrow X\) and \(V' \hookrightarrow Y\)). The same
  sections compute the left-hand side: the site-equivalence transport
  \(e\) of Lemma~\ref{mk-lem:pushforwardPushforwardEquivalence_mathlib} relabels the open \(W\) of \(V'\)
  to the corresponding open \(W_0\) of \(V\) and reads off \(H.\operatorname{over} V\) there. The two
  presheaves of sections therefore agree, compatibly with restriction maps, and the comparison is a
  morphism of sheaves of modules over \(V'\); being an isomorphism on every section group it is an
  isomorphism. This is the open-restriction analogue of the \(\operatorname{bind}\) construction's
  per-piece transport, with the homeomorphism-induced opens-site equivalence in place of the
  iterated-slice equivalence.
\end{proof}

\begin{lemma}[Quasi-coherence of an iso-pushforward reduces to its slices]
  \label{mk-lem:pushforward_iso_qcoh_of_slice_qcoh}
  \lean{AlgebraicGeometry.pushforward_iso_qcoh_of_slice_qcoh,
    AlgebraicGeometry.coversTop_preimage_of_iso}
  \dcref{custom}

  \uses{mk-lem:isQuasicoherent_of_coversTop_mathlib, mk-lem:nonempty_quasicoherentData_mathlib}
  Let \(\varphi : X \cong Y\) be an isomorphism of schemes,
  \(\Phi = \operatorname{pushforwardEquivOfIso}\varphi\) the induced equivalence, \(H :
  X.\mathrm{Modules}\), and \((U_i)_{i \in I}\) a family of opens covering \(X\). Then quasi-coherence
  of the pushforward \(\Phi.\mathrm{functor}.\mathrm{obj}\, H\) reduces to quasi-coherence of its
  restrictions to the preimage-transported cover: if for every \(i\) the slice
  \[
    \bigl(\Phi.\mathrm{functor}.\mathrm{obj}\, H\bigr).\operatorname{over}
      \bigl(\varphi.\mathrm{inv}^{-1}\,U_i\bigr)
  \]
  is quasi-coherent, then \(\Phi.\mathrm{functor}.\mathrm{obj}\, H\) is quasi-coherent.
\end{lemma}
\begin{proof}
  \uses{mk-lem:isQuasicoherent_of_coversTop_mathlib, mk-lem:nonempty_quasicoherentData_mathlib}
  The preimages \(V_i = \varphi.\mathrm{inv}^{-1}\,U_i\) again cover \(Y\): the underlying map of
  \(\varphi.\mathrm{inv}\) is a homeomorphism, so a covering family of \(X\) pulls back to a covering
  family of \(Y\). Concretely, membership in \(\operatorname{Opens.grothendieckTopology}\) is preserved
  under the iso by transporting the covering sieve \(\operatorname{Sieve.ofObjects}\) along
  \(\varphi\) (the cover-transport step). With \((V_i)\) covering \(Y\) and each slice
  \((\Phi.\mathrm{functor}.\mathrm{obj}\, H).\operatorname{over} V_i\) quasi-coherent by hypothesis,
  quasi-coherence of \(\Phi.\mathrm{functor}.\mathrm{obj}\, H\) follows from the cover criterion
  (Lemma~\ref{mk-lem:isQuasicoherent_of_coversTop_mathlib}); the family \((U_i)\) is supplied by the
  quasi-coherence datum of \(H\) (Lemma~\ref{mk-lem:nonempty_quasicoherentData_mathlib}).
\end{proof}

\begin{lemma}[The continuous-functor pullback of sheaves of modules]
  \label{mk-lem:sheafOfModules_pullback_mathlib}
  \lean{SheafOfModules.pullback, SheafOfModules.instIsLeftAdjointPullback}
  \mathlibok
  \dcref{custom}

  Let \(F : C \to D\) be a continuous functor of sites (\(F.\operatorname{IsContinuous} J K\)), let
  \(S : \operatorname{Sheaf} J\,\operatorname{RingCat}\), \(R : \operatorname{Sheaf}
  K\,\operatorname{RingCat}\), and let \(\varphi : S \to (F.\operatorname{sheafPushforwardContinuous}
  \operatorname{RingCat} J\,K).\mathrm{obj}\, R\) be a morphism of sheaves of rings such that the induced
  pushforward \(\operatorname{SheafOfModules.pushforward}\varphi\) admits a right adjoint. Then there is a
  pullback functor \(\operatorname{pullback}\varphi : \operatorname{SheafOfModules} S \to
  \operatorname{SheafOfModules} R\), and it is a left adjoint (of the pushforward); in particular it
  preserves all colimits.
\end{lemma}

\begin{lemma}[The functor on sheaves of rings induced by a continuous functor]
  \label{mk-lem:sheafPushforwardContinuous_mathlib}
  \lean{CategoryTheory.Functor.sheafPushforwardContinuous}
  \mathlibok
  \dcref{custom}

  A continuous functor of sites \(F : C \to D\) (\(F.\operatorname{IsContinuous} J K\)) induces a
  pushforward functor \(F.\operatorname{sheafPushforwardContinuous} A\,J\,K : \operatorname{Sheaf}
  K\,A \to \operatorname{Sheaf} J\,A\) on \(A\)-valued sheaves. Taking \(A = \operatorname{RingCat}\),
  this is the receiving object \((F.\operatorname{sheafPushforwardContinuous} \operatorname{RingCat}
  J\,K).\mathrm{obj}\,R\) of the structure-sheaf comparison \(\psi_r\) below.
\end{lemma}

\begin{lemma}[Structure-sheaf comparison of a scheme morphism]
  \label{mk-lem:scheme_hom_toRingCatSheafHom_mathlib}
  \lean{AlgebraicGeometry.Scheme.Hom.toRingCatSheafHom}
  \mathlibok
  \dcref{custom}

  A morphism of schemes \(f : X \to Y\) carries an associated morphism of sheaves of (commutative)
  rings comparing the structure sheaves along \(f\); this is the sheaf-of-rings datum underlying the
  comorphism \(f^\sharp\), packaged as a single morphism in the sheaves-of-rings category.
\end{lemma}

\begin{lemma}[Beck--Chevalley identity for the slice post-composition functor]
  \label{mk-lem:over_post_forget_mathlib}
  \lean{CategoryTheory.Over.post_forget_eq_forget_comp}
  \mathlibok
  \dcref{custom}

  For a functor \(F : T \to D\) and an object \(X : T\), post-composition on slices commutes with the
  forgetful functors out of the slices: \(\operatorname{Over.post} F \circ \operatorname{Over.forget} =
  \operatorname{Over.forget} \circ F\) as an equality of functors \((\operatorname{Over} X) \to D\).
\end{lemma}

\begin{lemma}[The slice structure-sheaf ring map $\psi_r$]
  \label{mk-lem:slice_structureSheaf_hom}
  \lean{AlgebraicGeometry.sliceStructureSheafHom,
    AlgebraicGeometry.opensMapInvBase_isEquivalence,
    AlgebraicGeometry.overPost_slice_isContinuous,
    AlgebraicGeometry.sliceStructureSheafHom_pre_isRightAdjoint,
    AlgebraicGeometry.sliceStructureSheafHom_isRightAdjoint}
  \dcref{custom}

  \uses{mk-lem:sheafPushforwardContinuous_mathlib, mk-lem:scheme_hom_toRingCatSheafHom_mathlib,
    mk-lem:over_post_forget_mathlib, mk-lem:pushforwardPushforwardEquivalence_mathlib,
    mk-lem:overEquivalence_isContinuous}
  Let \(\varphi : X \cong Y\) be an isomorphism of schemes, \(U_i \subseteq X\) an (affine) open and
  \(V_i = \varphi.\mathrm{inv}^{-1}\,U_i \subseteq Y\) its image. Let \(F : \operatorname{Opens} U_i \to
  \operatorname{Opens} V_i\) be the continuous functor of opens-sites induced by the homeomorphism
  underlying \(\varphi\) restricted to \(U_i\); since \(\varphi\) is an isomorphism, \(F\) is an
  equivalence, so it is continuous (\(F.\operatorname{IsContinuous} J\,K\)) and final
  (\(F.\operatorname{Final}\)). Let \(\mathcal{O}_{U_i} : \operatorname{Sheaf} J\,\operatorname{RingCat}\)
  and \(\mathcal{O}_{V_i} : \operatorname{Sheaf} K\,\operatorname{RingCat}\) be the structure sheaves of
  the two sliced opens-sites, regarded as sheaves of rings. Then there is a morphism of sheaves of rings
  \[
    \psi_r : \mathcal{O}_{U_i} \longrightarrow
      \bigl(F.\operatorname{sheafPushforwardContinuous} \operatorname{RingCat} J\,K\bigr).\mathrm{obj}\,
      \mathcal{O}_{V_i},
  \]
  carrying the instance \((\operatorname{SheafOfModules.pushforward}\psi_r).\operatorname{IsRightAdjoint}\)
  — the right-adjointness obligation that makes the pullback functor along \(\psi_r\)
  (Lemma~\ref{mk-lem:sheafOfModules_pullback_mathlib}) available. This single cross-ring slice
  structure-sheaf ring map is the datum from which the pullback functor of
  Lemma~\ref{mk-lem:pushforward_slice_pullback_iso} is built.
\end{lemma}
\begin{proof}
  \uses{mk-lem:scheme_hom_toRingCatSheafHom_mathlib, mk-lem:over_post_forget_mathlib}
  The functor \(F\) and the two structure sheaves \(\mathcal{O}_{U_i}\), \(\mathcal{O}_{V_i}\) arise by
  restricting the scheme isomorphism \(\varphi\) to \(U_i\): the underlying homeomorphism restricts to a
  homeomorphism \(U_i \cong V_i\), inducing the opens-site equivalence \(F\). The map \(\psi_r\) is the
  structure-sheaf comparison of \(\varphi.\mathrm{hom}\) — its associated morphism of sheaves of rings
  (Lemma~\ref{mk-lem:scheme_hom_toRingCatSheafHom_mathlib}) — read along the opens-equivalence \(F\). The
  transport to the sliced sites uses the Beck--Chevalley identity that post-composition with a functor
  commutes with the forgetful functor out of the slice,
  \(\operatorname{Over.post} F \circ \operatorname{Over.forget} = \operatorname{Over.forget} \circ F\)
  (Lemma~\ref{mk-lem:over_post_forget_mathlib}, an equality of functors holding by definition),
  so the comparison descends to the over-categories with no further coherence datum. The
  right-adjointness instance for \(\operatorname{SheafOfModules.pushforward}\psi_r\) holds because \(F\)
  is an equivalence, whence its induced pushforward of sheaves of rings has a right adjoint. The result
  is the slice structure-sheaf ring map \(\psi_r\).
\end{proof}

\begin{lemma}[Uniqueness of left adjoints]
  \label{mk-lem:leftAdjointUniq_mathlib}
  \lean{CategoryTheory.Adjunction.leftAdjointUniq}
  \mathlibok
  \dcref{custom}

  If \(F, F' : C \to D\) are both left adjoint to a single functor \(G : D \to C\) (witnessed by
  adjunctions \(\mathrm{adj}_1 : F \dashv G\) and \(\mathrm{adj}_2 : F' \dashv G\)), then \(F\) and
  \(F'\) are canonically naturally isomorphic,
  \(\operatorname{leftAdjointUniq}\mathrm{adj}_1\,\mathrm{adj}_2 : F \cong F'\).
\end{lemma}

\begin{lemma}[Two pushforwards from an adjunction of sites are adjoint]
  \label{mk-lem:pushforwardPushforwardAdj_mathlib}
  \lean{SheafOfModules.pushforwardPushforwardAdj}
  \mathlibok
  \dcref{custom}

  Let \(F : C \to D\) and \(G : D \to C\) be continuous functors of sites carrying an adjunction
  \(\mathrm{adj} : F \dashv G\); let \(S : \operatorname{Sheaf} J\,\operatorname{RingCat}\), \(R :
  \operatorname{Sheaf} K\,\operatorname{RingCat}\), and let \(\varphi : S \to
  (F.\operatorname{sheafPushforwardContinuous} \operatorname{RingCat} J\,K).\mathrm{obj}\, R\) and
  \(\psi : R \to (G.\operatorname{sheafPushforwardContinuous} \operatorname{RingCat}
  K\,J).\mathrm{obj}\, S\) be morphisms of sheaves of rings satisfying the two compatibility squares
  \(H_1\) (relating \(\psi\) and \(\varphi\) to \(\mathrm{adj}\)'s counit) and \(H_2\) (the unit
  triangle). Then the induced module pushforwards are adjoint:
  \(\operatorname{SheafOfModules.pushforward}\varphi \dashv
  \operatorname{SheafOfModules.pushforward}\psi\).
\end{lemma}

\begin{lemma}[Post-composition equivalence on slices, with its correction term]
  \label{mk-lem:over_postEquiv_mathlib}
  \lean{CategoryTheory.Over.postEquiv}
  \mathlibok
  \dcref{custom}

  An equivalence of categories \(F : T \simeq D\) induces an equivalence of slice categories
  \(\operatorname{Over.postEquiv} X\,F : \operatorname{Over} X \simeq \operatorname{Over}
  (F.\mathrm{functor}.\mathrm{obj}\, X)\). Its forward functor is \(\operatorname{Over.post}
  F.\mathrm{functor}\); crucially its \emph{inverse} is not the bare \(\operatorname{Over.post}
  F.\mathrm{inverse}\) but the corrected composite \(\operatorname{Over.post} F.\mathrm{inverse}
  \circ \operatorname{Over.map}\bigl(F.\mathrm{unitIso}.\mathrm{inv}.\mathrm{app}\, X\bigr)\). The
  \(\operatorname{Over.map}\) factor repairs the fact that \(F.\mathrm{inverse}.\mathrm{obj}
  (F.\mathrm{functor}.\mathrm{obj}\, X)\) is only canonically isomorphic to \(X\), not equal --- the
  source of all the correction bookkeeping in the route below.
\end{lemma}

\begin{lemma}[Opens functor of a space isomorphism]
  \label{mk-lem:opens_mapMapIso_mathlib}
  \lean{TopologicalSpace.Opens.mapMapIso}
  \mathlibok
  \dcref{custom}

  An isomorphism \(h : T \cong T'\) of topological spaces induces an equivalence of the lattices of
  opens \(\operatorname{Opens.mapMapIso} h : \operatorname{Opens} T \simeq \operatorname{Opens} T'\),
  assembled from the \(\operatorname{Opens.map}\) functors of the two components of \(h\).
\end{lemma}

\begin{lemma}[Over-map functors are continuous]
  \label{mk-lem:instIsContinuousOverMapOver_mathlib}
  \lean{CategoryTheory.GrothendieckTopology.instIsContinuousOverMapOver}
  \mathlibok
  \dcref{custom}

  For a Grothendieck topology \(J\) on \(\mathcal{C}\) and a morphism \(f : X \to X'\) in
  \(\mathcal{C}\), the slice functor \(\operatorname{Over.map} f : \operatorname{Over} X \to
  \operatorname{Over} X'\) is continuous for the induced over-topologies \(J.\mathrm{over}\,X\) and
  \(J.\mathrm{over}\,X'\).
\end{lemma}

\begin{lemma}[Continuity is closed under composition]
  \label{mk-lem:functor_isContinuous_comp_mathlib}
  \lean{CategoryTheory.Functor.isContinuous_comp}
  \mathlibok
  \dcref{custom}

  If \(F : (\mathcal{C}, J) \to (\mathcal{D}, K)\) and \(G : (\mathcal{D}, K) \to (\mathcal{E}, L)\)
  are continuous functors of sites, then their composite \(F \circ G\) is continuous for \(J\) and
  \(L\).
\end{lemma}

\begin{lemma}[Cover-preservation lifts to slices]
  \label{mk-lem:coverPreserving_overPost_mathlib}
  \lean{CategoryTheory.CoverPreserving.overPost}
  \mathlibok
  \dcref{custom}

  If \(F : \mathcal{C} \to \mathcal{D}\) is cover-preserving for topologies \(J, K\), then for any
  \(X : \mathcal{C}\) the post-composition functor \(\operatorname{Over.post} F : \operatorname{Over} X
  \to \operatorname{Over}(F.\mathrm{obj}\, X)\) is cover-preserving for the induced over-topologies.
\end{lemma}

\begin{lemma}[Pullback is left adjoint to pushforward for sheaves of modules]
  \label{mk-lem:pullbackPushforwardAdjunction_mathlib}
  \lean{SheafOfModules.pullbackPushforwardAdjunction}
  \mathlibok
  \dcref{custom}

  For a morphism of sheaves of rings \(\varphi\) along a continuous functor of sites whose induced
  pushforward admits a right adjoint, the pullback functor of
  Lemma~\ref{mk-lem:sheafOfModules_pullback_mathlib} is left adjoint to the pushforward:
  \(\operatorname{pullback}\varphi \dashv \operatorname{pushforward}\varphi\).
\end{lemma}

\begin{lemma}[Pullback is left adjoint to pushforward for modules on schemes]
  \label{mk-lem:scheme_pullbackPushforwardAdjunction_mathlib}
  \lean{AlgebraicGeometry.Scheme.Modules.pullbackPushforwardAdjunction}
  \mathlibok
  \dcref{custom}

  For a morphism of schemes \(f : X \to Y\), the pullback functor of sheaves of modules is left
  adjoint to the pushforward,
  \[
    \operatorname{pullback} f \;\dashv\; \operatorname{pushforward} f
    \qquad (\operatorname{pullback} f : Y.\mathrm{Modules} \to X.\mathrm{Modules},\
      \operatorname{pushforward} f : X.\mathrm{Modules} \to Y.\mathrm{Modules}).
  \]
  In particular \(\operatorname{pushforward} f\) is a right adjoint, hence preserves finite limits
  (is left exact).
\end{lemma}

\begin{lemma}[Pushforward of modules composes for a composite of schemes]
  \label{mk-lem:scheme_pushforwardComp_mathlib}
  \lean{AlgebraicGeometry.Scheme.Modules.pushforwardComp}
  \mathlibok
  \dcref{custom}

  For morphisms of schemes \(f : X \to Y\) and \(g : Y \to Z\), the composition of the two
  pushforward functors on sheaves of modules identifies with the pushforward along the composite:
  \[
    \operatorname{pushforward} f \circ \operatorname{pushforward} g
      \;\cong\; \operatorname{pushforward}(f \circ g).
  \]
\end{lemma}

\begin{lemma}[Continuity of the slice equivalence of a scheme isomorphism]
  \label{mk-lem:slice_overs_equiv_continuity}
  \lean{AlgebraicGeometry.opensMapHomBase_isEquivalence,
    AlgebraicGeometry.opensEquivOfIso,
    AlgebraicGeometry.sliceOversEquiv,
    AlgebraicGeometry.sliceOversEquiv_functor_isContinuous,
    AlgebraicGeometry.overPost_slice_inverse_isContinuous,
    AlgebraicGeometry.sliceOversEquiv_inverse_isContinuous}
  \dcref{custom}

  \uses{mk-lem:opens_mapMapIso_mathlib, mk-lem:over_postEquiv_mathlib,
    mk-lem:instIsContinuousOverMapOver_mathlib, mk-lem:functor_isContinuous_comp_mathlib,
    mk-lem:coverPreserving_overPost_mathlib}
  Let \(\varphi : X \cong Y\) be an isomorphism of schemes, \(U_i \subseteq X\) an open with image
  \(V_i = \varphi.\mathrm{inv}^{-1}\,U_i\), and write \(J_X, J_Y\) for the opens-Grothendieck-topologies
  of \(X, Y\). Then:
  \begin{itemize}
    \item the opens-lattice functor \(\operatorname{Opens.map}\varphi.\mathrm{hom}.\mathrm{base}\) is an
      equivalence;
    \item \(\varphi\) induces an equivalence of opens-sites \(\operatorname{opensEquivOfIso}\varphi :
      \operatorname{Opens} X \simeq \operatorname{Opens} Y\), whose forward functor is definitionally
      \(\operatorname{Opens.map}\varphi.\mathrm{inv}.\mathrm{base}\) and whose inverse is
      \(\operatorname{Opens.map}\varphi.\mathrm{hom}.\mathrm{base}\);
    \item the slice equivalence \(\operatorname{sliceOversEquiv}\varphi\,U_i = \operatorname{Over.postEquiv}
      U_i\,(\operatorname{opensEquivOfIso}\varphi) : \operatorname{Over} U_i \simeq \operatorname{Over} V_i\)
      has both its functor and its inverse continuous for the over-Grothendieck-topologies.
  \end{itemize}
\end{lemma}
\begin{proof}
  \uses{mk-lem:opens_mapMapIso_mathlib, mk-lem:over_postEquiv_mathlib,
    mk-lem:instIsContinuousOverMapOver_mathlib, mk-lem:functor_isContinuous_comp_mathlib,
    mk-lem:coverPreserving_overPost_mathlib}
  The opens-equivalence is \(\operatorname{opensEquivOfIso}\varphi = \operatorname{Opens.mapMapIso}
  (\operatorname{forgetToTop}.\mathrm{mapIso}\,\varphi).\mathrm{symm}\)
  (Lemma~\ref{mk-lem:opens_mapMapIso_mathlib}); its functor is definitionally
  \(\operatorname{Opens.map}\varphi.\mathrm{inv}.\mathrm{base}\) and its inverse
  \(\operatorname{Opens.map}\varphi.\mathrm{hom}.\mathrm{base}\), so the latter is an equivalence.
  Forward continuity of the slice functor is the cover-preservation of post-composition along the
  opens-map (Lemma~\ref{mk-lem:coverPreserving_overPost_mathlib}, applied to
  \(\operatorname{Opens.map}\varphi.\mathrm{inv}.\mathrm{base}\), which is cover-preserving since the
  opens-map of a morphism is). For the inverse, read it through \(\operatorname{Over.postEquiv}\)'s
  inverse (Lemma~\ref{mk-lem:over_postEquiv_mathlib}): it is the composite of
  \(\operatorname{Over.post}(\operatorname{Opens.map}\varphi.\mathrm{hom}.\mathrm{base})\) --- continuous
  by cover-preservation of the opens-map (Lemma~\ref{mk-lem:coverPreserving_overPost_mathlib}) --- with the
  correction \(\operatorname{Over.map}(\mathrm{unitIso}.\mathrm{inv})\), which is always continuous
  (Lemma~\ref{mk-lem:instIsContinuousOverMapOver_mathlib}). Continuity of a composite of continuous
  functors (Lemma~\ref{mk-lem:functor_isContinuous_comp_mathlib}) gives the inverse's continuity.
\end{proof}

\begin{lemma}[The reverse slice ring map $\varphi''$]\label{mk-lem:slice_reverse_ring_map}
  \lean{AlgebraicGeometry.sliceReverseRingMap}
  \leanok
  \dcref{custom}

  \uses{mk-lem:slice_structureSheaf_hom, mk-lem:slice_overs_equiv_continuity}
  In the setting of Lemma~\ref{mk-lem:slice_overs_equiv_continuity}, write \(\mathrm{eqv} =
  \operatorname{sliceOversEquiv}\varphi\,U_i\). There is a morphism of sheaves of rings
  \[
    \varphi'' : (Y.\mathrm{ringCatSheaf}.\operatorname{over} V_i) \longrightarrow
      \bigl(\mathrm{eqv}.\mathrm{inverse}.\operatorname{sheafPushforwardContinuous} \operatorname{RingCat}\,
        (J_Y.\mathrm{over}\,V_i)\,(J_X.\mathrm{over}\,U_i)\bigr).\mathrm{obj}\,
      (X.\mathrm{ringCatSheaf}.\operatorname{over} U_i),
  \]
  along the \emph{correction-carrying} inverse \(\mathrm{eqv}.\mathrm{inverse} =
  \operatorname{Over.post}(\operatorname{Opens.map}\varphi.\mathrm{hom}.\mathrm{base}) \circ
  \operatorname{Over.map}(\mathrm{unitIso}.\mathrm{inv})\). The map \(\varphi''\) is
  \emph{object-level correction-free}: over any \(W : (\operatorname{Over} V_i)^{\mathrm{op}}\) the
  codomain section is the section of \(X.\mathrm{ringCatSheaf}.\operatorname{over} U_i\) over
  \(\mathrm{eqv}.\mathrm{inverse}.\mathrm{obj}\, W\), whose underlying open is \(\varphi.\mathrm{hom}^{-1}
  \,W.\mathrm{left}\) --- because \(\operatorname{Over.map}\) leaves the underlying open
  \((\,\cdot\,).\mathrm{left}\) unchanged, only post-composing the structure arrow. Thus at the level of
  sections \(\varphi''\) is exactly the slice restriction of the structure-sheaf comparison
  \(\varphi.\mathrm{hom}.\operatorname{toRingCatSheafHom}\); the \(\operatorname{Over.map}\) correction
  affects only the action of \(\mathrm{eqv}.\mathrm{inverse}\) on morphisms.
\end{lemma}
\begin{proof}
  \leanok
  \uses{mk-lem:slice_structureSheaf_hom, mk-lem:slice_overs_equiv_continuity}
  Define \(\varphi''\) to be the forward slice structure-sheaf comparison of
  Lemma~\ref{mk-lem:slice_structureSheaf_hom} for the inverse isomorphism \(\varphi^{-1} = \varphi.
  \mathrm{symm}\) over the open \(\varphi.\mathrm{inv}^{-1}\,U_i = V_i\), that is
  \[
    \varphi'' := \operatorname{sliceStructureSheafHom}(\varphi.\mathrm{symm})\,(\varphi.\mathrm{inv}^{-1}\,U_i).
  \]
  This morphism is constructed exactly as the forward slice ring map, but for \(\varphi.\mathrm{symm}\);
  no separate codomain-bridge construction is required, because the codomain it naturally carries ---
  along the \emph{bare} composite \(\operatorname{Over.post}(\operatorname{Opens.map}\varphi.\mathrm{symm}.
  \mathrm{inv}.\mathrm{base})\) --- is \emph{definitionally equal} to the codomain demanded in the
  statement, namely the one along the correction-carrying inverse slice functor
  \(\mathrm{eqv}.\mathrm{inverse} = \operatorname{Over.post}(\operatorname{Opens.map}\varphi.\mathrm{hom}.
  \mathrm{base}) \circ \operatorname{Over.map}(\mathrm{unitIso}.\mathrm{inv}_{U_i})\).

  The reconciliation of the two codomains is absorbed by definitional equality and needs no explicit
  isomorphism. Two coherences make this so. First, the composition law for pushforward-continuous
  functors, \(\operatorname{Functor.sheafPushforwardContinuousComp}\) for the two factors of
  \(\mathrm{eqv}.\mathrm{inverse}\), is the identity isomorphism:
  composing pushforward-continuous functors along a composite of continuous functors agrees on the nose
  with the pushforward along the composite. Second, the relabelling
  \(\operatorname{Over.map}(g) \circ \operatorname{forget} = \operatorname{forget}\)
  (\(\operatorname{Over.mapForget}\)) holds by reflexivity, since \(\operatorname{Over.map}\) leaves
  the underlying open \((\,\cdot\,).\mathrm{left}\) unchanged and only post-composes the structure
  arrow. As both coherences are definitional, the would-be codomain bridge reduces to the identity, and
  \(\varphi''\) is simply \(\operatorname{sliceStructureSheafHom}(\varphi.\mathrm{symm})\,V_i\) retyped
  along the (definitionally equal) corrected-inverse codomain. At the level of sections \(\varphi''\) is therefore
  exactly the slice restriction of the structure-sheaf comparison of \(\varphi^{-1}\); the
  \(\operatorname{Over.map}\) correction affects only the action on morphisms, which is what the counit
  and unit compatibility squares \(H_1, H_2\) absorb.
\end{proof}

\begin{lemma}[Counit compatibility square $H_1$ for the slice adjunction]\label{mk-lem:pushforward_slice_adjunction_h1}
  \lean{AlgebraicGeometry.pushforwardSliceAdjunctionH1}
  \leanok
  \dcref{custom}

  \uses{mk-lem:slice_reverse_ring_map, mk-lem:slice_structureSheaf_hom}
  In the setting of Lemma~\ref{mk-lem:slice_reverse_ring_map}, let \(\mathrm{adj} =
  (\operatorname{sliceOversEquiv}\varphi\,U_i).\mathrm{symm}.\operatorname{toAdjunction}\) be the
  adjunction underlying the slice equivalence. Then the reverse ring map \(\varphi''\) and the forward
  slice ring map \(\psi_r\) satisfy the counit-naturality compatibility square \(H_1\) required by
  Mathlib's two-pushforward adjunction (Lemma~\ref{mk-lem:pushforwardPushforwardAdj_mathlib}), absorbing
  the \(\operatorname{Over.map}(\mathrm{unitIso}.\mathrm{inv})\) correction of the inverse functor.
\end{lemma}
\begin{proof}
  \leanok
  \uses{mk-lem:slice_reverse_ring_map, mk-lem:slice_structureSheaf_hom}
  The square relates \(\varphi''\), \(\psi_r\) and the counit of \(\mathrm{adj}\). Since
  \(\varphi.\mathrm{hom}.\operatorname{toRingCatSheafHom}\) and
  \(\varphi.\mathrm{inv}.\operatorname{toRingCatSheafHom}\) are mutually inverse structure-sheaf
  comparisons, and the two slice ring maps are their respective over-pullbacks, naturality collapses the
  square to an equality of canonical equality-transport isomorphisms along the open identity
  \(\varphi.\mathrm{hom}^{-1}(\varphi.\mathrm{inv}^{-1} U_i) = U_i\), which holds by
  proof-irrelevance.
\end{proof}

\begin{lemma}[Unit compatibility square $H_2$ for the slice adjunction]\label{mk-lem:pushforward_slice_adjunction_h2}
  \lean{AlgebraicGeometry.pushforwardSliceAdjunctionH2}
  \leanok
  \dcref{custom}

  \uses{mk-lem:slice_reverse_ring_map, mk-lem:slice_structureSheaf_hom}
  In the setting of Lemma~\ref{mk-lem:slice_reverse_ring_map}, with \(\mathrm{adj}\) as in
  Lemma~\ref{mk-lem:pushforward_slice_adjunction_h1}, the ring maps \(\varphi''\) and \(\psi_r\) satisfy
  the unit-triangle compatibility square \(H_2\) required by
  Lemma~\ref{mk-lem:pushforwardPushforwardAdj_mathlib}.
\end{lemma}
\begin{proof}
  \leanok
  \uses{mk-lem:slice_reverse_ring_map, mk-lem:slice_structureSheaf_hom}
  Identical in shape to \(H_1\) (Lemma~\ref{mk-lem:pushforward_slice_adjunction_h1}): the unit triangle
  reduces, via the mutually inverse structure-sheaf comparisons of \(\varphi.\mathrm{hom}\) and
  \(\varphi.\mathrm{inv}\), to an equality of equality-transport morphisms along the same open identity,
  true by proof-irrelevance.
\end{proof}

\begin{lemma}[The two-pushforward adjunction for the slice ring maps]\label{mk-lem:pushforward_slice_two_adjunction}
  \lean{AlgebraicGeometry.pushforwardSliceTwoAdjunction}
  \leanok
  \dcref{custom}

  \uses{mk-lem:slice_structureSheaf_hom, mk-lem:slice_reverse_ring_map,
    mk-lem:pushforward_slice_adjunction_h1, mk-lem:pushforward_slice_adjunction_h2,
    mk-lem:slice_overs_equiv_continuity, mk-lem:pushforwardPushforwardAdj_mathlib,
    mk-lem:over_postEquiv_mathlib}
  Let \(\varphi : X \cong Y\), \(U_i \subseteq X\), \(V_i = \varphi.\mathrm{inv}^{-1}\,U_i\). Let
  \(\psi_r = \operatorname{sliceStructureSheafHom}\varphi\,U_i\) be the slice structure-sheaf ring map
  of Lemma~\ref{mk-lem:slice_structureSheaf_hom}, and let \(\varphi''\) be the reverse slice ring map of
  Lemma~\ref{mk-lem:slice_reverse_ring_map} --- the object-level correction-free comparison along the
  corrected inverse slice functor \((\operatorname{sliceOversEquiv}\varphi\,U_i).\mathrm{inverse}\).
  Then the two module pushforwards are adjoint:
  \[
    \operatorname{SheafOfModules.pushforward}\varphi'' \;\dashv\;
    \operatorname{SheafOfModules.pushforward}\psi_r .
  \]
\end{lemma}
\begin{proof}
  \leanok
  \uses{mk-lem:slice_structureSheaf_hom, mk-lem:slice_reverse_ring_map,
    mk-lem:pushforward_slice_adjunction_h1, mk-lem:pushforward_slice_adjunction_h2,
    mk-lem:slice_overs_equiv_continuity, mk-lem:pushforwardPushforwardAdj_mathlib,
    mk-lem:over_postEquiv_mathlib}
    Apply Mathlib's two-pushforward adjunction (Lemma~\ref{mk-lem:pushforwardPushforwardAdj_mathlib}) to the
  adjunction \(\mathrm{adj} = (\operatorname{sliceOversEquiv}\varphi\,U_i).\mathrm{symm}.
  \operatorname{toAdjunction}\) underlying the slice equivalence
  (Lemma~\ref{mk-lem:over_postEquiv_mathlib}). Its left adjoint is the corrected inverse
  \((\operatorname{sliceOversEquiv}\varphi\,U_i).\mathrm{inverse} =
  \operatorname{Over.post}(\operatorname{Opens.map}\varphi.\mathrm{hom}.\mathrm{base}) \circ
  \operatorname{Over.map}(\mathrm{unitIso}.\mathrm{inv})\), and its right adjoint is the forward functor
  \(\operatorname{Over.post}(\operatorname{Opens.map}\varphi.\mathrm{inv}.\mathrm{base})\) along which
  \(\psi_r\) lives. The required continuity of both functors is
  Lemma~\ref{mk-lem:slice_overs_equiv_continuity}. Feed Mathlib the two ring maps \(\varphi''\)
  (Lemma~\ref{mk-lem:slice_reverse_ring_map}) and \(\psi_r\) (Lemma~\ref{mk-lem:slice_structureSheaf_hom}),
  together with the two compatibility squares \(H_1\) (Lemma~\ref{mk-lem:pushforward_slice_adjunction_h1})
  and \(H_2\) (Lemma~\ref{mk-lem:pushforward_slice_adjunction_h2}); the output is the asserted adjunction
  \(\operatorname{pushforward}\varphi'' \dashv \operatorname{pushforward}\psi_r\).

  The substance isolated by this decomposition is the
  \(\operatorname{Over.map}(\mathrm{unitIso}.\mathrm{inv})\) correction carried by the inverse of
  \(\operatorname{Over.postEquiv}\) (Lemma~\ref{mk-lem:over_postEquiv_mathlib}). Because the open identity
  \(\varphi.\mathrm{hom}^{-1}(\varphi.\mathrm{inv}^{-1} U_i) = U_i\) does \emph{not} hold
  definitionally, the bare functor
  \(\operatorname{Over.post}(\operatorname{Opens.map}\varphi.\mathrm{hom}.\mathrm{base})\) lands in
  \(\operatorname{Over}(\varphi.\mathrm{hom}^{-1} V_i)\) rather than \(\operatorname{Over} U_i\), and the
  two slices are reconciled only up to the canonical isomorphism \(\mathrm{unitIso}.\mathrm{inv}\).
  Consequently the reverse ring map must be taken along the \emph{corrected} inverse slice functor
  (Lemma~\ref{mk-lem:slice_reverse_ring_map}) --- it is \emph{not}
  \(\operatorname{sliceStructureSheafHom}\varphi^{-1}\,V_i\), which lives along the bare post-composition
  and lands in the wrong slice --- and the same correction is absorbed by the two squares \(H_1, H_2\),
  each reducing to a proof-irrelevant equality-transport identity.
\end{proof}

\begin{lemma}[Pullback along $\psi_r$ realizes the per-slice pushforward]\label{mk-lem:pushforward_slice_pullback_iso}
  \lean{AlgebraicGeometry.pushforwardSlicePullbackIso}
  \leanok
  \dcref{custom}

  \uses{mk-lem:slice_structureSheaf_hom, mk-lem:slice_reverse_ring_map, mk-lem:sheafOfModules_pullback_mathlib,
    mk-lem:pushforward_slice_two_adjunction, mk-lem:leftAdjointUniq_mathlib,
    mk-lem:pullbackPushforwardAdjunction_mathlib, mk-lem:pushforward_obj_obj_mathlib}
  Let \(\varphi : X \cong Y\), \(U_i \subseteq X\), \(V_i = \varphi.\mathrm{inv}^{-1}\,U_i\), let
  \(\psi_r\) be the slice structure-sheaf ring map of Lemma~\ref{mk-lem:slice_structureSheaf_hom}, and write
  \(\Phi = \operatorname{pushforwardEquivOfIso}\varphi\). For every \(\mathcal{O}_X\)-module \(H\) the
  pullback functor \(\operatorname{pullback}\psi_r : \operatorname{SheafOfModules}\mathcal{O}_{U_i} \to
  \operatorname{SheafOfModules}\mathcal{O}_{V_i}\) applied directly to the slice
  \(H.\operatorname{over} U_i\) (which already lies in its domain) is isomorphic to the restricted
  pushforward:
  \[
    \bigl(\operatorname{pullback}\psi_r\bigr).\mathrm{obj}\bigl(H.\operatorname{over} U_i\bigr)
      \;\cong\;
    \bigl(\Phi.\mathrm{functor}.\mathrm{obj}\, H\bigr).\operatorname{over} V_i .
  \]
\end{lemma}
\begin{proof}
  \leanok
  \uses{mk-lem:slice_reverse_ring_map, mk-lem:pushforward_slice_two_adjunction, mk-lem:leftAdjointUniq_mathlib,
    mk-lem:pullbackPushforwardAdjunction_mathlib, mk-lem:pushforward_obj_obj_mathlib}
    Assemble the isomorphism in two steps. The unit/free-module comparison
  \(\operatorname{pullbackObjUnitToUnit}\) is \emph{not} used: it identifies the pullback only on the
  unit (rank-one free) module, whereas \(H.\operatorname{over} U_i\) is an arbitrary module. Instead we
  identify the whole functor \(\operatorname{pullback}\psi_r\) by uniqueness of left adjoints, which is
  valid for every \(H\) at once.

  \emph{Step 1 --- identify the pullback functor with a pushforward.} The pullback
  \(\operatorname{pullback}\psi_r\) is left adjoint to \(\operatorname{pushforward}\psi_r\)
  (Lemma~\ref{mk-lem:pullbackPushforwardAdjunction_mathlib}). By
  Lemma~\ref{mk-lem:pushforward_slice_two_adjunction} the pushforward
  \(\operatorname{pushforward}\varphi''\) along the reverse slice ring map \(\varphi''\) of
  Lemma~\ref{mk-lem:slice_reverse_ring_map} is \emph{also} left adjoint to the same
  \(\operatorname{pushforward}\psi_r\). Two left adjoints of one functor are canonically isomorphic
  (Lemma~\ref{mk-lem:leftAdjointUniq_mathlib}); applied to these two adjunctions it yields a natural
  isomorphism of functors
  \[
    \operatorname{leftAdjointUniq}\bigl(\operatorname{pullbackPushforwardAdjunction}\psi_r\bigr)\,
      \mathrm{adj}_A : \operatorname{pullback}\psi_r \;\cong\; \operatorname{pushforward}\varphi'',
  \]
  where \(\mathrm{adj}_A\) is the two-pushforward adjunction of
  Lemma~\ref{mk-lem:pushforward_slice_two_adjunction}.

  \emph{Step 2 --- the section identity.} Evaluate this functor isomorphism at \(H.\operatorname{over}
  U_i\):
  \[
    (\operatorname{pullback}\psi_r).\mathrm{obj}\bigl(H.\operatorname{over} U_i\bigr) \;\cong\;
    (\operatorname{pushforward}\varphi'').\mathrm{obj}\bigl(H.\operatorname{over} U_i\bigr).
  \]
  The right-hand object is \(\bigl(\Phi.\mathrm{functor}.\mathrm{obj}\, H\bigr).\operatorname{over}
  V_i\) by direct definitional unfolding: pushforward sections are computed by
  preimage (Lemma~\ref{mk-lem:pushforward_obj_obj_mathlib}), so over an open \(W\) of the slice both
  objects have the very same sections \(\Gamma(H, \varphi.\mathrm{hom}^{-1}\, W.\mathrm{left})\) ---
  the slice pushforward along \(\varphi''\) and the slice of the pushforward
  \(\Phi.\mathrm{functor}.\mathrm{obj}\, H\) over \(V_i\) unfold to one and the same sections functor.
  Composing the functor isomorphism of Step~1, evaluated at \(H.\operatorname{over} U_i\), with this
  definitional identification yields the stated isomorphism.
\end{proof}

\begin{lemma}[Pushforward along an iso preserves quasi-coherence]\label{mk-lem:pushforward_iso_preserves_qcoh}
  \lean{AlgebraicGeometry.pushforward_iso_preserves_qcoh}
  \leanok
  \dcref{custom}

  \uses{mk-lem:slice_structureSheaf_hom, mk-lem:pushforward_slice_pullback_iso,
    mk-lem:pushforward_iso_qcoh_of_slice_qcoh, mk-lem:presentation_map_mathlib,
    mk-lem:presentation_ofIsIso_mathlib, mk-lem:nonempty_quasicoherentData_mathlib,
    mk-lem:isAffineOpen_image_of_iso_mathlib}
  \textit{Source: Stacks Project, Schemes, ``Functoriality for quasi-coherent modules''.}
  Let \(\varphi : X \cong Y\) be an isomorphism of schemes and \(H : X.\mathrm{Modules}\) a
  quasi-coherent module. Then the pushforward \(\Phi.\mathrm{functor}.\mathrm{obj}\, H\) along
  \(\varphi\) is again quasi-coherent.
\end{lemma}
\begin{proof}
  \leanok
  \uses{mk-lem:slice_structureSheaf_hom, mk-lem:pushforward_slice_pullback_iso,
    mk-lem:pushforward_iso_qcoh_of_slice_qcoh, mk-lem:presentation_map_mathlib,
    mk-lem:presentation_ofIsIso_mathlib, mk-lem:nonempty_quasicoherentData_mathlib,
    mk-lem:isAffineOpen_image_of_iso_mathlib}
    Quasi-coherence is a \emph{local} datum: it transports across the equivalence \(\Phi\) by moving
  \(H\)'s local presentations to the image cover, member by member. We never seek a global presentation
  of \(H\) (it has none — \(H\) is quasi-coherent only on the abstract source); we move the local
  presentations one open at a time. The carrying functor need not be the full
  restricted-module equivalence: a single \emph{left adjoint} pullback functor along one cross-ring
  slice structure-sheaf ring map suffices, which keeps the colimit-preservation that
  presentation-transport requires while collapsing the heavy equivalence bookkeeping to one ring map,
  one unit isomorphism, and one comparison isomorphism.

  \emph{The cover.} Extract \(H\)'s quasi-coherence datum
  (Lemma~\ref{mk-lem:nonempty_quasicoherentData_mathlib}): a family \((U_i)_{i \in I}\) of opens covering
  \(X\) together with, for each \(i\), a presentation \(p_i\) of \(H.\operatorname{over} U_i\). Form the
  image cover \(V_i = \varphi.\mathrm{inv}^{-1}\,U_i\) of \(Y\) (the image of an affine open under the
  scheme isomorphism is again affine, Lemma~\ref{mk-lem:isAffineOpen_image_of_iso_mathlib}, when one wants
  the cover to remain affine). By Lemma~\ref{mk-lem:pushforward_iso_qcoh_of_slice_qcoh} it suffices to show
  that \((\Phi.\mathrm{functor}.\mathrm{obj}\, H).\operatorname{over} V_i\) is quasi-coherent for every
  \(i\).

  \emph{Per-member transport.} Fix \(i\) and write \(W = U_i\), \(V = V_i =
  \varphi.\mathrm{inv}^{-1}\,W\). Let \(\psi_r\) be the slice structure-sheaf ring map of
  Lemma~\ref{mk-lem:slice_structureSheaf_hom}, and consider the pullback functor
  \(\operatorname{pullback}\psi_r\). Being a left adjoint, it preserves colimits, so by
  Lemma~\ref{mk-lem:presentation_map_mathlib} the presentation \(p_i\) of \(H.\operatorname{over} W\)
  (transported across the opens-equivalence) transports to a presentation of
  \((\operatorname{pullback}\psi_r).\mathrm{obj}(H.\operatorname{over} W \text{ transported})\). The
  comparison isomorphism of Lemma~\ref{mk-lem:pushforward_slice_pullback_iso},
  \[
    \bigl(\operatorname{pullback}\psi_r\bigr).\mathrm{obj}\bigl(H.\operatorname{over} W
      \text{ transported}\bigr)
      \;\cong\;
    \bigl(\Phi.\mathrm{functor}.\mathrm{obj}\, H\bigr).\operatorname{over} V,
  \]
  then carries that presentation, via Lemma~\ref{mk-lem:presentation_ofIsIso_mathlib}, to a presentation of
  \((\Phi.\mathrm{functor}.\mathrm{obj}\, H).\operatorname{over} V\). A sheaf of modules with a (global)
  presentation is quasi-coherent, so each member of the image cover is quasi-coherent, and the claim
  follows by Lemma~\ref{mk-lem:pushforward_iso_qcoh_of_slice_qcoh}.

  This is the isomorphism case of the Stacks functoriality remark (``it is in general not the case that
  the pushforward of a quasi-coherent module is quasi-coherent''; here it holds because \(\varphi\) is an
  isomorphism, equivalently \(\varphi_* = (\varphi^{-1})^*\) is a pullback).
\end{proof}

\begin{lemma}[The image of an affine open under an open immersion is affine]
  \label{mk-lem:isAffineOpen_image_of_iso_mathlib}
  \lean{AlgebraicGeometry.Scheme.Hom.isAffineOpen_iff_of_isOpenImmersion}
  \mathlibok
  \dcref{custom}

  Let \(f : X \to Y\) be an open immersion of schemes (in particular an isomorphism) and \(U\) an open
  of \(X\). Then the image open \(f \,{''}\, U\) is affine if and only if \(U\) is affine. In
  particular the image \(V'\) of an affine open \(V\) under a scheme isomorphism is an affine open.
\end{lemma}

\begin{lemma}
[Ext transport along the whole-scheme spectrum equivalence (Need~\#1)]
  \label{mk-lem:modules_isoSpec_ext_transport}
  \lean{AlgebraicGeometry.modulesIsoSpecExtTransport,
    AlgebraicGeometry.Scheme.Modules.pushforwardEquivOfIso,
    AlgebraicGeometry.pushforwardEquivOfIso_functor_additive,
    AlgebraicGeometry.Scheme.Modules.pushforwardExtAddEquiv}
  \dcref{custom}

  \uses{mk-lem:isoSpec_scheme_mathlib, mk-lem:ext_mapExactFunctor_mathlib,
    mk-lem:modules_pushforward_mathlib, mk-lem:jshriek_transport_along_iso,
    mk-lem:pushforward_iso_preserves_qcoh, mk-lem:isAffineOpen_image_of_iso_mathlib}
  Let \(U\) be a scheme with \([\,\mathrm{IsAffine}\ U\,]\), let \(V \subseteq U\) be an affine open,
  and let \(\mathcal{H}\) be a quasi-coherent \(\mathcal{O}_U\)-module. The whole-scheme isomorphism
  \(U \cong \operatorname{Spec}\Gamma(U, \mathcal{O}_U)\) (Lemma~\ref{mk-lem:isoSpec_scheme_mathlib})
  induces, for every \(q\), an isomorphism
  \[
    \operatorname{Ext}^q_{U}\bigl(j_!\mathcal{O}_V,\, \mathcal{H}\bigr)
      \;\cong\;
    \operatorname{Ext}^q_{\operatorname{Spec}\Gamma U}\bigl(j_!\mathcal{O}_{V'},\, \mathcal{H}'\bigr),
  \]
  where \(V'\) is the affine open image of \(V\) under the isomorphism and \(\mathcal{H}'\) is the
  transported (still quasi-coherent) coefficient. In particular the residual over \(U\) vanishes iff
  its transport over \(\operatorname{Spec}\Gamma U\) does.
\end{lemma}
\begin{proof}
  \uses{mk-lem:isoSpec_scheme_mathlib, mk-lem:ext_mapExactFunctor_mathlib,
    mk-lem:modules_pushforward_mathlib, mk-lem:jshriek_transport_along_iso,
    mk-lem:pushforward_iso_preserves_qcoh, mk-lem:isAffineOpen_image_of_iso_mathlib}
  Since \(U\) is affine, \(\operatorname{Scheme.isoSpec}\) (Lemma~\ref{mk-lem:isoSpec_scheme_mathlib})
  gives a genuine isomorphism of schemes \(\varphi : U \cong \operatorname{Spec}\Gamma(U, \mathcal{O}_U)\).
  Assemble from the pushforward functors and their coherences
  (Lemma~\ref{mk-lem:modules_pushforward_mathlib}) the equivalence of module categories
  \[
    \Phi = \operatorname{pushforwardEquivOfIso}\varphi :
      U.\mathrm{Modules} \;\simeq\; (\operatorname{Spec}\Gamma U).\mathrm{Modules},
  \]
  with \(\Phi.\mathrm{functor}\) the pushforward along \(\varphi.\mathrm{hom}\) and its quasi-inverse the
  pushforward along \(\varphi.\mathrm{inv}\); the unit and counit are supplied by the composition and
  identity coherences. An equivalence is exact and additive, so
  \(\operatorname{Ext.mapExactFunctor}\) (Lemma~\ref{mk-lem:ext_mapExactFunctor_mathlib}) applies and
  yields a map on \(\operatorname{Ext}\); since \(\Phi.\mathrm{functor}\), being one half of an
  equivalence, preserves injective objects, its induced \(\operatorname{Ext}\)-map is bijective
  (using enough injectives on \(U.\mathrm{Modules}\)), giving an additive isomorphism. This gives
  the displayed isomorphism once the two arguments are identified. For the first
  argument, a scheme isomorphism carries the corepresenting object along
  (Lemma~\ref{mk-lem:jshriek_transport_along_iso}): \(\Phi.\mathrm{functor}(j_!\mathcal{O}_V) \cong
  j_!\mathcal{O}_{V'}\), where \(V' = \varphi.\mathrm{hom}\,{''}\,V\) is the image, an affine open of
  \(\operatorname{Spec}\Gamma U\) by Lemma~\ref{mk-lem:isAffineOpen_image_of_iso_mathlib}. For the second,
  \(\Phi.\mathrm{functor}(\mathcal{H}) = \mathcal{H}'\) is again quasi-coherent because pushforward along
  an isomorphism preserves quasi-coherence (Lemma~\ref{mk-lem:pushforward_iso_preserves_qcoh}).
\end{proof}

\begin{lemma}[Serre Ext-vanishing over an affine scheme via the spectrum transport]\label{mk-lem:ext_jShriekOU_eq_zero_of_specIso}
  \lean{AlgebraicGeometry.ext_jShriekOU_eq_zero_of_specIso}
  \leanok
  \dcref{custom}

  \uses{mk-lem:modules_isoSpec_ext_transport, mk-lem:jshriek_transport_along_iso,
    mk-lem:pushforward_iso_preserves_qcoh, mk-lem:affine_serre_vanishing_general_open,
    mk-lem:enoughInjectives_of_hasInjectiveResolutions, mk-lem:subsingleton_ext_of_iso_fst}
  Let \(U\) be a scheme with chosen injective resolutions on \(U.\mathrm{Modules}\), let
  \(\varphi : U \cong \operatorname{Spec} R\) be an isomorphism of schemes,
  \(\Phi = \operatorname{pushforwardEquivOfIso}\varphi\) the induced equivalence of module categories,
  \(V \subseteq U\) an open and \(\mathcal{H}\) an \(\mathcal{O}_U\)-module. Suppose
  \(\Phi.\mathrm{functor}(j_!\mathcal{O}_V) \cong j_!\mathcal{O}_{V'}\) for some \emph{affine} open
  \(V' \subseteq \operatorname{Spec} R\) (Lemma~\ref{mk-lem:jshriek_transport_along_iso}) and that
  \(\Phi.\mathrm{functor}(\mathcal{H})\) is quasi-coherent
  (Lemma~\ref{mk-lem:pushforward_iso_preserves_qcoh}). Then for every \(q \geq 1\) every class
  \(e \in \operatorname{Ext}^q_U(j_!\mathcal{O}_V, \mathcal{H})\) vanishes.
\end{lemma}
\begin{proof}
  \leanok
  \uses{mk-lem:modules_isoSpec_ext_transport, mk-lem:jshriek_transport_along_iso,
    mk-lem:pushforward_iso_preserves_qcoh, mk-lem:affine_serre_vanishing_general_open,
    mk-lem:enoughInjectives_of_hasInjectiveResolutions, mk-lem:subsingleton_ext_of_iso_fst}
    The chosen injective resolutions provide enough injectives on \(U.\mathrm{Modules}\)
  (Lemma~\ref{mk-lem:enoughInjectives_of_hasInjectiveResolutions}), which transport across the equivalence
  \(\Phi\) to enough injectives on \((\operatorname{Spec} R).\mathrm{Modules}\). Over
  \(\operatorname{Spec} R\) the general-affine-open Serre vanishing
  (Lemma~\ref{mk-lem:affine_serre_vanishing_general_open}), applied to the affine open \(V'\) and the
  quasi-coherent coefficient \(\Phi.\mathrm{functor}(\mathcal{H})\), makes
  \(\operatorname{Ext}^q(j_!\mathcal{O}_{V'},\, \Phi.\mathrm{functor}\,\mathcal{H})\) a subsingleton.
  Transferring along the transport isomorphism
  \(\Phi.\mathrm{functor}(j_!\mathcal{O}_V) \cong j_!\mathcal{O}_{V'}\) in the first Ext-argument
  (Lemma~\ref{mk-lem:subsingleton_ext_of_iso_fst}) shows
  \(\operatorname{Ext}^q(\Phi.\mathrm{functor}(j_!\mathcal{O}_V),\, \Phi.\mathrm{functor}\,\mathcal{H})\)
  is a subsingleton as well. The spectrum Ext-transport additive isomorphism
  \[
    \operatorname{Ext}^q_U(j_!\mathcal{O}_V, \mathcal{H})
      \;\cong\;
    \operatorname{Ext}^q_{\operatorname{Spec} R}\bigl(\Phi.\mathrm{functor}(j_!\mathcal{O}_V),\,
      \Phi.\mathrm{functor}\,\mathcal{H}\bigr)
  \]
  (Lemma~\ref{mk-lem:modules_isoSpec_ext_transport}) is injective, so it carries \(e\) into a subsingleton
  and forces \(e = 0\).
\end{proof}

\begin{lemma}[Sheafification of a sectionwise-locally-zero presheaf vanishes]\label{mk-lem:isZero_presheafToSheaf_of_sections_locally_zero}
  \lean{CategoryTheory.GrothendieckTopology.isZero_presheafToSheaf_of_sections_locally_zero}
  \leanok
  \dcref{custom}

  \uses{mk-lem:sheafify_kills_locally_zero}
  Let \(J\) be a Grothendieck topology on a site \(\mathcal{C}\) and \(Q\) a presheaf of
  abelian groups. Suppose that for every object \(U\) and every section \(s \in Q(U)\) there
  is a covering sieve \(S \in J(U)\) on which \(s\) restricts to \(0\) --- i.e.\ \(Q(g)(s) =
  0\) for every arrow \(g : V \to U\) in \(S\). Then the sheafification of \(Q\) is the zero
  sheaf.

  This is the \emph{sectionwise} (local-injectivity) strengthening of
  Lemma~\ref{mk-lem:isZero_presheafToSheaf_of_locally_isZero}. That lemma requires the
  \emph{objectwise} vanishing \(Q(V) = 0\) for every member \(V\) of a single covering sieve,
  which forces the vanishing to hold on a downward-closed family. The present hypothesis only
  asks each individual section to die on \emph{some} covering sieve, so the witnessing sieve
  may be generated by a cofinal \emph{basis} of good opens (such as the affine opens, where
  the relevant section group already vanishes) even though that basis is not downward closed
  --- exactly the situation of the relative-affine vanishing argument, where the affine opens
  are not downward closed in the opens topology.
\end{lemma}
\begin{proof}
  \leanok
  \uses{mk-lem:sheafify_kills_locally_zero}
  Compare \(Q\) with the constant zero presheaf \(Z\) via the zero map \(0 : Q \to Z\). The
  hypothesis, applied to a difference of sections \(x - y\), shows \(0\) is locally injective;
  local surjectivity is automatic because the target is objectwise a zero object. Hence \(0 :
  Q \to Z\) is locally bijective, i.e.\ lies in the local-isomorphism class \(J.W\), and
  Lemma~\ref{mk-lem:sheafify_kills_locally_zero} collapses the sheafification of \(Q\) to the
  zero sheaf.
\end{proof}

\begin{lemma}
[Pullback of a quasi-coherent module along an open immersion is quasi-coherent]
  \label{mk-lem:isQuasicoherent_pullback_opens}
  \lean{AlgebraicGeometry.isQuasicoherent_pullback_opens,
    AlgebraicGeometry.modulesOverOpensEquivalence,
    AlgebraicGeometry.opens_ι_image_overEquivalence_functor,
    AlgebraicGeometry.overEquivalence_functor_isContinuous_opens,
    AlgebraicGeometry.overEquivalence_inverse_isContinuous_opens,
    AlgebraicGeometry.overOpensForgetIso,
    AlgebraicGeometry.overOpensForgetInvIso,
    AlgebraicGeometry.overOpensRingHom,
    AlgebraicGeometry.overOpensRingInvHom,
    AlgebraicGeometry.overOpensIsoRestrict,
    AlgebraicGeometry.overOpensFunctorUnitIso,
    AlgebraicGeometry.overOpensInverseUnitIso,
    AlgebraicGeometry.presentationOverOpens,
    AlgebraicGeometry.presentationRestrictOfOver,
    AlgebraicGeometry.presentationRestrictSliceOfOver,
    AlgebraicGeometry.pullbackObjUnitToUnit_isIso_of_isIso,
    AlgebraicGeometry.restrictIsoUnitIso,
    AlgebraicGeometry.pushforward_overOpensRingHom_isRightAdjoint,
    AlgebraicGeometry.pushforward_overOpensRingInvHom_isRightAdjoint}
  \dcref{custom}

  \uses{mk-lem:isQuasicoherent_restrict_basicOpen, mk-lem:restrict_over_compat,
    mk-lem:restrictFunctorIsoPullback_mathlib}
  Let \(\iota : U \hookrightarrow X\) be the inclusion of an open subscheme and let
  \(\mathcal{F}\) be a quasi-coherent \(\mathcal{O}_X\)-module. Then the pullback
  \(\iota^*\mathcal{F}\) --- equivalently, by
  Lemma~\ref{mk-lem:restrictFunctorIsoPullback_mathlib}, the restriction of \(\mathcal{F}\)
  to \(U\) --- is a quasi-coherent \(\mathcal{O}_U\)-module.
\end{lemma}
\begin{proof}
  \uses{mk-lem:isQuasicoherent_restrict_basicOpen, mk-lem:restrict_over_compat,
    mk-lem:restrictFunctorIsoPullback_mathlib}
  This is the general-open-immersion analogue of Lemma~\ref{mk-lem:restrict_over_compat}
  and its quasi-coherence corollary Lemma~\ref{mk-lem:isQuasicoherent_restrict_basicOpen},
  where the affine open \(D(g)\) is replaced by an arbitrary open \(U \subseteq X\).
  Quasi-coherence is local on the base, so it suffices to verify it on the members of
  an open cover; we cover \(X\) by affine opens and check \(\iota^*\mathcal{F}\) on
  each, then conclude by the cover-to-top criterion. On an affine chart, the
  site-theoretic equivalence between \(\mathcal{O}_U\)-modules and sheaves over the
  localized slice category \((\operatorname{Opens} X)/U\) allows the quasi-coherence
  presentation of \(\mathcal{F}\) to be restricted: the induced over-presentation of
  \(\mathcal{F}\) over \(U\) transports, slice by slice, to a presentation of
  \(\mathcal{F}|_U\) over the chart, exhibiting it as quasi-coherent there. The
  canonical isomorphism between \(\mathcal{F}|_U\) and the pullback
  \(\iota^*\mathcal{F}\) supplied by
  Lemma~\ref{mk-lem:restrictFunctorIsoPullback_mathlib}, together with coherence of the
  relevant unit maps, makes the transport well-defined; transporting the
  quasi-coherence property across this isomorphism yields quasi-coherence of
  \(\iota^*\mathcal{F}\).
\end{proof}

\begin{lemma}
[Each {\v C}ech term is $f_*$-acyclic]
  \label{mk-lem:cech_term_pushforward_acyclic}
  \lean{AlgebraicGeometry.cechTerm_pushforward_acyclic,
    AlgebraicGeometry.rightDerived_additive,
    AlgebraicGeometry.injectiveResolutions_additive,
    AlgebraicGeometry.isRightAcyclic_of_iso,
    AlgebraicGeometry.isZero_biproduct,
    AlgebraicGeometry.higherDirectImage_affineHom_acyclic,
    AlgebraicGeometry.isAffineHom_of_isAffine_of_isSeparated,
    AlgebraicGeometry.isAffineOpen_coverInterOpen,
    AlgebraicGeometry.isAffineOpen_iInf_fin,
    AlgebraicGeometry.pushPullObj_opens_pushforward_acyclic}
  \dcref{01SG,custom}

  \uses{mk-lem:affine_serre_vanishing, mk-def:cech_nerve, mk-lem:higher_direct_image_presheaf, mk-lem:isQuasicoherent_pullback_opens, mk-lem:pushPull_sigma_iso, mk-lem:rightAcyclic_finite_prod}
  \textit{Source: Stacks Project, Cohomology of Schemes,
  \texttt{lemma-relative-affine-vanishing} (relative affine vanishing).}
  Let \(f : X \to S\) be a quasi-compact morphism of schemes with \emph{both} the
  source \(X\) and the target \(S\) separated, let \(\mathcal{F}\) be a
  quasi-coherent \(\mathcal{O}_X\)-module, and let \(\mathfrak{U}\) be a finite
  affine open cover of \(X\) (so, by separatedness of \(X\), every intersection
  \(U_{i_0 \ldots i_p}\) is again affine). Assume moreover that for every
  multi-index \(\sigma : \{0, \ldots, p\} \to I\) the category of
  \(\mathcal{O}\)-modules on the intersection subscheme
  \(U_\sigma = U_{\sigma(0) \ldots \sigma(p)}\) admits injective resolutions
  (hypothesis \(\mathrm{hres}\)). Then for every \(p \geq 0\) the degree-\(p\)
  {\v C}ech term \(\mathcal{C}^p =
  \prod_{i_0 \ldots i_p} (j_{i_0 \ldots i_p})_*(\mathcal{F}|_{U_{i_0 \ldots i_p}})\)
  is right \((\,\cdot\,)_*\)-acyclic for the pushforward \(f_*\):
  \[
    R^k f_*(\mathcal{C}^p) = 0
      \qquad \text{for all } k \geq 1.
  \]

  Two hypotheses beyond the bare ``\(f\) separated and quasi-compact'' are genuinely
  required. First, the \emph{target} \(S\) must be separated: the structure map
  \(f \circ j_\sigma : U_\sigma \to S\) from the affine \(U_\sigma\) is an affine
  morphism only because a morphism from an affine scheme to a separated scheme is
  affine (Stacks Project, Tag~01SG). Over a non-separated base this fails --- for the
  affine line with doubled origin \(S\) and \(f = \operatorname{id}\), the inclusion
  of a standard affine chart is not an affine morphism, and termwise acyclicity
  breaks. Second, the higher direct images are taken in the derived-functor sense,
  which presupposes enough injectives; since this is not available for
  \(\operatorname{QCoh}\) in general, the existence of injective resolutions on each
  intersection subscheme \(U_\sigma\) is carried as the explicit hypothesis
  \(\mathrm{hres}\) (the relative analogue of asking that
  \(\operatorname{QCoh}(X)\) have injective resolutions).
\end{lemma}
\begin{proof}
  \uses{mk-lem:affine_serre_vanishing, mk-def:cech_nerve, mk-lem:higher_direct_image_presheaf, mk-lem:isQuasicoherent_pullback_opens, mk-lem:pushPull_sigma_iso, mk-lem:rightAcyclic_finite_prod}
  The degree-\(p\) {\v C}ech term is identified, via the push--pull product decomposition
  of Lemma~\ref{mk-lem:pushPull_sigma_iso}, with the finite product
  \(\mathcal{C}^p \cong \prod_{\sigma} (j_\sigma)_*(\mathcal{F}|_{U_\sigma})\) over the
  multi-indices \(\sigma\), with each \(U_\sigma\) affine. A finite product of right
  \(f_*\)-acyclic objects is right \(f_*\)-acyclic
  (Lemma~\ref{mk-lem:rightAcyclic_finite_prod}), so it suffices to treat a single factor
  \((j_s)_*(\mathcal{F}|_{U_s})\), where \(s = (i_0, \ldots, i_p)\), \(U_s =
  U_{i_0 \ldots i_p}\) is affine, and \(j_s : U_s \hookrightarrow X\) is the open
  immersion. Acyclicity for \(f_*\) is a local question on \(S\): the derived
  pushforward \(R^k f_*\, \mathcal{G}\) is the sheaf associated to the presheaf
  \(V \mapsto H^k(f^{-1}(V), \mathcal{G}|_{f^{-1}(V)})\), so it vanishes for
  \(k \geq 1\) iff that cohomology vanishes for all affine \(V \subseteq S\).
  Restricting to such a \(V = \operatorname{Spec}(A)\), the base change \(f_V :
  f^{-1}(V) \to V\) is again separated and quasi-compact over the separated base
  \(V\), hence \(f^{-1}(V)\) is separated and the affine \(U_s\) meets it in the
  affine open \(U_s \cap f^{-1}(V)\).

  Now decompose the structure map \(g_s := f \circ j_s : U_s \to S\). The open
  immersion \(j_s : U_s \hookrightarrow X\) is the inclusion of an affine open into
  the separated scheme \(X\), hence an affine morphism. Composing, \(g_s\) is a
  morphism from the affine scheme \(U_s\) to the separated scheme \(S\), and a
  morphism from an affine scheme to a separated scheme is affine (Stacks Project,
  Tag~01SG); this is exactly the point at which separatedness of the \emph{target}
  \(S\) enters, and it fails over a non-separated base. The single factor of the
  product is the pushforward \((j_s)_*(\mathcal{F}|_{U_s})\), whose restriction
  \(\mathcal{F}|_{U_s}\) is quasi-coherent on \(U_s\) by
  Lemma~\ref{mk-lem:isQuasicoherent_pullback_opens}; identifying its higher direct
  image under \(f_*\) with the higher direct image of the quasi-coherent sheaf
  \(\mathcal{F}|_{U_s}\) along the affine morphism \(g_s\), the relative affine
  vanishing gives
  \[
    R^k f_*\bigl((j_s)_*(\mathcal{F}|_{U_s})\bigr)
      \;\cong\;
    R^k (g_s)_*(\mathcal{F}|_{U_s}) = 0
      \qquad (k \geq 1).
  \]
  That relative vanishing is itself a consequence of Serre vanishing
  (Lemma~\ref{mk-lem:affine_serre_vanishing}): \(R^k(g_s)_*\mathcal{G}\) is the
  sheafification of \(V \mapsto H^k(g_s^{-1}(V), \mathcal{G})\)
  (Lemma~\ref{mk-lem:higher_direct_image_presheaf}), and for affine \(V\) the preimage
  \(g_s^{-1}(V)\) is affine --- precisely because \(g_s\) is an affine morphism ---
  on which higher cohomology of a quasi-coherent sheaf vanishes. The derived-functor
  formation of these higher direct images uses the injective resolutions supplied by
  the hypothesis \(\mathrm{hres}\) on each intersection subscheme. Combining,
  \(R^k f_*(\mathcal{C}^p) = 0\) for all \(k \geq 1\).
\end{proof}

\begin{lemma}[{\v C}ech complex computes the higher direct images]\label{mk-lem:cech_computes_cohomology}
  \lean{AlgebraicGeometry.cech_computes_higherDirectImage}
  \leanok
  \dcref{02KE,custom}

  \uses{mk-lem:cech_augmented_resolution, mk-lem:cech_term_pushforward_acyclic, mk-lem:acyclic_resolution_computes_derived, mk-lem:pushforward_mapHC_cechComplexOnX, mk-lem:cechAugmented_to_acyclicResolutionInput, mk-def:cech_complex, mk-def:higher_direct_image}
  \textit{Source: Stacks Project, Cohomology of Schemes, Tag 02KE
  (\texttt{lemma-cech-cohomology-quasi-coherent}) and
  \texttt{lemma-quasi-coherence-higher-direct-images-application}.}
  Let \(f : X \to S\) be a quasi-compact \emph{separated} morphism of schemes with
  both the target \(S\) and the source \(X\) separated, let \(\mathcal{F}\) be a
  quasi-coherent \(\mathcal{O}_X\)-module, and let
  \(\mathfrak{U} : X = \bigcup_{i \in I} U_i\) be a finite open cover all of whose
  members \(U_i\) are affine (the hypothesis \(\mathrm{h}\mathfrak{U}\); so, by
  separatedness of \(X\), every intersection \(U_{i_0 \ldots i_p}\) is affine).
  Assume the category of \(\mathcal{O}_X\)-modules admits injective resolutions, and
  that for every degree \(n\) and every multi-index \(\sigma : \{0, \ldots, n\} \to I\)
  the category of \(\mathcal{O}\)-modules on the intersection subscheme \(U_\sigma\)
  admits injective resolutions (the hypothesis \(\mathrm{hres}\) inherited from
  Lemma~\ref{mk-lem:cech_term_pushforward_acyclic}). Then the cohomology sheaves of the
  relative {\v C}ech complex \(\check{\mathcal{C}}^\bullet(\mathfrak{U}, \mathcal{F})\)
  of Definition~\ref{mk-def:cech_complex} compute the higher direct images: for every
  \(i \geq 0\) there is a canonical isomorphism of \(\mathcal{O}_S\)-modules
  \[
    H^i\bigl(\check{\mathcal{C}}^\bullet(\mathfrak{U}, \mathcal{F})\bigr)
      \;\cong\; R^i f_*\mathcal{F}.
  \]
  In particular, over an affine base \(S = \operatorname{Spec}(A)\), taking global
  sections gives \(H^i(X, \mathcal{F}) = \check{H}^i(\mathfrak{U}, \mathcal{F}) =
  H^0(S, R^i f_*\mathcal{F})\) as \(A\)-modules for all \(i\).

  Two points about the precise form of the comparison should be noted. First, the
  derived-functor higher direct image \(R^i f_*\mathcal{F}\) on the right-hand side
  is the one obtained from injective resolutions, which is available only when the
  category of \(\mathcal{O}_X\)-modules has enough injectives; accordingly this
  comparison is asserted under the hypothesis that injective resolutions exist in
  \(\operatorname{QCoh}(X)\), and it compares the (always-defined) {\v C}ech higher
  direct image \(H^i(\check{\mathcal{C}}^\bullet(\mathfrak{U}, \mathcal{F}))\)
  against the derived-functor one only in that case. Second, we assert the comparison in the
  weak form of the \emph{existence} of an isomorphism, i.e.\ as
  \(\bigl(H^i(\check{\mathcal{C}}^\bullet(\mathfrak{U}, \mathcal{F})) \cong
  R^i f_*\mathcal{F}\bigr)\) without singling out a chosen natural isomorphism.
  This existence statement is all the downstream consumers (representability of the
  relative Picard functor, local triviality of the comparison) require: they use
  only that the unconditional {\v C}ech \(R^i f_*\) and the derived-functor
  \(R^i f_*\) agree as objects, never a specific identification, so committing to a
  canonical natural transformation would add coherence obligations that play no
  role in the applications.
\end{lemma}

\begin{proof}
  \leanok
  \uses{mk-lem:cech_augmented_resolution, mk-lem:cech_term_pushforward_acyclic, mk-lem:acyclic_resolution_computes_derived, mk-def:cech_complex, mk-def:higher_direct_image}
  Apply the abstract Lemma~\ref{mk-lem:acyclic_resolution_computes_derived} with the
  right-derivable additive functor \(G = f_*\)
  (the pushforward \(\operatorname{QCoh}(X) \to \operatorname{QCoh}(S)\)), the object
  \(A = \mathcal{F}\), and the resolution \(J^\bullet = \mathcal{C}^\bullet\) given by
  the {\v C}ech nerve of Definition~\ref{mk-def:cech_nerve}. The two hypotheses of the
  abstract lemma are exactly the geometric sub-lemmas established above:
  \begin{enumerate}
    \item By Lemma~\ref{mk-lem:cech_augmented_resolution} the augmented {\v C}ech
      complex \(0 \to \mathcal{F} \to \mathcal{C}^0 \to \mathcal{C}^1 \to \cdots\) is
      exact, i.e.\ \(\mathcal{C}^\bullet\) is a resolution of \(\mathcal{F}\) in
      \(\operatorname{QCoh}(X)\).
    \item By Lemma~\ref{mk-lem:cech_term_pushforward_acyclic} each term
      \(\mathcal{C}^p\) is right \(f_*\)-acyclic:
      \((f_*).\mathrm{rightDerived}\,k\,(\mathcal{C}^p) = 0\) for \(k \geq 1\).
  \end{enumerate}
  Lemma~\ref{mk-lem:acyclic_resolution_computes_derived} then gives, for every \(i\), an
  isomorphism between the \(i\)-th derived functor evaluated at \(\mathcal{F}\) and
  the \(i\)-th cohomology of the complex obtained by applying \(f_*\) to the
  resolution,
  \[
    (f_*).\mathrm{rightDerived}\,i\,(\mathcal{F})
      \;\cong\;
    H^i\bigl(f_*\,\mathcal{C}^\bullet\bigr).
  \]
  By construction (Definition~\ref{mk-def:cech_complex}) the complex \(f_*\,
  \mathcal{C}^\bullet\) is precisely the relative {\v C}ech complex
  \(\check{\mathcal{C}}^\bullet(\mathfrak{U}, \mathcal{F})\), so its \(i\)-th
  cohomology is \(H^i\bigl(\check{\mathcal{C}}^\bullet(\mathfrak{U}, \mathcal{F})
  \bigr)\); and \((f_*).\mathrm{rightDerived}\,i\,(\mathcal{F})\) is the
  derived-functor higher direct image \(R^i f_*\mathcal{F}\) of
  Definition~\ref{mk-def:higher_direct_image}. This yields the isomorphism
  \(H^i\bigl(\check{\mathcal{C}}^\bullet(\mathfrak{U}, \mathcal{F})\bigr) \cong
  R^i f_*\mathcal{F}\) asserted by the lemma; taking it in the existence
  (\(\operatorname{Nonempty}\)) form discards the choice of resolution, as required.
  Over an affine base \(S = \operatorname{Spec}(A)\), passing to global sections
  recovers \(H^i(X, \mathcal{F}) = \check{H}^i(\mathfrak{U}, \mathcal{F}) = H^0(S,
  R^i f_*\mathcal{F})\).
\end{proof}

\begin{lemma}[A finite product of right-\(f_*\)-acyclic objects is right-\(f_*\)-acyclic]\label{mk-lem:rightAcyclic_finite_prod}
  \lean{AlgebraicGeometry.rightAcyclic_finite_prod}
  \leanok
  \dcref{custom}

  \textit{Project-local.}
  Let \(G : \mathcal{A} \to \mathcal{B}\) be an additive functor between abelian
  categories that admits a right-derived functor, and let \((X_s)_{s \in \Sigma}\)
  be a \emph{finite} family of objects of \(\mathcal{A}\), each right
  \(G\)-acyclic --- \(G.\mathrm{rightDerived}\,k\,(X_s) = 0\) for all \(k \geq 1\)
  and all \(s\). Then the product \(\prod_{s \in \Sigma} X_s\) is right
  \(G\)-acyclic: \(G.\mathrm{rightDerived}\,k\,\bigl(\prod_s X_s\bigr) = 0\) for all
  \(k \geq 1\).
\end{lemma}
\begin{proof}
  \leanok

A finite product in an abelian category is a finite biproduct, and an additive
  right-derived functor preserves finite biproducts; hence for every \(k\) there is
  an isomorphism
  \[
    G.\mathrm{rightDerived}\,k\,\Bigl(\prod_{s} X_s\Bigr)
      \;\cong\;
    \prod_{s} G.\mathrm{rightDerived}\,k\,(X_s).
  \]
  For \(k \geq 1\) each factor on the right vanishes by the acyclicity hypothesis, so
  the product vanishes, and the isomorphism forces the left-hand side to vanish as well.
\end{proof}

\begin{lemma}
[Pushforward commutes with the {\v C}ech complex functor]
  \label{mk-lem:pushforward_mapHC_cechComplexOnX}
  \lean{AlgebraicGeometry.pushforward_mapHomologicalComplex_cechComplexOnX,
    AlgebraicGeometry.mapAlternatingCofaceMapComplexIso}
  \dcref{custom}

  \uses{mk-def:cech_complex_on_X, mk-def:cech_complex}
  \textit{Project-local.}
  Let \(f : X \to S\) be a morphism of schemes, \(\mathcal{F}\) a quasi-coherent
  \(\mathcal{O}_X\)-module and \(\mathfrak{U}\) a finite affine open cover of \(X\).
  There is a canonical isomorphism of cochain complexes in \(S.\mathrm{Modules}\)
  \[
    (f_*).\mathrm{mapHomologicalComplex}\bigl(\mathcal{C}^\bullet(\mathfrak{U},
      \mathcal{F})\bigr)
      \;\cong\;
    \check{\mathcal{C}}^\bullet(\mathfrak{U}, \mathcal{F})
  \]
  between the termwise image under the direct image \(f_*\) of the un-augmented
  {\v C}ech complex on \(X\) of Definition~\ref{mk-def:cech_complex_on_X} and the
  relative {\v C}ech complex of Definition~\ref{mk-def:cech_complex}.
\end{lemma}
\begin{proof}
  \uses{mk-def:cech_complex_on_X, mk-def:cech_complex}
  The direct image \(f_*\) is an additive functor, so it commutes with the functor
  that assembles a cosimplicial module into its alternating-coface cochain complex.
  Concretely, in each degree \(p\) both complexes have the same object, the image
  \(f_*(\mathcal{C}^p)\) of the degree-\(p\) {\v C}ech term; the degree-\(p\)
  differential of the relative {\v C}ech complex is exactly the image under \(f_*\) of
  the degree-\(p\) differential of the complex on \(X\). The latter differential is
  the alternating sum \(\sum_{j} (-1)^j \delta_j\) of the cofaces \(\delta_j\), and
  additivity of \(f_*\) --- its preservation of finite sums, of integer scalar
  multiples and of negation --- gives
  \[
    f_*\Bigl(\sum_j (-1)^j \delta_j\Bigr)
      \;=\;
    \sum_j (-1)^j\, f_*(\delta_j).
  \]
  Assembling these degreewise data --- the identity on objects together with the
  commuting squares expressing the differential identities --- produces the asserted
  isomorphism of cochain complexes.
\end{proof}

\begin{lemma}[From augmented exactness to the acyclic-resolution input data]\label{mk-lem:cechAugmented_to_acyclicResolutionInput}
  \lean{AlgebraicGeometry.cechAugmented_to_acyclicResolutionInput}
  \leanok
  \dcref{custom}

  \uses{mk-lem:cech_augmented_resolution, mk-def:cech_complex_on_X, mk-def:cech_augmented_complex}
  \textit{Project-local.}
  Suppose the augmented {\v C}ech complex of
  Definition~\ref{mk-def:cech_augmented_complex} is exact (its homology vanishes in every
  degree; Lemma~\ref{mk-lem:cech_augmented_resolution}). Then the un-augmented complex
  \(\mathcal{C}^\bullet\) on
  \(X\) of Definition~\ref{mk-def:cech_complex_on_X} carries exactly the data required to
  feed the abstract acyclic-resolution lemma:
  \begin{enumerate}
    \item an isomorphism \(e : \mathcal{F} \cong (\mathcal{C}^\bullet).\mathrm{cycles}\,0\)
      identifying \(\mathcal{F}\) with the \(0\)-cocycles, and
    \item exactness in every positive degree, \((\mathcal{C}^\bullet).\mathrm{ExactAt}\,(n+1)\)
      for all \(n\).
  \end{enumerate}
\end{lemma}
\begin{proof}
  \leanok
  \uses{mk-lem:cech_augmented_resolution, mk-def:cech_complex_on_X, mk-def:cech_augmented_complex}
  The augmented complex is the complex obtained by prepending \(\mathcal{F}\) to
  \(\mathcal{C}^\bullet\) along the augmentation \(\varepsilon\): its degree-\(0\) term
  is \(\mathcal{F}\) and its degree-\((p+1)\) term is \(\mathcal{C}^p\). This
  augmentation shifts homological degree by one, so the homology of the augmented
  complex at degree \(n+2\) coincides with the homology of the un-augmented complex
  \(\mathcal{C}^\bullet\) at degree \(n+1\). Vanishing of the augmented homology in
  these degrees therefore gives \((\mathcal{C}^\bullet).\mathrm{ExactAt}\,(n+1)\) for
  every \(n\), which is (2).

  For (1) consider the bottom of the augmented complex,
  \(0 \to \mathcal{F} \xrightarrow{\varepsilon} \mathcal{C}^0 \xrightarrow{d^0}
  \mathcal{C}^1\). Vanishing of the augmented homology in degree \(0\) says that
  \(\varepsilon\) has trivial kernel, i.e.\ \(\varepsilon\) is a monomorphism;
  vanishing in degree \(1\) says that the image of \(\varepsilon\) is exactly the
  kernel of \(d^0\), i.e.\ \(\varepsilon\) is an epimorphism onto the \(0\)-cocycles
  \((\mathcal{C}^\bullet).\mathrm{cycles}\,0 = \ker d^0\). A monomorphism that is also
  an epimorphism onto \(\mathrm{cycles}\,0\) is an isomorphism, giving the desired
  \(e : \mathcal{F} \cong (\mathcal{C}^\bullet).\mathrm{cycles}\,0\).
\end{proof}

\section{The unconditional higher direct image}

\begin{definition}[Unconditional higher direct image via {\v C}ech]\label{mk-def:cech_higher_direct_image}
  \lean{AlgebraicGeometry.cechHigherDirectImage}
  \leanok
  \dcref{custom}

  \uses{mk-lem:cech_computes_cohomology}
  \textit{Source: Stacks Project, Cohomology of Schemes,
  \texttt{lemma-quasi-coherence-higher-direct-images-application}.}
  Let \(f : X \to S\) be a separated quasi-compact morphism and \(\mathcal{F}\) a
  quasi-coherent \(\mathcal{O}_X\)-module. Fix a finite affine open cover
  \(\mathfrak{U}\) of \(X\). The \emph{(unconditional) \(i\)-th higher direct
  image} is defined, for each \(i \geq 0\), as the \(i\)-th cohomology sheaf of the
  relative {\v C}ech complex,
  \[
    R^i f_*\mathcal{F} \;:=\;
      H^i\bigl(\check{\mathcal{C}}^\bullet(\mathfrak{U}, \mathcal{F})\bigr)
      \in \operatorname{QCoh}(S).
  \]
  This requires \emph{no} hypothesis that the category of \(\mathcal{O}_X\)-modules
  have enough injectives: the right-hand side is the cohomology of an explicit
  complex of quasi-coherent sheaves. By
  Lemma~\ref{mk-lem:cech_computes_cohomology} the result is canonically isomorphic to
  the derived-functor higher direct image wherever the latter is defined, and in
  particular is independent of the chosen affine cover \(\mathfrak{U}\) up to
  canonical isomorphism. For \(i = 0\) one recovers the ordinary
  pushforward \(R^0 f_*\mathcal{F} = f_*\mathcal{F}\).
\end{definition}

\section{The relative {\v C}ech-to-derived-pushforward spectral sequence}
\label{mk-sec:cech_to_derived_pushforward_ss}

This section assembles the relative {\v C}ech-to-derived-pushforward spectral
sequence
\[
  E_2^{p,q} = \check{H}^p\bigl(\mathfrak{U}, R^q f_* \mathcal{F}\bigr)
  \;\Longrightarrow\;
  R^{p+q} f_* \mathcal{F},
\]
functorial in \(\mathcal{F}\), for a separated quasi-compact morphism
\(f : X \to S\) and a finite affine open cover \(\mathfrak{U}\) of \(X\). It is the
relative (over the base \(S\)) version of the absolute {\v C}ech-to-cohomology
spectral sequence; the absolute statement is cited verbatim in
\ref{mk-lem:cech_to_derived_pushforward_ss}, the relativization replacing the global
sections functor by \(f_*\) and absolute cohomology \(H^q\) by the derived
pushforward \(R^q f_*\).

Mathlib supplies only the \emph{abstract} spectral-sequence scaffold (the
total-complex and spectral-object API recorded as Mathlib anchors below); it does
\emph{not} contain the Grothendieck spectral sequence as a black box. The
construction is therefore decomposed into the bicomplex of a {\v C}ech complex of an
injective resolution, its total complex, the two filtrations whose associated graded
pages give respectively \(R^n f_*\mathcal{F}\) and
\(\check{H}^p(\mathfrak{U}, R^q f_*\mathcal{F})\), and the final assembly through the
Mathlib bridge.

\subsection{Mathlib anchors for the abstract spectral-sequence scaffold}

\begin{lemma}[Total complex of a bicomplex]
  \label{mk-lem:total_complex_mathlib}
  \lean{HomologicalComplex₂.total, HomologicalComplex₂.total.mapIso}
  \mathlibok
  \dcref{custom}

  For a double complex \(K\) in a preadditive category with the requisite coproducts,
  Mathlib forms the associated total complex \(\operatorname{Tot}(K)\) and, for an
  isomorphism \(e : K \cong L\) of double complexes, a functorial isomorphism
  \(\operatorname{Tot}(K) \cong \operatorname{Tot}(L)\) of the totals. This is the
  total-complex half of the abstract scaffold on which the spectral object below is
  built.
\end{lemma}

\begin{lemma}[Spectral sequence of a spectral object]
  \label{mk-lem:spectralObject_spectralSequence_mathlib}
  \lean{CategoryTheory.Abelian.SpectralObject,
        CategoryTheory.Abelian.SpectralObject.spectralSequence,
        CategoryTheory.CohomologicalSpectralSequenceNat}
  \mathlibok
  \dcref{custom}

  A \emph{spectral object} in an abelian category \(C\) (Mathlib's
  \(\operatorname{SpectralObject}\)) packages the filtration data of a filtered
  cochain complex; from one, together with a page datum, Mathlib produces a genuine
  spectral sequence via \(\operatorname{SpectralObject.spectralSequence}\), and in the
  cohomological-quadrant indexing it is a
  \(\operatorname{CohomologicalSpectralSequenceNat}\). Applied to the two filtrations
  of a total complex this is exactly the abstract machine that turns a filtered total
  complex into a convergent first-quadrant spectral sequence; the project supplies the
  filtered total complex (\ref{mk-def:cech_resolution_total_complex}) and the page
  identifications (\ref{mk-lem:total_complex_first_filtration_degenerates},
  \ref{mk-lem:total_complex_second_filtration_E2}).
\end{lemma}

\subsection{The bicomplex and its total complex}

\begin{definition}[{\v C}ech bicomplex of an injective resolution]
  \label{mk-def:cech_resolution_bicomplex}
  \dcref{03OW,custom}
  \uses{mk-def:relative_cech_complex_of_nerve, mk-def:cech_free_presheaf_complex,
        mk-lem:enoughInjectives_of_hasInjectiveResolutions}
  \textit{Source: Stacks Project, Cohomology, Tag 03OW
  (\texttt{lemma-cech-spectral-sequence}); the bicomplex is the explicit model of the
  Grothendieck spectral sequence used in its proof.}
  Choose an injective resolution \(\mathcal{F} \to \mathcal{I}^\bullet\) in the
  category of \(\mathcal{O}_X\)-modules (it exists by
  \ref{mk-lem:enoughInjectives_of_hasInjectiveResolutions}). Applying the relative
  {\v C}ech-complex construction \ref{mk-def:relative_cech_complex_of_nerve} of the cover
  \(\mathfrak{U}\) to each \(\mathcal{I}^q\) yields a double complex
  \[
    C^{p,q} \;=\; \check{\mathcal{C}}^p\bigl(\mathfrak{U}, \mathcal{I}^q\bigr)
    \in \operatorname{QCoh}(S),
  \]
  with horizontal differential the {\v C}ech (alternating-coface) differential and
  vertical differential induced by \(\mathcal{I}^\bullet\). Each row \(C^{\bullet,q}\)
  is the relative {\v C}ech complex of \(\mathcal{I}^q\); each column \(C^{p,\bullet}\)
  is \(f_*\) of the push--pull nerve term at level \(p\) applied to the resolution
  \(\mathcal{I}^\bullet\).
\end{definition}

\begin{definition}[Total complex of the {\v C}ech bicomplex]
  \label{mk-def:cech_resolution_total_complex}
  \dcref{custom}
  \uses{mk-def:cech_resolution_bicomplex, mk-lem:total_complex_mathlib}
  The total complex \(\operatorname{Tot}^\bullet(C^{\bullet,\bullet})\) of the
  bicomplex \ref{mk-def:cech_resolution_bicomplex}, formed via Mathlib's
  \(\texttt{HomologicalComplex}_2.\texttt{total}\)
  (\ref{mk-lem:total_complex_mathlib}). Carrying both the row filtration (by the
  {\v C}ech degree \(p\)) and the column filtration (by the resolution degree \(q\)),
  it is the single filtered complex whose two filtrations produce the two pages of the
  spectral sequence.
\end{definition}

\subsection{The two filtrations}

\begin{lemma}[First filtration: degeneration onto \(R^n f_*\mathcal{F}\)]
  \label{mk-lem:total_complex_first_filtration_degenerates}
  \dcref{custom}
  \uses{mk-def:cech_resolution_total_complex, mk-def:cech_higher_direct_image,
        mk-lem:cech_term_pushforward_acyclic, mk-lem:map_hc_homology_iso}
  \textit{Source: Stacks Project, Cohomology,
  \texttt{lemma-cech-spectral-sequence-application} (the degeneration corollary).}
  The filtration of \(\operatorname{Tot}^\bullet(C^{\bullet,\bullet})\) by the
  resolution degree \(q\) has associated spectral object whose first page is computed
  by the columns. Because each {\v C}ech term is a finite product of pushforwards of
  injectives along affine intersections, it is \(f_*\)-acyclic
  (\ref{mk-lem:cech_term_pushforward_acyclic}); hence taking column cohomology in the
  \(q\)-direction is concentrated in the resolution direction and recovers
  \(\mathcal{I}^\bullet\) up to homotopy, so the total complex computes the derived
  pushforward: \(H^n\bigl(\operatorname{Tot}^\bullet\bigr) \cong R^n f_*\mathcal{F}\)
  with \(R^n f_*\mathcal{F}\) as in \ref{mk-def:cech_higher_direct_image}. The
  exact-functor/homology interchange \ref{mk-lem:map_hc_homology_iso} supplies the
  termwise identifications.
\end{lemma}

\begin{lemma}[Second filtration: \(E_2 = \check{H}^p(\mathfrak{U}, R^q f_*\mathcal{F})\)]
  \label{mk-lem:total_complex_second_filtration_E2}
  \dcref{custom}
  \uses{mk-def:cech_resolution_total_complex, mk-lem:injective_cech_acyclic,
        mk-lem:cech_complex_hom_identification, mk-lem:mod_pmod_adjunction,
        mk-def:cech_higher_direct_image}
  \textit{Source: Stacks Project, Cohomology,
  \texttt{lemma-cech-cohomology-derived-presheaves} and
  \texttt{lemma-injective-trivial-cech}.}
  The filtration of \(\operatorname{Tot}^\bullet(C^{\bullet,\bullet})\) by the
  {\v C}ech degree \(p\) has \(E_2\) page
  \[
    E_2^{p,q} \;=\; \check{H}^p\bigl(\mathfrak{U}, R^q f_*\mathcal{F}\bigr).
  \]
  Taking row cohomology first computes, in each column \(q\), the {\v C}ech cohomology
  of the resolution sheaf \(\mathcal{I}^q\); since injectives are {\v C}ech-acyclic
  (\ref{mk-lem:injective_cech_acyclic}, the Grothendieck-SS hypothesis), the
  presheaf-level {\v C}ech functor \(\check{H}^0(\mathfrak{U}, -)\) is computed by the
  {\v C}ech complex \ref{mk-lem:cech_complex_hom_identification}, and its right derived
  functors are the {\v C}ech cohomology functors
  \(\check{H}^p(\mathfrak{U}, -) = R^p\check{H}^0(\mathfrak{U}, -)\) on presheaves
  (using the exact inclusion \(i\) of \ref{mk-lem:mod_pmod_adjunction}). Applying this
  with the column cohomology \(R^q f_*\mathcal{F}\)
  (\ref{mk-def:cech_higher_direct_image}) in the coefficient slot identifies the second
  page with \(\check{H}^p(\mathfrak{U}, R^q f_*\mathcal{F})\).
\end{lemma}

\subsection{Assembly of the spectral sequence}

\begin{lemma}[Relative {\v C}ech-to-derived-pushforward spectral sequence]
  \label{mk-lem:cech_to_derived_pushforward_ss}
  \dcref{03OW,custom}
  \uses{mk-lem:total_complex_first_filtration_degenerates,
        mk-lem:total_complex_second_filtration_E2,
        mk-lem:spectralObject_spectralSequence_mathlib,
        mk-def:cech_higher_direct_image}
  \textit{Source: Stacks Project, Cohomology, Tag 03OW
  (\texttt{lemma-cech-spectral-sequence}, ``{\v C}ech-to-cohomology spectral
  sequence''); stated here in its relative form over the base \(S\).}
  Let \(f : X \to S\) be separated and quasi-compact, \(\mathcal{F}\) a quasi-coherent
  \(\mathcal{O}_X\)-module, and \(\mathfrak{U}\) a finite affine open cover of \(X\).
  There is a first-quadrant cohomological spectral sequence, functorial in
  \(\mathcal{F}\),
  \[
    E_2^{p,q} = \check{H}^p\bigl(\mathfrak{U}, R^q f_*\mathcal{F}\bigr)
    \;\Longrightarrow\; R^{p+q} f_*\mathcal{F},
  \]
  with \(R^\bullet f_*\) the unconditional {\v C}ech higher direct images of
  \ref{mk-def:cech_higher_direct_image}.
\end{lemma}
\begin{proof}
  \uses{mk-def:cech_resolution_total_complex,
        mk-lem:total_complex_first_filtration_degenerates,
        mk-lem:total_complex_second_filtration_E2,
        mk-lem:spectralObject_spectralSequence_mathlib}
  Form the {\v C}ech bicomplex \(C^{\bullet,\bullet}\) of an injective resolution
  \(\mathcal{F} \to \mathcal{I}^\bullet\) (\ref{mk-def:cech_resolution_bicomplex}) and its
  total complex \(\operatorname{Tot}^\bullet\) (\ref{mk-def:cech_resolution_total_complex}).
  Feeding the two filtrations of \(\operatorname{Tot}^\bullet\) into Mathlib's
  spectral-object machine (\ref{mk-lem:spectralObject_spectralSequence_mathlib}) yields
  two spectral sequences with the same abutment
  \(H^{p+q}(\operatorname{Tot}^\bullet)\). The column (resolution-degree) filtration
  degenerates at its second page onto \(R^n f_*\mathcal{F}\)
  (\ref{mk-lem:total_complex_first_filtration_degenerates}), identifying the abutment with
  \(R^{p+q} f_*\mathcal{F}\). The row ({\v C}ech-degree) filtration has second page
  \(\check{H}^p(\mathfrak{U}, R^q f_*\mathcal{F})\)
  (\ref{mk-lem:total_complex_second_filtration_E2}). The resulting spectral sequence is
  the asserted one; functoriality in \(\mathcal{F}\) is inherited from functoriality of
  the injective resolution, the {\v C}ech complex, and the total-complex/spectral-object
  constructions. This is the relative incarnation of the Grothendieck spectral sequence
  for the composite \(\check{H}^0(\mathfrak{U}, -) \circ i\) recorded in the source.
\end{proof}

\section{Flat base change for the {\v C}ech higher direct images}

\begin{lemma}[Flat base change of $R^i f_*$ via {\v C}ech]
  \label{mk-lem:cech_flat_base_change}
  \lean{AlgebraicGeometry.cech_flatBaseChange}
  \dcref{02KH,custom}

  \uses{mk-def:cech_higher_direct_image, mk-lem:cech_computes_cohomology}
  \textit{Source: Stacks Project, Cohomology of Schemes, Tag 02KH
  (\texttt{lemma-flat-base-change-cohomology}, ``Flat base change'').}
  Consider the cartesian square
  \[
    \begin{array}{ccc}
      X' & \xrightarrow{\;g'\;} & X \\[2pt]
      {\scriptstyle f'}\big\downarrow & & \big\downarrow{\scriptstyle f} \\[2pt]
      S' & \xrightarrow{\;g\;} & S
    \end{array}
  \]
  with \(f : X \to S\) separated and quasi-compact, \(\mathcal{F}\) a
  quasi-coherent \(\mathcal{O}_X\)-module, and
  \(\mathcal{F}' = (g')^*\mathcal{F}\). Assume \(g\) is flat. Fix a finite affine
  open cover \(\mathfrak{U}\) of \(X\); the cover \(\mathfrak{U}'\) of \(X'\) is
  required to be the \emph{base change of \(\mathfrak{U}\) along \(g'\)} (the cover
  by the preimages \((g')^{-1}(U_i)\)) — the isomorphism is constructed only for
  this matched pair, not for an unrelated cover \(\mathfrak{U}'\). Then for every
  \(i \geq 0\) the canonical base-change map between the unconditional {\v C}ech
  higher direct images of Definition~\ref{mk-def:cech_higher_direct_image} is an
  isomorphism
  \[
    g^*\bigl(R^i f_*\mathcal{F}\bigr) \;\xrightarrow{\;\sim\;}\;
      R^i f'_*\bigl((g')^*\mathcal{F}\bigr) = R^i f'_*\mathcal{F}'.
  \]
  Equivalently, when \(S = \operatorname{Spec}(A)\) and
  \(S' = \operatorname{Spec}(B)\) with \(A \to B\) flat, the comparison map of
  \(B\)-modules \(H^i(X, \mathcal{F}) \otimes_A B \to H^i(X', \mathcal{F}')\) is an
  isomorphism.
\end{lemma}

\begin{proof}
  \uses{mk-lem:pullback_mapHC_homologyIso, mk-def:cech_complex_base_change_iso,
        mk-lem:cech_to_derived_pushforward_ss}
  This is the \emph{separated} case of Stacks Tag 02KH. For a separated quasi-compact
  \(f\) the relative {\v C}ech complex already computes \(R^\bullet f_*\) on the nose,
  so the separated case needs \emph{no} spectral sequence; the
  {\v C}ech-to-derived-pushforward spectral sequence
  (\ref{mk-lem:cech_to_derived_pushforward_ss}) enters only in the
  separated~\(\to\) general quasi-separated promotion. Write
  \(\check{\mathcal{C}}^\bullet = \check{\mathcal{C}}^\bullet(\mathfrak{U},\mathcal{F})\)
  for the relative {\v C}ech complex in \(\operatorname{QCoh}(S)\), so that
  \(R^i f_*\mathcal{F} = H^i(\check{\mathcal{C}}^\bullet)\) by
  Definition~\ref{mk-def:cech_higher_direct_image}. The proof is the two-step composite
  \[
    g^*\bigl(H^i(\check{\mathcal{C}}^\bullet)\bigr)
      \;\xrightarrow[\ \sim\ ]{(1)}\;
    H^i\bigl(g^*\check{\mathcal{C}}^\bullet\bigr)
      \;\xrightarrow[\ \sim\ ]{(2)}\;
    H^i\bigl(\check{\mathcal{C}}^\bullet(\mathfrak{U}', \mathcal{F}')\bigr)
      = R^i f'_*\mathcal{F}'.
  \]
  Step~(1) is Lemma~\ref{mk-lem:pullback_mapHC_homologyIso}: since \(g\) is flat the
  pullback \(g^*\) is exact, hence commutes with taking the \(i\)-th cohomology of a
  complex. Step~(2) is Definition~\ref{mk-def:cech_complex_base_change_iso}: applying
  \(g^*\) degreewise to \(\check{\mathcal{C}}^\bullet\) recovers the {\v C}ech complex
  of the base-changed data \((\mathfrak{U}', \mathcal{F}')\), the {\v C}ech complex
  transforming tensorially under base change (Stacks Tag 02KG). Functoriality of
  \(H^i\) (\texttt{HomologicalComplex.homologyMapIso}) turns the complex isomorphism
  of~(2) into the homology isomorphism. Composing the two gives the claim.
\end{proof}

\begin{remark}[Lift to the general quasi-separated case]
  \label{mk-rem:cech_flat_base_change_general}
  \dcref{02KH,custom}
  \uses{mk-lem:cech_to_derived_pushforward_ss, mk-lem:cech_flat_base_change,
        mk-def:cech_complex_base_change_iso}
  The same comparison extends from separated to quasi-separated morphisms by
  functoriality of the spectral sequence \ref{mk-lem:cech_to_derived_pushforward_ss}.
\end{remark}

\section{Homology-side machinery for {\v C}ech flat base change}

The {\v C}ech flat-base-change lemma factors through two ingredients: exactness of the
flat pullback \(g^*\) (so it commutes with cohomology) and the tensorial base change of
the {\v C}ech complex itself. We record both here. The homology-side pieces are
already proved; the genuine open content is isolated in the two clearly-flagged gaps below.

\begin{lemma}[Exact functor commutes with complex homology]\label{mk-lem:map_hc_homology_iso}
  \lean{AlgebraicGeometry.mapHomologicalComplexHomologyIso}
  \leanok
  \dcref{custom}

  Let \(F : \mathcal{C} \to \mathcal{D}\) be an additive functor between categories
  with homology that preserves homology (i.e.\ preserves kernels and cokernels). Then
  for every cochain complex \(K\) and index \(i\) there is a canonical isomorphism
  \[
    H^i\bigl((F\text{-applied degreewise to } K)\bigr) \;\cong\; F\bigl(H^i(K)\bigr).
  \]
  Indeed, the short complex defining degree-\(i\) homology after applying \(F\) is the
  image under \(F\) of the short complex defining \(H^i(K)\).
\end{lemma}

\begin{lemma}[Flat pullback preserves finite colimits]\label{mk-lem:pullback_preserves_finite_colimits}
  \lean{AlgebraicGeometry.pullback_preservesFiniteColimits}
  \leanok
  \dcref{custom}

  For any morphism \(g : S' \to S\) the pullback \(g^* : \operatorname{QCoh}(S) \to
  \operatorname{QCoh}(S')\) preserves finite colimits, being a left adjoint.
\end{lemma}

\subsection{Left exactness of flat pullback}
\label{mk-sec:pullback_pfl_decomposition}

Left exactness of \(g^*\) is not a statement about three functors but about one property:
monomorphisms. For an additive functor between abelian categories, right exactness together
with preservation of monomorphisms already forces left exactness
(Lemma~\ref{mk-lem:right_exact_mono_left_exact}); \(g^*\) is a left adjoint, so it is right exact
for free, and the entire flatness-dependent content is that \(g^*\) preserves injections
(Lemma~\ref{mk-lem:pullback_preserves_monos}). This is not the fallacy ``right exact implies
exact'': the mono hypothesis is exactly where flatness enters, and without flatness the
statement is false --- tensoring with a non-flat module is right exact and not left exact.

Monomorphisms of \(\mathcal{O}_X\)-modules are sectionwise
(Lemma~\ref{mk-lem:mono_iff_injective_sections}), which settles the case of an open immersion
outright (Lemma~\ref{mk-lem:pullback_preserves_monos_openImmersion}), because there \(g^*\) is the
restriction functor and restriction does not move sections. For a general flat \(g\) no
sectionwise formula for \(g^*\) is available and the argument must be stalkwise; the record of
that gap is Remark~\ref{mk-rem:pullback_mono_missing_stalk_model}.

\begin{lemma}[Right exact and mono-preserving implies left exact]\label{mk-lem:right_exact_mono_left_exact}
  \lean{AlgebraicGeometry.preservesFiniteLimits_of_preservesMonomorphisms}
  \leanok
  \dcref{custom}

  Let \(F : \mathcal{C} \to \mathcal{D}\) be an additive functor between abelian categories
  which preserves finite colimits and monomorphisms. Then \(F\) preserves finite limits.
\end{lemma}

\begin{proof}
  \leanokPreservation of finite limits is equivalent to: for every short exact sequence
  \(0 \to A \to B \to C \to 0\), the image \(F A \to F B \to F C\) is exact and \(F A \to F B\)
  is a monomorphism. The second clause is the hypothesis. For the first, \(B \to C\) is an
  epimorphism, so \(C\) is the cokernel of \(A \to B\); a right exact functor carries that
  cokernel presentation to a cokernel presentation, which is exactness of the image at \(F B\).
\end{proof}

\begin{lemma}[Monomorphisms of \(\mathcal{O}_X\)-modules are sectionwise]\label{mk-lem:mono_iff_injective_sections}
  \lean{AlgebraicGeometry.Modules.injective_app_of_mono}
  \leanok
  \dcref{custom}

  A morphism \(\varphi : M \to N\) of \(\mathcal{O}_X\)-modules is a monomorphism if and only if
  \(\Gamma(U, \varphi)\) is injective for every open \(U \subseteq X\).
\end{lemma}

\begin{proof}
  \leanokThe forgetful functor to abelian presheaves preserves limits, hence monomorphisms, and is
  faithful, hence reflects them; a morphism of abelian presheaves is a monomorphism precisely
  when it is one in every section, evaluation being limit-preserving; and a morphism of abelian
  groups is a monomorphism precisely when it is injective.
\end{proof}

\begin{lemma}[Flat pullback preserves monomorphisms]
  \label{mk-lem:pullback_preserves_monos}
  \lean{AlgebraicGeometry.pullback_preservesMonomorphisms}
  \dcref{custom}

  \uses{mk-lem:mono_iff_injective_sections}
  If \(g : S' \to S\) is flat then \(g^*\) preserves monomorphisms of
  \(\mathcal{O}_S\)-modules --- for arbitrary modules, not only quasi-coherent ones.
\end{lemma}

\begin{proof}
  \uses{mk-lem:mono_iff_injective_sections}
  Stalkwise. For \(x \in S'\) the stalk of \(g^* M\) is
  \(M_{g(x)} \otimes_{\mathcal{O}_{S, g(x)}} \mathcal{O}_{S', x}\); the stalk map
  \(\mathcal{O}_{S, g(x)} \to \mathcal{O}_{S', x}\) of a flat morphism is a flat ring map, and a
  flat base change preserves injections, so each stalk map of \(g^* \varphi\) is injective.
  Monomorphy of a morphism of sheaves of modules is detected on stalks.
\end{proof}

\begin{lemma}[Pullback along an open immersion preserves monomorphisms]\label{mk-lem:pullback_preserves_monos_openImmersion}
  \lean{AlgebraicGeometry.Modules.pullback_preservesMonomorphisms_of_isOpenImmersion}
  \leanok
  \dcref{custom}

  \uses{mk-lem:mono_iff_injective_sections}
  If \(j : U \to X\) is an open immersion then \(j^*\) preserves monomorphisms.
\end{lemma}

\begin{proof}
  \leanokAlong an open immersion the pullback is isomorphic to the restriction functor, and restriction
  does not move sections: \(\Gamma(V, M|_U) = \Gamma(j(V), M)\), with the restricted morphism
  acting as \(\Gamma(j(V), \varphi)\). So injectivity on sections transfers verbatim and
  Lemma~\ref{mk-lem:mono_iff_injective_sections} concludes in both directions.
\end{proof}

\begin{lemma}[Cancellation of preserved limits]\label{mk-lem:preserves_limit_comp_cancel}
  \lean{AlgebraicGeometry.preservesLimit_comp_cancel}
  \leanok
  \dcref{custom}

  Let \(K : \mathcal{J} \to \mathcal{C}\) be a diagram admitting a limit, and let
  \(F : \mathcal{C} \to \mathcal{D}\), \(G : \mathcal{D} \to \mathcal{E}\) be functors. If \(F\)
  preserves the limit of \(K\) and \(F \circ G\) preserves the limit of \(K\), then \(G\)
  preserves the limit of the transported diagram \(K \circ F\).
\end{lemma}

\begin{proof}
  \leanokLet \(c\) be a limit cone of \(K\). Since \(F\) preserves that limit, \(F c\) is a limit cone of
  \(K \circ F\); a functor preserves the limit of a diagram as soon as it carries one limit cone
  of that diagram to a limit cone, so it suffices to check \(G\) on \(F c\). That \(G (F c)\) is a
  limit cone is precisely the hypothesis that \(F \circ G\) preserves the limit of \(K\).
\end{proof}

\begin{lemma}[Flat affine pullback is exact on quasi-coherent modules]\label{mk-lem:flat_pullback_kernels_qcoh}
  \lean{AlgebraicGeometry.pullback_preservesKernel_of_isQuasicoherent}
  \leanok
  \dcref{custom}

  \uses{mk-lem:preserves_limit_comp_cancel, mk-lem:tilde_preserves_kernels}
  Let \(g : S' \to S\) be a flat morphism of affine schemes and \(\psi : A \to B\) a morphism of
  quasi-coherent \(\mathcal{O}_S\)-modules. Then \(g^*\) preserves \(\ker \psi\).
\end{lemma}

\begin{proof}
  \leanok\uses{mk-lem:preserves_limit_comp_cancel, mk-lem:tilde_preserves_kernels}
  First let \(S = \operatorname{Spec} R\), \(S' = \operatorname{Spec} R'\) and \(g\) be
  \(\operatorname{Spec}\) of a flat ring map \(\varphi\). The composite \(M \mapsto g^*(\widetilde
  M)\) is exact, and \(\widetilde{(-)}\) is exact, so by
  Lemma~\ref{mk-lem:preserves_limit_comp_cancel} the functor \(g^*\) preserves the limit of any
  finite diagram of the form \(K \circ \widetilde{(-)}\); taking \(K\) a parallel pair gives that
  \(g^*\) preserves \(\ker(\widetilde f)\) for every morphism \(f\) of \(R\)-modules. Now
  \(\widetilde{(-)}\) is fully faithful, so a morphism between two objects of its essential image
  becomes, after conjugating by the two witnessing isomorphisms, equal to \(\widetilde f\) for an
  actual morphism \(f\) of modules; kernels are preserved along an isomorphism of parallel pairs,
  so the claim holds for every morphism of quasi-coherent modules, these being exactly the objects
  of the essential image over an affine scheme.

  For general affine \(S, S'\) transport along the isomorphisms \(S \cong \operatorname{Spec}
  \Gamma(S, \mathcal{O}_S)\): the section map \(\Gamma(g)\) is flat, quasi-coherence is stable
  under pullback, and the conjugating functors are pullbacks along isomorphisms, hence
  equivalences. An equivalence preserves the limit and, being fully faithful, also reflects it, so
  preservation of the kernel descends from the conjugate to \(g^*\) itself.
\end{proof}

\begin{lemma}[Homology comparison from a single kernel]\label{mk-lem:homology_iso_of_preserves_kernel}
  \lean{AlgebraicGeometry.mapHomologicalComplexHomologyIso_of_preservesKernel}
  \leanok
  \dcref{custom}

  Let \(F\) be an additive, right exact functor between abelian categories, \(K\) a homological
  complex and \(i\) a degree. If \(F\) preserves the kernel of the differential out of \(K_i\),
  then \(H^i(F K) \cong F(H^i K)\).
\end{lemma}

\begin{proof}
  \leanokThe degree-\(i\) homology of \(K\) is that of the short complex
  \(K_{i-1} \to K_i \to K_{i+1}\), and the corresponding short complex of \(F K\) is the image of
  that one under \(F\). A left homology datum of a short complex is preserved by \(F\) as soon as
  \(F\) preserves the kernel of its right-hand map and the cokernel of the induced map from the
  left-hand object; the latter is automatic for a right exact \(F\). Preservation of one left
  homology datum gives the canonical comparison isomorphism.
\end{proof}

\begin{remark}[Flat exactness is needed only on quasi-coherent modules]
  \label{mk-rem:flat_exactness_qcoh_suffices}
  \dcref{custom}
  \uses{mk-lem:flat_pullback_kernels_qcoh, mk-lem:homology_iso_of_preserves_kernel}
  Lemma~\ref{mk-lem:pullback_preserves_monos}, which quantifies over arbitrary
  \(\mathcal{O}_S\)-modules, is not needed for flat base change of the \v{C}ech higher direct
  images. The higher direct images are the homology of a complex whose terms are quasi-coherent,
  and by Lemma~\ref{mk-lem:homology_iso_of_preserves_kernel} the comparison of homology with \(g^*\)
  requires only that \(g^*\) preserve one kernel in each degree --- which
  Lemma~\ref{mk-lem:flat_pullback_kernels_qcoh} supplies. The general statement remains open, and
  Remark~\ref{mk-rem:pullback_mono_missing_stalk_model} records why; the point here is that it is not
  on the path to Lemma~\ref{mk-lem:cech_flat_base_change}.
\end{remark}

\begin{remark}[No stalk model of the module pullback]
  \label{mk-rem:pullback_mono_missing_stalk_model}
  \dcref{custom}
  The stalk formula used in the proof of Lemma~\ref{mk-lem:pullback_preserves_monos} is the one
  piece of the flat-base-change chapter with no counterpart in the ambient library. The pullback
  of sheaves of modules is characterised there as a left adjoint of the pushforward, with no
  pointwise description; the factorisation of
  Lemma~\ref{mk-lem:sheafOfModules_pullbackIso} does not help, since its middle factor
  (presheaf-level pullback) is characterised the same way, while its outer two factors are
  unconditionally left exact and so were never the obstruction. The affine section formula of the
  Picard chapter (\texttt{Scheme.Modules.pullback\_app\_isoTensor}) computes
  \(\Gamma(U, g^* N)\) as a tensor product but only for quasi-coherent \(N\), whereas the
  statement above quantifies over all modules; over
  an affine base the category of \(\mathcal{O}_S\)-modules is strictly larger than the category
  of \(\Gamma(S, \mathcal{O}_S)\)-modules, so quasi-coherence cannot be dropped from that route.
\end{remark}

The three factors of the pullback, recorded because the older route through them is still the
one a reader is likely to try.

\begin{lemma}[Factorisation of the sheaf-of-modules pullback]
  \label{mk-lem:sheafOfModules_pullbackIso}
  \lean{SheafOfModules.pullbackIso}
  \mathlibok
  \dcref{custom}

  For a continuous morphism of ringed sites carried by a ring-sheaf hom
  \(\varphi\), the pullback of sheaves of modules factors as the composite of three
  functors
  \[
    \texttt{SheafOfModules.pullback}\,\varphi \;\cong\;
      \operatorname{forget} \;\circ\;
      \bigl(\texttt{PresheafOfModules.pullback}\,\varphi_{\mathrm{hom}}\bigr) \;\circ\;
      \texttt{PresheafOfModules.sheafification},
  \]
  i.e.\ forget the sheaf condition, take the presheaf-level pullback, then
  re-sheafify.
\end{lemma}

\begin{lemma}[Forgetful functor is left-exact]
  \label{mk-lem:forget_sheafOfModules_pfl_mathlib}
  \lean{SheafOfModules.forget}
  \mathlibok
  \dcref{custom}

  The forgetful functor \(\operatorname{forget}\) from sheaves of modules to the
  underlying sheaves of abelian groups preserves finite limits.
\end{lemma}

\begin{lemma}[Sheafification of presheaves of modules is left-exact]
  \label{mk-lem:sheafification_pfl_mathlib}
  \lean{PresheafOfModules.sheafification}
  \mathlibok
  \dcref{custom}

  Sheafification of presheaves of modules is a left-exact reflector, hence preserves
  finite limits.
\end{lemma}

\begin{lemma}[Flat extension of scalars is left-exact]
  \label{mk-lem:extendScalars_pfl_flat_mathlib}
  \lean{ModuleCat.preservesFiniteLimits_extendScalars_of_flat}
  \mathlibok
  \dcref{custom}

  If \(A \to B\) is a flat ring homomorphism then the extension-of-scalars functor
  \(\operatorname{extendScalars} : \operatorname{Mod}_A \to \operatorname{Mod}_B\),
  \(M \mapsto B \otimes_A M\), preserves finite limits.
\end{lemma}

The two outer factors are left exact by
Lemmas~\ref{mk-lem:forget_sheafOfModules_pfl_mathlib}
and~\ref{mk-lem:sheafification_pfl_mathlib}. It remains to treat presheaf pullback.

\begin{definition}[Inverse-image model of presheaf pullback]
  \label{mk-def:presheaf_pullback_inverseImage_model}
  \dcref{custom}
  The inverse-image model for presheaf-level pullback along \(\varphi_{\mathrm{hom}}\) is
  the inverse-image functor \(g^{-1}\) on presheaves of modules followed by degreewise
  extension of scalars along the underlying flat ring map, i.e.\ the composite
  \(g^{-1} \;\circ\; \operatorname{extendScalars}\), assembled sectionwise from the
  presheaf inverse image and extension of scalars.
\end{definition}

\begin{lemma}[The inverse-image model computes presheaf pullback]
  \label{mk-lem:presheaf_pullback_inverseImage_adjoint}
  \dcref{custom}
  \uses{mk-def:presheaf_pullback_inverseImage_model}
  The model \(g^{-1} \circ \operatorname{extendScalars}\) of
  Definition~\ref{mk-def:presheaf_pullback_inverseImage_model} is left adjoint to
  \(\texttt{PresheafOfModules.pushforward}\,\varphi_{\mathrm{hom}}\); by uniqueness of
  adjoints it is therefore canonically isomorphic to
  \(\texttt{PresheafOfModules.pullback}\,\varphi_{\mathrm{hom}}
  = (\texttt{PresheafOfModules.pushforward}\,\varphi_{\mathrm{hom}})^{\mathrm{L}}\).
  The unit and counit are obtained by composing, sectionwise, the
  inverse-image/direct-image adjunction with the extension-of-scalars/restriction-of-scalars
  adjunction. Their triangle identities follow from those of the constituent
  adjunctions, and uniqueness of left adjoints gives the stated isomorphism.
\end{lemma}

\begin{lemma}[Presheaf inverse image is exact]
  \label{mk-lem:presheaf_inverseImage_exact}
  \dcref{custom}
  The inverse-image functor \(g^{-1}\) on presheaves of modules is exact, hence
  preserves finite limits. It is computed by filtered colimits over the relevant comma
  categories. Filtered colimits of modules are exact, so they commute with the finite
  limits and colimits defining exact sequences.
\end{lemma}

\begin{lemma}[Presheaf pullback is left exact under flatness]
  \label{mk-lem:presheaf_pullback_pfl_flat}
  \dcref{custom}
  \uses{mk-lem:presheaf_pullback_inverseImage_adjoint, mk-lem:presheaf_inverseImage_exact,
        mk-lem:extendScalars_pfl_flat_mathlib}
  Let \(\varphi_{\mathrm{hom}}\) be the ring-sheaf hom underlying a flat morphism
  \(g\). Then \(\texttt{PresheafOfModules.pullback}\,\varphi_{\mathrm{hom}}\) preserves
  finite limits.
\end{lemma}

\begin{proof}
  \uses{mk-lem:presheaf_pullback_inverseImage_adjoint, mk-lem:presheaf_inverseImage_exact,
        mk-lem:extendScalars_pfl_flat_mathlib}
  By Lemma~\ref{mk-lem:presheaf_pullback_inverseImage_adjoint} it suffices to prove that
  the model \(g^{-1} \circ \operatorname{extendScalars}\) preserves finite limits. The
  inverse image \(g^{-1}\) is exact by
  Lemma~\ref{mk-lem:presheaf_inverseImage_exact}, hence preserves finite limits; the
  extension of scalars along the flat ring map preserves finite limits by
  Lemma~\ref{mk-lem:extendScalars_pfl_flat_mathlib}. A composite of two
  finite-limit-preserving functors preserves finite limits, and a functor isomorphic
  to such a composite does as well.
\end{proof}

\begin{lemma}[Flat pullback is left exact]
  \label{mk-lem:pullback_preserves_finite_limits}
  \lean{AlgebraicGeometry.pullback_preservesFiniteLimits}
  \dcref{custom}

  \uses{mk-lem:right_exact_mono_left_exact, mk-lem:pullback_preserves_finite_colimits,
        mk-lem:pullback_preserves_monos}
  If \(g\) is flat then \(g^*\) preserves finite limits.
\end{lemma}

\begin{proof}
  \uses{mk-lem:right_exact_mono_left_exact, mk-lem:pullback_preserves_finite_colimits,
        mk-lem:pullback_preserves_monos}
  \(g^*\) is additive and preserves finite colimits, being a left adjoint
  (Lemma~\ref{mk-lem:pullback_preserves_finite_colimits}), and it preserves monomorphisms because
  \(g\) is flat (Lemma~\ref{mk-lem:pullback_preserves_monos}). Apply
  Lemma~\ref{mk-lem:right_exact_mono_left_exact}.
\end{proof}

\begin{remark}[The older three-factor route, and why it is not the one taken]
  \label{mk-rem:pullback_pfl_three_factor_route}
  \dcref{custom}
  \uses{mk-lem:sheafOfModules_pullbackIso, mk-lem:forget_sheafOfModules_pfl_mathlib,
        mk-lem:sheafification_pfl_mathlib, mk-lem:presheaf_pullback_pfl_flat}
  By Lemma~\ref{mk-lem:sheafOfModules_pullbackIso}, on the genuine scheme site \(g^*\)
  factors as the composite
  \(\operatorname{forget} \circ
  (\texttt{PresheafOfModules.pullback}\,\varphi_{\mathrm{hom}}) \circ
  \texttt{PresheafOfModules.sheafification}\). The outer two factors preserve finite
  limits unconditionally (Lemmas~\ref{mk-lem:forget_sheafOfModules_pfl_mathlib}
  and~\ref{mk-lem:sheafification_pfl_mathlib}, both supplied by Mathlib on the scheme
  site). The middle factor preserves finite limits because \(g\) is flat
  (Lemma~\ref{mk-lem:presheaf_pullback_pfl_flat}). A composite of
  finite-limit-preserving functors preserves finite limits, so \(g^*\) does.

  This is a correct argument and it is not the one used, because its middle input is no easier
  than the conclusion: presheaf-level pullback is characterised as a left adjoint with no
  pointwise description, so Lemma~\ref{mk-lem:presheaf_pullback_pfl_flat} needs the same stalk
  model that Remark~\ref{mk-rem:pullback_mono_missing_stalk_model} records as missing, while
  additionally carrying two sheafification hypotheses. The mono route reaches the same place
  through one statement instead of four.
\end{remark}

\begin{lemma}[Flat pullback preserves homology]
  \label{mk-lem:pullback_preserves_homology}
  \lean{AlgebraicGeometry.pullback_preservesHomology}
  \dcref{custom}

  \uses{mk-lem:pullback_preserves_finite_colimits, mk-lem:pullback_preserves_finite_limits}
  For \(g\) flat, \(g^*\) preserves homology. Derived from left-exactness and
  left-adjointness via \texttt{Functor.preservesHomologyOfExact}.
\end{lemma}

\begin{lemma}[Flat pullback commutes with {\v C}ech homology]\label{mk-lem:pullback_mapHC_homologyIso}
  \lean{AlgebraicGeometry.pullback_mapHC_homologyIso}
  \leanok
  \dcref{custom}

  \uses{mk-lem:map_hc_homology_iso, mk-lem:pullback_preserves_homology}
  For \(g\) flat, applying \(g^*\) to the \(i\)-th cohomology of a cochain complex of
  \(\mathcal{O}_S\)-modules agrees with the \(i\)-th cohomology of the degreewise
  pullback complex: \(g^*(H^i(K)) \cong H^i(g^* K)\). A direct specialisation of
  Lemma~\ref{mk-lem:map_hc_homology_iso} to \(g^*\).
\end{lemma}

\subsection{The {\v C}ech base-change isomorphism via a cosimplicial natural iso}
\label{mk-sec:cech_base_change_construction}

By Definition~\ref{mk-def:relative_cech_complex_of_nerve},
\(\check{\mathcal{C}}^\bullet(\mathfrak{U}, \mathcal{F}) =
\operatorname{alternatingCofaceMapComplex}\bigl((\operatorname{pushforward} f) \circ
\operatorname{CechNerve}(\mathfrak{U}, \mathcal{F})\bigr)\), whose degree-\(p\) term is
\((\operatorname{pushforward} f)\) of a push--pull nerve object and whose differentials
are \((\operatorname{pushforward} f)\) of cosimplicial coface maps. A natural
isomorphism of the underlying cosimplicial objects therefore induces an isomorphism
of {\v C}ech complexes by functoriality of
\(\operatorname{alternatingCofaceMapComplex}\).

\begin{lemma}[Additive functors commute with the alternating-coface differential]\label{mk-lem:map_alternatingCofaceMapComplex_objD}
  \lean{AlgebraicGeometry.map_alternatingCofaceMapComplex_objD}
  \leanok
  \dcref{custom}

  For an additive functor \(F\) between preadditive categories, \(F\) applied to the
  degree-\(p\) alternating-coface differential \(\operatorname{objD}\) of a cosimplicial
  object \(C\) agrees with the alternating-coface differential of \(F \circ C\): the
  differential is an alternating sum of coface maps, and \(F\) preserves finite sums and
  carries each coface map to the corresponding coface map of \(F \circ C\), so the two
  differentials coincide termwise.
\end{lemma}

\begin{lemma}[Flat pullback commutes with the alternating-coface complex]\label{mk-lem:pullback_commutes_alternatingCofaceMapComplex}
  \lean{AlgebraicGeometry.mapAlternatingCofaceMapComplexIso}
  \leanok
  \dcref{custom}

  \uses{mk-def:relative_cech_complex_of_nerve, mk-lem:map_alternatingCofaceMapComplex_objD}
  For an additive functor \(F\) between preadditive categories, \(F\) applied
  degreewise commutes with the alternating-coface-map cochain complex functor: there is
  a natural isomorphism
  \(F\bigl(\operatorname{alternatingCofaceMapComplex}(C)\bigr) \cong
  \operatorname{alternatingCofaceMapComplex}(F \circ C)\), natural in the cosimplicial
  object \(C\). In particular this holds for \(F = g^*\), which is additive.
  The degree-\(p\) terms agree canonically (\(F\) applied to the
  degree-\(p\) term of \(C\)); the differential is an alternating sum of coface maps,
  and since \(F\) is additive it commutes with finite sums and with each coface map, so
  the term-wise identifications assemble into a chain isomorphism.
\end{lemma}

\begin{lemma}[Degreewise affine reduction of the {\v C}ech base change]\label{mk-lem:cech_degree_affine_baseChange}
  \lean{AlgebraicGeometry.cech_degree_affine_baseChange}
  \leanok
  \dcref{custom}

  \uses{mk-lem:affine_pushforward_pullback_baseChange, mk-lem:pushforward_spec_tilde_iso,
        mk-lem:pullback_spec_tilde_iso, mk-lem:cancelBaseChange_mathlib,
        mk-def:cech_nerve_cosimplicial}
  Fix a cosimplicial degree \(p\). The \((p+1)\)-fold fibre power
  \(U_{i_0\dots i_p} = U_{i_0} \times_X \dots \times_X U_{i_p}\) appearing in the
  {\v C}ech nerve is, on the standard affine model of the cover, \(\operatorname{Spec}\)
  of the finite affine intersection \(A_{i_0} \otimes_R \dots \otimes_R A_{i_p}\) of the
  coordinate rings of the members \(U_{i_j} = \operatorname{Spec} A_{i_j}\). On that
  affine model the X-level cartesian square defining the base change along \(g'\)
  restricts to the affine pushout (tensor) square of rings, so the degreewise
  Beck--Chevalley comparison
  \(g'^*(p_* p^* F) \cong p'_* p'^*(g'^* F)\) over \(U_{i_0\dots i_p}\) IS the affine
  termwise base change \(\operatorname{affinePushforwardPullbackBaseChange}\)
  (\ref{mk-lem:affine_pushforward_pullback_baseChange}). These affine identifications are
  natural with respect to the index-omission maps that generate the cosimplicial
  structure of the nerve.
\end{lemma}

\begin{proof}
  \uses{mk-lem:affine_pushforward_pullback_baseChange, mk-lem:pushforward_spec_tilde_iso,
        mk-lem:pullback_spec_tilde_iso, mk-lem:cancelBaseChange_mathlib,
        mk-def:cech_nerve_cosimplicial}
  \emph{Spec-reduction.} On the standard affine model the members of the cover are
  \(U_{i_j} = \operatorname{Spec} A_{i_j}\) over the affine base \(\operatorname{Spec} R\),
  and the intersection \(U_{i_0\dots i_p}\) --- a fibre power over \(X\) --- is, by the
  affine description of fibre products, \(\operatorname{Spec}\) of the iterated tensor
  product \(A_{i_0} \otimes_R \dots \otimes_R A_{i_p}\). The base change \(g'\) is
  induced by a ring map \(R \to R'\); pulling back along \(g'\) replaces this affine
  intersection by \(\operatorname{Spec}\) of \((A_{i_0} \otimes_R \dots \otimes_R A_{i_p})
  \otimes_R R'\). Thus the X-level cartesian square restricts, over each intersection, to
  the affine pushout square of rings, and the degreewise base-change isomorphism is by
  construction \(\operatorname{affinePushforwardPullbackBaseChange}\)
  (\ref{mk-lem:affine_pushforward_pullback_baseChange}); the latter is assembled from the
  tilde dictionaries \ref{mk-lem:pushforward_spec_tilde_iso}, \ref{mk-lem:pullback_spec_tilde_iso}
  and the cancellation \(\operatorname{cancelBaseChange}\)
  (\ref{mk-lem:cancelBaseChange_mathlib}), \emph{not} from the canonical adjoint mate.

  \emph{Index-omission naturality.} Each coface map of the {\v C}ech nerve
  (\ref{mk-def:cech_nerve_cosimplicial}) omits one index \(i_j\), realised on rings as the
  inclusion \(A_{i_0} \otimes_R \dots \widehat{A_{i_j}} \dots \otimes_R A_{i_p} \to
  A_{i_0} \otimes_R \dots \otimes_R A_{i_p}\) (insertion of the missing tensor factor);
  codegeneracies repeat an index dually. The affine base-change isomorphism
  \(\operatorname{affinePushforwardPullbackBaseChange}\) is natural in the ring under
  consideration, so the degreewise comparisons commute with these omission/insertion
  maps and assemble into a morphism of cosimplicial objects. This is exactly the
  per-degree data that both \ref{mk-lem:cech_pushforward_baseChange_natIso} and the
  twisted-nerve identification consume.
\end{proof}

\subsection{The abstract-to-affine bridge for the base-change leaves}
\label{mk-sec:pushpull_tilde_bridge}

The brick \ref{mk-lem:cech_degree_affine_baseChange} speaks the affine
\(\widetilde{(-)}\)-model over \(\operatorname{Spec}\): it compares
\(g^*\bigl((\operatorname{Spec}\varphi)_*\widetilde{N}\bigr)\) with
\((\operatorname{Spec}\varphi')_*\bigl(g'^*\widetilde{N}\bigr)\) for an explicit module
\(N\). The two leaves below, however, present each degree as the \emph{abstract}
push--pull sheaf \(\operatorname{pushPullObj} \mathcal{F}\, Y_n = (p_n)_*\,p_n^*\mathcal{F}\)
(\ref{mk-def:push_pull_obj}) over a fibre power \(Y_n\), with no \(\widetilde{(-)}\) in
sight. The present subsection supplies the missing identification, at the two well-typed
altitudes at which it can be stated, and it is what lets each leaf hand its degreewise
square to the affine brick. The bridge rests on four already-available ingredients,
recorded first: the affine structure theorem in counit form, preservation of
quasi-coherence by pullback and by affine pushforward, and the affineness of fibre powers
of an affine open-immersion cover of a separated scheme.

\begin{lemma}[The counit-iso criterion in localizing form]
  \label{mk-lem:isIso_fromTildeGamma_iff_isLocalizing_mathlib}
  \lean{AlgebraicGeometry.isIso_fromTildeΓ_iff_isLocalizing}
  \mathlibok
  \dcref{custom}

  For an \(\mathcal{O}_{\operatorname{Spec} R}\)-module \(\mathcal{F}\), the tilde--\(\Gamma\)
  counit (the canonical map \(\widetilde{\Gamma(\mathcal{F})} \to \mathcal{F}\)) is an
  isomorphism if and only if \(\mathcal{F}\) is \emph{localizing}: for every \(f \in R\)
  the section-restriction map
  \(\Gamma(\operatorname{Spec} R, \mathcal{F}) \to \Gamma(D(f), \mathcal{F})\) exhibits its
  target as the localization of the global sections at the powers of \(f\). Combined with
  the keystone \ref{mk-lem:qcoh_section_isLocalizedModule}, this yields the unconditional
  01I8 instance \ref{mk-lem:qcoh_isIso_fromTildeGamma} directly.
\end{lemma}

\begin{lemma}[Affine pushforward preserves the localizing property]
  \label{mk-lem:isLocalizing_pushforward_mathlib}
  \lean{AlgebraicGeometry.isLocalizing_pushforward_of_isLocalizing}
  \mathlibok
  \dcref{custom}

  Let \(\varphi : R \to A\) be a ring homomorphism and let \(\mathcal{G}\) be a localizing
  (equivalently, \(\widetilde{(-)}\)-essential-image) \(\mathcal{O}_{\operatorname{Spec} A}\)-module.
  Then the pushforward \((\operatorname{Spec}\varphi)_*\mathcal{G}\) along the affine
  morphism \(\operatorname{Spec}\varphi : \operatorname{Spec} A \to \operatorname{Spec} R\)
  is again localizing; equivalently its tilde--\(\Gamma\) counit is an isomorphism, so
  \((\operatorname{Spec}\varphi)_*\widetilde{N} \cong \widetilde{N|_R}\) with \(N|_R\) the
  restriction of scalars of \(N\) along \(\varphi\).
\end{lemma}

\begin{lemma}[Finite intersections of affine opens of a separated scheme are affine]
  \label{mk-lem:isAffineOpen_biInf_mathlib}
  \lean{AlgebraicGeometry.IsAffineOpen.biInf}
  \mathlibok
  \dcref{custom}

  Let \(X\) be a scheme whose diagonal is affine --- in particular any \emph{separated}
  scheme. Then any finite
  intersection \(U_{i_0} \cap \dots \cap U_{i_p}\) of affine open subschemes of \(X\) is
  again an affine open. Consequently, for an affine open-immersion cover
  \(\mathfrak{U} = \{U_i = \operatorname{Spec} A_i\}\) of a separated scheme, every
  \((p+1)\)-fold fibre power \(Y_n = U_{i_0} \times_X \dots \times_X U_{i_p} =
  U_{i_0} \cap \dots \cap U_{i_p}\) of the {\v C}ech nerve is affine,
  \(Y_n = \operatorname{Spec} A_n\).
\end{lemma}

\begin{lemma}[{\v C}ech intersection opens are affine, separated case]\label{mk-lem:cech_interopen_isaffine}
  \lean{AlgebraicGeometry.coverInterOpen_isAffine}
  \leanok
  \dcref{custom}

  \uses{mk-lem:isAffineOpen_biInf_mathlib}
  Let \(f : X \to S\) be a \emph{separated} morphism with affine base
  \(S = \operatorname{Spec} R\), and let \(\mathfrak{U} = \{U_i = \operatorname{Spec} A_i\}\)
  be an affine open-immersion cover of \(X\). Then for every \emph{finite, nonempty} index
  family \(\sigma : \kappa \to I\) the {\v C}ech intersection open
  \(\operatorname{coverInterOpen}\mathfrak{U}\,\sigma = \bigcap_k (U_{\sigma(k)})\) --- the
  meet of the ranges of the cover members indexed by \(\sigma\) --- is an affine open
  subscheme of \(X\).
\end{lemma}
\begin{proof}
  \uses{mk-lem:isAffineOpen_biInf_mathlib}
  Since \(S = \operatorname{Spec} R\) is affine it is a separated scheme, so its structure
  morphism to the terminal scheme is separated. The structure morphism of \(X\) factors as
  \(f\) followed by that of \(S\); a composite of separated morphisms is separated, so the
  structure morphism of \(X\) is separated, i.e.\ \(X\) is a separated scheme. Separatedness
  says the diagonal \(\Delta_X : X \to X \times X\) is a closed immersion, in particular an
  affine morphism, so \(X\) has affine diagonal. Consequently any finite \emph{nonempty}
  intersection of affine opens of \(X\) is again affine --- this is the indexed-meet form of
  \ref{mk-lem:isAffineOpen_biInf_mathlib}, where each successive meet is an affine open because
  it is cut out by the affine diagonal. Applying it to the nonempty finite family
  \(\{U_{\sigma(k)}\}_k\) shows \(\operatorname{coverInterOpen}\mathfrak{U}\,\sigma\) is
  affine. (Nonemptiness of \(\kappa\) is load-bearing: an empty meet would be all of \(X\),
  which need not be affine.)
\end{proof}

\begin{lemma}[Pullback preserves quasi-coherence (open-immersion case)]\label{mk-lem:pullback_preserves_quasicoherent}
  \lean{AlgebraicGeometry.pullback_isQuasicoherent}
  \leanok
  \dcref{01BG,custom}

  \uses{mk-lem:isQuasicoherent_pullback_opens}
  \textit{Source: Stacks Project, Modules on Sites, Tag 01BG, \texttt{lemma-pullback-quasi-coherent}.}
  Let \(\iota_V : V \hookrightarrow X\) be the open immersion of an open subscheme \(V\) of \(X\),
  and let \(\mathcal{F}\) be a quasi-coherent \(\mathcal{O}_X\)-module. Then the restriction
  pullback \(\iota_V^*\mathcal{F}\) is a quasi-coherent \(\mathcal{O}_V\)-module. In particular,
  each fibre power \(Y_n = U_{i_0} \cap \dots \cap U_{i_p}\) of the {\v C}ech nerve embeds in
  \(X\) by an open immersion \(p_n : Y_n \hookrightarrow X\), so \(p_n^*\mathcal{F}\) is
  quasi-coherent on \(Y_n\).
\end{lemma}
\begin{proof}
  \uses{mk-lem:isQuasicoherent_pullback_opens}
  This is the open-immersion instance of \ref{mk-lem:isQuasicoherent_pullback_opens}: the
  preimages of a quasi-coherence cover of \(X\) cover \(V\), and quasi-coherence descends
  along that cover. The {\v C}ech case follows since each fibre-power inclusion
  \(p_n : Y_n \hookrightarrow X\) is an open immersion.
\end{proof}

\begin{lemma}[The opens-preimage comma posets of a continuous map are filtered]\label{mk-lem:opensMap_structuredArrow_isFiltered}
  \lean{AlgebraicGeometry.opensMapStructuredArrow_isFiltered}
  \leanok
  \dcref{custom}

  Let \(\varphi : T \to T'\) be a continuous map of topological spaces and let \(d\) be an
  open subset of \(T\). The poset of open subsets \(U \subseteq T'\) with
  \(d \subseteq \varphi^{-1}(U)\) is filtered.
\end{lemma}
\begin{proof}
  \leanokThe poset is nonempty: \(T'\) itself satisfies \(d \subseteq \varphi^{-1}(T') = T\). It is
  directed: for \(U_1, U_2\) in the poset the union \(U_1 \cup U_2\) again lies in it, since
  \(d \subseteq \varphi^{-1}(U_1) \subseteq \varphi^{-1}(U_1 \cup U_2)\), and it dominates
  both. A poset is thin, so any parallel pair of maps is already equal, and a nonempty
  directed thin category is filtered.
\end{proof}

\begin{lemma}[The opens-preimage functor of a continuous map is final]\label{mk-lem:opensMap_final}
  \lean{AlgebraicGeometry.opensMap_final}
  \leanok
  \dcref{custom}

  \uses{mk-lem:opensMap_structuredArrow_isFiltered}
  For any continuous map \(\varphi : T \to T'\), the preimage functor
  \(U \mapsto \varphi^{-1}(U)\), from open subsets of \(T'\) to open subsets of \(T\), is a
  final functor.
\end{lemma}
\begin{proof}
  \leanok\uses{mk-lem:opensMap_structuredArrow_isFiltered}
  A functor is final as soon as, for every object \(d\) of the target, the comma category of
  maps from \(d\) into its image is connected; by
  \ref{mk-lem:opensMap_structuredArrow_isFiltered} these comma categories are filtered, and
  filtered categories are connected.
\end{proof}

\begin{lemma}[Pullback of the structure sheaf]\label{mk-lem:module_pullback_unit_iso}
  \lean{AlgebraicGeometry.pullbackObjUnitToUnit_isIso_hom, AlgebraicGeometry.pullbackUnitIso}
  \leanok
  \dcref{custom}

  \uses{mk-lem:opensMap_final}
  For any morphism of schemes \(g : Y \to X\), the canonical comparison morphism
  \(g^*\mathcal{O}_X \to \mathcal{O}_Y\) of sheaves of \(\mathcal{O}_Y\)-modules is an
  isomorphism.
\end{lemma}
\begin{proof}
  \leanok\uses{mk-lem:opensMap_final}
  By the pullback--pushforward adjunction, morphisms \(g^*\mathcal{O}_X \to \mathcal{M}\)
  correspond to morphisms \(\mathcal{O}_X \to g_*\mathcal{M}\), i.e.\ to global sections of
  \(g_*\mathcal{M}\). Since the opens-preimage functor underlying \(g\) is final
  (\ref{mk-lem:opensMap_final}), global sections of \(g_*\mathcal{M}\) coincide with global
  sections of \(\mathcal{M}\), which correspond to morphisms \(\mathcal{O}_Y \to \mathcal{M}\).
  Thus \(g^*\mathcal{O}_X\) and \(\mathcal{O}_Y\) corepresent the same functor, and the
  comparison morphism realizes this identification, hence is an isomorphism.
\end{proof}

\begin{lemma}[Per-slice presentation of a pullback]\label{mk-lem:presentation_pullback_slice_of_over}
  \lean{AlgebraicGeometry.presentationPullbackSliceOfOver}
  \leanok
  \dcref{custom}

  \uses{mk-lem:module_pullback_unit_iso}
  Let \(g : Y \to X\) be a morphism of schemes, \(\mathcal{F}\) an \(\mathcal{O}_X\)-module,
  and \(A \subseteq X\) an open subset over which \(\mathcal{F}\) admits a presentation
  \(\bigoplus_{\iota}\mathcal{O} \to \bigoplus_{\sigma}\mathcal{O} \to \mathcal{F}|_A \to 0\).
  Then the pullback \(g^*\mathcal{F}\) admits a presentation over the preimage open
  \(W = g^{-1}(A)\).
\end{lemma}
\begin{proof}
  \leanok\uses{mk-lem:module_pullback_unit_iso, mk-lem:isQuasicoherent_pullback_opens}
  Let \(g_A : W \to A\) denote the restriction of \(g\), so that
  \(\iota_W \circ g = g_A \circ \iota_A\) with \(\iota_A, \iota_W\) the open inclusions.
  Pseudofunctoriality of the module pullback applied to this square identifies
  \((g^*\mathcal{F})|_W \cong g_A^*(\mathcal{F}|_A)\). The functor \(g_A^*\) is a left
  adjoint, hence preserves colimits, and it carries the structure sheaf to the structure
  sheaf (\ref{mk-lem:module_pullback_unit_iso}), hence free modules to free modules; applying it to
  the given presentation of \(\mathcal{F}|_A\) therefore yields a presentation of
  \(g_A^*(\mathcal{F}|_A) \cong (g^*\mathcal{F})|_W\). Carrying this presentation of the
  restriction back to the slice over \(W\) is the same restrict--over bridge used in
  \ref{mk-lem:isQuasicoherent_pullback_opens}.
\end{proof}

\begin{lemma}[Pullback preserves quasi-coherence (general morphism)]\label{mk-lem:pullback_isQuasicoherent_hom}
  \lean{AlgebraicGeometry.pullback_isQuasicoherent_hom}
  \leanok
  \dcref{01BG,custom}

  \uses{mk-lem:presentation_pullback_slice_of_over}
  \textit{Source: Stacks Project, Tag 01BG, \texttt{lemma-pullback-quasi-coherent}.}
  Let \(g : Y \to X\) be a morphism of schemes and \(\mathcal{F}\) a quasi-coherent
  \(\mathcal{O}_X\)-module. Then \(g^*\mathcal{F}\) is a quasi-coherent
  \(\mathcal{O}_Y\)-module.
\end{lemma}
\begin{proof}
  \leanok\uses{mk-lem:presentation_pullback_slice_of_over}
  Choose an open cover \(\{A_i\}\) of \(X\) over whose members \(\mathcal{F}\) admits
  presentations. The preimages \(g^{-1}(A_i)\) cover \(Y\), and over each of them
  \(g^*\mathcal{F}\) admits a presentation by
  \ref{mk-lem:presentation_pullback_slice_of_over}. Quasi-coherence is local on the base, so
  \(g^*\mathcal{F}\) is quasi-coherent.
\end{proof}

\begin{lemma}
[Abstract push--pull is the affine tilde-model on a fibre power]
  \label{mk-lem:pushPullObj_iso_tilde}
  \lean{AlgebraicGeometry.pullbackRestrict_iso_tilde,
        AlgebraicGeometry.pushPullObj_pushforward_iso_tilde}
  \dcref{01I8,custom}

  \uses{mk-def:push_pull_obj, mk-lem:qcoh_iso_tilde_sections, mk-lem:qcoh_isIso_fromTildeGamma,
        mk-lem:pullback_preserves_quasicoherent, mk-lem:pushforward_iso_preserves_qcoh,
        mk-lem:isLocalizing_pushforward_mathlib,
        mk-lem:isAffineOpen_biInf_mathlib, mk-lem:scheme_pushforwardComp_mathlib,
        mk-lem:isoSpec_scheme_mathlib}
  \textit{Source: Stacks Project, Schemes, Tag 01I8, \texttt{lemma-quasi-coherent-affine}.}
  Let \(f : X \to S\) be separated with affine base \(S = \operatorname{Spec} R\), let
  \(\mathfrak{U} = \{U_i = \operatorname{Spec} A_i\}\) be an affine open-immersion cover of
  \(X\), and fix a cosimplicial degree, writing \(Y_n\) for the corresponding fibre power
  with projection \(p_n : Y_n \to X\). By \ref{mk-lem:isAffineOpen_biInf_mathlib} the
  separatedness makes \(Y_n = \operatorname{Spec} A_n\) affine. There are two well-typed
  identifications of the abstract push--pull data with the affine \(\widetilde{(-)}\)-model.

  \emph{Altitude 1 (restriction level).} Since the formation of \(\widetilde{(-)}\) requires a
  genuine \(\operatorname{Spec}\), and the open subscheme \(Y_n\) is not literally the spectrum
  of a ring, we transport along the whole-scheme affine identification
  \(\sigma_n : Y_n \xrightarrow{\ \sim\ } \operatorname{Spec}\Gamma(X, Y_n)\). Pushing the
  restriction \(p_n^*\mathcal{F}\) forward along \(\sigma_n\) onto the genuine affine
  \(\operatorname{Spec}\Gamma(X, Y_n)\), and setting
  \(N_n := \Gamma\bigl(\operatorname{Spec}\Gamma(X, Y_n),\ (\sigma_n)_*\,p_n^*\mathcal{F}\bigr)\),
  \[
    (\sigma_n)_*\,p_n^*\mathcal{F} \;\cong\; \widetilde{N_n}
    \qquad\text{(over } \operatorname{Spec}\Gamma(X, Y_n)\text{).}
  \]

  \emph{Altitude 2 (assembled pushed-forward level).} Over the affine base \(S\), with
  \(\varphi_n : R \to A_n\) the ring map presenting the composite \(f \circ p_n\) as
  \(\operatorname{Spec}\varphi_n\),
  \[
    (\operatorname{pushforward} f).\!\operatorname{obj}\bigl(\operatorname{pushPullObj}
      \mathcal{F}\, Y_n\bigr)
    \;\cong\;
    (\operatorname{pushforward}(\operatorname{Spec}\varphi_n)).\!\operatorname{obj}\bigl(
      \widetilde{N_n}\bigr).
  \]
  The left-hand side is exactly the term \(f_*\bigl((p_n)_* p_n^*\mathcal{F}\bigr)\) whose
  pullback along \(g\) appears in the degree-\(n\) Beck--Chevalley square; the right-hand
  side is the affine pushforward of a \(\widetilde{(-)}\)-module, the form to which
  \ref{mk-lem:affine_pushforward_pullback_baseChange} applies.
\end{lemma}
\begin{proof}
  \uses{mk-def:push_pull_obj, mk-lem:qcoh_iso_tilde_sections, mk-lem:qcoh_isIso_fromTildeGamma,
        mk-lem:pullback_preserves_quasicoherent, mk-lem:pushforward_iso_preserves_qcoh,
        mk-lem:isLocalizing_pushforward_mathlib,
        mk-lem:isAffineOpen_biInf_mathlib, mk-lem:scheme_pushforwardComp_mathlib,
        mk-lem:isoSpec_scheme_mathlib}
  \emph{Altitude 1.} The restriction \(p_n^*\mathcal{F}\) is quasi-coherent on \(Y_n\)
  (\ref{mk-lem:pullback_preserves_quasicoherent}). Pushing it forward along the whole-scheme
  iso \(\sigma_n : Y_n \cong \operatorname{Spec}\Gamma(X, Y_n)\) (\ref{mk-lem:isoSpec_scheme_mathlib})
  lands it on a genuine affine scheme, where \(\widetilde{(-)}\) is defined, and pushforward
  along an isomorphism preserves quasi-coherence (\ref{mk-lem:pushforward_iso_preserves_qcoh}).
  The affine structure theorem 01I8 then applies to the quasi-coherent
  \((\sigma_n)_*\,p_n^*\mathcal{F}\): the tilde--\(\Gamma\) counit is an isomorphism
  (\ref{mk-lem:qcoh_isIso_fromTildeGamma}), so
  \((\sigma_n)_*\,p_n^*\mathcal{F} \cong \widetilde{N_n}\) with
  \(N_n = \Gamma\bigl(\operatorname{Spec}\Gamma(X, Y_n),\ (\sigma_n)_*\,p_n^*\mathcal{F}\bigr)\)
  by \ref{mk-lem:qcoh_iso_tilde_sections}.

  \emph{Altitude 2.} Collapse \(f_* \circ (p_n)_* \cong (f \circ p_n)_*\) by
  \ref{mk-lem:scheme_pushforwardComp_mathlib}; thus
  \((\operatorname{pushforward} f).\operatorname{obj}(\operatorname{pushPullObj}\mathcal{F}\,
  Y_n) = (f \circ p_n)_*\,p_n^*\mathcal{F}\). As \(Y_n\) is affine, the composite
  \(f \circ p_n : Y_n \to S\) of schemes with affine source and target is presented as
  \(\operatorname{Spec}\varphi_n\) for \(\varphi_n = \Gamma(f \circ p_n) : R \to A_n\),
  using the whole-scheme identification \(Y_n \cong \operatorname{Spec}\Gamma(Y_n)\)
  (\ref{mk-lem:isoSpec_scheme_mathlib}). Applying the
  pushforward of this presented morphism to the altitude-1 iso gives
  \((f \circ p_n)_*\,p_n^*\mathcal{F} \cong (\operatorname{Spec}\varphi_n)_*\widetilde{N_n}\),
  which is the asserted altitude-2 form; that the affine pushforward of a
  \(\widetilde{(-)}\)-module is again of \(\widetilde{(-)}\)-form
  is \ref{mk-lem:isLocalizing_pushforward_mathlib}.
\end{proof}

\begin{lemma}[Abstract push--pull tilde-model over an abstract affine base]\label{mk-lem:pushPullObj_iso_tilde_affine}
  \lean{AlgebraicGeometry.pushPullObj_pushforward_iso_tilde_affine}
  \leanok
  \dcref{custom}

  \uses{mk-lem:pushPullObj_iso_tilde, mk-lem:isoSpec_scheme_mathlib,
        mk-lem:scheme_pushforwardComp_mathlib, mk-def:push_pull_obj}
  Let \(f : X \to S\) be \emph{separated} with \(S\) an \emph{abstract} affine scheme
  (\([\,\mathrm{IsAffine}\ S\,]\), so \(S\) need not be a literal \(\operatorname{Spec}\)), and
  write \(e_S : S \xrightarrow{\ \sim\ } \operatorname{Spec}\Gamma(S, \mathcal{O}_S) =
  \operatorname{Spec} R\) for the canonical affine identification \(\operatorname{Scheme.isoSpec}\)
  (\ref{mk-lem:isoSpec_scheme_mathlib}), \(R = \Gamma(S, \mathcal{O}_S)\). Fix an affine
  open-immersion cover \(\mathfrak{U} = \{U_i\}\) of \(X\) and a cosimplicial degree, writing
  \(Y_n\) for the corresponding fibre power and \(N_n\), \(\varphi_n : R \to A_n\) for the data of
  \ref{mk-lem:pushPullObj_iso_tilde}. Then the altitude-2 identification holds over the abstract
  base \(S\), with the affine \(\widetilde{(-)}\)-form transported back along \(e_S\) so that it
  lands in \(\mathcal{O}_S\)-modules rather than \(\mathcal{O}_{\operatorname{Spec} R}\)-modules:
  \[
    (\operatorname{pushforward} f).\!\operatorname{obj}\bigl(\operatorname{pushPullObj}
      \mathcal{F}\, Y_n\bigr)
    \;\cong\;
    (\operatorname{pushforward} e_S^{-1}).\!\operatorname{obj}\Bigl(
      (\operatorname{pushforward}(\operatorname{Spec}\varphi_n)).\!\operatorname{obj}\bigl(
        \widetilde{N_n}\bigr)\Bigr).
  \]
\end{lemma}
\begin{proof}
  \uses{mk-lem:pushPullObj_iso_tilde, mk-lem:isoSpec_scheme_mathlib,
        mk-lem:scheme_pushforwardComp_mathlib, mk-def:push_pull_obj}
  Let \(e_S : S \cong \operatorname{Spec}\Gamma(S, \mathcal{O}_S)\) be the affine identification
  (\ref{mk-lem:isoSpec_scheme_mathlib}). The composite
  \(e_S \circ f : X \to \operatorname{Spec}\Gamma(S, \mathcal{O}_S)\) is separated with a
  \emph{literal} affine base \(\operatorname{Spec} R\), \(R = \Gamma(S, \mathcal{O}_S)\), so the
  literal-base bridge \ref{mk-lem:pushPullObj_iso_tilde} at \emph{altitude 2} applies to it,
  presenting \((e_S \circ f) \circ p_n\) as \(\operatorname{Spec}\varphi_n\) and giving
  \[
    (\operatorname{pushforward}(e_S \circ f)).\!\operatorname{obj}\bigl(\operatorname{pushPullObj}
      \mathcal{F}\, Y_n\bigr)
    \;\cong\;
    (\operatorname{pushforward}(\operatorname{Spec}\varphi_n)).\!\operatorname{obj}\bigl(
      \widetilde{N_n}\bigr).
  \]
  Collapsing the pushforward of the composite,
  \(\operatorname{pushforward}(e_S \circ f) \cong (\operatorname{pushforward} e_S) \circ
  (\operatorname{pushforward} f)\) (\ref{mk-lem:scheme_pushforwardComp_mathlib}). Applying
  \(\operatorname{pushforward} e_S^{-1}\) to both sides and cancelling
  \(\operatorname{pushforward} e_S^{-1} \circ \operatorname{pushforward} e_S \cong
  \operatorname{pushforward}(e_S \circ e_S^{-1}) = \operatorname{pushforward}\operatorname{id} =
  \operatorname{id}\) on the left (again \ref{mk-lem:scheme_pushforwardComp_mathlib}) yields the
  asserted isomorphism, with the right-hand side transported by
  \(\operatorname{pushforward} e_S^{-1}\). The single affine identification
  \(e_S = \operatorname{Scheme.isoSpec}\) is reused on both the \(f\)-side and the
  presented-morphism side.
\end{proof}

\begin{lemma}[Spec-cartesian commutative squares are pushouts of rings]\label{mk-lem:commRingCat_isPushout_iff_mathlib}
  \lean{AlgebraicGeometry.isPushout_of_isPullback_SpecMap}
  \leanok
  \dcref{custom}

  A commutative square of commutative rings \(\varphi : A \to B\), \(\psi : A \to C\),
  \(\rho : B \to P\), \(\sigma : C \to P\) whose \(\operatorname{Spec}\)-square
  \[
    \begin{array}{ccc}
      \operatorname{Spec} P & \xrightarrow{\ \operatorname{Spec}\rho\ } & \operatorname{Spec} B\\
      {\scriptstyle \operatorname{Spec}\sigma}\big\downarrow & &
        \big\downarrow{\scriptstyle \operatorname{Spec}\varphi}\\
      \operatorname{Spec} C & \xrightarrow{\ \operatorname{Spec}\psi\ } & \operatorname{Spec} A
    \end{array}
  \]
  is cartesian in the category of schemes is a pushout (cocartesian) square in the category of
  commutative rings; equivalently, the corner ring is the tensor product \(P = B \otimes_A C\)
  of the two legs over the base \(A\). (This is the converse used by the ring-pushout carve;
  the forward direction is Mathlib's
  \(\operatorname{isPullback\_SpecMap\_of\_isPushout}\).)
\end{lemma}
\begin{proof}
  \leanok
  The functor \(\operatorname{Spec} : \mathrm{CommRing}^{\mathrm{op}} \to \mathrm{Sch}\) is
  fully faithful, and a fully faithful functor reflects limits: the cartesian scheme square
  of the four \(\operatorname{Spec}\)-maps reflects to a pullback square in
  \(\mathrm{CommRing}^{\mathrm{op}}\), i.e.\ the original square is a pushout in
  \(\mathrm{CommRing}\).
\end{proof}

\begin{lemma}[Restriction of the cartesian base-change square over an intersection open is cartesian]\label{mk-lem:restricted_cartesian_affine_pushout}
  \lean{AlgebraicGeometry.restrictedCartesianAffinePushout}
  \leanok
  \dcref{custom}

  Let \(g' : X' \to X\) be an arbitrary morphism of schemes and let
  \(V = \operatorname{coverInterOpen}\mathfrak{U}\,\sigma = \bigcap_k U_{\sigma(k)}\) be a
  {\v C}ech intersection open of an affine open-immersion cover \(\mathfrak{U}\) of \(X\),
  with open immersion \(j_\sigma : V \hookrightarrow X\). Form the fibre product
  \(X' \times_X V\) and the cartesian square
  \[
    \begin{array}{ccc}
      X' \times_X V & \longrightarrow & V\\
      \big\downarrow & & \big\downarrow{\scriptstyle j_\sigma}\\
      X' & \xrightarrow{\ g'\ } & X
    \end{array}
  \]
  in which the top map is the projection \(\operatorname{pullback.snd}\,g'\,j_\sigma\), the
  left map is \(\operatorname{pullback.fst}\,g'\,j_\sigma\), the right map is the open
  immersion \(j_\sigma\), and the bottom map is \(g'\). This square is cartesian: it is a
  pullback square of schemes.
\end{lemma}
\begin{proof}
  \leanok
  The fibre product \(X' \times_X V\) exists and its defining projections fit, by
  construction, into a cartesian square. Reading off that pullback square with the projections
  in the stated positions gives exactly the asserted cartesian square; no further hypotheses
  are required.
\end{proof}

\begin{lemma}[LHS abstract \(\to\) tilde for a single intersection open]\label{mk-lem:coverinter_lhs_iso_tilde}
  \lean{AlgebraicGeometry.pushPullObj_coverInter_pushforward_iso_tilde}
  \leanok
  \dcref{custom}

  \uses{mk-lem:pushPullObj_iso_tilde, mk-lem:cech_interopen_isaffine, mk-def:push_pull_obj}
  Let \(f : X \to S\) be \emph{separated} with affine base \(S = \operatorname{Spec} R\), fix an
  affine open-immersion cover \(\mathfrak{U} = \{U_i = \operatorname{Spec} A_i\}\) of \(X\) and a
  \emph{finite, nonempty} multi-index \(\sigma : \kappa \to I\), with intersection open
  \(V = \operatorname{coverInterOpen}\mathfrak{U}\,\sigma = \bigcap_k U_{\sigma(k)}\) and open
  immersion \(j_\sigma : V \hookrightarrow X\). By \ref{mk-lem:cech_interopen_isaffine} the open
  \(V = \operatorname{Spec} A_\sigma\) is affine; write \(\varphi : R \to A_\sigma\) for the ring
  map presenting the composite \(f \circ j_\sigma : V \to S\) as \(\operatorname{Spec}\varphi\),
  and \(N = \Gamma\bigl(V, (j_\sigma)^*\mathcal{F}\bigr)\) for the \(A_\sigma\)-module of sections
  of the (quasi-coherent) restriction. Then there is an isomorphism of
  \(\mathcal{O}_S\)-modules
  \[
    f_*\bigl(\operatorname{pushPullObj}\mathcal{F}\,(\operatorname{Over.mk} j_\sigma)\bigr)
      \;=\; f_*\bigl((j_\sigma)_* (j_\sigma)^*\mathcal{F}\bigr)
      \;\cong\; (\operatorname{Spec}\varphi)_*\,\widetilde{N}.
  \]
\end{lemma}
\begin{proof}
  \uses{mk-lem:pushPullObj_iso_tilde, mk-lem:cech_interopen_isaffine, mk-def:push_pull_obj}
  The intersection open \(V\) is affine, \(V = \operatorname{Spec} A_\sigma\), by
  \ref{mk-lem:cech_interopen_isaffine} (the separated-diagonal instance). The single open
  immersion \(j_\sigma : V \hookrightarrow X\) is the projection of the one-fold fibre power of
  the cover indexed by \(\sigma\), so \(\operatorname{Over.mk} j_\sigma\) is the degenerate fibre
  power \(Y\) to which \ref{mk-lem:pushPullObj_iso_tilde} applies. Reading that bridge at
  \emph{altitude 2} over the affine base \(S\) collapses
  \(f_* \circ (j_\sigma)_* \cong (f \circ j_\sigma)_*\) and presents the composite as
  \(\operatorname{Spec}\varphi\) with \(\varphi = \Gamma(f \circ j_\sigma) : R \to A_\sigma\),
  yielding exactly
  \(f_*\bigl((j_\sigma)_*(j_\sigma)^*\mathcal{F}\bigr) \cong
  (\operatorname{Spec}\varphi)_*\widetilde{N}\) with
  \(N = \Gamma(V, (j_\sigma)^*\mathcal{F})\).
\end{proof}

\begin{lemma}[Base-change cancellation: the restricted corner module is the tensor base change of \(N\)]\label{mk-lem:coverinter_baseChanged_module_iso_tensor}
  \lean{AlgebraicGeometry.coverInter_baseChanged_sections_iso_tensor}
  \leanok
  \dcref{custom}

  \uses{mk-lem:baseChangeCancelModuleIso, mk-lem:cancelBaseChange_mathlib,
        mk-lem:commRingCat_isPushout_iff_mathlib}
  Let \(\varphi : R \to A_\sigma\), \(\rho : A_\sigma \to B\), \(\psi : R \to R'\),
  \(\sigma' : R' \to B\) be commutative-ring homomorphisms forming a pushout (cocartesian)
  square in \(\mathrm{CommRing}\), so that the corner ring is the tensor product
  \(B = A_\sigma \otimes_R R'\) of the two legs over the base \(R\)
  (\ref{mk-lem:commRingCat_isPushout_iff_mathlib}). Let \(N\) be an arbitrary
  \(A_\sigma\)-module. Then there is a canonical isomorphism of \(R'\)-modules
  \[
    \operatorname{restr}_{\sigma'}\bigl(\operatorname{extendScalars}_\rho N\bigr)
      \;\cong\;
    \operatorname{extendScalars}_\psi\bigl(\operatorname{restr}_\varphi N\bigr),
    \qquad\text{i.e.}\qquad
    \operatorname{restr}_{\sigma'}\bigl(B \otimes_{A_\sigma} N\bigr)
      \;\cong\; R' \otimes_R N.
  \]
  This is the pure commutative-algebra cancellation that underlies the geometric corner
  identification of \ref{mk-lem:coverinter_rhs_iso_tilde}; it carries no sheaf, pullback, or
  \(\Gamma\)-content.
\end{lemma}
\begin{proof}
  \uses{mk-lem:baseChangeCancelModuleIso, mk-lem:cancelBaseChange_mathlib,
        mk-lem:commRingCat_isPushout_iff_mathlib}
  This is exactly the module-level base-change cancellation \ref{mk-lem:baseChangeCancelModuleIso}
  read for the pushout square \((\varphi, \rho, \psi, \sigma')\) and the module \(N\): its
  inverse realises the Mathlib cancellation \(\operatorname{cancelBaseChange}\)
  (\ref{mk-lem:cancelBaseChange_mathlib}),
  \((R' \otimes_R A_\sigma) \otimes_{A_\sigma} N \cong R' \otimes_R N\), which after the
  identification \(B \cong A_\sigma \otimes_R R'\) of the corner ring
  (\ref{mk-lem:commRingCat_isPushout_iff_mathlib}) reads as the asserted isomorphism
  \(\operatorname{restr}_{\sigma'}(B \otimes_{A_\sigma} N) \cong R' \otimes_R N\). No flatness
  hypothesis is required.
\end{proof}

\begin{lemma}[RHS abstract \(\to\) tilde for a single intersection open (base-changed side)]\label{mk-lem:coverinter_rhs_iso_tilde}
  \lean{AlgebraicGeometry.pushPullObj_coverInter_baseChanged_pushforward_iso_tilde}
  \leanok
  \dcref{custom}

  \uses{mk-lem:twisted_cech_nerve_per_sigma, mk-lem:coverinter_lhs_iso_tilde,
        mk-lem:pullback_isQuasicoherent_hom, mk-lem:cech_interopen_isaffine,
        mk-def:push_pull_obj}
  In the cartesian base-change setting of \ref{mk-lem:pushpullobj_coverinter_basechange}, let
  \(f : X \to \operatorname{Spec} R\) be \emph{separated} over the affine base
  \(\operatorname{Spec} R\), let \(g : \operatorname{Spec} R' \to \operatorname{Spec} R\) be the
  base-change morphism, and let \(f' : X' \to \operatorname{Spec} R'\) be the \emph{separated}
  (\([\,\mathrm{IsSeparated}\ f'\,]\)) pulled-back morphism with \(g' : X' \to X\) the top map of
  the cartesian square. Assume the pulled-back affine open-immersion cover
  \(\mathfrak{U}' = \{\mathfrak{U}'.X\, i\}\) of \(X'\) consists of affine opens
  (\([\,\forall i,\,\mathrm{IsAffine}(\mathfrak{U}'.X\,i)\,]\)). Fix the affine intersection
  open \(V = \operatorname{Spec} A_\sigma\) (\ref{mk-lem:cech_interopen_isaffine}) with
  \(j_\sigma : V \hookrightarrow X\) and \(N = \Gamma(V, (j_\sigma)^*\mathcal{F})\); let
  \(j'_\sigma : V' \hookrightarrow X'\) be the pullback of \(j_\sigma\) along \(g'\), so that
  \(V' = X' \times_X V = \operatorname{Spec}(A_\sigma \otimes_R R')\), and set
  \[
    N' \;=\; \Gamma\bigl(V', (j'_\sigma)^*((g')^*\mathcal{F})\bigr)
  \]
  for the raw base-changed corner module. Write \(\psi : R' \to A_\sigma \otimes_R R'\) for the
  ring map presenting the composite \(f' \circ j'_\sigma : V' \to \operatorname{Spec} R'\) as
  \(\operatorname{Spec}\psi\). Then there is an isomorphism of
  \(\mathcal{O}_{\operatorname{Spec} R'}\)-modules
  \[
    f'_*\Bigl((g')^*\bigl(\operatorname{pushPullObj}\mathcal{F}\,
      (\operatorname{Over.mk} j_\sigma)\bigr)\Bigr)
      \;\cong\; (\operatorname{Spec}\psi)_*\,\widetilde{N'}.
  \]
  (The further identification \(N' \cong N \otimes_R R'\) is
  \ref{mk-lem:coverinter_rhs_tensor_rewrite}.)
\end{lemma}
\begin{proof}
  \leanok

\uses{mk-lem:twisted_cech_nerve_per_sigma, mk-lem:coverinter_lhs_iso_tilde,
        mk-lem:pullback_isQuasicoherent_hom, mk-lem:cech_interopen_isaffine,
        mk-def:push_pull_obj}
  First reorganise the pulled-back data by the per-intersection-open \(X\)-level
  Beck--Chevalley isomorphism \ref{mk-lem:twisted_cech_nerve_per_sigma}: since \(j_\sigma\) is an
  open immersion whose pullback along \(g'\) is the open immersion
  \(j'_\sigma : V' \hookrightarrow X'\) onto the preimage \((g')^{-1}(V)\),
  \[
    (g')^*\bigl(\operatorname{pushPullObj}\mathcal{F}\,(\operatorname{Over.mk} j_\sigma)\bigr)
      = (g')^*\bigl((j_\sigma)_*(j_\sigma)^*\mathcal{F}\bigr)
      \;\cong\; (j'_\sigma)_*(j'_\sigma)^*\bigl((g')^*\mathcal{F}\bigr)
      = \operatorname{pushPullObj}\bigl((g')^*\mathcal{F}\bigr)\,(\operatorname{Over.mk}
        j'_\sigma).
  \]
  Pushing forward along \(f'\) reduces the left side of the claim to
  \(f'_*\bigl(\operatorname{pushPullObj}((g')^*\mathcal{F})\,(\operatorname{Over.mk}
  j'_\sigma)\bigr)\). The base-changed module \((g')^*\mathcal{F}\) is quasi-coherent by
  \ref{mk-lem:pullback_isQuasicoherent_hom}, and the intersection open
  \(V' = \operatorname{coverInterOpen}\mathfrak{U}'\,\sigma\) of the base-changed cover is
  affine for the separated \(f' : X' \to \operatorname{Spec} R'\) over the affine base
  \(\operatorname{Spec} R'\) (\ref{mk-lem:cech_interopen_isaffine}). Hence the LHS-side bridge
  \ref{mk-lem:coverinter_lhs_iso_tilde}, applied to the base-changed data
  \((f', \mathfrak{U}', (g')^*\mathcal{F})\) at the index tuple \(\sigma\), gives
  \[
    f'_*\bigl(\operatorname{pushPullObj}((g')^*\mathcal{F})\,(\operatorname{Over.mk}
      j'_\sigma)\bigr)
      \;\cong\; (\operatorname{Spec}\psi)_*\,\widetilde{N'},
    \qquad N' = \Gamma\bigl(V', (j'_\sigma)^*((g')^*\mathcal{F})\bigr),
  \]
  with \(\psi = \Gamma(f' \circ j'_\sigma) : R' \to A_\sigma \otimes_R R'\). This is the asserted
  isomorphism.
\end{proof}

\begin{lemma}[Corner identification of the base-changed section module]\label{mk-lem:coverinter_rhs_tensor_rewrite}
  \lean{AlgebraicGeometry.coverInter_baseChanged_sections_tensor_rewrite}
  \leanok
  \dcref{custom}

  \uses{mk-lem:restricted_cartesian_affine_pushout, mk-lem:cech_interopen_isaffine,
        mk-lem:coverinter_baseChanged_module_iso_tensor, mk-lem:pullback_spec_tilde_iso,
        mk-lem:openimm_beckchevalley}
  In the setting of \ref{mk-lem:coverinter_rhs_iso_tilde}, the raw base-changed corner module
  \(N' = \Gamma\bigl(V', (j'_\sigma)^*((g')^*\mathcal{F})\bigr)\) satisfies
  \[
    N' \;\cong\; N \otimes_R R'
  \]
  naturally in \(\mathcal{F}\), so the right-hand side of
  \ref{mk-lem:coverinter_rhs_iso_tilde} may equivalently be written
  \((\operatorname{Spec}\psi)_*\,\widetilde{N \otimes_R R'}\).
\end{lemma}
\begin{proof}
  \leanok
  \uses{mk-lem:restricted_cartesian_affine_pushout, mk-lem:cech_interopen_isaffine,
        mk-lem:coverinter_baseChanged_module_iso_tensor, mk-lem:pullback_spec_tilde_iso,
        mk-lem:openimm_beckchevalley}
  The restricted square
  \(V' \to V\), \(V' \to X'\), \(g' : X' \to X\), \(j_\sigma\) is cartesian by
  \ref{mk-lem:restricted_cartesian_affine_pushout}, and supplies the canonical pullback
  identification \((j'_\sigma)^*\bigl((g')^*\mathcal{F}\bigr) \cong
  (g'_V)^*\bigl((j_\sigma)^*\mathcal{F}\bigr)\) (the same comparison of the cartesian square used
  in \ref{mk-lem:openimm_beckchevalley}), where \(g'_V : V' \to V\) is the top map of that square.
  Over the affine \(V = \operatorname{Spec} A_\sigma\) the restriction \((j_\sigma)^*\mathcal{F}\)
  is quasi-coherent, \(\cong \widetilde{N}\). The top map \(g'_V\) is the base change of the
  base-change morphism \(g\) over the affine \(V\), hence the affine morphism
  \(\operatorname{Spec}\rho\) for the corner ring map \(\rho : A_\sigma \to A_\sigma \otimes_R R'\);
  pulling \(\widetilde{N}\) back along it is, by the pullback tilde dictionary
  \ref{mk-lem:pullback_spec_tilde_iso},
  \((g'_V)^*\widetilde{N} \cong \widetilde{N \otimes_{A_\sigma}(A_\sigma \otimes_R R')}
  = \widetilde{N \otimes_R R'}\). Taking global sections over
  \(V' = \operatorname{Spec}(A_\sigma \otimes_R R')\) gives
  \(N' \cong N \otimes_{A_\sigma}(A_\sigma \otimes_R R') = N \otimes_R R'\), naturally in
  \(\mathcal{F}\); the underlying \(R'\)-module cancellation
  \(\operatorname{restr}_{\sigma'}(B \otimes_{A_\sigma} N) \cong R' \otimes_R N\) (with
  \(B = A_\sigma \otimes_R R'\)) is the pure algebraic core
  \ref{mk-lem:coverinter_baseChanged_module_iso_tensor}. A consumer rewrites the right-hand side
  \((\operatorname{Spec}\psi)_*\widetilde{N'}\) as
  \((\operatorname{Spec}\psi)_*\widetilde{N \otimes_R R'}\) along this identification.
\end{proof}

\begin{lemma}[The restricted square carves into a ring pushout]\label{mk-lem:coverinter_ring_pushout}
  \lean{AlgebraicGeometry.coverInter_ring_isPushout}
  \leanok
  \dcref{custom}

  \uses{mk-lem:cech_interopen_isaffine, mk-lem:restricted_cartesian_affine_pushout,
        mk-lem:commRingCat_isPushout_iff_mathlib}
  In the setting of \ref{mk-lem:pushpullobj_coverinter_basechange}, write
  \(V = \operatorname{Spec} A_\sigma\) for the affine intersection open
  (\ref{mk-lem:cech_interopen_isaffine}), \(\varphi : R \to A_\sigma\) for the ring map
  presenting \(f \circ j_\sigma\), \(\gamma : R \to R'\) for the ring map presenting \(g\),
  \(\rho : A_\sigma \to B\) for the corner ring map obtained by applying global sections to
  the top map \(V' \to V\) of the restricted square, and \(\psi : R' \to B\) for the ring map
  presenting \(f' \circ j'_\sigma\), where \(B = \Gamma(X', V'_\sigma)\). Then the square
  \[
    \begin{array}{ccc}
      R & \xrightarrow{\ \varphi\ } & A_\sigma\\
      {\scriptstyle \gamma}\big\downarrow & & \big\downarrow{\scriptstyle \rho}\\
      R' & \xrightarrow{\ \psi\ } & B
    \end{array}
  \]
  is a pushout of commutative rings; in particular \(B \cong A_\sigma \otimes_R R'\).
\end{lemma}
\begin{proof}
  \leanok
  \uses{mk-lem:cech_interopen_isaffine, mk-lem:restricted_cartesian_affine_pushout,
        mk-lem:commRingCat_isPushout_iff_mathlib}
  Pasting the restricted cartesian square of
  \ref{mk-lem:restricted_cartesian_affine_pushout} vertically onto the ambient cartesian
  base-change square exhibits \(X' \times_X V\) as the fibre product of
  \(V \xrightarrow{j_\sigma} X \xrightarrow{f} S\) and \(g : S' \to S\). The categorical
  pullback \(X' \times_X V\) is identified with the base-changed intersection open
  \(V' = V'_\sigma\): the first projection is an open immersion with range
  \((g')^{-1}(V)\), which equals \(V'\) because the base-changed cover's member opens are
  the \(g'\)-preimages of the original member opens. Under this identification and the
  affine identifications \(V = \operatorname{Spec} A_\sigma\),
  \(V' = \operatorname{Spec} B\), all four corners of the resulting cartesian square are
  affine and all four maps are \(\operatorname{Spec}\)-maps of the stated ring maps; by
  \ref{mk-lem:commRingCat_isPushout_iff_mathlib} the ring square is a pushout, and the corner
  of a ring pushout is the tensor product \(A_\sigma \otimes_R R'\).
\end{proof}

\begin{lemma}[Per-intersection-open base change of the push--pull object]\label{mk-lem:pushpullobj_coverinter_basechange}
  \lean{AlgebraicGeometry.pushPullObj_coverInter_baseChange}
  \leanok
  \dcref{custom}

  \uses{mk-lem:coverinter_lhs_iso_tilde, mk-lem:coverinter_rhs_iso_tilde,
        mk-lem:cech_degree_affine_baseChange,
        mk-lem:affine_pushforward_pullback_baseChange, mk-lem:isAffineOpen_biInf_mathlib,
        mk-lem:cech_interopen_isaffine, mk-def:push_pull_obj}
  Let \(f : X \to S\) be \emph{separated} with affine base \(S = \operatorname{Spec} R\), let
  \(g : S' \to S\) be a base-change morphism with \(S' = \operatorname{Spec} R'\) affine, and
  let
  \[
    \begin{array}{ccc}
      X' & \xrightarrow{\ g'\ } & X\\
      {\scriptstyle f'}\big\downarrow & & \big\downarrow{\scriptstyle f}\\
      S' & \xrightarrow{\ g\ } & S
    \end{array}
  \]
  be the associated cartesian square. Fix an affine open-immersion cover
  \(\mathfrak{U} = \{U_i = \operatorname{Spec} A_i\}\) of \(X\) and a multi-index
  \(\sigma : \kappa \to I\), with intersection open
  \(V = \operatorname{coverInterOpen}\mathfrak{U}\,\sigma = \bigcap_k U_{\sigma(k)}\) and open
  immersion \(j_\sigma : V \hookrightarrow X\). Then there is an isomorphism of
  \(\mathcal{O}_{S'}\)-modules
  \[
    g^*\Bigl(f_*\bigl((j_\sigma)_* (j_\sigma)^*\mathcal{F}\bigr)\Bigr)
      \;\cong\;
    f'_*\Bigl((g')^*\bigl((j_\sigma)_* (j_\sigma)^*\mathcal{F}\bigr)\Bigr),
  \]
  i.e.\ \(g^*\bigl(f_*(\operatorname{pushPullObj}\mathcal{F}\,(\operatorname{Over.mk}
  j_\sigma))\bigr) \cong f'_*\bigl((g')^*(\operatorname{pushPullObj}\mathcal{F}\,
  (\operatorname{Over.mk} j_\sigma))\bigr)\). This is the per-intersection-open instance of the
  Beck--Chevalley comparison; it is the genuine open content that survives the
  coproduct/product decomposition of the degreewise {\v C}ech nerve.
\end{lemma}
\begin{proof}
  \leanok
  \uses{mk-lem:coverinter_lhs_iso_tilde, mk-lem:coverinter_rhs_iso_tilde,
        mk-lem:coverinter_rhs_tensor_rewrite,
        mk-lem:cech_degree_affine_baseChange, mk-lem:affine_pushforward_pullback_baseChange,
        mk-lem:isAffineOpen_biInf_mathlib, mk-lem:cech_interopen_isaffine,
        mk-lem:restricted_cartesian_affine_pushout, mk-lem:commRingCat_isPushout_iff_mathlib,
        mk-lem:coverinter_ring_pushout}
  The two sides of the comparison are routed independently to the common affine
  \((\operatorname{Spec}\cdot)_*\widetilde{(-)}\) form, where the affine termwise base change
  closes the gap between them.
  \begin{enumerate}
    \item \emph{\(V\) is affine.} The intersection open
      \(V = \bigcap_k U_{\sigma(k)}\) is a finite nonempty intersection of affine opens of the
      separated scheme \(X\), hence affine by \ref{mk-lem:cech_interopen_isaffine} (the
      separated-diagonal instance of \ref{mk-lem:isAffineOpen_biInf_mathlib}),
      \(V = \operatorname{Spec} A_\sigma\); write \(\varphi : R \to A_\sigma\) for the ring map
      presenting \(f \circ j_\sigma\) as \(\operatorname{Spec}\varphi\) and
      \(N = \Gamma(V, (j_\sigma)^*\mathcal{F})\).
    \item \emph{LHS \(\to\) tilde.} By \ref{mk-lem:coverinter_lhs_iso_tilde} the left abstract
      term is the affine pushforward of \(\widetilde{N}\) over \(S\),
      \[
        f_*\bigl(\operatorname{pushPullObj}\mathcal{F}\,(\operatorname{Over.mk} j_\sigma)\bigr)
          \;\cong\; (\operatorname{Spec}\varphi)_*\,\widetilde{N},
      \]
      so applying \(g^*\) gives
      \(g^*\bigl(f_*(\operatorname{pushPullObj}\mathcal{F}\,(\operatorname{Over.mk}
      j_\sigma))\bigr) \cong g^*\bigl((\operatorname{Spec}\varphi)_*\widetilde{N}\bigr)\).
    \item \emph{RHS \(\to\) tilde.} By \ref{mk-lem:coverinter_rhs_iso_tilde} the right abstract
      term is likewise the affine pushforward of the base-changed module
      \(N \otimes_R R'\) over \(S'\),
      \[
        f'_*\Bigl((g')^*\bigl(\operatorname{pushPullObj}\mathcal{F}\,
          (\operatorname{Over.mk} j_\sigma)\bigr)\Bigr)
          \;\cong\; (\operatorname{Spec}\psi)_*\,\widetilde{N \otimes_R R'},
      \]
      where \(\psi : R' \to A_\sigma \otimes_R R'\) presents
      \(f' \circ j'_\sigma\). Both comparison sides are now in
      \((\operatorname{Spec}\cdot)_*\widetilde{(-)}\) form.
    \item \emph{The affine pushout square closes the gap.} Restrict the cartesian base-change
      square over \(j_\sigma : V \hookrightarrow X\). By
      \ref{mk-lem:restricted_cartesian_affine_pushout} the resulting square
      \(X' \times_X V \to V\), \(X' \times_X V \to X'\), \(g'\), \(j_\sigma\) is cartesian, and
      with \(V = \operatorname{Spec} A_\sigma\) (step~1), \(S = \operatorname{Spec} R\),
      \(S' = \operatorname{Spec} R'\) affine it is a cartesian square of affine schemes,
      \(\operatorname{Spec}(A_\sigma \otimes_R R') \to \operatorname{Spec} A_\sigma\),
      \(\operatorname{Spec}(A_\sigma \otimes_R R') \to \operatorname{Spec} R'\). Applying global
      sections, the affine-pullback \(\leftrightarrow\) ring-pushout equivalence
      \ref{mk-lem:commRingCat_isPushout_iff_mathlib} converts it into the cocartesian square of
      commutative rings
      \[
        \begin{array}{ccc}
          R & \xrightarrow{\ \varphi\ } & A_\sigma\\
          \big\downarrow & & \big\downarrow\\
          R' & \xrightarrow{\ \psi\ } & A_\sigma \otimes_R R'
        \end{array}
      \]
      with corner \(A_\sigma \otimes_R R'\). For this pushout square the affine termwise base
      change \ref{mk-lem:cech_degree_affine_baseChange}, the per-open instance of
      \ref{mk-lem:affine_pushforward_pullback_baseChange}, supplies
      \[
        g^*\bigl((\operatorname{Spec}\varphi)_*\widetilde{N}\bigr)
          \;\cong\; (\operatorname{Spec}\psi)_*\,\widetilde{N \otimes_R R'}.
      \]
      Composing step~2 (and \(g^*\)), this affine isomorphism, and the inverse of step~3 yields
      the asserted comparison
      \(g^*\bigl(f_*(\operatorname{pushPullObj}\mathcal{F}\,(\operatorname{Over.mk}
      j_\sigma))\bigr) \cong f'_*\bigl((g')^*(\operatorname{pushPullObj}\mathcal{F}\,
      (\operatorname{Over.mk} j_\sigma))\bigr)\).
  \end{enumerate}
\end{proof}

\begin{lemma}[Affine pushforward preserves quasi-coherence]\label{mk-lem:isQuasicoherent_pushforward_of_isAffine}
  \lean{AlgebraicGeometry.isQuasicoherent_pushforward_of_isAffine}
  \leanok
  \dcref{custom}

  Let \(q : A \to B\) be a morphism of affine schemes and \(\mathcal{M}\) a quasi-coherent
  \(\mathcal{O}_A\)-module. Then \(q_*\mathcal{M}\) is quasi-coherent.
\end{lemma}

\begin{proof}
  \leanokFor a literal \(q = \operatorname{Spec}\varphi\) this is the statement that an affine
  pushforward preserves the \(\widetilde{(-)}\) model: \(\mathcal{M} \cong \widetilde{N}\)
  for an \(A\)-module \(N\), and \((\operatorname{Spec}\varphi)_*\widetilde{N}\) is again a
  tilde, hence quasi-coherent. A general \(q\) between affine schemes is conjugate to
  \(\operatorname{Spec}\Gamma(q)\) by the canonical identifications
  \(A \cong \operatorname{Spec}\Gamma(A,\top)\), \(B \cong \operatorname{Spec}\Gamma(B,\top)\),
  and pushforward along an isomorphism preserves quasi-coherence in both directions.
\end{proof}

\begin{lemma}[Every {\v C}ech term is quasi-coherent]\label{mk-lem:isQuasicoherent_cechComplex_X}
  \lean{AlgebraicGeometry.isQuasicoherent_cechComplex_X}
  \leanok
  \dcref{custom}

  \uses{mk-lem:isQuasicoherent_pushforward_of_isAffine, mk-lem:pushPull_sigma_iso,
        mk-lem:isAffineOpen_biInf_mathlib}
  Let \(f : X \to S\) be separated with \(X\) separated and \(S\) affine, let
  \(\mathfrak{U}\) be a finite affine open cover of \(X\), and let \(\mathcal{F}\) be
  quasi-coherent. Then every term \(\check{C}^p(\mathfrak{U}, \mathcal{F})\) of the relative
  {\v C}ech complex is quasi-coherent.
\end{lemma}

\begin{proof}
  \leanok\uses{mk-lem:isQuasicoherent_pushforward_of_isAffine, mk-lem:pushPull_sigma_iso,
        mk-lem:isAffineOpen_biInf_mathlib, mk-lem:isQuasicoherent_pullback_opens}
  The degree-\(p\) term is \(f_*\bigl(\operatorname{pushPullObj}\mathcal{F}\,Y_p\bigr)\) for the
  degree-\(p\) {\v C}ech backbone \(Y_p\). By \ref{mk-lem:pushPull_sigma_iso} this push--pull object
  is the finite product \(\prod_\sigma (j_\sigma)_*(j_\sigma)^*\mathcal{F}\) over index tuples
  \(\sigma\), and \(f_*\) preserves that finite product, so it suffices to treat one factor and
  then reassemble: over an affine base a finite product of quasi-coherent modules is
  quasi-coherent, since each factor is a tilde, \(\widetilde{(-)}\) preserves finite products, and
  quasi-coherence is closed under isomorphism.

  For a single factor, \(f_*\bigl((j_\sigma)_* \mathcal{N}\bigr) \cong (j_\sigma \circ f)_*
  \mathcal{N}\). Each intersection open \(U_\sigma\) is affine
  (\ref{mk-lem:isAffineOpen_biInf_mathlib}, using that \(X\) is separated and the cover is affine)
  and \(S\) is affine, so \(j_\sigma \circ f : U_\sigma \to S\) is a morphism of affine schemes and
  \ref{mk-lem:isQuasicoherent_pushforward_of_isAffine} applies to
  \(\mathcal{N} = (j_\sigma)^*\mathcal{F}\), which is quasi-coherent by
  \ref{mk-lem:isQuasicoherent_pullback_opens}.
\end{proof}

\begin{lemma}[Beck--Chevalley natural iso through the {\v C}ech nerve]\label{mk-lem:cech_pushforward_baseChange_natIso}
  \lean{AlgebraicGeometry.cech_pushforward_baseChange_natIso_flat}
  \leanok
  \dcref{custom}

  \uses{mk-lem:affine_pushforward_pullback_baseChange, mk-lem:pushPullObj_iso_tilde,
        mk-def:cech_nerve_cosimplicial, mk-def:push_pull_functor}
  Let \(f : X \to S\) be \emph{separated} with affine base \(S = \operatorname{Spec} R\),
  fix an affine open-immersion cover \(\mathfrak{U} = \{U_i = \operatorname{Spec} A_i\}\) of
  \(X\) (so that, by \ref{mk-lem:isAffineOpen_biInf_mathlib}, every fibre power \(Y_n\) of the
  {\v C}ech nerve is affine). Whiskered through the {\v C}ech nerve, there is a natural
  isomorphism of cosimplicial \(\mathcal{O}_{S'}\)-modules
  \[
    g^* \circ (\operatorname{pushforward} f) \;\cong\;
      (\operatorname{pushforward} f') \circ (g')^*,
  \]
  the Beck--Chevalley comparison for the cartesian square, valid at every cosimplicial
  degree.
\end{lemma}

\begin{proof}
  \uses{mk-lem:cech_degree_affine_baseChange, mk-lem:affine_pushforward_pullback_baseChange,
        mk-lem:pushPullObj_iso_tilde, mk-lem:isAffineOpen_biInf_mathlib,
        mk-lem:pushforward_spec_tilde_iso,
        mk-lem:pullback_spec_tilde_iso, mk-lem:cancelBaseChange_mathlib,
        mk-lem:pushPull_sigma_iso, mk-lem:pushpullobj_coverinter_basechange,
        mk-def:cech_nerve_cosimplicial, mk-def:push_pull_functor,
        mk-thm:quot_canonical_basechange_isIso, mk-lem:openimm_beckchevalley}
  The natural isomorphism is built degreewise. Fix a cosimplicial degree \(n\) and work over
  the affine base \(S = \operatorname{Spec} R\). The degree-\(n\) comparison is between
  \(g^*\bigl(f_*(\operatorname{pushPullObj}\mathcal{F}\, Y_n)\bigr)\) and
  \(f'_*\bigl((g')^*(\operatorname{pushPullObj}\mathcal{F}\, Y_n)\bigr)\), where \(Y_n\) is the
  degree-\(n\) fibre power of the {\v C}ech nerve.
  \begin{enumerate}
    \item \emph{Product decomposition.} The degree-\(n\) fibre power is the coproduct
      \(Y_n = \coprod_\sigma U_\sigma\) over index tuples \(\sigma : \mathrm{Fin}(n+1) \to I\)
      of the intersection opens \(U_\sigma = \operatorname{coverInterOpen}\mathfrak{U}\,\sigma\),
      so by \ref{mk-lem:pushPull_sigma_iso} (which requires the cover to be finite) the
      push--pull object decomposes as the finite product \(\operatorname{pushPullObj}\mathcal{F}\,
      Y_n \cong \prod_\sigma \operatorname{pushPullObj}\mathcal{F}\,(\operatorname{Over.mk}
      j_\sigma)\). Both \(\operatorname{pushforward} f\) and \(g^*\) (and likewise \(f'_*\),
      \((g')^*\)) preserve this finite product, so the degree-\(n\) comparison reduces termwise
      to the per-\(\sigma\) single-open base change.
    \item \emph{The per-\(\sigma\) factor.} Each summand is exactly the per-intersection-open
      base-change isomorphism \(g^*\bigl(f_*((j_\sigma)_*(j_\sigma)^*\mathcal{F})\bigr) \cong
      f'_*\bigl((g')^*((j_\sigma)_*(j_\sigma)^*\mathcal{F})\bigr)\) of
      \ref{mk-lem:pushpullobj_coverinter_basechange}. There the intersection open
      \(U_\sigma = U_{i_0} \cap \dots \cap U_{i_p}\) is affine by separatedness
      (\ref{mk-lem:isAffineOpen_biInf_mathlib}); the abstract pushed-forward term is identified
      with the affine \(\widetilde{(-)}\)-form \((\operatorname{Spec}\varphi)_*\widetilde{N}\)
      over \(S\) by \ref{mk-lem:pushPullObj_iso_tilde} (\emph{altitude 2}); and on that affine
      model the comparison IS the affine termwise base change
      \ref{mk-lem:affine_pushforward_pullback_baseChange} (via
      \ref{mk-lem:cech_degree_affine_baseChange}), assembled from the concrete tilde
      dictionaries \ref{mk-lem:pushforward_spec_tilde_iso} and
      \ref{mk-lem:pullback_spec_tilde_iso} with the commutative-algebra cancellation
      \(\operatorname{cancelBaseChange}\) (\ref{mk-lem:cancelBaseChange_mathlib}) --- not the
      canonical adjoint mate.
    \item \emph{Naturality.} The naturality condition is the index-omission
      restriction: each coface map is a restriction along an inclusion of finite affine
      intersections, and the termwise isomorphisms are natural in the intersection;
      naturality in the base-change morphism \(g\) is inherited termwise from the affine
      base change.
  \end{enumerate}

  A second, shorter proof avoids the naturality verification entirely, and it is the one the
  the construction follows. The comparison is the \emph{mate} \(\alpha : f_* \circ g^* \to (g')^*
  \circ f'_*\) of the canonical pullback \(2\)-isomorphism attached to the cartesian square
  (\ref{mk-lem:openimm_beckchevalley}, mate form, instantiated at the outer square rather than a
  restricted one). Whiskering a natural transformation with the {\v C}ech nerve produces a map of
  cosimplicial objects whose compatibility with the cofaces is \(\alpha\)'s own naturality, so no
  square has to be checked; and \(\alpha\) is invertible degreewise because, for \(g\) flat and
  \(f\) quasi-compact and quasi-separated, it is invertible at every quasi-coherent module
  (\ref{mk-thm:quot_canonical_basechange_isIso}), while each degree of the nerve is a finite product of
  the quasi-coherent objects \((j_\sigma)_*(j_\sigma)^*\mathcal{F}\) and a natural transformation
  invertible at each factor of a preserved finite product is invertible at the product.
\end{proof}

\begin{lemma}[The base-changed cover intersection is the preimage of the intersection]\label{mk-lem:coverinteropen_basechange_eq}
  \lean{AlgebraicGeometry.coverInterOpen_baseChange_eq}
  \leanok
  \dcref{custom}

  \uses{mk-def:cover_cech_nerve}
  Let \(g' : X' \to X\) be the pullback of a base change \(g : S' \to S\) along a separated
  \(f : X \to S\), and let \(\mathfrak{U}'\) be the base change along \(g'\) of the affine
  open-immersion cover \(\mathfrak{U}\) of \(X\). Then for every finite index family
  \(\sigma\) the {\v C}ech intersection open of the base-changed cover is the preimage of the
  intersection open of the original cover:
  \[
    \operatorname{coverInterOpen}\mathfrak{U}'\,\sigma
      \;=\; (g')^{-1}\bigl(\operatorname{coverInterOpen}\mathfrak{U}\,\sigma\bigr).
  \]
\end{lemma}
\begin{proof}
  \uses{mk-def:cover_cech_nerve}
  The base-changed cover member \(\mathfrak{U}'.X\,i\) is by definition the pullback
  \((g')^{-1}(U_i)\), so its range in \(X'\) is the preimage of the range of \(U_i\). Taking
  preimages of opens commutes with finite meets (the inverse-image map on the frame of opens
  preserves intersections), whence
  \(\operatorname{coverInterOpen}\mathfrak{U}'\,\sigma = \bigcap_k (g')^{-1}(U_{\sigma(k)})
  = (g')^{-1}\bigl(\bigcap_k U_{\sigma(k)}\bigr)
  = (g')^{-1}\bigl(\operatorname{coverInterOpen}\mathfrak{U}\,\sigma\bigr)\).
\end{proof}

\begin{lemma}[Open-immersion Beck--Chevalley over a restricted cartesian square]\label{mk-lem:openimm_beckchevalley}
  \lean{AlgebraicGeometry.openImmersion_beckChevalley}
  \leanok
  \dcref{custom}

  \uses{mk-def:push_pull_obj}
  Let
  \[
    \begin{array}{ccc}
      V' & \xrightarrow{\ g'_V\ } & V\\
      {\scriptstyle p'}\big\downarrow & & \big\downarrow{\scriptstyle p}\\
      X' & \xrightarrow{\ g'\ } & X
    \end{array}
  \]
  be the cartesian square obtained by restricting \(g' : X' \to X\) over an \emph{open
  immersion} \(p : V \hookrightarrow X\) \emph{with \(V\) affine} (in the formal proof \(p\) is taken to be
  the inclusion \(V_0.\iota\) of an affine open \(V_0 : X.\mathrm{Opens}\),
  \(\mathrm{IsAffineOpen}\,V_0\) --- the form in which every consumer instantiates the square),
  so that \(p' : V' \hookrightarrow X'\) is again an open immersion (the pullback of \(p\) along
  \(g'\), with image the preimage \((g')^{-1}(V)\)). Assume in addition that \(X\) is
  \emph{separated} (so that the affine \(V\) meets every affine open of \(X\) in an affine;
  in the formal proof this is a separatedness hypothesis on \(X\) --- e.g.\ \([\mathrm{IsSeparated}\,
  (\mathrm{terminal.from}\,X)]\), or any equivalent input guaranteeing affine
  intersections, at the prover's choice). Then for any \emph{quasi-coherent}
  \(\mathcal{O}_X\)-module \(\mathcal{F}\) (\(\mathtt{hF} : \mathcal{F}.\mathrm{IsQuasicoherent}\))
  there is a Beck--Chevalley isomorphism
  \[
    (g')^*\bigl(p_*\,p^*\mathcal{F}\bigr)
      \;\cong\;
    p'_*\,p'^*\bigl((g')^*\mathcal{F}\bigr),
  \]
  i.e.\ \((g')^*(\operatorname{pushPullObj}\mathcal{F}\,(\operatorname{Over.mk} p)) \cong
  \operatorname{pushPullObj}((g')^*\mathcal{F})\,(\operatorname{Over.mk} p')\). This is the
  open-immersion (in particular flat) case of Beck--Chevalley. The affineness of \(V\)
  together with the quasi-coherence of \(\mathcal{F}\) is what reduces it, after refining
  \(X'\) by an affine open cover, to the affine termwise base-change isomorphism
  \ref{mk-lem:affine_pushforward_pullback_baseChange}. The separatedness hypothesis on \(X\) is
  what makes the affine \(V\) meet each affine open of \(X\) in an affine; in the application
  the consumer discharges it from its own separatedness hypothesis through
  \ref{mk-lem:cech_interopen_isaffine}.
\end{lemma}
\begin{proof}
  \leanok\uses{mk-def:push_pull_obj, mk-lem:openimm_bareBC, mk-lem:openimm_bc_telescope,
        mk-lem:restrictFunctorIsoPullback_mathlib, mk-lem:openimm_bareBC_isIso}
  Throughout write \(G := p^*\mathcal{F}\) for the restricted module on \(V\). A warning on the
  geometry first: \(p\) being an open immersion does \emph{not} make the \emph{right}-adjoint
  direct image \(p_*\) extension-by-zero. Extension-by-zero is the \emph{left} adjoint \(p_!\);
  the right adjoint \(p_*\) has genuinely nonzero sections off \(V\) (over an open \(U \subseteq
  X\) its sections are \(\Gamma(U \cap V, p^*\mathcal{F})\), not \(0\) for \(U \not\subseteq V\)).
  So there is no ``both sides vanish away from \(V'\)'' support argument. The off-\(V'\) data is
  real and is exactly the content of the comparison; the correct tool is a sectionwise criterion.
  The orientation is the open-immersion (flat) special case of base change, Stacks Project
  Tag~02KH (\texttt{lemma-flat-base-change-cohomology}); the section-level input is
  supplied, after refining \(X'\) by an affine open cover and using that \(V\) is affine and
  \(\mathcal{F}\) quasi-coherent, by the affine termwise base-change brick
  \ref{mk-lem:affine_pushforward_pullback_baseChange} (Stacks Project Tag~02KG), which already
  packages the affine tilde dictionaries and the cancellation
  \(\operatorname{cancelBaseChange}\).

  \emph{Stage 1 (telescope; no affineness required).} Since \(p\) and \(p'\) are open immersions,
  pullback along them \emph{is} restriction: the restriction functor and the module pullback
  agree, \(\operatorname{restrict} p \cong p^*\) and \(\operatorname{restrict} p' \cong p'^*\)
  (\ref{mk-lem:restrictFunctorIsoPullback_mathlib}). Using the pseudofunctor structure of pullback
  --- the comparison \((q \circ r)^* \cong r^* \circ q^*\) for composable morphisms, together with
  reindexing of \((-)^*\) along an equality of morphisms --- and the commuting square
  \(g'_V \circ p' = p \circ g'\) (in the sense \(p' \cdot g' = g'_V \cdot p\) as maps
  \(V' \to X\)), one telescopes
  \[
    p'^*\,(g')^*
      \;\cong\; (p' \cdot g')^*
      \;=\; (g'_V \cdot p)^*
      \;\cong\; (g'_V)^*\,p^* .
  \]
  Applied to \(\mathcal{F}\) this identifies \(p'^*(g')^*\mathcal{F} \cong (g'_V)^* G\), hence the
  right-hand side rewrites as
  \(p'_*\,p'^*(g')^*\mathcal{F} \cong p'_*\,(g'_V)^* G = (\operatorname{pushforward} p')
  ((g'_V)^* G)\), while the left-hand side is
  \((g')^*\bigl(p_*\,G\bigr) = (g')^*\bigl((\operatorname{pushforward} p) G\bigr)\). The whole
  telescope is formal and uses no affineness. It reduces the lemma, with the single fixed
  coefficient \(G = p^*\mathcal{F}\), to the assertion that the \emph{bare base-change map}
  \[
    \operatorname{bareBC} \colon\quad
    (g')^*\bigl((\operatorname{pushforward} p)\,G\bigr)
      \;\longrightarrow\;
    (\operatorname{pushforward} p')\bigl((g'_V)^* G\bigr)
  \]
  --- the canonical mate of \(\operatorname{pushforward} p \circ (g')^* \to (g'_V)^* \circ
  \operatorname{pushforward} p'\) attached to the cartesian square --- is an isomorphism. (This
  canonical mate is the local comparison morphism; it is stated from first principles,
  independently of any relative-spectrum base-change theory.)

  \emph{Stage 2 (the bare comparison is an isomorphism).} It remains only to prove that
  \(\operatorname{bareBC}\) is an isomorphism. This is exactly \ref{mk-lem:openimm_bareBC_isIso}.
  By the mate formula, \(\operatorname{bareBC}\) at \(G\) is the \(p'\)-unit at
  \(M = (g')^*(p_*\,G)\) followed by isomorphisms (the telescope components and the \(p\)-counit,
  invertible for the open immersion \(p\) by \ref{mk-lem:openimm_pullback_counit_isIso}); since
  \(p'_*\) is fully faithful, the unit is an isomorphism iff \(M\) lies in the essential image of
  \(p'_*\), and that membership is cover-local (\ref{mk-lem:essimage_pushforward_cover_local}).
  Check it on an affine open cover \(\{W_j\}\) of \(X'\) refining
  \((g')^{-1}(\text{affine opens of } X)\) (\ref{mk-lem:openimm_refining_affine_cover}): over each
  member \(W_j = \operatorname{Spec} B\), with \(g'(W_j) \subseteq U_j = \operatorname{Spec} A\),
  --- because \(V\) is affine and \(X\) is separated --- the intersection \(V \cap U_j\) is
  affine, the restricted local square is a \emph{fully-affine} pullback, and the
  affine termwise base-change isomorphism \ref{mk-lem:affine_pushforward_pullback_baseChange},
  applied with the open immersion \(V \cap U_j \hookrightarrow U_j\) in the role of the affine
  map and presenting module \(M' = \Gamma(V \cap U_j,\ p^*\mathcal{F})\) (quasi-coherent by
  \(\mathtt{hF}\)), exhibits \(M|_{W_j}\) as a pushforward from \(W_j \cap (g')^{-1}(V)\)
  (\ref{mk-lem:openimm_bareBC_app_isIso_affine}). Hence \(\operatorname{bareBC}\) is an
  isomorphism. Combined with the Stage~1 reduction this proves the lemma.

  \emph{Remark (the hypotheses are essential).} The affineness of \(V\), the separatedness of
  \(X\), and the quasi-coherence of \(\mathcal{F}\) are not cosmetic. Separatedness is what makes
  the affine \(V\) meet each affine open \(U_j\) of \(X\) in an \emph{affine} intersection
  \(V \cap U_j\); without it the local square over \(W_j\) need not be fully affine and the
  reduction to the affine brick breaks. The outer base pullback \((g')^*\) is the sheafification
  of the presheaf tensor \(U \mapsto (p_*\,p^*\mathcal{F})(U) \otimes_{\mathcal{O}_X(U)}
  \mathcal{O}_{X'}(W)\); this tensor description agrees with the sections of
  \(p'_*\,(g'_V)^*(p^*\mathcal{F})\) only on the \emph{quasi-coherent} locus, so for a general
  \(\mathcal{O}_X\)-module the comparison fails --- quasi-coherence of \(\mathcal{F}\) is genuinely
  required. Affineness of \(V\) is what makes the local square over each \(W_j\) fully affine and
  hence reducible to the affine brick; in the application it is supplied by the consumer, where
  each fibre power \(U_\sigma\) of the {\v C}ech nerve is affine because the cover is affine and
  the structure morphism is separated (\ref{mk-lem:cech_interopen_isaffine}). In particular the
  bare-mate route never forms the global push--pull \(p_*\,p^*\mathcal{F}\), so no
  quasi-compactness of \(p\) is needed; there is no claim of a hypothesis-free open-immersion
  Beck--Chevalley.
\end{proof}

\begin{lemma}[Affine cover of \(X'\) refining preimages of affine opens of \(X\)]\label{mk-lem:openimm_refining_affine_cover}
  \lean{AlgebraicGeometry.openImmersion_refiningAffineCover}
  \leanok
  \dcref{custom}

  \uses{mk-lem:isoSpec_scheme_mathlib}
  Let \(g' : X' \to X\) be any morphism of schemes. Then there is an affine open cover
  \(\{W_j\}\) of \(X'\) together with, for each \(j\), an affine open \(U_j \subseteq X\) such that
  \(g'(W_j) \subseteq U_j\); equivalently the restricted map \(g'|_{W_j} : W_j \to U_j\) is a map
  of affine schemes, hence a \(\operatorname{Spec}\)-map \(g'|_{W_j} = \operatorname{Spec}\varphi_j\)
  for a ring map \(\varphi_j : \Gamma(U_j, \mathcal{O}_X) \to \Gamma(W_j, \mathcal{O}_{X'})\).
\end{lemma}
\begin{proof}
  \uses{mk-lem:isoSpec_scheme_mathlib}
  Take the standard affine open cover \(\mathfrak{A} = \{A_i = \operatorname{Spec} R_i\}\) of \(X\)
  and pull it back along \(g'\): the base-changed cover \(g'^{-1}\mathfrak{A}\) is an open cover of
  \(X'\) whose member over index \(i\) maps into the affine \(A_i\). Its members \(g'^{-1}(A_i)\)
  need not be affine, so refine each of them by its own affine open cover and combine the results
  into a single cover of \(X'\). The resulting cover \(\{W_j\}\) is affine, and each \(W_j\) sits
  inside some \(g'^{-1}(A_i)\), whence \(g'(W_j) \subseteq A_i =: U_j\). A morphism between two
  affine schemes is induced by a ring map on global sections (\ref{mk-lem:isoSpec_scheme_mathlib}),
  so \(g'|_{W_j} = \operatorname{Spec}\varphi_j\) with \(\varphi_j = \Gamma(g'|_{W_j})\).
\end{proof}

\begin{lemma}[The geometric counit at an open immersion is invertible]\label{mk-lem:openimm_pullback_counit_isIso}
  \lean{AlgebraicGeometry.openImmersion_pullback_counit_isIso}
  \leanok
  \dcref{custom}

  \uses{mk-lem:restrictFunctorIsoPullback_mathlib}
  Let \(q : V \hookrightarrow X\) be an open immersion of schemes.  Then for every
  \(\mathcal{O}_V\)-module \(c\) the counit \(q^*(q_*\,c) \to c\) of the pullback--pushforward
  adjunction \(q^* \dashv q_*\) is an isomorphism.  (Equivalently: \(q_*\) is fully faithful for
  an open immersion \(q\).)
\end{lemma}
\begin{proof}
  \leanok\uses{mk-lem:restrictFunctorIsoPullback_mathlib}
  For an open immersion the module pullback \(q^*\) is naturally isomorphic to the restriction
  functor (\ref{mk-lem:restrictFunctorIsoPullback_mathlib}), and the counit of the
  restriction--pushforward adjunction is invertible: over an open \(O \subseteq V\) it is the
  identification \(\Gamma(O,\ \operatorname{restrict}_q (q_*\,c)) = \Gamma(q^{-1}(q(O)),\ c) =
  \Gamma(O,\ c)\), using that \(q\) is an open embedding so \(q^{-1}(q(O)) = O\).  Counits of
  adjunctions with a common right adjoint correspond under the canonical comparison of left
  adjoints, so the geometric counit is the composite of an isomorphism (the comparison, evaluated
  at \(q_*\,c\)) with an isomorphism (the restriction counit), hence an isomorphism.
\end{proof}

\begin{lemma}[Restriction maps of a pushforward-presented module are isomorphisms]\label{mk-lem:restrictionMap_isIso_of_essImage}
  \lean{AlgebraicGeometry.restrictionMap_isIso_of_essImage}
  \leanok
  \dcref{custom}

  Let \(W\) be a scheme, \(O_0 \subseteq W\) an open with inclusion \(\iota : O_0 \hookrightarrow
  W\), and \(N\) an \(\mathcal{O}_W\)-module lying in the essential image of \(\iota_*\), say
  \(N \cong \iota_*\,K\).  Then for every open \(O \subseteq W\) the presheaf restriction map
  \[
    N(O) \longrightarrow N(O \cap O_0)
  \]
  is an isomorphism.
\end{lemma}
\begin{proof}
  \leanokTransport along the isomorphism \(N \cong \iota_*\,K\) (naturality of its components in the
  open).  On the pushforward side the map in question is the restriction
  \((\iota_*K)(O) = K(\iota^{-1}(O)) \to K(\iota^{-1}(O \cap O_0)) = (\iota_*K)(O \cap O_0)\),
  and \(\iota^{-1}(O \cap O_0) = \iota^{-1}(O) \cap \iota^{-1}(O_0) = \iota^{-1}(O)\) since
  \(\iota^{-1}(O_0)\) is everything; so it is \(K\) applied to an identity inclusion, hence an
  isomorphism.
\end{proof}

\begin{lemma}[Essential-image membership for \(p'_*\) is cover local]\label{mk-lem:essimage_pushforward_cover_local}
  \lean{AlgebraicGeometry.essImage_pushforward_of_openCover,
        AlgebraicGeometry.restrictFunctor_map_app}
  \leanok
  \dcref{custom}

  \uses{mk-lem:isIso_iff_isIso_restrict, mk-lem:restrictionMap_isIso_of_essImage,
        mk-lem:restrictFunctorIsoPullback_mathlib}
  Let \(p' : V' \hookrightarrow X'\) be an open immersion with open range \(R = p'(V')\), let
  \(M\) be an \(\mathcal{O}_{X'}\)-module, and let \(\{W_j\}\) be an open cover of \(X'\).  If
  for every \(j\) the restriction \(M|_{W_j}\) lies in the essential image of the pushforward
  along the open inclusion \(W_j \cap R \hookrightarrow W_j\), then \(M\) lies in the essential
  image of \(p'_*\).
\end{lemma}
\begin{proof}
  \leanok\uses{mk-lem:isIso_iff_isIso_restrict, mk-lem:restrictionMap_isIso_of_essImage,
        mk-lem:restrictFunctorIsoPullback_mathlib}
  Since \(p'_*\) is fully faithful (\ref{mk-lem:openimm_pullback_counit_isIso}), \(M\) lies in its
  essential image iff the unit \(M \to p'_*\,p'^*\,M\) of \(p'^* \dashv p'_*\) is an isomorphism.
  Under the identification of \(p'^*\) with restriction
  (\ref{mk-lem:restrictFunctorIsoPullback_mathlib}) the unit has, over an open \(U \subseteq X'\),
  the plain presheaf restriction component \(M(U) \to M(U \cap R)\).  Invertibility of this
  morphism of sheaves is local on \(X'\) (\ref{mk-lem:isIso_iff_isIso_restrict}): it suffices that
  each restriction to \(W_j\) be invertible, and — because restriction along an open immersion
  has definitional components — the restricted morphism has, over an open \(O \subseteq W_j\),
  again the plain component \(M(w_j(O)) \to M(p'(p'^{-1}(w_j(O))))\).  By the lattice identity
  \(w_j(O \cap w_j^{-1}(R)) = w_j(O) \cap R = p'(p'^{-1}(w_j(O)))\) this component is, up to the
  induced identification of the target sections, the restriction map
  \(M|_{W_j}(O) \to M|_{W_j}(O \cap w_j^{-1}(R))\), which is an isomorphism by
  \ref{mk-lem:restrictionMap_isIso_of_essImage} applied to the member datum.  Hence the unit is an
  isomorphism and \(M\) lies in the essential image of \(p'_*\).
\end{proof}

\begin{lemma}[Local essential-image comparison over an affine cover member]\label{mk-lem:openimm_bareBC_app_isIso_affine}
  \lean{AlgebraicGeometry.openImmersion_pushPull_essImage_member,
        AlgebraicGeometry.restrict_pullback_pushforward_essImage,
        AlgebraicGeometry.pullback_pushforward_affineOpen_essImage,
        AlgebraicGeometry.pushforwardRestrictOpensIso,
        AlgebraicGeometry.opensInter_image_eq,
        AlgebraicGeometry.opensInter_appIso_compat,
        AlgebraicGeometry.Scheme.Modules.pullbackIsoPushforwardInv}
  \leanok
  \dcref{custom}

  \uses{mk-lem:affine_pushforward_pullback_baseChange,
        mk-lem:pullback_preserves_quasicoherent, mk-lem:pullback_spec_tilde_iso,
        mk-lem:openimm_refining_affine_cover, mk-lem:restrictFunctorIsoPullback_mathlib}
  Keep the data of \ref{mk-lem:openimm_beckchevalley}. Assume that \(V\) is an affine
  open of \(X\), that \(X\) is separated, and that \(\mathcal F\) is quasi-coherent.
  Write \(M := (g')^*(p_*\,p^*\mathcal{F})\).  Let \(W_j = \operatorname{Spec} B \subseteq X'\),
  \(U_j = \operatorname{Spec} A \subseteq X\) be as in \ref{mk-lem:openimm_refining_affine_cover},
  with \(g'|_{W_j} = \operatorname{Spec}\varphi\), \(\varphi : A \to B\). Since \(V\) is affine and
  \(X\) is separated, the intersection \(V \cap U_j\) is affine, say \(V \cap U_j =
  \operatorname{Spec} A'\), and the open immersion \(V \cap U_j \hookrightarrow U_j\) is
  \(\operatorname{Spec}\) of a ring map \(A \to A'\). Then the restriction \(M|_{W_j}\) lies in
  the essential image of the pushforward along the open inclusion
  \(W_j \cap (g')^{-1}(V) \hookrightarrow W_j\).
\end{lemma}
\begin{proof}
  \leanok\uses{mk-lem:affine_pushforward_pullback_baseChange,
        mk-lem:pullback_preserves_quasicoherent, mk-lem:pullback_spec_tilde_iso,
        mk-lem:openimm_refining_affine_cover, mk-lem:restrictFunctorIsoPullback_mathlib}
  By the pullback pseudofunctor along \(w_j \cdot g' = g'|_{W_j} \cdot (U_j \hookrightarrow X)\),
  the restriction \(M|_{W_j}\) is the pullback along \(g'|_{W_j} = \operatorname{Spec}\varphi\) of
  the restriction \((p_*\,p^*\mathcal{F})|_{U_j}\); the open-restriction identity
  \((p_*\,p^*\mathcal{F})|_{U_j} = (V \cap U_j \hookrightarrow U_j)_*\,(p^*\mathcal{F}|_{V \cap U_j})\)
  (both sides are site-level pushforwards along the same opens functor) presents that restriction
  as \(f_*\widetilde{M'}\) for \(f = \operatorname{Spec}(A \to A')\) and the
  \(A'\)-module \(M' = \Gamma(V \cap U_j,\ p^*\mathcal{F})\). This \(M'\) exists because
  \(p^*\mathcal{F}\) is quasi-coherent (the restriction \(p^*\) of the quasi-coherent
  \(\mathcal{F}\), \ref{mk-lem:pullback_preserves_quasicoherent}) and a quasi-coherent module over an
  affine is the \(\widetilde{(-)}\) of its sections (\ref{mk-lem:pullback_spec_tilde_iso}).
  Applying \(\operatorname{Spec}\) to the ring pushout
  \[
    \begin{array}{ccc}
      A & \xrightarrow{\ \varphi\ } & B\\
      \downarrow & & \downarrow\\
      A' & \longrightarrow & B \otimes_A A'
    \end{array}
  \]
  (cut out by the affine map \(\varphi : A \to B\) presenting \(g'|_{W_j}\) and the affine open
  immersion \(A \to A'\) presenting \(V \cap U_j \hookrightarrow U_j\)) gives a fully-affine
  cartesian square, and the affine termwise base-change isomorphism
  \[
    g^*\bigl(f_*\widetilde{M'}\bigr) \;\cong\; f'_*\bigl((g')^*\widetilde{M'}\bigr)
  \]
  of \ref{mk-lem:affine_pushforward_pullback_baseChange} for this pushout square (with
  \(g = \operatorname{Spec}\varphi\)) exhibits \(M|_{W_j}\) as the pushforward of
  \((g')^*\widetilde{M'}\) along \(f' = \operatorname{Spec}(B \to B \otimes_A A')\), the base
  change of the open immersion \(V \cap U_j \hookrightarrow U_j\); its range is
  \(g'|_{W_j}^{-1}(V \cap U_j) = W_j \cap (g')^{-1}(V)\).  Transporting the pushforward along
  this range identification, \(M|_{W_j}\) lies in the essential image of the pushforward along
  \(W_j \cap (g')^{-1}(V) \hookrightarrow W_j\), as claimed.
\end{proof}

\begin{lemma}[Stage 2, assembly: \(\operatorname{bareBC}\) is an isomorphism]\label{mk-lem:openimm_bareBC_isIso}
  \lean{AlgebraicGeometry.openImmersion_bareBC_app_eq,
        AlgebraicGeometry.openImmersion_bareBC_app_isIso_of_unit,
        AlgebraicGeometry.openImmersion_unit_isIso_of_essImage,
        AlgebraicGeometry.openImmersion_pushPull_essImage,
        AlgebraicGeometry.openImmersion_pushPull_unit_isIso}
  \leanok
  \dcref{custom}

  \uses{mk-lem:openimm_pullback_counit_isIso, mk-lem:essimage_pushforward_cover_local,
        mk-lem:openimm_refining_affine_cover, mk-lem:openimm_bareBC_app_isIso_affine,
        mk-lem:openimm_bareBC, mk-lem:openimm_bc_telescope,
        mk-lem:isOpenImmersion_of_isPullback_left}
  For the cartesian square \(p' \cdot g' = g'_V \cdot p\) of \ref{mk-lem:openimm_beckchevalley} with
  \(p\) an open immersion, \(V\) affine (\([\mathrm{IsAffine}\,V]\)), \(X\) separated, and
  \(\mathcal{F}\) a \emph{quasi-coherent} \(\mathcal{O}_X\)-module
  (\(\mathtt{hF} : \mathcal{F}.\mathrm{IsQuasicoherent}\)), the bare base-change map
  \(\operatorname{bareBC}\) (\ref{mk-lem:openimm_bareBC}), evaluated at \(G = p^*\mathcal{F}\), is an
  isomorphism. (These hypotheses are essential: the proof applies the affine member node
  \ref{mk-lem:openimm_bareBC_app_isIso_affine}, which needs each \(V \cap U_j\) affine ---
  hence \(V\) affine and \(X\) separated --- and needs \(p^*\mathcal{F}\) quasi-coherent ---
  hence \(\mathtt{hF}\). The arbitrary-\(\mathcal{F}\) form is false.)
\end{lemma}
\begin{proof}
  \leanok\uses{mk-lem:openimm_pullback_counit_isIso, mk-lem:essimage_pushforward_cover_local,
        mk-lem:openimm_refining_affine_cover, mk-lem:openimm_bareBC_app_isIso_affine,
        mk-lem:openimm_bareBC, mk-lem:openimm_bc_telescope,
        mk-lem:isOpenImmersion_of_isPullback_left}
  Write \(M := (g')^*(p_*\,G)\).  \emph{Mate factorization.}  By the componentwise mate formula,
  \(\operatorname{bareBC}\) evaluated at \(G\) is the composite of the unit
  \(M \to p'_*\,p'^*\,M\) of \(p'^* \dashv p'_*\), followed by \(p'_*\) applied to the pullback
  telescope components (\ref{mk-lem:openimm_bc_telescope}, isomorphisms), followed by
  \(p'_*\,(g'_V)^*\) applied to the counit \(p^*(p_*\,G) \to G\) --- an isomorphism because
  \(p\) is an open immersion (\ref{mk-lem:openimm_pullback_counit_isIso}).  Hence
  \(\operatorname{bareBC}\) is an isomorphism as soon as the leading unit is.

  \emph{The unit is an isomorphism.}  Since \(p'\) is an open immersion
  (\ref{mk-lem:isOpenImmersion_of_isPullback_left}), \(p'_*\) is fully faithful
  (\ref{mk-lem:openimm_pullback_counit_isIso}), so the unit at \(M\) is an isomorphism iff \(M\)
  lies in the essential image of \(p'_*\).  That membership is cover-local
  (\ref{mk-lem:essimage_pushforward_cover_local}); check it on the affine open cover \(\{W_j\}\) of
  \(X'\) refining \((g')^{-1}(\text{affine opens of } X)\) supplied by
  \ref{mk-lem:openimm_refining_affine_cover}: over each member the restriction \(M|_{W_j}\) is a
  pushforward from \(W_j \cap (g')^{-1}(V)\) by the member node
  \ref{mk-lem:openimm_bareBC_app_isIso_affine}.  Hence \(M\) is in the essential image, the unit is
  an isomorphism, and \(\operatorname{bareBC}\) is an isomorphism.
\end{proof}

\begin{lemma}
[Open immersions are stable under base change]
  \label{mk-lem:isOpenImmersion_stableUnderBaseChange_mathlib}
  \lean{AlgebraicGeometry.isOpenImmersion_stableUnderBaseChange}
  \mathlibok
  \dcref{custom}

  The class of open immersions of schemes is stable under base change: in any pullback square,
  the base change of an open immersion is again an open immersion.
\end{lemma}

\begin{lemma}
[Left base change of an open immersion is an open immersion]
  \label{mk-lem:isOpenImmersion_of_isPullback_left}
  \lean{AlgebraicGeometry.isOpenImmersion_of_isPullback_left}
  \dcref{custom}

  \uses{mk-lem:isOpenImmersion_stableUnderBaseChange_mathlib}
  Given a cartesian square with maps \(g_V, p', p, g'\) satisfying \(p' \cdot g_V = g'
  \cdot p\), if \(p\) is an open immersion then so is its base change \(p'\).

  This supplies the open-immersion property of \(p'\) needed for the target-side restriction in
  Stage~2 of \ref{mk-lem:openimm_beckchevalley}.
\end{lemma}
\begin{proof}
  \uses{mk-lem:isOpenImmersion_stableUnderBaseChange_mathlib}
  Open immersions are stable under base change
  (\ref{mk-lem:isOpenImmersion_stableUnderBaseChange_mathlib}); applying that stability to the given
  cartesian square turns the open immersion \(p\) into the open immersion \(p'\).
\end{proof}

\begin{lemma}[Bare Beck--Chevalley mate of the restricted square]\label{mk-lem:openimm_bareBC}
  \lean{AlgebraicGeometry.openImmersion_bareBC}
  \leanok
  \dcref{custom}

  \uses{mk-lem:openimm_bc_telescope}
  For the cartesian square \(p' \cdot g' = g'_V \cdot p\) of \ref{mk-lem:openimm_beckchevalley} there
  is a canonical natural transformation --- the \emph{bare base-change map} ---
  \[
    \operatorname{bareBC} \colon\quad
    \operatorname{pushforward} p \circ (g')^*
      \;\Longrightarrow\;
    (g'_V)^* \circ \operatorname{pushforward} p',
  \]
  namely the Beck--Chevalley \emph{mate}, under the two pullback--pushforward adjunctions
  \(p^* \dashv p_*\) and \(p'^* \dashv p'_*\), of the canonical pullback \(2\)-isomorphism
  \((g')^* \circ p'^* \cong p^* \circ (g'_V)^*\) attached to the square. It exists for \emph{any}
  morphisms forming the cartesian square --- no flatness, no open-immersion hypothesis --- and its
  being an isomorphism is the genuine Beck--Chevalley content discharged in Stage~2 of
  \ref{mk-lem:openimm_beckchevalley}.
\end{lemma}
\begin{proof}
  Assemble the pullback \(2\)-isomorphism
  \((g')^* p'^* \cong (p' \cdot g')^* = (g'_V \cdot p)^* \cong (g'_V)^* p^*\) by combining
  the pseudofunctor composition isomorphism for \(p' \cdot g'\), the reindexing along the square
  equation \(p' \cdot g' = g'_V \cdot p\), and the inverse composition isomorphism for
  \(g'_V \cdot p\) (the content of \ref{mk-lem:openimm_bc_telescope}). The mate construction for the adjunctions \(p^* \dashv p_*\) and
  \(p'^* \dashv p'_*\) turns this \(2\)-isomorphism into the stated natural transformation. The
  construction is formal and uses no hypothesis on the maps.
\end{proof}

\begin{lemma}[Pullback telescope across the restricted square]\label{mk-lem:openimm_bc_telescope}
  \lean{AlgebraicGeometry.openImmersion_bc_telescope}
  \leanok
  \dcref{custom}

  In the cartesian square of \ref{mk-lem:openimm_beckchevalley}, the iterated pullback of any
  \(\mathcal{O}_X\)-module \(\mathcal{F}\) telescopes:
  \[
    p'^*\bigl((g')^*\mathcal{F}\bigr)
      \;\cong\; (p' \cdot g')^*\mathcal{F}
      \;=\; (g'_V \cdot p)^*\mathcal{F}
      \;\cong\; (g'_V)^*\bigl(p^*\mathcal{F}\bigr),
  \]
  using only the pseudofunctor comparison \((q \cdot r)^* \cong r^* \circ q^*\) and reindexing of
  \((-)^*\) along the square equation \(p' \cdot g' = g'_V \cdot p\). No flatness and no
  affineness hypotheses are needed.
\end{lemma}
\begin{proof}
  Compose the pseudofunctor composition isomorphism for \(p' \cdot g'\), the reindexing
  isomorphism along \(p' \cdot g' = g'_V \cdot p\), and the inverse composition isomorphism for
  \(g'_V \cdot p\), then evaluate the resulting natural isomorphism at \(\mathcal{F}\).
\end{proof}

\begin{lemma}[Per-intersection-open X-level Beck--Chevalley for the twisted nerve]\label{mk-lem:twisted_cech_nerve_per_sigma}
  \lean{AlgebraicGeometry.twisted_cech_nerve_per_sigma}
  \leanok
  \dcref{custom}

  \uses{mk-lem:coverinteropen_basechange_eq, mk-lem:openimm_beckchevalley, mk-lem:pushPullObj_iso_tilde}
  In the setting of \ref{mk-lem:twisted_cech_nerve_iso}, fix an index family \(\sigma\) with
  intersection open \(U_\sigma = \operatorname{coverInterOpen}\mathfrak{U}\,\sigma\),
  \(j_\sigma : U_\sigma \hookrightarrow X\), and write \(U'_\sigma\),
  \(j'_\sigma : U'_\sigma \hookrightarrow X'\) for the corresponding data of the base-changed
  cover \(\mathfrak{U}'\). Then there is the per-intersection-open \(X\)-level Beck--Chevalley
  isomorphism
  \[
    (g')^*\bigl((j_\sigma)_* (j_\sigma)^*\mathcal{F}\bigr)
      \;\cong\;
    (j'_\sigma)_* (j'_\sigma)^*\bigl((g')^*\mathcal{F}\bigr),
  \]
  over the base-changed intersection open \(U'_\sigma\); this is the single-open factor of the
  twisted-nerve identification. The index family \(\sigma\) is required to be \emph{finite}, so
  that \ref{mk-lem:coverinteropen_basechange_eq} applies to identify
  \((g')^{-1}(U_\sigma)\) with \(\operatorname{coverInterOpen}\mathfrak{U}'\,\sigma\) (the
  preimage commutes with the cover intersection only for finite index families); in the
  {\v C}ech application \(\sigma\) ranges over the finite index sets of the nerve.
\end{lemma}
\begin{proof}
  \leanok\uses{mk-lem:coverinteropen_basechange_eq, mk-lem:openimm_beckchevalley, mk-lem:pushPullObj_iso_tilde,
        mk-lem:cech_interopen_isaffine, mk-lem:restricted_cartesian_affine_pushout}
  The restriction \(j_\sigma : U_\sigma \hookrightarrow X\) is an open immersion, and its
  pullback along \(g'\) is the open immersion \(j'_\sigma : U'_\sigma \hookrightarrow X'\) onto
  the preimage \((g')^{-1}(U_\sigma)\), which by \ref{mk-lem:coverinteropen_basechange_eq} is
  exactly the base-changed intersection open
  \(U'_\sigma = \operatorname{coverInterOpen}\mathfrak{U}'\,\sigma\). Restricting the cartesian
  square \(g' : X' \to X\) over \(j_\sigma\) is therefore the restricted cartesian square of
  \ref{mk-lem:openimm_beckchevalley}, whose open-immersion Beck--Chevalley isomorphism is the
  asserted comparison. When the affine \(\widetilde{(-)}\)-model is needed downstream, both
  sides are identified with their affine forms by \ref{mk-lem:pushPullObj_iso_tilde} at
  altitude~1.
\end{proof}

\begin{lemma}[The base-changed nerve is the nerve of the base-changed data]
  \label{mk-lem:twisted_cech_nerve_iso}
  \lean{AlgebraicGeometry.twisted_cech_nerve_iso}
  \dcref{custom}

  \uses{mk-def:cech_nerve, mk-def:cover_cech_nerve, mk-lem:pushPullObj_iso_tilde}
  Let \(f : X \to S\) be \emph{separated} and fix an affine open-immersion cover
  \(\mathfrak{U} = \{U_i = \operatorname{Spec} A_i\}\) of \(X\) (so every fibre power is
  affine, \ref{mk-lem:isAffineOpen_biInf_mathlib}). Applying \((g')^*\) to the {\v C}ech
  nerve of \((\mathfrak{U}, \mathcal{F})\) yields the {\v C}ech nerve of the base-changed
  data \((\mathfrak{U}', (g')^*\mathcal{F})\), where \(\mathfrak{U}'\) is the base change of
  \(\mathfrak{U}\) along \(g'\): a natural isomorphism
  \((g')^* \circ \operatorname{CechNerve}(\mathfrak{U}, -) \cong
  \operatorname{CechNerve}(\mathfrak{U}', (g')^*(-))\).
\end{lemma}

\begin{proof}
  \uses{mk-def:cech_nerve, mk-def:cover_cech_nerve, mk-lem:cech_degree_affine_baseChange,
        mk-lem:affine_pushforward_pullback_baseChange, mk-lem:pushPullObj_iso_tilde,
        mk-lem:isAffineOpen_biInf_mathlib, mk-lem:pushforward_spec_tilde_iso,
        mk-lem:pullback_spec_tilde_iso, mk-lem:cancelBaseChange_mathlib,
        mk-lem:pushPull_sigma_iso, mk-lem:pushpullobj_coverinter_basechange,
        mk-lem:twisted_cech_nerve_per_sigma, mk-lem:cech_interopen_isaffine,
        mk-lem:coverinteropen_basechange_eq}
  \emph{Geometric backbone (first half).} The cover {\v C}ech nerve
  \(\operatorname{coverCechNerve}\) of \(\mathfrak{U}\) base-changes to that of
  \(\mathfrak{U}'\): the fibre powers \(U_{i_0} \times_X \dots \times_X U_{i_p}\) of the
  cover commute with the base change \(g'\) (pullback preserves fibre products), so the
  preimages \((g')^{-1}(U_{i_0\dots i_p})\) are exactly the corresponding intersections
  \(U'_{i_0\dots i_p}\) of \(\mathfrak{U}'\). This identifies the underlying diagrams of
  open subschemes. It does \emph{not} by itself produce the nerve isomorphism: the
  {\v C}ech nerve is valued not in open subschemes but in the push--pull functor
  \(p_* p^*\) (for \(p\) the projection of the relevant fibre power), so the comparison
  carries genuine sheaf-theoretic content beyond the identification of supports.

  \emph{Genuine content (second half, Beck--Chevalley).} The degree-\(n\) component
  is the \(X\)-level square along \(g'\): the asserted identification is the sheaf-level
  Beck--Chevalley isomorphism \((g')^*(\operatorname{pushPullObj}\mathcal{F}\, Y_n) \cong
  \operatorname{pushPullObj}((g')^*\mathcal{F})\, Y'_n\) (that is, \(g'^*(p_* p^* F) \cong
  p'_* p'^*(g'^* F)\)), where \(Y_n\), \(Y'_n\) are the degree-\(n\) fibre powers of the
  cover {\v C}ech nerves of \(\mathfrak{U}\) and \(\mathfrak{U}'\).
  \begin{enumerate}
    \item \emph{Product decomposition (LHS).} The degree-\(n\) fibre power is the coproduct
      \(Y_n = \coprod_\sigma U_\sigma\) over index tuples \(\sigma : \mathrm{Fin}(n+1) \to I\),
      so by \ref{mk-lem:pushPull_sigma_iso} (which requires the cover to be finite) the push--pull
      object decomposes as the finite product \(\operatorname{pushPullObj}\mathcal{F}\, Y_n \cong
      \prod_\sigma \operatorname{pushPullObj}\mathcal{F}\,(\operatorname{Over.mk} j_\sigma)\),
      and \((g')^*\) preserves this finite product. The left-hand side thus reduces to
      \(\prod_\sigma (g')^*\bigl((j_\sigma)_*(j_\sigma)^*\mathcal{F}\bigr)\).
    \item \emph{Per-\(\sigma\) component.} Each factor is the per-intersection-open \(X\)-level
      Beck--Chevalley iso \((g')^*\bigl((j_\sigma)_*(j_\sigma)^*\mathcal{F}\bigr) \cong
      (j'_\sigma)_*(j'_\sigma)^*\bigl((g')^*\mathcal{F}\bigr)\) of
      \ref{mk-lem:twisted_cech_nerve_per_sigma} (the \(X\)-level analogue of
      \ref{mk-lem:pushpullobj_coverinter_basechange} for the restricted cartesian square over
      \(U_\sigma\)). Here each intersection \(U_\sigma = U_{i_0} \cap \dots \cap U_{i_p}\) is
      affine by separatedness (\ref{mk-lem:cech_interopen_isaffine}); the abstract data is
      identified with the affine \(\widetilde{(-)}\)-form by \ref{mk-lem:pushPullObj_iso_tilde}
      at \emph{altitude 1} (the restriction level \(p_n^*\mathcal{F} \cong \widetilde{N_n}\)
      over the affine \(Y_n\)), pushed forward along \(g'\); and on that model the comparison
      IS the affine termwise base change \ref{mk-lem:affine_pushforward_pullback_baseChange}
      (via \ref{mk-lem:cech_degree_affine_baseChange}), assembled from the concrete tilde
      dictionaries \ref{mk-lem:pushforward_spec_tilde_iso} and
      \ref{mk-lem:pullback_spec_tilde_iso} with the commutative-algebra cancellation
      \(\operatorname{cancelBaseChange}\) (\ref{mk-lem:cancelBaseChange_mathlib}), \emph{not}
      the canonical adjoint mate.
    \item \emph{RHS reassembly (the residual cover-base-change obligation).} To recognise the
      product \(\prod_\sigma (j'_\sigma)_*(j'_\sigma)^*((g')^*\mathcal{F})\) as the push--pull
      object \(\operatorname{pushPullObj}((g')^*\mathcal{F})\, Y'_n\) of the base-changed cover
      one applies \ref{mk-lem:pushPull_sigma_iso} \emph{in reverse} to \(\mathfrak{U}'\). This
      requires the base-changed cover \(\mathfrak{U}'\) to be finite with affine members,
      together with the cover-base-change identification
      \(\operatorname{coverInterOpen}\mathfrak{U}'\,\sigma = (g')^{-1}
      (\operatorname{coverInterOpen}\mathfrak{U}\,\sigma)\) of
      \ref{mk-lem:coverinteropen_basechange_eq} (the preimages of the intersections
      of \(\mathfrak{U}\) are the intersections of \(\mathfrak{U}'\)). This geometric matching
      of the base-changed cover with the preimage cover is the Beck--Chevalley content of the
      argument.
    \item \emph{Naturality.} The naturality condition is verified by the index-omission
      restriction: each coface and codegeneracy is induced by an index omission, hence by
      an inclusion of finite affine intersections, and the termwise comparisons are natural
      in those inclusions.
  \end{enumerate}
\end{proof}

\begin{definition}[Degreewise pullback of the {\v C}ech complex as an alternating-coface complex]\label{mk-def:pullback_cechComplex_alternatingIso}
  \lean{AlgebraicGeometry.pullback_cechComplex_alternatingIso}
  \leanok
  \dcref{custom}

  \uses{mk-lem:pullback_commutes_alternatingCofaceMapComplex,
        mk-def:relative_cech_complex_of_nerve}
  Specialising \ref{mk-lem:pullback_commutes_alternatingCofaceMapComplex} to \(F = g^*\)
  and to the push--pull nerve underlying \(\check{\mathcal{C}}^\bullet(\mathfrak{U},
  \mathcal{F})\): applying \(g^*\) degreewise to the relative {\v C}ech complex is the
  alternating-coface complex of \(g^* \circ (\operatorname{pushforward} f) \circ
  \operatorname{CechNerve}(\mathfrak{U}, \mathcal{F})\). This is step~(1) of the
  base-change construction below.
\end{definition}

\begin{definition}[Reduction of the {\v C}ech base change to a cosimplicial iso]\label{mk-def:cech_complex_base_change_iso_of_cosimplicialIso}
  \lean{AlgebraicGeometry.cechComplex_baseChange_iso_of_cosimplicialIso}
  \leanok
  \dcref{02KG,custom}

  \uses{mk-def:pullback_cechComplex_alternatingIso,
        mk-lem:pullback_commutes_alternatingCofaceMapComplex,
        mk-lem:cech_pushforward_baseChange_natIso, mk-lem:twisted_cech_nerve_iso,
        mk-def:cech_complex_base_change_cosimplicialIso}
  Given a cosimplicial natural isomorphism
  \[
    e : g^* \circ (\operatorname{pushforward} f) \circ
        \operatorname{CechNerve}(\mathfrak{U}, \mathcal{F})
      \;\cong\;
      (\operatorname{pushforward} f') \circ
        \operatorname{CechNerve}(\mathfrak{U}', (g')^*\mathcal{F}),
  \]
  the {\v C}ech base-change complex isomorphism follows mechanically: it is the composite
  \[
    \operatorname{pullback\_cechComplex\_alternatingIso}
    \;\circ\;
    \operatorname{alternatingCofaceMapComplex}.\!\operatorname{mapIso}(e),
  \]
  i.e.\ the degreewise pullback identification
  (\ref{mk-def:pullback_cechComplex_alternatingIso}, step~(1)) followed by
  \(\operatorname{Functor.mapIso}\) of \(e\) through the alternating-coface complex
  functor. No per-degree {\v C}ech differential square is checked by hand:
  functoriality of \(\operatorname{alternatingCofaceMapComplex}\) supplies the
  differential compatibility automatically.

  The hypothesis \(e\) is precisely the residual obligation. It is constructed in
  \ref{mk-def:cech_complex_base_change_cosimplicialIso} as the whiskered composite of the
  Beck--Chevalley natural iso \ref{mk-lem:cech_pushforward_baseChange_natIso} with the
  twisted-nerve identification \ref{mk-lem:twisted_cech_nerve_iso}, both pushed through the
  {\v C}ech nerve. Supplying \(e\) is the sole remaining open obligation (Stacks Tag 02KG);
  \ref{mk-def:cech_complex_base_change_iso} below obtains the complex
  iso by feeding \(e\) into the present reduction.
\end{definition}

\begin{definition}[The {\v C}ech base-change cosimplicial isomorphism \(e\)]
  \label{mk-def:cech_complex_base_change_cosimplicialIso}
  \lean{AlgebraicGeometry.cechComplex_baseChange_cosimplicialIso_flat}
  \dcref{custom}

  \uses{mk-lem:cech_pushforward_baseChange_natIso, mk-lem:twisted_cech_nerve_iso}
  The cosimplicial natural isomorphism \(e\) required by
  \ref{mk-def:cech_complex_base_change_iso_of_cosimplicialIso} is assembled from the two
  base-change leaves. It is the composite
  \[
    e \;=\;
      \bigl(\ref{mk-lem:cech_pushforward_baseChange_natIso}\bigr)
      \mathbin{;}
      (\operatorname{pushforward} f')\!.\operatorname{mapIso}
        \bigl(\ref{mk-lem:twisted_cech_nerve_iso}\bigr),
  \]
  i.e.\ first the Beck--Chevalley natural iso
  \(g^* \circ (\operatorname{pushforward} f) \cong (\operatorname{pushforward} f')
  \circ (g')^*\) whiskered through the {\v C}ech nerve, followed by the image under the
  functor \(\operatorname{pushforward} f'\) (via \(\operatorname{Functor.mapIso}\)) of the
  twisted-nerve identification
  \((g')^* \circ \operatorname{CechNerve}(\mathfrak{U}, -) \cong
  \operatorname{CechNerve}(\mathfrak{U}', (g')^*(-))\). Both factors are isomorphisms of
  cosimplicial objects, so \(e\) is one as well; it is precisely the hypothesis consumed
  by \ref{mk-def:cech_complex_base_change_iso_of_cosimplicialIso}.
\end{definition}

\begin{definition}[Tensorial base change of the {\v C}ech complex \textnormal{(Stacks 02KG, OPEN)}]
  \label{mk-def:cech_complex_base_change_iso}
  \lean{AlgebraicGeometry.cechComplex_baseChange_iso}
  \dcref{02KG,custom}

  \uses{mk-def:cech_complex, mk-def:relative_cech_complex_of_nerve,
        mk-def:cech_nerve_cosimplicial}
  \textit{Source: Stacks Project, Cohomology of Schemes, Tag 02KG
  (\texttt{lemma-affine-base-change}); the {\v C}ech-complex relation in the proof of
  Tag 02KH (\texttt{lemma-flat-base-change-cohomology}).}
  Applying \(g^*\) degreewise to the relative {\v C}ech complex
  \(\check{\mathcal{C}}^\bullet(\mathfrak{U}, \mathcal{F})\) yields a {\v C}ech complex
  naturally isomorphic to the relative {\v C}ech complex
  \(\check{\mathcal{C}}^\bullet(\mathfrak{U}', (g')^*\mathcal{F})\) of the base-changed
  data, with \(\mathfrak{U}'\) the base change of \(\mathfrak{U}\) along \(g'\). Only the
  \emph{existence} of such a natural isomorphism is asserted --- the downstream
  assembly Lemma~\ref{mk-lem:cech_flat_base_change} concludes \(\mathrm{Nonempty}\) of an
  isomorphism --- so no identification with any canonical base-change map is required.
\end{definition}
\begin{proof}
  \uses{mk-def:cech_complex_base_change_iso_of_cosimplicialIso,
        mk-def:cech_complex_base_change_cosimplicialIso,
        mk-lem:pullback_commutes_alternatingCofaceMapComplex,
        mk-lem:cech_pushforward_baseChange_natIso, mk-lem:twisted_cech_nerve_iso,
        mk-lem:cech_degree_affine_baseChange,
        mk-lem:affine_pushforward_pullback_baseChange, mk-lem:pushforward_spec_tilde_iso,
        mk-lem:pullback_spec_tilde_iso, mk-lem:cancelBaseChange_mathlib}
  The construction factors through
  \ref{mk-def:cech_complex_base_change_iso_of_cosimplicialIso}: it suffices to supply the
  cosimplicial iso \(e\) and feed it into that reduction. We bypass the canonical
  base-change map entirely. The sole genuine input is the \emph{affine termwise} base
  change \ref{mk-lem:affine_pushforward_pullback_baseChange}, which is assembled from the
  concrete tilde identifications \ref{mk-lem:pushforward_spec_tilde_iso} and
  \ref{mk-lem:pullback_spec_tilde_iso} together with the commutative-algebra cancellation
  \(\operatorname{cancelBaseChange}\) (\ref{mk-lem:cancelBaseChange_mathlib}) --- not the
  canonical adjoint mate \(\operatorname{pushforwardBaseChangeMap}\). Unwinding the
  reduction, the construction is the composite,
  pushed through the functor \(\operatorname{alternatingCofaceMapComplex}\):
  \begin{enumerate}
    \item By Lemma~\ref{mk-lem:pullback_commutes_alternatingCofaceMapComplex}, applying
      \(g^*\) degreewise commutes with \(\operatorname{alternatingCofaceMapComplex}\),
      so \(g^*\) of the {\v C}ech complex is the alternating-coface complex of
      \(g^* \circ (\operatorname{pushforward} f) \circ
      \operatorname{CechNerve}(\mathfrak{U}, \mathcal{F})\).
    \item The Beck--Chevalley natural iso
      (Lemma~\ref{mk-lem:cech_pushforward_baseChange_natIso}) rewrites the cosimplicial
      object as \((\operatorname{pushforward} f') \circ (g')^* \circ
      \operatorname{CechNerve}(\mathfrak{U}, \mathcal{F})\); since this is a natural iso
      of cosimplicial objects, feeding it through the \emph{functor}
      \(\operatorname{alternatingCofaceMapComplex}\) via \(\operatorname{Functor.mapIso}\)
      makes compatibility with the {\v C}ech differentials automatic --- no per-degree
      naturality square need be checked by hand.
    \item Lemma~\ref{mk-lem:twisted_cech_nerve_iso} identifies the \((g')^*\)-twisted nerve
      with \(\operatorname{CechNerve}(\mathfrak{U}', (g')^*\mathcal{F})\), so the
      alternating-coface complex of the rewritten object is exactly
      \(\check{\mathcal{C}}^\bullet(\mathfrak{U}', (g')^*\mathcal{F})\).
  \end{enumerate}
  Composing (1)--(3) yields the required isomorphism of complexes; the homology
  isomorphism that Lemma~\ref{mk-lem:cech_flat_base_change} consumes is then
  \(\operatorname{homologyMapIso}\) of this complex iso, exactly as in that lemma's
  assembly. Following Stacks Tag 02KH, since \(g\) is flat the relation persists on
  taking cohomology; flatness is \emph{not} needed for the present complex iso, the
  affine base change being flatness-free.
\end{proof}
